% (User input window)

\subsubsection{尺度核的质量塔分解与粗反卷积：有限阶局部算子、尾核矩与带宽良定性}\label{subsubsec:xi-scale-kernel-mass-tower-deconvolution}
以下条目把 logistic 尺度核在 Fourier 侧的 \(\sinh\) 乘积结构实现为一个可执行的质量塔卷积模型，
并由此导出一族有限阶局部反卷积算子及其在带宽受控制度下的稳定性证书；
尾核的可计算矩与四阶累积量给出最小非高斯偏差的可观测刻度，从而把 \(q=4\) 的首触发机制提升为结构必然性。

\begin{conclusion}[尺度核的无穷质量塔分解：logistic 核等价于独立 Laplace 级联系统]\label{con:xi-terminal-scale-kernel-mass-tower-laplace}
令尺度单峰核为
\[
\phi(s):=\frac{1}{4\cosh^{2}(s/2)}=\frac{e^{s}}{(1+e^{s})^{2}}\qquad(s\in\RR).
\]
对每个整数 \(n\ge 1\) 定义 Laplace 核
\[
g_n(s):=\frac{n}{2}e^{-n|s|},
\qquad
\widehat g_n(\xi)=\frac{1}{1+\xi^{2}/n^{2}}.
\]
则存在一族独立随机变量 \((X_n)_{n\ge 1}\)（\(X_n\) 的密度为 \(g_n\)），使得级数
\[
Y:=\sum_{n\ge 1}X_n
\]
在 \(L^{2}\) 与几乎处处意义下收敛，并且其和的密度严格等于 \(\phi\)。
等价地，令 \(\phi_{\le N}:=g_1*\cdots*g_N\)，则 \(\phi_{\le N}\to\phi\)（\(L^2\) 且弱收敛），从而在分布意义下可记为
\[
\phi=g_1*g_2*\cdots .
\]
该分解把 \(\phi\) 解释为一条“整数质量谱”\(\{n\}_{n\ge 1}\) 的逐层卷积封口：每一层 \(g_n\) 是一维局部算子 \((1-\frac{1}{n^{2}}\partial_s^{2})^{-1}\) 的 Green 核，整条质量塔的乘积谱恰闭合为 \(\phi\)。
\end{conclusion}
\begin{proof}[证明要点]
由命题 \ref{prop:xi-bulk-curvature-tomography-identifiability} 的 Fourier 闭式 \(\widehat\phi(\xi)=\pi\xi/\sinh(\pi\xi)\) 及 \(\sinh\) 的 Hadamard 乘积
\(
\sinh(\pi\xi)=\pi\xi\prod_{n\ge 1}(1+\xi^{2}/n^{2})
\)
得
\(
\widehat\phi(\xi)=\prod_{n\ge 1}\widehat g_n(\xi)
\)。
取独立随机变量 \(X_n\)（密度为 \(g_n\)），则其特征函数为 \(\widehat g_n\)。
令部分和 \(Y_N=\sum_{n=1}^{N}X_n\)，则 \(Y_N\) 的特征函数为 \(\prod_{n=1}^{N}\widehat g_n\)，故 \(Y_N\) 的密度为 \(\phi_{\le N}\)。
又因 \(\sum_{n\ge 1}\Var(X_n)=2\sum_{n\ge 1}n^{-2}<\infty\)，得 \((Y_N)\) 在 \(L^2\) 中 Cauchy，且由 Kolmogorov 判据可得几乎处处收敛。
极限变量 \(Y\) 的特征函数为 \(\widehat\phi\)，从而其密度为 \(\phi\)。证毕。
\end{proof}

\begin{conclusion}[部分反卷积即分辨率旋钮：有限阶局部算子把无限反卷积降为尾核卷积]\label{con:xi-terminal-scale-kernel-partial-deconvolution-resolution}
对 \(N\ge 1\) 定义有限阶局部算子（以 Fourier 乘子理解）
\[
\mathcal D_N:=\prod_{n=1}^{N}\Bigl(1-\frac{1}{n^2}\partial_s^2\Bigr),
\qquad
\widehat{(\mathcal D_N f)}(\xi)=M_N(\xi)\,\widehat f(\xi),
\qquad
M_N(\xi):=\prod_{n=1}^{N}(1+\xi^{2}/n^{2}).
\]
并令尾核 \(\phi_{>N}\) 为结论 \ref{con:xi-terminal-scale-kernel-mass-tower-laplace} 中尾部和
\(
Y_{>N}:=\sum_{n>N}X_n
\)
的密度（亦即 \(\widehat{\phi_{>N}}(\xi)=\prod_{n>N}(1+\xi^{2}/n^{2})^{-1}\)）。
则对任意有限符号测度 \(\nu\)，若 \(p=\phi*\nu\)，则成立严格恒等式
\[
\boxed{\;
\mathcal D_N p=\phi_{>N}*\nu.\;}
\]
因而 \(\mathcal D_N\) 在制度上实现“去除前 \(N\) 层质量壳”的粗反卷积：当 \(N\to\infty\) 时 \(\phi_{>N}\Rightarrow\delta_0\)，从而 \(\mathcal D_N p\Rightarrow \nu\) 于分布意义下成立。
该结论给出一条可执行的连续族：反卷积由 \(N\) 明确调控的分辨率流替代了“要么精确要么崩溃”的二元机制。
\end{conclusion}
\begin{proof}[证明要点]
由结论 \ref{con:xi-terminal-scale-kernel-mass-tower-laplace}，\(\widehat\phi(\xi)=M_N(\xi)^{-1}\widehat{\phi_{>N}}(\xi)\)。
对 \(p=\phi*\nu\) 取 Fourier 变换得 \(\widehat p=\widehat\phi\cdot\widehat\nu\)，从而
\[
\widehat{(\mathcal D_N p)}(\xi)=M_N(\xi)\widehat p(\xi)=\widehat{\phi_{>N}}(\xi)\widehat\nu(\xi)=\widehat{(\phi_{>N}*\nu)}(\xi),
\]
得到恒等式。
又因 \(\Var(Y_{>N})\to 0\)，故 \(\phi_{>N}\Rightarrow\delta_0\)，从而 \(\phi_{>N}*\nu\Rightarrow\nu\)。证毕。
\end{proof}

\begin{conclusion}[尾核的矩完全可计算：分辨率尺度 \texorpdfstring{$\sigma_N\sim \sqrt{2/N}$}{sigma\_N ~ sqrt(2/N)} 与 Wasserstein 控制]\label{con:xi-terminal-scale-kernel-tail-moments-w2}
沿用结论 \ref{con:xi-terminal-scale-kernel-mass-tower-laplace} 的随机变量 \(X_n\)，令 \(Y_{>N}:=\sum_{n>N}X_n\)，其密度为 \(\phi_{>N}\)。
则有可加矩证书
\[
\Var(Y_{>N})=\sum_{n>N}\Var(X_n)=2\sum_{n>N}\frac{1}{n^2}=: \sigma_N^2,
\qquad
\kappa_4(Y_{>N})=\sum_{n>N}\kappa_4(X_n)=12\sum_{n>N}\frac{1}{n^4}.
\]
并且对任意具有有限二阶矩的概率测度 \(\nu\) 有可证书化的 \(W_2\) 上界
\[
\boxed{\;
W_2\bigl(\phi_{>N}*\nu,\nu\bigr)\ \le\ \sigma_N
\ =\ \sqrt{2\sum_{n>N}n^{-2}}
\ \sim\ \sqrt{\frac{2}{N}}\quad(N\to\infty).\;}
\]
因此若允许分辨率 \(\sigma>0\)，取 \(N\asymp 2/\sigma^{2}\) 即可保证 \(W_2\)-误差 \(\le \sigma\)；
在该制度下，分辨率成本从指数型降为幂律型。
\end{conclusion}
\begin{proof}[证明要点]
对 \(X_n\sim g_n\) 直接计算得到 \(\Var(X_n)=2/n^2\)、\(\kappa_4(X_n)=12/n^4\)，累积量对独立和可加，从而得矩证书。
对 \(W_2\) 上界，取耦合：令 \(U\sim \nu\)，令 \(V\sim \phi_{>N}\) 且独立，则 \((U+V,U)\) 是 \((\phi_{>N}*\nu,\nu)\) 的一个耦合，故
\(
W_2^2(\phi_{>N}*\nu,\nu)\le \EE[|V|^2]=\Var(V)=\sigma_N^2
\)。
渐近 \(\sum_{n>N}n^{-2}\sim 1/N\) 由积分判别法得到。证毕。
\end{proof}

\begin{conclusion}[尾核的高斯化与四阶首破缺：非高斯性由 \texorpdfstring{$q=4$}{q=4} 量级唯一主导]\label{con:xi-terminal-scale-kernel-tail-gaussianization-kurtosis}
令 \(\sigma_N^2=\Var(Y_{>N})\)，并定义标准化尾变量 \(Z_N:=Y_{>N}/\sigma_N\)。
则有弱收敛 \(Z_N\Rightarrow \mathcal N(0,1)\)，且其标准化超额峰度（excess kurtosis）具有闭式
\[
\boxed{\;
\gamma_2(N):=\frac{\kappa_4(Y_{>N})}{\sigma_N^4}
=3\,\frac{\sum_{n>N}n^{-4}}{\bigl(\sum_{n>N}n^{-2}\bigr)^2}
\sim \frac{1}{N}\quad(N\to\infty).\;}
\]
并且所有奇阶累积量严格为零。
因此尾核的首个非高斯偏差不可避免地发生在四阶：在尺度流的可观测校正中，最小破缺阶被结构性锁定为 \(4\)，并以 \(N^{-1}\) 的强度进入。
\end{conclusion}
\begin{proof}[证明要点]
由于 \(X_n\) 对称，故全部奇阶累积量为零，并在独立和下保持为零。
对 \(Z_N\) 的中心极限定理可由 Lindeberg 条件给出：\(\max_{n>N}\Var(X_n)/\sigma_N^2\to 0\) 且 \(\sum_{n>N}\Var(X_n)=\sigma_N^2\)。
峰度公式由结论 \ref{con:xi-terminal-scale-kernel-tail-moments-w2} 的累积量证书代入并化简得到；
渐近由 \(\sum_{n>N}n^{-4}\sim (3N^3)^{-1}\)、\(\sum_{n>N}n^{-2}\sim N^{-1}\) 给出。证毕。
\end{proof}

\begin{conclusion}[带宽约束下的条件数上界：粗反卷积在任何固定带宽内均一致良定]\label{con:xi-terminal-scale-kernel-bandwidth-conditioned}
设观测为 \(\widetilde p=p+\eta\)，并假设噪声在频带 \(|\xi|\le\Omega\) 内满足
\[
\sup_{|\xi|\le\Omega}|\widehat\eta(\xi)|\le\varepsilon.
\]
则对任意 \(N\ge 1\)，有限阶算子输出误差在同一频带内满足
\[
\sup_{|\xi|\le\Omega}\bigl|\widehat{(\mathcal D_N\eta)}(\xi)\bigr|
\le
\varepsilon\cdot \sup_{|\xi|\le\Omega}M_N(\xi)
\le
\varepsilon\cdot \exp\!\Bigl(\Omega^2\sum_{n=1}^{\infty}\frac{1}{n^2}\Bigr)
=\varepsilon\cdot \exp\!\Bigl(\frac{\pi^2}{6}\Omega^2\Bigr).
\]
因此在任何固定带宽 \(\Omega\) 内，\(\mathcal D_N\) 的噪声放大因子对 \(N\) 一致有界；
粗反卷积的病态性并非来自“阶数增大”，而来自“试图在无限带宽上实现 \(\delta_0\) 极限”。
该结论为尺度流的可计算相图提供严格基础：在带宽受控的制度下，分辨率可按 \(N\sim\sigma^{-2}\) 幂律提升而不触发指数型条件数爆炸。
\end{conclusion}
\begin{proof}[证明要点]
由结论 \ref{con:xi-terminal-scale-kernel-partial-deconvolution-resolution} 的 Fourier 乘子表示，
\(
\widehat{(\mathcal D_N\eta)}=M_N\widehat\eta
\)
于 \(|\xi|\le\Omega\) 上成立。
并且对 \(x\ge 0\) 有 \(\log(1+x)\le x\)，故
\[
\log M_N(\xi)=\sum_{n=1}^{N}\log\Bigl(1+\frac{\xi^2}{n^2}\Bigr)\le \xi^2\sum_{n=1}^{N}\frac{1}{n^2}
\le \Omega^2\sum_{n=1}^{\infty}\frac{1}{n^2}=\frac{\pi^2}{6}\Omega^2,
\]
从而得所示一致上界。证毕。
\end{proof}

\begin{conclusion}[有限阶谱乘子的 \texorpdfstring{$\Gamma$}{Gamma}-闭式与小频率展开：局部算子的解析完备性]\label{con:xi-terminal-scale-kernel-gamma-closed-form}
记有限阶谱乘子
\[
M_N(\xi):=\prod_{n=1}^{N}(1+\xi^2/n^2).
\]
则 \(M_N\) 具有严格的 \(\Gamma\)-比闭式
\[
\boxed{\;
M_N(\xi)=\frac{\Gamma(N+1+\mathrm{i}\xi)\Gamma(N+1-\mathrm{i}\xi)}{(N!)^2\,\Gamma(1+\mathrm{i}\xi)\Gamma(1-\mathrm{i}\xi)}.\;}
\]
从而在 \(|\xi|=o(N)\) 下有可审计展开
\[
\log M_N(\xi)=\xi^2\sum_{n=1}^{N}\frac{1}{n^2}-\frac{\xi^4}{2}\sum_{n=1}^{N}\frac{1}{n^4}+O\!\Bigl(\xi^6\sum_{n=1}^{N}\frac{1}{n^6}\Bigr),
\]
揭示出“二阶（\(\zeta(2)\)）决定主放大、四阶（\(\zeta(4)\)）决定首修正”的严格层级。
该闭式使 \(\mathcal D_N\) 的解析结构可被有限步复算，并允许在带宽与精度给定时对 \(N\) 作可证伪的最优截断。
\end{conclusion}
\begin{proof}[证明要点]
将 \(M_N(\xi)=\prod_{n=1}^{N}(n^2+\xi^2)/n^2=\prod_{n=1}^{N}(n+\mathrm{i}\xi)(n-\mathrm{i}\xi)/(N!)^2\)，再用有限乘积恒等式
\(
\prod_{n=1}^{N}(n+z)=\Gamma(N+1+z)/\Gamma(1+z)
\)
即得 \(\Gamma\)-比闭式。
对小频率展开，使用 \(\log(1+u)=u-\frac{u^2}{2}+O(u^3)\) 并对 \(u=\xi^2/n^2\) 求和即可。证毕。
\end{proof}

\begin{conclusion}[尺度流的两阶段可恢复性：幂律分辨率与指数精确性的制度分离]\label{con:xi-terminal-scale-kernel-two-regime-recoverability}
对任意 \(\sigma>0\)，取 \(N(\sigma)\) 使 \(\sigma_{N(\sigma)}\le\sigma\)（其中 \(\sigma_N\) 如结论 \ref{con:xi-terminal-scale-kernel-tail-moments-w2}）。
则由结论 \ref{con:xi-terminal-scale-kernel-partial-deconvolution-resolution} 有
\[
\mathcal D_{N(\sigma)}p=\phi_{>N(\sigma)}*\nu,
\]
给出一条分辨率为 \(\sigma\) 的稳定重建；其代价为局部算子阶数 \(2N(\sigma)\asymp\sigma^{-2}\)，且在固定带宽下条件数一致可控（结论 \ref{con:xi-terminal-scale-kernel-bandwidth-conditioned}）。
\par
相反，若要求 \(\sigma\downarrow 0\) 的精确极限（即 \(\phi_{>N}\Rightarrow\delta_0\) 的无带宽极限），则必然进入指数不适定区：尾核逼近 \(\delta_0\) 要求同时控制任意高频，制度上必须允许带宽 \(\Omega\to\infty\)。
更精确地，由 \(\sinh(\pi\xi)=\pi\xi\prod_{n\ge 1}(1+\xi^2/n^2)\) 得极限乘子
\[
M_\infty(\xi):=\lim_{N\to\infty}M_N(\xi)=\frac{\sinh(\pi\xi)}{\pi\xi},
\qquad
M_\infty(\xi)\sim \frac{1}{2\pi|\xi|}\,e^{\pi|\xi|}\quad(|\xi|\to\infty).
\]
因此若把“精确”理解为需要在 \(|\xi|\lesssim \Omega\) 上逼近 \(\delta_0\)，则噪声放大随 \(\Omega\) 指数增长；
若进一步把制度分辨率 \(\sigma\) 视为要求可见带宽 \(\Omega\asymp \sigma^{-1}\)，则放大因子量级为 \(\exp(\pi/\sigma)\)（仅差多项式因子），从而精度位数随 \(1/\sigma\) 线性增长。
因此“幂律可控”与“指数不可控”并非矛盾，而是由尺度选择（是否接受 \(\sigma\)-分辨率）所分离的两种制度。
\end{conclusion}
\begin{proof}[证明要点]
稳定重建部分由结论 \ref{con:xi-terminal-scale-kernel-partial-deconvolution-resolution} 的恒等式与结论 \ref{con:xi-terminal-scale-kernel-tail-moments-w2} 的 \(\sigma_N\) 控制给出；
在带宽受控制度下的良定性由结论 \ref{con:xi-terminal-scale-kernel-bandwidth-conditioned} 给出。
精确极限要求 \(\Omega\) 不受界，从而噪声放大必随 \(\Omega\) 发散，给出指数不适定区。证毕。
\end{proof}

\begin{conclusion}[\texorpdfstring{$(q=4)$}{(q=4)} 的正则化地位：四阶破缺决定可见相变的最小矩阶]\label{con:xi-terminal-scale-kernel-q4-minimal-visible-order}
在结论 \ref{con:xi-terminal-scale-kernel-mass-tower-laplace}--\ref{con:xi-terminal-scale-kernel-tail-gaussianization-kurtosis} 的质量塔模型中，尾核的全部奇阶累积量严格为零，而首个非高斯项由四阶累积量 \(\kappa_4\) 给出并以
\(\gamma_2(N)\sim N^{-1}\)
控制。
故任何以“二阶近似（Gaussian）+ 带宽受控的粗反卷积”为基底的可审计推断体系，其最小系统性偏差必出现在四阶；
换言之，若仅允许有限矩阶作为统计可见接口，则最小能够检测“非高斯边界/相变临界”之矩阶必为 \(4\)。
该事实为 \(q=4\) 的首越线现象提供完全结构性的解释：它并非偶然数值，而是尺度核尾部的对称性与质量塔叠加所共同强制的最小破缺阶。
\end{conclusion}
\begin{proof}[证明要点]
由结论 \ref{con:xi-terminal-scale-kernel-tail-gaussianization-kurtosis}，标准化尾变量的奇阶累积量为零，且四阶累积量给出首个非高斯偏差并以 \(N^{-1}\) 量级衰减。
在以二阶高斯近似为基底的制度中，任何有限阶矩统计量的系统偏差由最低非零累积量阶决定，故最小可见矩阶为 \(4\)。证毕。
\end{proof}

\subsubsection{行列式体积熵的树/环分层：Chow--Liu 投影、条件互信息与四阶显影}\label{subsubsec:xi-det-entropy-tree-loop-hierarchy}
在以 Gram/相关矩阵承载“体积势”的口径下，\(-\log\det\) 自然给出一个多体体积熵势；
本段把该势的可审计下界压缩为 Chow--Liu 树投影，并将剩余缺口解释为由环结构承载的不可逆体积亏损。
在一维排序与单调距离核制度中，最优树选择刚性退化为相邻链，从而把“树承载的体积熵”外部化为相邻间隔账本；
而在弱相干区，环熵的首个系统性贡献由条件互信息控制并从四阶开始显影，与本文反复出现的 \(q=4\) 首触发结构同型。

\begin{conclusion}[\texorpdfstring{$(\det)$}{(det)}-熵势的 Chow--Liu 树下界：行列式体积由最大生成树控制]\label{con:xi-det-entropy-chow-liu-tree-lower-bound}
令 \(R\in\RR^{\kappa\times\kappa}\) 为相关矩阵（\(R\succ0\)、\(\mathrm{diag}(R)=\mathbf 1\)），定义行列式体积熵势
\[
S_{\det}(R):=-\log\det R.
\]
对每条无向边 \((i,j)\) 定义二体互信息权重
\[
w_{ij}:=-\log(1-R_{ij}^2).
\]
则对任意生成树 \(T\) 都有
\[
S_{\det}(R)\ \ge\ \sum_{(i,j)\in T} w_{ij},
\qquad\text{从而}\qquad
S_{\det}(R)\ \ge\ \max_{T}\sum_{(i,j)\in T} w_{ij}.
\]
并且对给定 \(T\)，等号成立当且仅当 \(R^{-1}\) 的非对角非零支撑包含于 \(T\)（等价于 \(\mathcal N(0,R)\) 在图 \(T\) 上为树形高斯马尔可夫场）。
\end{conclusion}
\begin{proof}[证明要点]
取 \(X\sim\mathcal N(0,R)\)。
在固定树 \(T\) 下，Chow--Liu 树投影 \(\mathcal N_T\)（见结论 \ref{con:xi-det-entropy-loop-entropy-kl-gap} 的设定）满足
\(
2D_{\mathrm{KL}}(\mathcal N(0,R)\,\|\,\mathcal N_T)
=S_{\det}(R)-\sum_{(i,j)\in T}w_{ij}\ge 0
\)，
从而得下界。
等号等价于 \(D_{\mathrm{KL}}=0\)，即 \(\mathcal N(0,R)=\mathcal N_T\)；
对高斯族这等价于精度矩阵仅在 \(T\) 的边上允许出现非零项。证毕。
\end{proof}

\begin{conclusion}[环熵即树近似的 \texorpdfstring{$D_{\mathrm{KL}}$}{DKL} 缺口：体积亏损的精确树分解]\label{con:xi-det-entropy-loop-entropy-kl-gap}
沿用结论 \ref{con:xi-det-entropy-chow-liu-tree-lower-bound} 的设定。
对给定生成树 \(T\) 定义树环熵
\[
\mathcal L_T(R):=S_{\det}(R)-\sum_{(i,j)\in T}w_{ij}
=-\log\det R+\sum_{(i,j)\in T}\log(1-R_{ij}^2).
\]
则 \(\mathcal L_T(R)\ge 0\) 恒成立。
若 \(X\sim\mathcal N(0,R)\)，并令 \(\mathcal N_T\) 为以 \(T\) 为图的 Chow--Liu 树形高斯投影（其一维边缘与树边二元边缘与 \(\mathcal N(0,R)\) 一致），则成立严格恒等式
\[
\mathcal L_T(R)=2D_{\mathrm{KL}}\!\bigl(\mathcal N(0,R)\,\|\,\mathcal N_T\bigr).
\]
因此
\(
\arg\min_T\mathcal L_T(R)=\arg\max_T\sum_{(i,j)\in T}w_{ij}
\)
给出最优树近似。
\end{conclusion}
\begin{proof}[证明要点]
令 \(R_T\) 为 \(\mathcal N_T\) 的相关矩阵。
树形投影 \(\mathcal N_T\) 的协方差由树边相关唯一决定，并满足
\[
\det R_T=\prod_{(i,j)\in T}(1-R_{ij}^2)
\]
（树形高斯的逐边条件方差乘积公式）。
另一方面，零均值高斯族的 KL 闭式给出
\[
2D_{\mathrm{KL}}(\mathcal N(0,R)\,\|\,\mathcal N(0,R_T))
=\tr(R_T^{-1}R)-\kappa+\log\frac{\det R_T}{\det R}.
\]
由于 \(R_T^{-1}\) 的非对角支撑包含于 \(T\)，且 \(R\) 与 \(R_T\) 在对角与树边上取值一致，故
\(
\tr(R_T^{-1}R)=\tr(R_T^{-1}R_T)=\kappa
\)，
从而
\(
2D_{\mathrm{KL}}=\log\frac{\det R_T}{\det R}
\)。
代入行列式因子化即得所示恒等式。证毕。
\end{proof}

\begin{conclusion}[单调距离核的树选择刚性：一维排序点的最优树必为相邻链]\label{con:xi-det-entropy-monotone-kernel-chain-mst}
设存在严格单调递减函数 \(A:(0,\infty)\to(0,1)\) 使
\[
R_{ij}=A(|s_i-s_j|),
\qquad
s_1<\cdots<s_\kappa.
\]
则边权 \(w_{ij}=-\log(1-A(|s_i-s_j|)^2)\) 亦严格随距离递减。
因此最大生成树唯一为路径图（相邻链）
\[
T_{\mathrm{chain}}:=\{(1,2),(2,3),\dots,(\kappa-1,\kappa)\},
\]
从而得到最强树下界
\[
S_{\det}(R)\ \ge\ \sum_{j=1}^{\kappa-1}-\log\!\bigl(1-A(\Delta_j)^2\bigr),
\qquad
\Delta_j:=s_{j+1}-s_j.
\]
\end{conclusion}
\begin{proof}[证明要点]
对任意割 \(\{1,\dots,j\}\,\cup\,\{j+1,\dots,\kappa\}\)，跨割的距离最小边唯一为 \((j,j+1)\)，且其权重最大。
由最大生成树的割性质（cut property），每条 \((j,j+1)\) 必属于最大生成树；合并所有 \(j\) 即得唯一性与链结构。证毕。
\end{proof}

\begin{conclusion}[三点局部环熵的闭式：链缺口等于条件互信息，且为部分相关的对数势]\label{con:xi-det-entropy-three-point-loop-entropy-partial-corr}
取任意 \(a<b<c\)，记 \(\rho_{ab}=R_{ab}\)、\(\rho_{bc}=R_{bc}\)、\(\rho_{ac}=R_{ac}\)，并令 \(R_{abc}\) 为对应三元相关子矩阵。
相对于链树 \((a,b)\)--\((b,c)\) 定义三点局部环熵
\[
\mathcal L^{(3)}_{a,b,c}
:=-\log\det R_{abc}+\log(1-\rho_{ab}^2)+\log(1-\rho_{bc}^2).
\]
则有恒等式
\[
\mathcal L^{(3)}_{a,b,c}
=-\log\bigl(1-\rho_{ac\mid b}^2\bigr),
\qquad
\rho_{ac\mid b}:=\frac{\rho_{ac}-\rho_{ab}\rho_{bc}}{\sqrt{(1-\rho_{ab}^2)(1-\rho_{bc}^2)}}.
\]
并且若 \(X\sim\mathcal N(0,R)\)，则
\[
\mathcal L^{(3)}_{a,b,c}=2I(X_a;X_c\mid X_b).
\]
\end{conclusion}
\begin{proof}[证明要点]
由三元相关矩阵行列式与部分相关的标准恒等式
\(
1-\rho_{ac\mid b}^2=\det R_{abc}/((1-\rho_{ab}^2)(1-\rho_{bc}^2))
\)
得到第一式。
对高斯族，条件互信息满足
\(
2I(X_a;X_c\mid X_b)=-\log(1-\rho_{ac\mid b}^2)
\)，
从而得第二式。证毕。
\end{proof}

\begin{conclusion}[四阶首触发律：远距衰减型弱相干区的环熵从四阶开始显影]\label{con:xi-det-entropy-q4-trigger-weak-coherence}
在结论 \ref{con:xi-det-entropy-three-point-loop-entropy-partial-corr} 的三点设定下，若
\(
|\rho_{ab}|,|\rho_{bc}|,|\rho_{ac}|\le\varepsilon\ll1
\)，
则有一致展开
\[
\rho_{ac\mid b}=(\rho_{ac}-\rho_{ab}\rho_{bc})+O(\varepsilon^3),
\qquad
\mathcal L^{(3)}_{a,b,c}=\rho_{ac\mid b}^2+O(\varepsilon^4).
\]
特别地，当满足远距衰减型关系 \(\rho_{ac}=\rho_{ab}\rho_{bc}+O(\varepsilon^3)\) 时，\(\rho_{ac\mid b}=O(\varepsilon^2)\) 且
\[
\mathcal L^{(3)}_{a,b,c}=O(\varepsilon^4),
\]
即局部环熵具有“最小四阶可见性”：二体相关 \(O(\varepsilon)\) 在该制度下不足以触发链缺口，缺口的首个系统性量级必为四阶。
\end{conclusion}
\begin{proof}[证明要点]
分母 \(\sqrt{(1-\rho_{ab}^2)(1-\rho_{bc}^2)}=1+O(\varepsilon^2)\)，故部分相关可由分子线性化得到。
又由 \(-\log(1-x)=x+O(x^2)\)（\(x\to0\)）并取 \(x=\rho_{ac\mid b}^2=O(\varepsilon^2)\)，得所示展开。证毕。
\end{proof}

\begin{conclusion}[指数尾尺度核下的显式四阶标度：\texorpdfstring{$\mathcal L^{(3)}$}{L(3)} 的可审计渐近]\label{con:xi-det-entropy-exp-tail-kernel-audit-asymptotics}
取特定单调距离核
\[
A(\Delta):=\frac{\Delta}{2\sinh(\Delta/2)}\qquad(\Delta\ge0),
\qquad
A(0):=1,
\]
并令 \(\Delta_1=s_b-s_a\)、\(\Delta_2=s_c-s_b\)。
当 \(\Delta_1,\Delta_2\to\infty\) 时有
\[
\rho_{ab}\sim \Delta_1e^{-\Delta_1/2},\quad
\rho_{bc}\sim \Delta_2e^{-\Delta_2/2},\quad
\rho_{ac}\sim (\Delta_1+\Delta_2)e^{-(\Delta_1+\Delta_2)/2}.
\]
从而
\[
\rho_{ac}-\rho_{ab}\rho_{bc}
=e^{-(\Delta_1+\Delta_2)/2}\Bigl((\Delta_1+\Delta_2)-\Delta_1\Delta_2+o(\Delta_1\Delta_2)\Bigr),
\]
并在弱相干区得到可审计的四阶主导项
\[
\mathcal L^{(3)}_{a,b,c}
=\bigl(\rho_{ac}-\rho_{ab}\rho_{bc}\bigr)^2\,(1+o(1))
\asymp (\Delta_1\Delta_2)^2\,e^{-(\Delta_1+\Delta_2)}.
\]
\end{conclusion}
\begin{proof}[证明要点]
由 \(A(\Delta)\sim \Delta e^{-\Delta/2}\)（\(\Delta\to\infty\)）得三条相关渐近。
再用结论 \ref{con:xi-det-entropy-q4-trigger-weak-coherence} 将 \(\mathcal L^{(3)}\) 线性化为 \(\rho_{ac\mid b}^2\)，并注意分母 \(1+o(1)\)，即可得到主导项与标度。证毕。
\end{proof}

\begin{conclusion}[链树环熵的链式分解：全局缺口等于远过去条件互信息之和]\label{con:xi-det-entropy-chain-loop-entropy-chain-rule}
令 \(X=(X_1,\dots,X_\kappa)\sim\mathcal N(0,R)\)，并取相邻链树
\(
T_{\mathrm{chain}}=\{(1,2),\dots,(\kappa-1,\kappa)\}
\)。
定义全局链环熵
\[
\mathcal L_{\mathrm{chain}}(R):=\mathcal L_{T_{\mathrm{chain}}}(R)
=-\log\det R+\sum_{j=1}^{\kappa-1}\log(1-R_{j,j+1}^2).
\]
则存在精确链式分解
\[
\mathcal L_{\mathrm{chain}}(R)
=2\sum_{j=3}^{\kappa} I\!\bigl(X_j;X_{1:j-2}\mid X_{j-1}\bigr),
\]
从而
\[
\mathcal L_{\mathrm{chain}}(R)\ \ge\ 2\sum_{j=3}^{\kappa} I\!\bigl(X_j;X_{j-2}\mid X_{j-1}\bigr)
=\sum_{j=3}^{\kappa}\mathcal L^{(3)}_{j-2,j-1,j}.
\]
\end{conclusion}
\begin{proof}[证明要点]
由结论 \ref{con:xi-det-entropy-loop-entropy-kl-gap}，\(\mathcal L_{\mathrm{chain}}/2\) 等于 \(\mathcal N(0,R)\) 相对于链形 Chow--Liu 投影的 KL 缺口；
对链形图该缺口等于一列条件互信息的和（链式法则）。
下界由条件互信息的逐项非负分解
\(
I(X_j;X_{1:j-2}\mid X_{j-1})=\sum_{i=1}^{j-2}I(X_j;X_i\mid X_{i+1:j-1})\ge I(X_j;X_{j-2}\mid X_{j-1})
\)
给出。证毕。
\end{proof}

\begin{conclusion}[弱相干链的 \(q=4\) 体积律：环熵主导项为二阶差分偏离的平方和]\label{con:xi-det-entropy-weak-coherence-chain-q4-law}
在结论 \ref{con:xi-det-entropy-monotone-kernel-chain-mst} 的相邻排序设定下，记 \(\rho_{j,j+1}=R_{j,j+1}\) 并定义二阶差分偏离
\[
\eta_j:=R_{j-2,j}-R_{j-2,j-1}R_{j-1,j}\qquad(3\le j\le\kappa).
\]
若整体处于弱相干区（例如对指数尾核取全部相邻间隔 \(\Delta_j\) 足够大，使所有 \(|R_{ij}|\) 一致足够小），则有一致近似
\[
\mathcal L_{\mathrm{chain}}(R)
=\sum_{j=3}^{\kappa}\frac{\eta_j^2}{(1-\rho_{j-2,j-1}^2)(1-\rho_{j-1,j}^2)}
\ +\ o\!\Bigl(\sum_{j=3}^{\kappa}\eta_j^2\Bigr).
\]
因此链环熵的首个系统性贡献由“二阶差分偏离”的平方和控制，并在远距衰减制度下呈现四阶主导量级。
\end{conclusion}
\begin{proof}[证明要点]
由结论 \ref{con:xi-det-entropy-chain-loop-entropy-chain-rule}，\(\mathcal L_{\mathrm{chain}}\) 的首项由一列条件互信息给出。
在弱相干区，远过去对 \(X_j\) 的条件影响高阶可忽略，主导项由 \(I(X_j;X_{j-2}\mid X_{j-1})\) 给出；
并由结论 \ref{con:xi-det-entropy-three-point-loop-entropy-partial-corr} 得
\(
\mathcal L^{(3)}_{j-2,j-1,j}=-\log(1-\rho_{j-2,j\mid j-1}^2)
\)
与
\(
\rho_{j-2,j\mid j-1}^2=\eta_j^2/((1-\rho_{j-2,j-1}^2)(1-\rho_{j-1,j}^2))
\)。
再用 \(-\log(1-x)=x+o(x)\) 线性化即可。证毕。
\end{proof}

\begin{conclusion}[马尔可夫体积与真实体积的比值：链环熵是精确的体积亏损对数]\label{con:xi-det-entropy-markov-volume-ratio}
定义链马尔可夫相关矩阵 \(R^{\mathrm{MC}}\) 为满足
\[
R^{\mathrm{MC}}_{jj}=1,\qquad
R^{\mathrm{MC}}_{j,j+1}=R_{j,j+1},
\qquad
(R^{\mathrm{MC}})^{-1}\ \text{为三对角},
\]
的唯一相关矩阵（即与 \(T_{\mathrm{chain}}\) 相容的树形高斯投影）。
则成立行列式闭式
\[
\det R^{\mathrm{MC}}=\prod_{j=1}^{\kappa-1}\bigl(1-R_{j,j+1}^2\bigr),
\qquad
\mathcal L_{\mathrm{chain}}(R)= -\log\frac{\det R}{\det R^{\mathrm{MC}}}.
\]
因此链环熵是“真实体积”相对于“马尔可夫体积”的精确对数亏损，提供一种无需外部拟合的体积比较标尺。
\end{conclusion}
\begin{proof}[证明要点]
链形树投影对应一列逐步回归表示
\(
X_{j+1}=R_{j,j+1}X_j+\sqrt{1-R_{j,j+1}^2}\,\varepsilon_{j+1}
\)
（\(\varepsilon_{j+1}\) 独立标准高斯），从而协方差行列式为条件方差乘积，给出 \(\det R^{\mathrm{MC}}\) 的因子化；
再由结论 \ref{con:xi-det-entropy-loop-entropy-kl-gap} 对链树取 \(T=T_{\mathrm{chain}}\) 即得体积比恒等式。证毕。
\end{proof}

\begin{conclusion}[任意树的基环条件互信息证书：非树边的代价受全局环熵支配]\label{con:xi-det-entropy-tree-basecycle-cmi-certificate}
对任意生成树 \(T\) 与任意非树边 \(e=(u,v)\notin T\)，令 \(P_T(u,v)\) 为树上连接 \(u\) 与 \(v\) 的唯一路径顶点集合，并记分离集
\(
S_e:=P_T(u,v)\setminus\{u,v\}
\)。
定义基环条件互信息势
\[
\mathcal I_e:=2I\!\bigl(X_u;X_v\mid X_{S_e}\bigr)
=-\log\bigl(1-\rho_{uv\mid S_e}^2\bigr),
\]
其中 \(\rho_{uv\mid S_e}\) 为在高斯族下对 \(S_e\) 条件化的部分相关系数。
则对高斯族恒有逐边支配
\[
\mathcal I_e\ \le\ \mathcal L_T(R)\qquad\text{对所有 }e\notin T.
\]
因此任意非树边的条件相关强度都被同一全局环熵所控制；在弱相干制度（非树边相关由路径相关乘积主导）下，每个 \(\mathcal I_e\) 的首项由四阶量级给出，从而 \(\mathcal L_T\) 的最小可见阶在该制度下被结构性锁定为 \(4\)。
\end{conclusion}
\begin{proof}[证明要点]
树形投影 \(\mathcal N_T\) 在图 \(T\) 上满足分离性：对任意 \(e=(u,v)\notin T\)，必有 \(X_u\perp X_v\mid X_{S_e}\) 于 \(\mathcal N_T\) 下成立。
设 \(\mathcal Q_e\) 为所有满足该条件独立性的分布族，则
\(
\mathcal N_T\in\mathcal Q_e
\)。
另一方面，\(D_{\mathrm{KL}}(\mathcal N(0,R)\,\|\,\mathcal Q_e)\) 的最小值恰为 \(I(X_u;X_v\mid X_{S_e})\)；
因此
\(
D_{\mathrm{KL}}(\mathcal N(0,R)\,\|\,\mathcal N_T)\ge I(X_u;X_v\mid X_{S_e})
\)。
乘以 \(2\) 并用结论 \ref{con:xi-det-entropy-loop-entropy-kl-gap} 的恒等式即得所示支配。证毕。
\end{proof}

\begin{conclusion}[等距弱相干相的每步环熵标度：\texorpdfstring{$\Delta\to\infty$}{Delta->inf} 时 \texorpdfstring{$\mathcal L_{\mathrm{chain}}/\kappa$}{L_chain/kappa} 的显式衰减]\label{con:xi-det-entropy-equispaced-loop-entropy-scaling}
取等距深度 \(s_j=s_0+(j-1)\Delta\) 并令 \(\Delta\to\infty\)。
记
\(
a:=A(\Delta)\sim \Delta e^{-\Delta/2}
\)、
\(
b:=A(2\Delta)\sim 2\Delta e^{-\Delta}
\)，
则每个相邻三点的二阶差分偏离满足
\[
\eta=b-a^2\sim-\Delta^2e^{-\Delta}.
\]
从而链环熵的每步强度满足
\[
\frac{1}{\kappa}\mathcal L_{\mathrm{chain}}(R)
\sim \eta^2
\asymp \Delta^{4}e^{-2\Delta},
\qquad (\Delta\to\infty,\ \kappa\ \text{大}).
\]
该式把弱相干下的非马尔可夫性外部化为完全显式的指数律：指数率为 \(2\Delta\)，并以多项式 \(\Delta^4\) 作为前因子修正。
\end{conclusion}
\begin{proof}[证明要点]
由结论 \ref{con:xi-det-entropy-exp-tail-kernel-audit-asymptotics} 的核渐近得 \(a\)、\(b\) 的指数尾标度，并计算 \(\eta=b-a^2\)。
再由结论 \ref{con:xi-det-entropy-weak-coherence-chain-q4-law}，弱相干区 \(\mathcal L_{\mathrm{chain}}\) 的主导项为 \(\sum \eta_j^2\)；等距情形各项同阶，除以 \(\kappa\) 得每步标度。证毕。
\end{proof}

\begin{conclusion}[体积—维度—熵的树/环分层：边界维度由树项给出，分形粗糙由环项给出]\label{con:xi-det-entropy-tree-loop-stratification}
对任意相关矩阵 \(R\) 定义
\[
S_{\mathrm{tree}}(R):=\max_{T}\sum_{(i,j)\in T}-\log(1-R_{ij}^2),
\qquad
S_{\mathrm{loop}}(R):=S_{\det}(R)-S_{\mathrm{tree}}(R)\ (\ge0).
\]
则 \(S_{\mathrm{tree}}\) 给出“可由树结构承载的最强体积熵”，而 \(S_{\mathrm{loop}}\) 给出“必须由环结构承载的额外体积熵”。
在单调距离核的一维排序情形，\(S_{\mathrm{tree}}\) 由结论 \ref{con:xi-det-entropy-monotone-kernel-chain-mst} 刚性退化为相邻间隔账本，并给出边界侧的主导体积熵贡献；
而 \(S_{\mathrm{loop}}\) 在弱相干区由结论 \ref{con:xi-det-entropy-q4-trigger-weak-coherence}--\ref{con:xi-det-entropy-weak-coherence-chain-q4-law} 从四阶开始显影，从而可被视为刻画“分形粗糙度/相变临界性”的最小可见体积指标。
\end{conclusion}
\begin{proof}[证明要点]
非负性由结论 \ref{con:xi-det-entropy-chow-liu-tree-lower-bound} 立即推出。
一维排序下的链刚性由结论 \ref{con:xi-det-entropy-monotone-kernel-chain-mst} 给出；
弱相干区的四阶显影由结论 \ref{con:xi-det-entropy-q4-trigger-weak-coherence} 与结论 \ref{con:xi-det-entropy-weak-coherence-chain-q4-law} 给出。
因此树项与环项分别对应两类结构性承载。证毕。
\end{proof}

\subsubsection{树/环分层后的跨模块推进：矩阵核、算子封口与时序账本}\label{subsubsec:xi-det-entropy-cross-module-propagation}

\paragraph{有限否证的可定位读出与跨口径不变量：Toeplitz 铅笔、负平方指数、端点预算与算术护栏}
本段给出一组与正文/附录可审计链条一致的条目化归约：有限级 Toeplitz 失效不仅提供二值否证器，而且可被提升为候选点生成、端点深度读出与指数不变量的闭合语义；并在扭曲图 \(\zeta\) 与 \(p\)-幂同余口径中给出可枚举的障碍分类与零成本审计护栏。

\begin{conclusion}[有限否证器的构造性读出：Toeplitz 失效可算法化为离线零点候选输出]\label{con:xi-finite-falsifier-offline-fiber-readout}
沿用视界 Carath\'eodory 系数 \((c_n)\) 与 Toeplitz 主子矩阵 \(\mathbf{T}_N=(c_{j-k})_{0\le j,k\le N}\) 的记号（定理 \ref{thm:app-horizon-toeplitz-lmi}）。
若存在某个 \(N\) 使 \(\mathbf{T}_N\nsucceq 0\)，则可由一阶移位 Toeplitz 矩阵
\[
\mathbf{T}_N^{(1)}:=(c_{j-k+1})_{0\le j,k\le N}
\]
形成 Toeplitz 铅笔 \((\mathbf{T}_N,\mathbf{T}_N^{(1)})\)，并由其广义谱在单位圆盘内读出有限共振候选点集
\[
\Lambda_N:=\mathrm{Spec}(\mathbf{T}_N,\mathbf{T}_N^{(1)})\cap\mathbb{D}
=\{\lambda^{(N)}_j\}_{j=1}^{\kappa_N}\subset\mathbb{D},
\qquad \kappa_N:=\nu_-(\mathbf{T}_N)\ge 1.
\]
对每个 \(\lambda^{(N)}_j\) 取 Cayley 逆映射 \(t^{(N)}_j:=\mathcal{C}^{-1}(\lambda^{(N)}_j)\in\mathbb{H}\)，并写作
\(
t^{(N)}_j=\gamma^{(N)}_j+\mathrm{i}\delta^{(N)}_j
\)
（\(\gamma^{(N)}_j\in\RR\)、\(\delta^{(N)}_j>0\)）。
则候选离线零点纤维可被显式生成为
\[
\boxed{\;
\rho^{(N)}_j:=\frac12+\mathrm{i}\gamma^{(N)}_j-\delta^{(N)}_j\qquad(1\le j\le \kappa_N).\;}
\]
在“有限型且边界无极点”且负惯性指标在高阶截断上稳定的情形下，该候选集合在 \(N\to\infty\) 极限中收敛到真实盘内极点集合（等价于离线零点纤维），从而将有限级 Toeplitz 否证从二值失败见证提升为可定位的反例几何读出。
\end{conclusion}
\begin{proof}[证明要点]
候选点集 \(\Lambda_N\) 的谱读出由定理 \ref{thm:xi-resonance-pencil-reconstruction} 给出；纤维构造与有限型收敛口径由定理 \ref{thm:xi-finite-witness-to-offline-fiber} 给出。
负惯性稳定性与有限型等价接口见定理 \ref{thm:xi-toeplitz-inertia-stabilizes-to-kappa} 与推论 \ref{cor:xi-inertia-stable-counts-disk-zeros}。
\end{proof}

\begin{conclusion}[Toeplitz 负平方指数的稳定律：负惯性最终等于 Blaschke 缺陷维数]\label{con:xi-toeplitz-negative-square-stability-law}
在有限缺陷设定中（\(F_{\zeta}\) 为 bounded type、边界无极点且其 Blaschke 因子为有限乘积），记
\(\kappa:=\deg(B_{F_{\zeta}})=\dim K_{B_{F_{\zeta}}}\)。
则圆盘核 \(K_{\mathscr{C}_{\zeta}}\) 的 Pontryagin 负平方数恰为 \(\kappa\)，并且 Toeplitz 主子矩阵族的负惯性指标满足
\[
\nu_-(\mathbf{T}_N)\le \kappa\qquad(\forall\,N),
\qquad
\nu_-(\mathbf{T}_N)=\kappa\qquad(N\ \text{充分大}).
\]
因此 RH 等价于 \(\kappa=0\)；任何偏离在有限层级上必表现为最终稳定的有限维负方向，从而把 Toeplitz 失效的复杂度压缩为一个有限整数不变量。
\end{conclusion}
\begin{proof}[证明要点]
这是一条“有限极点 \(\Longleftrightarrow\) 有限负平方”的标准对齐：见定理 \ref{thm:xi-toeplitz-negative-squares-equals-blaschke-defect} 与定理 \ref{thm:xi-toeplitz-inertia-stabilizes-to-kappa}（并与 \S\ref{subsubsec:xi-negative-square-index-pontryagin-realization} 的负平方指数口径一致）。
\end{proof}

\begin{conclusion}[视界证书阶数的端点可见性下界：最小 Toeplitz 失败阶服从普适尺度律]\label{con:xi-endpoint-visibility-lower-bound-toeplitz-order}
若存在离临界线零点 \(\rho=\beta+\mathrm{i}\gamma\)（\(\beta<\tfrac12\)、\(\delta=\tfrac12-\beta>0\)），记其圆盘像 \(w_\rho=\mathcal{C}(\gamma+\mathrm{i}\delta)\in\mathbb{D}\)，并定义端点边界层深度
\[
\boxed{\ h(\rho):=1-|w_\rho|^2=\frac{4\delta}{\gamma^2+(1+\delta)^2}\ }.
\]
则任何试图在有限阶 Toeplitz--PSD 层级上检出该缺陷的方案，其最小证书阶数 \(N_{\min}\) 必满足定量门槛
\[
N_{\min}\ \gtrsim\ \frac{1}{h(\rho)}\,\log\!\bigl(1+\delta\,h(\rho)^{-1}\bigr),
\qquad\text{从而特别地}\qquad
N_{\min}\ \gtrsim\ h(\rho)^{-1}.
\]
在 \(|\gamma|\to\infty\) 且 \(\delta\) 有界时，
\(
h(\rho)^{-1}\asymp \gamma^2/\delta
\)
给出普适尺度律：端点二次径向压缩将离切片缺陷强制推入端点边界层，并把其可见化成本锁定在 \(h(\rho)^{-1}\) 的量纲上，不能被任何有限深度证书规避。
\end{conclusion}
\begin{proof}[证明要点]
证书阶数下界见命题 \ref{con:xi-toeplitz-length-depth-duality} 与附录定理 \ref{thm:app-horizon-toeplitz-detection-threshold}；
深度闭式恒等式见定理 \ref{thm:xi-offcritical-quadratic-radial-compression}，并由其高高度渐近给出 \(h(\rho)^{-1}\asymp \gamma^2/\delta\)。
\end{proof}

\begin{proposition}[Toeplitz--PSD 失效的 Fibonacci 局部化]\label{prop:xi-toeplitz-psd-fibonacci-localization}
\leanverified{paper\_xi\_toeplitz\_psd\_fibonacci\_localization}
沿用 \(\mathbf{T}_N=(c_{j-k})_{0\le j,k\le N}\) 的 Toeplitz 主子矩阵记号（定理 \ref{thm:app-horizon-toeplitz-lmi}）。
若存在最小失效阶
\[
N_0:=\min\{N\ge 0:\ \mathbf{T}_N\nsucceq 0\}<\infty,
\]
并令 \(F_1=F_2=1\) 为 Fibonacci 数列，取
\[
k_0:=\min\{k\ge 1:\ F_k\ge N_0\},
\]
则 \(\mathbf{T}_{F_{k_0}}\nsucceq 0\)，并且有尺度界
\[
\boxed{\;F_{k_0}<\varphi\,N_0,\qquad \varphi=\frac{1+\sqrt5}{2}.\;}
\]
因此当 PSD 失效存在时，仅在 Fibonacci 子序列上检验即可在一个普适常数因子内定位到某个失效层级。
\end{proposition}
\begin{proof}
由于 \(\mathbf{T}_{N_0}\) 是 \(\mathbf{T}_M\) 的左上主子矩阵（\(M\ge N_0\)），若 \(\mathbf{T}_M\succeq 0\) 则必有 \(\mathbf{T}_{N_0}\succeq 0\)，与 \(N_0\) 的定义矛盾。
故对一切 \(M\ge N_0\) 都有 \(\mathbf{T}_M\nsucceq 0\)，取 \(M=F_{k_0}\) 得 \(\mathbf{T}_{F_{k_0}}\nsucceq 0\)。
\par
由 \(k_0\) 的最小性，\(F_{k_0-1}<N_0\) 且 \(F_{k_0}=F_{k_0-1}+F_{k_0-2}\)。
又对所有 \(k\ge 2\) 有经典不等式 \(F_k/F_{k-1}<\varphi\)，从而
\(
F_{k_0}<\varphi F_{k_0-1}<\varphi N_0
\)。
证毕。
\end{proof}

\begin{conclusion}[Adams 粗粒化的否证回推：折叠尺度与证书阶数存在乘积下界]\label{con:xi-adams-coarsening-backprop-budget-product}
设 \(\mathscr{C}(w)=c_0+2\sum_{n\ge 1}c_n w^n\) 为 Carath\'eodory 视界函数的矩序列口径（并约定 \(c_{-n}=\overline{c_n}\)）。
对整数 \(m\ge 1\) 定义推前粗粒化矩
\(
c^{(m)}_n:=c_{mn}
\)
与派生 Toeplitz 层级
\[
\mathbf{T}^{(m)}_N:=\bigl(c_{m(j-k)}\bigr)_{0\le j,k\le N}.
\]
则对任意 \(m,N\ge 1\) 有蕴含关系
\[
\boxed{\ \mathbf{T}_{mN}\succeq 0\ \Longrightarrow\ \mathbf{T}^{(m)}_{N}\succeq 0\ },
\qquad\text{其逆否命题为}\qquad
\mathbf{T}^{(m)}_{N}\nsucceq 0\ \Longrightarrow\ \mathbf{T}_{mN}\nsucceq 0.
\]
因此任何派生层级的否证都可严格回推到原层级的否证，从而“更深层级支配其全部粗粒化派生”。
在端点边界层机制下，若缺陷深度为 \(h(\rho)\)，则由检出门槛 \(mN\gtrsim h(\rho)^{-1}\) 得到一条可审计的时间--证书预算乘积下界：折叠尺度 \(m\) 与视界阶数 \(N\) 的混合策略必须支付 \(h(\rho)^{-1}\) 量纲的资源成本。
\end{conclusion}
\begin{proof}[证明要点]
派生层级 \(\mathbf{T}^{(m)}_N\succeq 0\) 的 pushforward 解释与 \(\mathbf{T}_{mN}\Rightarrow \mathbf{T}^{(m)}_{N}\) 的结构蕴含见 \eqref{eq:xi-adams-toeplitz-psd}；
乘积下界由命题 \ref{con:xi-toeplitz-length-depth-duality} 对 \(mN\) 阶原层级适用并结合 \(h(\rho)\) 的端点压缩闭式得到。
\end{proof}

\paragraph{二端点二项式 Toeplitz 见证的极点闭式与尺度交换}
在结论 \ref{con:xi-adams-coarsening-backprop-budget-product} 的稀疏 Toeplitz 层级与定理 \ref{thm:xi-endpoint-binomial-falsifier} 的端点二项式否证器之间，还存在一条更细的闭合链：负端点与正端点二项式模板分别给出主导极点在 \(-1\) 与 \(+1\) 端点处的闭式响应；二者的增长率联合起来可代数恢复主导极点的模长与实部；再把同一模板施加到 Adams 稀疏子序列 \(\{c_{mn}\}\) 上，则主导项的增长指数乘以 \(m\)，从而把 Toeplitz 检出深度压缩到 \((m\Lambda_{\mathrm{mom}})^{-1}\) 量级。共动图表进一步把这一尺度律中的高度依赖完全剥离，留下只由偏移 \(\delta\) 控制的一参数选择原则。

\begin{theorem}[负端点二项式 Toeplitz 见证对主导极点的闭式响应]\label{thm:xi-binomial-toeplitz-dominant-pole-response}
\leanverified{paper\_xi\_binomial\_toeplitz\_dominant\_pole\_response}
沿用定理 \ref{thm:xi-endpoint-binomial-falsifier} 的记号，并进一步假设最深极点 \(w_\ast\in\mathbb D\) 为单极，且存在 \(r_2>|w_\ast|\) 与常数 \(M>0\)，使得
\[
\mathscr C_{\zeta}(w)=\frac{R_\ast}{w-w_\ast}+H(w),
\qquad
R_\ast\neq 0,
\]
其中 \(H\) 在 \(|w|<r_2\) 内全纯，并满足系数分解
\[
c_n=-\frac{R_\ast}{2\,w_\ast^{\,n+1}}+h_n,
\qquad
|h_n|\le M r_2^{-n}
\qquad(n\ge 1).
\]
对每个 \(N\ge 1\)，定义负端点二项式向量
\[
b_-^{(N)}:=\bigl((-1)^k\binom Nk\bigr)_{k=0}^{N},
\qquad
Q_N^-:=\bigl(b_-^{(N)}\bigr)^{\mathsf T}\mathbf T_N\,b_-^{(N)}.
\]
则存在常数 \(C=C(M,r_2)>0\) 使得
\begin{equation}\label{eq:xi-binomial-toeplitz-negative-endpoint-response}
\boxed{\;
Q_N^-
=
\Re\!\Bigl(
-\frac{R_\ast}{2}\,
(-1)^N\,
(w_\ast-1)^{2N}\,
w_\ast^{-(N+1)}
\Bigr)
+E_N^-,
\qquad
|E_N^-|\le C\,4^N r_2^{-N}.
\;}
\end{equation}
\end{theorem}

\begin{proof}
由定理 \ref{thm:xi-endpoint-binomial-falsifier} 的组合恒等式，
\[
Q_N^-=
\sum_{n=-N}^{N}(-1)^n\binom{2N}{N+n}c_n.
\]
将 \(c_n\) 分解为主导极点项与解析余项，得到
\[
Q_N^-=
-\frac{R_\ast}{2}w_\ast^{-1}
\sum_{n=-N}^{N}(-1)^n\binom{2N}{N+n}w_\ast^{-n}
+
\sum_{n=-N}^{N}(-1)^n\binom{2N}{N+n}h_n.
\]
对主项使用二项式恒等式
\[
\sum_{n=-N}^{N}(-1)^n\binom{2N}{N+n}z^n
=
(-1)^N z^{-N}(1-z)^{2N},
\]
再代入 \(z=w_\ast^{-1}\)，即得
\[
-\frac{R_\ast}{2}w_\ast^{-1}
\sum_{n=-N}^{N}(-1)^n\binom{2N}{N+n}w_\ast^{-n}
=
-\frac{R_\ast}{2}
(-1)^N
(w_\ast-1)^{2N}
w_\ast^{-(N+1)}.
\]
由于 \(Q_N^-\) 为实数，取实部即得到 \eqref{eq:xi-binomial-toeplitz-negative-endpoint-response} 的主项。
另一方面，由 Cauchy 估计得 \(|h_n|\le M r_2^{-n}\)，并由 \(\sum_{n=-N}^{N}\binom{2N}{N+n}=4^N\) 可得
\[
|E_N^-|
\le
\sum_{n=-N}^{N}\binom{2N}{N+n}|h_n|
\le
C\,4^N r_2^{-N},
\]
其中 \(C\) 吸收有限个低阶项的统一常数。证毕。
\end{proof}

\begin{theorem}[双端点增长率的主导极点反演]\label{thm:xi-binomial-toeplitz-two-endpoint-pole-inversion}
在定理 \ref{thm:xi-binomial-toeplitz-dominant-pole-response} 的设定下，再定义正端点二项式向量
\[
b_+^{(N)}:=\bigl(\binom Nk\bigr)_{k=0}^{N},
\qquad
Q_N^+:=\bigl(b_+^{(N)}\bigr)^{\mathsf T}\mathbf T_N\,b_+^{(N)}.
\]
则存在常数 \(C'=C'(M,r_2)>0\) 使得
\begin{equation}\label{eq:xi-binomial-toeplitz-positive-endpoint-response}
\boxed{\;
Q_N^+
=
\Re\!\Bigl(
-\frac{R_\ast}{2}\,
(w_\ast+1)^{2N}\,
w_\ast^{-(N+1)}
\Bigr)
+E_N^+,
\qquad
|E_N^+|\le C'4^N r_2^{-N}.
\;}
\end{equation}
若进一步满足：
\begin{enumerate}
  \item 模长最小的极点仅由共轭对 \(\{w_\ast,\overline{w_\ast}\}\) 构成；
  \item 两端点响应均严格主导解析余项，即
  \[
  \frac{|w_\ast-1|^2}{|w_\ast|}>\frac{4}{r_2},
  \qquad
  \frac{|w_\ast+1|^2}{|w_\ast|}>\frac{4}{r_2},
  \]
\end{enumerate}
则极限上确界
\[
A:=\limsup_{N\to\infty}|Q_N^-|^{1/N},
\qquad
B:=\limsup_{N\to\infty}|Q_N^+|^{1/N}
\]
满足
\begin{equation}\label{eq:xi-binomial-toeplitz-two-endpoint-growth-rates}
\boxed{\;
A=\frac{|w_\ast-1|^2}{|w_\ast|},
\qquad
B=\frac{|w_\ast+1|^2}{|w_\ast|}.
\;}
\end{equation}
若写 \(w_\ast=x+\mathrm{i}y\)、\(r_\ast:=|w_\ast|\in(0,1)\)，则
\begin{equation}\label{eq:xi-binomial-toeplitz-radius-recovery}
\boxed{\;
2r^2-(A+B)r+2=0
\quad\text{在 }(0,1)\text{ 中的唯一根即为 }r_\ast,
\;}
\end{equation}
并且
\begin{equation}\label{eq:xi-binomial-toeplitz-realpart-recovery}
\boxed{\;
x=-\frac{(A-B)r_\ast}{4},
\qquad
|y|=\sqrt{r_\ast^2-x^2}.
\;}
\end{equation}
\end{theorem}

\begin{proof}
由
\[
\sum_{n=-N}^{N}\binom{2N}{N+n}z^n
=
z^{-N}(1+z)^{2N}
\]
与定理 \ref{thm:xi-binomial-toeplitz-dominant-pole-response} 的同一分解，立即得到
\eqref{eq:xi-binomial-toeplitz-positive-endpoint-response}。
在假设 (1)--(2) 下，两式主项的指数增长率分别为
\[
\frac{|w_\ast-1|^2}{|w_\ast|},
\qquad
\frac{|w_\ast+1|^2}{|w_\ast|},
\]
且解析余项的指数率严格更小，故 \eqref{eq:xi-binomial-toeplitz-two-endpoint-growth-rates} 成立。
再由
\[
|w_\ast\mp 1|^2=r_\ast^2\mp 2x+1
\]
得到
\[
Ar_\ast=r_\ast^2-2x+1,
\qquad
Br_\ast=r_\ast^2+2x+1.
\]
将两式相加、相减，分别得到
\[
(A+B)r_\ast=2(r_\ast^2+1),
\qquad
(A-B)r_\ast=-4x,
\]
即 \eqref{eq:xi-binomial-toeplitz-radius-recovery} 与 \eqref{eq:xi-binomial-toeplitz-realpart-recovery}。
由于二次方程两根之积为 \(1\)，其中恰有一根落在 \((0,1)\)，故 \(r_\ast\) 唯一。最后由 \(r_\ast^2=x^2+y^2\) 得 \(|y|\) 的表达式。证毕。
\end{proof}

\begin{theorem}[稀疏 Toeplitz 见证的尺度--深度交换律]\label{thm:xi-binomial-toeplitz-scale-depth-exchange}
在结论 \ref{con:xi-adams-coarsening-backprop-budget-product} 的记号下，对每个 \(m\ge 1\) 定义稀疏 Toeplitz 层级
\[
\mathbf T_N^{(m)}:=\bigl(c_{m(j-k)}\bigr)_{0\le j,k\le N},
\qquad
Q_{N,m}^{\pm}:=\bigl(b_{\pm}^{(N)}\bigr)^{\mathsf T}\mathbf T_N^{(m)}\,b_{\pm}^{(N)}.
\]
则存在 \(m_0\in\mathbb N\) 及常数 \(C_0,N_0>0\)，使得对所有 \(m\ge m_0\)：
\begin{enumerate}
  \item \(\{c_{mn}\}_{n\ge 1}\) 的主导极点为 \(w_\ast^m\)，其矩增长指数变为
  \[
  \boxed{\;
  \Lambda_{\mathrm{mom}}^{(m)}
  :=-\log|w_\ast^m|
  =m\,\Lambda_{\mathrm{mom}}.
  \;}
  \]
  \item 两端点增长率满足
  \[
  \limsup_{N\to\infty}|Q_{N,m}^-|^{1/N}
  =\frac{|w_\ast^m-1|^2}{|w_\ast|^m},
  \qquad
  \limsup_{N\to\infty}|Q_{N,m}^+|^{1/N}
  =\frac{|w_\ast^m+1|^2}{|w_\ast|^m}.
  \]
  \item 若记 \(N_{\mathrm{dom}}(m)\) 为主导项在尺度 \(m\) 上压过解析余项的最小层级，则
  \begin{equation}\label{eq:xi-binomial-toeplitz-scale-depth-exchange-law}
  \boxed{\;
  N_{\mathrm{dom}}(m)\le N_0+\frac{C_0}{m\,\Lambda_{\mathrm{mom}}}.
  \;}
  \end{equation}
\end{enumerate}
特别地，在定理 \ref{thm:xi-endpoint-binomial-falsifier} 的相位非退化口径下，某个端点见证在尺度 \(m\) 上的检出深度被压缩到 \((m\Lambda_{\mathrm{mom}})^{-1}\) 量级。
\end{theorem}

\begin{proof}
把定理 \ref{thm:xi-binomial-toeplitz-dominant-pole-response} 与定理 \ref{thm:xi-binomial-toeplitz-two-endpoint-pole-inversion} 逐字应用到稀疏子序列 \(\{c_{mn}\}_{n\ge 1}\)。
其主导极点由 \(w_\ast\) 变为 \(w_\ast^m\)，故
\[
\Lambda_{\mathrm{mom}}^{(m)}=-\log|w_\ast^m|=m\Lambda_{\mathrm{mom}}.
\]
又因 \(|w_\ast|<1\)，当 \(m\to\infty\) 时 \(w_\ast^m\to 0\)，从而
\[
|w_\ast^m\pm 1|\to 1.
\]
于是对充分大的 \(m\)，定理 \ref{thm:xi-binomial-toeplitz-two-endpoint-pole-inversion} 中的端点主导条件自动成立，得到第二条。
\par
对第三条，只需比较尺度 \(m\) 上主项与余项的指数率。
主项的增长因子为 \(e^{m\Lambda_{\mathrm{mom}}N}\) 乘以有界的端点因子 \(|w_\ast^m\pm 1|^{2N}\)，而解析余项仍至多保留 \(4^N\) 级的 Toeplitz 组合增长。
因此一旦
\[
e^{m\Lambda_{\mathrm{mom}}N}\gg 1,
\]
主导项便压过余项；把隐含常数写出即可得到 \eqref{eq:xi-binomial-toeplitz-scale-depth-exchange-law}。
最后，定理 \ref{thm:xi-endpoint-binomial-falsifier} 的相位振荡机制在尺度 \(m\) 上仍然适用，故主导项一旦进入支配区，便会沿某个无穷子序列给出严格负值。证毕。
\end{proof}

\begin{corollary}[二次红移压缩下的常数深度检出]\label{cor:xi-binomial-toeplitz-constant-depth-under-quadratic-redshift}
若最深极点由某个离线零点 \(\rho=\beta+\mathrm{i}\gamma\)（\(\delta:=\tfrac12-\beta>0\)）所诱导，则由 \eqref{eq:xi-horizon-Lambda-closed}--\eqref{eq:xi-horizon-Lambda-asymptotic} 有
\[
\Lambda_{\mathrm{mom}}
=
\frac12\log\!\left(\frac{\gamma^2+(1+\delta)^2}{\gamma^2+(1-\delta)^2}\right)
=
\frac{2\delta}{\gamma^2}+O(\gamma^{-4}).
\]
因此对任意固定常数 \(c>1\)，若取
\[
m=\left\lceil \frac{c}{\Lambda_{\mathrm{mom}}}\right\rceil,
\]
则存在只依赖于解析余项结构的常数 \(N_\sharp\)，使得
\[
N_{\mathrm{dom}}(m)\le N_\sharp.
\]
在定理 \ref{thm:xi-endpoint-binomial-falsifier} 的相位非退化口径下，某个尺度 \(m\) 的端点稀疏 Toeplitz 见证因而在统一常数深度内出现严格负值，并可无穷次重复出现。
\end{corollary}

\begin{proof}
由定理 \ref{thm:xi-binomial-toeplitz-scale-depth-exchange}，
\[
N_{\mathrm{dom}}(m)
\le
N_0+\frac{C_0}{m\Lambda_{\mathrm{mom}}}.
\]
而 \(m\Lambda_{\mathrm{mom}}\ge c\)，故
\[
N_{\mathrm{dom}}(m)\le N_0+\frac{C_0}{c}.
\]
令 \(N_\sharp:=N_0+\lceil C_0/c\rceil\) 即得结论。证毕。
\end{proof}

\begin{theorem}[共动图表中的高度解耦尺度律]\label{thm:xi-binomial-toeplitz-comoving-delta-scale-law}
在共动选择 \(x_0=\gamma\) 下，由 \eqref{eq:xi-comoving-pure-delta} 有
\[
w_{\rho,\gamma}=\frac{1-\delta}{1+\delta}\in(0,1).
\]
定义共动矩增长指数
\[
\Lambda_{\mathrm{com}}(\delta)
:=
-\log|w_{\rho,\gamma}|
=
\log\!\left(\frac{1+\delta}{1-\delta}\right).
\]
对任意固定 \(c>1\)，取
\[
m_{\mathrm{com}}(\delta):=
\left\lceil \frac{c}{\Lambda_{\mathrm{com}}(\delta)}\right\rceil.
\]
则当 \(\delta\downarrow 0\) 时有
\[
\boxed{\;
m_{\mathrm{com}}(\delta)\asymp \delta^{-1},
\;}
\]
并且存在只依赖于 \(c\) 与解析余项结构的常数 \(N_{\mathrm{com}}\)，使得共动尺度 \(m_{\mathrm{com}}(\delta)\) 把端点稀疏 Toeplitz 见证的主导深度压缩到
\[
N_{\mathrm{dom}}\!\bigl(m_{\mathrm{com}}(\delta)\bigr)\le N_{\mathrm{com}},
\]
该结论对所有高度 \(\gamma\) 一致有效。
\end{theorem}

\begin{proof}
由 \eqref{eq:xi-comoving-pure-delta}，
\[
\Lambda_{\mathrm{com}}(\delta)
=
\log\!\left(\frac{1+\delta}{1-\delta}\right)
=
2\delta+O(\delta^2)
\qquad(\delta\downarrow 0),
\]
故
\[
m_{\mathrm{com}}(\delta)
\asymp \frac{1}{\Lambda_{\mathrm{com}}(\delta)}
\asymp \delta^{-1}.
\]
再将定理 \ref{thm:xi-binomial-toeplitz-scale-depth-exchange} 中的
\(\Lambda_{\mathrm{mom}}\) 替换为 \(\Lambda_{\mathrm{com}}(\delta)\)，并注意该量不含 \(\gamma\)，便得到高度无关的常数深度界。
最后一段与推论 \ref{cor:xi-binomial-toeplitz-constant-depth-under-quadratic-redshift} 同理。证毕。
\end{proof}

\endinput


\begin{conclusion}[算子值视界正性与 strictification 的统一：完全正性失败等价于有限维否证器残留]\label{con:xi-operator-valued-horizon-cp-falsifier}
将标量 Carath\'eodory 正性提升为取值于有限维 \(C^\ast\)-代数 \(\mathcal{A}\) 的算子值 Carath\'eodory 类时，通过 Peter--Weyl 分解
\(
\mathcal{A}\cong\bigoplus_{\pi}M_{d_\pi}(\CC)
\)
可将完全正性严格对齐为有限个表示块上的 Toeplitz--PSD 条件：算子值函数 \(F:\mathbb{D}\to\mathcal{A}\) 满足 \(\Re F\succeq 0\) 当且仅当其块 Toeplitz 截断 \(\mathbf{T}^{\mathcal{A}}_N(F)\succeq 0\) 对所有 \(N\) 成立。
在 strictification 语境中，最大可见商给出完全正可见实现的最小封装；真共振或不可 strictify 的残差被严格化为完全正性的有限维否证器，从而把多通道情形下的可证伪点压缩为有限表示分量上的 PSD 失败。
\end{conclusion}
\begin{proof}[证明要点]
完全正性与块 Toeplitz--PSD 的等价由定理 \ref{thm:xi-operator-herglotz-toeplitz-completeness} 给出；
表示分解下的逐通道对齐见定理 \ref{thm:xi-fourier-pontryagin-block-toeplitz}；
strictification 的完全正可见实现解释见定理 \ref{thm:xi-gbc-maximal-cp-visible-realization}。
\end{proof}

\paragraph{模障碍、cyclotomic 因子与相对谱隙：共振角色群的有限可计算化}

\begin{conclusion}[模障碍的可枚举化：非平凡扭曲通道达半径当且仅当出现 cyclotomic 因子]\label{con:xi-mod-obstruction-cyclotomic-factor}
对 Ihara/动力学 \(\zeta\) 的扭曲通道
\[
Z_K(u,t;\rho):=\det(I-tB_\rho(u))^{-1},
\]
存在非平凡模障碍当且仅当某个 \(\rho\neq \mathbf{1}\) 满足 \(\rho(B_\rho(u))=\lambda_1(u)\)（与平凡通道同半径）；
该条件等价于 \(\det(I-tB_\rho(u))\) 含 cyclotomic 因子（即单位圆上的单位根因子）。
因此障碍来源可被约化为有限维整系数多项式的 cyclotomic 因子分类问题；该判别在一般扭曲图 \(\zeta\) 口径中同样成立。
\end{conclusion}
\begin{proof}[证明要点]
见结论 \ref{con:xi-terminal-mod-obstruction-cyclotomic} 及其引用的分类定理 \ref{thm:conclusion75-ihara-cyclotomic-factor-classification}。
\end{proof}

\begin{conclusion}[障碍剥离后的相对谱隙：共振仅占有限维可识别子空间]\label{con:xi-relative-gap-after-obstruction-finite-resonant-subspace}
令同调障碍诱导闭子群 \(\mathcal{H}\subset\widehat{\ZZ}\)，并记扭曲转移算子族为 \(A(\chi)\)。
则对任意角色 \(\chi\) 满足 \(\chi|_{\mathcal{H}}\neq 1\) 时，存在统一常数 \(\delta>0\) 使
\[
\boxed{\ \rho(A(\chi))\ \le\ (1-\delta)\rho(A(1))\ }.
\]
因此即便存在有限模障碍，在障碍角色张成子空间的正交补上仍获得统一扭曲谱隙与相对改进；共振被严格限制为一个可计算的有限阿贝尔群，其结构可由同调 torsion 与 Smith 标准形显式刻画。
\end{conclusion}
\begin{proof}[证明要点]
相对谱隙与正交补改进见定理 \ref{thm:gm-relative-gap-after-obstruction}；
共振角色群的同调刻画与 Smith 形可计算性分别见定理 \ref{thm:gm-resonant-character-group-homology} 与命题 \ref{prop:gm-smith-compute-obstructions}。
\end{proof}

\paragraph{\texorpdfstring{$p$}{p}-幂塔的同余压缩：Frobenius 稀疏化、Chebyshev 闭合与导数滤波}

\begin{conclusion}[\texorpdfstring{$p^k$}{p^k}-稀疏化的普适 Frobenius 规律：模 \texorpdfstring{$p^\ell$}{p^ell} 下只“看见”\texorpdfstring{$u^{p^{k-\ell+1}}$}{u^{p^{k-l+1}}}]\label{con:xi-witt-frobenius-sparsification-universal}
对周期迹多项式 \(a_{p^k}(u)\in\ZZ[u]\)，Dwork 型同余导出：在商环 \((\ZZ/p^\ell\ZZ)[u]\) 中，\(a_{p^k}(u)\) 仅依赖于变量 \(u^{p^{k-\ell+1}}\)，等价地
\[
\boxed{\ a_{p^k}(u)\bmod p^\ell\in(\ZZ/p^\ell\ZZ)\bigl[u^{p^{k-\ell+1}}\bigr]\ }.
\]
因此 \(p\)-幂塔的高层模信息可被确定性地推回低层数据，实现模审计复杂度随层级对数增长的结构性压缩。
\end{conclusion}
\begin{proof}[证明要点]
见命题 \ref{prop:witt-frobenius-iterated-descent}（由 Dwork 同余逐级降模迭代得到）。
\end{proof}

\begin{conclusion}[\texorpdfstring{$p=2$}{p=2} 的 Chebyshev--Dwork 闭合：矩阵幂同余递推降维为单变量代入 \texorpdfstring{$t\mapsto t^2-2$}{t->t^2-2}]\label{con:xi-chebyshev-dwork-closure-p2}
在不变量坐标 \(t=u+u^{-1}\) 下，偶次回文迹 \(a_{2^k}(u)=u^{2^{k-1}}A_{2^k}(t)\) 满足同余闭合递推
\[
\boxed{\ A_{2^k}(t)\equiv A_{2^{k-1}}(t^2-2)\pmod{2^k}\qquad(k\ge 2)\ }.
\]
从而对固定模 \(2^k\) 的计算可将对象从多项式矩阵幂 \(B(u)^{2^k}\) 归约为“单变量代入 + 取模”的对数深度迭代；该闭合把二幂塔上的 Frobenius 推前显式识别为 Chebyshev 动力系统。
\end{conclusion}
\begin{proof}[证明要点]
见定理 \ref{thm:chebyshev-dwork-congruence-chain} 与推论 \ref{cor:chebyshev-dwork-matrixfree}。
\end{proof}

\begin{conclusion}[\texorpdfstring{$\theta$}{theta}-导数的 \texorpdfstring{$p$}{p}-进滤波：矩满足严格缩放同余与估值下界]\label{con:xi-theta-derivative-padic-filter}
令 \(u=e^\theta\)，并定义导数矩
\[
A_{n}^{(r)}:=\left.\frac{\mathrm{d}^r}{\mathrm{d}\theta^r}\,a_n(e^\theta)\right|_{\theta=0}\in\ZZ\qquad(r\ge 0).
\]
则 Dwork 同余蕴含逐阶缩放同余
\[
\boxed{\ A_{p^k}^{(r)}\equiv p^r\,A_{p^{k-1}}^{(r)}\pmod{p^k}\ }\qquad(r\ge 0),
\]
特别地当 \(0\le r<k\) 时得到 \(p\)-进估值下界
\[
\boxed{\ p^{\,k-r}\ \big|\ A_{p^k}^{(r)}\ }.
\]
因此在 \(p\)-幂塔上，所有高阶导数矩自动落入严格的 \(p\)-进滤波窗，给出一个无需额外结构即可核验的 \(p\)-进刚性证书。
\end{conclusion}
\begin{proof}[证明要点]
见命题 \ref{prop:witt-theta-derivative-filter}。
\end{proof}

\paragraph{自然边界与自由度预算：素数层级数的不可延拓约束与零点密度门槛}

\begin{conclusion}[primitive--Dirichlet 级数的自然边界：\texorpdfstring{$\Re(s)=0$}{Re(s)=0} 上稠密分支奇点强制不可延拓]\label{con:xi-primedirichlet-natural-boundary}
对 primitive 轨道 Dirichlet 级数 \(\mathcal{P}(s)=\sum_{n\ge 1}p_n\lambda^{-ns}\)，可通过 M\"obius--Abel 型表示在 \(\Re(s)>0\) 上解析延拓；
但对每个奇素数 \(p\) 与任意 \(m\in\ZZ\)，点
\[
s_{p,m}:=\frac{1}{p}+\frac{2\pi \mathrm{i} m}{p\log\lambda}
\]
为对数型分支奇点，且 \(\{s_{p,m}\}\) 在直线 \(\Re(s)=0\) 上稠密。
因此 \(\Re(s)=0\) 为自然边界：不存在任何穿越该直线的局部解析延拓。
该约束表明组织“素数影子”误差项时必须优先通过 \(\zeta\)（或其对数导数/显式公式）而非直接 primitive--Dirichlet 级数来实现可控叠加。
\end{conclusion}
\begin{proof}[证明要点]
见定理 \ref{thm:real-input-40-primedirichlet-dense-branch}（稠密分支奇点机制与自然边界推断）。
\end{proof}

\begin{conclusion}[有限维 Dirichlet 化的 \texorpdfstring{$\Theta(\log T)$}{Theta(log T)} 门槛：固定带宽范式在零点密度阶上结构性失败]\label{con:xi-finite-dimension-dirichlet-logT-threshold}
若目标是逼近 \(\Xi_\zeta\) 型零点计数主阶 \(\widetilde N_\zeta(T)\sim \tfrac{T}{\pi}\log T\)，则任何长期停留在“固定有限带宽/固定有限尺度窗/固定有限维行列式口径”的实现，其对称零点计数至多为 \(O(T)\)，从而在密度阶上必然失败。
因此必须释放 \(\Theta(\log T)\) 级自由度预算（例如使指数型 \(\tau(T)\) 随 \(T\) 增长并满足 \(\tau(T)\gtrsim \log T\)，进入增长维数逼近序列或无穷维 Fredholm 行列式口径）。
\end{conclusion}
\begin{proof}[证明要点]
Cartwright 密度上界给出 \(\widetilde N_T(T)\le \frac{2}{\pi}\tau(T)T+o(\tau(T)T)\)；
若要匹配 \(\widetilde N_\zeta(T)\sim \frac{T}{\pi}\log T\)，必须有 \(\tau(T)\gtrsim \frac12\log T\)。
见结论 \ref{con:xi-terminal-zero-density-product-lower-bound}。
\end{proof}

\paragraph{相位闭包与单位圆根证书链：固定切片、Jensen 缺陷与 Schur--Cohn/LMI}

\begin{conclusion}[相位化闭包的唯一固定切片：拒绝外置径向寄存器时零点被制度性锁定在 \texorpdfstring{$\Re(s)=\tfrac12$}{Re(s)=1/2}]\label{con:xi-unitary-slice-unique-fixed-slice}
在 unitary-slice 读出制度中，若核验对象可表示为
\(
D(s)=\det(I-u_L(s)U)
\)
（\(U\) 为酉算子，\(u_L(s)=L^{2s-1}\)），且制度拒绝在模式轴外置径向/幅度寄存器，则任何可被寻址的零点事件必满足 \(|u_L(s)|=1\)，从而 \(\Re(s)=\tfrac12\)。
等价地，\(\Re(s)=\tfrac12\) 是相位化闭包在完成化坐标中的唯一固定切片，为有限型 Hilbert--P{\'o}lya 路线提供制度内的强约束框架。
\end{conclusion}
\begin{proof}[证明要点]
由 \(u_L(s)=L^{2s-1}\) 得 \(\Re(s)=\tfrac12\Longleftrightarrow |u_L(s)|=1\)；
在可见域取 \(\mathsf{Phase}\) 且拒绝外置径向寄存器时，\(|u_L(s)|\neq 1\) 的点按协议读出为 \(\mathsf{NULL}\)，从而可寻址零点被锁定于临界线。
见 \S\ref{subsec:xi-pw-type-safety-null}（特别是定理 \ref{thm:xi-offset-null-type-safety}）。
\end{proof}

\begin{conclusion}[Jensen 缺陷作为视界熵距离：偏离可被连续量化并可核验]\label{con:xi-jensen-defect-horizon-entropy-distance}
对 self-inversive 完成化多项式 \(P\)（满足 \(y^dP(1/y)=e^{\mathrm{i}\phi}P(y)\)，并假设 \(P(0)\neq 0\)），对 \(0<\varrho<1\) 定义 Jensen 缺陷
\[
D_P(\varrho):=\frac{1}{2\pi}\int_{0}^{2\pi}\log|P(\varrho e^{\mathrm{i}\theta})|\,\dd\theta-\log|P(0)|.
\]
则“全部零点在单位圆”与 \(D_P(\varrho)\equiv 0\)（\(0<\varrho<1\)）等价，亦等价于
\(\lim_{\varrho\uparrow 1}D_P(\varrho)=0\)。
因此 \(D_P(\varrho)\) 可被解释为对径向/幅度泄漏的连续缺陷函数：在保持 self-inversive（完成化自对偶）结构不变时，它给出从“严格 Ramanujan/RH-sharp”到“可观缺陷”的可审计距离尺度。
\end{conclusion}
\begin{proof}[证明要点]
对 \(P\) 应用 Jensen 公式得 \(D_P(\varrho)=\sum_{|a_k|<\varrho}\log(\varrho/|a_k|)\ge 0\)；
因此 \(D_P(\varrho)\equiv 0\) 等价于对任意 \(\varrho<1\) 无盘内零点。
若 \(P\) self-inversive，则零点关于 \(y\mapsto 1/\overline y\) 成对，故无盘内零点等价于全部零点落在单位圆。
本稿的函数型口径见命题 \ref{prop:app-jensen-formula-defect} 与定理 \ref{thm:app-rh-iff-disk-zero-free}。
\end{proof}

\begin{conclusion}[Cayley 超双曲化与 Cohn 导数盘内性：单位圆根判据压缩为可编译的半正定证书链]\label{con:xi-cayley-cohn-lmi-certificate-chain}
对自倒数实系数多项式 \(\Phi_L\) 的“单位圆全根”条件，可经 Cayley 变换等价转写为一元实多项式 \(\Psi_L\) 的全实根（超双曲性）；
同时由 Cohn 判据，单位圆全根又等价于 \(\Phi_L\) self-inversive 且 \(\Phi_L'\) 的全部根落在闭单位圆盘内。
因此单位圆根几何可被压缩为两条可实现的审计路径：实轴超双曲性路径与导数盘内性路径；二者均可进一步对接 Schur--Cohn/Toeplitz 型半正定证书，从而把根位置断言制度化为可核验的 LMI 链条。
\end{conclusion}
\begin{proof}[证明要点]
Cayley 超双曲化见命题 \ref{prop:xi-rh-cayley-hyperbolicity}；
Cohn 型导数盘内性判别见命题 \ref{prop:xi-rh-cohn-derivative}；
Schur--Cohn 半正定接口与 Toeplitz 证书的对齐见注 \ref{rem:xi-schur-cohn-sdp-interface} 与定理 \ref{thm:xi-toeplitz-schur-cohn-lift}。
\end{proof}

\begin{conclusion}[阈值共振的区间化压缩：\texorpdfstring{$\QQ[u+u^{-1}]$}{Q[u+u^{-1}]} 系数下的区间 PSD/一元正性证书]\label{con:xi-threshold-resonance-interval-psd-compression}
在 unitary slice（\(|u|=1\)）上令
\[
S:=u+u^{-1}\in[-2,2].
\]
若证书对象（例如 Schur--Cohn Hermitian 形式或 Toeplitz/Carath\'eodory 证书所依赖的系数）落在 \(\QQ[u+u^{-1}]=\QQ[S]\)，则“单位圆根/阈值共振”条件可被压缩为关于 \(S\in[-2,2]\) 的一元区间判定（区间根或区间正性）。
在一元情形，区间正性可通过 Gram 矩阵的半正定性给出，从而把高维 Schur--Cohn/Toeplitz 型 LMI 证书进一步降为区间 PSD 可行性（等价地：一元正性证书）。
结合完成化迹坐标的有理参数化 \(S(x)=\frac{2(1-x^2)}{1+x^2}\)（式 \eqref{eq:phase-completion-trace-rational}）与 \(\tau_d\) 的全有理半群实现（命题 \ref{prop:chebyshev-adams-fully-rational}、推论 \ref{cor:chebyshev-adams-rational-semigroup}），上述区间证书可被编译为有理函数复合与整系数递推，从而提供一条无需高维 SDP 的最简可审计通道。
\end{conclusion}

\paragraph{自对偶完成化、有限型 Hilbert--P{\'o}lya 与素相位环流}
在相位化闭包的固定切片、Jensen 缺陷的径向量化、单位圆根的 Cayley/Schur--Cohn/Toeplitz 证书链以及有限窗 CMV 酉谱实现已经建立之后，还可进一步把 self-dual completion、unitary slice、视界正性与 prime-phase lifting 收束为一组闭合后果。它们分别刻画九重等价闭环与双倍自伴压缩、实越界根诱导的双塔算术结构及计数律、素相位流在素因子环面上的稠密唯一遍历、RH 证书成立时 Toeplitz 主子矩阵的秩冻结，以及潜在 prime-phase 表示所诱导的 Bohr 几乎周期零点几何。

\begin{theorem}[有限型 Hilbert--P{\'o}lya 的九重等价闭包与双倍自伴压缩]\label{thm:xi-hilbert-polya-ninefold-equivalence-doubled-selfadjoint-compression}
\leanverified{paper\_xi\_hilbert\_polya\_ninefold\_equivalence\_doubled\_selfadjoint\_compression}
沿用命题 \ref{prop:xi-rh-interval-criterion}、推论 \ref{cor:xi-rh-unit-circle-reciprocal}、命题 \ref{prop:xi-rh-cayley-hyperbolicity}、定理 \ref{thm:xi-slice-horizon-carath-toeplitz}、推论 \ref{cor:xi-slice-weyl-herglotz} 与命题 \ref{prop:xi-hilbert-polya-cmv} 的记号。固定 \(L>1\)，令
\[
\widetilde P_L(S):=\frac{P_L(S)}{\operatorname{lc}(P_L)},
\qquad
d:=\deg P_L,
\]
\[
\widetilde\Phi_L(y):=y^d\widetilde P_L(y+y^{-1}),
\]
并令 \(\Psi_L\) 为命题 \ref{prop:xi-rh-cayley-hyperbolicity} 中的 Cayley 多项式，\(\mathscr{C}_{\Phi_L}\) 为定理 \ref{thm:xi-slice-horizon-carath-toeplitz} 中的视界 Carath\'eodory 函数，\(\mathbf{T}_N\) 为其 Toeplitz 主子矩阵，\(m_{\Delta,L}\) 为推论 \ref{cor:xi-slice-weyl-herglotz} 中的 Weyl/Herglotz 对数导数，且 \(U_L^{\mathrm{CMV}}\) 为命题 \ref{prop:xi-hilbert-polya-cmv} 所给的 CMV-like 伴随矩阵。则下列九条断言等价：
\begin{enumerate}
  \item \(\Xi_{\Delta,L}\) 的全部零点满足 \(\Re s=\tfrac12\)；
  \item \(\widetilde P_L\) 的全部根落在 \([-2,2]\)；
  \item \(\widetilde\Phi_L\) 的全部根落在单位圆 \(\TT\)；
  \item \(\Psi_L\) 的全部根都落在 \(\overline{\RR}=\RR\cup\{\infty\}\)；
  \item \(\mathscr{C}_{\Phi_L}\) 属于 Carath\'eodory 类；
  \item 对所有 \(N\ge 0\)，有 \(\mathbf{T}_N\succeq 0\)；
  \item \(m_{\Delta,L}\) 在 \(\mathbb H\) 上为 Herglotz 函数；
  \item \(U_L^{\mathrm{CMV}}\) 为酉矩阵；
  \item 存在一个 \(2d\times 2d\) 自伴矩阵 \(A_L\) 使
  \[
  \operatorname{Spec}(A_L)\subseteq[-2,2],
  \qquad
  \det(SI_{2d}-A_L)=\widetilde P_L(S)^2.
  \]
\end{enumerate}
\end{theorem}
\begin{proof}
由定理 \ref{thm:xi-slice-horizon-carath-toeplitz}，可得
\[
(1)\Longleftrightarrow(2)\Longleftrightarrow(3)\Longleftrightarrow(5)\Longleftrightarrow(6).
\]
由命题 \ref{prop:xi-rh-cayley-hyperbolicity}，有
\[
(3)\Longleftrightarrow(4).
\]
由推论 \ref{cor:xi-slice-weyl-herglotz}，有
\[
(1)\Longleftrightarrow(7).
\]
由命题 \ref{prop:xi-hilbert-polya-cmv}，有
\[
(3)\Longleftrightarrow(8).
\]
因此只需证明 \((8)\Longleftrightarrow(9)\)。

先证 \((8)\Rightarrow(9)\)。若 \(U:=U_L^{\mathrm{CMV}}\) 为酉矩阵，则令
\[
A_L:=U+U^{-1}=U+U^\ast.
\]
于是 \(A_L\) 自伴；并由定理 \ref{thm:phase-lift-spectral-bound} 的谱映射结论可知
\[
\operatorname{Spec}(A_L)\subseteq[-2,2].
\]
另一方面，由命题 \ref{prop:xi-hilbert-polya-cmv}，
\[
\widetilde\Phi_L(y)=\det(yI_{2d}-U),
\]
且
\[
\widetilde\Phi_L(y)=y^d\widetilde P_L(y+y^{-1}).
\]
若把 \(\widetilde P_L\) 的根按重数记为 \(S_1,\dots,S_d\)，则存在相应的 \(y_{j,\pm}\) 满足
\[
y_{j,+}y_{j,-}=1,
\qquad
y_{j,+}+y_{j,-}=S_j,
\]
并且
\[
\widetilde\Phi_L(y)=\prod_{j=1}^{d}(y-y_{j,+})(y-y_{j,-}).
\]
故 \(U\) 的特征值恰为 \(\{y_{j,+},y_{j,-}\}_{j=1}^{d}\)，而 \(A_L=U+U^{-1}\) 的对应特征值为
\[
y_{j,\pm}+y_{j,\pm}^{-1}=S_j.
\]
因此每个 \(S_j\) 在 \(A_L\) 中恰出现两次，从而
\[
\det(SI_{2d}-A_L)=\prod_{j=1}^{d}(S-S_j)^2=\widetilde P_L(S)^2.
\]
这就得到第 \((9)\) 条。

再证 \((9)\Rightarrow(2)\)。若存在满足 \((9)\) 的自伴矩阵 \(A_L\)，则其谱包含于 \([-2,2]\)。而 \(\widetilde P_L(S)^2\) 作为 \(A_L\) 的特征多项式，意味着 \(\widetilde P_L\) 的全部根都是 \(A_L\) 的特征值，因而也都属于 \([-2,2]\)。故第 \((2)\) 条成立。证毕。
\end{proof}

\begin{theorem}[实 Lee--Yang 越界根产生恰好两条垂直零塔，偏移量由 \texorpdfstring{$\operatorname{arcosh}$}{arcosh} 严格给出]\label{thm:xi-real-outlier-root-two-vertical-zero-towers-arcosh}
沿用命题 \ref{prop:xi-rh-interval-criterion} 的记号。设 \(S_\ast\in\RR\) 是 \(P_L(S)\) 的一个实根，代数重数为 \(\mu_\ast\ge 1\)，并满足 \(|S_\ast|>2\)。记
\[
\eta_\ast:=\operatorname{arcosh}\!\Bigl(\frac{|S_\ast|}{2}\Bigr)>0,
\qquad
\delta_\ast:=
\begin{cases}
0,& S_\ast>2,\\
1,& S_\ast<-2.
\end{cases}
\]
\leanverified{paper\_xi\_real\_outlier\_root\_two\_vertical\_zero\_towers\_arcosh}
则由该根贡献的全部零点恰为两条垂直算术塔
\[
s_{k,\pm}(S_\ast)
=
\frac12\pm\frac{\eta_\ast}{\log L}
+\frac{(2k+\delta_\ast)\pi\,\mathrm{i}}{\log L},
\qquad k\in\ZZ,
\]
并且每一点都以重数 \(\mu_\ast\) 出现。

若记
\[
M_L^{\mathrm{out}}
:=
\sum_{\substack{S_\ast\in\RR\\ |S_\ast|>2}}\mu_\ast,
\]
并令 \(N^{\mathrm{off}}_{L,\mathrm{real}}(T)\) 表示所有由实越界根产生且满足 \(|\Im s|\le T\) 的零点总数（按重数计），则有
\[
N^{\mathrm{off}}_{L,\mathrm{real}}(T)
=
\frac{2\log L}{\pi}\,M_L^{\mathrm{out}}\,T+O(1)
\qquad(T\to\infty).
\]
\end{theorem}
\begin{proof}
由推论 \ref{cor:xi-single-scale-radial-deviation-arcosh}，若 \(S_\ast\) 为 \(P_L\) 的实根，则相应二次方程
\[
y+y^{-1}=S_\ast
\]
的两根为
\[
y_\pm=\frac{S_\ast\pm\sqrt{S_\ast^2-4}}{2},
\]
并且全部零点满足
\[
s=\frac12+\frac{\log y_\pm+2\pi\mathrm{i}k}{\log L},
\qquad k\in\ZZ.
\]
当 \(S_\ast>2\) 时，可写 \(S_\ast=2\cosh\eta_\ast\)，于是
\[
y_\pm=e^{\pm\eta_\ast},
\]
从而
\[
s_{k,\pm}(S_\ast)=\frac12\pm\frac{\eta_\ast}{\log L}+\frac{2k\pi\mathrm{i}}{\log L}.
\]
当 \(S_\ast<-2\) 时，可写 \(S_\ast=-2\cosh\eta_\ast\)，于是
\[
y_\pm=-e^{\pm\eta_\ast}=e^{\pm\eta_\ast+(2k+1)\pi\mathrm{i}},
\]
从而
\[
s_{k,\pm}(S_\ast)=\frac12\pm\frac{\eta_\ast}{\log L}+\frac{(2k+1)\pi\mathrm{i}}{\log L}.
\]
两种情形合并即得所示统一公式。重数 \(\mu_\ast\) 的保持来自复合方程 \(\Xi_{\Delta,L}(s)=P_L(L^{s-1/2}+L^{1/2-s})\) 中多项式根重数的直接继承。

对计数公式，只需对单条塔
\[
\frac12+\sigma+\frac{(2k+\delta)\pi\mathrm{i}}{\log L},
\qquad k\in\ZZ,
\]
计算满足 \(|\Im s|\le T\) 的整数 \(k\) 个数。该个数为
\[
\frac{\log L}{\pi}T+O(1).
\]
每个实越界根贡献两条这样的塔，且按重数 \(\mu_\ast\) 计入，故对全部实越界根求和即得
\[
N^{\mathrm{off}}_{L,\mathrm{real}}(T)
=
2\cdot\frac{\log L}{\pi}T
\sum_{\substack{S_\ast\in\RR\\ |S_\ast|>2}}\mu_\ast+O(1)
=
\frac{2\log L}{\pi}\,M_L^{\mathrm{out}}\,T+O(1).
\]
证毕。
\end{proof}

\begin{theorem}[unitary slice 的素相位提升在素因子环面上稠密且唯一遍历]\label{thm:xi-prime-phase-torus-dense-uniquely-ergodic}
沿用定理 \ref{thm:xi-unitary-slice-prime-phase-lift-to-arithmetic-singular-ring} 的记号。对整数 \(L\ge 2\)，记
\[
S(L):=\{p:\ v_p(L)>0\},
\qquad
\omega_{p,L}:=2\,v_p(L)\log p.
\]
并定义素相位轨道
\[
\vartheta_L:\RR\longrightarrow \TT^{S(L)},
\qquad
\vartheta_L(t):=\bigl(u_{p,L}(t)\bigr)_{p\in S(L)}
=\bigl(e^{\mathrm{i}\omega_{p,L}t}\bigr)_{p\in S(L)}.
\]
则 \(\vartheta_L(\RR)\) 在整个环面 \(\TT^{S(L)}\) 上稠密，且该线性流唯一遍历。等价地，对任意连续函数 \(F\in C(\TT^{S(L)})\)，都有
\[
\lim_{T\to\infty}\frac1T\int_0^T F\!\bigl(\vartheta_L(t)\bigr)\,dt
=
\int_{\TT^{S(L)}}F\,dm,
\]
其中 \(m\) 为 \(\TT^{S(L)}\) 上的 Haar 概率测度。
\end{theorem}
\begin{proof}
只需证明频率向量
\[
\Bigl(\frac{\omega_{p,L}}{2\pi}\Bigr)_{p\in S(L)}
=
\Bigl(\frac{v_p(L)\log p}{\pi}\Bigr)_{p\in S(L)}
\]
在 \(\QQ\) 上线性无关。设存在有理数 \(q_p\) 使
\[
\sum_{p\in S(L)} q_p\,\omega_{p,L}=0.
\]
取公共分母 \(N\) 并令 \(n_p:=Nq_p\in\ZZ\)，则
\[
\sum_{p\in S(L)} n_p\,v_p(L)\log p=0,
\]
从而
\[
\prod_{p\in S(L)}p^{\,n_p v_p(L)}=1.
\]
由算术基本定理，必有 \(n_p v_p(L)=0\) 对所有 \(p\in S(L)\) 成立；又因 \(v_p(L)>0\)，故 \(n_p=0\) 对所有 \(p\) 成立，进而 \(q_p=0\)。这就证明了 \(\QQ\)-线性无关。

于是 Kronecker--Weyl 定理直接给出：线性流 \(t\mapsto (e^{\mathrm{i}\omega_{p,L}t})_{p\in S(L)}\) 在 \(\TT^{S(L)}\) 上稠密且唯一遍历。定理 \ref{thm:xi-unitary-slice-prime-phase-lift-to-arithmetic-singular-ring} 说明 \(\vartheta_L\) 正是 \(\tau_L^{\mathrm{prime}}\) 的 \(\TT^{S(L)}\)-分量，因此结论同时刻画了素相位 lifting 的环面动力学。证毕。
\leanverified{paper\_xi\_prime\_phase\_torus\_dense\_uniquely\_ergodic}
\end{proof}

\begin{theorem}[RH 证书成立时 Toeplitz 主子矩阵的秩在有限步后精确稳定]\label{thm:xi-toeplitz-rank-stabilizes-to-distinct-unitcircle-roots}
沿用定理 \ref{thm:xi-hilbert-polya-ninefold-equivalence-doubled-selfadjoint-compression} 的记号，并假设其等价条件成立。于是 \(\widetilde\Phi_L\) 的全部根都在单位圆上。记其不同根为
\[
\xi_1,\dots,\xi_s\in\TT,
\]
对应重数为
\[
m_1,\dots,m_s,
\qquad
\sum_{\ell=1}^{s}m_\ell=D,
\qquad
D:=\deg\widetilde\Phi_L=2d.
\]
定义有限原子概率测度
\[
\mu_{\Phi_L}
:=
\sum_{\ell=1}^{s}\frac{m_\ell}{D}\,\delta_{\xi_\ell},
\]
并令其 Fourier 系数为
\[
c_n:=\int_{\TT}\xi^{-n}\,d\mu_{\Phi_L}(\xi),
\qquad
\mathbf{T}_N:=(c_{j-k})_{0\le j,k\le N}.
\]
则对每个 \(N\ge 0\) 都有
\[
\operatorname{rank}\mathbf{T}_N=\min(N+1,s).
\]
因此当 \(N\ge s-1\) 时，
\[
\operatorname{rank}\mathbf{T}_N=s
\]
精确稳定。
\end{theorem}
\begin{proof}
由注 \ref{rem:xi-slice-horizon-measure-reflection-positivity}，在单位圆根条件成立时，\(\mathscr{C}_{\Phi_L}\) 的 Herglotz 表示可由有限原子测度 \(\mu_{\Phi_L}\) 给出。于是
\[
c_{j-k}
=
\sum_{\ell=1}^{s}\frac{m_\ell}{D}\,\xi_\ell^{-(j-k)}
=
\sum_{\ell=1}^{s}\frac{m_\ell}{D}\,\xi_\ell^{-j}\,\overline{\xi_\ell^{-k}},
\]
因为 \(|\xi_\ell|=1\)。令
\[
V_N:=(\xi_\ell^{-j})_{0\le j\le N,\ 1\le \ell\le s}\in M_{(N+1)\times s}(\CC),
\qquad
W:=\operatorname{diag}\!\Bigl(\frac{m_1}{D},\dots,\frac{m_s}{D}\Bigr).
\]
则
\[
\mathbf{T}_N=V_N W V_N^\ast.
\]
由于 \(W\) 为正定对角矩阵，故
\[
\operatorname{rank}\mathbf{T}_N=\operatorname{rank}V_N.
\]

若 \(N+1\ge s\)，则 \(V_N\) 的某个 \(s\times s\) Vandermonde 主子式为
\[
\det(\xi_\ell^{-j})_{0\le j\le s-1,\ 1\le \ell\le s}
=
\prod_{1\le r<\ell\le s}(\xi_\ell^{-1}-\xi_r^{-1}),
\]
因 \(\xi_1,\dots,\xi_s\) 两两不同而非零，故该子式非零，从而 \(\operatorname{rank}V_N=s\)。

若 \(N+1<s\)，则 \(V_N\) 至多有秩 \(N+1\)。另一方面，取任意 \(N+1\) 个不同原子，例如 \(\xi_1,\dots,\xi_{N+1}\)，对应的 \((N+1)\times(N+1)\) Vandermonde 子式同样非零，故 \(\operatorname{rank}V_N=N+1\)。

综上，
\[
\operatorname{rank}\mathbf{T}_N=\operatorname{rank}V_N=\min(N+1,s),
\]
证毕。
\end{proof}
\leanverified{paper\_xi\_toeplitz\_rank\_stabilizes\_to\_distinct\_unitcircle\_roots}

\begin{corollary}[simple-spectrum 情形下首个奇异 Toeplitz 主子矩阵恰在次数层出现]\label{cor:xi-toeplitz-first-singular-principal-minor-simple-spectrum}
在定理 \ref{thm:xi-toeplitz-rank-stabilizes-to-distinct-unitcircle-roots} 的设定下，若 \(\widetilde\Phi_L\) 的全部根均为单根，则 \(s=D\)。于是
\[
\mathbf{T}_N\succ 0
\qquad(0\le N\le D-1),
\]
而
\[
\mathbf{T}_N\ \text{奇异}
\qquad(N\ge D).
\]
特别地，第一块奇异 Toeplitz 主子矩阵恰在阶数 \(D+1\) 处出现。
\end{corollary}
\begin{proof}
若全部根均为单根，则不同根数 \(s\) 恰等于总次数 \(D\)。由定理 \ref{thm:xi-toeplitz-rank-stabilizes-to-distinct-unitcircle-roots}，
\[
\operatorname{rank}\mathbf{T}_N=\min(N+1,D).
\]
对 \(0\le N\le D-1\)，有 \(\operatorname{rank}\mathbf{T}_N=N+1\)，而 \(\mathbf{T}_N\) 又由定理 \ref{thm:xi-hilbert-polya-ninefold-equivalence-doubled-selfadjoint-compression} 的第 \((6)\) 条知为半正定，故必为正定。

对 \(N\ge D\)，有 \(\operatorname{rank}\mathbf{T}_N=D<N+1\)，故 \(\mathbf{T}_N\) 必奇异。于是首个奇异主子矩阵恰是大小 \(D+1\) 的那一块。证毕。
\end{proof}
\leanverified{paper\_xi\_toeplitz\_first\_singular\_principal\_minor\_simple\_spectrum}

\begin{conjecture}[prime-phase 表示诱导的 Bohr 几乎周期零点几何]\label{conj:xi-prime-phase-bohr-almost-periodic-zero-geometry}
沿用猜想 \ref{conj:xi-unitary-slice-arithmetic-singular-ring-representation} 的记号，设存在与 unitary slice 兼容的连续酉表示
\[
\rho_L:\mathfrak{R}_{S(L)}\to U(D_L),
\qquad
U_L(2t\log L)=\rho_L\bigl(\tau_L^{\mathrm{prime}}(t)\bigr).
\]
定义
\[
f_L(t):=\det\!\bigl(I-\rho_L(\tau_L^{\mathrm{prime}}(t))\bigr).
\]
则 \(f_L\) 应为一个实解析的 Bohr 几乎周期函数，其频率模满足
\[
\operatorname{Mod}(f_L)
=
\mathbf Z\text{-span}\bigl\{\,2v_p(L)\log p:\ p\mid L\,\bigr\}.
\]
从而，若 \(f_L\not\equiv 0\) 且零集非空，则由该表示在 unitary slice 上给出的临界线零点集应当是一个相对稠密的 Bohr 几乎周期集合；其统计性质应由 prime-phase torus \(\TT^{S(L)}\) 上行列式零除子的 Haar 几何完全控制。
\end{conjecture}

\noindent
由此，自对偶 completion 所固定的 unitary slice 不再只是一个单一边界约束，而在区间根、单位圆根、Cayley 超双曲、Carath\'eodory/Toeplitz/Herglotz 正性、CMV 酉谱、自伴 Joukowski 压缩与素相位环流之间形成同一台有限型 Hilbert--P{\'o}lya 机器：临界线条件既可由有限维谱实现完全读出，也可由素因子环面上的无共振流作动力学组织，而 Toeplitz 证书的秩冻结则把这一结构进一步压缩为有限原子相位数的显式计数。

\endinput


\paragraph{完成化截面的实根--Jacobi 证书、条带梳齿与有限多尺度去周期化}
在有限型 Hilbert--P{\'o}lya 的九重等价、实越界根双塔、条带禁区窗口与双底去周期化已经建立之后，还可把完成化截面 \(\Xi_{\Delta,L}\) 的实根几何进一步压缩为下列五条彼此衔接的后果。它们分别给出全部零点塔的显式分类、开临界条带无离线零点的 Jacobi 四次半定判据、条带离线零点的精确密度常数、有限多尺度叠加下的完全去周期化，以及禁区两端的 sharp onset law。

\begin{theorem}[完成化实根诱导的零点塔完全分类]\label{thm:xi-time-part28b-completed-realroot-zero-tower-classification}
\leanverified{paper\_xi\_time\_part28b\_completed\_realroot\_zero\_tower\_classification}
沿用命题 \ref{prop:xi-rh-interval-criterion} 的记号，设
\[
\Xi_{\Delta,L}(s)=P_L\!\bigl(S_L(s)\bigr),
\qquad
S_L(s):=L^{s-\frac12}+L^{\frac12-s},
\]
并假设 \(P_L\in\RR[S]\) 的全部根均为实数。取 \(S_0\) 为 \(P_L\) 的一个实根，则由该根诱导的零点集合可按下列五类完全分类：
\begin{enumerate}
  \item 若 \(S_0=2\cos\theta\in(-2,2)\)（\(\theta\in(0,\pi)\)），则对应两条临界线梳齿
  \[
  s=\frac12+\frac{i}{\log L}\bigl(\pm\theta+2\pi k\bigr),
  \qquad k\in\ZZ.
  \]
  \item 若 \(S_0=2\)，则对应单条临界线梳齿
  \[
  s=\frac12+\frac{2\pi i k}{\log L},
  \qquad k\in\ZZ.
  \]
  \item 若 \(S_0=-2\)，则对应单条半格偏移的临界线梳齿
  \[
  s=\frac12+\frac{(2k+1)\pi i}{\log L},
  \qquad k\in\ZZ.
  \]
  \item 若 \(S_0>2\)，并记
  \[
  \sigma_L(S_0):=\frac{\operatorname{arcosh}(S_0/2)}{\log L}>0,
  \]
  则对应两条对称竖线梳齿
  \[
  s=\frac12\pm \sigma_L(S_0)+\frac{2\pi i k}{\log L},
  \qquad k\in\ZZ.
  \]
  \item 若 \(S_0<-2\)，并记
  \[
  \sigma_L(S_0):=\frac{\operatorname{arcosh}(|S_0|/2)}{\log L}>0,
  \]
  则对应两条带半周期偏移的对称竖线梳齿
  \[
  s=\frac12\pm \sigma_L(S_0)+\frac{(2k+1)\pi i}{\log L},
  \qquad k\in\ZZ.
  \]
\end{enumerate}
这五类给出了 \(\Xi_{\Delta,L}\) 的全部零点作为点集的完整分类。更精确地，若 \(S_0\neq\pm2\) 是 \(P_L\) 的 \(m\) 重根，则上述每个零点都以重数 \(m\) 出现；若 \(S_0=\pm2\) 是 \(m\) 重根，则相应梳齿上的每个零点都以重数 \(2m\) 出现。
\end{theorem}
\begin{proof}
令
\[
y:=L^{s-\frac12},
\]
则
\[
\Xi_{\Delta,L}(s)=0
\iff
P_L(y+y^{-1})=0.
\]
因此对固定根 \(S_0\) 而言，零点条件化为
\[
y+y^{-1}=S_0.
\]
若 \(|S_0|<2\)，写 \(S_0=2\cos\theta\)（\(\theta\in(0,\pi)\)），则二次方程的两根为
\[
y=e^{\pm i\theta},
\]
从而
\[
L^{s-\frac12}=e^{\pm i\theta}
\iff
s=\frac12+\frac{i}{\log L}\bigl(\pm\theta+2\pi k\bigr).
\]
这给出第一类。

若 \(S_0=2\)，则
\[
y+y^{-1}-2=y^{-1}(y-1)^2,
\]
故零点满足 \(y=1\)，即
\[
s=\frac12+\frac{2\pi i k}{\log L}.
\]
同理，若 \(S_0=-2\)，则
\[
y+y^{-1}+2=y^{-1}(y+1)^2,
\]
故零点满足 \(y=-1\)，即
\[
s=\frac12+\frac{(2k+1)\pi i}{\log L}.
\]
这分别给出第二、三类。

若 \(S_0>2\)，写 \(S_0=2\cosh\tau\)（\(\tau>0\)），则两根为
\[
y_\pm=\frac{S_0\pm\sqrt{S_0^2-4}}{2}=e^{\pm\tau},
\qquad
\tau=\operatorname{arcosh}(S_0/2),
\]
从而
\[
s=\frac12\pm \frac{\tau}{\log L}+\frac{2\pi i k}{\log L}.
\]
若 \(S_0<-2\)，写 \(S_0=-2\cosh\tau\)（\(\tau>0\)），则两根为
\[
y_\pm=-e^{\pm\tau}=e^{\pm\tau+(2k+1)\pi i},
\]
于是
\[
s=\frac12\pm \frac{\tau}{\log L}+\frac{(2k+1)\pi i}{\log L}.
\]
这就得到第四、五类。

最后说明重数。若 \(S_0\neq\pm2\) 是 \(m\) 重根，则
\[
P_L(S)=(S-S_0)^m R(S),\qquad R(S_0)\neq 0.
\]
在相应零点处，
\[
\frac{d}{ds}S_L(s)=\log L\,(y-y^{-1})\neq 0,
\]
故 \(S_L(s)-S_0\) 为单零，从而 \(\Xi_{\Delta,L}(s)\) 的零点重数恰为 \(m\)。

若 \(S_0=2\)，则
\[
S_L(s)-2=y^{-1}(y-1)^2,
\]
而 \(y-1\) 在每个对应点上都是单零，故 \(S_L(s)-2\) 为二重零；于是 \(m\) 重根 \(S_0=2\) 产生 \(2m\) 重零点。对 \(S_0=-2\) 完全同理。证毕。
\end{proof}

\begin{theorem}[开临界条带无离线零点的 Jacobi 四次半定判据]\label{thm:xi-time-part28b-strip-jacobi-quartic-psd}
沿用命题 \ref{prop:xi-hilbert-polya-jacobi} 与命题 \ref{prop:xi-strip-forbidden-real} 的记号。设 \(P_L\) 全部根均为实数，并令
\[
\widetilde P_L(S)=\det(SI-J_L)
\]
为命题 \ref{prop:xi-hilbert-polya-jacobi} 给出的实对称三对角实现。记
\[
c_L:=\sqrt L+\frac{1}{\sqrt L}>2,
\qquad
\mathcal Q_L^{\mathrm{strip}}(x):=(x^2-4)(x^2-c_L^2).
\]
再记按重数计的禁区根数
\[
\nu_{\mathrm{forb}}(L)
:=
\#\Bigl(\mathrm{Roots}(P_L)\cap\bigl((2,c_L)\cup(-c_L,-2)\bigr)\Bigr),
\]
并记 \(\nu_-(H)\) 为 Hermitian 矩阵 \(H\) 的负惯性指标。则下述三条等价：
\begin{enumerate}
  \item \(\Xi_{\Delta,L}\) 在开临界条带
  \[
  \bigl\{s\in\CC:\ 0<\Re(s)<1,\ \Re(s)\neq\tfrac12\bigr\}
  \]
  中无零点；
  \item \(P_L\) 在禁区 \((2,c_L)\cup(-c_L,-2)\) 上无根；
  \item
  \[
  \mathcal Q_L^{\mathrm{strip}}(J_L)\succeq 0.
  \]
\end{enumerate}
并且有精确计数恒等式
\[
\boxed{\;
\nu_{\mathrm{forb}}(L)=\nu_-\!\bigl(\mathcal Q_L^{\mathrm{strip}}(J_L)\bigr).\;}
\]
\end{theorem}
\begin{proof}
由命题 \ref{prop:xi-strip-forbidden-real}，第 \((1)\) 条与第 \((2)\) 条等价。

再证第 \((2)\) 条与第 \((3)\) 条等价。由于 \(J_L\) 为实对称矩阵，谱定理给出
\[
\mathcal Q_L^{\mathrm{strip}}(J_L)\succeq 0
\iff
\mathcal Q_L^{\mathrm{strip}}(\lambda)\ge 0
\quad\text{对所有 }\lambda\in\mathrm{Spec}(J_L).
\]
而由命题 \ref{prop:xi-hilbert-polya-jacobi}，
\[
\mathrm{Spec}(J_L)=\mathrm{Roots}(P_L)
\]
按重数一致。于是只需分析标量函数 \(\mathcal Q_L^{\mathrm{strip}}\) 的符号。由定义，
\[
\mathcal Q_L^{\mathrm{strip}}(x)=(x^2-4)(x^2-c_L^2).
\]
当 \(|x|\le 2\) 时，两因子都非正，故 \(\mathcal Q_L^{\mathrm{strip}}(x)\ge 0\)；当 \(2<|x|<c_L\) 时，第一因子为正、第二因子为负，故 \(\mathcal Q_L^{\mathrm{strip}}(x)<0\)；当 \(|x|\ge c_L\) 时，两因子都非负，故 \(\mathcal Q_L^{\mathrm{strip}}(x)\ge 0\)。因此
\[
\mathcal Q_L^{\mathrm{strip}}(\lambda)\ge 0\ \text{对全部谱点成立}
\iff
\mathrm{Spec}(J_L)\cap\bigl((2,c_L)\cup(-c_L,-2)\bigr)=\varnothing,
\]
即得第 \((2)\) 条与第 \((3)\) 条的等价。

最后，\(\mathcal Q_L^{\mathrm{strip}}(J_L)\) 的特征值正是
\[
\mathcal Q_L^{\mathrm{strip}}(\lambda),
\qquad \lambda\in\mathrm{Spec}(J_L),
\]
按重数计。故其负惯性指标恰等于落在 \((2,c_L)\cup(-c_L,-2)\) 中的谱点个数，也就是 \(\nu_{\mathrm{forb}}(L)\)。证毕。
\end{proof}

\begin{theorem}[开临界条带离线零点的精确密度常数]\label{thm:xi-time-part28b-strip-offline-density-exact}
在定理 \ref{thm:xi-time-part28b-strip-jacobi-quartic-psd} 的设定下，记
\[
N_{\mathrm{off}}(T;L)
:=
\#\Bigl\{s:\ \Xi_{\Delta,L}(s)=0,\ 0<\Re(s)<1,\ \Re(s)\neq\tfrac12,\ |\Im(s)|\le T\Bigr\},
\]
按重数计。则当 \(T\to\infty\) 时有
\[
\boxed{\;
N_{\mathrm{off}}(T;L)
=
\frac{2\log L}{\pi}\,\nu_{\mathrm{forb}}(L)\,T+O(1).\;}
\]
更精确地，每个禁区根 \(S_0\in(2,c_L)\cup(-c_L,-2)\) 都恰贡献两条竖直算术梳齿，其虚部步长均为 \(2\pi/\log L\)。
\end{theorem}
\begin{proof}
由定理 \ref{thm:xi-time-part28b-completed-realroot-zero-tower-classification}，当且仅当 \(S_0\in(2,c_L)\cup(-c_L,-2)\) 时，相应零点位于开临界条带内但偏离临界线；并且每个这类根都贡献两条竖线梳齿。若 \(S_0\in(2,c_L)\)，则对应零点为
\[
s=\frac12\pm \sigma_L(S_0)+\frac{2\pi i k}{\log L},
\qquad k\in\ZZ,
\]
其中 \(0<\sigma_L(S_0)<\tfrac12\)。若 \(S_0\in(-c_L,-2)\)，则对应零点为
\[
s=\frac12\pm \sigma_L(S_0)+\frac{(2k+1)\pi i}{\log L},
\qquad k\in\ZZ,
\]
同样满足 \(0<\sigma_L(S_0)<\tfrac12\)。

对任意一条形如
\[
\Im(s)=\tau_0+\frac{2\pi k}{\log L},
\qquad k\in\ZZ
\]
的竖线梳齿，满足 \(|\Im(s)|\le T\) 的点数为
\[
\frac{\log L}{\pi}\,T+O(1).
\]
因此每个禁区根贡献
\[
2\cdot \frac{\log L}{\pi}\,T+O(1)
\]
个条带内离线零点。对全部 \(\nu_{\mathrm{forb}}(L)\) 个禁区根按重数求和，即得
\[
N_{\mathrm{off}}(T;L)
=
\frac{2\log L}{\pi}\,\nu_{\mathrm{forb}}(L)\,T+O(1).
\]
证毕。
\end{proof}

\leanverified{paper\_xi\_time\_part28b\_finite\_multiscale\_complete\_deperiodization}
\begin{theorem}[有限多尺度完成化的完全去周期化]\label{thm:xi-time-part28b-finite-multiscale-complete-deperiodization}
定理 \ref{thm:xi-two-base-deperiodization-no-imag-period} 的双底去周期化结论可推广到任意有限多尺度。具体而言，设 \(r\ge 2\)，取 \(L_1,\dots,L_r>1\) 满足
\[
\frac{\log L_i}{\log L_j}\notin\QQ
\qquad (i\neq j).
\]
对每个 \(1\le j\le r\)，设 \(g_j\) 为一变量亚纯函数，并满足“无非平凡单位根乘法不变性”：
\[
g_j(\omega z)\equiv g_j(z),\ \omega^n=1
\quad\Longrightarrow\quad
\omega=1.
\]
定义有限多尺度完成化
\[
F(s):=\sum_{j=1}^{r} g_j(L_j^{-s}).
\]
则若存在 \(T\in i\RR\) 使
\[
F(s+T)\equiv F(s),
\]
则必有 \(T=0\)。
\end{theorem}
\begin{proof}
反设存在非零纯虚周期 \(T\in i\RR\setminus\{0\}\)。令
\[
\omega_j:=e^{-T\log L_j}\in\TT,
\qquad 1\le j\le r.
\]
则恒等式
\[
F(s+T)-F(s)
=
\sum_{j=1}^{r}\bigl(g_j(\omega_j L_j^{-s})-g_j(L_j^{-s})\bigr)\equiv 0
\]
成立。取 \(\sigma_0\) 足够大，使得对每个 \(j\) 而言，当 \(\Re(s)\ge \sigma_0\) 时，\(L_j^{-s}\) 都落在 \(g_j\) 某个穿孔圆盘的 Laurent 展开域内。于是可写
\[
g_j(z)=\sum_{n\ge -M_j} a_{j,n}z^n,
\qquad 0<|z|<r_j.
\]
代入得
\[
0
=
\sum_{j=1}^{r}\sum_{n\ge -M_j}
a_{j,n}(\omega_j^n-1)L_j^{-ns}.
\]
由于对 \(n\neq 0\)，所有实数 \(n\log L_j\) 两两不同，否则将推出某个 \(\log L_i/\log L_j\in\QQ\)，与假设矛盾，故上式是由互异指数构成的 Dirichlet 型恒等式。由唯一性可知，对每个 \(j,n\) 都有
\[
a_{j,n}(\omega_j^n-1)=0.
\]
因此在各自的 Laurent 展开域内
\[
g_j(\omega_j z)=g_j(z).
\]
再由亚纯延拓得该恒等式全局成立。由“无非平凡单位根乘法不变性”假设，遂得
\[
\omega_j=1
\qquad (1\le j\le r).
\]
于是
\[
T\in \frac{2\pi i}{\log L_j}\ZZ
\qquad (1\le j\le r).
\]
任取 \(i\neq j\)，若 \(T\neq 0\)，则上式将推出 \(\log L_i/\log L_j\in\QQ\)，与假设矛盾。故只能有 \(T=0\)。证毕。
\end{proof}

\begin{theorem}[条带禁区两端的 onset law]\label{thm:xi-time-part28b-strip-onset-law}
沿用定理 \ref{thm:xi-time-part28b-strip-jacobi-quartic-psd} 的记号，令
\[
c_L:=\sqrt L+\frac{1}{\sqrt L},
\qquad
\sigma_L(S):=\frac{\operatorname{arcosh}(S/2)}{\log L}
\qquad (2<S<c_L).
\]
则成立两条端点渐近：
\begin{enumerate}
  \item 当 \(\varepsilon\downarrow 0\) 时，
  \[
  \boxed{\;
  \sigma_L(2+\varepsilon)
  =
  \frac{\sqrt{\varepsilon}}{\log L}
  +O(\varepsilon^{3/2}).\;}
  \]
  \item 当 \(\delta\downarrow 0\) 时，
  \[
  \boxed{\;
  \frac12-\sigma_L(c_L-\delta)
  =
  \frac{\delta}{2(\log L)\sinh\!\bigl(\frac12\log L\bigr)}
  +O(\delta^2).\;}
  \]
\end{enumerate}
因此，实根刚越过 \(2\) 时，相应离线竖线以平方根速度离开临界线；而当实根逼近禁区端点 \(c_L\) 时，相应竖线以线性速度贴近条带边界 \(\Re(s)=0,1\)。对负侧禁区 \((-c_L,-2)\) 完全同样的结论在 \(|S|\) 变量下成立。
\end{theorem}
\begin{proof}
第一条直接来自标准展开
\[
\operatorname{arcosh}(1+x)=\sqrt{2x}+O(x^{3/2})
\qquad(x\downarrow 0).
\]
取 \(x=\varepsilon/2\)，便得到
\[
\operatorname{arcosh}\!\Bigl(\frac{2+\varepsilon}{2}\Bigr)
=
\sqrt{\varepsilon}+O(\varepsilon^{3/2}),
\]
再除以 \(\log L\) 即得所示公式。

第二条中令
\[
\lambda:=\frac12\log L,
\qquad
c_L=2\cosh\lambda.
\]
对 \(S=c_L-\delta\)，设
\[
\operatorname{arcosh}(S/2)=\lambda-\eta.
\]
则
\[
2\cosh(\lambda-\eta)=2\cosh\lambda-\delta.
\]
在 \(\eta=0\) 处作 Taylor 展开，
\[
2\cosh(\lambda-\eta)
=
2\cosh\lambda-2\eta\sinh\lambda+O(\eta^2).
\]
比较系数得
\[
\delta=2\eta\sinh\lambda+O(\eta^2),
\qquad
\eta=\frac{\delta}{2\sinh\lambda}+O(\delta^2).
\]
而
\[
\sigma_L(c_L-\delta)
=
\frac{\lambda-\eta}{\log L}
=
\frac12-\frac{\eta}{\log L},
\]
故
\[
\frac12-\sigma_L(c_L-\delta)
=
\frac{\delta}{2(\log L)\sinh\lambda}+O(\delta^2)
=
\frac{\delta}{2(\log L)\sinh\!\bigl(\frac12\log L\bigr)}+O(\delta^2).
\]
证毕。
\end{proof}

\noindent
由此，完成化截面的实根结构不再只是“区间根/单位圆根”式的静态判据，而被进一步组织为一条显式动力几何链：一方面，全部零点塔可由实根位置完全分类，自伴三对角证书又把开临界条带的离线排斥压缩为单个四次多项式在 \(J_L\) 上的半定性；另一方面，条带内离线零点的密度常数与禁区根数严格相等，而一旦把单尺度梳齿推广到有限多个两两对数不可公度的底，严格虚周期即被彻底摧毁。于是，实根、Jacobi、梳齿与去周期化四者在同一完成化接口上获得了闭合。

\endinput


\paragraph{CMV 前缀刚性、端点通量稳定泛函与 Artin 条带障碍算子的统一证书链}
在有限型 Hilbert--P\'olya 的 CMV 规范型、有限 CMV 证书的点态稳定化、有限窗谱矩冻结、端点通量的有限证书指数逼近，以及 Artin 通道化的完成化/条带判据已经分别建立之后，还可把这几条主线继续压缩为如下五条结果。它们分别刻画端点通量在 Verblunsky 参数空间上的指数圆柱连续性、点态稳定化族的端点通量极限与显式 Cauchy 模、Verblunsky 前缀的严格谱矩刚性、Artin 通道条带离线零点计数的严格可加性，以及 Jacobi--CMV 双规范在多通道直和下的单一障碍算子证书。

\begin{theorem}[端点通量在 Verblunsky 前缀空间上的指数圆柱连续性]\label{thm:xi-time-part28c-endpoint-flux-exponential-cylinder-continuity}
\leanverified{paper\_xi\_time\_part28c\_endpoint\_flux\_exponential\_cylinder\_continuity}
沿用命题 \ref{con:xi-endpoint-flux-computable-by-finite-cmv} 的记号。设 \(B_\alpha,B_\beta\) 为该命题适用范围内的两个 Blaschke 对象，其 Verblunsky 系列分别记为
\[
\alpha=(\alpha_n)_{n\ge 0},
\qquad
\beta=(\beta_n)_{n\ge 0}.
\]
若对某个 \(N\ge 1\) 有
\[
\alpha_j=\beta_j
\qquad (0\le j\le N-1),
\]
则存在常数 \(C>0\)、\(q\in(0,1)\)（只依赖于 \(B_\alpha,B_\beta\)）使
\[
\boxed{\;
\bigl|\Phi_{-1}(B_\alpha)-\Phi_{-1}(B_\beta)\bigr|
\le 2Cq^N.\;}
\]
\end{theorem}
\begin{proof}
对 \(B_\alpha\) 与 \(B_\beta\) 分别应用命题 \ref{con:xi-endpoint-flux-computable-by-finite-cmv}，可得同一显式有限前缀函数族 \(\mathcal F_N\) 以及常数 \(C_\alpha>0,q_\alpha\in(0,1)\)、\(C_\beta>0,q_\beta\in(0,1)\)，满足
\[
\bigl|\Phi_{-1}(B_\alpha)-\mathcal F_N(\alpha_0,\dots,\alpha_{N-1})\bigr|
\le C_\alpha q_\alpha^N,
\]
\[
\bigl|\Phi_{-1}(B_\beta)-\mathcal F_N(\beta_0,\dots,\beta_{N-1})\bigr|
\le C_\beta q_\beta^N.
\]
由于两组 Verblunsky 前缀完全一致，
\[
\mathcal F_N(\alpha_0,\dots,\alpha_{N-1})
=
\mathcal F_N(\beta_0,\dots,\beta_{N-1}).
\]
令
\[
C:=\max\{C_\alpha,C_\beta\},
\qquad
q:=\max\{q_\alpha,q_\beta\}\in(0,1),
\]
并应用三角不等式，即得
\[
\bigl|\Phi_{-1}(B_\alpha)-\Phi_{-1}(B_\beta)\bigr|
\le C_\alpha q_\alpha^N+C_\beta q_\beta^N
\le 2Cq^N.
\]
证毕。
\end{proof}

\begin{corollary}[点态稳定化族的端点通量极限与显式 Cauchy 模]\label{cor:xi-time-part28c-stabilized-endpoint-flux-cauchy-modulus}
沿用定理 \ref{thm:xi-cmv-stabilization-infinite} 与结论 \ref{con:xi-cmv-finite-window-freeze} 的设定。设有限 CMV 证书族 \(\{U_L^{\mathrm{CMV}}\}_L\) 的 Verblunsky 系数满足：对每个固定 \(n\ge 0\)，当 \(L\) 充分大时
\[
\alpha_n(L)=\alpha_n
\]
稳定为与 \(L\) 无关的值。再设极限序列 \((\alpha_n)_{n\ge 0}\) 对应的对象 \(B_\infty\) 落在命题 \ref{con:xi-endpoint-flux-computable-by-finite-cmv} 的适用范围内。定义
\[
f_N:=\mathcal F_N(\alpha_0,\dots,\alpha_{N-1}),
\]
其中 \(\mathcal F_N\) 为命题 \ref{con:xi-endpoint-flux-computable-by-finite-cmv} 中的显式有限前缀函数。则：
\begin{enumerate}
  \item 对每个固定 \(N\ge 1\)，存在 \(L_\ast(N)\) 使得当 \(L\ge L_\ast(N)\) 时，
  \[
  \mathcal F_N(\alpha_0(L),\dots,\alpha_{N-1}(L))=f_N.
  \]
  \item 存在常数 \(C>0\)、\(q\in(0,1)\) 使
  \[
  \boxed{\;
  |f_N-\Phi_{-1}(B_\infty)|\le Cq^N.\;}
  \]
  \item 序列 \((f_N)_{N\ge 1}\) 具有显式 Cauchy 模：
  \[
  \boxed{\;
  |f_{N+1}-f_N|\le C(1+q)q^N.\;}
  \]
\end{enumerate}
\end{corollary}
\begin{proof}
对第 \(1\) 条，只需取
\[
L_\ast(N):=\max_{0\le j\le N-1}L_0(j),
\]
其中 \(L_0(j)\) 是使 \(\alpha_j(L)=\alpha_j\) 稳定的阈值。则对一切 \(L\ge L_\ast(N)\)，前 \(N\) 个 Verblunsky 系数完全固定，故有限前缀函数 \(\mathcal F_N\) 的值亦固定。

对第 \(2\) 条，把命题 \ref{con:xi-endpoint-flux-computable-by-finite-cmv} 直接应用于极限对象 \(B_\infty\) 即得。

对第 \(3\) 条，由三角不等式，
\[
|f_{N+1}-f_N|
\le |f_{N+1}-\Phi_{-1}(B_\infty)|
+|\Phi_{-1}(B_\infty)-f_N|
\le Cq^{N+1}+Cq^N,
\]
即
\[
|f_{N+1}-f_N|\le C(1+q)q^N.
\]
证毕。
\end{proof}

\begin{theorem}[Verblunsky 前缀的严格谱矩刚性与 Carath\'eodory 前缀刚性]\label{thm:xi-time-part28c-verblunsky-prefix-exact-moment-rigidity}
设 \(U_\alpha,U_\beta\) 为两组半无限 CMV 幺正算子，其 Verblunsky 参数分别为
\[
\alpha=(\alpha_n)_{n\ge 0},
\qquad
\beta=(\beta_n)_{n\ge 0},
\qquad
\alpha_n,\beta_n\in\mathbb D.
\]
若对某个整数 \(m\ge 0\) 有
\[
\alpha_j=\beta_j
\qquad (0\le j\le 2m+1),
\]
则对一切 \(0\le r\le m\) 都有
\[
U_\alpha^r e_0=U_\beta^r e_0,
\]
因而
\[
\boxed{\;
\langle e_0,U_\alpha^r e_0\rangle
=
\langle e_0,U_\beta^r e_0\rangle
\qquad (0\le r\le m).\;}
\]
等价地，相应 Carath\'eodory 函数
\[
F_\alpha(z):=\langle e_0,(U_\alpha+zI)(U_\alpha-zI)^{-1}e_0\rangle,
\qquad
F_\beta(z):=\langle e_0,(U_\beta+zI)(U_\beta-zI)^{-1}e_0\rangle
\]
在 \(z=0\) 处的前 \(m\) 阶 Taylor 系数完全一致。

若进一步 \(B_\alpha,B_\beta\) 为与这两列参数对应且落在命题 \ref{con:xi-endpoint-flux-computable-by-finite-cmv} 适用范围内的对象，则存在常数 \(C>0\)、\(q\in(0,1)\)（只依赖于 \(B_\alpha,B_\beta\)）使
\[
\boxed{\;
\bigl|\Phi_{-1}(B_\alpha)-\Phi_{-1}(B_\beta)\bigr|
\le 2Cq^{2m+2}.\;}
\]
\end{theorem}
\begin{proof}
CMV 矩阵是五对角带状算子。由定理 \ref{thm:xi-cmv-stabilization-infinite} 的证明可知：对每个 \(r\ge 0\)，向量 \(U^r e_0\) 的支撑包含在
\[
\mathrm{span}\{e_0,\dots,e_{2r}\},
\]
并且这段有限窗口上的矩阵元只依赖于 Verblunsky 前缀
\[
\alpha_0,\dots,\alpha_{2r+1}.
\]
因此当 \(0\le r\le m\) 时，假设中的前缀一致性保证 \(U_\alpha\) 与 \(U_\beta\) 在生成 \(U^r e_0\) 所需的全部有限窗口上完全一致，从而
\[
U_\alpha^r e_0=U_\beta^r e_0
\qquad (0\le r\le m).
\]
取与 \(e_0\) 的内积即得谱矩相等。

另一方面，对半无限 CMV 幺正算子的 Carath\'eodory 函数有幂级数展开
\[
F(z)=1+2\sum_{r\ge 1}\langle e_0,U^{-r}e_0\rangle z^r,
\qquad |z|<1.
\]
而 \(U\) 为酉算子时
\[
\langle e_0,U^{-r}e_0\rangle
=
\overline{\langle e_0,U^r e_0\rangle}.
\]
故前 \(m\) 阶 Taylor 系数也完全一致。

最后，把定理 \ref{thm:xi-time-part28c-endpoint-flux-exponential-cylinder-continuity} 应用于
\[
N=2m+2
\]
即得
\[
\bigl|\Phi_{-1}(B_\alpha)-\Phi_{-1}(B_\beta)\bigr|
\le 2Cq^{2m+2}.
\]
证毕。
\end{proof}

\begin{theorem}[Artin 通道条带离线零点计数的严格可加性]\label{thm:xi-time-part28c-artin-strip-offline-count-additivity}
\leanverified{paper\_xi\_time\_part28c\_artin\_strip\_offline\_count\_additivity}
沿用注 \ref{rem:xi-artin-channel-pipeline} 的通道化设定。设 \(G\) 为有限群，\(d_\rho:=\dim\rho\)，并定义
\[
\Delta_G(z,u):=\prod_{\rho\in\widehat G}\Delta_\rho(z,u)^{d_\rho},
\qquad
\Xi_{G,L}(s):=\prod_{\rho\in\widehat G}\Xi_{\rho,L}(s)^{d_\rho}.
\]
假设每个通道多项式
\[
P_{L,\rho}(S):=\widehat\Delta_\rho(L^{-1/2},S)\in\mathbb R[S]
\]
都具有全实根。记
\[
c_L:=\sqrt L+\frac{1}{\sqrt L},
\qquad
I_L:=(2,c_L)\cup(-c_L,-2),
\]
并按重数定义禁区根数
\[
\nu_\rho(L):=\#\bigl(\mathrm{Roots}(P_{L,\rho})\cap I_L\bigr).
\]
再定义总条带离线零点计数函数
\[
N_{G,L}^{\mathrm{off}}(T)
:=
\#\Bigl\{
s:\ \Xi_{G,L}(s)=0,\ 0<\Re(s)<1,\ \Re(s)\neq\tfrac12,\ |\Im s|\le T
\Bigr\},
\]
按重数计。则当 \(T\to\infty\) 时有
\[
\boxed{\;
N_{G,L}^{\mathrm{off}}(T)
=
\frac{2\log L}{\pi}
\Bigl(\sum_{\rho\in\widehat G} d_\rho\,\nu_\rho(L)\Bigr)T
+O(1).\;}
\]
特别地，
\[
\boxed{\;
\Xi_{G,L}\ \text{在开临界条带内无离线零点}
\iff
\sum_{\rho\in\widehat G} d_\rho\,\nu_\rho(L)=0.\;}
\]
由于各项均非负，上式又等价于对所有 \(\rho\in\widehat G\) 都有 \(\nu_\rho(L)=0\)。
\end{theorem}
\begin{proof}
由注 \ref{rem:xi-artin-channel-pipeline}，每个通道截面 \(\Xi_{\rho,L}\) 都具有与单通道完成化截面完全相同的形式
\[
\Xi_{\rho,L}(s)=P_{L,\rho}\!\bigl(S_L(s)\bigr),
\qquad
S_L(s)=L^{s-1/2}+L^{1/2-s}.
\]
又由命题 \ref{prop:xi-strip-forbidden-real}，\(\Xi_{\rho,L}\) 在开临界条带内出现离线零点，当且仅当 \(P_{L,\rho}\) 在禁区 \(I_L\) 内有根。对每个禁区根 \(S_0\in I_L\)，定理 \ref{thm:xi-time-part28b-completed-realroot-zero-tower-classification} 给出恰好两条竖直算术梳齿，其虚部步长均为 \(2\pi/\log L\)。因此该根在 \(|\Im s|\le T\) 中贡献
\[
2\cdot \frac{\log L}{\pi}\,T+O(1)
\]
个离线零点；对通道 \(\rho\) 的全部禁区根按重数求和，得到
\[
N_{\rho,L}^{\mathrm{off}}(T)
=
\frac{2\log L}{\pi}\,\nu_\rho(L)\,T+O(1).
\]

现在考察全局乘积
\[
\Xi_{G,L}(s)=\prod_{\rho\in\widehat G}\Xi_{\rho,L}(s)^{d_\rho}.
\]
零点除子的乘法规则表明：若某个点 \(s_0\) 在第 \(\rho\) 个通道中以重数 \(m\) 出现，则它在全局乘积中的贡献重数为 \(d_\rho m\)；若不同通道在同一点重合，则总重数按通道求和。因此总计数满足
\[
N_{G,L}^{\mathrm{off}}(T)
=
\sum_{\rho\in\widehat G} d_\rho\,N_{\rho,L}^{\mathrm{off}}(T)
=
\frac{2\log L}{\pi}
\Bigl(\sum_{\rho\in\widehat G} d_\rho\,\nu_\rho(L)\Bigr)T
+O(1).
\]
特别地，开临界条带内无离线零点当且仅当右端主项系数为零；而所有 \(d_\rho\nu_\rho(L)\) 都非负，故这又等价于逐通道 \(\nu_\rho(L)=0\)。证毕。
\end{proof}

\begin{theorem}[Artin 通道的 Jacobi--CMV 双规范可平行直和，并给出单一障碍算子]\label{thm:xi-time-part28c-artin-jacobi-cmv-directsum-obstruction}
在定理 \ref{thm:xi-time-part28c-artin-strip-offline-count-additivity} 的设定下，对每个 \(\rho\in\widehat G\) 令
\[
P_{L,\rho}(S)=\det(SI-J_{L,\rho})
\]
为命题 \ref{prop:xi-hilbert-polya-jacobi} 的逐通道 Jacobi 三对角实现。再定义
\[
c_L:=\sqrt L+\frac{1}{\sqrt L},
\qquad
Q_L(x):=(x^2-4)(x^2-c_L^2),
\]
\[
\mathcal J_{G,L}
:=
\bigoplus_{\rho\in\widehat G}
\bigl(I_{d_\rho}\otimes J_{L,\rho}\bigr),
\qquad
\mathfrak O_{G,L}:=\sum_{\rho\in\widehat G} d_\rho\,\nu_\rho(L).
\]
则：
\begin{enumerate}
  \item 总通道障碍指数可由单一自伴算子读出：
  \[
  \boxed{\;
  \mathfrak O_{G,L}
  =
  \nu_-\!\bigl(Q_L(\mathcal J_{G,L})\bigr).\;}
  \]
  \item 因而
  \[
  \boxed{\;
  \Xi_{G,L}\ \text{在开临界条带内无离线零点}
  \iff
  Q_L(\mathcal J_{G,L})\succeq 0.\;}
  \]
  \item 若对每个通道 \(\rho\) 的有限 CMV 证书族 \(\{U_{\rho,L}^{\mathrm{CMV}}\}_L\) 都满足定理 \ref{thm:xi-cmv-stabilization-infinite} 的点态稳定化假设，则唯一的无限维直和
  \[
  \mathcal U_{G,\infty}
  :=
  \bigoplus_{\rho\in\widehat G}
  \bigl(I_{d_\rho}\otimes U_{\rho,\infty}\bigr)
  \]
  存在，并且其任意有限窗口上的有限阶谱矩在某一有限级别处全部冻结。若进一步每个通道对象 \(B_\rho\) 都落在命题 \ref{con:xi-endpoint-flux-computable-by-finite-cmv} 的适用范围内，并记
  \[
  \Phi_{G,-1}
  :=
  \Phi_{-1}\!\Bigl(\prod_{\rho\in\widehat G} B_\rho^{d_\rho}\Bigr)
  =
  \sum_{\rho\in\widehat G} d_\rho\,\Phi_{-1}(B_\rho),
  \]
  则存在显式有限前缀近似
  \[
  \Phi_{G,-1}^{(N)}:=\sum_{\rho\in\widehat G} d_\rho\,f_{\rho,N}
  \]
  及常数 \(C_G>0,\ q_G\in(0,1)\) 使
  \[
  |\Phi_{G,-1}^{(N)}-\Phi_{G,-1}|\le C_G q_G^N,
  \qquad
  |\Phi_{G,-1}^{(N+1)}-\Phi_{G,-1}^{(N)}|
  \le C_G(1+q_G)q_G^N.
  \]
\end{enumerate}
\end{theorem}
\begin{proof}
对每个通道 \(\rho\)，定理 \ref{thm:xi-time-part28b-strip-jacobi-quartic-psd} 的逐通道版本给出
\[
\nu_\rho(L)=\nu_-\!\bigl(Q_L(J_{L,\rho})\bigr).
\]
另一方面，
\[
Q_L(\mathcal J_{G,L})
=
\bigoplus_{\rho\in\widehat G}
\bigl(I_{d_\rho}\otimes Q_L(J_{L,\rho})\bigr).
\]
负惯性指数对直和严格可加，且张量上 \(I_{d_\rho}\) 只把每个负方向复制 \(d_\rho\) 次，因此
\[
\nu_-\!\bigl(Q_L(\mathcal J_{G,L})\bigr)
=
\sum_{\rho\in\widehat G} d_\rho\,\nu_-\!\bigl(Q_L(J_{L,\rho})\bigr)
=
\sum_{\rho\in\widehat G} d_\rho\,\nu_\rho(L)
=
\mathfrak O_{G,L}.
\]
这证明了第 \(1\) 条。

第 \(2\) 条由第 \(1\) 条与定理 \ref{thm:xi-time-part28c-artin-strip-offline-count-additivity} 的零点判据立即得到。

对第 \(3\) 条，逐通道应用定理 \ref{thm:xi-cmv-stabilization-infinite} 可得每个 \(U_{\rho,\infty}\) 的唯一存在；再由结论 \ref{con:xi-cmv-finite-window-freeze}，对每个固定通道、固定窗口与固定矩阶，有限窗谱矩都在某一有限级别处完全冻结。由于 \(\widehat G\) 为有限集，有限多个通道的阈值取最大值即可，于是直和算子 \(\mathcal U_{G,\infty}\) 存在且具有同样的有限窗冻结性质。

若再给定通道对象 \(B_\rho\) 并对每个 \(\rho\) 应用推论 \ref{cor:xi-time-part28c-stabilized-endpoint-flux-cauchy-modulus}，则存在近似 \(f_{\rho,N}\) 以及常数 \(C_\rho>0,\ q_\rho\in(0,1)\) 使
\[
|f_{\rho,N}-\Phi_{-1}(B_\rho)|\le C_\rho q_\rho^N,
\qquad
|f_{\rho,N+1}-f_{\rho,N}|\le C_\rho(1+q_\rho)q_\rho^N.
\]
又由结论 \ref{con:xi-endpoint-flux-artin-additivity}，
\[
\Phi_{G,-1}
=
\sum_{\rho\in\widehat G} d_\rho\,\Phi_{-1}(B_\rho).
\]
令
\[
C_G:=\sum_{\rho\in\widehat G} d_\rho C_\rho,
\qquad
q_G:=\max_{\rho\in\widehat G} q_\rho\in(0,1),
\]
则
\[
|\Phi_{G,-1}^{(N)}-\Phi_{G,-1}|
\le
\sum_{\rho\in\widehat G} d_\rho\,|f_{\rho,N}-\Phi_{-1}(B_\rho)|
\le
\sum_{\rho\in\widehat G} d_\rho C_\rho q_\rho^N
\le C_G q_G^N,
\]
并且
\[
|\Phi_{G,-1}^{(N+1)}-\Phi_{G,-1}^{(N)}|
\le
\sum_{\rho\in\widehat G} d_\rho\,|f_{\rho,N+1}-f_{\rho,N}|
\le
C_G(1+q_G)q_G^N.
\]
证毕。
\end{proof}

\noindent
以上五条结果把有限 CMV 证书、端点通量与 Artin 条带障碍收束为同一条有限可审计主线：Verblunsky 前缀一方面以严格谱矩刚性冻结低阶谱矩并控制端点通量的指数圆柱模，另一方面经 Artin 通道化与 Jacobi 规范化把开临界条带内的离线障碍压缩为单一四次多项式作用在自伴直和算子上的负惯性指数；而点态稳定化又把这些有限前缀证书稳定地推送到无限维 CMV 极限，从而在不离开既有证明链语言的前提下，把“有限证书—端点通量—条带障碍”三者组织为同一组可复核对象。

\endinput


\paragraph{有限样本可检出的谱指纹与瞬态深度：Edgeworth 常数、幂零指数与 Witt--Dwork 继承}

\begin{conclusion}[规范差涨落的三阶谱指纹：Edgeworth 系数给出普适负偏度常数并锁定 Jordan 缺陷]\label{con:xi-gauge-anomaly-third-order-spectral-fingerprint}
对规范差过程 \(G_m\) 的中心极限定理尺度归一化
\[
Z_m:=\frac{G_m-\frac49 m}{\sqrt{\frac{118}{243}m}},
\]
其一阶 Edgeworth 修正项系数具有完全闭式
\[
\boxed{\;
\gamma_1^{(G)}
=\frac{P_G^{(3)}(0)}{(P_G''(0))^{3/2}}
=\frac{-\frac{1174}{2187}}{\left(\frac{118}{243}\right)^{3/2}}\;<\;0\;}
\]
从而偏度为严格负且由有限维谱/Jordan 数据唯一锁定。
该三阶常数把 \((-1/2)\) 的 Jordan 缺陷从结构叙事提升为有限样本即可检出的数值指纹。
\end{conclusion}
\begin{proof}[证明要点]
三阶压力导数闭式见命题 \ref{prop:fold-gauge-anomaly-pressure}；
Edgeworth 一阶封口与系数表达见推论 \ref{cor:fold-gauge-anomaly-edgeworth1}。
\end{proof}

\begin{conclusion}[幂零指数即瞬态记忆深度：核维数链提供严格上界并与复位词长度同一]\label{con:xi-nilpotent-index-transient-memory-depth}
对真实输入 \(40\) 状态核的全矩阵 \(M\)，核维数链 \(\dim\ker(M^k)\) 的稳定点给出 \(0\) 特征值处最大 Jordan 块大小 \(r_{\mathrm{full}}=5\)；
因此对任意固定向量 \(v,w\)，序列 \(f_m=v^\ast M^m w\) 在 \(m\ge r_{\mathrm{full}}\) 后与 \(w\) 在 \(0\)-谱瞬态子空间上的分量无关（瞬态贡献被幂零部分彻底湮灭）。
同理本质核的对应指数为 \(3\)。
该不变量把在线延迟与复位词长度提升为可由线性代数链条一键复算的硬上界。
\end{conclusion}
\begin{proof}[证明要点]
见命题 \ref{prop:real-input-40-nilpotent-index} 与推论 \ref{cor:real-input-40-nilpotent-memory-depth}（核维数链与 Jordan 块计数的可复现输出见注 \ref{rem:real-input-40-nilpotent-index-repro}）。
\end{proof}

\begin{conclusion}[Big-Witt/Dwork 护栏的零成本继承：多变量升级不改变整除与同余结构]\label{con:xi-big-witt-dwork-guardrail-multivariable}
二维/三维加权矩阵 \(B(u,v)\)、\(B(u,v,w)\) 的迹多项式 \(a_n\) 仍满足 necklace 整除律与 Dwork 型同余的逐字推广：
\[
\sum_{d\mid n}\mu(d)\,a_{n/d}(u^d,v^d)\in n\ZZ[u,v],
\qquad
a_{p^k}(u,v)\equiv a_{p^{k-1}}(u^p,v^p)\pmod{p^k},
\]
三维情形同理。
因此 primitive--periodic 的 Witt Euler 乘积与 \(p\)-幂同余在多维能量分辨下保持不变，提供一条实现自审计护栏：仅需追踪迹序列即可完成整除/同余一致性检查，而无需重新进入矩阵幂计算的高成本域。
\end{conclusion}
\begin{proof}[证明要点]
Witt Euler 乘积与 Dwork 同余的单变量模板见结论 \ref{con:xi-terminal-big-witt-functoriality} 与推论 \ref{cor:witt-dwork-congruence}；
多变量推广由“逐变量代入保持整系数与模同余”的函子性与 ghost 坐标的逐项定义直接推出（亦与 \S\ref{subsec:operator-cyclic-traceclass-witt-rigidity} 的迹序列口径一致）。
\end{proof}


\paragraph{稳定负子空间、端点谱迹与时间纤维：Toeplitz 否证的可编译接口}
本段将“有限否证器”的负谱结构进一步压缩为可复算的模型空间截断像与端点读出不变量，并把 window-\(6\) 的最小热化闭合、Witt--Fourier 双桥与同步核中心六重聚合纳入同一套自动编号的结论口径。

\begin{conclusion}[Toeplitz 否证器的可编译性：稳定负子空间是模型空间的截断像]\label{con:xi-toeplitz-compiled-model-space-truncation}
在有限 Blaschke 缺陷设定下，令 \(\kappa:=\deg(B_{F_{\zeta}})=\dim K_{B_{F_{\zeta}}}\)。
则存在 \(N_0\) 与一族随 \(N\) 稳定的显式线性嵌入
\[
\iota_N:K_{B_{F_{\zeta}}}\hookrightarrow \CC^{N+1}\qquad(N\ge N_0)
\]
（例如取 Hardy 基 \(\{1,w,w^2,\dots\}\) 的前 \(N{+}1\) 个 Taylor 系数坐标），使 Toeplitz 二次型
\(
x\mapsto x^\ast\mathbf{T}_N x
\)
在 \(\iota_N(K_{B_{F_{\zeta}}})\) 上的限制具有\textbf{恰好 \(\kappa\) 维负定部分}，且该负定部分在 \(N\to\infty\) 的稳定区间内保持不变。
因此一旦观测到 \(\mathbf{T}_N\) 的稳定负子空间，即可在有限维上实现一个与压缩移位 \(A_{B_{F_\zeta}}\) 的点谱残差轴同型的“残差轴”读出，并以其点谱数据反演圆盘内缺陷点（按重数计）。
\end{conclusion}
\begin{proof}[证明要点]
稳定嵌入与“负子空间 \(=\) 模型空间截断像”的同型性见定理 \ref{thm:xi-toeplitz-witness-generated-by-model-space}；
点谱残差轴与 Blaschke 零点的对应见定理 \ref{thm:blaschke-point-spectrum-correspondence}。
结合谱读出接口（例如定理 \ref{thm:xi-resonance-pencil-reconstruction} 及结论 \ref{con:xi-finite-falsifier-offline-fiber-readout}）即可把稳定负子空间提升为有限维可审计的缺陷位置反演。
\end{proof}

\begin{conclusion}[失败见证的极小生成律：否证由点谱残差轴强制生成]\label{con:xi-toeplitz-witness-minimal-generator-axis}
沿用上条结论的有限缺陷设定。
当 \(N\) 充分大时，Toeplitz 层级的一切负方向都可由点谱残差轴
\(
\mathcal H_{\mathrm{pp}}(A_{B_{F_\zeta}})=K_{B_{F_\zeta}}
\)
的点模态经 \(\iota_N\) 的像生成；从而 Toeplitz 否证的“极小生成维数”严格等于
\(
\deg(B_{F_\zeta})=\kappa
\)。
因此有限否证器是可审计的有限维生成对象，而非数值不稳定的弥散副产物。
\end{conclusion}
\begin{proof}[证明要点]
负惯性指标在高阶截断上的稳定性见定理 \ref{thm:xi-toeplitz-inertia-stabilizes-to-kappa}；
点模态生成性与极小维数结论由定理 \ref{thm:xi-toeplitz-witness-generated-by-model-space} 直接给出。
\end{proof}

\begin{conclusion}[端点通量的算子化同一：Poisson--Busemann 权重等于 Julia 指标与端点评价算子的迹]\label{con:xi-endpoint-flux-operator-trace-identity}
设 \(B\) 为有限 Blaschke 乘积，零点多重集为 \(\{a_k\}_{k=1}^{\kappa}\subset\mathbb{D}\)，并记端点 Poisson 权重
\[
P_a(-1)=\frac{1-|a|^2}{|1+a|^2}.
\]
定义端点通量（吸收系数）
\[
\Phi_{-1}(B):=\sum_{k=1}^{\kappa}P_{a_k}(-1),
\]
并令 \(K_B:=H^2\ominus BH^2\)、端点评价泛函 \(L_{-1}(f):=f(-1)\)。
则有严格恒等式
\[
\boxed{\ \Phi_{-1}(B)=\mathcal{J}_{-1}(B)=\mathrm{tr}\!\bigl(L_{-1}^\ast L_{-1}\bigr)\ },
\]
从而端点吸收强度被提升为“端点评价算子”的可谱化迹不变量，并与点模态边界耦合平方和完全等价。
\end{conclusion}
\begin{proof}[证明要点]
这是定理 \ref{thm:xi-endpoint-flux-pointmode-coupling} 的结论（并用到 \eqref{eq:app-poisson-kernel-endpoint-minus1} 的 Poisson 端点闭式与 Julia 指标接口）。
\end{proof}

\begin{conclusion}[Artin 分解下的端点通量可加性：通道归一化行列式因子的 Poisson 计数律]\label{con:xi-endpoint-flux-artin-additivity}
设 \(G\) 为有限群，并沿用命题 \ref{prop:ihara-artin-factorization} 的群扩张 Hashimoto 记号。
固定 \(u>0\)，记 \(\lambda_1(u):=\rho(B_{\mathbf 1}(u))\)，并对每个不可约表示 \(\varrho\in\widehat G\) 定义归一化行列式因子
\[
F_\varrho(t):=\det\!\Bigl(I-(t/\lambda_1(u))B_\varrho(u)\Bigr).
\]
若对每个 \(\varrho\)，\(F_\varrho\) 在单位圆上为 inner（等价地：其零点全部位于 \(\mathbb D\) 且边界模长恒为 \(1\)，从而 \(F_\varrho\)（除去单位模常数后）为有限 Blaschke 乘积），则端点通量满足严格表示论可加性：
\[
\boxed{\;
\Phi_{-1}\!\Bigl(\prod_{\varrho\in\widehat G} F_\varrho(t)^{\dim\varrho}\Bigr)
=
\sum_{\varrho\in\widehat G}(\dim\varrho)\,\Phi_{-1}(F_\varrho)\; }.
\]
在 \(G=S_4\) 且取 \(\varrho=\mathrm{sgn}\) 的情形，记
\[
R_{\det}(u,t):=\frac{F_{\mathrm{sgn}}(t)}{F_{\mathbf 1}(t)}.
\]
则可定义跨支端点通量差
\[
\Phi_{-1}(R_{\det}):=\Phi_{-1}(F_{\mathrm{sgn}})-\Phi_{-1}(F_{\mathbf 1}),
\]
其由零点 Poisson 权重完全确定，并可借由结论 \ref{con:xi-endpoint-flux-operator-trace-identity} 的迹表示提升为可谱化不变量。
\end{conclusion}
\begin{proof}[证明要点]
inner 假设把每个 \(F_\varrho\)（除去单位模常数后）识别为有限 Blaschke 乘积，其零点多重集按乘积并合且幂次引入重数因子。
由定义 \(\Phi_{-1}=\sum_k P_{a_k}(-1)\) 对乘积严格可加，立即得到所示恒等式；\(\Phi_{-1}(R_{\det})\) 为两通道通量之差。证毕。
\end{proof}

\begin{conclusion}[端点 Poisson 通量的可审计迹公式：\texorpdfstring{$t=-1$}{t=-1} 处 resolvent 的通道差]\label{con:xi-endpoint-flux-resolvent-trace-channel-difference}
令 \(u\mapsto M^{\pm}(u)\) 为某个复邻域内解析的有限维算子族，并假设 \(I+M^{\pm}(u)\) 可逆。
定义行列式比值
\[
R(u,t):=\frac{\det(I-tM^{+}(u))}{\det(I-tM^{-}(u))}.
\]
若进一步 \(R(u,\cdot)\) 在 \(t=-1\) 处解析且 \(R(u,-1)\neq 0\)，则可定义端点通量
\[
\Phi_{-1}(R(u,\cdot))
:=
-\Re\left(\frac{\dd}{\dd t}\log R(u,t)\Big|_{t=-1}\right).
\]
此时成立严格迹恒等式
\[
\boxed{\;
\Phi_{-1}(R(u,\cdot))
=
\Re\,\mathrm{tr}\!\Bigl((I+M^{+}(u))^{-1}M^{+}(u)-(I+M^{-}(u))^{-1}M^{-}(u)\Bigr)\; }.
\]
进一步，若存在整数 \(A_{\mathrm{flip}}\ge 1\) 以及算子 \(M_0,B\) 使得在 \(u\downarrow 0\) 下有
\[
M^{\pm}(u)=M_0\pm u^{A_{\mathrm{flip}}}B+O\bigl(u^{A_{\mathrm{flip}}+1}\bigr),
\qquad
\det(I+M_0)\neq 0,
\]
则端点通量差具有首项展开
\[
\boxed{\;
\Phi_{-1}(R(u,\cdot))
=
2u^{A_{\mathrm{flip}}}\Re\,\mathrm{tr}\!\Bigl((I+M_0)^{-1}B(I+M_0)^{-1}\Bigr)
+O\bigl(u^{A_{\mathrm{flip}}+1}\bigr)\; }.
\]
特别地，若 \(\Re\,\mathrm{tr}\!\bigl((I+M_0)^{-1}B(I+M_0)^{-1}\bigr)\neq 0\)，则 \(\operatorname{ord}_u \Phi_{-1}(R(u,\cdot))=A_{\mathrm{flip}}\)。
\end{conclusion}
\begin{proof}[证明要点]
由恒等式
\(
\frac{\dd}{\dd t}\log\det(I-tM)=-\mathrm{tr}\!\bigl((I-tM)^{-1}M\bigr)
\)
得
\[
\frac{\dd}{\dd t}\log R(u,t)
=
-\mathrm{tr}\!\bigl((I-tM^{+}(u))^{-1}M^{+}(u)\bigr)
+\mathrm{tr}\!\bigl((I-tM^{-}(u))^{-1}M^{-}(u)\bigr),
\]
代入 \(t=-1\) 并取实部即得第一式。
\par
对首项展开，令 \(K:=(I+M_0)^{-1}\)。
由 Fr\'echet 展开
\(
(I+M_0\pm u^{A_{\mathrm{flip}}}B)^{-1}=K\mp u^{A_{\mathrm{flip}}}KBK+O(u^{A_{\mathrm{flip}}+1})
\)，
并用恒等式 \((I+M)^{-1}M=I-(I+M)^{-1}\)，
得到
\[
(I+M^{+}(u))^{-1}M^{+}(u)-(I+M^{-}(u))^{-1}M^{-}(u)
=(I+M^{-}(u))^{-1}-(I+M^{+}(u))^{-1}
=2u^{A_{\mathrm{flip}}}KBK+O(u^{A_{\mathrm{flip}}+1}),
\]
取迹并取实部即得所示首项系数。
\end{proof}

\begin{conclusion}[端点通量的两条稀疏证书：深度计数上界与端点读出能量预算]\label{con:xi-endpoint-flux-sparsity-and-readout-budget}
沿用结论 \ref{con:xi-endpoint-flux-operator-trace-identity} 的有限 Blaschke 设定。
对任意阈值 \(\delta>0\) 定义端点深度计数
\[
N_\delta(B):=\#\{1\le k\le \kappa:\ P_{a_k}(-1)\ge \delta\}.
\]
则有严格上界
\[
\boxed{\ N_\delta(B)\ \le\ \frac{\Phi_{-1}(B)}{\delta}\ }.
\]
此外，对任意 \(f\in K_B\) 有端点读出能量预算
\[
\boxed{\ |f(-1)|^2\ \le\ \Phi_{-1}(B)\,\|f\|^2\ }.
\]
并且常数 \(\Phi_{-1}(B)\) 为最优：存在 \(f_\star\in K_B\setminus\{0\}\) 使等号成立。
\end{conclusion}
\begin{proof}[证明要点]
第一条由定义直接得到：若 \(P_{a_k}(-1)\ge \delta\) 的项有 \(N_\delta(B)\) 个，则
\(
\Phi_{-1}(B)=\sum_k P_{a_k}(-1)\ge N_\delta(B)\delta
\)，从而 \(N_\delta(B)\le \Phi_{-1}(B)/\delta\)。
\par
第二条由 Cauchy--Schwarz 给出：\(|f(-1)|\le \|L_{-1}\|\,\|f\|\)。
在有限 Blaschke 设定下 \(\dim K_B=\kappa<\infty\)，算子 \(L_{-1}^\ast L_{-1}\) 为正的秩一算子，故
\[
\|L_{-1}\|^2=\mathrm{tr}(L_{-1}^\ast L_{-1})=\Phi_{-1}(B)
\]
（结论 \ref{con:xi-endpoint-flux-operator-trace-identity}）。
代回即可得 \(|f(-1)|^2\le \Phi_{-1}(B)\|f\|^2\)。
最优性由再生核向量的取等条件（或等价的秩一算子谱分解）给出。
\end{proof}

\begin{conclusion}[时间纤维的内生性：上半平面虚部即端点 \texorpdfstring{$-1$}{-1} 的 Poisson--Busemann 密度]\label{con:xi-timefiber-poisson-busemann-identity}
令 \(t=x+\mathrm{i}y\in\mathbb{H}\)，并取 Cayley 变换
\[
w=\mathcal{C}(t)=\frac{1+\mathrm{i}t}{1-\mathrm{i}t}\in\mathbb{D}.
\]
则成立严格恒等式
\[
\boxed{\ y=B^{-1}(w)=\frac{1-|w|^2}{|1+w|^2}\ }.
\]
因此在端点紧化口径下，“时间纤维坐标”并非外加连续轴，而被严格同一为端点 \(-1\) 的谐测度密度（Poisson--Busemann 权重）；高高度信息被强制压缩为 \(w\approx-1\) 的边界层贡献，从而形成可核验的端点局部化原则。
\end{conclusion}
\begin{proof}[证明要点]
见推论 \ref{cor:dephys-poisson-busemann-timefiber}（并引用命题 \ref{prop:app-busemann-poisson-minus1} 的 Busemann/Poisson 定义接口）。
\end{proof}

\paragraph{window-\texorpdfstring{$6$}{6} 的最小热化闭合：推送核、谱隙与 lumpability 障碍}

\begin{conclusion}[window-\texorpdfstring{$6$}{6} 的最小热化闭合：推送核的可逆性、细致平衡与显式谱隙]\label{con:xi-window6-p6-minimal-thermalization-closure}
以推送邻接计数 \(A_6\) 定义 \(X_6\) 上的推送核 \(P_6\) 与平稳分布 \(\pi_6\)（生成片段 \eqref{eq:window6_pushforward_markov_kernel}）：
\[
\pi_6(x)=\frac{d^{\mathrm{bin}}_6(x)}{64},\qquad
P_6(x,x')=\frac{A_6(x,x')}{6\,d^{\mathrm{bin}}_6(x)}.
\]
则 \(P_6\) 满足细致平衡
\(
\pi_6(x)P_6(x,x')=\pi_6(x')P_6(x',x)
\)，
因而可逆并以 \(\pi_6\) 为平稳分布。
谱隙审计给出可复核常数（生成片段 \eqref{eq:window6_pushforward_markov_spectral_gap}）
\[
\lambda_2(P_6)\approx 0.4841207858,\qquad
\lambda_\star(P_6)\approx 0.6030939754,
\]
从而对任意初始分布 \(\mu\) 存在常数 \(C(\mu)\) 使对所有 \(t\ge 0\) 有指数混合上界
\[
\|\mu P_6^t-\pi_6\|_{\mathrm{TV}}\le C(\mu)\,\lambda_\star(P_6)^t.
\]
\end{conclusion}
\begin{proof}[证明要点]
核与细致平衡见定理 \ref{thm:terminal-foldbin6-pushforward-markov} 与 \eqref{eq:window6_pushforward_markov_kernel}；
谱隙常数与 TV 上界见推论 \ref{cor:terminal-foldbin6-minimal-markov-closure-gap} 及 \eqref{eq:window6_pushforward_markov_spectral_gap}。
\end{proof}

\begin{conclusion}[从谱隙到混合时间的硬常数：window-\texorpdfstring{$6$}{6} 的热化时间尺度为 \texorpdfstring{$1/(-\log\lambda_\star)$}{1/(-log lambda*)}]\label{con:xi-window6-mixing-time-hard-constant}
沿用上条结论的记号。
若对给定初始分布 \(\mu\) 定义收敛时间
\[
t(\varepsilon;\mu):=\min\bigl\{t\ge 0:\ \|\mu P_6^t-\pi_6\|_{\mathrm{TV}}\le \varepsilon\bigr\},
\]
则由 \(\|\mu P_6^t-\pi_6\|_{\mathrm{TV}}\le C(\mu)\lambda_\star(P_6)^t\) 得到显式上界
\[
\boxed{\ t(\varepsilon;\mu)\ \le\ \frac{\log(C(\mu)/\varepsilon)}{-\log\lambda_\star(P_6)}\ }
\approx 1.9775264\cdot \log\!\bigl(C(\mu)/\varepsilon\bigr).
\]
因此 window-\(6\) 的无记忆热化闭合在时间尺度上携带一个实现无关的协议级硬常数 \(1/(-\log\lambda_\star(P_6))\)，并由谱证书直接输出而不依赖额外拟合。
\end{conclusion}
\begin{proof}[证明要点]
由推论 \ref{cor:terminal-foldbin6-minimal-markov-closure-gap} 的 TV 指数收敛式，取对数即得；
数值常数由 \eqref{eq:window6_pushforward_markov_spectral_gap} 中的 \(\lambda_\star(P_6)\) 代入计算得到。
\end{proof}

\begin{conclusion}[强可并失败的结构性后果：稳定类型过程在一般初态下不可能为一阶马尔可夫闭合]\label{con:xi-window6-strong-lumpability-failure-structural}
由纤维划分 \(\{F_6^{-1}(x)\}_{x\in X_6}\) 诱导的粗粒化不满足 strong lumpability：
存在同一纤维内两点 \(\omega,\omega'\) 使其一步邻域在稳定类型上的计数分布不同（显式反例见 \eqref{eq:foldbin6_strong_lumpability_counterexample}）。
因此不存在任何仅依赖 \(X_t\) 的齐次转移核 \(P\) 能在\emph{所有}初始微态分布下使 \((X_t)\) 闭合为 \(X_6\) 上的无记忆马尔可夫链。
\end{conclusion}
\begin{proof}[证明要点]
这是命题 \ref{prop:terminal-foldbin6-strong-lumpability-fails}；反例由生成片段 \eqref{eq:foldbin6_strong_lumpability_counterexample} 给出。
\end{proof}

\begin{conclusion}[最小无记忆闭包的强迫性：纤维内瞬时热化与推送核闭合严格等价]\label{con:xi-window6-instant-thermalization-equivalence}
在强可并性失败的背景下，若引入最小的无记忆闭包机制：每一步立方体翻转后，条件于当前稳定类型 \(X_t=x\) 将微态立即均匀化到纤维 \(F_6^{-1}(x)\)（纤维内瞬时热化），
则稳定类型过程闭合为 \(X_6\) 上的可逆马尔可夫链，其转移核必为推送核 \(P_6\)，平衡测度必为 \(\pi_6\propto d^{\mathrm{bin}}_6\)。
等价地，在该协议口径下，一阶无记忆闭合的存在性与“纤维内瞬时热化”严格等价。
\end{conclusion}
\begin{proof}[证明要点]
闭合的唯一性与可逆性由定理 \ref{thm:terminal-foldbin6-pushforward-markov} 与推论 \ref{cor:terminal-foldbin6-minimal-markov-closure-gap} 给出；其强迫性由命题 \ref{prop:terminal-foldbin6-strong-lumpability-fails} 的否定性结论推出。
\end{proof}

\paragraph{Witt--Fourier 双桥与算子化 Big-Witt：统一幽灵标尺与 \texorpdfstring{$p$}{p}-典型护栏}

\begin{conclusion}[Witt--Fourier 双桥的严格粘接：同一角色轴 \texorpdfstring{$\mathrm{Tw}_\chi$}{Tw\_chi} 同时对角化“卷积--碰撞”与“迹--Euler”]\label{con:xi-witt-fourier-commuting-cube}
在有限环 \(G_m\) 及其角色群 \(\widehat G_m\) 上，对每个角色 \(\chi\) 定义纤维侧扭曲
\[
\mathrm{Tw}_\chi(d_m):=\sum_{r\in G_m}d_m(r)\chi(r)=\widehat d_m(\chi),
\]
并在轨道侧以同一 \(\chi\) 定义 Dirichlet 扭曲核 \(T_\chi\)。
则 Fourier 对角化与 Witt/Euler 交换图在同一条 \(\mathrm{Tw}_\chi\) 轴上严格相容，二者沿角色方向粘接形成严格交换立方体；
并且在 \(\Fold_m\) 口径下有完全因子化
\[
\mathrm{Tw}_\chi(d_m)=\prod_{j=1}^m\bigl(1+\chi(F_{j+1})\bigr).
\]
\end{conclusion}
\begin{proof}[证明要点]
交换立方体与扭曲轴一致性见注 \ref{rem:twist-commuting-cube}；
乘积分解见定理 \ref{thm:fold-spectrum-factorization}（并在同一注中给出对齐方式）。
\end{proof}

\begin{conclusion}[统一幽灵标尺：二阶碰撞缺口同时度量 Fourier 非平凡频谱能量与 Witt/Euler 的偏离强度]\label{con:xi-collision-gap-ghost-yardstick}
Witt 侧的“迹 \(\to\) primitive \(\to\) Euler”桥包与 Fourier 侧的 Parseval 能量桥包共享同一个可计算的标量指纹——二阶碰撞缺口（等价于非平凡频谱能量的 \(\ell^2\)-范数）：
\[
\boxed{\ F_{m+2}\,\mathsf{Col}_m-1=\frac{1}{4^m}\sum_{t\ne 0}\bigl|\widehat c_m(t)\bigr|^2\ }.
\]
因此 prime-like（Euler 乘积）侧与谱/迹侧的偏离强度可被压缩到同一个可审计标量上，形成跨通道一致性核验的最短核验回路。
\end{conclusion}
\begin{proof}[证明要点]
见注 \ref{rem:fw-bridge}（并由定理 \ref{thm:fold-collision-spectrum} 的 Parseval 恒等式给出该标尺的严格推导；其与 Rényi-2/均匀性缺口的等价关系见推论 \ref{cor:pom-renyi2-near-uniform}）。
\end{proof}

\begin{conclusion}[碰撞概率的谱迹公式与 Riesz--Parseval 恒等式]\label{con:xi-fold-collision-spectral-trace-parseval}
令 \(M:=F_{m+2}\)，并把折叠多重度剖面视为模环上的函数
\(
d_m:\ZZ/(M\ZZ)\to\NN
\)。
定义归一化分布
\[
\mu_m(r):=\frac{d_m(r)}{2^m},
\qquad
u(r):=\frac{1}{M}\qquad\bigl(r\in\ZZ/(M\ZZ)\bigr),
\]
以及碰撞概率
\[
\mathsf{Col}_m:=\sum_{r\in\ZZ/(M\ZZ)}\mu_m(r)^2.
\]
则有严格谱迹恒等式
\[
\boxed{\;
\mathsf{Col}_m
=\frac{1}{M}\sum_{k\in\ZZ/(M\ZZ)}
\prod_{j=1}^{m}\cos^2\!\Bigl(\pi k\,\frac{F_{j+1}}{M}\Bigr)\; }.
\]
因此
\[
\mathsf{Col}_m-\frac1M
=\frac{1}{M}\sum_{k\ne 0}
\prod_{j=1}^{m}\cos^2\!\Bigl(\pi k\,\frac{F_{j+1}}{M}\Bigr).
\]
并且 \(\chi^2\)-散度与碰撞缺口满足精确等价
\[
\boxed{\;
\chi^2(\mu_m\mid u)
:=\sum_{r}\frac{\bigl(\mu_m(r)-u(r)\bigr)^2}{u(r)}
=M\sum_{r}\Bigl(\mu_m(r)-\frac1M\Bigr)^2
=M\mathsf{Col}_m-1
=M\Bigl(\mathsf{Col}_m-\frac1M\Bigr).\;}
\]
\end{conclusion}
\begin{proof}[证明要点]
由定理 \ref{thm:fold-collision-spectrum} 的 Parseval 恒等式，
\(
\mathsf{Col}_m=\frac{1}{M}\sum_k|\widehat\mu_m(k)|^2
\)。
再由定理 \ref{thm:fold-spectrum-factorization} 的乘积分解得到
\(
|\widehat\mu_m(k)|^2=\prod_{j=1}^m\cos^2(\pi kF_{j+1}/M)
\)，代回即得谱迹公式。
\(\chi^2\) 恒等式则由
\(
M\sum_r(\mu_m(r)-1/M)^2=M\sum_r\mu_m(r)^2-1
\)
直接化简得到。证毕。
\end{proof}

\begin{conclusion}[群环卷积算子的零点--核维数等价]\label{con:xi-fold-group-ring-convolution-kernel-dimension}
在群环
\[
\CC[\ZZ/(M\ZZ)]\cong\CC[x]/(x^M-1)
\]
中令生成函数
\[
D_m(x):=\sum_{r=0}^{M-1}d_m(r)x^r,
\]
并定义卷积/乘法算子
\[
T_m:\CC[x]/(x^M-1)\to\CC[x]/(x^M-1),\qquad f\longmapsto D_m\cdot f\ \ (\mathrm{mod}\ x^M-1).
\]
记谱零点集合
\[
\mathcal Z_m:=\{k\in\ZZ/(M\ZZ):\ \widehat d_m(k)=0\}.
\]
则核维数与零点数严格一致：
\[
\boxed{\;
\dim_\CC\ker(T_m)=|\mathcal Z_m|,
\qquad
\mathrm{rank}(T_m)=M-|\mathcal Z_m|.\;}
\]
\end{conclusion}
\begin{proof}[证明要点]
取标准 Fourier 基（或等价地：把 \(\CC[x]/(x^M-1)\) 以 \(x\mapsto e^{-2\pi i k/M}\) 的特征分解对角化）。
在该基下，乘法算子 \(T_m\) 对角化，其特征值为
\(
D_m(e^{-2\pi i k/M})=\widehat d_m(k)
\)（见命题 \ref{prop:fold-multiplicity-group-algebra} 与定理 \ref{thm:fold-spectrum-factorization}）。
故核由所有 \(\widehat d_m(k)=0\) 的特征方向张成，维数等于零点数目。证毕。
\end{proof}

\begin{conclusion}[谱零点的同余方程刻画与余类分层]\label{con:xi-fold-spectrum-zero-congruence-coset-stratification}
令 \(M:=F_{m+2}\) 且 \(k\in\ZZ/(M\ZZ)\)。
则
\[
k\in\mathcal Z_m
\quad\Longleftrightarrow\quad
\exists j\in\{1,\dots,m\}\ \text{s.t.}\ e^{-2\pi i kF_{j+1}/M}=-1.
\]
等价地（因此必有 \(M\) 偶）：
\[
k\in\mathcal Z_m
\quad\Longleftrightarrow\quad
\bigl(M\ \text{为偶数}\bigr)\ \wedge\
\exists j\in\{1,\dots,m\}\ \text{s.t.}\ kF_{j+1}\equiv \frac{M}{2}\pmod M.
\]
并且对固定的 \(j\)，令 \(g_j:=\gcd(F_{j+1},M)\)。
则同余
\(
kF_{j+1}\equiv M/2\ (\mathrm{mod}\ M)
\)
可解当且仅当 \(g_j\mid M/2\)；
在可解时解集为 \(\ZZ/(M\ZZ)\) 的一个仿射陪集，且其元素个数恰为 \(g_j\)。
\end{conclusion}
\begin{proof}[证明要点]
第一段由定理 \ref{thm:fold-spectrum-factorization}：\(\widehat d_m(k)=\prod_{j=1}^m(1+e^{-2\pi i kF_{j+1}/M})\)，故为零当且仅当某因子为零。
第二段与第三段是线性同余方程的标准判别，并在本文记号下由定理 \ref{thm:fold-zero-coset-rigidity} 给出显式陪集与其大小。证毕。
\end{proof}

\begin{conclusion}[六倍窗的核维数下界与线性不可逆的半阶崩塌]\label{con:xi-fold-window6-kernel-half-index-lower-bound}
若 \(m+2\equiv 0\pmod 6\)，则存在一个半阶余类块满足
\[
|\mathcal Z_m|\ \ge\ F_{(m+2)/2}.
\]
因此由结论 \ref{con:xi-fold-group-ring-convolution-kernel-dimension} 立刻推出
\[
\dim_\CC\ker(T_m)\ge F_{(m+2)/2},
\qquad
\mathrm{rank}(T_m)\le M-F_{(m+2)/2}.
\]
\end{conclusion}
\begin{proof}[证明要点]
谱零点下界为推论 \ref{cor:fold-zero-half-index-multiple6}；
核维数等价由结论 \ref{con:xi-fold-group-ring-convolution-kernel-dimension} 得到。证毕。
\end{proof}

\begin{conclusion}[\texorpdfstring{$\chi^2$}{chi^2} 距离的常数级刚性下界]\label{con:xi-fold-chi2-constant-residual-energy}
令 \(M:=F_{m+2}\)，并沿用结论 \ref{con:xi-fold-collision-spectral-trace-parseval} 的 \(\mu_m,u\)。
设临界共振常数 \(C_\varphi>0\) 如定理 \ref{thm:fold-critical-resonance-constant}。
则有常数级下界
\[
\boxed{\;
\liminf_{m\to\infty}\chi^2(\mu_m\mid u)\ \ge\ 2C_\varphi^2\;}
\]
等价地，
\[
\mathsf{Col}_m-\frac1M\ \ge\ \frac{2C_\varphi^2+o(1)}{M}.
\]
\end{conclusion}
\begin{proof}[证明要点]
由结论 \ref{con:xi-fold-collision-spectral-trace-parseval}，
\(
\chi^2(\mu_m\mid u)=M(\mathsf{Col}_m-1/M)
\)。
而推论 \ref{cor:fold-critical-resonance-gap} 给出
\(
\liminf_{m\to\infty}M(\mathsf{Col}_m-1/M)\ge 2C_\varphi^2
\)。
合并即得。证毕。
\end{proof}

\begin{conclusion}[高阶碰撞核 Perron 根的端点斜率由最大纤维指数锁定]\label{con:xi-fold-high-q-perron-endpoint-slope-maxfiber}
记最大纤维 \(M_m:=\max_{x\in X_m}d_m(x)\)，并令 \(\alpha_m:=\frac1m\log M_m\)。
若极限 \(\alpha_\ast:=\lim_{m\to\infty}\alpha_m\) 存在且有限，则对每个固定整数 \(q\ge 1\)，令
\[
S_q(m):=\sum_{x\in X_m}d_m(x)^q,
\qquad
\log r_q:=\lim_{m\to\infty}\frac{1}{m}\log S_q(m)
\quad(\text{若极限存在}),
\]
则恒有双边夹逼
\[
q\alpha_\ast\ \le\ \log r_q\ \le\ q\alpha_\ast+\log\varphi,
\]
并且端点斜率满足
\[
\lim_{q\to\infty}\frac{1}{q}\log r_q=\alpha_\ast.
\]
\end{conclusion}
\begin{proof}[证明要点]
由平凡夹逼 \(M_m^q\le S_q(m)\le |X_m|\,M_m^q\) 得
\[
q\alpha_m\ \le\ \frac{1}{m}\log S_q(m)\ \le\ q\alpha_m+\frac{1}{m}\log|X_m|.
\]
利用 \(|X_m|=F_{m+2}=\Theta(\varphi^m)\) 与 \(\alpha_m\to\alpha_\ast\) 取极限得到 \(\log r_q\) 的夹逼；
再令 \(q\to\infty\) 得端点斜率。证毕。
\end{proof}

\begin{conclusion}[Rényi 指纹率的零温极限]\label{con:xi-fold-renyi-fingerprint-rate-zero-temp}
沿用结论 \ref{con:xi-fold-high-q-perron-endpoint-slope-maxfiber} 的 \(r_q\) 与 \(\alpha_\ast\)。
定义每阶 \(q\) 的指纹率
\[
h_R(q):=\frac{q\log 2-\log r_q}{q-1}.
\]
则存在零温极限
\[
\boxed{\;\lim_{q\to\infty}h_R(q)=\log 2-\alpha_\ast.\;}
\]
\end{conclusion}
\begin{proof}[证明要点]
由结论 \ref{con:xi-fold-high-q-perron-endpoint-slope-maxfiber}，有 \(\log r_q=q\alpha_\ast+O(1)\)，代入定义并令 \(q\to\infty\) 得结论。证毕。
\end{proof}

\begin{conclusion}[规范体积密度的余维一常数：离散体积与折叠缺陷熵的稳定平移]\label{con:xi-fold-gauge-volume-codim1-constant}
对任意折叠系统，令 \(d_m(w)\) 为纤维大小函数并定义规范群体积
\[
|G_m|:=\prod_{w\in X_m} d_m(w)!,
\qquad
g_m:=\frac{1}{2^m}\log|G_m|,
\qquad
\kappa_m:=\frac{1}{2^m}\sum_{w\in X_m} d_m(w)\log d_m(w).
\]
在计数守恒 \(\sum_w d_m(w)=2^m\) 下，Stirling 展开给出普适差常数：存在与 \(m\) 无关的绝对常数使
\[
g_m=\kappa_m-1+O\!\Bigl(\frac{|X_m|\,m}{2^m}\Bigr).
\]
在本文给定的 Fibonacci 稳定类型尺度 \(|X_m|=F_{m+2}=\Theta(\varphi^m)\) 下，误差为指数级小量，从而
\[
\kappa_m-g_m=1+O\!\bigl(m(\varphi/2)^m\bigr)=1+o(1).
\]
\end{conclusion}
\begin{proof}[证明要点]
见定理 \ref{thm:fold-bin-gauge-volume-stirling-second-order}（并结合命题 \ref{prop:fold-bin-gauge-volume} 的记号设定与余项界）。证毕。
\end{proof}
\paragraph{规范体积密度与输出熵的定量镜像}
在 bin-fold 口径下，纤维信息项 \(\kappa_m\)、输出 Shannon 熵与规范体积密度 \(g_m\) 并非三条彼此独立的复杂度账本，而是由同一条带有显式常数项的 Stirling--entropy 恒等式锁定。

\begin{theorem}[规范体积密度与输出 Shannon 熵之间的一纳特定律]\label{thm:xi-foldbin-gauge-entropy-one-nat-law}
\leanverified{paper\_xi\_foldbin\_gauge\_entropy\_one\_nat\_law}
令
\[
\mu_m(w):=\frac{d_m^{\mathrm{bin}}(w)}{2^m},
\qquad
w\in X_m,
\]
并记
\[
\kappa_m:=\mathbb E_{\mu_m}\!\bigl[\log d_m^{\mathrm{bin}}(W_m)\bigr],
\qquad
g_m:=2^{-m}\log|\mathcal G_m^{\mathrm{bin}}|,
\]
其中 \(W_m\sim \mu_m\)。
则有精确恒等式
\[
\kappa_m=m\log 2-H(\mu_m),
\]
其中
\[
H(\mu_m):=-\sum_{w\in X_m}\mu_m(w)\log\mu_m(w).
\]
进一步，
\[
g_m
=
m\log 2-H(\mu_m)-1
+\frac{(\varphi/2)^m}{2\sqrt5}
\Bigl(
\varphi^2 m\log\frac{2}{\varphi}
+\varphi^2\log(2\pi)-\log\varphi
\Bigr)
+O\!\bigl((\varphi/2)^m\bigr),
\]
从而
\[
m\log 2-H(\mu_m)-g_m\longrightarrow 1.
\]
\end{theorem}
\begin{proof}
由 \(\mu_m(w)=d_m^{\mathrm{bin}}(w)/2^m\) 得
\[
\log d_m^{\mathrm{bin}}(w)=m\log 2+\log\mu_m(w).
\]
对 \(\mu_m\) 取期望可得
\[
\kappa_m
=
\sum_{w\in X_m}\mu_m(w)\log d_m^{\mathrm{bin}}(w)
=
m\log 2+\sum_{w\in X_m}\mu_m(w)\log\mu_m(w)
=
m\log 2-H(\mu_m).
\]

另一方面，命题 \ref{prop:fold-bin-gauge-density-two-scale} 给出
\[
g_m=\kappa_m-1
+\frac{(\varphi/2)^m}{2\sqrt5}
\Bigl(
\varphi^2 m\log\frac{2}{\varphi}
+\varphi^2\log(2\pi)-\log\varphi
\Bigr)
+O\!\bigl((\varphi/2)^m\bigr).
\]
将 \(\kappa_m=m\log 2-H(\mu_m)\) 代入即得所示展开。
最后，由结论 \ref{con:xi-fold-gauge-volume-codim1-constant} 已知
\(
\kappa_m-g_m=1+O(m(\varphi/2)^m)
\)，故
\[
m\log 2-H(\mu_m)-g_m
=
\kappa_m-g_m
\longrightarrow 1.
\]
证毕。
\end{proof}

\begin{corollary}[bin-fold 输出 Shannon 熵与 gauge 密度的显式常数项镜像]\label{cor:xi-foldbin-shannon-gauge-constant-mirror}
\leanverified{paper\_xi\_foldbin\_shannon\_gauge\_constant\_mirror}
沿用定理 \ref{thm:xi-foldbin-gauge-entropy-one-nat-law} 的记号。则有展开
\[
\kappa_m
=
m\log\frac{2}{\varphi}
-\frac{\log\varphi}{1+\varphi^2}
+O\!\bigl((\varphi/2)^m\bigr)
=
m\log\frac{2}{\varphi}
-\frac{\log\varphi}{\varphi\sqrt5}
+O\!\bigl((\varphi/2)^m\bigr),
\]
从而
\[
H(\mu_m)
=
m\log\varphi
+\frac{\log\varphi}{\varphi\sqrt5}
+O\!\bigl((\varphi/2)^m\bigr),
\]
以及
\[
g_m
=
m\log\frac{2}{\varphi}
-1
-\frac{\log\varphi}{\varphi\sqrt5}
+O\!\bigl(m(\varphi/2)^m\bigr).
\]
\end{corollary}
\begin{proof}
当 \(q=1\) 时，
\(
\pi_{m,1}(w)=d_m^{\mathrm{bin}}(w)/S_1^{\mathrm{bin}}(m)=d_m^{\mathrm{bin}}(w)/2^m=\mu_m(w)
\)，
故命题 \ref{prop:fold-bin-escort-lastbit} 在 \(q=1\) 时给出
\[
\kappa_m
=
m\log\frac{2}{\varphi}
-\frac{\log\varphi}{1+\varphi^2}
+O\!\bigl((\varphi/2)^m\bigr).
\]
由恒等式 \(1+\varphi^2=\varphi\sqrt5\) 得第二种写法。

再由定理 \ref{thm:xi-foldbin-gauge-entropy-one-nat-law}，
\[
H(\mu_m)=m\log 2-\kappa_m.
\]
代入上式即得 Shannon 熵展开。

最后，由结论 \ref{con:xi-fold-gauge-volume-codim1-constant}，
\[
g_m=\kappa_m-1+O\!\bigl(m(\varphi/2)^m\bigr),
\]
将 \(\kappa_m\) 的展开代入即得 \(g_m\) 的常数项公式。证毕。
\end{proof}


\begin{conclusion}[Big-Witt 运算的算子编译：原子 Euler 数据的半环可同态嵌入迹类算子组合]\label{con:xi-atomic-witt-operator-compilation}
令 \(\mathcal{W}_{\mathrm{atom}}\) 为由原子数据 \((\beta,n,m)\) 生成的 Big-Witt 型半环，令 \(\mathrm{TC}_1\) 为迹类算子在酉等价类下的交换半环（加法为直和、乘法为张量积）。
则映射
\[
(\beta,n,m)\longmapsto \bigoplus_{i=1}^{m} C(n,\alpha),\qquad \alpha^n=\beta,
\]
延拓为构造性的半环同态
\[
\boxed{\ \mathcal{W}_{\mathrm{atom}}\ \longrightarrow\ \mathrm{TC}_1\ }.
\]
因此 Big-Witt 的抽象 Euler/plethystic 运算可被编译为迹类算子的直和/张量积组合运算，并与 Fredholm \(\zeta\) 接口严格相容，形成可执行的 Euler 代数。
\end{conclusion}
\begin{proof}[证明要点]
见定理 \ref{thm:atomic-witt-into-tc1}（乘法的 \(\gcd/\mathrm{lcm}\) 规则由命题 \ref{prop:cyclic-block-tensor-gcd-lcm} 保证）。
\end{proof}

\begin{conclusion}[\texorpdfstring{$p$}{p}-典型 Frobenius--Verschiebung 的迹层实现：\texorpdfstring{$\mathsf F_p\circ\mathsf V_p$}{Fp o Vp} 等价于“直和乘以 \texorpdfstring{$p$}{p}”]\label{con:xi-ptypical-frobenius-verschiebung-trace}
在 \(p\)-典型分量上，Frobenius 与 Verschiebung 的算子实现满足 ghost 兼容性与 Verschiebung--Frobenius 关系：
对任意由 \(p\)-幂循环块有限直和得到的算子 \(L\)，对一切 \(n\ge 1\) 有
\[
\boxed{\ \Tr\!\bigl(\mathsf F_p(L)^n\bigr)=\Tr\!\bigl(L^{pn}\bigr)\ },
\qquad
\boxed{\ \mathsf F_p\!\bigl(\mathsf V_p(L)\bigr)\ \simeq\ L^{\oplus p}\ }.
\]
该关系把 Dwork 型同余从“多项式/序列”提升为“迹类算子组合”的结构恒等式，使 \(p\)-进护栏以算子函子方式零成本继承。
\end{conclusion}
\begin{proof}[证明要点]
见命题 \ref{prop:cyclic-p-typical-frobenius-verschiebung}。
\end{proof}

\paragraph{模障碍、cyclotomic 因子与 Smith 形：共振角色表的有限终止审计}

\begin{conclusion}[模障碍与 cyclotomic 因子的完全等价：共振来源归约为单位根因子分类]\label{con:xi-mod-obstruction-cyclotomic-equivalence}
对 Ihara/动力学 \(\zeta\) 的扭曲通道 \(\det(I-tB_\rho(u))^{-1}\)，存在非平凡模障碍当且仅当某个 \(\rho\neq \mathbf{1}\) 使相应谱半径达到平凡通道；
且该条件当且仅当 \(\det(I-tB_\rho(u))\) 含 cyclotomic 因子（归一化后在单位圆上出现单位根因子）。
因此“共振来源”可被等价归约为有限维整系数多项式的 cyclotomic 因子分类问题。
\end{conclusion}
\begin{proof}[证明要点]
见结论 \ref{con:xi-mod-obstruction-cyclotomic-factor} 及其所引分类定理 \ref{thm:conclusion75-ihara-cyclotomic-factor-classification}。
\end{proof}

\begin{conclusion}[障碍群的同调分类与多项式时间可计算性：全共振角色表可编译为一次 SNF 计算]\label{con:xi-resonant-character-snf-computable}
共振角色集合 \(X_{\mathrm{res}}\) 构成有限阿贝尔群，且可由同调商的 torsion 部分刻画；
其不变因子分解与生成角色可由相应整数矩阵的 Smith 正规形直接读出。
更强地，存在确定性算法在时间 \(\mathrm{poly}(|V|,|E|,\log W)\) 内输出障碍不变量与 \(X_{\mathrm{res}}\) 的不变因子分解及生成角色，从而“是否存在 cyclotomic 因子/共振通道”的全角色审计可在有限步内终止。
\end{conclusion}
\begin{proof}[证明要点]
同调刻画见定理 \ref{thm:gm-resonant-character-group-homology}；Smith 形与多项式时间可计算性见命题 \ref{prop:gm-smith-compute-obstructions}。
\end{proof}

\subsubsection{中心六重聚合与 Perron 分支的六分支惯性：解析半径与代数次数约束}

\begin{conclusion}[中心六重聚合的几何后果：中心唯一实交与最近复奇点距离给出解析半径]\label{con:xi-rate-center-sixfold-unique-real-rstar}
同步核速率曲线的规范化消元多项式 \(R(\alpha,u)\) 具有对偶自对称结构（注 \ref{rem:rate-algebraic-elimination-structure}），并在中心截面满足分解
\[
R\!\left(\tfrac12,u\right)=-\frac{(u-1)^6}{64}\,Q(u).
\]
其中 \(Q(u)\) 系数回文，故 \(Q(u)=u^{16}Q(1/u)\)，零点以互逆对出现（并与复共轭合并为四元组）；并且 Sturm 计数给出 \(Q\) 在实轴无零点，从而 \(Q(u)>0\) 对所有 \(u\in\RR\) 成立（命题 \ref{prop:rate-center-Q-reciprocal-no-real}）。
因此 \(R(\tfrac12,u)=0\) 在实参数上当且仅当 \(u=1\)，且重数恰为 \(6\)（推论 \ref{cor:rate-center-unique-real}）。
进一步定义最近复奇点距离
\[
R_\star:=\min\{|u-1|:\ Q(u)=0\},
\]
则 \(R_\star\) 给出中心分支的解析半径并可数值审计（推论 \ref{cor:rate-center-Rstar}）。
\end{conclusion}
\begin{proof}[证明要点]
自对称与中心分解的可复算证书见注 \ref{rem:rate-algebraic-elimination-structure}；
\(Q\) 的自倒易性与实根空性见命题 \ref{prop:rate-center-Q-reciprocal-no-real}；
中心唯一实交与六重性见推论 \ref{cor:rate-center-unique-real}；
最近复奇点距离的数值证书见推论 \ref{cor:rate-center-Rstar}。
\end{proof}

\begin{conclusion}[惯性群强迫的“六分支单元”：Perron 分支代数次数必被 \texorpdfstring{$6$}{6} 整除]\label{con:xi-rate-center-perron-degree-multiple-6}
沿用中心六重聚合证书 \(R(\tfrac12,u)=-(u-1)^6Q(u)/64\)，并假设一般横截条件 \(\partial_\alpha R(\tfrac12,1)\neq 0\)。
令 \(F(\alpha,u)\in\QQ[\alpha,u]\) 为 \(R(\alpha,u)\) 在 \(\QQ(\alpha)[u]\) 中包含点 \((\tfrac12,1)\) 的不可约因子（例如包含 Perron 分支的那一因子）。
则在函数域扩张 \(\QQ(\alpha)\subset \QQ(\alpha,u)\ (F(\alpha,u)=0)\) 中，\(\alpha=\tfrac12\) 处的分歧指数为 \(e=6\)，从而 \(\deg_uF\)（亦即该分支对 \(\QQ(\alpha)\) 的代数次数）必被 \(6\) 整除。
等价地，中心点邻域的 Perron 分支具有不可约的六分支惯性结构，任何试图以单值解析支“穿越中心点”的方案在代数层面被排除。
\end{conclusion}
\begin{proof}[证明要点]
见推论 \ref{cor:rate-center-perron-degree-multiple-6}；惯性群语义对齐见注 \ref{rem:rate-center-inertia-c6}。
\end{proof}

\paragraph{中心六重临界分歧向有限节点反演病态的传递}
中心六重聚合把参数扰动压缩到六次根尺度之后，Prony/Hankel 型有限窗口反演的不适定性会沿同一尺度被放大为显式幂律爆炸。

\begin{theorem}[六重临界分歧强制 \texorpdfstring{$\kappa$}{kappa}-节点谱反演的普适爆炸率]\label{thm:xi-center-sixfold-kappa-node-inversion-blowup}
沿用结论 \ref{con:xi-rate-center-perron-degree-multiple-6} 的记号，并在 \((\alpha,u)=(\tfrac12,1)\) 的邻域内取包含 Perron 分支的局部解析支。
假设存在常数 \(c_1,c_2>0\) 使得当 \(\alpha\to \tfrac12\) 时有
\[
c_1\left|\alpha-\frac12\right|^{1/6}
\le
|u-1|
\le
c_2\left|\alpha-\frac12\right|^{1/6}.
\]

固定 \(\kappa\ge 2\)，并考虑由长度 \(2\kappa-1\) 的 Prony/Hankel 数据恢复 \(\kappa\) 个谱节点
\[
r_1(u),\dots,r_\kappa(u)
\]
的局部反演问题。
记最小分离度为
\[
\operatorname{sep}(u):=\min_{j\neq k}|r_j(u)-r_k(u)|.
\]
若各成对节点在 \(u=1\) 处都横截裂分，即对每个 \(j\neq k\) 都存在 \(c_{jk}\neq 0\) 使
\[
r_j(u)-r_k(u)=c_{jk}(u-1)+O\bigl((u-1)^2\bigr),
\]
则任意局部 Lipschitz 右逆的 Lipschitz 常数 \(L(\alpha)\) 都满足下界
\[
L(\alpha)\ge C_\kappa\left|\alpha-\frac12\right|^{-(\kappa-1)/6},
\]
其中 \(C_\kappa>0\) 仅依赖于 \(\kappa\)、采样规范以及局部横截常数。
\end{theorem}
\begin{proof}
由各成对差值的横截裂分，存在常数 \(c_3,c_4>0\) 使当 \(u\) 充分接近 \(1\) 时，
\[
c_3|u-1|
\le
\operatorname{sep}(u)
\le
c_4|u-1|.
\]
再与假设的六重临界分歧尺度比较合并，得到
\[
\operatorname{sep}(u)\asymp \left|\alpha-\frac12\right|^{1/6}.
\]

另一方面，命题 \ref{con:xi-collision-illposedness-sep-blowup} 已给出任意基于长度 \(2\kappa-1\) 有限样本的 Prony/Hankel 型局部反演都满足
\[
L\ge C'_\kappa\,\operatorname{sep}(u)^{-(\kappa-1)},
\]
其中 \(C'_\kappa>0\) 仅依赖于 \(\kappa\) 与采样规范。
将上式中的 \(\operatorname{sep}(u)\) 以 \(|\alpha-\tfrac12|^{1/6}\) 代换，并把隐含比较常数吸收入新的常数 \(C_\kappa\)，即得
\[
L(\alpha)\ge C_\kappa\left|\alpha-\frac12\right|^{-(\kappa-1)/6}.
\]
证毕。
\end{proof}
\endinput


\endinput


\begin{conclusion}[六倍窗碰撞补余的 Lucas 屏障]\label{con:xi-fold-window6-lucas-barrier-collision-gap}
若 \(m+2\equiv 0\pmod 6\)，并记
\[
d:=\frac{m+2}{2},
\qquad
M:=F_{m+2}=F_{2d}.
\]
则有严格下界
\[
\boxed{\;
1-\mathsf{Col}_m
\ge \frac{|\mathcal Z_m|}{M}
\ge \frac{F_d}{F_{2d}}
=\frac{1}{L_d}\;},
\]
其中 \(L_d:=F_{d-1}+F_{d+1}\) 为第 \(d\) 个 Lucas 数。
特别地，
\[
\boxed{\;
\liminf_{\substack{m\to\infty\\ m+2\equiv 0\ (\mathrm{mod}\ 6)}}
\varphi^{(m+2)/2}\bigl(1-\mathsf{Col}_m\bigr)\ge 1\; }.
\]
\end{conclusion}
\begin{proof}[证明要点]
由定理 \ref{thm:fold-collision-zero-reduction}，
\[
1-\mathsf{Col}_m\ge \frac{|\mathcal Z_m|}{M}.
\]
另一方面，结论 \ref{con:xi-fold-window6-kernel-half-index-lower-bound} 给出
\[
|\mathcal Z_m|\ge F_d.
\]
两式合并得到
\[
1-\mathsf{Col}_m\ge \frac{F_d}{F_{2d}}.
\]
再由 Fibonacci--Lucas 恒等式
\[
F_{2d}=F_dL_d
\]
可化为
\[
1-\mathsf{Col}_m\ge \frac{1}{L_d}.
\]
最后由 Binet 公式
\[
L_d=\varphi^d+(-\varphi)^{-d}
=\varphi^d\bigl(1+O(\varphi^{-2d})\bigr)
\]
可得
\[
\frac{1}{L_d}
=\varphi^{-d}\bigl(1+O(\varphi^{-2d})\bigr).
\]
将 \(d=(m+2)/2\) 代回并取下极限，即得最后一式。证毕。
\end{proof}


\begin{conclusion}[临界共振常数强制纤维谱的统一二次下界与峰值盈余]\label{con:xi-fold-critical-resonance-quadratic-and-peak-floor}
沿用定理 \ref{thm:fold-collision-spectrum}、定义 \ref{def:fold-normalized-amplitude} 与定理 \ref{thm:fold-max-fiber-fourier} 的记号，令
\[
\overline d_m:=\frac{2^m}{F_{m+2}},
\qquad
\mathsf{Col}_m:=\frac{1}{2^{2m}}\sum_{r} d_m(r)^2,
\qquad
M_m:=\max_r d_m(r).
\]
设临界共振常数 \(C_\varphi>0\) 如定理 \ref{thm:fold-critical-resonance-constant}。则有统一下界
\[
\liminf_{m\to\infty}
\frac{F_{m+2}}{2^{2m}}
\sum_r\bigl(d_m(r)-\overline d_m\bigr)^2
\ge 2C_\varphi^2,
\]
\[
\liminf_{m\to\infty}F_{m+2}\,\mathsf{Col}_m
\ge 1+2C_\varphi^2,
\]
\[
\liminf_{m\to\infty}\frac{F_{m+2}}{2^m}\,M_m
\ge 1+C_\varphi.
\]
因此，仅由临界共振模 \(k=F_m,F_{m+1}\) 即可强制一个与 \(m\) 无关的统一碰撞盈余与统一峰值盈余；均匀纤维极限在二次能量与 \(L^\infty\) 两个层面同时被永久排除。
\end{conclusion}
\begin{proof}[证明]
由定理 \ref{thm:fold-collision-spectrum} 的 Parseval 恒等式，
\[
\frac{F_{m+2}}{2^{2m}}
\sum_r\bigl(d_m(r)-\overline d_m\bigr)^2
=\sum_{k\ne 0} a_m(k)^2.
\]
右端至少包含两项 \(a_m(F_m)^2+a_m(F_{m+1})^2\)，而定理 \ref{thm:fold-critical-resonance-constant} 给出
\(
a_m(F_m),a_m(F_{m+1})\to C_\varphi
\)；
故第一条不等式成立。

再由同一 Parseval 恒等式与 \(a_m(0)=1\) 可得
\[
F_{m+2}\,\mathsf{Col}_m
=\sum_k a_m(k)^2
=1+\frac{F_{m+2}}{2^{2m}}
\sum_r\bigl(d_m(r)-\overline d_m\bigr)^2,
\]
故第二条不等式由第一条立即推出。

最后，由定理 \ref{thm:fold-max-fiber-fourier}，
\[
M_m\ge \overline d_m+\frac{S_m}{F_{m+2}},
\qquad
S_m:=\max_{k\ne 0}|\widehat d_m(k)|.
\]
又因为
\[
S_m\ge 2^m\max\{a_m(F_m),a_m(F_{m+1})\},
\]
故
\[
\frac{F_{m+2}}{2^m}M_m
\ge 1+\max\{a_m(F_m),a_m(F_{m+1})\}.
\]
令 \(m\to\infty\) 并再次使用定理 \ref{thm:fold-critical-resonance-constant} 即得第三条。证毕。
\end{proof}


\subsubsection{负谱体积势、伸缩平移与两态退化：Pick--Poisson 不变量、Hankel 同源病态与可复算 DP 核}
本段把 Pick--Poisson 负谱体积势的完备不变量、统一伸缩下的 \(\log m\) 平移化、Hankel/Toeplitz 两类反演的同源病态，以及 window-\(6\)/bin-fold 侧的离散选择、两态闭式化与尾谱驱动的精确反演门槛，组织为一组可直接引用的条目；其证明入口均可由本文既有定理/命题精确定位。

\begin{conclusion}[负谱行列式的完备不变量：Poisson 权重与伪双曲 Vandermonde 的乘积极限]\label{con:xi-terminal-negative-spectrum-det-invariant}
设点状缺陷由有限 Blaschke 零点多重集 \(\{a_j\}_{j=1}^{\kappa}\subset\mathbb{D}\) 给出，并令 \(\mathsf P\) 为相应 Pick--Poisson Gram 矩阵 \eqref{eq:xi-pick-poisson-gram}。
令伪双曲距离
\[
\rho(a,b):=\left|\frac{a-b}{1-\overline b\,a}\right|\qquad(a,b\in\mathbb{D}).
\]
则行列式具有闭式分解
\[
\boxed{\;
\det\mathsf P
=
\Bigl(\prod_{j=1}^{\kappa}\frac{1-|a_j|^2}{|1+a_j|^2}\Bigr)
\Bigl(\prod_{1\le j<k\le\kappa}\rho(a_j,a_k)^2\Bigr)\; }.
\]
并且当 Toeplitz 主子矩阵 \(\mathbf{T}_N\) 进入稳定显影区间、其负特征值（按重数）为 \(\{\lambda_j^{-}(\mathbf{T}_N)\}_{j=1}^{\kappa}\) 时，有严格乘积极限
\[
\boxed{\;
\lim_{N\to\infty}\ \prod_{j=1}^{\kappa}\bigl(-\lambda_j^{-}(\mathbf{T}_N)\bigr)=\det\mathsf P\; }.
\]
该不变量对贴边（\(|a_j|\uparrow 1\)）与碰撞（\(\rho(a_j,a_k)\downarrow 0\)）同时敏感，因而在半经典意义下完整编码点状缺陷几何。
\end{conclusion}
\begin{proof}[证明要点]
见命题 \ref{prop:xi-pick-poisson-det-factorization}。
\end{proof}

\begin{conclusion}[统一伸缩的协变律：显影复杂度发生 \texorpdfstring{$\log m$}{log m} 平移]\label{con:xi-terminal-visibility-complexity-logm-translation}
沿用上一条结论的有限缺陷设定，并令 \(\phi_m(w):=\mathcal{C}(m\,\mathcal{C}^{-1}(w))\)（\(m>0\)，固定端点 \(-1\)）为统一伸缩。
令 \(a_j(m):=\phi_m(a_j)\)，并令 \(\mathsf P(m)\) 为由 \(\{a_j(m)\}\) 通过 \eqref{eq:xi-pick-poisson-gram} 生成的 Pick--Poisson Gram 矩阵。
定义边界熵势、边界体积势与显影复杂度
\[
S_{\partial}:=-\log\det\mathsf P,
\qquad
V_{\partial}:=\log\mathrm{tr}(\mathsf P),
\qquad
\mathscr{C}_{\partial}:=S_{\partial}+(\kappa-1)V_{\partial},
\]
并以同式定义 \(S_{\partial}(m),V_{\partial}(m),\mathscr{C}_{\partial}(m)\)。
则对一切 \(m>0\) 成立严格协变律
\[
\boxed{\;
\mathscr{C}_{\partial}(m)=\mathscr{C}_{\partial}(1)+\log m\; }.
\]
因此尺度参数在 \(\mathscr{C}_{\partial}\) 上被线性化为加法坐标；该平移化与 Witt 伸缩在 Koenigs 深度坐标 \(y\) 上的平移恒等式 \(y(ms)=y(s)+\log m\)（定理 \ref{thm:koenigs-linearization-witt-dilation}）在机制上对齐。
\end{conclusion}
\begin{proof}[证明要点]
由命题 \ref{prop:xi-pick-poisson-det-scaling-law}，\(\rho\) 在 \(\phi_m\) 下不变且 \(P_{\phi_m(a)}(-1)=m^{-1}P_a(-1)\)，从而 \(\det\mathsf P(m)=m^{-\kappa}\det\mathsf P\) 且 \(\mathrm{tr}(\mathsf P(m))=m^{-1}\mathrm{tr}(\mathsf P)\)。
代入 \(\mathscr{C}_{\partial}=S_{\partial}+(\kappa-1)V_{\partial}\) 即得结论；详见推论 \ref{cor:xi-toeplitz-visibility-complexity-anomaly}。
\end{proof}

\begin{conclusion}[统一伸缩下 Pick--Poisson Gram 谱的严格同伦缩放]\label{con:xi-terminal-pick-poisson-homothetic-spectrum-scaling}
沿用结论 \ref{con:xi-terminal-visibility-complexity-logm-translation} 的统一伸缩 \(\phi_m\) 与 Gram 矩阵 \(\mathsf P(m)\) 的记号，并令 \(\chi_{\mathsf P}(\lambda):=\det(\lambda I_\kappa-\mathsf P)\)。
则特征多项式满足严格协变律
\[
\boxed{\;
\chi_{\mathsf P(m)}(\lambda)=m^{-\kappa}\chi_{\mathsf P}(m\lambda)\; } .
\]
从而按重数计有谱严格缩放
\[
\boxed{\;\mathrm{Spec}(\mathsf P(m))=m^{-1}\mathrm{Spec}(\mathsf P).\;}
\]
\end{conclusion}
\begin{proof}
由命题 \ref{prop:xi-pick-poisson-det-scaling-law}，\(\rho\) 在 \(\phi_m\) 下不变且端点 Poisson 权重满足 \(p_j(m)=m^{-1}p_j\)。
因此对任意指标子集 \(I\subset[\kappa]\)，命题 \ref{prop:xi-pick-poisson-principal-minors-partition} 的闭式给出
\(
\det(\mathsf P_I(m))=m^{-|I|}\det(\mathsf P_I)
\)。
令 \(Z_r\) 与 \(Z_r(m)\) 为命题 \ref{prop:xi-pick-poisson-charpoly-coefficients} 中的 \(r\)-阶主子式配分函数，则
\(
Z_r(m)=m^{-r}Z_r
\)。
代入 \(\chi_{\mathsf P}(\lambda)=\sum_{r=0}^\kappa (-1)^r Z_r \lambda^{\kappa-r}\) 得
\[
\chi_{\mathsf P(m)}(\lambda)
=\sum_{r=0}^\kappa (-1)^r m^{-r}Z_r\,\lambda^{\kappa-r}
=m^{-\kappa}\sum_{r=0}^\kappa(-1)^r Z_r\,(m\lambda)^{\kappa-r}
=m^{-\kappa}\chi_{\mathsf P}(m\lambda).
\]
根随变量按 \(m^{-1}\) 缩放，遂得谱缩放结论。证毕。
\end{proof}

\begin{conclusion}[全部谱函数的重整化恒等式]\label{con:xi-terminal-pick-poisson-spectral-function-renormalization}
在结论 \ref{con:xi-terminal-pick-poisson-homothetic-spectrum-scaling} 的设定下，对任意在 \(\mathrm{Spec}(\mathsf P)\) 的邻域内解析的函数 \(f\)，有
\[
\boxed{\;
\mathrm{tr}\,f(\mathsf P(m))=\mathrm{tr}\,f(m^{-1}\mathsf P)\;}
\]
并且对任意复参数 \(z\) 有
\[
\boxed{\;
\det(I+z\mathsf P(m))=\det\!\bigl(I+(z/m)\mathsf P\bigr)\; } .
\]
\end{conclusion}
\begin{proof}
由谱缩放，存在重排使 \(\lambda_j(\mathsf P(m))=m^{-1}\lambda_j(\mathsf P)\)。
于是 \(\mathrm{tr}\,f(\mathsf P(m))=\sum_j f(m^{-1}\lambda_j(\mathsf P))=\mathrm{tr}\,f(m^{-1}\mathsf P)\)。
行列式恒等式同理由 \(\det(I+z\mathsf P(m))=\prod_j(1+z\lambda_j(\mathsf P(m)))\) 得到。证毕。
\end{proof}

\begin{conclusion}[尺度不变量与平移型势的规范唯一性]\label{con:xi-terminal-pick-poisson-scale-invariant-and-affine-uniqueness}
沿用结论 \ref{con:xi-terminal-visibility-complexity-logm-translation} 中的
\(
S_{\partial}=-\log\det\mathsf P
\)、
\(
V_{\partial}=\log\mathrm{tr}(\mathsf P)
\)
与 \(\phi_m\) 的记号。
定义尺度不变量
\[
\boxed{\;
I_{\partial}:=S_{\partial}+\kappa V_{\partial}
=-\log\!\left(\frac{\det\mathsf P}{(\mathrm{tr}\,\mathsf P)^{\kappa}}\right)\; } .
\]
则对一切 \(m>0\) 有严格不变性 \(I_{\partial}(m)\equiv I_{\partial}(1)\)。
\par
更一般地，若线性组合 \(F:=aS_{\partial}+bV_{\partial}\) 满足对任意配置与任意 \(m>0\) 皆有
\[
F(m)=F(1)+c\log m\qquad(c\ \text{与配置无关}),
\]
则必有 \(c=a\kappa-b\)。
因此：
\begin{enumerate}
  \item \(c=0\) 当且仅当 \((a,b)\) 与 \((1,\kappa)\) 成比例（即 \(F\) 与 \(I_{\partial}\) 成比例）；
  \item \(c=a\) 当且仅当 \((a,b)\) 与 \((1,\kappa-1)\) 成比例（即 \(F\) 与 \(\mathscr C_{\partial}=S_{\partial}+(\kappa-1)V_{\partial}\) 成比例）。
\end{enumerate}
\end{conclusion}
\begin{proof}
由命题 \ref{prop:xi-pick-poisson-det-scaling-law} 得
\(
S_{\partial}(m)=S_{\partial}(1)+\kappa\log m
\)、
\(
V_{\partial}(m)=V_{\partial}(1)-\log m
\)。
故
\(
I_{\partial}(m)=I_{\partial}(1)
\)
且
\(
F(m)=F(1)+(a\kappa-b)\log m
\)，从而 \(c=a\kappa-b\)。
两条“当且仅当”分别对应 \(a\kappa-b=0\) 与 \(a\kappa-b=a\)。证毕。
\end{proof}

\begin{conclusion}[特征行列式的齐次性与一阶重整化方程]\label{con:xi-terminal-pick-poisson-characteristic-homogeneity-pde}
设
\[
\Psi(z,m):=\det(I-z\mathsf P(m)).
\]
则有严格齐次性
\[
\boxed{\;\Psi(z,m)=\Psi(z/m,1)\;}
\]
并因此满足一阶重整化方程
\[
\boxed{\;
\bigl(m\partial_m+z\partial_z\bigr)\Psi(z,m)=0,\qquad
\bigl(m\partial_m+z\partial_z\bigr)\log\Psi(z,m)=0\; } .
\]
\end{conclusion}
\begin{proof}
由结论 \ref{con:xi-terminal-pick-poisson-spectral-function-renormalization} 取 \(f(x)=\log(1-zx)\)（或直接对谱乘积）得
\(
\det(I-z\mathsf P(m))=\det(I-(z/m)\mathsf P)
\)，即 \(\Psi(z,m)=\Psi(z/m,1)\)。
对该恒等式对 \(m,z\) 求导并消去 \(\partial_{(\cdot)}\Psi(z/m,1)\) 即得所示偏微分方程。证毕。
\end{proof}

\begin{corollary}[零点集的严格缩放]\label{cor:xi-pick-poisson-characteristic-zero-set-scaling}
在结论 \ref{con:xi-terminal-pick-poisson-characteristic-homogeneity-pde} 的设定下，定义零点集
\[
Z(m):=\{z\in\CC:\ \Psi(z,m)=0\}.
\]
则对一切 \(m>0\) 有严格比例律
\[
Z(m)=m\cdot Z(1).
\]
\end{corollary}
\begin{proof}
由齐次性 \(\Psi(z,m)=\Psi(z/m,1)\)，有
\[
z\in Z(m)\ \Longleftrightarrow\ \Psi(z,m)=0
\Longleftrightarrow\ \Psi(z/m,1)=0
\Longleftrightarrow\ z/m\in Z(1),
\]
故 \(Z(m)=mZ(1)\)。证毕。
\end{proof}

\begin{corollary}[谱倒数的对数平移不变量]\label{cor:xi-pick-poisson-spectrum-log-translation-invariant}
在上述设定下，令 \(\{\lambda_j(m)\}_{j=1}^{\kappa}=\mathrm{Spec}(\mathsf P(m))\)（计重数）。
则按重数计有
\[
\Bigl\{\frac{1}{\lambda_j(m)}\Bigr\}_{j=1}^{\kappa}=Z(m)=mZ(1)
=\Bigl\{\frac{m}{\lambda_j(1)}\Bigr\}_{j=1}^{\kappa},
\]
从而
\[
\lambda_j(m)=\frac{1}{m}\lambda_j(1)\qquad(1\le j\le\kappa),
\]
并且谱的对数计数测度
\[
\mu_m:=\frac{1}{\kappa}\sum_{j=1}^{\kappa}\delta_{\log \lambda_j(m)}
\]
满足推送恒等式
\[
\mu_m=\Bigl(x\mapsto x-\log m\Bigr)_{\#}\mu_1.
\]
\end{corollary}
\begin{proof}
由 \(\Psi(z,m)=\det(I-z\mathsf P(m))=\prod_{j=1}^{\kappa}(1-z\lambda_j(m))\)，其零点（计重数）恰为 \(\{1/\lambda_j(m)\}_{j=1}^{\kappa}\)。
结合推论 \ref{cor:xi-pick-poisson-characteristic-zero-set-scaling} 得 \(\{1/\lambda_j(m)\}=m\{1/\lambda_j(1)\}\)，逐项取倒数即得 \(\lambda_j(m)=m^{-1}\lambda_j(1)\)；
对数推送公式由 \(\log\lambda_j(m)=\log\lambda_j(1)-\log m\) 直接推出。证毕。
\end{proof}

\begin{conclusion}[Schur 补账本的逐项同伦缩放]\label{con:xi-terminal-pick-poisson-schur-ledger-homothetic-scaling}
沿用结论 \ref{con:xi-terminal-pick-poisson-schur-ledger} 中的 Schur 补通量 \(\pi_m\)。
在统一伸缩 \(\phi_m\) 下，用同一顺序定义 \(\pi_r(m)\)。
则对每个 \(r=1,\dots,\kappa\) 有严格缩放律
\[
\boxed{\;\pi_r(m)=m^{-1}\pi_r(1),\qquad -\log\pi_r(m)=-\log\pi_r(1)+\log m.\;}
\]
\end{conclusion}
\begin{proof}
由 Schur 补链式分解，
\(
\pi_r=\det\mathsf P^{(r)}/\det\mathsf P^{(r-1)}
\)
（其中 \(\mathsf P^{(r)}\) 为前 \(r\) 个点的主子矩阵）。
而命题 \ref{prop:xi-pick-poisson-det-scaling-law} 的协变律同样适用于任意主子矩阵，给出
\(
\det\mathsf P^{(r)}(m)=m^{-r}\det\mathsf P^{(r)}
\)。
代入比值即得 \(\pi_r(m)=m^{-1}\pi_r\)，取对数得到第二式。证毕。
\end{proof}

\begin{conclusion}[尺度中心化账本的严格不变量分解]\label{con:xi-terminal-pick-poisson-centered-ledger-invariants}
定义尺度中心化通量增量
\[
\eta_r:=-\log\pi_r+V_{\partial}\qquad(1\le r\le\kappa).
\]
则对一切 \(m>0\) 与一切 \(r\) 有严格尺度不变性 \(\eta_r(m)=\eta_r(1)\)。
并且满足严格求和分解
\[
\boxed{\;
\sum_{r=1}^{\kappa}\eta_r
=\sum_{r=1}^{\kappa}(-\log\pi_r)+\kappa V_{\partial}
=S_{\partial}+\kappa V_{\partial}
=I_{\partial}\; } .
\]
\end{conclusion}
\begin{proof}
由结论 \ref{con:xi-terminal-pick-poisson-schur-ledger-homothetic-scaling}，
\(
-\log\pi_r(m)=-\log\pi_r(1)+\log m
\)；
由命题 \ref{prop:xi-pick-poisson-det-scaling-law}，
\(
V_{\partial}(m)=V_{\partial}(1)-\log m
\)。
两式相加即得 \(\eta_r(m)=\eta_r(1)\)。
求和恒等式由 \(\det\mathsf P=\prod_r\pi_r\) 与 \(I_{\partial}=S_{\partial}+\kappa V_{\partial}\) 直接推出。证毕。
\end{proof}

\begin{conclusion}[尺度不变量的谱型 KL 表示与等号刻画]\label{con:xi-terminal-pick-poisson-scale-invariant-spectral-kl}
设 \(\mathsf P\succ 0\) 的特征值为 \(\lambda_1,\dots,\lambda_\kappa>0\)，并令
\[
p_\Sigma:=\mathrm{tr}(\mathsf P)=\sum_{j=1}^{\kappa}\lambda_j,
\qquad
\mu_j:=\frac{\lambda_j}{p_\Sigma}\ \ (1\le j\le\kappa).
\]
记 \(U_\kappa\) 为 \(\{1,\dots,\kappa\}\) 上的均匀分布，则尺度不变量满足严格恒等式
\[
\boxed{\;
I_{\partial}
\;=\;
\kappa\log\kappa+\kappa\,D_{\mathrm{KL}}(U_\kappa\Vert \mu)\; } ,
\]
其中
\[
D_{\mathrm{KL}}(U_\kappa\Vert \mu)
:=\sum_{j=1}^{\kappa}\frac1\kappa\log\frac{1/\kappa}{\mu_j}\ \ge\ 0.
\]
从而得到谱不等式
\[
\boxed{\;
\det\mathsf P\ \le\ \left(\frac{\mathrm{tr}(\mathsf P)}{\kappa}\right)^{\kappa}\; } ,
\]
且等号当且仅当 \(\mu_1=\cdots=\mu_\kappa=\frac1\kappa\)，亦即
\[
\mathsf P=\frac{\mathrm{tr}(\mathsf P)}{\kappa}\,I_\kappa .
\]
\end{conclusion}
\begin{proof}
由 \(\det\mathsf P=\prod_{j=1}^{\kappa}\lambda_j=(p_\Sigma)^\kappa\prod_{j=1}^{\kappa}\mu_j\) 得
\[
I_{\partial}
:=-\log\!\left(\frac{\det\mathsf P}{(\mathrm{tr}\,\mathsf P)^{\kappa}}\right)
=-\sum_{j=1}^{\kappa}\log\mu_j.
\]
另一方面，
\[
D_{\mathrm{KL}}(U_\kappa\Vert \mu)
=\sum_{j=1}^{\kappa}\frac1\kappa\log\frac{1/\kappa}{\mu_j}
=-\log\kappa-\frac1\kappa\sum_{j=1}^{\kappa}\log\mu_j,
\]
故 \(-\sum_j\log\mu_j=\kappa\log\kappa+\kappa D_{\mathrm{KL}}(U_\kappa\Vert \mu)\)，代回即得所示恒等式。
由 \(D_{\mathrm{KL}}\ge 0\) 得 \(-\sum_j\log\mu_j\ge \kappa\log\kappa\)，等价于 \(\prod_j\mu_j\le \kappa^{-\kappa}\)，从而得到 \(\det\mathsf P\le (\mathrm{tr}(\mathsf P)/\kappa)^\kappa\)。
等号条件 \(D_{\mathrm{KL}}=0\) 等价于 \(\mu\) 为均匀分布，亦即 \(\mathsf P\) 的谱完全平坦，从而 \(\mathsf P=(\mathrm{tr}(\mathsf P)/\kappa)I_\kappa\)。证毕。
\end{proof}

\begin{conclusion}[显影复杂度的 KL--互能--标度三项分解]\label{con:xi-terminal-pick-poisson-visibility-complexity-kl-energy-split}
记端点 Poisson 权重
\(
p_j:=\mathsf P_{jj}=P_{a_j}(-1)
\)
并令 \(p_\Sigma=\sum_{j=1}^{\kappa}p_j=\mathrm{tr}(\mathsf P)\)，再定义归一化权重分布
\[
q_j:=\frac{p_j}{p_\Sigma}\qquad(1\le j\le\kappa).
\]
则显影复杂度
\(
\mathscr{C}_{\partial}=S_{\partial}+(\kappa-1)V_{\partial}
\)
满足严格恒等式
\[
\boxed{\;
\mathscr{C}_{\partial}
=\kappa\log\kappa-\log p_\Sigma
\;+\;\kappa\,D_{\mathrm{KL}}(U_\kappa\Vert q)
\;+\;2\sum_{1\le j<k\le\kappa}\bigl(-\log\rho(a_j,a_k)\bigr)\; } .
\]
其中第一项为纯标度项，第二项为端点权重的均匀性偏差，第三项为伪双曲互能。
\end{conclusion}
\begin{proof}
由结论 \ref{con:xi-terminal-negative-spectrum-det-invariant} 的行列式分解
\[
\det\mathsf P=\Bigl(\prod_{j=1}^{\kappa}p_j\Bigr)\Bigl(\prod_{j<k}\rho(a_j,a_k)^2\Bigr)
\]
得
\[
S_{\partial}=-\log\det\mathsf P
=-\sum_{j=1}^{\kappa}\log p_j
2\sum_{j<k}\bigl(-\log\rho(a_j,a_k)\bigr).
\]
于是
\[
\mathscr{C}_{\partial}
=S_{\partial}+(\kappa-1)\log p_\Sigma
=-\sum_{j=1}^{\kappa}\log p_j+(\kappa-1)\log p_\Sigma
2\sum_{j<k}\bigl(-\log\rho(a_j,a_k)\bigr).
\]
又 \(-\sum_j\log p_j=-\kappa\log p_\Sigma-\sum_j\log q_j\)，故
\[
\mathscr{C}_{\partial}
=-\log p_\Sigma-\sum_{j=1}^{\kappa}\log q_j
2\sum_{j<k}\bigl(-\log\rho(a_j,a_k)\bigr).
\]
最后
\[
D_{\mathrm{KL}}(U_\kappa\Vert q)
=-\log\kappa-\frac1\kappa\sum_{j=1}^{\kappa}\log q_j
\]
给出 \(-\sum_j\log q_j=\kappa\log\kappa+\kappa D_{\mathrm{KL}}(U_\kappa\Vert q)\)，代回即得结论。证毕。
\end{proof}

\begin{conclusion}[Schur 外部化的归一化乘积恒等式与有限步塌缩见证]\label{con:xi-terminal-pick-poisson-schur-vandermonde-externalization}
对任意排列 \(\sigma\in\mathfrak S_\kappa\)，令 \(\mathsf P^{(m)}_\sigma\) 为按顺序 \(\sigma(1),\dots,\sigma(m)\) 取出的主子矩阵，并定义相应 Schur 补通量
\[
\pi_m(\sigma)
:=\det(\mathsf P^{(m)}_\sigma)\Big/\det(\mathsf P^{(m-1)}_\sigma)
\qquad(1\le m\le\kappa),
\]
其中约定 \(\det(\mathsf P^{(0)}_\sigma)=1\)。
则满足严格归一化恒等式
\[
\boxed{\;
\prod_{m=1}^{\kappa}\frac{\pi_m(\sigma)}{p_{\sigma(m)}}
=\prod_{1\le j<k\le\kappa}\rho(a_j,a_k)^2\; } .
\]
\end{conclusion}
\begin{proof}
由 Schur 补链式分解可得 \(\det\mathsf P=\prod_{m=1}^{\kappa}\pi_m(\sigma)\)（对置换后的矩阵施行同一分解）。
另一方面，结论 \ref{con:xi-terminal-negative-spectrum-det-invariant} 给出
\(
\det\mathsf P=(\prod_j p_j)(\prod_{j<k}\rho(a_j,a_k)^2)
\)。
将两式相除并注意 \(\prod_m p_{\sigma(m)}=\prod_j p_j\)，得到归一化乘积恒等式。
\par
证毕。
\end{proof}

\begin{proposition}[Schur 消元熵的排列不变性与瓶颈步界]\label{prop:xi-schur-elimination-entropy-permutation-invariant-bottleneck}
\leanverified{paper\_xi\_schur\_elimination\_entropy\_permutation\_invariant\_bottleneck}
沿用结论 \ref{con:xi-terminal-pick-poisson-schur-vandermonde-externalization} 的设定与记号。
对任意排列 \(\sigma\in\mathfrak S_\kappa\)，定义 \textbf{Schur 消元熵}
\[
\mathcal E_{\mathrm{elim}}
:=\frac1\kappa\sum_{m=1}^{\kappa}\log\frac{p_{\sigma(m)}}{\pi_m(\sigma)}.
\]
则：
\begin{enumerate}
  \item \(\mathcal E_{\mathrm{elim}}\) 与排列 \(\sigma\) 无关，并满足闭式
  \[
  \mathcal E_{\mathrm{elim}}
  =\frac{-1}{\kappa}\log\!\Bigl(\prod_{1\le j<k\le\kappa}\rho(a_j,a_k)^2\Bigr)
  =\frac{2}{\kappa}\sum_{1\le j<k\le\kappa}\log\frac{1}{\rho(a_j,a_k)}.
  \]
  \item 对任意排列 \(\sigma\)，必存在某一步 \(m_\star\in\{1,\dots,\kappa\}\) 使
  \[
  \frac{\pi_{m_\star}(\sigma)}{p_{\sigma(m_\star)}}
  \ \le\
  \exp\!\bigl(-\mathcal E_{\mathrm{elim}}\bigr)
  =\Bigl(\prod_{1\le j<k\le\kappa}\rho(a_j,a_k)^2\Bigr)^{1/\kappa}.
  \]
\end{enumerate}
\end{proposition}
\begin{proof}
由结论 \ref{con:xi-terminal-pick-poisson-schur-vandermonde-externalization} 的乘积恒等式，
\(
\prod_{m=1}^{\kappa}\frac{\pi_m(\sigma)}{p_{\sigma(m)}}=\prod_{j<k}\rho(a_j,a_k)^2
\)；
取对数并除以 \(\kappa\) 得
\[
\frac1\kappa\sum_{m=1}^{\kappa}\log\frac{p_{\sigma(m)}}{\pi_m(\sigma)}
=\frac{-1}{\kappa}\log\!\Bigl(\prod_{j<k}\rho(a_j,a_k)^2\Bigr)
=\frac{2}{\kappa}\sum_{j<k}\log\frac1{\rho(a_j,a_k)},
\]
从而 \(\mathcal E_{\mathrm{elim}}\) 与 \(\sigma\) 无关并满足闭式。
\par
对数列 \(\{\log(p_{\sigma(m)}/\pi_m(\sigma))\}_{m=1}^{\kappa}\) 取平均，存在某个 \(m_\star\) 使该项不小于平均值 \(\mathcal E_{\mathrm{elim}}\)；
指数化即得所示瓶颈步界。证毕。
\end{proof}

\begin{conclusion}[Toeplitz 负谱几何均值的严格同伦缩放]\label{con:xi-terminal-toeplitz-negative-geometric-mean-homothety}
设 Toeplitz 主子矩阵 \(\mathbf{T}_N\) 进入稳定显影区间且其负特征值（按重数）为 \(\{\lambda_j^{-}(\mathbf{T}_N)\}_{j=1}^{\kappa}\)，则
\[
\lim_{N\to\infty}\ \prod_{j=1}^{\kappa}\bigl(-\lambda_j^{-}(\mathbf{T}_N)\bigr)=\det\mathsf P.
\]
对统一伸缩后的缺陷数据，记相应 Toeplitz 系列为 \(\mathbf{T}_N(m)\)，其负特征值为 \(\{\lambda_j^{-}(\mathbf{T}_N(m))\}_{j=1}^{\kappa}\)。
则负谱几何均值满足严格缩放律
\[
\boxed{\;
\lim_{N\to\infty}\left(\prod_{j=1}^{\kappa}\bigl(-\lambda_j^{-}(\mathbf{T}_N(m))\bigr)\right)^{1/\kappa}
=m^{-1}\,
\lim_{N\to\infty}\left(\prod_{j=1}^{\kappa}\bigl(-\lambda_j^{-}(\mathbf{T}_N)\bigr)\right)^{1/\kappa}\; } .
\]
\end{conclusion}
\begin{proof}
对伸缩后的系统同样有乘积极限，其极限为 \(\det\mathsf P(m)\)（命题 \ref{prop:xi-pick-poisson-det-factorization} 的逐配置应用）。
而命题 \ref{prop:xi-pick-poisson-det-scaling-law} 给出 \(\det\mathsf P(m)=m^{-\kappa}\det\mathsf P\)。
故乘积极限相差因子 \(m^{-\kappa}\)，两边取 \(1/\kappa\) 次方即得结论。证毕。
\end{proof}

\begin{proposition}[簇分离下交叉互能的显式指数上界]\label{prop:xi-pick-poisson-cross-energy-exp-bound}
设缺陷多重集分解为不交并 \(A\sqcup B\)。
令双曲距离
\[
d_{\mathbb{D}}(a,b):=2\,\operatorname{artanh}\rho(a,b)\qquad(a,b\in\mathbb D),
\]
并设簇间最小分离度
\[
L:=\inf\{d_{\mathbb D}(a,b):a\in A,\ b\in B\}.
\]
则当 \(L\ge \log 2\) 时有显式上界
\[
\boxed{\;
0\le -\log\det\mathsf P(A\cup B)+\log\det\mathsf P(A)+\log\det\mathsf P(B)
\le 8|A||B|\,e^{-L}\; } .
\]
\end{proposition}
\begin{proof}
由命题 \ref{prop:xi-pick-poisson-union-submultiplicativity-cross-energy}，
交叉缺口等于
\(
2\sum_{a\in A}\sum_{b\in B}(-\log\rho(a,b))
\)。
由 \(\rho=\tanh(d_{\mathbb D}/2)\)，令 \(q:=e^{-d_{\mathbb D}}\in(0,1)\)，则
\[
-\log\rho=-\log\tanh\!\Bigl(\frac{d_{\mathbb D}}{2}\Bigr)=\log\frac{1+q}{1-q}.
\]
若 \(d_{\mathbb D}\ge L\ge\log 2\)，则 \(q\le e^{-\log2}=1/2\)，从而
\[
\log\frac{1+q}{1-q}\le q+\frac{q}{1-q}\le 4q\le 4e^{-L}.
\]
代回求和得到
\(
2\sum_{a,b}(-\log\rho(a,b))\le 2|A||B|\cdot 4e^{-L}
\)，即得结论。证毕。
\end{proof}

\begin{corollary}[簇分离下尺度不变量的近可加性]\label{cor:xi-pick-poisson-scale-invariant-near-additivity}
在命题 \ref{prop:xi-pick-poisson-cross-energy-exp-bound} 的条件下，尺度不变量
\(
I_{\partial}(X)=S_{\partial}(X)+|X|\,V_{\partial}(X)
\)
满足估计
\[
\boxed{\;
\bigl|I_{\partial}(A\cup B)-I_{\partial}(A)-I_{\partial}(B)\bigr|
\le 8|A||B|\,e^{-L}
\;+\;(|A|+|B|)\left|\log\frac{\mathrm{tr}\,\mathsf P(A)\,\mathrm{tr}\,\mathsf P(B)}{\mathrm{tr}\,\mathsf P(A\cup B)}\right|\; } .
\]
特别地，若进一步满足 \(\mathrm{tr}\,\mathsf P(A\cup B)=\mathrm{tr}\,\mathsf P(A)+\mathrm{tr}\,\mathsf P(B)\)，则右端第二项为零，\(I_{\partial}\) 的偏离仅由 \(e^{-L}\) 控制。
\end{corollary}
\begin{proof}
由定义
\(
I_{\partial}=S_{\partial}+\kappa V_{\partial}
\)
且 \(|A\cup B|=|A|+|B|\)，有恒等分解
\[
I_{\partial}(A\cup B)-I_{\partial}(A)-I_{\partial}(B)
=\bigl[S_{\partial}(A\cup B)-S_{\partial}(A)-S_{\partial}(B)\bigr]
\]
\[
\qquad
+(|A|+|B|)\log\mathrm{tr}\,\mathsf P(A\cup B)-|A|\log\mathrm{tr}\,\mathsf P(A)-|B|\log\mathrm{tr}\,\mathsf P(B).
\]
第一括号由命题 \ref{prop:xi-pick-poisson-cross-energy-exp-bound} 上界控制；
第二项取绝对值并提取 \((|A|+|B|)\) 即得所示估计。证毕。
\end{proof}

\endinput


\begin{conclusion}[负谱体积势的静电学刻画：\texorpdfstring{$-\log\det\mathsf P$}{-log det P} 等价于上半平面 Dirichlet Green 镜像能量]\label{con:xi-terminal-pick-poisson-green-mirror-energy}
在上半平面坐标 \(z_j=\mathcal{C}^{-1}(a_j)\in\mathbb{H}\) 下，令
\(
t_j:=-1/z_j\in\mathbb{H}
\)。
则端点 Poisson 权重与伪双曲距离满足
\[
P_{a_j}(-1)=\Im(t_j),
\qquad
\rho(a_j,a_k)=\left|\frac{t_j-t_k}{t_j-\overline t_k}\right|.
\]
定义上半平面 Dirichlet Green 函数
\[
G_{\mathbb{H}}(t,s):=\log\left|\frac{t-\overline s}{t-s}\right|.
\]
则广义缺陷熵（负谱体积势）
\(
S_{\mathrm{gen}}:=-\log\det\mathsf P
\)
具有严格等价式
\[
\boxed{\;
S_{\mathrm{gen}}
=\sum_{j=1}^{\kappa}\bigl(-\log\Im(t_j)\bigr)
+2\sum_{1\le j<k\le\kappa}G_{\mathbb{H}}(t_j,t_k)\; }.
\]
并且若缺陷多重集分解为不交并 \(A\sqcup B\)，则交叉能满足恒等式
\[
\boxed{\;
-\log\det\mathsf P(A\cup B)
=-\log\det\mathsf P(A)-\log\det\mathsf P(B)
+2\sum_{a\in A}\sum_{b\in B}\bigl(-\log\rho(a,b)\bigr)\; }.
\]
从而推出次乘法律 \(\det\mathsf P(A\cup B)\le \det\mathsf P(A)\det\mathsf P(B)\)，且交叉项完全由伪双曲分离度控制。
\end{conclusion}
\begin{proof}[证明要点]
像荷镜像能量表示见命题 \ref{prop:xi-pick-poisson-image-charge-green-energy}；
并集交叉互能恒等式与次乘性见命题 \ref{prop:xi-pick-poisson-union-submultiplicativity-cross-energy}。
\end{proof}

\begin{conclusion}[Schur 补链式外部化记账：逐点加入的新增能量等于条件通量损失]\label{con:xi-terminal-pick-poisson-schur-ledger}
对固定点序 \(1,\dots,\kappa\)，记 \(\mathsf P^{(m)}\) 为前 \(m\) 个点诱导的主子矩阵，并定义 Schur 补（条件通量）
\[
\pi_m
:=\mathsf P_{mm}-\mathsf P_{m,1:m-1}\,(\mathsf P^{(m-1)})^{-1}\,\mathsf P_{1:m-1,m},
\qquad
m=1,\dots,\kappa.
\]
则有严格分解
\[
\boxed{\;
\det\mathsf P=\prod_{m=1}^{\kappa}\pi_m,
\qquad
S_{\mathrm{gen}}=-\sum_{m=1}^{\kappa}\log\pi_m\; }.
\]
并且每一 \(\pi_m\) 满足 \(0<\pi_m\le \mathsf P_{mm}=P_{a_m}(-1)\)。
因此负谱体积势具有逐步记账结构：第 \(m\) 个缺陷的新增代价精确等于其在既有缺陷张成空间上的条件通量损失，并被端点自能显式上界控制。
\end{conclusion}
\begin{proof}[证明要点]
见命题 \ref{prop:xi-pick-poisson-schur-flux-ledger}。
\end{proof}

\paragraph{Schur 补账本的份额化与交叉能量外部化}
在结论 \ref{con:xi-terminal-pick-poisson-schur-ledger} 与结论 \ref{con:xi-terminal-pick-poisson-centered-ledger-invariants} 给出的 Schur 补记账与尺度中心化分解基础上，归一化份额提供一组与可审计证书链直接兼容的标量接口。

\setcounter{auditthm}{582}

\begin{auditthm}[Schur 补账本的次概率刻画与离对角缺口]\label{thm:xi-schur-ledger-subprobability-offdiagonal-gap}
令 \(\mathsf P\in\CC^{\kappa\times\kappa}\) 为由缺陷点列 \((a_1,\dots,a_\kappa)\) 生成的 Pick--Poisson Gram 矩阵 \eqref{eq:xi-pick-poisson-gram}，并记端点通量
\[
p_\Sigma:=\mathrm{tr}(\mathsf P)=\sum_{j=1}^{\kappa}p_j,
\qquad
p_j:=\mathsf P_{jj}=P_{a_j}(-1).
\]
对固定顺序 \(1,\dots,\kappa\)，令 \((\pi_m)_{m=1}^{\kappa}\) 为结论 \ref{con:xi-terminal-pick-poisson-schur-ledger} 中的 Schur 补通量，使得
\(
\det\mathsf P=\prod_{m=1}^{\kappa}\pi_m
\)
且 \(0<\pi_m\le p_m\)。
定义尺度中心化增量（结论 \ref{con:xi-terminal-pick-poisson-centered-ledger-invariants}）
\[
\eta_m:=-\log\pi_m+V_\partial,
\qquad
V_\partial:=\log p_\Sigma,
\]
并定义归一化 Schur 份额
\[
s_m:=\frac{\pi_m}{p_\Sigma}=e^{-\eta_m}\qquad(1\le m\le\kappa).
\]
则成立：
\begin{enumerate}
  \item 在统一伸缩下 \(s_m\) 与 \(\eta_m\) 严格不变。
  \item 次概率约束成立：
  \[
  \sum_{m=1}^{\kappa}s_m=\frac{\sum_{m=1}^{\kappa}\pi_m}{p_\Sigma}\le 1,
  \]
  且等号成立当且仅当 \(\mathsf P\) 为对角矩阵（等价地，对一切 \(m\ge 2\) 有 \(\mathsf P_{m,1:m-1}=0\)）。
  \item 若 \(\mathsf P\succ 0\) 且 \(\mathsf P=LL^\ast\) 为与该顺序相容的 Cholesky 分解（\(L\) 为下三角），则
  \[
  \pi_m=\abs{L_{mm}}^2,
  \qquad
  1-\sum_{m=1}^{\kappa}s_m=\frac{\sum_{1\le j<i\le\kappa}\abs{L_{ij}}^2}{\|L\|_F^2}.
  \]
  \item 记
  \(
  I_\partial:=-\log\!\left(\frac{\det\mathsf P}{(\mathrm{tr}\,\mathsf P)^{\kappa}}\right)
  \)
  （结论 \ref{con:xi-terminal-pick-poisson-scale-invariant-and-affine-uniqueness}），则熵缺口满足下界
  \[
  I_\partial
  =-\sum_{m=1}^{\kappa}\log s_m
  \ \ge\
  \kappa\log\!\Bigl(\frac{\kappa}{\sum_{m=1}^{\kappa}s_m}\Bigr)
  \ \ge\
  \kappa\log\kappa.
  \]
\end{enumerate}
\end{auditthm}
\begin{proof}
由结论 \ref{con:xi-terminal-pick-poisson-schur-ledger-homothetic-scaling} 得 \(\pi_m\) 在统一伸缩下按 \(m^{-1}\) 缩放；
由 \(\mathrm{tr}(\mathsf P(m))=m^{-1}\mathrm{tr}(\mathsf P)\)（命题 \ref{prop:xi-pick-poisson-det-scaling-law} 的迹协变）得 \(p_\Sigma\) 同阶缩放，从而 \(s_m=\pi_m/p_\Sigma\) 不变。
结论 \ref{con:xi-terminal-pick-poisson-centered-ledger-invariants} 已给出 \(\eta_m\) 的尺度不变性。
\par
由结论 \ref{con:xi-terminal-pick-poisson-schur-ledger} 的逐项上界 \(\pi_m\le p_m\) 求和得 \(\sum_m\pi_m\le\sum_m p_m=p_\Sigma\)，从而 \(\sum_m s_m\le 1\)。
若取等号，则每个 \(m\) 必有 \(\pi_m=p_m\)，即
\(
\mathsf P_{m,1:m-1}(\mathsf P^{(m-1)})^{-1}\mathsf P_{1:m-1,m}=0
\)；
由 \(\mathsf P^{(m-1)}\succ 0\) 推出 \(\mathsf P_{m,1:m-1}=0\)，归纳得 \(\mathsf P\) 对角。
反向显然。
\par
对正定 Hermitian 矩阵的 Cholesky--Schur 关系给出 \(\pi_m=\abs{L_{mm}}^2\)，并且
\(
p_\Sigma=\mathrm{tr}(\mathsf P)=\mathrm{tr}(LL^\ast)=\|L\|_F^2
\)。
因此
\[
1-\sum_m s_m
=\frac{p_\Sigma-\sum_m\pi_m}{p_\Sigma}
=\frac{\|L\|_F^2-\sum_m\abs{L_{mm}}^2}{\|L\|_F^2}
=\frac{\sum_{1\le j<i\le\kappa}\abs{L_{ij}}^2}{\|L\|_F^2}.
\]
\par
由 \(\prod_m s_m=\frac{\det\mathsf P}{(\mathrm{tr}\,\mathsf P)^{\kappa}}\) 与 \(I_\partial=-\log\prod_m s_m\) 得 \(I_\partial=-\sum_m\log s_m\)。
对正数 \((s_m)\) 用 AM--GM 不等式得
\(
\prod_m s_m\le (\frac{\sum_m s_m}{\kappa})^\kappa
\)，取负对数得到第一条下界；再用 \(\sum_m s_m\le 1\) 得第二条下界。证毕。
\end{proof}

\begin{auditcor}[Schur 份额瓶颈由尺度不变量强制]\label{cor:xi-schur-ledger-bottleneck-share}
在定理 \ref{thm:xi-schur-ledger-subprobability-offdiagonal-gap} 的设定下，存在索引 \(m_\star\in\{1,\dots,\kappa\}\) 使
\[
s_{m_\star}\ \le\ e^{-I_\partial/\kappa},
\qquad\text{等价地}\qquad
\pi_{m_\star}\ \le\ (\mathrm{tr}\,\mathsf P)\,e^{-I_\partial/\kappa}.
\]
并且在统一伸缩下，上述不等式两侧均严格不变。
\end{auditcor}
\begin{proof}
由 \(I_\partial=-\sum_{m=1}^{\kappa}\log s_m\) 得
\(
\prod_{m=1}^{\kappa}s_m=e^{-I_\partial}
\)。
因此
\(
\min_m s_m\le (\prod_m s_m)^{1/\kappa}=e^{-I_\partial/\kappa}
\)，取达到最小值的 \(m_\star\) 即得第一式；第二式由 \(s_m=\pi_m/(\mathrm{tr}\,\mathsf P)\) 得到。
尺度不变性由定理 \ref{thm:xi-schur-ledger-subprobability-offdiagonal-gap}(1) 直接推出。证毕。
\end{proof}

\begin{auditcor}[条件数的熵上界]\label{cor:xi-schur-ledger-condition-number-entropy-bound}
在定理 \ref{thm:xi-schur-ledger-subprobability-offdiagonal-gap} 的设定下，若 \(\mathsf P\succ 0\) 的谱条件数为
\(
\kappa(\mathsf P):=\lambda_{\max}(\mathsf P)/\lambda_{\min}(\mathsf P)
\)，则有
\[
\log \kappa(\mathsf P)\le I_\partial.
\]
\end{auditcor}
\begin{proof}
有 \(\lambda_{\max}(\mathsf P)\le \mathrm{tr}(\mathsf P)\)。
另一方面，对 \(\lambda_{\min}\) 有粗下界
\(
\det\mathsf P=\lambda_{\min}\prod_{j<\kappa}\lambda_j\le \lambda_{\min}(\mathrm{tr}\,\mathsf P)^{\kappa-1}
\)，故 \(\lambda_{\min}\ge \det\mathsf P/(\mathrm{tr}\,\mathsf P)^{\kappa-1}\)。
于是
\[
\kappa(\mathsf P)\le \frac{\mathrm{tr}(\mathsf P)}{\det\mathsf P/(\mathrm{tr}\,\mathsf P)^{\kappa-1}}
=\frac{(\mathrm{tr}\,\mathsf P)^{\kappa}}{\det\mathsf P}
=\exp(I_\partial),
\]
取对数即得结论。证毕。
\end{proof}

\begin{auditthm}[特征行列式齐次性蕴含 Newton 和与谱矩的精确缩放]\label{thm:xi-schur-ledger-newton-sums-homogeneity}
设
\(
\Psi(z,m):=\det(I-z\mathsf P(m))
\)
满足齐次性 \(\Psi(z,m)=\Psi(z/m,1)\)（结论 \ref{con:xi-terminal-pick-poisson-characteristic-homogeneity-pde}）。
令 \(\mathsf P:=\mathsf P(1)\)。
则对任意整数 \(n\ge 1\) 有
\[
\mathrm{tr}\bigl(\mathsf P(m)^n\bigr)=m^{-n}\mathrm{tr}(\mathsf P^n).
\]
并且对任意 \(k\le \kappa\)，基本对称多项式满足
\[
e_k(\mathsf P(m))=m^{-k}e_k(\mathsf P).
\]
\end{auditthm}
\begin{proof}
在 \(|z|\) 充分小的邻域，\(\log\Psi(z,m)=\log\det(I-z\mathsf P(m))\) 具有展开
\[
\log\Psi(z,m)=-\sum_{n\ge 1}\frac{z^n}{n}\,\mathrm{tr}\bigl(\mathsf P(m)^n\bigr).
\]
由齐次性 \(\Psi(z,m)=\Psi(z/m,1)\) 得
\[
\log\Psi(z,m)=\log\Psi(z/m,1)=-\sum_{n\ge 1}\frac{(z/m)^n}{n}\,\mathrm{tr}(\mathsf P^n).
\]
比较 \(z^n\) 系数即得 \(\mathrm{tr}(\mathsf P(m)^n)=m^{-n}\mathrm{tr}(\mathsf P^n)\)。
另一方面，\(\Psi(z,m)\) 的系数由 \(\mathsf P(m)\) 的主子式配分函数确定（见结论 \ref{con:xi-terminal-pick-poisson-homothetic-spectrum-scaling} 的证明结构），齐次性等价于各阶系数按 \(m^{-k}\) 缩放，故 \(e_k(\mathsf P(m))=m^{-k}e_k(\mathsf P)\)。证毕。
\end{proof}

\begin{auditthm}[\texorpdfstring{$I_\partial$}{Ipartial} 的权重--几何分解与平均分离度约束]\label{thm:xi-schur-ledger-ipartial-weight-geometry-split}
在结论 \ref{con:xi-terminal-negative-spectrum-det-invariant} 的设定下，令
\[
q_j:=\frac{p_j}{p_\Sigma}\qquad(1\le j\le\kappa),
\qquad
\rho_{jk}:=\rho(a_j,a_k)\qquad(1\le j<k\le\kappa).
\]
则尺度不变量 \(I_\partial\) 具有分解
\[
I_\partial
=\kappa\log\kappa+\kappa D_{\mathrm{KL}}(U_\kappa\Vert q)
\;+\;
2\sum_{1\le j<k\le\kappa}\bigl(-\log\rho_{jk}\bigr).
\]
进一步，定义交叉几何平均分离度
\[
\bar\rho:=\left(\prod_{1\le j<k\le\kappa}\rho_{jk}\right)^{2/(\kappa(\kappa-1))}\in(0,1],
\]
则有
\[
-\log\bar\rho
=\frac{I_\partial-\kappa\log\kappa-\kappa D_{\mathrm{KL}}(U_\kappa\Vert q)}{\kappa(\kappa-1)}
\le
\frac{I_\partial-\kappa\log\kappa}{\kappa(\kappa-1)}.
\]
从而得到平均分离度下界
\[
\bar\rho\ge 1-\frac{I_\partial-\kappa\log\kappa}{\kappa(\kappa-1)}.
\]
\end{auditthm}
\begin{proof}
由结论 \ref{con:xi-terminal-negative-spectrum-det-invariant} 得
\(
\det\mathsf P=(\prod_j p_j)(\prod_{j<k}\rho_{jk}^2)
\)，故
\[
\frac{\det\mathsf P}{(\mathrm{tr}\,\mathsf P)^\kappa}
=\left(\prod_{j=1}^{\kappa}q_j\right)\left(\prod_{j<k}\rho_{jk}^2\right).
\]
取负对数并用 \(I_\partial=-\log\!\left(\frac{\det\mathsf P}{(\mathrm{tr}\,\mathsf P)^\kappa}\right)\) 得
\(
I_\partial=-\sum_j\log q_j+2\sum_{j<k}(-\log\rho_{jk})
\)。
再由
\(
D_{\mathrm{KL}}(U_\kappa\Vert q)=-\log\kappa-\frac1\kappa\sum_j\log q_j
\)
得 \(-\sum_j\log q_j=\kappa\log\kappa+\kappa D_{\mathrm{KL}}(U_\kappa\Vert q)\)，从而得到分解式。
\par
由 \(\bar\rho\) 的定义直接得到
\(
-\log\bar\rho=\frac{2}{\kappa(\kappa-1)}\sum_{j<k}(-\log\rho_{jk})
\)，与上式比较即得等式与不等式。
最后由 \(1-x\le -\log x\)（\(x\in(0,1]\)）推出 \(\bar\rho\ge 1-(-\log\bar\rho)\)，代入上界即得下界。证毕。
\end{proof}

\begin{auditprop}[并集交叉能量的行列式比值外部化]\label{prop:xi-schur-ledger-union-det-ratio-geometric-mean}
设缺陷多重集分解为不交并 \(A\bigsqcup B\)，并以同一构造得到 \(\mathsf P(A),\mathsf P(B),\mathsf P(A\cup B)\)。
则交叉能量恒等式（结论 \ref{con:xi-terminal-pick-poisson-green-mirror-energy}）等价于乘法式
\[
\frac{\det\mathsf P(A\cup B)}{\det\mathsf P(A)\det\mathsf P(B)}
=\prod_{a\in A}\prod_{b\in B}\rho(a,b)^2.
\]
从而交叉伪双曲几何平均
\[
\bar\rho(A,B)
:=\left(\prod_{a\in A}\prod_{b\in B}\rho(a,b)\right)^{1/(\card{A}\card{B})}\in(0,1]
\]
可由三个行列式标量恢复为
\[
\bar\rho(A,B)
=\exp\!\left(
\frac{1}{2\card{A}\card{B}}
\log\frac{\det\mathsf P(A\cup B)}{\det\mathsf P(A)\det\mathsf P(B)}
\right).
\]
\end{auditprop}
\begin{proof}
由结论 \ref{con:xi-terminal-pick-poisson-green-mirror-energy} 的并集交叉能量恒等式，
\[
-\log\det\mathsf P(A\cup B)
=-\log\det\mathsf P(A)-\log\det\mathsf P(B)
+2\sum_{a\in A}\sum_{b\in B}\bigl(-\log\rho(a,b)\bigr).
\]
移项并指数化即得行列式比值公式。
再对 \(\card{A}\card{B}\) 次几何平均取对数即可得到 \(\bar\rho(A,B)\) 的恢复式。证毕。
\end{proof}

\endinput


\begin{conclusion}[谱裕度与病态封口：\texorpdfstring{$\lambda_{\min}(\mathsf P)$}{lambda_min(P)} 与条件数在贴边/碰撞极限下同步塌缩]\label{con:xi-terminal-pick-poisson-conditioning-collapse}
对扫描谱点集诱导的 Pick--Poisson 矩阵 \(\mathsf P(t)\)（命题 \ref{prop:xi-scan-toeplitz-conditioning-by-separation}），记端点权重 \(p_j(t):=P_{a_j(t)}(-1)\) 与伪双曲距离 \(\rho_{jk}(t):=\rho(a_j(t),a_k(t))\)。
则其最小特征值满足显式夹逼：一方面
\[
\boxed{\;
\lambda_{\min}(\mathsf P(t))
\ \ge\
\frac{\Bigl(\prod_{j=1}^{\kappa}p_j(t)\Bigr)\Bigl(\prod_{1\le j<k\le\kappa}\rho_{jk}(t)^2\Bigr)}
{\Bigl(\sum_{j=1}^{\kappa}p_j(t)\Bigr)^{\kappa-1}}\; } ;
\]
另一方面令 \(p_{\max}(t):=\max_j p_j(t)\)、\(\rho_{\min}(t):=\min_{j<k}\rho_{jk}(t)\)，则有二次上界
\[
\boxed{\;
\lambda_{\min}(\mathsf P(t))\ \le\ p_{\max}(t)\,\rho_{\min}(t)^2\; }.
\]
再由 \(\lambda_{\max}(\mathsf P)\le \mathrm{tr}(\mathsf P)\) 与推论 \ref{cor:xi-pick-poisson-gap-lower-bound} 的下界，可得条件数封口
\[
\boxed{\;
\kappa(\mathsf P)
:=\frac{\lambda_{\max}(\mathsf P)}{\lambda_{\min}(\mathsf P)}
\ \le\
\frac{\bigl(\mathrm{tr}\,\mathsf P\bigr)^{\kappa}}
{\Bigl(\prod_{j}P_{a_j}(-1)\Bigr)\Bigl(\prod_{j<k}\rho(a_j,a_k)^2\Bigr)}\; }.
\]
因此当 \(\rho_{\min}(t)\downarrow 0\)（碰撞）或 \(p_j(t)\downarrow 0\)（贴边）时，Pick/Toeplitz 侧的谱裕度必按至少二次速率塌缩，形成可审计的病态证书。
\end{conclusion}
\begin{proof}[证明要点]
上下界见命题 \ref{prop:xi-scan-toeplitz-conditioning-by-separation}；
条件数分式封口见推论 \ref{cor:xi-pick-poisson-condition-number-certificate}（并用 \(\det\mathsf P\) 的闭式分解）。
\end{proof}

\begin{conclusion}[Hankel--Toeplitz 病态同源：两类反演的最小稳定性常数共享 Vandermonde 平方律]\label{con:xi-terminal-hankel-toeplitz-common-illconditioning}
在有限原子矩口径 \(\mu=\sum_{j=1}^{\kappa}w_j\delta_{a_j}\) 下，记 \(b_j:=a_j^{-1}\in(0,1)\)、\(\widetilde w_j:=w_j b_j\)，并令 Hankel 主块 \(\mathsf H_{\kappa-1}:=(\mu_{p+q})_{0\le p,q\le \kappa-1}\)。
则有行列式闭式
\[
\boxed{\;
\det \mathsf H_{\kappa-1}
=\left(\prod_{j=1}^{\kappa}\widetilde w_j\right)
\left(\prod_{1\le j<k\le \kappa}(b_k-b_j)^2\right)\; }.
\]
因此任意基于 Hankel 反演（Prony/Pad\'e/Hankel 核空间法）的误差放大因子不可能在 \(b\)-碰撞极限保持一致有界，至少按 \(\prod_{j<k}|b_k-b_j|^{-2}\) 发散。
该 Vandermonde 平方律与
\(
\det\mathsf P=(\prod_j P_{a_j}(-1))(\prod_{j<k}\rho(a_j,a_k)^2)
\)
在退化机制上同型，表明“矩证书通道（Toeplitz/Pick）”与“层析反演通道（Hankel/Vandermonde）”在碰撞参数上发生同源病态，并可由行列式闭式给出统一的下界型稳定性护栏。
\end{conclusion}
\begin{proof}[证明要点]
Hankel 行列式的 Vandermonde 分解与病态发散见命题 \ref{prop:xi-depth-hankel-determinant-vandermonde-square}；
Pick--Poisson 的同型乘积律见命题 \ref{prop:xi-pick-poisson-det-factorization}。
\end{proof}

\begin{theorem}[Prony 矩映射的 \texorpdfstring{$\Delta^4$}{Delta^4} 雅可比律]\label{thm:xi-prony-moment-map-jacobian-delta4}
设 \(k\) 为特征不为 \(2\) 的域，取两两不同的 \(a_1,\dots,a_\kappa\in k\) 与非零权 \(\omega_1,\dots,\omega_\kappa\in k^{\times}\)。
定义矩量
\[
s_m:=\sum_{i=1}^{\kappa}\omega_i a_i^{m}\qquad(m\ge 0),
\]
并考虑 Prony 矩映射
\[
\Phi_{\kappa}:k^{2\kappa}\longrightarrow k^{2\kappa},\qquad
(a_1,\dots,a_\kappa,\omega_1,\dots,\omega_\kappa)\longmapsto (s_0,s_1,\dots,s_{2\kappa-1}).
\]
则雅可比行列式满足恒等式
\[
\det\!\left(\frac{\partial(s_0,\dots,s_{2\kappa-1})}{\partial(\omega_1,\dots,\omega_\kappa,a_1,\dots,a_\kappa)}\right)
=
(-1)^{\frac{\kappa(\kappa-1)}2}
\Bigl(\prod_{i=1}^{\kappa}\omega_i\Bigr)
\Bigl(\prod_{1\le i<j\le \kappa}(a_j-a_i)^4\Bigr).
\]
\leanverified{paper\_xi\_prony\_moment\_map\_jacobian\_delta4}
\end{theorem}
\begin{proof}
记 \(J\) 为上述 \((2\kappa\times 2\kappa)\) 雅可比矩阵。
对 \(m=0,\dots,2\kappa-1\) 与 \(i=1,\dots,\kappa\) 有
\[
\frac{\partial s_m}{\partial \omega_i}=a_i^{m},
\qquad
\frac{\partial s_m}{\partial a_i}=\omega_i\,m\,a_i^{m-1}.
\]
将后 \(\kappa\) 列分别提出 \(\omega_i\) 因子，得
\[
\det(J)=\Bigl(\prod_{i=1}^{\kappa}\omega_i\Bigr)\det(C),
\]
其中 \(C\) 的列由 \(\{a_i^{m}\}_{m=0}^{2\kappa-1}\) 与 \(\{m a_i^{m-1}\}_{m=0}^{2\kappa-1}\) 并置构成，即每个节点 \(a_i\) 的“值—导数”汇合矩阵。
该 \(C\) 为每个节点重数为 \(2\) 的汇合 Vandermonde（confluent Vandermonde）矩阵，其行列式的闭式为
\[
\det(C)=(-1)^{\frac{\kappa(\kappa-1)}2}\prod_{1\le i<j\le \kappa}(a_j-a_i)^{2\cdot 2},
\]
从而
\(
\det(C)=(-1)^{\frac{\kappa(\kappa-1)}2}\prod_{i<j}(a_j-a_i)^4
\)。
代回即得结论。证毕。
\end{proof}

\begin{corollary}[Prony 谱的判定与重建门槛：一般位置下相差一项采样]\label{cor:xi-prony-decision-vs-reconstruction-one-sample-gap}
\leanverified{paper\_xi\_prony\_decision\_vs\_reconstruction\_one\_sample\_gap}
在定理 \ref{thm:xi-prony-moment-map-jacobian-delta4} 的设定下，令 \((s_m)_{m\ge 0}\) 由
\(
s_m=\sum_{i=1}^{\kappa}\omega_i a_i^{m}
\)
给出，并假设处于一般位置：\(\{a_i\}\) 两两不同且 \(\prod_i\omega_i\neq 0\)。
则成立以下三条互相兼容的最小性结论：
\begin{enumerate}
  \item \textbf{最小可判定阈值.}
        记 Hankel 主块
        \[
        H_{\kappa-1}:=\bigl(s_{i+j}\bigr)_{0\le i,j\le \kappa-1}\in M_{\kappa}(k).
        \]
        则在一般位置下 \(H_{\kappa-1}\) 可逆，并且其行列式对最高阶样本 \(s_{2\kappa-2}\) 为仿射函数：
        对固定前缀 \((s_0,\dots,s_{2\kappa-3})\) 有
        \[
        \det(H_{\kappa-1})=\alpha\,s_{2\kappa-2}+\beta,
        \qquad
        \alpha=\det\!\bigl(s_{i+j}\bigr)_{0\le i,j\le \kappa-2}.
        \]
        因而在 \(\alpha\neq 0\) 的一般位置上，仅凭更短前缀 \((s_0,\dots,s_{2\kappa-3})\) 不能判定 \(\det(H_{\kappa-1})\) 的消失/非消失，
        而给出长度 \(2\kappa-1\) 的前缀 \((s_0,\dots,s_{2\kappa-2})\) 则可直接完成该判定。
  \item \textbf{截断纤维维数律.}
        记截断矩映射
        \[
        \Phi_{\kappa}^{\langle 2\kappa-1\rangle}:(a_1,\dots,a_\kappa,\omega_1,\dots,\omega_\kappa)\longmapsto (s_0,\dots,s_{2\kappa-2})\in k^{2\kappa-1}.
        \]
        则在上述一般位置的 Zariski 开集上，\(\Phi_{\kappa}^{\langle 2\kappa-1\rangle}\) 的纤维为一维光滑代数簇；
        换言之，长度 \(2\kappa-1\) 的截断数据在一般位置下尚残留一个自由参数。
  \item \textbf{最小可重建阈值.}
        加入一个额外采样 \(s_{2\kappa-1}\)（即给出长度 \(2\kappa\) 的前缀 \((s_0,\dots,s_{2\kappa-1})\)），
        则在一般位置下可唯一重建节点集 \(\{a_i\}\) 与权重集 \(\{\omega_i\}\)（至多差一个置换），
        等价地，Prony 矩映射 \(\Phi_{\kappa}\) 在一般位置上为局部双射并在置换商后无歧义。
\end{enumerate}
\end{corollary}
\begin{proof}
\textbf{(1)}
对 \(\det(H_{\kappa-1})\) 沿右下角元素 \(s_{2\kappa-2}\) 作 Laplace 展开即得仿射表达，其中系数 \(\alpha\) 为对应余子式。
在一般位置下 \(H_{\kappa-2}\) 可逆从而 \(\alpha\neq 0\)，故可通过选择 \(s_{2\kappa-2}\) 使 \(\det(H_{\kappa-1})\) 取零或取非零，得到“更短前缀不可判定”；而给出 \(s_{2\kappa-2}\) 后可直接计算 \(\det(H_{\kappa-1})\) 完成判定。
\par
\textbf{(2)}
由定理 \ref{thm:xi-prony-moment-map-jacobian-delta4}，在一般位置下 \(\det J\neq 0\)，故 \(\Phi_\kappa\) 在该点附近为局部解析同胚。
将 \(\Phi_\kappa\) 与坐标投影 \((s_0,\dots,s_{2\kappa-1})\mapsto(s_0,\dots,s_{2\kappa-2})\) 复合，可知 \(\Phi_{\kappa}^{\langle 2\kappa-1\rangle}\) 在该点的微分秩为 \(2\kappa-1\)，
从而由隐函数定理得其局部纤维为一维流形；代数意义下即为一般位置上的一维光滑纤维。
\par
\textbf{(3)}
仍由定理 \ref{thm:xi-prony-moment-map-jacobian-delta4}，在一般位置下 \(\Phi_\kappa\) 为局部双射。
由于 \(\Phi_\kappa\) 对指标置换不变，故节点与权重仅可能差一个置换；在取置换商后得到无歧义重建。证毕。
\end{proof}

\begin{corollary}[外部 Hankel 证书与参数见证之间的平方间隙]\label{cor:xi-hankel-vs-prony-square-gap}
在定理 \ref{thm:xi-prony-moment-map-jacobian-delta4} 的设定下，令
\[
H_{\kappa}:=(s_{i+j})_{0\le i,j\le \kappa-1},
\qquad
D_{\kappa}:=\det(H_{\kappa}).
\]
则有 Vandermonde 平方律
\[
D_{\kappa}=\Bigl(\prod_{i=1}^{\kappa}\omega_i\Bigr)\Bigl(\prod_{1\le i<j\le \kappa}(a_j-a_i)^2\Bigr),
\]
并且雅可比行列式满足
\[
\det(J)=(-1)^{\frac{\kappa(\kappa-1)}2}\,D_{\kappa}\cdot\Bigl(\prod_{1\le i<j\le \kappa}(a_j-a_i)^2\Bigr),
\]
其中 \(J\) 为定理 \ref{thm:xi-prony-moment-map-jacobian-delta4} 中的雅可比矩阵。
等价地，若记
\(
\Delta(a):=\prod_{1\le i<j\le \kappa}(a_j-a_i)
\)，
则
\[
\det(J)=(-1)^{\frac{\kappa(\kappa-1)}2}\,D_{\kappa}\,\Delta(a)^2.
\]
特别地，对任意离散赋值 \(v\)（例如 \(p\)-进赋值），有精确恒等式
\[
v(\det(J))=v(D_{\kappa})+2\sum_{1\le i<j\le \kappa} v(a_j-a_i).
\]
\leanverified{hankel2}
\leanverified{hankel2\_vandermonde\_square}
\leanverified{hankel2\_vandermonde\_square\_symm}
\leanverified{hankel3DetScalar}
\leanverified{hankelMoment0}
\leanverified{hankelMoment1}
\leanverified{hankelMoment2}
\leanverified{hankelMoment3}
\leanverified{hankelMoment4}
\leanverified{hankel3\_vandermonde\_square\_scalar}
\leanverified{hankel4}
\leanverified{hankel4\_vandermonde\_square}
\leanverified{paper\_hankel\_vandermonde\_package}
\leanverified{hankel2\_eq\_zero\_of\_collision}
\leanverified{hankel2\_ne\_zero\_iff}
\leanverified{hankel3\_eq\_zero\_of\_any\_collision}
\leanverified{paper\_hankel\_collision\_degeneracy}
\leanverified{hankel4\_eq\_zero\_of\_collision\_12}
\leanverified{hankel4\_eq\_zero\_of\_collision\_34}
\leanverified{hankel4\_eq\_zero\_of\_any\_collision}
\leanverified{paper\_hankel4\_collision\_degeneracy}
\leanverified{hankel3\_matrix\_eq\_zero\_of\_a1\_a2}
\leanverified{hankel3\_matrix\_eq\_zero\_of\_a2\_a3}
\leanverified{hankel3\_matrix\_eq\_zero\_of\_any\_collision}
\leanverified{paper\_hankel3\_matrix\_collision\_degeneracy}
\leanverified{hankel2\_swap\_symmetry}
\leanverified{hankel2\_zero\_first\_weight}
\leanverified{hankel2\_zero\_second\_weight}
\leanverified{paper\_hankel2\_symmetry\_zero\_weight\_package}
\leanverified{paper\_xi\_hankel\_vs\_prony\_square\_gap}
\end{corollary}
\begin{proof}
令 Vandermonde 矩阵 \(V=(a_j^{i})_{0\le i\le \kappa-1,\ 1\le j\le \kappa}\)。
则
\[
H_{\kappa}=V\,\mathrm{diag}(\omega_1,\dots,\omega_\kappa)\,V^{\mathsf T},
\]
从而
\[
D_{\kappa}=\det(V)^2\prod_{i=1}^{\kappa}\omega_i
=\Bigl(\prod_{1\le i<j\le \kappa}(a_j-a_i)^2\Bigr)\prod_{i=1}^{\kappa}\omega_i,
\]
得到第一式。
将其与定理 \ref{thm:xi-prony-moment-map-jacobian-delta4} 的 \(\det(J)\) 表达式相除即得第二式与等价形式。
最后的赋值恒等式由乘法性与 \(v(\Delta(a)^2)=2\sum_{i<j}v(a_j-a_i)\) 直接推出。证毕。
\end{proof}

\begin{corollary}[非阿基米德可识别性的单位判别：单位 Jacobian \texorpdfstring{$\Longleftrightarrow$}{<->} 单位 Hankel 且无 \texorpdfstring{$p$}{p}-进碰撞]\label{cor:xi-nonarchimedean-prony-unit-jacobian-criterion}
\leanverified{paper\_xi\_nonarchimedean\_prony\_unit\_jacobian\_criterion}
在推论 \ref{cor:xi-hankel-vs-prony-square-gap} 的设定下，进一步设 \(k=K\) 为完备离散赋值域并以 \(v\) 记其赋值。
若 \(a_i\in\mathcal O_K\) 两两不同且 \(\omega_i\in\mathcal O_K^\times\)，则以下条件等价：
\begin{enumerate}
  \item \(v(\det(J))=0\)；
  \item \(v(D_\kappa)=0\) 且对所有 \(i<j\) 有 \(v(a_j-a_i)=0\)。
\end{enumerate}
并且当上述条件成立时，Prony 矩映射 \(\Phi_\kappa\) 在该点邻域内为双射解析同胚（从而局部可逆且解唯一），见定理 \ref{thm:xi-nonarchimedean-prony-hensel-newton-inversion}。
\end{corollary}
\begin{proof}
由推论 \ref{cor:xi-hankel-vs-prony-square-gap} 的赋值分解
\(
v(\det(J))=v(D_\kappa)+2\sum_{i<j}v(a_j-a_i)
\)，且两项均为非负整数。
因此 \(v(\det(J))=0\) 当且仅当 \(v(D_\kappa)=0\) 且 \(v(a_j-a_i)=0\) 对所有 \(i<j\) 同时成立，得到等价性。
最后一条由定理 \ref{thm:xi-nonarchimedean-prony-hensel-newton-inversion} 的第一点直接推出。证毕。
\end{proof}

\begin{theorem}[非阿基米德局部可识别性：\texorpdfstring{$\Delta^4$}{Delta^4} 门槛的 Hensel--Newton 反演]\label{thm:xi-nonarchimedean-prony-hensel-newton-inversion}
\leanverified{paper\_xi\_nonarchimedean\_prony\_hensel\_newton\_inversion}
令 \(K\) 为完备离散赋值域，\(\mathcal O_K\) 为其赋值环，\(\pi\) 为均值元，赋值记为 \(v=v_\pi\)。
取两两不同的 \(a_1,\dots,a_\kappa\in\mathcal O_K\) 与非零权 \(\omega_1,\dots,\omega_\kappa\in\mathcal O_K^\times\)，并令
\[
s_m:=\sum_{i=1}^{\kappa}\omega_i a_i^{m}\qquad(m\ge 0).
\]
定义 Prony 矩映射
\[
\Phi_{\kappa}:\mathcal O_K^{2\kappa}\longrightarrow \mathcal O_K^{2\kappa},\qquad
x=(a_1,\dots,a_\kappa,\omega_1,\dots,\omega_\kappa)\longmapsto (s_0,\dots,s_{2\kappa-1}),
\]
并设其雅可比矩阵为 \(J(x)\)。
则成立：
\begin{enumerate}
  \item 若 \(v(\det J(x))=0\)，则存在邻域 \(U\ni x\)、\(V\ni \Phi_\kappa(x)\)，使得 \(\Phi_\kappa:U\to V\) 为双射解析同胚；并且对任意 \(y'\in V\)，方程 \(\Phi_\kappa(z)=y'\) 在 \(U\) 内有唯一解。
  \item 更一般地，令 \(y\in\mathcal O_K^{2\kappa}\)，并取近似解 \(x_0\in\mathcal O_K^{2\kappa}\)。
  记向量的赋值为 \(v(u):=\min_i v(u_i)\)。
  若
  \[
  v\!\bigl(\Phi_\kappa(x_0)-y\bigr)\ \ge\ 2\,v\!\bigl(\det J(x_0)\bigr)+1,
  \]
  则存在唯一 \(x\in\mathcal O_K^{2\kappa}\) 使 \(\Phi_\kappa(x)=y\)，并且满足误差界
  \[
  v(x-x_0)\ \ge\ v\!\bigl(\Phi_\kappa(x_0)-y\bigr)-v\!\bigl(\det J(x_0)\bigr).
  \]
\end{enumerate}
\end{theorem}
\begin{proof}
将 \(F(z):=\Phi_\kappa(z)-y\) 视作 \(\mathcal O_K\)-解析映射。
当 \(v(\det J(x))=0\) 时，\(J(x)\) 在 \(\mathcal O_K\) 上可逆，非阿基米德逆函数定理给出局部解析同胚与唯一解。
\par
一般情形由多变量 Hensel 引理的标准形式给出：条件
\(
v(F(x_0))\ge 2v(\det J(x_0))+1
\)
保证 Newton 迭代收敛并产生唯一提升解，同时误差界来自一次线性化与 Cramer/伴随估计。证毕。
\end{proof}


\begin{conclusion}[Hankel 刚性理想与非阿基米德极小化原理]\label{con:xi-terminal-hankel-rigidity-ideal-nonarchimedean}
设 \(R\) 为 Dedekind 域，商域为 \(K\)。
取两两不同的元 \(a_1,\dots,a_r\in R\) 与非零权 \(\omega_1,\dots,\omega_r\in R\setminus\{0\}\)。
定义矩量
\[
s_k:=\sum_{i=1}^{r}\omega_i a_i^k\in R\qquad(k\ge 0),
\]
以及阶数 \(n\ge 1\) 的 Hankel 主块与主 Hankel 行列式
\[
\mathsf H_n:=\bigl(s_{i+j}\bigr)_{0\le i,j\le n-1}\in M_n(R),\qquad
D_n:=\det(\mathsf H_n)\in R.
\]
对任意 \(I\subseteq\{1,\dots,r\}\) 且 \(|I|=n\)，记
\[
\omega_I:=\prod_{i\in I}\omega_i,\qquad
\Delta_I(a):=\prod_{\substack{i<j\\ i,j\in I}}(a_j-a_i)\in R
\]
为权重乘积与 Vandermonde 乘积，并定义第 \(n\) 阶 Hankel 刚性理想
\[
\mathfrak H_n(a,\omega)\ :=\ \big\langle\ \omega_I\Delta_I(a)^2:\ I\subseteq\{1,\dots,r\},\ |I|=n\ \big\rangle\ \subset R.
\]
则成立：
\begin{enumerate}
  \item \(D_n\in\mathfrak H_n(a,\omega)\)。因此对任意非零素理想 \(\mathfrak p\subset R\) 都有
  \[
  v_{\mathfrak p}(D_n)\ \ge\ v_{\mathfrak p}\!\bigl(\mathfrak H_n(a,\omega)\bigr),
  \]
  其中 \(v_{\mathfrak p}(\mathfrak a)\) 表示 \(\mathfrak a\) 在 \(\mathfrak p\) 处的指数（亦即 \(\mathfrak a\subseteq \mathfrak p^{e}\) 的最大 \(e\)）。
  \item 对任意 \(\mathfrak p\)，有精确等式
  \[
  \boxed{\;
  v_{\mathfrak p}\!\bigl(\mathfrak H_n(a,\omega)\bigr)
  =\min_{|I|=n} v_{\mathfrak p}\!\bigl(\omega_I\Delta_I(a)^2\bigr)
  =\min_{|I|=n}\Bigl(\sum_{i\in I}v_{\mathfrak p}(\omega_i)+2\!\sum_{\substack{i<j\\ i,j\in I}}v_{\mathfrak p}(a_j-a_i)\Bigr)\; }.
  \]
  \item 设 \(R_{\mathfrak p}\) 为局部化并取其均值元 \(\pi\)，剩余域记为 \(\kappa:=R_{\mathfrak p}/(\pi)\)。
  令
  \[
  e:=v_{\mathfrak p}\!\bigl(\mathfrak H_n(a,\omega)\bigr),\qquad
  \mathcal M:=\Bigl\{I:\ |I|=n,\ v_{\mathfrak p}\!\bigl(\omega_I\Delta_I(a)^2\bigr)=e\Bigr\}.
  \]
  则在 \(\kappa\) 中有同余恒等式
  \[
  \frac{D_n}{\pi^e}\ \equiv\ \sum_{I\in\mathcal M} \overline{\Bigl(\frac{\omega_I\Delta_I(a)^2}{\pi^e}\Bigr)}\pmod{\pi}.
  \]
  尤其若右端在 \(\kappa\) 中非零，则 \(v_{\mathfrak p}(D_n)=e\)。
\end{enumerate}
\end{conclusion}
\begin{proof}[证明要点]
令 Vandermonde 矩阵 \(\mathsf V=(a_i^j)_{1\le i\le r,\ 0\le j\le n-1}\) 且 \(\mathsf W=\mathrm{diag}(\omega_1,\dots,\omega_r)\)。
则
\(
(\mathsf V^{\mathsf T}\mathsf W\mathsf V)_{ij}=\sum_{k=1}^r\omega_k a_k^{i+j}=s_{i+j}
\)，从而 \(\mathsf H_n=\mathsf V^{\mathsf T}\mathsf W\mathsf V\)。
由 Cauchy--Binet 公式
\[
D_n=\det(\mathsf V^{\mathsf T}\mathsf W\mathsf V)
=\sum_{|I|=n}\det(\mathsf V_I)^2\prod_{i\in I}\omega_i
=\sum_{|I|=n}\omega_I\,\Delta_I(a)^2,
\]
其中 \(\mathsf V_I\) 为取行集 \(I\) 的 \(n\times n\) 子 Vandermonde 矩阵且 \(\det(\mathsf V_I)=\pm\Delta_I(a)\)。
于是 \(D_n\) 为 \(\{\omega_I\Delta_I(a)^2\}\) 的和，故 \(D_n\in\mathfrak H_n(a,\omega)\)，并得到第一条。
\par
对任意素理想 \(\mathfrak p\)，在离散赋值环 \(R_{\mathfrak p}\) 中，理想由生成元张成时其 \(\mathfrak p\)-指数等于生成元赋值的最小值，故
\(
v_{\mathfrak p}(\mathfrak H_n)=\min_{|I|=n}v_{\mathfrak p}(\omega_I\Delta_I^2)
\)；
又由 \(v_{\mathfrak p}(\Delta_I)=\sum_{i<j}v_{\mathfrak p}(a_j-a_i)\) 得第二条。
第三条将 \(D_n\) 写成 \(\pi^e\) 乘以单位和之和并模 \(\pi\) 取像即得。证毕。
\end{proof}

\begin{conclusion}[\texorpdfstring{$p$}{p}-进层状分解与刚性指数的树形动态规划表达]\label{con:xi-terminal-hankel-rigidity-dp}
取 \(R=\ZZ_p\) 并以 \(v_p(p)=1\) 归一化 \(p\)-进赋值。
在结论 \ref{con:xi-terminal-hankel-rigidity-ideal-nonarchimedean} 的设定下，
以同余关系 \(a_i\equiv a_j\ (\mathrm{mod}\ p^k)\)（\(k\ge 1\)）将指标集 \(\{1,\dots,r\}\) 分层划分为有限树 \(\mathcal T\)：
根为整集（层 \(0\)），层 \(k\) 的节点为模 \(p^k\) 的同余类及其包含关系，叶为单点类。
对叶 \(i\) 定义 \(\alpha_i:=v_p(\omega_i)\)。
对任意节点 \(u\in\mathcal T\) 记其叶子集合为 \(L(u)\)、容量为 \(c(u):=|L(u)|\)。
定义函数
\(
F_u:\{0,1,\dots,c(u)\}\to\ZZ_{\ge 0}\cup\{\infty\}
\)
如下：
\begin{itemize}
  \item 若 \(u\) 为叶 \(i\)，则 \(F_u(0)=0\)、\(F_u(1)=\alpha_i\)，其余取 \(\infty\)；
  \item 若 \(u\) 为非根节点，子节点为 \(u_1,\dots,u_m\)，则对 \(0\le t\le c(u)\) 令
  \[
  F_u(t)\ :=\ 2\binom{t}{2}\ +\ \min_{\substack{t_1+\cdots+t_m=t\\ 0\le t_j\le c(u_j)}}\ \sum_{j=1}^m F_{u_j}(t_j);
  \]
  \item 若 \(u=\mathrm{root}\)（层 \(0\)），则不加 \(2\binom{t}{2}\)：
  \[
  F_{\mathrm{root}}(t)\ :=\ \min_{\substack{t_1+\cdots+t_m=t\\ 0\le t_j\le c(u_j)}}\ \sum_{j=1}^m F_{u_j}(t_j).
  \]
\end{itemize}
则对每个 \(0\le n\le r\) 有
\[
\boxed{\;v_p\!\bigl(\mathfrak H_n(a,\omega)\bigr)\ =\ F_{\mathrm{root}}(n)\; }.
\]
\end{conclusion}
\begin{proof}[证明要点]
对任意 \(x,y\in\ZZ_p\) 有恒等式
\(
v_p(x-y)=\sum_{k\ge 1}\mathbf{1}_{x\equiv y\ (\mathrm{mod}\ p^k)}
\)。
对任意 \(I\subseteq\{1,\dots,r\}\)（\(|I|=n\)）将两两差的赋值求和得
\[
\sum_{\substack{i<j\\ i,j\in I}}v_p(a_j-a_i)
=\sum_{k\ge 1}\sum_{u\in\mathcal T_k}\binom{|I\cap L(u)|}{2},
\]
其中 \(\mathcal T_k\) 为层 \(k\) 的节点集。
于是
\[
v_p\!\bigl(\omega_I\Delta_I(a)^2\bigr)
=\sum_{i\in I}\alpha_i
2\sum_{k\ge 1}\sum_{u\in\mathcal T_k}\binom{|I\cap L(u)|}{2}.
\]
由结论 \ref{con:xi-terminal-hankel-rigidity-ideal-nonarchimedean}，
\(
v_p(\mathfrak H_n)=\min_{|I|=n}v_p(\omega_I\Delta_I^2)
\)；
而右端最小化问题对树 \(\mathcal T\) 具有严格的最优子结构：对任一节点 \(u\)，若在其叶子中选取 \(t\) 个，则该层贡献恰为 \(2\binom{t}{2}\)，并且最优代价等于子节点分配的最优代价之和加上该层贡献（根层对应 \(k=0\)，不出现 \(2\binom{t}{2}\)）。
因此自底向上的动态规划即给出全局最小值，并得到所示等式。
DP 与暴力枚举在固定样例窗内的一致性证书可由脚本 \texttt{scripts/exp\_xi\_hankel\_rigidity\_dp\_audit.py} 复核。证毕。
\end{proof}

\begin{conclusion}[无权刚性理想的平方结构与 \texorpdfstring{$p$}{p}-ordering 刻画]\label{con:xi-terminal-hankel-rigidity-unweighted-bhargava}
设 \(R\) 为 Dedekind 域，取两两不同的 \(a_1,\dots,a_r\in R\)。
对 \(n\le r\) 定义 Vandermonde 理想与无权刚性理想
\[
J_n(a)\ :=\ \big\langle\ \Delta_I(a):\ |I|=n\ \big\rangle\subset R,
\qquad
\mathfrak H_n(a,1)\ :=\ \big\langle\ \Delta_I(a)^2:\ |I|=n\ \big\rangle\subset R.
\]
则有理想恒等式
\[
\boxed{\;\mathfrak H_n(a,1)\ =\ J_n(a)^2.\;}
\]
当 \(R=\ZZ\) 且固定素数 \(p\) 时，令 \(S:=\{a_1,\dots,a_r\}\) 并取其一个 \(p\)-ordering（Bhargava \cite{Bhargava1997POrderings}），记 \(V_k(S,p)\) 为相应的 \(p\)-序列。
则对任意 \(n\le r\) 有
\[
v_p\bigl(J_n(a)\bigr)\ =\ \sum_{k=1}^{n-1} v_p\bigl(V_k(S,p)\bigr),
\qquad
v_p\bigl(\mathfrak H_n(a,1)\bigr)\ =\ 2\sum_{k=1}^{n-1} v_p\bigl(V_k(S,p)\bigr).
\]
\end{conclusion}
\begin{proof}[证明要点]
对任意素理想 \(\mathfrak p\) 有
\(
v_{\mathfrak p}(\mathfrak H_n(a,1))=\min_{|I|=n}2v_{\mathfrak p}(\Delta_I)
=2\min_{|I|=n}v_{\mathfrak p}(\Delta_I)
=2v_{\mathfrak p}(J_n(a))
=v_{\mathfrak p}(J_n(a)^2)
\)，
而 Dedekind 域中非零理想由所有素理想指数唯一确定，故 \(\mathfrak H_n(a,1)=J_n(a)^2\)。
\par
当 \(R=\ZZ\) 时，Bhargava 的 \(p\)-ordering 逐步最小化
\(
v_p\!\bigl(\prod_{j<k}(b_k-b_j)\bigr)
\)（\(b_k\in S\)），并证明所得到的 \(p\)-序列与 \(p\)-ordering 的选取无关。
又 Vandermonde 乘积分解
\(
\Delta(b_0,\dots,b_{n-1})=\prod_{k=1}^{n-1}\prod_{j=0}^{k-1}(b_k-b_j)
\)
将赋值写为步骤赋值之和，从而给出 \(v_p(J_n)\) 的最小值表达；与 \(\mathfrak H_n=J_n^2\) 合并即得第二式。证毕。
\end{proof}

\begin{conclusion}[刚性指数的离散凸性]\label{con:xi-terminal-hankel-rigidity-discrete-convexity}
在结论 \ref{con:xi-terminal-hankel-rigidity-dp} 的设定下，定义
\[
f(n):=v_p\!\bigl(\mathfrak H_n(a,\omega)\bigr)=F_{\mathrm{root}}(n)\qquad(0\le n\le r).
\]
则 \(f\) 为离散凸函数：其二阶差分处处非负，
\[
\boxed{\ f(n+1)-2f(n)+f(n-1)\ \ge\ 0\qquad(1\le n\le r-1).\ }
\]
\end{conclusion}
\begin{proof}[证明要点]
叶节点的代价函数 \(t\mapsto F_u(t)\) 显然离散凸（其有效支撑仅在 \(t=0,1\)）。
若各子节点代价函数离散凸，则其下确界卷积
\(
t\mapsto \min_{t_1+\cdots+t_m=t}\sum_j F_{u_j}(t_j)
\)
仍离散凸；而 \(2\binom{t}{2}\) 的二阶差分为常数 \(2\)。
因此对非根节点有
\(
F_u(t)=2\binom{t}{2}+(\text{下确界卷积})(t)
\)
保持离散凸性；根节点仅为下确界卷积，仍保持离散凸性。
自叶向根归纳即得结论。证毕。
\end{proof}

\begin{conclusion}[仿射变换下的刚性理想同变性]\label{con:xi-terminal-hankel-rigidity-affine-covariance}
设 \(R\) 为 Dedekind 域，取 \(u\in R\setminus\{0\}\)、\(v\in R\)，并令 \(a_i':=ua_i+v\)。
则对每个 \(n\le r\) 有
\[
\boxed{\;\mathfrak H_n(a',\omega)\ =\ u^{n(n-1)}\,\mathfrak H_n(a,\omega).\;}
\]
特别地若 \(u\in R^{\times}\) 为单位，则 \(\mathfrak H_n(a',\omega)=\mathfrak H_n(a,\omega)\)。
\end{conclusion}
\begin{proof}[证明要点]
对任意 \(|I|=n\)，有
\(
\Delta_I(a')=\prod_{i<j}(ua_j+v-ua_i-v)=u^{\binom{n}{2}}\Delta_I(a)
\)，
故
\(
\omega_I\Delta_I(a')^2=u^{n(n-1)}\omega_I\Delta_I(a)^2
\)。
由于 \(\mathfrak H_n\) 由这些生成元张成，结论成立。证毕。
\end{proof}

\begin{conclusion}[刚性下界的有限活跃素集]\label{con:xi-terminal-hankel-rigidity-finite-active-primes}
取 \(R=\ZZ\) 并固定数据 \((a_i,\omega_i)_{i=1}^r\)。
令
\[
\Xi:=\Bigl(\prod_{i=1}^r \omega_i\Bigr)\cdot\Bigl(\prod_{1\le i<j\le r}(a_j-a_i)\Bigr)\in\ZZ.
\]
则对任意素数 \(p\nmid \Xi\) 与任意 \(n\le r\) 有
\[
\boxed{\;v_p\!\bigl(\mathfrak H_n(a,\omega)\bigr)=0.\;}
\]
因此一切由几何数据 \((\omega_i,a_i)\) 强制产生的非平凡 \(p\)-进刚性下界，只可能发生在有限集合 \(\{p:\ p\mid \Xi\}\) 上。
\end{conclusion}
\begin{proof}[证明要点]
若 \(p\nmid \Xi\)，则每个 \(\omega_i\) 与每个差 \(a_j-a_i\) 皆为 \(p\)-进单位，因而任意生成元 \(\omega_I\Delta_I(a)^2\) 亦为 \(p\)-进单位。
故 \(\mathfrak H_n(a,\omega)\) 的 \(p\)-进指数为 \(0\)，即 \(v_p(\mathfrak H_n)=0\)。证毕。
\end{proof}

\begin{conjecture}[Hankel 原始性猜想]\label{conj:xi-terminal-hankel-primitiveness}
在结论 \ref{con:xi-terminal-hankel-rigidity-ideal-nonarchimedean} 的设定下，对每个 \(1\le n\le r\) 猜想恒有主理想恒等式
\[
\boxed{\ (D_n)\ =\ \mathfrak H_n(a,\omega).\ }
\]
等价地，对每个非零素理想 \(\mathfrak p\subset R\) 都有
\[
v_{\mathfrak p}(D_n)=v_{\mathfrak p}\!\bigl(\mathfrak H_n(a,\omega)\bigr).
\]
再等价地，令 \(R_{\mathfrak p}\) 的均值元为 \(\pi\)，并记
\(
e=v_{\mathfrak p}\!\bigl(\mathfrak H_n(a,\omega)\bigr)
=\min_{|I|=n}v_{\mathfrak p}(\omega_I\Delta_I(a)^2)
\)。
取极小集
\(
\mathcal M=\{I:\ |I|=n,\ v_{\mathfrak p}(\omega_I\Delta_I(a)^2)=e\}
\)，
则应有非消去性
\[
\sum_{I\in\mathcal M}\overline{\Bigl(\frac{\omega_I\Delta_I(a)^2}{\pi^{e}}\Bigr)}\ \neq\ 0
\qquad\text{于}\quad
R_{\mathfrak p}/(\pi)
\]
成立，从而 \(v_{\mathfrak p}(D_n)=e\)。
\end{conjecture}

\begin{proposition}[刚性下界的泛型锐性：例外集为真 Zariski 闭集]\label{prop:xi-terminal-hankel-rigidity-generic-sharpness-zariski}
在结论 \ref{con:xi-terminal-hankel-rigidity-ideal-nonarchimedean} 的记号下，固定 \(n\) 与素理想 \(\mathfrak p\)，并取局部离散赋值环 \(R_{\mathfrak p}\) 的均值元 \(\pi\)、剩余域 \(\kappa:=R_{\mathfrak p}/(\pi)\)。
令
\[
e:=v_{\mathfrak p}\!\bigl(\mathfrak H_n(a,\omega)\bigr)=\min_{|I|=n}v_{\mathfrak p}\!\bigl(\omega_I\Delta_I(a)^2\bigr),
\qquad
\mathcal M:=\Bigl\{I:\ |I|=n,\ v_{\mathfrak p}\!\bigl(\omega_I\Delta_I(a)^2\bigr)=e\Bigr\}.
\]
将 \(\overline{\omega_i}\in \kappa^\times\) 视为参数，并令
\[
F(\overline\omega_1,\dots,\overline\omega_r)
:=
\sum_{I\in\mathcal M}\overline{\left(\frac{\omega_I\Delta_I(a)^2}{\pi^{e}}\right)}
\ \in\ \kappa[\overline\omega_1^{\pm 1},\dots,\overline\omega_r^{\pm 1}]
\]
为由 \(\mathcal M\) 指标出的 Laurent 多项式。
则 \(F\not\equiv 0\)，从而零点集
\[
\mathcal Z:=\{(\overline\omega_1,\dots,\overline\omega_r)\in(\kappa^\times)^r:\ F(\overline\omega_1,\dots,\overline\omega_r)=0\}
\]
是 \((\kappa^\times)^r\) 的真 Zariski 闭子集。
特别地，在 Zariski 开稠密意义下有
\[
v_{\mathfrak p}(D_n)=e,
\]
即刚性下界在泛型权重参数下达到锐等号。
\end{proposition}
\leanverified{paper\_xi\_terminal\_hankel\_rigidity\_generic\_sharpness\_zariski}
\begin{proof}
由结论 \ref{con:xi-terminal-hankel-rigidity-ideal-nonarchimedean} 的同余式，
\(
v_{\mathfrak p}(D_n)>e
\)
当且仅当 \(F(\overline\omega_1,\dots,\overline\omega_r)=0\) 于 \(\kappa\)。
\par
对任意 \(I\in\mathcal M\)，因
\(
v_{\mathfrak p}(\omega_I\Delta_I(a)^2)=e
\)，故 \(\omega_I\Delta_I(a)^2/\pi^e\) 在 \(R_{\mathfrak p}\) 中为单位，其像
\(
\overline{\bigl(\omega_I\Delta_I(a)^2/\pi^e\bigr)}
\in \kappa^\times
\)
非零。
并且每一项对应的单项式恰为 \(\prod_{i\in I}\overline\omega_i\)，对不同 \(I\) 彼此不同。
因此 \(F\) 为具有非零系数的有限和，不能恒为零。
\par
非零 Laurent 多项式在 \((\kappa^\times)^r\) 上的零点集为真 Zariski 闭集，结论成立。证毕。
\end{proof}

\begin{proposition}[模 \texorpdfstring{$\pi^t$}{pi^t} 的不可逆性：Hankel 奇异迫使多解纤维]\label{prop:xi-terminal-hankel-modpi-noninvertibility-fiber}
保持结论 \ref{con:xi-terminal-hankel-rigidity-ideal-nonarchimedean} 的设定，并取 \(R_{\mathfrak p}\) 与 \(\pi\) 如上。
若对某个整数 \(t\ge 1\) 有
\[
v_{\mathfrak p}(D_n)\ \ge\ t,
\]
则 \(\mathsf H_n\) 在模 \(\pi^t\) 上不可逆：线性映射
\[
(\mathsf H_n\bmod \pi^t):\ (R_{\mathfrak p}/\pi^t)^{n}\to (R_{\mathfrak p}/\pi^t)^{n}
\]
具有非平凡核。
从而对任何依赖于在模 \(\pi^t\) 上求解 Hankel 线性系统的解码/重构步骤，一旦存在某个解，则该解在模 \(\pi^t\) 意义下并不唯一，而至少形成一个由 \(\ker(\mathsf H_n\bmod\pi^t)\) 平移得到的仿射多解纤维。
\leanverified{paper\_xi\_terminal\_hankel\_modpi\_noninvertibility\_fiber}
\end{proposition}
\begin{proof}
条件 \(v_{\mathfrak p}(D_n)\ge t\) 等价于 \(\det(\mathsf H_n)\in(\pi^t)\subset R_{\mathfrak p}\)，故 \(\det(\mathsf H_n\bmod\pi^t)=0\)，从而 \(\mathsf H_n\bmod\pi^t\) 不可逆并具有非零核。
若 \((\mathsf H_n\bmod\pi^t)x=b\) 存在解 \(x\)，则对任意 \(u\in\ker(\mathsf H_n\bmod\pi^t)\) 有 \((\mathsf H_n\bmod\pi^t)(x+u)=b\)，故解集为 \(x+\ker(\mathsf H_n\bmod\pi^t)\) 的陪集。证毕。
\end{proof}

\begin{theorem}[\texorpdfstring{$p$}{p}-进精度损耗律：Hankel 标准形的最小可恢复精度]\label{thm:xi-terminal-hankel-standard-form-padic-loss}
\leanverified{paper\_xi\_terminal\_hankel\_standard\_form\_padic\_loss}
设 \((a_n)_{n\ge 0}\subset\ZZ\) 具有 Hankel 秩 \(d\)。
取任意窗口偏移 \(\ell\ge 0\)，并令
\[
H_0:=\bigl(a_{\ell+i+j}\bigr)_{0\le i,j\le d-1},\qquad
H_1:=\bigl(a_{\ell+i+j+1}\bigr)_{0\le i,j\le d-1}.
\]
假设 \(\det(H_0)\neq 0\)，并定义 Hankel 标准形矩阵
\[
A:=H_0^{-1}H_1\in\QQ^{d\times d}.
\]
固定素数 \(p\)，令
\(
s:=v_p(\det H_0)
\)。
若观测到 \(H_0,H_1\) 的全部条目在模 \(p^E\) 意义下确定（等价于窗口数据 \(a_{\ell},\dots,a_{\ell+2d-1}\) 在模 \(p^E\) 意义下确定），
则 \(A\) 在 \(\QQ_p\) 中被确定到模 \(p^{E-s}\) 的精度；
并且一般情形下该损耗 \(s\) 不可改进。
\end{theorem}
\begin{proof}
由伴随矩阵恒等式
\[
A=H_0^{-1}H_1=\frac{\mathrm{adj}(H_0)\,H_1}{\det(H_0)},
\]
每个矩阵元 \(A_{ij}\) 可由 Cramer 法则表示为“分子行列式除以 \(\det(H_0)\)”。
分子是 \(H_0,H_1\) 条目的整系数多项式，因此当 \(H_0,H_1\) 在模 \(p^E\) 下确定时，各分子在模 \(p^E\) 下确定。
而分母 \(\det(H_0)\) 的 \(p\)-进赋值为 \(s\)，故在 \(\QQ_p\) 中除法至多损失 \(s\) 位精度，从而 \(A\) 在模 \(p^{E-s}\) 意义下确定。
\par
不可改进性：在泛型情形下可取某个 \(A_{ij}\) 使其分子行列式对 \(p\) 为单位（\(v_p=0\)）。
对 \(H_1\) 中相应条目施加 \(p^E\) 级别扰动将引起分子同阶扰动；除以 \(\det(H_0)\) 后恰产生 \(p^{E-s}\) 级别的 \(A_{ij}\) 扰动，
因此无法保证优于 \(E-s\) 的恢复精度。证毕。
\end{proof}

\leanverified{paper\_xi\_terminal\_hankel\_standard\_form\_p5\_threshold}
\begin{corollary}[\texorpdfstring{$p=5$}{p=5} 的刚性阈值：七位损耗的不可消去性]\label{cor:xi-terminal-hankel-standard-form-p5-threshold}
在表 \ref{tab:fold_collision_moment_hankel_rank} 的 \(k=5\) Hankel 见证块 \(H_{k=5}\) 中，
\[
\det(H_{k=5})=2^{13}\cdot 3^5\cdot 5^7,
\]
故 \(v_5(\det(H_{k=5}))=7\)。
因此若仅以 \(5\)-进精度 \(5^E\) 观测到对应的 \(2d\) 项窗口数据，则由定理 \ref{thm:xi-terminal-hankel-standard-form-padic-loss}，
相应 Hankel 标准形矩阵至多可恢复到 \(5^{E-7}\) 的精度；该七位损耗一般不可消去。
\end{corollary}

\begin{theorem}[有限精度势的 Smith 计数律：\texorpdfstring{$\ker/\mathrm{coker}$}{ker/coker} 的阶与分支计数]\label{thm:xi-hankel-finite-precision-potential-smith-kernel-coker}
保持定理 \ref{thm:xi-terminal-hankel-standard-form-padic-loss} 的设定，并假设 \(H_0\in M_d(\ZZ)\) 满秩。
取其 Smith 正规形
\[
UH_0V=\mathrm{diag}(s_1,\dots,s_d),\qquad
s_1\mid s_2\mid\cdots\mid s_d,\quad U,V\in \mathrm{GL}_d(\ZZ).
\]
对素数 \(p\) 与整数 \(k\ge 1\) 定义有限精度势
\[
H_p(k):=\sum_{i=1}^{d}\min\bigl(v_p(s_i),k\bigr).
\]
则有完全显式恒等式
\[
\#\ker\!\Bigl(H_0:(\ZZ/p^k\ZZ)^d\to(\ZZ/p^k\ZZ)^d\Bigr)=p^{H_p(k)},
\]
并且同样
\[
\#\mathrm{coker}\!\Bigl(H_0:(\ZZ/p^k\ZZ)^d\to(\ZZ/p^k\ZZ)^d\Bigr)=p^{H_p(k)}.
\]
特别地，
\[
\lim_{k\to\infty}H_p(k)=\sum_{i=1}^{d}v_p(s_i)=v_p(\det H_0).
\]
\end{theorem}
\begin{proof}
在 \(\ZZ/p^k\ZZ\) 上，\(U,V\) 仍为可逆变换，故核与余核的阶不变；因此可将 \(H_0\) 替换为对角矩阵 \(\mathrm{diag}(s_i)\)。
对单坐标映射 \(x\mapsto s_ix\) 于 \(\ZZ/p^k\ZZ\) 上，其核大小为 \(p^{\min(v_p(s_i),k)}\)；
核在直积上相乘即得第一式。
对满秩方阵在有限环上，核与余核的阶相等；在对角情形下可逐坐标直接验证，从而得到第二式。
极限式由 \(\min(v,k)\to v\) 逐项收敛与 \(\prod_i s_i=\pm\det(H_0)\) 得到。证毕。
\end{proof}
\leanverified{paper\_xi\_hankel\_finite\_precision\_potential\_smith\_kernel\_coker}

\begin{corollary}[模 \texorpdfstring{$p^k$}{p^k} 的矩阵重建多解度：右核的 \texorpdfstring{$d$}{d} 次直和]\label{cor:xi-hankel-lifting-branch-count-affine-solution-space}
令 \(H_1'\in M_d(\ZZ)\) 给定，并考虑线性方程 \(H_0B=H_1'\) 在模 \(p^k\) 下的解集
\[
\mathcal S_k:=\Bigl\{X\in M_d(\ZZ/p^k\ZZ):\ (H_0\bmod p^k)\,X\equiv (H_1'\bmod p^k)\Bigr\}.
\]
若 \(\mathcal S_k\neq\varnothing\)，则 \(\mathcal S_k\) 为一个仿射空间，其平移空间为右核 \(\ker(H_0\bmod p^k)\) 的 \(d\) 次直和；
并且其基数满足
\[
\#\mathcal S_k=p^{d\,H_p(k)},
\]
其中 \(H_p(k)\) 如定理 \ref{thm:xi-hankel-finite-precision-potential-smith-kernel-coker}。
\end{corollary}
\begin{proof}
考虑线性映射
\[
\Phi_k:M_d(\ZZ/p^k\ZZ)\to M_d(\ZZ/p^k\ZZ),\qquad
\Phi_k(X)=(H_0\bmod p^k)\,X.
\]
其核由每一列独立满足 \((H_0\bmod p^k)x=0\) 的条件组成，故
\(
\#\ker\Phi_k=(\#\ker(H_0\bmod p^k))^d
\)。
由定理 \ref{thm:xi-hankel-finite-precision-potential-smith-kernel-coker} 得 \(\#\ker(H_0\bmod p^k)=p^{H_p(k)}\)，从而 \(\#\ker\Phi_k=p^{dH_p(k)}\)。
若 \(\mathcal S_k\) 非空，则它是任一特解加上 \(\ker\Phi_k\) 的陪集，故其基数等于 \(\#\ker\Phi_k\)。证毕。
\end{proof}
\leanverified{paper\_xi\_hankel\_lifting\_branch\_count\_affine\_solution\_space}


\begin{theorem}[Hankel 确定子理想的根等价]\label{thm:xi-hankel-determinantal-radical-eq-rigidity}
\leanverified{paper\_xi\_hankel\_determinantal\_radical\_eq\_rigidity}
令 \(k\) 为代数闭域。
取形式变量 \(a_1,\dots,a_r,\omega_1,\dots,\omega_r\) 并定义矩量
\[
s_m:=\sum_{i=1}^{r}\omega_i a_i^{m}\qquad(m\ge 0),
\]
以及无穷 Hankel 矩阵 \(H=(s_{i+j})_{i,j\ge 0}\)。
给定整数 \(n\le r\)，令 \(I_n\subset k[a_1,\dots,a_r,\omega_1,\dots,\omega_r]\) 为由 \(H\) 的所有 \(n\times n\) 子式生成的确定子理想。
另一方面，令
\[
\mathfrak H_n(a,\omega):=\big\langle\,\omega_I\Delta_I(a)^2:\ I\subseteq\{1,\dots,r\},\ |I|=n\,\big\rangle,
\qquad
\omega_I:=\prod_{i\in I}\omega_i,\ \ 
\Delta_I(a):=\prod_{\substack{i<j\\ i,j\in I}}(a_j-a_i).
\]
则有根理想恒等式
\[
\sqrt{I_n}=\sqrt{\mathfrak H_n(a,\omega)}.
\]
\end{theorem}
\begin{proof}
先证 \(I_n\subseteq \mathfrak H_n(a,\omega)\)。
对任意有限行列索引集 \(P,Q\subset \ZZ_{\ge 0}\) 且 \(|P|=|Q|=n\)，记相应 Hankel 子矩阵为 \(H_{P,Q}\)。
令 \(V\) 为无穷 Vandermonde 矩阵 \(V_{t,i}:=a_i^{t}\)，并令 \(W=\mathrm{diag}(\omega_1,\dots,\omega_r)\)，则有分解 \(H=V W V^{\mathsf T}\)。
由 Cauchy--Binet，
\[
\det(H_{P,Q})
=\sum_{|I|=n}\omega_I\,\det(V_{P,I})\,\det(V_{Q,I}).
\]
对固定 \(I\)，\(\det(V_{P,I})\) 与 \(\det(V_{Q,I})\) 作为 \(a_I\) 的交错多项式，均被 \(\Delta_I(a)\) 整除，从而每一项都含因子 \(\omega_I\Delta_I(a)^2\)，故 \(\det(H_{P,Q})\in \mathfrak H_n(a,\omega)\)。
于是 \(I_n\subseteq \mathfrak H_n(a,\omega)\)，从而 \(\sqrt{I_n}\subseteq \sqrt{\mathfrak H_n(a,\omega)}\)。
\par
再证 \(V(\mathfrak H_n(a,\omega))\subseteq V(I_n)\)。
设一点 \((a,\omega)\) 满足对所有 \(|I|=n\) 有 \(\omega_I\Delta_I(a)^2=0\)。
令 \(T\) 为满足 \(\omega_i\neq 0\) 且 \(a_i\) 两两不同的指标集合。
若 \(|T|\ge n\)，则可取 \(I\subseteq T\) 且 \(|I|=n\)，此时 \(\omega_I\neq 0\) 且 \(\Delta_I(a)\neq 0\)，矛盾。
故必有 \(|T|\le n-1\)，从而矩序列 \((s_m)\) 实际上是至多 \(n-1\) 个互异指数项的有限和，Hankel 秩不超过 \(n-1\)，因此所有 \(n\times n\) Hankel 子式均为零，即点属于 \(V(I_n)\)。
于是 \(\sqrt{\mathfrak H_n(a,\omega)}\subseteq \sqrt{I_n}\)。
两侧互含即得结论。证毕。
\end{proof}

\begin{corollary}[首 Hankel 主块行列式的判别式--权重分解]\label{cor:xi-hankel-first-block-determinant-discriminant-weight}
\leanverified{paper\_xi\_hankel\_first\_block\_determinant\_discriminant\_weight}
在定理 \ref{thm:xi-hankel-determinantal-radical-eq-rigidity} 的设定下，取 \(n=r\) 并令首 Hankel 主块
\[
H_0:=\bigl(s_{i+j}\bigr)_{0\le i,j\le r-1}.
\]
则在 \(k\) 中成立显式分解
\[
\det(H_0)
=
\Bigl(\prod_{i=1}^{r}\omega_i\Bigr)\cdot
\prod_{1\le i<j\le r}(a_j-a_i)^2.
\]
特别地，若在某一赋值环 \(\mathcal O\) 中取 \(a_i,\omega_i\in\mathcal O\) 并把上式视为分裂域中的恒等式，则任一使得模 \(\mathfrak p\) 约化后主块秩发生下降的素理想 \(\mathfrak p\) 必整除 \(\prod_i\omega_i\) 或 Vandermonde 判别式 \(\prod_{i<j}(a_j-a_i)^2\)。
\end{corollary}
\begin{proof}
对 \(H_0\) 取 \(P=Q=\{0,1,\dots,r-1\}\)，并应用定理 \ref{thm:xi-hankel-determinantal-radical-eq-rigidity} 证明中给出的 Cauchy--Binet 展开式。
当 \(n=r\) 时，求和中只出现唯一集合 \(I=\{1,\dots,r\}\)。
此时 \(V_{P,I}\) 与 \(V_{Q,I}\) 均为标准 Vandermonde 矩阵，\(\det(V_{P,I})=\det(V_{Q,I})=\prod_{i<j}(a_j-a_i)\)，故
\(
\det(H_0)=\omega_{I}\det(V_{P,I})\det(V_{Q,I})
\)
即为所示分解。
末尾断言由“秩下降 \(\Rightarrow\) 相应主子式为零”与分解式直接推出。证毕。
\end{proof}

\input{cor__xi-fiber-spectrum-hankel-conductor-prony-padic}

\begin{proposition}[无权刚性理想的四次方结构与 \texorpdfstring{$p$}{p}-ordering 赋值律]\label{prop:xi-unweighted-rigidity-quartic-ideal-pordering}
设 \(R\) 为 Dedekind 域，取两两不同的 \(a_1,\dots,a_r\in R\)。
对 \(n\le r\) 定义
\[
J_n(a):=\big\langle\,\Delta_I(a):\ |I|=n\,\big\rangle\subset R,
\qquad
\mathfrak H_n(a,1):=\big\langle\,\Delta_I(a)^2:\ |I|=n\,\big\rangle\subset R,
\]
以及四次理想
\[
K_n(a):=\big\langle\,\Delta_I(a)^4:\ |I|=n\,\big\rangle\subset R.
\]
则有理想恒等式
\[
\boxed{\;K_n(a)=J_n(a)^4=\mathfrak H_n(a,1)^2.\;}
\]
当 \(R=\ZZ\) 且固定素数 \(p\) 时，令 \(S=\{a_1,\dots,a_r\}\)，取任意 \(p\)-ordering（Bhargava）并记相应 \(p\)-序列为 \(V_k(S,p)\)。
则对任意 \(n\le r\) 有赋值公式
\[
v_p\!\bigl(K_n(a)\bigr)=4\sum_{k=1}^{n-1} v_p\!\bigl(V_k(S,p)\bigr).
\]
\end{proposition}
\begin{proof}
对任意素理想 \(\mathfrak p\subset R\)，局部化 \(R_{\mathfrak p}\) 为 DVR。
设
\(
e:=\min_{|I|=n} v_{\mathfrak p}(\Delta_I(a))
\)，
则在 \(R_{\mathfrak p}\) 中
\[
J_n(a)R_{\mathfrak p}=\mathfrak p^{e},
\qquad
K_n(a)R_{\mathfrak p}=\langle \Delta_I(a)^4\rangle R_{\mathfrak p}=\mathfrak p^{4e}=\bigl(J_n(a)R_{\mathfrak p}\bigr)^4.
\]
由于 Dedekind 域中非零理想由各素理想指数唯一确定，得到 \(K_n(a)=J_n(a)^4\)。
再由结论 \ref{con:xi-terminal-hankel-rigidity-unweighted-bhargava} 的 \(\mathfrak H_n(a,1)=J_n(a)^2\) 得 \(K_n(a)=\mathfrak H_n(a,1)^2\)。
\par
当 \(R=\ZZ\) 时，结论 \ref{con:xi-terminal-hankel-rigidity-unweighted-bhargava} 已给出
\(
v_p(J_n(a))=\sum_{k=1}^{n-1}v_p(V_k(S,p))
\)，
故 \(v_p(K_n(a))=4v_p(J_n(a))\)，得到所示公式。证毕。
\end{proof}


\begin{corollary}[POM 最小实现与谱见证的 \texorpdfstring{$p$}{p}-进预算分离]\label{cor:xi-pom-minreal-vs-spectral-budget-separation}
设 \((a_m)_{m\ge 0}\subset\ZZ\) 的 Hankel 秩为 \(d\)，并且主块
\[
H_0:=\bigl(a_{i+j}\bigr)_{0\le i,j\le d-1}
\]
可逆（等价于窗口数据 \((a_0,\dots,a_{2d-1})\) 处于非退化区间）。
则由定理 \ref{thm:xi-terminal-hankel-standard-form-padic-loss}，最小实现的 Hankel 标准形
\(
A=H_0^{-1}H_1
\)
在 \(p\)-进意义下的最小可恢复精度门槛由 \(v_p(\det H_0)\) 控制。
\par
进一步，若同一窗口数据来自一组互异谱参数的指数和表示
\[
a_m=\sum_{i=1}^{d}\omega_i\lambda_i^{m},
\qquad
\lambda_i\neq \lambda_j,\ \ \omega_i\neq 0,
\]
并将 \((a_0,\dots,a_{2d-1})\) 视作 Prony 矩映射 \(\Phi_d\) 的输出，则谱参数 \((\lambda_i,\omega_i)\) 的局部稳定恢复所需 \(p\)-进精度门槛至少为
\[
v_p\!\bigl(\det J\bigr)
\;=\;
v_p(\det H_0)+2\sum_{1\le i<j\le d} v_p(\lambda_j-\lambda_i),
\]
其中 \(J\) 为 \(\Phi_d\) 的雅可比矩阵。
因此当谱在模 \(p\) 意义下发生碰撞（即 \(\sum_{i<j} v_p(\lambda_j-\lambda_i)>0\)）时，谱见证提取的预算严格强于最小实现证书恢复。
\end{corollary}
\begin{proof}
第一部分即定理 \ref{thm:xi-terminal-hankel-standard-form-padic-loss}。
对谱参数部分，将 \(a_i:=\lambda_i\)、\(\omega_i\) 代入推论 \ref{cor:xi-hankel-vs-prony-square-gap}，并注意 \(D_d=\det H_0\)（对应 \(d\times d\) Hankel 主块），得到
\(
v_p(\det J)=v_p(\det H_0)+2\sum_{i<j}v_p(\lambda_j-\lambda_i)
\)。
谱参数的局部反演阈值由定理 \ref{thm:xi-nonarchimedean-prony-hensel-newton-inversion} 的 Hensel 条件给出，从而得出预算分离。证毕。
\end{proof}


\begin{conjecture}[非阿基米德极小子集唯一性：\texorpdfstring{$p$}{p}-骨架的单极小化]\label{conj:xi-nonarchimedean-minimizer-skeleton-uniqueness}
在结论 \ref{con:xi-terminal-hankel-rigidity-ideal-nonarchimedean} 的设定下，
对任意非零素理想 \(\mathfrak p\subset R\) 与任意 \(1\le n\le r\)，令
\[
\mathcal M(\mathfrak p,n)
:=\Bigl\{\,I\subseteq\{1,\dots,r\}:\ |I|=n,\ v_{\mathfrak p}\!\bigl(\omega_I\Delta_I(a)^2\bigr)=v_{\mathfrak p}\!\bigl(\mathfrak H_n(a,\omega)\bigr)\,\Bigr\}.
\]
猜想：在 Zariski 开稠密的参数集合上，对所有 \((\mathfrak p,n)\) 均有 \(\bigl|\mathcal M(\mathfrak p,n)\bigr|=1\)。
等价地，每个 \((\mathfrak p,n)\) 存在唯一的达到刚性指数极小值的 \(n\)-子骨架，且不会出现并列极小。
\end{conjecture}



\begin{conclusion}[Toeplitz--PSD 证书的最小失败定位：等价于 CMV 反射谱前缀越界]\label{con:xi-terminal-toeplitz-psd-failure-cmv-prefix}
令 \(\mathbf{T}_m\) 为由视界 Carath\'eodory 矩 \((c_n)\) 生成的 Toeplitz 主子矩阵，并令 \((\alpha_n)_{n\ge 0}\) 为相应 Schur 算法给出的反射系数（Verblunsky 参数）。
则 Toeplitz--PSD 层级的失效可由最小越界反射系数精确定位：存在最小 \(m^\ast\ge 1\) 使
\[
\boxed{\ |\alpha_{m^\ast-1}|\ge 1\ },
\]
且该 \(m^\ast\) 等价刻画“最小 PSD 失败阶”（在更低阶 \(m<m^\ast\) 时仍保持 \(\mathbf{T}_m\succeq 0\)）。
因此否证器出现的位置被压缩为可审计的最小前缀指标，而不依赖全局数值谱排序。
\end{conclusion}
\begin{proof}[证明要点]
见定理 \ref{thm:xi-toeplitz-det-verblunsky}（Toeplitz 行列式反射乘积分解与最小失败索引）及其后随的讨论。
\end{proof}

\begin{conclusion}[Chebyshev--Frobenius 的整系数 \texorpdfstring{$\delta_p$}{delta_p}-结构：Adams 倍角成为 Frobenius-lift]\label{con:xi-terminal-chebyshev-frobenius-pderivation}
在边界坐标 \(S=2\cos(\theta/2)\) 下，幂映射 \(w\mapsto w^{m}\) 等价于 \(S\mapsto C_m(S)\)（Chebyshev--Adams 多项式）。
对素数 \(p\) 定义差额多项式
\[
\Delta_p(S):=\frac{C_p(S)-S^p}{p}\in\ZZ[S],
\]
并对任意 \(f\in\ZZ[S]\) 定义
\[
\delta_p(f):=\frac{f(C_p(S))-\bigl(f(S)\bigr)^p}{p}.
\]
则有整系数闭包与 Frobenius-lift 恒等式
\[
\boxed{\ \delta_p(f)\in\ZZ[S],\qquad f(C_p(S))=\bigl(f(S)\bigr)^p+p\,\delta_p(f)\ },
\]
并且模 \(p\) 层面严格退化为 \(f(C_p(S))\equiv f(S)^p\pmod p\)。
因此倍角/幂半群的完成化作用可被提升为 \(p\)-进同余护栏的代数源头。
\end{conclusion}
\begin{proof}[证明要点]
见定义 \ref{def:chebyshev-frobenius-pderivation} 与命题 \ref{prop:chebyshev-frobenius-pderivation-integrality}。
\end{proof}

\begin{conclusion}[Witt 伸缩的 Koenigs 线性化：\texorpdfstring{$s\mapsto ms$}{s->ms} 在深度线 \texorpdfstring{$y$}{y} 上严格平移]\label{con:xi-terminal-witt-koenigs-logm-translation}
在 Dirichlet--Abel 接口 \(z=\lambda^{-s}\) 下，Witt 伸缩 \(s\mapsto ms\) 与幂替换 \(z\mapsto z^m\) 严格等价。
定义无参深度坐标
\[
y(s):=\log\!\bigl(s\log\lambda\bigr)=\log\!\bigl(-\log(\lambda^{-s})\bigr),
\]
则对任意整数 \(m\ge 1\) 有严格恒等式
\[
\boxed{\ y(ms)=y(s)+\log m\ }.
\]
因此 Witt 伸缩在 \(y\)-坐标中被完全对角化为平移半群；相应 Witt 算子在深度线 \(y\) 上成为由离散测度 \(\nu=\sum_{m\ge 1}\frac{1}{m}\delta_{\log m}\) 生成的平移叠加算子。
\end{conclusion}
\begin{proof}[证明要点]
见定理 \ref{thm:koenigs-linearization-witt-dilation} 与推论 \ref{cor:witt-depth-line-convolution}。
\end{proof}

\begin{conclusion}[window-\texorpdfstring{$6$}{6} 的局域 uplift 选择律：\texorpdfstring{$SO(10)$}{SO(10)} 分支为唯一保局域的顶层候选]\label{con:xi-terminal-window6-local-uplift-so10-unique}
在保持稳定类型标签 \(F_6\) 不变且由 \(\mathrm{Aut}(Q_6)\) 实现的强局域 uplift 口径下，除平凡位移外仅允许 \(\delta=\pm 34\)（推论 \ref{cor:terminal-window6-local-uplift-admissibility}）。
另一方面，window-\(6\) uplift 的尾部三分支集合为 \(\{21,34,55\}=\{F_8,F_9,F_{10}\}\)，并与 \((SU(5),SO(10),E_6)\) 的顶层 Zeckendorf 项分别对齐为 \((F_8,F_9,F_{10})\)（定理 \ref{thm:terminal-window6-tail-three-branch}）。
因此在附加施加“uplift 必须局域可实现且标签不变”的协议约束时，\(\delta=34\) 分支（\(F_9\)）是唯一与局域性相容的顶层补全；其余两支的顶层对齐只能通过非局域机制外置实现。
\end{conclusion}
\begin{proof}[证明要点]
局域可实现性与 \(\delta=\pm 34\) 的唯一性见推论 \ref{cor:terminal-window6-local-uplift-admissibility}；
三分支与 Zeckendorf--GUT 顶层项对齐见定理 \ref{thm:terminal-window6-tail-three-branch}（并与家族数锁定律 \ref{thm:terminal-family-uplift-lock} 的 \(F_9\) 选中一致）。
\end{proof}

\begin{conclusion}[bin-fold 推前分布的末位充分化：全变差塌缩与单参数指数倾斜]\label{con:xi-terminal-foldbin-lastbit-sufficient-statistic}
令 \(\mu_m(w):=d_m^{\mathrm{bin}}(w)/2^m\) 为 bin-fold 在 \(X_m\) 上的推前分布，并令末位指示 \(\ell(w):=w_m\in\{0,1\}\)。
将 \(X_m\) 按末位分块并令 \(\overline\mu_m\) 为块内均匀化近似，则存在绝对常数 \(C>0\) 使
\[
\boxed{\ D_{\mathrm{TV}}(\mu_m,\overline\mu_m)\ \le\ C\left(\frac{\varphi}{2}\right)^{m}\ }.
\]
并且存在 \(\beta_m\in\RR\) 使 \(\overline\mu_m\) 可写成稳定层均匀分布 \(u_m\) 的指数倾斜
\[
\overline\mu_m(w)=Z_m^{-1}u_m(w)e^{-\beta_m\ell(w)},
\qquad
\beta_m\to\log\varphi.
\]
因此末位 \(\ell(w)\) 在指数精度上成为充分统计量，bin-fold 输出分布在信息几何意义下被压缩为单参数指数族。
\end{conclusion}
\begin{proof}[证明要点]
见命题 \ref{prop:fold-bin-lastbit-sufficient-tv}。
\end{proof}

\begin{conclusion}[escort 温度族与散度谱的二态闭式化：温度间 KL 的 Bernoulli 退化与 Rényi 熵率塌缩]\label{con:xi-terminal-foldbin-escort-divergence-two-state-closure}
对 \(q\ge 0\) 定义 escort 分布 \(\pi_{m,q}(w)\propto (d_m^{\mathrm{bin}}(w))^{q}\)，并令
\(
\theta(q):=\lim_{m\to\infty}\mathbb{P}_{\pi_{m,q}}(w_m=1)
=\frac{1}{1+\varphi^{q+1}}
\)。
则对任意 \(p,q\ge 0\)，温度间 KL 极限存在且严格退化为一维 Bernoulli KL：
\[
\boxed{\;
\lim_{m\to\infty}D_{\mathrm{KL}}(\pi_{m,q}\Vert\pi_{m,p})
=\theta(q)\log\!\frac{\theta(q)}{\theta(p)}+(1-\theta(q))\log\!\frac{1-\theta(q)}{1-\theta(p)}\; }.
\]
同时，bin-fold 推前分布 \(\mu_m\) 相对稳定层均匀基线的 Rényi 散度极限对每个 \(\alpha>0\) 存在且给出闭式表达，并具有 \(O((\varphi/2)^m)\) 的指数收敛（定理 \ref{thm:fold-bin-renyi-divergence-limit}）。
并且 Rényi 熵密度满足统一塌缩：对一切 \(q>0\) 有
\[
\boxed{\ \lim_{m\to\infty}\frac{1}{m}H_q(\mu_m)=\log\varphi\ },
\]
从而“多温度/多散度/多分形”结构在极限中被制度性压缩为末位 Bernoulli 参数 \(\theta(q)\) 的一维谱。
\end{conclusion}
\begin{proof}[证明要点]
Bernoulli KL 退化见命题 \ref{prop:fold-bin-escort-escort-kl}；
Rényi 散度极限见定理 \ref{thm:fold-bin-renyi-divergence-limit}；
Rényi 熵率塌缩见定理 \ref{thm:fold-bin-renyi-rate-collapse}。
\end{proof}

\paragraph{escort 温度流的总变差压缩、散度闭式与 Fisher 平直化}
在末位充分化已经把 escort 温度族压缩到一维 Bernoulli 参数之后，总变差塌缩、Rényi--KL 散度几何与 Fisher 度量都可沿同一末位两块分解完全显式化。

\begin{theorem}[escort 温度族在总变差意义下塌缩为一维 Bernoulli 指数族]\label{thm:xi-foldbin-escort-tv-collapse-onebit-gibbs}
对任意 \(q\ge 0\)，定义 escort 分布
\[
\pi_{m,q}(w):=\frac{(d_m^{\mathrm{bin}}(w))^q}{S_q^{\mathrm{bin}}(m)},
\qquad
S_q^{\mathrm{bin}}(m):=\sum_{u\in X_m}(d_m^{\mathrm{bin}}(u))^q.
\]
并按末位将 \(X_m\) 分成
\[
A_0:=\{w\in X_m:\ w_m=0\},
\qquad
A_1:=\{w\in X_m:\ w_m=1\}.
\]
记
\[
\theta(q):=\frac{1}{1+\varphi^{q+1}},
\]
并定义块内均匀近似
\[
\bar\pi_{m,q}(w):=
\frac{1-\theta(q)}{|A_0|}\mathbf 1_{A_0}(w)
+
\frac{\theta(q)}{|A_1|}\mathbf 1_{A_1}(w).
\]
则对每个固定 \(q\ge 0\)，存在常数 \(C_q>0\) 使
\[
\|\pi_{m,q}-\bar\pi_{m,q}\|_{\mathrm{TV}}
\le
C_q\Bigl(\frac{\varphi}{2}\Bigr)^m.
\]
进一步，\(\bar\pi_{m,q}\) 可写成稳定层均匀分布 \(u_m(w):=1/|X_m|\) 的单参数指数倾斜
\[
\bar\pi_{m,q}(w)=Z_{m,q}^{-1}u_m(w)e^{-\beta_{m,q}w_m},
\]
且
\[
\beta_{m,q}\longrightarrow q\log\varphi.
\]
因此整个 escort 温度族在全变差精度下严格塌缩为由末位比特 \(w_m\) 支配的一维 Bernoulli 指数族。
\end{theorem}
\begin{proof}
由定理 \ref{thm:fold-bin-two-state-asymptotic}，对 \(i\in\{0,1\}\) 与 \(w\in A_i\) 一致有
\[
d_m^{\mathrm{bin}}(w)
=
\left(\frac{2}{\varphi}\right)^m
\varphi^{-i}
\Bigl(1+O\!\Bigl(\Bigl(\frac{\varphi}{2}\Bigr)^m\Bigr)\Bigr).
\]
对固定 \(q\ge 0\) 取 \(q\) 次幂可得
\[
(d_m^{\mathrm{bin}}(w))^q
=
\left(\frac{2}{\varphi}\right)^{qm}
\varphi^{-qi}
\Bigl(1+O_q\!\Bigl(\Bigl(\frac{\varphi}{2}\Bigr)^m\Bigr)\Bigr),
\]
且误差在块内一致。于是
\[
\sum_{w\in A_i}(d_m^{\mathrm{bin}}(w))^q
=
|A_i|
\left(\frac{2}{\varphi}\right)^{qm}
\varphi^{-qi}
\Bigl(1+O_q\!\Bigl(\Bigl(\frac{\varphi}{2}\Bigr)^m\Bigr)\Bigr).
\]
再用 \(|A_0|=F_{m+1}\)、\(|A_1|=F_m\) 以及
\(
F_m/F_{m+1}=\varphi^{-1}+O(\varphi^{-2m})
\)，得到
\[
\pi_{m,q}(A_1)=\theta(q)+O_q\!\Bigl(\Bigl(\frac{\varphi}{2}\Bigr)^m\Bigr),
\qquad
\pi_{m,q}(A_0)=1-\theta(q)+O_q\!\Bigl(\Bigl(\frac{\varphi}{2}\Bigr)^m\Bigr).
\]
另一方面，对 \(w\in A_i\) 有
\[
\pi_{m,q}(w)
=
\frac{\pi_{m,q}(A_i)}{|A_i|}
\Bigl(1+O_q\!\Bigl(\Bigl(\frac{\varphi}{2}\Bigr)^m\Bigr)\Bigr).
\]
将其与 \(\bar\pi_{m,q}(w)\) 比较，并在两块上分别求和，即得
\[
\sum_{w\in X_m}\bigl|\pi_{m,q}(w)-\bar\pi_{m,q}(w)\bigr|
=
O_q\!\Bigl(\Bigl(\frac{\varphi}{2}\Bigr)^m\Bigr),
\]
从而给出总变差界。

对指数倾斜表示，只需令
\[
e^{-\beta_{m,q}}
:=
\frac{\bar\pi_{m,q}(A_1)}{\bar\pi_{m,q}(A_0)}
\cdot
\frac{|A_0|}{|A_1|}
=
\frac{\theta(q)}{1-\theta(q)}
\cdot
\frac{F_{m+1}}{F_m}.
\]
由于
\(
\theta(q)/(1-\theta(q))=\varphi^{-(q+1)}
\)
且 \(F_{m+1}/F_m\to\varphi\)，故
\[
e^{-\beta_{m,q}}\to \varphi^{-q},
\]
即 \(\beta_{m,q}\to q\log\varphi\)。证毕。
\end{proof}
\leanverified{paper_xi_foldbin_escort_tv_collapse_onebit_gibbs}

\begin{theorem}[escort 参数在总变差上保持严格可辨识]\label{thm:xi-foldbin-escort-tv-identifiable-limit-distance}\leanverified{paper\_xi\_foldbin\_escort\_tv\_identifiable\_limit\_distance}
沿用定理 \ref{thm:xi-foldbin-escort-tv-collapse-onebit-gibbs} 的记号。则对任意固定 \(p,q\ge 0\)，有
\[
\|\pi_{m,q}-\pi_{m,p}\|_{\mathrm{TV}}
=
|\theta(q)-\theta(p)|
+O_{p,q}\!\Bigl(\Bigl(\frac{\varphi}{2}\Bigr)^m\Bigr),
\]
从而
\[
\lim_{m\to\infty}\|\pi_{m,q}-\pi_{m,p}\|_{\mathrm{TV}}
=
|\theta(q)-\theta(p)|.
\]
特别地，极限参数可由单一标量 \(\theta(q)\) 反演为
\[
q=\log_\varphi\!\left(\frac{1-\theta(q)}{\theta(q)}\right)-1.
\]
\end{theorem}
\begin{proof}
由三角不等式，
\[
\Bigl|
\|\pi_{m,q}-\pi_{m,p}\|_{\mathrm{TV}}
-\|\bar\pi_{m,q}-\bar\pi_{m,p}\|_{\mathrm{TV}}
\Bigr|
\le
\|\pi_{m,q}-\bar\pi_{m,q}\|_{\mathrm{TV}}
+\|\pi_{m,p}-\bar\pi_{m,p}\|_{\mathrm{TV}},
\]
而右端由定理 \ref{thm:xi-foldbin-escort-tv-collapse-onebit-gibbs} 为
\(
O_{p,q}((\varphi/2)^m)
\)。

另一方面，\(\bar\pi_{m,q}\) 与 \(\bar\pi_{m,p}\) 都在同一分块 \(A_0\sqcup A_1\) 上块内均匀，且两块质量分别为
\(
(1-\theta(q),\theta(q))
\)
与
\(
(1-\theta(p),\theta(p))
\)。
因此
\[
\|\bar\pi_{m,q}-\bar\pi_{m,p}\|_{\mathrm{TV}}
=
|\theta(q)-\theta(p)|.
\]
代回即得所示展开与极限。
最后由
\(
\theta(q)=1/(1+\varphi^{q+1})
\)
直接解得反演公式。证毕。
\end{proof}

\begin{theorem}[escort 温度族之间的 Rényi--KL 几何在极限中封闭为显式二点公式]\label{thm:xi-foldbin-escort-renyi-kl-two-point-closure}
沿用定理 \ref{thm:xi-foldbin-escort-tv-collapse-onebit-gibbs} 的记号。对 \(q,r\ge 0\) 与
\(\alpha>0\)、\(\alpha\neq 1\)，定义 Rényi 散度
\[
D_\alpha(\pi_{m,q}\Vert \pi_{m,r})
:=
\frac{1}{\alpha-1}
\log\sum_{w\in X_m}\pi_{m,q}(w)^\alpha \pi_{m,r}(w)^{1-\alpha}.
\]
则极限
\[
\lim_{m\to\infty}D_\alpha(\pi_{m,q}\Vert \pi_{m,r})
=:\mathcal D_\alpha(q,r)
\]
存在，且满足
\[
\mathcal D_\alpha(q,r)
=
\frac{1}{\alpha-1}
\log\!\Bigl(
(1-\theta(q))^\alpha(1-\theta(r))^{1-\alpha}
+
\theta(q)^\alpha\theta(r)^{1-\alpha}
\Bigr).
\]
等价地，
\[
\mathcal D_\alpha(q,r)
=
\frac{1}{\alpha-1}
\log
\frac{\varphi+\varphi^{-\alpha q-(1-\alpha)r}}
{(\varphi+\varphi^{-q})^\alpha(\varphi+\varphi^{-r})^{1-\alpha}}.
\]
进一步，令 \(\alpha\to 1\) 得到极限 KL 散度的闭式
\[
\lim_{m\to\infty}D_{\mathrm{KL}}(\pi_{m,q}\Vert \pi_{m,r})
=
\log\!\frac{1+\varphi^{r+1}}{1+\varphi^{q+1}}
+
\frac{\varphi^{q+1}}{1+\varphi^{q+1}}(q-r)\log\varphi.
\]
因此 escort 温度族的双参数散度几何在极限中完全退化为由 \(\theta(q)\) 决定的二点统计流形。
\end{theorem}
\begin{proof}
记
\(
\varepsilon_m:=(\varphi/2)^m
\)。
由定理 \ref{thm:xi-foldbin-escort-tv-collapse-onebit-gibbs} 的证明，在 \(A_i\) 上有一致估计
\[
\pi_{m,q}(w)
=
\frac{\pi_{m,q}(A_i)}{|A_i|}
\Bigl(1+O_q(\varepsilon_m)\Bigr),
\qquad
\pi_{m,r}(w)
=
\frac{\pi_{m,r}(A_i)}{|A_i|}
\Bigl(1+O_r(\varepsilon_m)\Bigr).
\]
于是
\[
\sum_{w\in A_i}\pi_{m,q}(w)^\alpha\pi_{m,r}(w)^{1-\alpha}
=
\pi_{m,q}(A_i)^\alpha\pi_{m,r}(A_i)^{1-\alpha}
\Bigl(1+O_{q,r,\alpha}(\varepsilon_m)\Bigr),
\]
因为块内的 \(|A_i|\) 因子恰好抵消。

另一方面，由定理 \ref{thm:xi-foldbin-escort-tv-collapse-onebit-gibbs}，
\[
\bigl|\pi_{m,q}(A_1)-\theta(q)\bigr|
\le
\|\pi_{m,q}-\bar\pi_{m,q}\|_{\mathrm{TV}}
=
O_q(\varepsilon_m),
\]
且同理
\(
\pi_{m,q}(A_0)=1-\theta(q)+O_q(\varepsilon_m)
\)；
对 \(r\) 亦然。
因此把 \(i=0,1\) 两块相加，并令 \(m\to\infty\)，即可得到
\[
\mathcal D_\alpha(q,r)
=
\frac{1}{\alpha-1}
\log\!\Bigl(
(1-\theta(q))^\alpha(1-\theta(r))^{1-\alpha}
+
\theta(q)^\alpha\theta(r)^{1-\alpha}
\Bigr).
\]
再代入
\(
\theta(q)=1/(1+\varphi^{q+1})
\)
与
\(
1-\theta(q)=\varphi^{q+1}/(1+\varphi^{q+1})
\)
即可化简为第二种写法。
最后令 \(\alpha\to 1\) 即得 KL 闭式。证毕。
\end{proof}

\begin{theorem}[escort 温度流形的 Fisher 度量是全局平直的]\label{thm:xi-foldbin-escort-logistic-fisher-geometry}
定义 escort 家族的 Fisher 信息
\[
\mathcal I_m(q)
:=
\sum_{w\in X_m}\pi_{m,q}(w)\Bigl(\partial_q\log \pi_{m,q}(w)\Bigr)^2.
\]
则对每个固定 \(q\ge 0\)，有极限
\[
\mathcal I_m(q)\longrightarrow \mathcal I_\infty(q)
:=
(\log\varphi)^2\frac{\varphi^{q+1}}{(1+\varphi^{q+1})^2}.
\]
更强地，极限度量
\[
ds^2=\mathcal I_\infty(q)\,\dd q^2
\]
在变量
\[
u(q):=2\arctan\!\bigl(\varphi^{(q+1)/2}\bigr)
\]
下化为 Euclidean 形式
\[
ds^2=du^2.
\]
因此对任意 \(q_1,q_2\ge 0\)，极限 Fisher 距离都可显式写为
\[
d_F(q_1,q_2)
=
\bigl|u(q_2)-u(q_1)\bigr|
=
2\left|
\arctan\!\bigl(\varphi^{(q_2+1)/2}\bigr)
-
\arctan\!\bigl(\varphi^{(q_1+1)/2}\bigr)
\right|.
\]
特别地，半轴 \(q\in[0,\infty)\) 的总 Fisher 长度为
\[
\pi-2\arctan\sqrt{\varphi}.
\]
\end{theorem}
\begin{proof}
由定义
\[
\partial_q\log\pi_{m,q}(w)
=
\log d_m^{\mathrm{bin}}(w)-\partial_q\log S_q^{\mathrm{bin}}(m),
\]
故
\[
\mathcal I_m(q)
=
\operatorname{Var}_{\pi_{m,q}}\!\bigl(\log d_m^{\mathrm{bin}}(W)\bigr),
\qquad
W\sim \pi_{m,q}.
\]
另一方面，由定理 \ref{thm:fold-bin-two-state-asymptotic} 的一致对数展开，
\[
\log d_m^{\mathrm{bin}}(w)
=
m\log\!\Bigl(\frac{2}{\varphi}\Bigr)-w_m\log\varphi
+O\!\Bigl(\Bigl(\frac{\varphi}{2}\Bigr)^m\Bigr)
\]
在 \(w\in X_m\) 上一致成立。
再由定理 \ref{thm:xi-foldbin-escort-tv-collapse-onebit-gibbs}，
\[
\pi_{m,q}(w_m=1)
=
\theta(q)+O_q\!\Bigl(\Bigl(\frac{\varphi}{2}\Bigr)^m\Bigr).
\]
因此
\[
\mathcal I_m(q)
=
(\log\varphi)^2\operatorname{Var}_{\pi_{m,q}}(w_m)
+O_q\!\Bigl(\Bigl(\frac{\varphi}{2}\Bigr)^m\Bigr)
\]
并进一步得到
\[
\mathcal I_m(q)
=
(\log\varphi)^2\theta(q)(1-\theta(q))
+O_q\!\Bigl(\Bigl(\frac{\varphi}{2}\Bigr)^m\Bigr),
\]
从而
\[
\mathcal I_\infty(q)
=
(\log\varphi)^2\frac{\varphi^{q+1}}{(1+\varphi^{q+1})^2}.
\]

再计算
\[
u'(q)
=
2\cdot \frac{1}{1+\varphi^{q+1}}
\cdot \frac{\log\varphi}{2}\varphi^{(q+1)/2}
=
(\log\varphi)\frac{\varphi^{(q+1)/2}}{1+\varphi^{q+1}},
\]
故
\[
\bigl(u'(q)\bigr)^2
=
(\log\varphi)^2\frac{\varphi^{q+1}}{(1+\varphi^{q+1})^2}
=
\mathcal I_\infty(q).
\]
于是
\(
ds^2=\mathcal I_\infty(q)\,\dd q^2=du^2
\)，距离公式立得。
最后令 \(q_1=0\)、\(q_2\to\infty\)，便得到总长度
\(
\pi-2\arctan\sqrt{\varphi}
\)。
证毕。
\end{proof}


\begin{conclusion}[bin-fold 纤维多重度的线性时间可计算性：Fibonacci-digit DP 给出可复算审计核]\label{con:xi-terminal-foldbin-digit-dp-linear-time}
给定 \(m\ge 1\) 与 \(w\in X_m\)，令 \(K=K(m)\)、\(V=V_m(w)\)、\(\mathrm{slack}:=2^m-1-V\)，并将 \(\mathrm{slack}\) 写为合法 Zeckendorf 位串 \(Z(\mathrm{slack})=(b_1,\dots,b_K)\)。
则 \(d_m^{\mathrm{bin}}(w)\) 可由状态
\(
\mathrm{DP}(k,\mathrm{prev},\mathrm{tight})
\)
的 digit DP 在 \(O(K-m)=O(m)\) 时间内精确计算；转移中仅需施加有限冲突约束（无相邻 \(11\)）、紧性约束与末位约束 \(w_m=1\Rightarrow c_{m+1}=0\)。
因此 bin-fold 的纤维多重度不仅具有可审计的渐近律，同时具备线性复杂度的精确计算内核，为大窗口的完全确定性复算提供算法闭合。
\end{conclusion}
\begin{proof}[证明要点]
见命题 \ref{prop:fold-bin-digit-dp}。
\end{proof}

\begin{proposition}[bin-fold 纤维计数的壁穿越表示：尾部谱 \texorpdfstring{$\Rightarrow$}{=>} 椭圆平均二值跳跃]\label{prop:xi-foldbin-fiber-wallcrossing-representation}
\leanverified{paper\_xi\_foldbin\_fiber\_wallcrossing\_representation}
固定 \(m\ge 1\)。
对每个末位 \(b\in\{0,1\}\)，存在一个有限集合 \(\mathcal{S}_{m,b}\subset\NN_{\ge 1}\)（称为该末位约束下的\emph{尾部权重谱}）使得对任意 \(w\in X_m\) 满足 \(w_m=b\)，令
\[
\mathrm{slack}(w):=2^m-1-V_m(w)\in\NN_{\ge 0},
\]
则 bin-fold 纤维基数满足
\[
d_m^{\mathrm{bin}}(w)
\;=\;
1+\#\{s\in\mathcal{S}_{m,b}:\ s\le \mathrm{slack}(w)\}.
\]
并且存在有理测试函数 \(f_{m,b}\in\CC(w)\)，使对任意 \(w\in X_m\) 满足 \(w_m=b\) 都有
\[
\boxed{\;
d_m^{\mathrm{bin}}(w)
\;=\;
1+\mathcal{L}_{e^{2(\mathrm{slack}(w)+1/2)}}\!\left(f_{m,b}\right)\; },
\]
其中 \(\mathcal{L}_L\) 为定理 \ref{thm:xi-singular-ring-ellipse-wallcrossing-rational} 的奇异环椭圆平均。
\end{proposition}
\begin{proof}
尾部谱的存在性与第一式来自命题 \ref{prop:fold-bin-digit-dp}：digit DP 计数可视为在固定有限约束下枚举“尾部权重和”不超过 \(\mathrm{slack}(w)\) 的可行解个数，并与零尾部（恒存在）相加得到纤维基数。
\par
令 \(\mathcal{A}_{m,b}:=\{e^s:\ s\in\mathcal{S}_{m,b}\}\subset(1,\infty)\)，并对每个 \(a=e^s\in\mathcal{A}_{m,b}\) 赋权 \(c_a:=1\)。
由定理 \ref{thm:xi-singular-ring-ellipse-wallcrossing-rational} 的留数壁穿越构造（命题 \ref{prop:xi-singular-ring-ellipse-wallcrossing-atomic-measure} 所刻画的径向原子测度口径），存在有理函数 \(f_{m,b}\in\CC(w)\) 使得对一切 \(r>1\) 且 \(r\notin\mathcal{A}_{m,b}\) 有
\[
\mathcal{L}_{r^2}(f_{m,b})=\#\{a\in\mathcal{A}_{m,b}:\ a<r\}.
\]
取 \(r=e^{\mathrm{slack}(w)+1/2}\)（从而 \(r\notin \mathcal{A}_{m,b}\)），则右端等于 \(\#\{s\in\mathcal{S}_{m,b}: s\le \mathrm{slack}(w)\}\)，代回即得所示公式。证毕。
\end{proof}

\begin{corollary}[离散差分恢复尾部权重谱]\label{cor:xi-foldbin-tail-spectrum-discrete-difference}
\leanverified{paper\_xi\_foldbin\_tail\_spectrum\_discrete\_difference}
在命题 \ref{prop:xi-foldbin-fiber-wallcrossing-representation} 的设定下，固定 \(m\) 与 \(b\in\{0,1\}\)，定义计数函数
\[
d_{m,b}(R):=\#\{s\in\mathcal{S}_{m,b}:\ s\le R\}\qquad(R\in\NN_{\ge 0}),
\]
并约定 \(d_{m,b}(-1)=0\)。
则对所有 \(R\ge 0\) 有
\[
d_{m,b}(R)-d_{m,b}(R-1)=\mathbf{1}_{\{R\in\mathcal{S}_{m,b}\}},
\]
从而 \(\mathcal{S}_{m,b}\) 可由单个计数函数 \(d_{m,b}\) 完全重构。
\end{corollary}
\begin{proof}
由定义，差分恰为计数在 \(R\) 处新增的点数，即 \(\#\{s\in\mathcal{S}_{m,b}:\ s=R\}\in\{0,1\}\)。证毕。
\end{proof}

\begin{theorem}[固定末位层的 bin-fold 纤维多重度是精确的单跳阶梯]\label{thm:xi-foldbin-lastbit-tail-spectrum-single-jump-staircase}
\leanverified{paper\_xi\_foldbin\_lastbit\_tail\_spectrum\_single\_jump\_staircase}
固定 \(m\ge 1\) 与 \(i\in\{0,1\}\)，记 \(K:=K(m)\)，并定义尾部谱
\[
\mathcal T_{m,i}
:=
\left\{
\sum_{k=m+1}^{K}c_kF_{k+1}\;:\;
c_k\in\{0,1\},\ 
c_kc_{k+1}=0,\ 
i\,c_{m+1}=0
\right\}.
\]
对任意满足 \(w_m=i\) 的稳定字 \(w\in X_m\)，记
\[
V:=V_m(w).
\]
则 bin-fold 纤维多重度满足精确公式
\[
\boxed{\;
d_m^{\mathrm{bin}}(w)
=
\#\Bigl\{t\in\mathcal T_{m,i}:\ t\le 2^m-1-V\Bigr\}.\;}
\]
因此若定义
\[
D_{m,i}(V)
:=
\#\Bigl\{t\in\mathcal T_{m,i}:\ t\le 2^m-1-V\Bigr\},
\]
则 \(V\mapsto D_{m,i}(V)\) 在其定义域上是严格非增的单跳阶梯函数，并且对一切相容的 \(V\) 都有
\[
\boxed{\;
D_{m,i}(V)-D_{m,i}(V+1)
=
\mathbf 1_{\{\,2^m-1-V\in\mathcal T_{m,i}\,\}}.\;}
\]
\end{theorem}
\begin{proof}
由命题 \ref{prop:fold-bin-digit-dp}，对任意 \(w\in X_m\)，纤维 \((\mathrm{Fold}^{\mathrm{bin}}_m)^{-1}(w)\) 与满足
\[
c_k\in\{0,1\},\qquad
c_kc_{k+1}=0,\qquad
w_m c_{m+1}=0,
\]
\[
V_m(w)+\sum_{k=m+1}^{K}c_kF_{k+1}\le 2^m-1
\]
的尾部位串一一对应。固定 \(w_m=i\) 后，前两类约束恰好给出 \(\mathcal T_{m,i}\) 的定义，而阈值约束等价于
\[
t:=\sum_{k=m+1}^{K}c_kF_{k+1}\le 2^m-1-V.
\]
故第一条公式立即成立。

还需说明不同尾部位串必对应不同的 \(t\)。这是 Zeckendorf 唯一性的直接推论：因为这些尾部位串本身已经满足无相邻 \(1\) 的合法语法，因而就是 Fibonacci 权重系下的规范表示；若两组不同尾部给出同一 \(t\)，便得到同一整数的两种合法 Zeckendorf 展开，矛盾。

于是 \(D_{m,i}(V)\) 仅仅是有限集合 \(\mathcal T_{m,i}\) 的初始段计数函数。随着 \(V\) 增大，阈值 \(2^m-1-V\) 每次只下降 \(1\)，故计数只可能在该阈值恰好穿过某个 \(t\in\mathcal T_{m,i}\) 时下降一次，且每次恰下降 \(1\)。这就给出单跳公式。证毕。
\end{proof}

\begin{theorem}[bin-fold 纤维多重度的 Zeckendorf 单调阶梯律]\label{thm:xi-foldbin-zeckendorf-monotone-staircase-law}
沿用定理 \ref{thm:xi-foldbin-lastbit-tail-spectrum-single-jump-staircase} 的记号。对每个固定 \(\varepsilon\in\{0,1\}\)，存在非增函数
\[
N_{m,\varepsilon}:\ZZ_{\ge 0}\to \ZZ_{\ge 0}
\]
使得对任意满足 \(w_m=\varepsilon\) 的稳定字 \(w\in X_m\)，都有
\[
d_m^{\mathrm{bin}}(w)=N_{m,\varepsilon}\!\bigl(V_m(w)\bigr).
\]
亦即，在固定末位之后，纤维多重度仅依赖于 Zeckendorf 值 \(V_m(w)\)；并且只要 \(V\le V'\)，便有
\[
\boxed{\;
N_{m,\varepsilon}(V)\ge N_{m,\varepsilon}(V').\;}
\]
特别地，当 \(m=6\) 时有显式三档公式
\[
d_6^{\mathrm{bin}}(w)=
\begin{cases}
2,& w_6=1,\\[4pt]
4,& w_6=0,\ V_6(w)\le 8,\\[4pt]
3,& w_6=0,\ V_6(w)>8.
\end{cases}
\]
\end{theorem}
\begin{proof}
对固定的 \(\varepsilon\in\{0,1\}\)，定义
\[
N_{m,\varepsilon}(V):=D_{m,\varepsilon}(V),
\]
其中 \(D_{m,\varepsilon}\) 为定理
\ref{thm:xi-foldbin-lastbit-tail-spectrum-single-jump-staircase} 中的初始段计数函数。该定理已经证明：当 \(w_m=\varepsilon\) 时，
\[
d_m^{\mathrm{bin}}(w)=D_{m,\varepsilon}(V_m(w)),
\]
并且 \(V\mapsto D_{m,\varepsilon}(V)\) 在定义域上严格非增。于是所述“只依赖于 \((V_m(w),w_m)\)”以及单调阶梯律立刻成立。

当 \(m=6\) 时，推论 \ref{cor:fold-bin-m6-fiber-hist} 给出纤维直方图
\[
2:8,\qquad 3:4,\qquad 4:9,
\]
而表 \ref{tab:fold6_bin_uplift_choice_collapse} 给出精确对应
\[
w_6=1 \Longleftrightarrow d_6^{\mathrm{bin}}(w)=2,
\]
\[
w_6=0,\ V_6(w)\le 8 \Longleftrightarrow d_6^{\mathrm{bin}}(w)=4,
\]
\[
w_6=0,\ V_6(w)>8 \Longleftrightarrow d_6^{\mathrm{bin}}(w)=3.
\]
合并即得三档公式。证毕。
\end{proof}

\begin{theorem}[bin-fold 精确反演的最小辅助信息位数具有闭式公式]\label{thm:xi-foldbin-exact-inversion-minimal-advice-bits}
\leanverified{paper\_xi\_foldbin\_exact\_inversion\_minimal\_advice\_bits}
记
\[
d_{\max}^{\mathrm{bin}}(m):=\max_{w\in X_m}d_m^{\mathrm{bin}}(w).
\]
定义 \(B_{\min}^{\mathrm{bin}}(m)\) 为使 \(\mathrm{Fold}^{\mathrm{bin}}_m\) 可逐纤维精确反演的最小辅助信息位数，即最小的整数 \(B\ge 0\)，使得存在映射
\[
\mathrm{enc}_m:\{0,\dots,2^m-1\}\to\{0,1\}^{B},
\qquad
\mathrm{dec}_m:X_m\times\{0,1\}^{B}\to\{0,\dots,2^m-1\},
\]
满足对每个 \(n\in\{0,\dots,2^m-1\}\) 都有
\[
\mathrm{dec}_m\bigl(\mathrm{Fold}^{\mathrm{bin}}_m(n),\mathrm{enc}_m(n)\bigr)=n,
\]
并且解码永远不越出输入 \(w\) 所对应的纤维。则有精确闭式
\[
\boxed{\;
B_{\min}^{\mathrm{bin}}(m)
=
\left\lceil \log_2 d_{\max}^{\mathrm{bin}}(m)\right\rceil.\;}
\]
\end{theorem}
\begin{proof}
必要性显然。设 \(w\in X_m\) 使
\[
d_m^{\mathrm{bin}}(w)=d_{\max}^{\mathrm{bin}}(m).
\]
若 \(2^B<d_{\max}^{\mathrm{bin}}(m)\)，则长度 \(B\) 的二进制标签不足以在该最大纤维内实现单射编码；于是存在不同的
\[
n,n'\in (\mathrm{Fold}^{\mathrm{bin}}_m)^{-1}(w)
\]
满足 \(\mathrm{enc}_m(n)=\mathrm{enc}_m(n')\)。这样一来，解码器在同一输入 \((w,\mathrm{enc}_m(n))\) 上不可能同时恢复 \(n\) 与 \(n'\)，矛盾。故必须有
\[
2^B\ge d_{\max}^{\mathrm{bin}}(m),
\]
即
\[
B\ge \left\lceil \log_2 d_{\max}^{\mathrm{bin}}(m)\right\rceil.
\]

反过来，若
\[
2^B\ge d_{\max}^{\mathrm{bin}}(m),
\]
则对每个纤维 \(F_w:=(\mathrm{Fold}^{\mathrm{bin}}_m)^{-1}(w)\) 都可任选一个单射
\[
\iota_w:F_w\hookrightarrow \{0,1\}^{B}.
\]
定义
\[
\mathrm{enc}_m(n):=\iota_{\mathrm{Fold}^{\mathrm{bin}}_m(n)}(n).
\]
再定义 \(\mathrm{dec}_m(w,b)\) 为：若 \(b\in \iota_w(F_w)\)，则输出 \(\iota_w^{-1}(b)\)；否则在 \(F_w\) 中固定任选一个默认元素输出。这样对每个真实输入 \(n\) 都有
\[
\mathrm{dec}_m\bigl(\mathrm{Fold}^{\mathrm{bin}}_m(n),\mathrm{enc}_m(n)\bigr)=n,
\]
且解码始终停留在纤维 \(F_w\) 内。故该下界可达，结论成立。证毕。
\end{proof}

\begin{theorem}[bin-fold 精确反演的辅助位率主项恰为 \texorpdfstring{$\log_2(2/\varphi)$}{log2(2/phi)}]\label{thm:xi-foldbin-exact-inversion-advice-rate-log-two-over-phi}
记
\[
L(m):=K(m)-m.
\]
则 bin-fold 精确反演的最小辅助信息位数满足双边界
\[
\boxed{\;
m-\log_2F_{m+2}
\;\le\;
B_{\min}^{\mathrm{bin}}(m)
\;\le\;
\left\lceil \log_2F_{L(m)+2}\right\rceil.\;}
\]
从而
\[
\boxed{\;
B_{\min}^{\mathrm{bin}}(m)
=
\log_2\!\Bigl(\frac{2}{\varphi}\Bigr)m+O(1).\;}
\]
\end{theorem}
\begin{proof}
下界来自平均纤维大小。由于
\[
|X_m|=F_{m+2},
\qquad
\sum_{w\in X_m}d_m^{\mathrm{bin}}(w)=2^m,
\]
必有
\[
d_{\max}^{\mathrm{bin}}(m)\ge \frac{2^m}{F_{m+2}}.
\]
再由定理 \ref{thm:xi-foldbin-exact-inversion-minimal-advice-bits}，
\[
B_{\min}^{\mathrm{bin}}(m)
=
\left\lceil \log_2 d_{\max}^{\mathrm{bin}}(m)\right\rceil
\ge
\log_2\!\Bigl(\frac{2^m}{F_{m+2}}\Bigr)
=
m-\log_2F_{m+2}.
\]

上界则由定理 \ref{thm:xi-foldbin-universal-fibonacci-upper-bound-lastbit-improvement} 给出：
\[
d_{\max}^{\mathrm{bin}}(m)\le F_{L(m)+2}.
\]
再次应用定理 \ref{thm:xi-foldbin-exact-inversion-minimal-advice-bits} 即得
\[
B_{\min}^{\mathrm{bin}}(m)\le \left\lceil \log_2F_{L(m)+2}\right\rceil.
\]

最后使用 H.55 的定义不等式
\[
F_{K(m)+1}\le 2^m-1<F_{K(m)+2}
\]
与 Binet 公式，得到
\[
K(m)=\frac{\log 2}{\log\varphi}\,m+O(1),
\qquad
L(m)=\left(\frac{\log 2}{\log\varphi}-1\right)m+O(1).
\]
再由
\[
\log F_n=n\log\varphi+O(1)
\]
可得
\[
\log_2F_{L(m)+2}
=
\frac{\log\varphi}{\log 2}\,L(m)+O(1)
=
\Bigl(1-\log_2\varphi\Bigr)m+O(1)
=
\log_2\!\Bigl(\frac{2}{\varphi}\Bigr)m+O(1).
\]
同理，
\[
m-\log_2F_{m+2}
=
m-\frac{(m+2)\log\varphi+O(1)}{\log 2}
=
\log_2\!\Bigl(\frac{2}{\varphi}\Bigr)m+O(1).
\]
双边夹逼即得所示主项。证毕。
\end{proof}

\begin{theorem}[bin-fold 精确反演阈值不仅可达，而且可在多项式时间内可达]\label{thm:xi-foldbin-rank-unrank-optimal-advice-polytime}
\leanverified{paper\_xi\_foldbin\_rank\_unrank\_optimal\_advice\_polytime}
存在确定性算法族
\[
\operatorname{Rank}_m,\qquad \operatorname{Unrank}_m
\]
满足：
\begin{enumerate}
  \item 对任意 \(w\in X_m\) 与 \(0\le j<d_m^{\mathrm{bin}}(w)\)，\(\operatorname{Unrank}_m(w,j)\) 输出纤维 \((\mathrm{Fold}^{\mathrm{bin}}_m)^{-1}(w)\) 中按自然大小排序的第 \(j\) 个整数；
  \item 对任意 \(n\in\{0,\dots,2^m-1\}\)，\(\operatorname{Rank}_m(n)\) 输出 \(n\) 在其所属纤维中的秩；
  \item 在单位成本 RAM 口径下，两者都可在 \(O(m)\) 次状态更新内完成；在比特口径下，两者都可在 \(O(m^2)\) 位运算内完成。
\end{enumerate}
因此，定理 \ref{thm:xi-foldbin-exact-inversion-minimal-advice-bits} 所给出的最小辅助信息阈值并非纯存在性界，而是一个可在多项式时间内达到的尖锐阈值：存在最优编码/解码器族，仅使用
\[
\boxed{\;
\left\lceil \log_2 d_{\max}^{\mathrm{bin}}(m)\right\rceil\;}
\]
位辅助信息即可逐纤维精确反演。
\end{theorem}
\begin{proof}
由命题 \ref{prop:fold-bin-digit-dp}，在固定 \(w\in X_m\) 后，纤维元素与尾部位串
\[
(c_{m+1},\dots,c_K)
\]
一一对应；其合法性完全由相邻约束、边界约束与阈值约束控制，而每个固定前缀的可扩展完成数都由状态
\[
\mathrm{DP}(k,\mathrm{prev},\mathrm{tight})
\]
精确给出。

先构造 \(\operatorname{Unrank}_m\)。给定 \((w,j)\)，令 \(K=K(m)\)、\(V=V_m(w)\) 与 \(\mathrm{slack}=2^m-1-V\)。随后按位从高到低扫描 \(k=K,K-1,\dots,m+1\)。在每一步，先尝试 \(c_k=0\)，并用相应 DP 状态求出在该前缀下可完成的尾部总数 \(N_0\)。若 \(j<N_0\)，就取 \(c_k=0\)；否则令
\[
j\leftarrow j-N_0
\]
并在不违反局部约束的前提下取 \(c_k=1\)。如此递推下去，便唯一恢复第 \(j\) 个合法尾部。再把它与前缀 \(w\) 拼回，即得唯一的整数原像 \(\operatorname{Unrank}_m(w,j)\)。

再构造 \(\operatorname{Rank}_m\)。给定 \(n\in\{0,\dots,2^m-1\}\)，先计算
\[
w:=\mathrm{Fold}^{\mathrm{bin}}_m(n)
\]
及其 Zeckendorf 尾部位串。随后从高位到低位扫描：每当实际数码在某一步取 \(1\) 时，就把“若该位改取 \(0\) 且其余低位任意合法完成”所对应的完成数加到秩上；这一完成数仍由同一 DP 状态值给出。最终所得即为 \(n\) 在纤维 \((\mathrm{Fold}^{\mathrm{bin}}_m)^{-1}(w)\) 内按自然大小排序的秩。

算法正确性来自两点：第一，H.56 的尾部参数化给出纤维与合法尾串的严格双射；第二，digit DP 精确计算了每个前缀状态的完成数，因此 lexicographic rank/unrank 的标准递推完全成立。

复杂度方面，尾长
\[
L(m)=K(m)-m=O(m).
\]
DP 状态数为 \(O(L(m))=O(m)\)，每个状态只有常数个转移。又由定理 \ref{thm:xi-foldbin-universal-fibonacci-upper-bound-lastbit-improvement}，
\[
d_m^{\mathrm{bin}}(w)\le F_{L(m)+2},
\]
故所有 DP 计数与秩参数的比特长度都是 \(O(m)\)。因此在单位成本 RAM 口径下，填表与扫描总共需要 \(O(m)\) 次状态更新；在比特口径下，每一步比较、加法或减法的成本为 \(O(m)\)，总成本为 \(O(m^2)\)。

最后，令
\[
B_m:=\left\lceil \log_2 d_{\max}^{\mathrm{bin}}(m)\right\rceil.
\]
把 \(\operatorname{Rank}_m(n)\) 的值编码成长度 \(B_m\) 的二进制串，作为最优编码器；解码器则在输入 \((w,b)\) 时先读出 \(j=\mathrm{val}(b)\)，若 \(j<d_m^{\mathrm{bin}}(w)\) 则返回 \(\operatorname{Unrank}_m(w,j)\)，否则在该纤维内返回固定默认元。于是对所有真实编码输入都能精确恢复，从而该阈值在多项式时间内可达。证毕。
\end{proof}

\endinput


\begin{conjecture}[dyadic--Fibonacci 归一化尾部谱的奇异连续极限]\label{conj:xi-foldbin-tail-spectrum-singular-limit}
固定末位 \(b\in\{0,1\}\)，并以命题 \ref{prop:xi-foldbin-fiber-wallcrossing-representation} 的尾部权重谱 \(\mathcal{S}_{m,b}\) 为基础，定义归一化点集
\[
\widetilde{\mathcal{S}}_{m,b}:=\left\{\frac{s}{2^m}:\ s\in\mathcal{S}_{m,b}\right\}\subset[0,1],
\]
以及对应的经验分布函数
\[
F_{m,b}(t):=\frac{1}{|\mathcal{S}_{m,b}|}\#\left\{s\in\mathcal{S}_{m,b}:\ \frac{s}{2^m}\le t\right\}\qquad(t\in[0,1]).
\]
猜想存在非平凡极限分布函数 \(F_b\) 与稠密子集 \(T\subset[0,1]\)，使得对每个 \(t\in T\) 有 \(F_{m,b}(t)\to F_b(t)\)，并且 \(F_b\) 严格递增但奇异连续（Lebesgue a.e. 导数为 \(0\)）。
此外，\(F_b\) 的多重分形标度由 dyadic--Fibonacci 比率 \(\log_\varphi 2\) 控制，并与本文中由压力函数/反演核诱导的谱对象保持一致性接口。
\end{conjecture}

\paragraph{\texorpdfstring{$q$}{q}-复制软核传递矩阵的整系数结构：多尺度转级数、补集自对偶反演与 Smith 形刚性}
沿用命题 \ref{prop:pom-multiplicity-composition-replica-softcore-transfer} 的状态空间 \(U_q=\{0,1\}^q\) 与软核传递矩阵 \(T^{(q)}\)（\(T^{(q)}_{u,v}=2^{-\mathbf{1}_{\{u\wedge v\neq 0\}}}\)），并记增长常数
\(
\lambda_q:=2\,\rho(T^{(q)})
\)。
将 \(U_q\) 与子集幂集 \(\mathcal P([q])\) 通过指示向量识别，并令 \(|u|\) 表示 Hamming 权。
再令 \(D^{(q)}:=K^{\otimes q}\)（定理 \ref{thm:pom-replica-softcore-tensor-eigenbasis}），则
\(
T^{(q)}=\frac12(\mathbf{1}\mathbf{1}^{\mathsf T}+D^{(q)})
\)，从而
\[
B^{(q)}:=2T^{(q)}=\mathbf{1}\mathbf{1}^{\mathsf T}+D^{(q)}\in \mathrm{M}_{2^q}(\ZZ).
\]

\begin{conclusion}[全谱的纤维分裂刚性：固定 Kronecker 体谱簇与 \texorpdfstring{$(q+1)$}{(q+1)} 维例外块]\label{con:xi-terminal-replica-softcore-fiber-splitting-full-spectrum}
固定整数 \(q\ge 2\)，沿用 \(U_q=\{0,1\}^q\)、\(D^{(q)}=K^{\otimes q}\) 与
\(
T^{(q)}=\frac12(\mathbf{1}\mathbf{1}^{\mathsf T}+D^{(q)})
\)
的记号（\S\ref{sec:pom}）。
定义固定谱簇数列
\[
\nu_k^{(q)}:=\frac12(-1)^k\varphi^{q-2k}\qquad(0\le k\le q),
\]
则 \(T^{(q)}\) 的谱由两部分完全刻画：
\begin{enumerate}
  \item \textbf{（固定谱簇）}
  对每个 \(0\le k\le q\)，数 \(\nu_k^{(q)}\) 作为 \(T^{(q)}\) 的特征值出现，其重数恰为 \(\binom{q}{k}-1\)；
  并约定 \(k=0,q\) 时该重数为 \(0\)（因 \(\binom{q}{0}=\binom{q}{q}=1\)）。
  \item \textbf{（例外谱）}
  除上述固定谱簇外，\(T^{(q)}\) 还恰有 \(q+1\) 个互异的实特征值；
  它们与 \(S_q\)-对称降维矩阵 \(A^{(q)}\) 的谱完全一致（定义 \eqref{eq:pom-multiplicity-composition-Aq-def}，命题 \ref{prop:pom-multiplicity-composition-sq-symmetry-reduction}），并且全部为单根。
\end{enumerate}
\end{conclusion}
\begin{proof}[证明要点]
由 \(D^{(q)}=K^{\otimes q}\) 的张量本征分解，\(D^{(q)}\) 的特征值为 \((-1)^k\varphi^{q-2k}\)（重数 \(\binom{q}{k}\)）。
并且 \(\mathbf{1}\mathbf{1}^{\mathsf T}\) 为秩一算子：若 \(x\perp\mathbf{1}\)，则 \((\mathbf{1}\mathbf{1}^{\mathsf T})x=0\)，从而 \(T^{(q)}x=\frac12D^{(q)}x\)。
\par
另一方面，结论 \ref{con:xi-terminal-replica-softcore-exceptional-secular-stieltjes} 的权重计算给出 \(P_k\mathbf{1}\neq 0\)（\(0\le k\le q\)），其中 \(P_k\) 为 \(D^{(q)}\) 对应第 \(k\) 个谱层的正交投影。
故每个谱层与 \(\mathbf{1}^\perp\) 的交维数为 \(\binom{q}{k}-1\)，得到固定谱簇的重数。
剩余维数为
\(
2^q-\sum_{k=0}^{q}(\binom{q}{k}-1)=q+1
\)，并由 \(S_q\)-对称子空间承载；
其限制矩阵即为 \(A^{(q)}\)（命题 \ref{prop:pom-multiplicity-composition-sq-symmetry-reduction}），从而例外谱为 \(\mathrm{spec}(A^{(q)})\)。
单根性由同一结论的“对角 \(+\) 秩一”世俗刻画与权重全非零推出。证毕。
\end{proof}

\begin{conclusion}[例外谱的世俗方程：双代数权重的 Stieltjes 型刻画]\label{con:xi-terminal-replica-softcore-exceptional-secular-stieltjes}
沿用结论 \ref{con:xi-terminal-replica-softcore-fiber-splitting-full-spectrum} 的记号。
令 \(K\) 的单位正交本征向量为 \((e_+,e_-)\)，对应本征值 \((\varphi,-\varphi^{-1})\)。
将 \(u:=(1,1)^{\mathsf T}\) 写为
\[
u=\alpha e_+ + \beta e_-,
\qquad
\alpha^2=1+\frac{2}{\sqrt5},\quad \beta^2=1-\frac{2}{\sqrt5}.
\]
对 \(0\le k\le q\) 定义权重
\[
w_k^{(q)}:=\binom{q}{k}\alpha^{2(q-k)}\beta^{2k}
=\binom{q}{k}\Bigl(1+\frac{2}{\sqrt5}\Bigr)^{q-k}\Bigl(1-\frac{2}{\sqrt5}\Bigr)^{k},
\]
并仍记固定谱簇对角值 \(\nu_k^{(q)}=\frac12(-1)^k\varphi^{q-2k}\)。
则例外谱 \(\{\lambda_i^{(q)}\}_{i=0}^{q}=\mathrm{spec}(A^{(q)})\) 恰为下列世俗方程的全部（且全为单根）实解：
\[
\boxed{\ \frac12\sum_{k=0}^{q}\frac{w_k^{(q)}}{\lambda-\nu_k^{(q)}}=1\ }.
\]
等价地，令
\[
P_q(\lambda)
:=\prod_{k=0}^{q}(\lambda-\nu_k^{(q)})
-\frac12\sum_{k=0}^{q}w_k^{(q)}\prod_{\substack{0\le j\le q\\ j\ne k}}(\lambda-\nu_j^{(q)}),
\]
则 \(\{\lambda_i^{(q)}\}\) 为 \(P_q(\lambda)=0\) 的根集；
并且 \(P_q(\lambda)\in\QQ[\lambda]\)，更强地 \(2^{q+1}P_q(\lambda)\in\ZZ[\lambda]\)。
\end{conclusion}
\begin{proof}
由 \(\mathbf{1}=u^{\otimes q}\) 与 \(D^{(q)}=K^{\otimes q}\) 的张量本征分解，
\[
\mathbf{1}
=\sum_{k=0}^{q}\alpha^{q-k}\beta^k\sum_{|S|=k}e_-^{\otimes S}\otimes e_+^{\otimes S^{\mathsf c}},
\]
其中对固定 \(k\) 的内层向量两两正交且位于 \(D^{(q)}\) 的第 \(k\) 个谱层。
故 \(P_k\mathbf{1}\neq 0\) 且
\[
\|P_k\mathbf{1}\|^2
=\binom{q}{k}\alpha^{2(q-k)}\beta^{2k}
=w_k^{(q)}.
\]
令 \(f_k:=P_k\mathbf{1}/\|P_k\mathbf{1}\|\) 并记 \(c_k:=\|P_k\mathbf{1}\|=\sqrt{w_k^{(q)}}\)。
在正交基 \(\{f_k\}_{k=0}^{q}\) 下，限制到 \(\mathrm{span}\{P_k\mathbf{1}\}_{k=0}^{q}\) 的算子为
\[
T^{(q)}\big|_{\mathrm{sym}}
=\mathrm{diag}(\nu_0^{(q)},\dots,\nu_q^{(q)})+\frac12\,cc^{\mathsf T}.
\]
对“对角 \(+\) 秩一”矩阵应用行列式引理得到
\[
0=\det\bigl(\lambda I-\mathrm{diag}(\nu_k^{(q)})-\tfrac12cc^{\mathsf T}\bigr)
=\det(\lambda I-\mathrm{diag}(\nu_k^{(q)}))
\Bigl(1-\frac12\sum_{k=0}^{q}\frac{c_k^2}{\lambda-\nu_k^{(q)}}\Bigr),
\]
从而得到所示世俗方程，并且由于 \(c_k^2=w_k^{(q)}>0\)，各根均为单根且与 \(\{\nu_k^{(q)}\}\) 严格交错。
\par
多项式表示由两边同乘 \(\prod_{k}(\lambda-\nu_k^{(q)})\) 得到。
至于有理/整性：例外谱亦等于 \(A^{(q)}\) 的谱（命题 \ref{prop:pom-multiplicity-composition-sq-symmetry-reduction}），而 \(A^{(q)}\) 在基 \(\{e_b\}_{b=0}^{q}\) 下为半整数矩阵（\eqref{eq:pom-multiplicity-composition-Aq-def}），故其特征多项式属于 \(\QQ[\lambda]\)，并且乘以 \(2^{q+1}\) 后系数皆为整数。证毕。
\end{proof}

\begin{conclusion}[例外块逆矩阵的 Pascal--反转分解与整数核闭式]\label{con:xi-terminal-replica-softcore-Aq-explicit-inverse-integrality}
固定整数 \(q\ge 1\)，并沿用 \(S_q\)-对称降维矩阵 \(A^{(q)}\) 的定义 \eqref{eq:pom-multiplicity-composition-Aq-def}。
记 Pascal 矩阵 \(P_q\in M_{q+1}(\ZZ)\) 与反转置换矩阵 \(R_q\in M_{q+1}(\ZZ)\) 为
\[
(P_q)_{i,j}:=\binom{i}{j}\mathbf 1_{\{i\ge j\}},
\qquad
(R_q)_{i,j}:=\mathbf 1_{\{i+j=q\}},
\]
并令 \(E_{00}:=e_0e_0^{\mathsf T}\)。
则
\[
\boxed{\;
(A^{(q)})^{-1}=2P_q^{-1}R_q-E_{00}\in M_{q+1}(\ZZ)\; }.
\]
等价地，对任意 \(0\le k,\ell\le q\)，有
\[
\boxed{\;
(A^{(q)})^{-1}_{k,\ell}
=2(-1)^{k+\ell-q}\binom{k}{q-\ell}\mathbf 1_{\{k+\ell\ge q\}}
-\delta_{k,0}\delta_{\ell,0}\; } .
\]
\end{conclusion}
\begin{proof}
令 \(b_\ell:=\binom{q}{\ell}\)（\(0\le \ell\le q\)）与 \(\mathbf 1:=(1,\dots,1)^{\mathsf T}\)。
并定义
\[
C^{(q)}_{k,\ell}:=\binom{q-k}{\ell},
\qquad
0\le k,\ell\le q.
\]
由 \(R_q\) 的反转作用直接得到 \(C^{(q)}=R_qP_q\)，从而
\[
A^{(q)}=\frac12\bigl(C^{(q)}+\mathbf 1\,b^{\mathsf T}\bigr)
=\frac12\bigl(R_qP_q+\mathbf 1\,b^{\mathsf T}\bigr).
\]
因为 \(R_q\) 为置换矩阵且 \(P_q\) 为单位下三角矩阵，故
\[
(R_qP_q)^{-1}=P_q^{-1}R_q.
\]
对秩一扰动应用 Sherman--Morrison 公式：
\[
\bigl(R_qP_q+\mathbf 1b^{\mathsf T}\bigr)^{-1}
=(R_qP_q)^{-1}
-\frac{(R_qP_q)^{-1}\mathbf 1\,b^{\mathsf T}(R_qP_q)^{-1}}
{1+b^{\mathsf T}(R_qP_q)^{-1}\mathbf 1}.
\]
由反二项恒等式 \(P_q^{-1}\mathbf 1=e_0\) 可得
\[
(R_qP_q)^{-1}\mathbf 1=P_q^{-1}R_q\mathbf 1=e_0.
\]
又 \(b^{\mathsf T}\) 是 \(P_q\) 的第 \(q\) 行，故
\[
b^{\mathsf T}(R_qP_q)^{-1}
=b^{\mathsf T}P_q^{-1}R_q
=e_q^{\mathsf T}R_q
=e_0^{\mathsf T}.
\]
于是分母为 \(1+e_0^{\mathsf T}e_0=2\)，并且
\[
\bigl(R_qP_q+\mathbf 1b^{\mathsf T}\bigr)^{-1}
=P_q^{-1}R_q-\frac12E_{00}.
\]
再由 \(A^{(q)}=\frac12(R_qP_q+\mathbf 1b^{\mathsf T})\) 得
\[
(A^{(q)})^{-1}=2P_q^{-1}R_q-E_{00}.
\]
最后用
\[
(P_q^{-1})_{i,j}=(-1)^{i-j}\binom{i}{j}\mathbf 1_{\{i\ge j\}}
\]
与 \(R_q\) 的定义即可得到逐元闭式。证毕。
\end{proof}

\begin{conclusion}[例外块逆矩阵的 Smith 正规形与模二秩塌缩]\label{con:xi-terminal-replica-softcore-Aq-inverse-smith-mod2-collapse}
对任意整数 \(q\ge 1\)，整数矩阵 \((A^{(q)})^{-1}\in M_{q+1}(\ZZ)\) 的 Smith 正规形满足
\[
\boxed{\;
\mathrm{SNF}\bigl((A^{(q)})^{-1}\bigr)
=\mathrm{diag}\bigl(1,\underbrace{2,\dots,2}_{q\text{ 次}}\bigr)\; }.
\]
因而
\[
\mathrm{coker}\!\bigl((A^{(q)})^{-1}:\ZZ^{q+1}\to\ZZ^{q+1}\bigr)
\cong (\ZZ/2\ZZ)^q.
\]
并且
\[
\boxed{\;
(A^{(q)})^{-1}\equiv E_{00}\pmod 2,\qquad
\mathrm{rank}_{\FF_2}\!\bigl((A^{(q)})^{-1}\bmod 2\bigr)=1\; }.
\]
\end{conclusion}
\begin{proof}
由结论 \ref{con:xi-terminal-replica-softcore-Aq-explicit-inverse-integrality}，
\[
P_q(A^{(q)})^{-1}R_q
=P_q\bigl(2P_q^{-1}R_q-E_{00}\bigr)R_q
=2I_{q+1}-P_qE_{00}R_q
=:M_q.
\]
注意 \(P_qE_{00}=(P_qe_0)e_0^{\mathsf T}=\mathbf 1e_0^{\mathsf T}\)，且 \(e_0^{\mathsf T}R_q=e_q^{\mathsf T}\)，故
\[
M_q=2I_{q+1}-\mathbf 1e_q^{\mathsf T}.
\]
由于 \(P_q,R_q\in\mathrm{GL}_{q+1}(\ZZ)\)，\(M_q\) 与 \((A^{(q)})^{-1}\) 在 \(\ZZ\) 上酉模等价，因而 Smith 正规形相同。
对 \(M_q\) 做如下初等变换：
\begin{enumerate}
  \item 列变换 \(C_q\leftarrow C_q+\sum_{j=0}^{q-1}C_j\)，将最后一列化为全 \(1\) 列；
  \item 对 \(0\le i\le q-1\)，行变换 \(R_i\leftarrow R_i-R_q\)，消去前 \(q\) 行最后一列。
\end{enumerate}
得到对角矩阵 \(\mathrm{diag}(2,\dots,2,1)\)，交换对角元顺序即得
\(\mathrm{diag}(1,2,\dots,2)\)。
于是余核同构与模二秩结论同时成立。证毕。
\end{proof}

\begin{corollary}[例外谱倒数一次幂和的 Fibonacci 闭式]\label{cor:xi-terminal-replica-softcore-exceptional-reciprocal-sum-fibonacci}
对一切整数 \(q\ge 1\)，令 \(\{\lambda_i^{(q)}\}_{i=0}^{q}=\mathrm{spec}(A^{(q)})\)。
则成立严格恒等式
\[
\boxed{\ \sum_{i=0}^{q}\frac{1}{\lambda_i^{(q)}}
:=\operatorname{tr}\bigl((A^{(q)})^{-1}\bigr)
=-1+2(-1)^qF_{q+1}\ } .
\]
\end{corollary}
\begin{proof}
记 \(U_q:=P_q^{-1}R_q\)。由结论 \ref{con:xi-terminal-replica-softcore-Aq-explicit-inverse-integrality}，
\[
\operatorname{tr}\bigl((A^{(q)})^{-1}\bigr)=2\operatorname{tr}(U_q)-1.
\]
并且
\[
(U_q)_{k,k}
=(P_q^{-1})_{k,q-k}
=(-1)^{k-(q-k)}\binom{k}{q-k}\mathbf 1_{\{k\ge q-k\}}
=(-1)^q\binom{k}{q-k}\mathbf 1_{\{k\ge \lceil q/2\rceil\}}.
\]
故
\[
\operatorname{tr}(U_q)
=(-1)^q\sum_{k=\lceil q/2\rceil}^{q}\binom{k}{q-k}
=(-1)^q\sum_{r=0}^{\lfloor q/2\rfloor}\binom{q-r}{r}
=(-1)^qF_{q+1}.
\]
代回即得结论。证毕。
\end{proof}

\begin{corollary}[例外谱倒数二次幂和的 Fibonacci 闭式]\label{cor:xi-terminal-replica-softcore-exceptional-reciprocal-square-sum-fibonacci}
对一切整数 \(q\ge 1\)，令 \(\{\lambda_i^{(q)}\}_{i=0}^{q}=\mathrm{spec}(A^{(q)})\)。
则成立严格恒等式
\[
\boxed{\ \sum_{i=0}^{q}\frac{1}{\bigl(\lambda_i^{(q)}\bigr)^2}
:=\operatorname{tr}\bigl((A^{(q)})^{-2}\bigr)
=4F_{2q+2}+1\ } .
\]
\end{corollary}
\begin{proof}
仍记 \(U_q:=P_q^{-1}R_q\)。由结论 \ref{con:xi-terminal-replica-softcore-Aq-explicit-inverse-integrality}，
\[
(A^{(q)})^{-1}=2U_q-E_{00},
\]
故
\[
\operatorname{tr}\bigl((A^{(q)})^{-2}\bigr)
=4\operatorname{tr}(U_q^2)-2\operatorname{tr}(U_qE_{00})-2\operatorname{tr}(E_{00}U_q)+1.
\]
当 \(q\ge1\) 时 \((U_q)_{0,0}=(P_q^{-1})_{0,q}=0\)，因此
\[
\operatorname{tr}\bigl((A^{(q)})^{-2}\bigr)=4\operatorname{tr}(U_q^2)+1.
\]
由上一结论中的逐元公式，
\[
(U_q)_{k,\ell}=(-1)^{k+\ell-q}\binom{k}{q-\ell}\mathbf 1_{\{k+\ell\ge q\}},
\]
从而
\[
\operatorname{tr}(U_q^2)
=\sum_{k,\ell=0}^{q}(U_q)_{k,\ell}(U_q)_{\ell,k}
=\sum_{\substack{k,\ell\ge0\\k+\ell\ge q}}
\binom{k}{q-\ell}\binom{\ell}{q-k}.
\]
令 \(k'=q-k,\ \ell'=q-\ell\)，得到
\[
\operatorname{tr}(U_q^2)
=\sum_{\substack{k',\ell'\ge0\\k'+\ell'\le q}}
\binom{q-k'}{\ell'}\binom{q-\ell'}{k'}.
\]
该双和与推论 \ref{cor:xi-terminal-replica-softcore-exceptional-power-sum-m2-closed-form}
证明中的组合项 \(T_4(q)\) 完全一致，因此
\(
\operatorname{tr}(U_q^2)=F_{2q+2}
\)。
代回即得结论。证毕。
\end{proof}

\begin{theorem}[例外谱块的二项核逆矩阵闭式]\label{thm:xi-exceptional-binomial-kernel-inverse}
\leanverified{paper\_xi\_exceptional\_binomial\_kernel\_inverse}
固定 \(q\ge 2\)。定义
\[
A^{(q)}_{k,\ell}:=\frac12\!\left(\binom{q}{\ell}+\binom{q-k}{\ell}\right),\qquad
C^{(q)}_{k,\ell}:=\binom{q-k}{\ell},
\qquad 0\le k,\ell\le q,
\]
并记
\[
b:=\bigl(\binom{q}{0},\binom{q}{1},\dots,\binom{q}{q}\bigr)^{\mathsf T},\qquad
\mathbf 1:=(1,\dots,1)^{\mathsf T},\qquad
E_{00}:=e_0e_0^{\mathsf T}.
\]
则
\[
A^{(q)}=\frac12\bigl(\mathbf 1\,b^{\mathsf T}+C^{(q)}\bigr),
\]
且 \(A^{(q)}\) 可逆并满足
\[
\boxed{\ \bigl(A^{(q)}\bigr)^{-1}=2\bigl(C^{(q)}\bigr)^{-1}-E_{00}\ }\in M_{q+1}(\ZZ).
\]
逐元公式可写为
\[
\bigl(A^{(q)}\bigr)^{-1}_{k,\ell}=
\begin{cases}
-1,&(k,\ell)=(0,0),\\[2mm]
2(-1)^{k+\ell-q}\binom{k}{q-\ell},&k+\ell\ge q,\ (k,\ell)\neq(0,0),\\[2mm]
0,&k+\ell<q.
\end{cases}
\]
\end{theorem}
\begin{proof}
由定义有 \(A^{(q)}=\frac12(C^{(q)}+\mathbf 1\,b^{\mathsf T})\)。
又 \(C^{(q)}e_0=\mathbf 1\)，故 \((C^{(q)})^{-1}\mathbf 1=e_0\)。
并且 \(C^{(q)}\) 的第 \(0\) 行正是 \(b^{\mathsf T}\)，即 \(b^{\mathsf T}=e_0^{\mathsf T}C^{(q)}\)，从而
\[
b^{\mathsf T}(C^{(q)})^{-1}=e_0^{\mathsf T},
\qquad
b^{\mathsf T}(C^{(q)})^{-1}\mathbf 1=e_0^{\mathsf T}e_0=1.
\]
对秩一扰动应用 Sherman--Morrison 公式：
\[
\bigl(C^{(q)}+\mathbf 1\,b^{\mathsf T}\bigr)^{-1}
=(C^{(q)})^{-1}
-\frac{(C^{(q)})^{-1}\mathbf 1\ b^{\mathsf T}(C^{(q)})^{-1}}
{1+b^{\mathsf T}(C^{(q)})^{-1}\mathbf 1}
=(C^{(q)})^{-1}-\frac12E_{00}.
\]
两侧乘以 \(2\) 即得
\(
(A^{(q)})^{-1}=2(C^{(q)})^{-1}-E_{00}
\)。
再由结论 \ref{con:xi-terminal-disjointness-bq-sympower-unimodular-inverse} 的逆矩阵闭式与 \(C^{(q)}=B_q^{\mathsf T}\) 得
\[
\bigl(C^{(q)}\bigr)^{-1}_{k,\ell}
=(-1)^{k+\ell-q}\binom{k}{q-\ell}\mathbf 1_{\{k+\ell\ge q\}},
\]
代回即得逐元表达。证毕。
\end{proof}

\begin{proposition}[整格像的首坐标二元判别与余核刚性]\label{prop:xi-exceptional-binomial-kernel-image-cokernel}
\leanverified{paper\_xi\_exceptional\_binomial\_kernel\_image\_cokernel}
在定理 \ref{thm:xi-exceptional-binomial-kernel-inverse} 记号下，令
\[
M^{(q)}:=2A^{(q)}=C^{(q)}+\mathbf 1\,b^{\mathsf T}\in M_{q+1}(\ZZ).
\]
则对任意 \(y\in\ZZ^{q+1}\)，方程 \(M^{(q)}x=y\) 存在整数解 \(x\in\ZZ^{q+1}\) 当且仅当 \(y_0\) 为偶数。等价地，
\[
M^{(q)}(\ZZ^{q+1})
=\{y\in\ZZ^{q+1}:y_0\equiv 0\pmod 2\},
\qquad
\mathrm{coker}_{\ZZ}(M^{(q)})\cong \ZZ/2\ZZ.
\]
并且在 \(\FF_2\) 上
\[
\ker\!\bigl(M^{(q)}\bmod 2\bigr)=\langle e_0\rangle.
\]
\end{proposition}
\begin{proof}
由上一定理，
\[
\bigl(M^{(q)}\bigr)^{-1}
=\bigl(C^{(q)}\bigr)^{-1}-\frac12E_{00}.
\]
故对 \(y\in\ZZ^{q+1}\) 有
\[
x=\bigl(M^{(q)}\bigr)^{-1}y
=\bigl(C^{(q)}\bigr)^{-1}y-\frac{y_0}{2}e_0.
\]
因 \((C^{(q)})^{-1}\in M_{q+1}(\ZZ)\)，第一项始终为整数向量，因此
\(
x\in\ZZ^{q+1}\iff y_0/2\in\ZZ
\)，即 \(y_0\) 为偶数。
这给出整像的精确描述，并立得余核为 \(\ZZ/2\ZZ\)。

对模 \(2\) 像空间，以上整像描述表明其恰为 \(\{z\in\FF_2^{q+1}:z_0=0\}\)，故秩为 \(q\)、核维数为 \(1\)。
又 \(M^{(q)}\) 的第 \(0\) 列全为 \(2\)，故 \(e_0\in\ker(M^{(q)}\bmod 2)\)；
由一维性可知核即 \(\langle e_0\rangle\)。证毕。
\end{proof}

\begin{theorem}[例外谱倒数幂和的 Fibonacci 闭式]\label{thm:xi-exceptional-reciprocal-moment-fibonacci-unified}
设 \(\{\lambda_i^{(q)}\}_{i=0}^{q}=\mathrm{spec}(A^{(q)})\)。
则对任意 \(q\ge 1\) 有
\[
\sum_{i=0}^{q}\frac{1}{\lambda_i^{(q)}}
=-1+2(-1)^qF_{q+1},
\]
\[
\sum_{i=0}^{q}\frac{1}{\bigl(\lambda_i^{(q)}\bigr)^2}
=4F_{2q+2}+1.
\]
\end{theorem}
\begin{proof}
第一式即推论 \ref{cor:xi-terminal-replica-softcore-exceptional-reciprocal-sum-fibonacci}，
第二式即推论 \ref{cor:xi-terminal-replica-softcore-exceptional-reciprocal-square-sum-fibonacci}。
两式可分别写为
\(
\operatorname{tr}((A^{(q)})^{-1})
\)
与
\(
\operatorname{tr}((A^{(q)})^{-2})
\)
的闭式。证毕。
\end{proof}

\begin{corollary}[全矩阵行列式的塌缩与整数逆核]\label{cor:xi-terminal-replica-softcore-Tq-det-and-integer-inverse-kernel}
设 \(q\ge 2\) 并记 \(N:=2^q\)。
则
\[
\boxed{\ \det T^{(q)}=2^{-(N-1)}\ }.
\]
并且 \(\bigl(T^{(q)}\bigr)^{-1}\) 为整数矩阵：对任意 \(u,v\in U_q\)，记 \(\mathbf{0}:=(0,\dots,0)\)、\(\mathbf{1}:=(1,\dots,1)\)，则
\[
\boxed{\ \bigl(T^{(q)}\bigr)^{-1}_{u,v}
=2(-1)^{|u\wedge v|}\,\mathbf{1}_{\{u\vee v=\mathbf{1}\}}
-\delta_{u,\mathbf{0}}\delta_{v,\mathbf{0}}\ } .
\]
\end{corollary}
\begin{proof}
由 \(B^{(q)}:=2T^{(q)}\) 与结论 \ref{con:xi-terminal-disjointness-bq-smith-snf-coker-parity} 的 \(\det(B^{(q)})=2\)，
\(
\det(T^{(q)})=2^{-N}\det(B^{(q)})=2^{-(N-1)}
\)。
另一方面，结论 \ref{con:xi-terminal-disjointness-bq-complement-selfdual-inverse-rankone} 给出
\(
(B^{(q)})^{-1}=(D^{(q)})^{-1}-\frac12E_{\varnothing,\varnothing}
\)
且 \((D^{(q)})^{-1}_{u,v}=(-1)^{|u\cap v|}\mathbf{1}_{\{u\cup v=[q]\}}\)；
将 \(U_q\cong\mathcal P([q])\) 复原为按位 \(\wedge,\vee\) 记号并用 \(T^{(q)}=\frac12B^{(q)}\) 即得所示整数核。证毕。
\end{proof}

\begin{corollary}[两层核的行列式闭式：仅由 \texorpdfstring{$(a,b,N)$}{(a,b,N)} 决定]\label{cor:xi-terminal-replica-softcore-twolevel-kernel-det}
设 \(q\ge 2\) 并记 \(N:=2^q\)。
对任意标量 \(a,b\in\CC\)，定义两层核矩阵 \(T^{(q)}(a,b)\in\CC^{U_q\times U_q}\) 为
\[
T^{(q)}(a,b)_{u,v}:=
\begin{cases}
a,&u\wedge v=\mathbf{0},\\[2pt]
b,&u\wedge v\neq \mathbf{0}.
\end{cases}
\]
则成立严格闭式
\[
\boxed{\ \det\bigl(T^{(q)}(a,b)\bigr)=a\,(a-b)^{N-1}\ }.
\]
\end{corollary}
\begin{proof}
当 \(a=b\) 时，\(T^{(q)}(a,a)=a\,\mathbf{1}\mathbf{1}^{\mathsf T}\) 为秩一矩阵且 \(N\ge 4\)，故行列式为 \(0\)，与右端一致。
以下设 \(a\neq b\)。
\par
由定义
\(
T^{(q)}(a,b)=b\,\mathbf{1}\mathbf{1}^{\mathsf T}+(a-b)D^{(q)}
\)。
由于 \(D^{(q)}=K^{\otimes q}\) 且 \(q\ge 2\)，有 \(\det(D^{(q)})=1\)（见结论 \ref{con:xi-terminal-disjointness-bq-smith-snf-coker-parity} 的证明）。
并且由 \(\bigl(D^{(q)}\bigr)^{-1}\mathbf{1}=e_{\varnothing}\) 得
\(\mathbf{1}^{\mathsf T}\bigl(D^{(q)}\bigr)^{-1}\mathbf{1}=1\)。
对矩阵行列式引理应用于
\(
A:=(a-b)D^{(q)}
\)
与秩一更新 \(b\,\mathbf{1}\mathbf{1}^{\mathsf T}\)，得到
\[
\det\bigl(T^{(q)}(a,b)\bigr)
=\det(A)\Bigl(1+b\,\mathbf{1}^{\mathsf T}A^{-1}\mathbf{1}\Bigr)
=(a-b)^{N}\det(D^{(q)})\Bigl(1+\frac{b}{a-b}\Bigr)
=a\,(a-b)^{N-1}.
\]
证毕。
\end{proof}

\begin{corollary}[逆核支撑的精确计数：\texorpdfstring{$3^q+1$}{3^q+1} 稀疏度与逐行分布]\label{cor:xi-terminal-replica-softcore-inverse-support-count}
沿用推论 \ref{cor:xi-terminal-replica-softcore-Tq-det-and-integer-inverse-kernel} 的显式逆核。
则 \(\bigl(T^{(q)}\bigr)^{-1}\) 的非零元总数满足
\[
\#\Bigl\{(u,v)\in U_q^2:\ \bigl(T^{(q)}\bigr)^{-1}_{u,v}\neq 0\Bigr\}=3^q+1.
\]
更精确地，若 \(u\neq \mathbf{0}\)，则第 \(u\) 行恰有 \(2^{|u|}\) 个非零元，且其取值集合包含于 \(\{\pm 2\}\)；
而第 \(\mathbf{0}\) 行恰有两个非零元：
\[
\bigl(T^{(q)}\bigr)^{-1}_{\mathbf{0},\mathbf{0}}=-1,
\qquad
\bigl(T^{(q)}\bigr)^{-1}_{\mathbf{0},\mathbf{1}}=2.
\]
\end{corollary}
\begin{proof}
由推论 \ref{cor:xi-terminal-replica-softcore-Tq-det-and-integer-inverse-kernel}，
当 \((u,v)\neq(\mathbf{0},\mathbf{0})\) 时有
\(
\bigl(T^{(q)}\bigr)^{-1}_{u,v}\neq 0
\iff
u\vee v=\mathbf{1}
\)，且此时取值必为 \(\pm 2\)。
因此对固定 \(u\neq \mathbf{0}\)，条件 \(u\vee v=\mathbf{1}\) 要求 \(v\) 在 \(\{j:\ u_j=0\}\) 上强制取 \(1\)，并在 \(\{j:\ u_j=1\}\) 上自由取 \(0/1\)，从而解的个数为 \(2^{|u|}\)。
对 \(u=\mathbf{0}\) 时仅有 \(v=\mathbf{1}\) 满足 \(u\vee v=\mathbf{1}\)，且 \((\mathbf{0},\mathbf{0})\) 处另有修正项 \(-1\)，得到所示两处非零值。
\par
将 \(2^{|u|}\) 对所有 \(u\in U_q\) 求和得
\[
\sum_{u\in U_q}2^{|u|}
=\sum_{k=0}^{q}\binom{q}{k}2^{k}
=(1+2)^q
=3^q,
\]
对应于 \(u\vee v=\mathbf{1}\) 产生的全部非零元；再加上额外的 \((\mathbf{0},\mathbf{0})\) 条目即得总数 \(3^q+1\)。
证毕。
\end{proof}

\begin{conclusion}[Fibonacci 幂矩恒等式与 Perron 根的收敛级数方程]\label{con:xi-terminal-replica-softcore-fib-moment-and-perron-series}
沿用结论 \ref{con:xi-terminal-replica-softcore-fiber-splitting-full-spectrum} 的记号。
对任意整数 \(m\ge 0\)，成立严格幂矩恒等式
\[
\boxed{\ \mathbf{1}^{\mathsf T}(D^{(q)})^{m}\mathbf{1}=F_{m+3}^{\,q}\ }.
\]
记 Perron 根（最大特征值）为 \(\rho(T^{(q)})\)。
则当 \(\lambda>\varphi^{q}/2\) 时有收敛的 resolvent 展开，并且 \(\lambda=\rho(T^{(q)})\) 满足严格方程
\[
\boxed{\ 1=\frac12\sum_{m\ge 0}\frac{F_{m+3}^{\,q}}{2^{m}\lambda^{m+1}}\ } .
\]
\end{conclusion}
\begin{proof}
由 \(D^{(q)}=K^{\otimes q}\) 与 \(\mathbf{1}=u^{\otimes q}\)（\(u=(1,1)^{\mathsf T}\)），
\[
\mathbf{1}^{\mathsf T}(D^{(q)})^{m}\mathbf{1}
=\bigl(u^{\mathsf T}K^{m}u\bigr)^q.
\]
而 \(K\) 为 Fibonacci \(Q\)-矩阵，满足
\(
K^{m}=\bigl(\begin{smallmatrix}F_{m+1}&F_m\\ F_m&F_{m-1}\end{smallmatrix}\bigr)
\)；
故 \(u^{\mathsf T}K^{m}u=F_{m+1}+2F_m+F_{m-1}=F_{m+3}\)，得到幂矩恒等式。
\par
对 Perron 根，取正特征向量 \(x\) 并令 \(\mathbf{1}^{\mathsf T}x\neq 0\)。
由
\(
(\lambda I-\tfrac12D^{(q)})x=\tfrac12(\mathbf{1}^{\mathsf T}x)\mathbf{1}
\)
消去 \(\mathbf{1}^{\mathsf T}x\) 得
\(
1=\tfrac12\,\mathbf{1}^{\mathsf T}(\lambda I-\tfrac12D^{(q)})^{-1}\mathbf{1}
\)。
当 \(\lambda>\|D^{(q)}/2\|=\varphi^{q}/2\) 时，Neumann 级数
\[
(\lambda I-\tfrac12D^{(q)})^{-1}
=\sum_{m\ge 0}\frac{(D^{(q)})^m}{2^{m}\lambda^{m+1}}
\]
收敛；代入并用幂矩恒等式即得所示级数方程。证毕。
\end{proof}

\input{subsubsec__xi-time-protocol-conclusions-part30a-boolean-kernel-rigidity}

\endinput


\begin{conclusion}[\texorpdfstring{$\lambda_q$}{lambda_q} 的多尺度转级数：整数系数与 Fibonacci 比率半群]\label{con:xi-terminal-replica-softcore-lambdaq-multiscale-transseries}
令 Fibonacci 数列满足 \(F_{n+2}=F_{n+1}+F_n\) 且 \(F_1=F_2=1\)，并对 \(k\ge 2\) 定义
\[
r_k:=\frac{F_{k+2}}{2^{k}}\in(0,1).
\]
则当 \(q\to\infty\) 时，\(\lambda_q\) 存在一个由乘法半群 \(\langle r_2,r_3,r_4,\dots\rangle\) 支撑的多尺度渐近展开；其前若干主导项可取为
\[
\begin{aligned}
\lambda_q
&=2^{q}
+\Big(\frac{3}{2}\Big)^{q}
+\Big(\frac{5}{4}\Big)^{q}
-\Big(\frac{9}{8}\Big)^{q}
+1
-3\Big(\frac{15}{16}\Big)^{q}
+2\Big(\frac{27}{32}\Big)^{q}
+\Big(\frac{13}{16}\Big)^{q}\\
&\quad
-2\Big(\frac{25}{32}\Big)^{q}
-4\Big(\frac{3}{4}\Big)^{q}
+10\Big(\frac{45}{64}\Big)^{q}
+O\!\Big(\Big(\frac{21}{32}\Big)^{q}\Big).
\end{aligned}
\]
其中各底数均可写为 \(2\cdot \prod_{i=1}^m r_{k_i}\) 的形式，且对应系数为整数（由固定点方程 \(A_q(\rho(q))=1\) 的 Lagrange 反演所诱导的组合计数给出）。
\end{conclusion}
\begin{proof}[证明要点]
由定义 \ref{def:pom-multiplicity-real-q-pressure} 的方程 \(A_q(\rho)=\sum_{k\ge 1}F_{k+2}^q\rho^k=1\)。
作换元 \(\rho=2^{-q}y\) 得
\(
1=y+\sum_{k\ge 2}r_k^q\,y^k
\)
且 \(y\to 1\)。
对该隐式方程应用 Lagrange 反演并在 \(\lambda_q=\rho^{-1}=2^q/y\) 上作展开即可得到所示多尺度项；
系数的整性由反演公式中的树形组合系数保证。
\end{proof}

\begin{conclusion}[\texorpdfstring{$B^{(q)}$}{B^(q)} 的补集自对偶反演：与 \texorpdfstring{$D^{(q)}$}{D^(q)} 的共轭逆仅差一秩修正]\label{con:xi-terminal-disjointness-bq-complement-selfdual-inverse-rankone}
令 \(C\) 为补集置换矩阵（\((Ce_u)=e_{[q]\setminus u}\)），令 \(S=\mathrm{diag}(({-1})^{|u|})\)，并记 \(E_{\varnothing,\varnothing}\) 为在 \((\varnothing,\varnothing)\) 处为 \(1\)、其余为 \(0\) 的矩阵。
则对一切 \(q\ge 2\) 成立严格恒等式
\[
\boxed{\;
\bigl(B^{(q)}\bigr)^{-1}
=(-1)^{q}\,C^{\mathsf T}S\,D^{(q)}\,S C\;-\;\frac12\,E_{\varnothing,\varnothing}\; }.
\]
从而
\(
\bigl(B^{(q)}\bigr)^{-1}+\frac12E_{\varnothing,\varnothing}\in \mathrm{M}_{2^q}(\ZZ)
\)，且该整矩阵与 \((-1)^qD^{(q)}\) 共轭等谱；
而 \(\bigl(B^{(q)}\bigr)^{-1}\) 本身则为该共轭像的单秩负扰动。
\end{conclusion}
\begin{proof}
将 \(U_q\) 与 \(\mathcal P([q])\) 识别，并记 \(D^{(q)}_{u,v}=\mathbf{1}_{\{u\wedge v=0\}}\)。
直接计算可得对一切 \(u,v\subseteq[q]\) 有
\[
\bigl(D^{(q)}\bigr)^{-1}_{u,v}
=(-1)^{|u\cap v|}\,\mathbf{1}_{\{u\cup v=[q]\}}.
\]
事实上，
\[
\sum_{w\subseteq[q]}D^{(q)}_{u,w}\,(-1)^{|w\cap v|}\mathbf{1}_{\{w\cup v=[q]\}}
=\sum_{t\subseteq v\setminus u}(-1)^{|t|}
=(1-1)^{|v\setminus u|}
=\mathbf{1}_{\{u=v\}},
\]
从而 \(D^{(q)}\bigl(D^{(q)}\bigr)^{-1}=I\)。
另一方面，若 \(u\cup v=[q]\)，则 \(|u|+|v|=q+|u\cap v|\)，故
\[
\bigl(C^{\mathsf T}S\,D^{(q)}\,S C\bigr)_{u,v}
=(-1)^{|u|+|v|}\,D^{(q)}_{[q]\setminus u,[q]\setminus v}
=(-1)^{q+|u\cap v|}\mathbf{1}_{\{u\cup v=[q]\}},
\]
从而 \(((-1)^q C^{\mathsf T}S D^{(q)} S C)=\bigl(D^{(q)}\bigr)^{-1}\)。
最后对 \(B^{(q)}=D^{(q)}+\mathbf{1}\mathbf{1}^{\mathsf T}\) 应用 Sherman--Morrison 公式并用
\(
\bigl(D^{(q)}\bigr)^{-1}\mathbf{1}=e_{\varnothing}
\)
与
\(
\mathbf{1}^{\mathsf T}\bigl(D^{(q)}\bigr)^{-1}\mathbf{1}=1
\)
即得
\(
\bigl(B^{(q)}\bigr)^{-1}
=\bigl(D^{(q)}\bigr)^{-1}-\frac12E_{\varnothing,\varnothing}
\)，从而结论成立。
\end{proof}

\begin{conclusion}[Smith 标准形的刚性：余核恒为 \texorpdfstring{$\mathbb Z/2\mathbb Z$}{Z/2Z}，且唯一可解障碍为 \texorpdfstring{$b_{\varnothing}$}{b_empty} 的奇偶性]\label{con:xi-terminal-disjointness-bq-smith-snf-coker-parity}
对任意 \(q\ge 2\)，整数矩阵 \(B^{(q)}\) 的 Smith 标准形为
\[
\mathrm{SNF}(B^{(q)})=\mathrm{diag}(1,\dots,1,2),
\]
从而 \(\det B^{(q)}=2\) 且
\(
\mathrm{coker}(B^{(q)}:\ZZ^{2^q}\to\ZZ^{2^q})\cong \ZZ/2\ZZ
\)。
等价地，对任意 \(b\in\ZZ^{2^q}\)，丢番图方程
\[
B^{(q)}x=b,\qquad x\in\ZZ^{2^q},
\]
有解当且仅当 \(b_{\varnothing}\) 为偶数；且在此条件下解空间为一个满秩格的陪集。
\end{conclusion}
\begin{proof}
由结论 \ref{con:xi-terminal-disjointness-bq-complement-selfdual-inverse-rankone}，
\(
\bigl(B^{(q)}\bigr)^{-1}=\mathsf M-\frac12E_{\varnothing,\varnothing}
\)
其中 \(\mathsf M\in \mathrm{M}_{2^q}(\ZZ)\)。
因此若 \(x=\bigl(B^{(q)}\bigr)^{-1}b\)，则除 \(\varnothing\)-坐标外各分量自动为整数，而
\(
x_{\varnothing}=\mathsf M_{\varnothing,\bullet}b-\frac12 b_{\varnothing}
\)
为整数当且仅当 \(b_{\varnothing}\) 为偶数；这给出丢番图可解性的充要条件，并推出余核为 \(\ZZ/2\ZZ\)。
另一方面，\(\det\!\bigl(D^{(q)}\bigr)=1\)（\(q\ge 2\)）且
\(
\det(B^{(q)})=\det(D^{(q)})\bigl(1+\mathbf{1}^{\mathsf T}(D^{(q)})^{-1}\mathbf{1}\bigr)=2
\)
（由 Sherman--Morrison 行列式公式）。
结合 \(\gcd\) 的 \(1\times 1\) 余子式为 \(1\) 且 \(|\det(B^{(q)})|=2\)，得到 \(\mathrm{SNF}(B^{(q)})=\mathrm{diag}(1,\dots,1,2)\)。
\end{proof}

\begin{conclusion}[伴随矩阵的稀疏闭式：所有余子式仅取 \texorpdfstring{$\{0,\pm 2\}$}{\{0,±2\}} 与唯一的 \texorpdfstring{$-1$}{-1}]\label{con:xi-terminal-disjointness-bq-adjugate-sparse-minors}
记 \(\mathrm{adj}(B^{(q)})\) 为伴随矩阵。
则对 \(q\ge 2\)，其条目满足完全稀疏闭式：
\[
\mathrm{adj}(B^{(q)})_{\varnothing,\varnothing}=-1,
\qquad
\mathrm{adj}(B^{(q)})_{u,v}=
\begin{cases}
2(-1)^{|u\cap v|},&u\cup v=[q],\\
0,&\text{其余情形}.
\end{cases}
\]
因此，所有 \((2^q-1)\times(2^q-1)\) 余子式除上述覆盖情形外全部消失，且非零者在绝对值意义下至多为 \(2\)。
\end{conclusion}
\begin{proof}
由结论 \ref{con:xi-terminal-disjointness-bq-smith-snf-coker-parity} 有 \(\det(B^{(q)})=2\)。
再由结论 \ref{con:xi-terminal-disjointness-bq-complement-selfdual-inverse-rankone} 得到
\[
\bigl(B^{(q)}\bigr)^{-1}_{\varnothing,\varnothing}=-\frac12,
\qquad
\bigl(B^{(q)}\bigr)^{-1}_{u,v}=(-1)^{|u\cap v|}\mathbf{1}_{\{u\cup v=[q]\}}\quad((u,v)\neq(\varnothing,\varnothing)).
\]
于是 \(\mathrm{adj}(B^{(q)})=\det(B^{(q)})\bigl(B^{(q)}\bigr)^{-1}=2\bigl(B^{(q)}\bigr)^{-1}\)，逐项即得所示闭式。
\end{proof}

\begin{conclusion}[例外谱平方和的纯整闭式：\texorpdfstring{$4^{q-1}+\frac12 3^q+\frac14F_{2q+2}$}{4^(q-1)+1/2 3^q+1/4 F_{2q+2}}]\label{con:xi-terminal-replica-softcore-exceptional-spectrum-square-sum}
令 \(\{\lambda_i^{(q)}\}_{i=0}^{q}\) 表示 \(T^{(q)}\) 的 \(q+1\) 个例外本征值（亦即 \(S_q\)-对称降维矩阵 \(A^{(q)}\) 的谱，命题 \ref{prop:pom-multiplicity-composition-sq-symmetry-reduction}）。
则对一切 \(q\ge 1\)，成立严格恒等式
\[
\boxed{\;
\sum_{i=0}^{q}\big(\lambda_i^{(q)}\big)^2
=4^{q-1}
\;+\;\frac12\,3^{q}
\;+\;\frac14\,F_{2q+2}\; }.
\]
该恒等式给出一个谱--组合桥：例外谱的二阶幂和在代数上塌缩为仅含 \(4^q\)、\(3^q\) 与 Fibonacci 指标 \(2q+2\) 的整系数表达。
\end{conclusion}
\begin{proof}
由于 \(T^{(q)}\) 对称，
\(
\Tr\!\bigl((T^{(q)})^2\bigr)=\sum_{u,v\in U_q}(T^{(q)}_{u,v})^2
\)。
其中 \(T^{(q)}_{u,v}=1\) 当且仅当 \(u\wedge v=0\)；逐坐标计数得该类对数为 \(3^q\)，其余为 \(4^q-3^q\)。
因此
\[
\Tr\!\bigl((T^{(q)})^2\bigr)
=3^q+\frac14(4^q-3^q)
=4^{q-1}+\frac34\,3^q.
\]
另一方面，由定理 \ref{thm:pom-replica-softcore-characteristic-factorization}，固定谱簇
\(
\{\frac{(-1)^j\varphi^{q-2j}}{2}\}_{j=0}^{q}
\)
在 \(T^{(q)}\) 中的重数为 \(\binom{q}{j}-1\)。
故其平方和贡献为
\[
\frac14\sum_{j=0}^{q}\bigl(\binom{q}{j}-1\bigr)\varphi^{2(q-2j)}
=\frac14\Bigl(\sum_{j=0}^{q}\binom{q}{j}\varphi^{2(q-2j)}-\sum_{j=0}^{q}\varphi^{2(q-2j)}\Bigr).
\]
利用恒等式 \(\varphi^2+\varphi^{-2}=3\) 得
\(
\sum_{j=0}^{q}\binom{q}{j}\varphi^{2(q-2j)}=(\varphi^2+\varphi^{-2})^q=3^q
\)；
又由等比求和与 \(\varphi^4-1=\sqrt5\,\varphi^2\) 得
\(
\sum_{j=0}^{q}\varphi^{2(q-2j)}=F_{2q+2}
\)。
因此固定谱簇的平方和为 \(\frac14(3^q-F_{2q+2})\)。
从 \(\Tr((T^{(q)})^2)=\sum_i(\lambda_i^{(q)})^2+\)（固定谱簇平方和）得出
\[
\sum_{i=0}^{q}\big(\lambda_i^{(q)}\big)^2
=\Bigl(4^{q-1}+\frac34\,3^q\Bigr)-\frac14(3^q-F_{2q+2})
=4^{q-1}+\frac12\,3^q+\frac14\,F_{2q+2},
\]
证毕。
\end{proof}

\endinput


\subsubsection{例外谱幂和与二元 necklace 的指数分解链}\label{subsubsec:xi-exceptional-spectrum-necklace-chain}

\subsubsection{例外谱幂和与低阶迹矩的指数分解：Lucas--Pell 纠正项、符号型递推与有限枚举核验}\label{subsubsec:xi-exceptional-power-sum-low-moment-decomposition}
令 \(q\ge 2\)，并取 \(U_q=\{0,1\}^q\)；将 \(U_q\) 与 \(\mathcal P([q])\) 通过指示向量识别，并以 \(u\wedge v\) 表示逐坐标与。
定义对称矩阵 \(T^{(q)}\in\RR^{U_q\times U_q}\) 为
\[
T^{(q)}_{u,v}=
\begin{cases}
1,&u\wedge v=0,\\[2pt]
\frac12,&u\wedge v\neq 0,
\end{cases}
\qquad (u,v\in U_q).
\]
由命题 \ref{prop:pom-multiplicity-composition-sq-symmetry-reduction}，\(T^{(q)}\) 在 \(S_q\)-对称子空间上的谱由显式 \((q+1)\times(q+1)\) 矩阵 \(A^{(q)}\) 给出；记该 \((q+1)\) 个单根特征值（例外谱）为 \(\{\lambda_i^{(q)}\}_{i=0}^{q}\)，并对整数 \(m\ge 1\) 定义幂和
\[
S_m(q):=\sum_{i=0}^{q}\big(\lambda_i^{(q)}\big)^m.
\]

\begin{conclusion}[例外谱一次幂和的整系数塌缩：\texorpdfstring{$S_1(q)=2^{q-1}+\frac12F_{q+1}$}{S1(q)=2^{q-1}+1/2 F_{q+1}}]\label{con:xi-terminal-replica-softcore-exceptional-power-sum-m1-integer-collapse}
令 Fibonacci 数列满足 \(F_0=0\)、\(F_1=1\)。
则对一切 \(q\ge 2\) 成立严格闭式
\[
\boxed{\;
S_1(q)=2^{q-1}+\frac12\,F_{q+1}\; }.
\]
因此例外谱总和在 \(\QQ(\sqrt5)\) 中出现的 \(\varphi\)-表示可完全消去，并在 \(\QQ\) 中塌缩为半整数序列。
\end{conclusion}
\begin{proof}
由于 \(S_1(q)=\Tr(A^{(q)})\)，由 \eqref{eq:pom-multiplicity-composition-Aq-def} 直接计算
\[
\Tr(A^{(q)})
=\frac12\sum_{k=0}^{q}\binom{q}{k}+\frac12\sum_{k=0}^{q}\binom{q-k}{k}
=2^{q-1}+\frac12\sum_{k\ge 0}\binom{q-k}{k}.
\]
再用经典恒等式 \(\sum_{k\ge 0}\binom{q-k}{k}=F_{q+1}\) 即得结论（亦与定理 \ref{thm:pom-replica-softcore-exceptional-spectrum-trace} 等价）。证毕。
\end{proof}

\begin{conclusion}[Lucas--Pell 纠正项的统一生成与符号型递推律]\label{con:xi-terminal-replica-softcore-lucas-pell-omega-unified}
令 Lucas 数列满足 \(L_0=2\)、\(L_1=1\)，并令 \(\varphi=\frac{1+\sqrt5}{2}\)。
对任意整数 \(m\ge 1\) 定义
\[
\Omega_m(q):=
\begin{cases}
\dfrac{\varphi^{m(q+2)}+(-1)^q\varphi^{-mq}}{\varphi^{2m}+1},& m\ \text{为奇数},\\[10pt]
\dfrac{\varphi^{m(q+2)}-\varphi^{-mq}}{\varphi^{2m}-1},& m\ \text{为偶数}.
\end{cases}
\]
则对每个 \(m\ge 1\) 与 \(q\ge 0\)，有 \(\Omega_m(q)\in\ZZ\)，并满足统一的二阶线性递推
\[
\boxed{\;
\Omega_m(q)=L_m\,\Omega_m(q-1)-(-1)^m\,\Omega_m(q-2),\qquad
\Omega_m(0)=1,\ \Omega_m(1)=L_m\; }.
\]
尤其地，\(\Omega_1(q)=F_{q+1}\)、\(\Omega_2(q)=F_{2q+2}\) 为该机制的低阶退化。
\end{conclusion}
\begin{proof}
记 \(\widehat\varphi:=-\varphi^{-1}\) 为 \(\QQ(\sqrt5)\) 的共轭根，则
\(
L_m=\varphi^m+\widehat\varphi^{\,m}=\varphi^m+(-1)^m\varphi^{-m}
\)
且
\(
\varphi^m\widehat\varphi^{\,m}=(-1)^m
\)。
因此数列 \(q\mapsto \varphi^{mq}\) 与 \(q\mapsto \widehat\varphi^{\,mq}\) 均满足线性递推
\(
x_q=L_m x_{q-1}-(-1)^m x_{q-2}
\)；
由定义可将 \(\Omega_m(q)\) 写成二者的 \(\QQ\)-线性组合，从而递推式成立。
初值由 \(q=0,1\) 的直接代入得到。
递推系数与初值均为整数，故 \(\Omega_m(q)\in\ZZ\) 由归纳法给出。
最后，\(m=1,2\) 的退化式可由相同递推与初值对齐（或由 \(\Omega_1\)、\(\Omega_2\) 的闭式与 Binet 公式直接化简）得到。证毕。
\end{proof}

\begin{conclusion}[低阶迹矩的完全指数分解：\texorpdfstring{$\Tr((T^{(q)})^m)$}{Tr((T^(q))^m)} 的闭式]\label{con:xi-terminal-replica-softcore-low-moment-trace-exponential-decomposition}
对任意 \(q\ge 2\)，\(T^{(q)}\) 的低阶迹矩满足严格闭式
\[
\boxed{\;
\Tr\!\bigl((T^{(q)})^3\bigr)=\frac18\Big(8^{q}+3\cdot 6^{q}+3\cdot 5^{q}+4^{q}\Big)\; } ,
\]
\[
\boxed{\;
\Tr\!\bigl((T^{(q)})^4\bigr)=\frac1{16}\Big(16^{q}+4\cdot 12^{q}+4\cdot 10^{q}+2\cdot 9^{q}+4\cdot 8^{q}+7^{q}\Big)\; } ,
\]
\[
\boxed{\;
\Tr\!\bigl((T^{(q)})^5\bigr)=\frac1{32}\Big(32^{q}+5\cdot 24^{q}+5\cdot 20^{q}+5\cdot 18^{q}+5\cdot 16^{q}+5\cdot 15^{q}+5\cdot 13^{q}+11^{q}\Big)\; } ,
\]
\[
\boxed{\;
\Tr\!\bigl((T^{(q)})^6\bigr)=\frac1{64}\Big(
64^{q}+6\cdot48^{q}+6\cdot40^{q}+9\cdot36^{q}+6\cdot32^{q}+12\cdot30^{q}
+2\cdot27^{q}+6\cdot26^{q}+3\cdot25^{q}+6\cdot24^{q}+6\cdot21^{q}+18^{q}\Big)\; } ,
\]
\[
\boxed{\;
\Tr\!\bigl((T^{(q)})^7\bigr)=\frac1{128}\Big(
128^{q}+7\cdot96^{q}+7\cdot80^{q}+14\cdot72^{q}+7\cdot64^{q}+21\cdot60^{q}
+7\cdot54^{q}+7\cdot52^{q}+7\cdot50^{q}+14\cdot48^{q}+7\cdot45^{q}
+7\cdot42^{q}+7\cdot40^{q}+7\cdot39^{q}+7\cdot34^{q}+29^{q}\Big)\; } .
\]
\end{conclusion}
\begin{proof}[证明要点]
由 \(T^{(q)}_{u,v}=\frac12\bigl(1+\mathbf{1}_{\{u\wedge v=0\}}\bigr)\) 展开
\(
\Tr((T^{(q)})^m)=\sum_{u_1,\dots,u_m}\prod_{k=1}^{m}T^{(q)}_{u_k,u_{k+1}}
\)
（其中 \(u_{m+1}=u_1\)），得到
\[
\Tr\!\bigl((T^{(q)})^m\bigr)
=2^{-m}\sum_{E'\subseteq E(C_m)}\ \sum_{u_1,\dots,u_m\in U_q}\ \prod_{(k,k+1)\in E'}\mathbf{1}_{\{u_k\wedge u_{k+1}=0\}}.
\]
对固定 \(E'\)，指示乘积在各坐标上分解，故内层求和按坐标完全因子化为 \(a(E')^q\)，其中 \(a(E')\) 为单坐标二元 \(m\)-环上满足相应“禁相邻 \(11\)”约束的可行赋值数。
当 \(m=3,4,5\) 时，\(E'\subseteq E(C_m)\) 的同构类型有限，可逐型枚举得到对应的 \(a(E')\) 分别为
\[
8,6,5,4;\qquad
16,12,10,9,8,7;\qquad
32,24,20,18,16,15,13,11,
\]
且其出现次数分别为
\[
1,3,3,1;\qquad
1,4,4,2,4,1;\qquad
1,5,5,5,5,5,5,1.
\]
当 \(m=6,7\) 时同理可枚举得到相应的 \(a(E')\) 与出现次数；
等价地，可由分拆闭式结论 \ref{con:xi-terminal-replica-softcore-partition-trace-closed-form} 在 \(m=6,7\) 的特化直接读出底数与重数。
代回并收集同底数项即得五条闭式。证毕。
\end{proof}

\begin{conclusion}[例外谱三次幂和的全闭式：\texorpdfstring{$\{8,6,5\}$}{\{8,6,5\}} 指数族与 Lucas--Pell 纠正项]\label{con:xi-terminal-replica-softcore-exceptional-power-sum-m3}
定义整数序列
\[
\Omega_3(q):=\frac{\varphi^{3q+6}+(-1)^q\varphi^{-3q}}{\varphi^{6}+1}\in\ZZ.
\]
则对一切 \(q\ge 2\) 有严格恒等式
\[
\boxed{\;
S_3(q)=\frac18\Big(8^{q}+3\cdot 6^{q}+3\cdot 5^{q}+\Omega_3(q)\Big)\; }.
\]
并且 \(\Omega_3\) 满足二阶线性递推
\[
\Omega_3(q)=4\,\Omega_3(q-1)+\Omega_3(q-2),\qquad \Omega_3(0)=1,\ \Omega_3(1)=4.
\]
\end{conclusion}
\begin{proof}[证明要点]
由定理 \ref{thm:pom-replica-softcore-characteristic-factorization}，\(T^{(q)}\) 的全谱由例外谱与固定谱簇并置给出。
因此
\(
\Tr((T^{(q)})^3)=S_3(q)+\sum_{j=0}^{q}(\binom{q}{j}-1)\bigl(\frac{(-1)^j\varphi^{q-2j}}{2}\bigr)^3
\)。
对固定谱簇幂和作二项式求和与几何级数求和可得
\[
\sum_{j=0}^{q}\Bigl(\binom{q}{j}-1\Bigr)\Bigl(\frac{(-1)^j\varphi^{q-2j}}{2}\Bigr)^3
=\frac18\bigl(4^q-\Omega_3(q)\bigr),
\]
其中 \(\varphi^3-\varphi^{-3}=4\) 被识别为 Lucas 数 \(L_3\)。
再与结论 \ref{con:xi-terminal-replica-softcore-low-moment-trace-exponential-decomposition} 的 \(\Tr((T^{(q)})^3)\) 闭式合并即得所示。
\(\Omega_3\) 的递推为结论 \ref{con:xi-terminal-replica-softcore-lucas-pell-omega-unified} 在 \(m=3\) 的特例。证毕。
\end{proof}

\begin{conclusion}[例外谱四次幂和的全闭式：\texorpdfstring{$\{16,12,10,9,8\}$}{\{16,12,10,9,8\}} 指数族与 Lucas--Pell 纠正项]\label{con:xi-terminal-replica-softcore-exceptional-power-sum-m4}
定义整数序列
\[
\Omega_4(q):=\frac{\varphi^{4q+8}-\varphi^{-4q}}{\varphi^{8}-1}\in\ZZ.
\]
则对一切 \(q\ge 2\) 有严格恒等式
\[
\boxed{\;
S_4(q)=\frac1{16}\Big(16^{q}+4\cdot 12^{q}+4\cdot 10^{q}+2\cdot 9^{q}+4\cdot 8^{q}+\Omega_4(q)\Big)\; }.
\]
并且 \(\Omega_4\) 满足二阶线性递推
\[
\Omega_4(q)=7\,\Omega_4(q-1)-\Omega_4(q-2),\qquad \Omega_4(0)=1,\ \Omega_4(1)=7.
\]
\end{conclusion}
\begin{proof}[证明要点]
同上，\(
\Tr((T^{(q)})^4)=S_4(q)+\sum_{j}(\binom{q}{j}-1)\bigl(\frac{(-1)^j\varphi^{q-2j}}{2}\bigr)^4
\)。
固定谱簇幂和化简为 \(\frac1{16}(7^q-\Omega_4(q))\)（其中 \(\varphi^4+\varphi^{-4}=7=L_4\)），再与结论 \ref{con:xi-terminal-replica-softcore-low-moment-trace-exponential-decomposition} 的 \(\Tr((T^{(q)})^4)\) 闭式合并即得所示。
递推为结论 \ref{con:xi-terminal-replica-softcore-lucas-pell-omega-unified} 在 \(m=4\) 的特例。证毕。
\end{proof}

\begin{conclusion}[例外谱五次幂和的全闭式：\texorpdfstring{$\{32,24,20,18,16,15,13\}$}{\{32,24,20,18,16,15,13\}} 指数族与 Lucas--Pell 纠正项]\label{con:xi-terminal-replica-softcore-exceptional-power-sum-m5}
定义整数序列
\[
\Omega_5(q):=\frac{\varphi^{5q+10}+(-1)^q\varphi^{-5q}}{\varphi^{10}+1}\in\ZZ.
\]
则对一切 \(q\ge 2\) 有严格恒等式
\[
\boxed{\;
S_5(q)=\frac1{32}\Big(32^{q}+5\cdot 24^{q}+5\cdot 20^{q}+5\cdot 18^{q}+5\cdot 16^{q}+5\cdot 15^{q}+5\cdot 13^{q}+\Omega_5(q)\Big)\; }.
\]
并且 \(\Omega_5\) 满足二阶线性递推
\[
\Omega_5(q)=11\,\Omega_5(q-1)+\Omega_5(q-2),\qquad \Omega_5(0)=1,\ \Omega_5(1)=11.
\]
\end{conclusion}
\begin{proof}[证明要点]
同理，固定谱簇五次幂和化简为 \(\frac1{32}(11^q-\Omega_5(q))\)（其中 \(\varphi^5-\varphi^{-5}=11=L_5\)），再与结论 \ref{con:xi-terminal-replica-softcore-low-moment-trace-exponential-decomposition} 的 \(\Tr((T^{(q)})^5)\) 闭式合并即可得到所示表达。
递推为结论 \ref{con:xi-terminal-replica-softcore-lucas-pell-omega-unified} 在 \(m=5\) 的特例。证毕。
\end{proof}

\begin{conclusion}[例外谱六次幂和的全闭式：\texorpdfstring{$\{64,48,40,36,32,30,27,26,25,24,21\}$}{\{64,48,40,36,32,30,27,26,25,24,21\}} 指数族与 Lucas--Pell 纠正项]\label{con:xi-terminal-replica-softcore-exceptional-power-sum-m6}
定义整数序列
\[
\Omega_6(q):=\frac{\varphi^{6q+12}-\varphi^{-6q}}{\varphi^{12}-1}\in\ZZ.
\]
则对一切 \(q\ge 2\) 有严格恒等式
\[
\boxed{\;
S_6(q)=\frac1{64}\Big(
64^{q}+6\cdot48^{q}+6\cdot40^{q}+9\cdot36^{q}+6\cdot32^{q}+12\cdot30^{q}
+2\cdot27^{q}+6\cdot26^{q}+3\cdot25^{q}+6\cdot24^{q}+6\cdot21^{q}+\Omega_6(q)\Big)\; } .
\]
并且 \(\Omega_6\) 满足二阶线性递推
\[
\Omega_6(q)=18\,\Omega_6(q-1)-\Omega_6(q-2),\qquad \Omega_6(0)=1,\ \Omega_6(1)=18.
\]
\end{conclusion}
\begin{proof}[证明要点]
同理，\(
\Tr((T^{(q)})^6)=S_6(q)+\sum_{j=0}^{q}(\binom{q}{j}-1)\bigl(\frac{(-1)^j\varphi^{q-2j}}{2}\bigr)^6
\)。
固定谱簇六次幂和化简为 \(\frac1{64}(18^q-\Omega_6(q))\)（其中 \(\varphi^6+\varphi^{-6}=18=L_6\)），再与结论 \ref{con:xi-terminal-replica-softcore-low-moment-trace-exponential-decomposition} 的 \(\Tr((T^{(q)})^6)\) 闭式合并即得所示表达。
递推为结论 \ref{con:xi-terminal-replica-softcore-lucas-pell-omega-unified} 在 \(m=6\) 的特例。证毕。
\end{proof}

\begin{conclusion}[例外谱七次幂和的全闭式：\texorpdfstring{$\{128,96,80,72,64,60,54,52,50,48,45,42,40,39,34\}$}{\{128,96,80,72,64,60,54,52,50,48,45,42,40,39,34\}} 指数族与 Lucas--Pell 纠正项]\label{con:xi-terminal-replica-softcore-exceptional-power-sum-m7}
定义整数序列
\[
\Omega_7(q):=\frac{\varphi^{7q+14}+(-1)^q\varphi^{-7q}}{\varphi^{14}+1}\in\ZZ.
\]
则对一切 \(q\ge 2\) 有严格恒等式
\[
\boxed{\;
S_7(q)=\frac1{128}\Big(
128^{q}+7\cdot96^{q}+7\cdot80^{q}+14\cdot72^{q}+7\cdot64^{q}+21\cdot60^{q}
+7\cdot54^{q}+7\cdot52^{q}+7\cdot50^{q}+14\cdot48^{q}+7\cdot45^{q}
+7\cdot42^{q}+7\cdot40^{q}+7\cdot39^{q}+7\cdot34^{q}+\Omega_7(q)\Big)\; } .
\]
并且 \(\Omega_7\) 满足二阶线性递推
\[
\Omega_7(q)=29\,\Omega_7(q-1)+\Omega_7(q-2),\qquad \Omega_7(0)=1,\ \Omega_7(1)=29.
\]
\end{conclusion}
\begin{proof}[证明要点]
同理，固定谱簇七次幂和化简为 \(\frac1{128}(29^q-\Omega_7(q))\)（其中 \(\varphi^7-\varphi^{-7}=29=L_7\)），再与结论 \ref{con:xi-terminal-replica-softcore-low-moment-trace-exponential-decomposition} 的 \(\Tr((T^{(q)})^7)\) 闭式合并即可得到所示表达。
递推为结论 \ref{con:xi-terminal-replica-softcore-lucas-pell-omega-unified} 在 \(m=7\) 的特例。证毕。
\end{proof}

\begin{conclusion}[例外谱八次幂和的全闭式：\texorpdfstring{$\{256,192,160,144,128,120,108,104,100,96,90,84,81,80,78,75,72,68,65,64,63,55\}$}{\{256,192,160,144,128,120,108,104,100,96,90,84,81,80,78,75,72,68,65,64,63,55\}} 指数族与 Lucas--Pell 纠正项]\label{con:xi-terminal-replica-softcore-exceptional-power-sum-m8}
定义整数序列
\[
\Omega_8(q):=\frac{\varphi^{8q+16}-\varphi^{-8q}}{\varphi^{16}-1}\in\ZZ.
\]
则对一切 \(q\ge 2\) 有严格恒等式
\[
\boxed{\;
\begin{aligned}
S_8(q)=\frac1{256}\Big(
&256^{q}+8\cdot 192^{q}+8\cdot 160^{q}+20\cdot 144^{q}+8\cdot 128^{q}+32\cdot 120^{q}+16\cdot 108^{q}\\
&+8\cdot 104^{q}+12\cdot 100^{q}+24\cdot 96^{q}+24\cdot 90^{q}+8\cdot 84^{q}+2\cdot 81^{q}+16\cdot 80^{q}\\
&+16\cdot 78^{q}+8\cdot 75^{q}+8\cdot 72^{q}+8\cdot 68^{q}+8\cdot 65^{q}+4\cdot 64^{q}+8\cdot 63^{q}+8\cdot 55^{q}
+\Omega_8(q)\Big)
\end{aligned}\; } .
\]
并且 \(\Omega_8\) 满足二阶线性递推
\[
\Omega_8(q)=47\,\Omega_8(q-1)-\Omega_8(q-2),\qquad \Omega_8(0)=1,\ \Omega_8(1)=47.
\]
\end{conclusion}
\begin{proof}[证明要点]
由结论 \ref{con:xi-terminal-replica-softcore-exceptional-power-sum-cycle-sector-replacement} 取 \(m=8\)，得到
\[
2^8S_8(q)=\Omega_8(q)+\sum_{\substack{\omega\in\{0,1\}^8\\ \omega\neq 0^8}}\tau(\omega)^{q},
\]
其中 \(\tau(\omega)\) 为结论 \ref{con:xi-terminal-replica-softcore-binary-necklace-trace-formula} 定义的二元词迹。
对全部 \(255\) 个非零词的 \(\tau(\omega)\) 进行有限枚举并按同值归并，即得到所示底数簇与重数。
递推式为结论 \ref{con:xi-terminal-replica-softcore-lucas-pell-omega-unified} 在 \(m=8\) 的特例。证毕。
\end{proof}

\endinput



\subsubsection{二元 word/necklace 核：迹矩与例外幂和的全阶离散指数展开}\label{subsubsec:xi-binary-word-necklace-full-order-expansion}
前述低阶闭式可由一个统一的二元 word/necklace 展开推出；下述条目给出该机制在任意阶 \(m\) 上的完全形式，并把 \(\Omega_m\) 与 Fibonacci 乘法权重精确对齐。

\begin{conclusion}[二元 necklace 张量迹公式：\texorpdfstring{$\Tr((T^{(q)})^{m})$}{Tr((T^(q))^m)} 的完全离散指数展开]\label{con:xi-terminal-replica-softcore-binary-necklace-trace-formula}
令
\[
J:=\begin{pmatrix}1&1\\[2pt]1&1\end{pmatrix},
\qquad
K:=\begin{pmatrix}1&1\\[2pt]1&0\end{pmatrix}.
\]
则在 \(U_q=\{0,1\}^q\) 的自然张量基下有严格分解
\[
T^{(q)}=\frac12\bigl(J^{\otimes q}+K^{\otimes q}\bigr).
\]
对任意整数 \(m\ge 1\) 与二元词 \(\omega=\omega_1\cdots\omega_m\in\{0,1\}^m\)，定义
\[
A_{\omega_t}:=\begin{cases}K,&\omega_t=0,\\ J,&\omega_t=1,\end{cases}
\qquad
\tau(\omega):=\Tr\!\bigl(A_{\omega_1}A_{\omega_2}\cdots A_{\omega_m}\bigr)\in\ZZ_{\ge 0}.
\]
则对一切 \(q\ge 1\) 成立精确恒等式
\[
\boxed{\;
\Tr\!\bigl((T^{(q)})^{m}\bigr)=2^{-m}\sum_{\omega\in\{0,1\}^m}\tau(\omega)^{q}\; }.
\]
进一步，由迹的循环不变性 \(\tau(\omega)\) 在循环移位下不变。
令 \(\mathrm{Neck}_m^{(2)}:=\{0,1\}^m/\!\sim\) 为长度 \(m\) 的二元 necklace 集（模循环移位等价），
对 \(\eta\in \mathrm{Neck}_m^{(2)}\) 记 \(\#\mathrm{orb}(\eta)\) 为其旋转轨道大小，则
\[
\boxed{\;
\Tr\!\bigl((T^{(q)})^{m}\bigr)=2^{-m}\sum_{\eta\in \mathrm{Neck}_m^{(2)}}\#\mathrm{orb}(\eta)\,\tau(\eta)^{q}\; } ,
\]
其中 \(\tau(\eta)\) 由任意代表元定义且良定。
并且 \(\tau(\eta)\) 具有完全的 Fibonacci--Lucas 乘法形：
\begin{enumerate}
  \item 若 \(\eta=0^m\)，则 \(\tau(\eta)=L_m\)；
  \item 若 \(\eta\) 含至少一个 \(1\)，将 \(\eta\) 以循环方式写成
  \[
  \eta\sim 1\,0^{r_1}\,1\,0^{r_2}\cdots 1\,0^{r_s},
  \qquad
  s\ge 1,\ r_i\ge 0,\ \sum_{i=1}^s r_i=m-s,
  \]
  则
  \[
  \tau(\eta)=\prod_{i=1}^{s}F_{r_i+3}.
  \]
\end{enumerate}
\end{conclusion}
\begin{proof}[证明要点]
由命题 \ref{prop:pom-multiplicity-composition-replica-explicit-eigs}，
\(
T^{(q)}=\frac12(\mathbf{1}\mathbf{1}^{\mathsf T}+K^{\otimes q})
\)；
注意 \(\mathbf{1}\mathbf{1}^{\mathsf T}=J^{\otimes q}\)，得到所示张量分解。
对 \(\Tr((T^{(q)})^m)\) 展开并用 \((A^{\otimes q})(B^{\otimes q})=(AB)^{\otimes q}\) 与 \(\Tr(M^{\otimes q})=(\Tr M)^q\) 得到 word 求和式。
循环移位不变性由 \(\Tr(X_1\cdots X_m)=\Tr(X_mX_1\cdots X_{m-1})\) 给出。
\par
若 \(\eta=0^m\)，则 \(\tau(0^m)=\Tr(K^m)=L_m\)（由 Fibonacci 矩阵恒等式 \(K^m=\bigl(\begin{smallmatrix}F_{m+1}&F_m\\ F_m&F_{m-1}\end{smallmatrix}\bigr)\) 得 \(\Tr(K^m)=F_{m+1}+F_{m-1}=L_m\)，见 \cite{FQFibonacciLucasChap11Matrices,MSEFibonacciMatrices784710}）。
若 \(\eta\) 含至少一个 \(1\)，令 \(u=(1,1)^{\top}\) 则 \(J=uu^{\top}\)，并有
\(
J K^{r} J=(u^{\top}K^{r}u)\,J
\)
与
\(
\Tr(JK^{r})=u^{\top}K^{r}u
\)。
再由 \(u^{\top}K^{r}u=F_{r+3}\) 得到 \(\tau(\eta)=\prod_i F_{r_i+3}\)（见 \cite{FQFibonacciLucasChap11Matrices}）。
\end{proof}

\paragraph{循环子图与独立集解释：迹矩的图计数展开}
固定整数 \(m\ge 1\)。
记 \(C_m\) 为顶点集 \(\ZZ/m\ZZ\) 上的循环图（当 \(m=1\) 允许自环、当 \(m=2\) 允许重边），边集
\[
E(C_m):=\{(t,t+1):t\in\ZZ/m\ZZ\}.
\]
对任意 \(S\subseteq E(C_m)\)，记 \(C_m[S]\) 为同顶点集的生成子图。
令 \(\Ind(G)\) 表示图 \(G\) 的独立集全体（含空集），并记
\[
I(G):=|\Ind(G)|.
\]

\begin{conclusion}[全阶迹矩的循环子图独立集展开]\label{con:xi-terminal-replica-softcore-trace-moment-cycle-subgraph-independent-sets}
对任意整数 \(q\ge 2\) 与 \(m\ge 1\)，成立
\[
\boxed{\;
\Tr\!\bigl((T^{(q)})^m\bigr)
=\frac1{2^m}\sum_{S\subseteq E(C_m)} I\!\bigl(C_m[S]\bigr)^{\,q}\; } .
\]
\end{conclusion}
\begin{proof}
按迹的路径展开，
\[
\Tr\!\bigl((T^{(q)})^m\bigr)
=\sum_{u_1,\dots,u_m\in U_q}\ \prod_{t\in\ZZ/m\ZZ}T^{(q)}_{u_t,u_{t+1}},
\qquad (u_{m+1}:=u_1),
\]
而 \(T^{(q)}_{u,v}=\frac12\bigl(1+\mathbf{1}_{\{u\wedge v=0\}}\bigr)\)（由定义）给出
\[
\Tr\!\bigl((T^{(q)})^m\bigr)
=2^{-m}\sum_{S\subseteq E(C_m)}\ \sum_{u_1,\dots,u_m\in U_q}\ \prod_{(t,t+1)\in S}\mathbf{1}_{\{u_t\wedge u_{t+1}=0\}}.
\]
对固定 \(S\)，将 \(u_t=(u_t^{(1)},\dots,u_t^{(q)})\) 按坐标分解。
由于
\[
\mathbf{1}_{\{u_t\wedge u_{t+1}=0\}}=\prod_{j=1}^{q}\mathbf{1}_{\{u_t^{(j)}u_{t+1}^{(j)}=0\}},
\]
内层求和按坐标完全因子化为
\[
\sum_{u_1,\dots,u_m\in U_q}\ \prod_{(t,t+1)\in S}\mathbf{1}_{\{u_t\wedge u_{t+1}=0\}}
=\left(\sum_{x\in\{0,1\}^{\ZZ/m\ZZ}}\ \prod_{(t,t+1)\in S}\mathbf{1}_{\{x_t x_{t+1}=0\}}\right)^{q}.
\]
而对二元赋值 \(x\)，令 \(A(x):=\{t\in\ZZ/m\ZZ:\ x_t=1\}\)。
乘积约束 \(\mathbf{1}_{\{x_t x_{t+1}=0\}}\)（对每条 \((t,t+1)\in S\)）等价于 \(A(x)\) 不含任何 \(C_m[S]\) 的边，
即 \(A(x)\in\Ind(C_m[S])\)。
因此括号内的计数恰为 \(I(C_m[S])\)，从而内层求和为 \(I(C_m[S])^{q}\)。
代回即得所示恒等式。证毕。
\end{proof}

\begin{conclusion}[真子图通道的 Fibonacci 乘法分解与分母刚性]\label{con:xi-terminal-replica-softcore-cycle-subgraph-fibonacci-factorization-denominator-rigidity}
固定整数 \(m\ge 1\)。对任意真子边集 \(S\subsetneq E(C_m)\)，图 \(C_m[S]\) 为森林，且每个连通分量均为路径。
设其非平凡路径分量的边数为 \(\ell_1,\dots,\ell_r\ge 1\)，并记孤立点数为
\[
s:=m-\sum_{i=1}^{r}(\ell_i+1).
\]
则有严格分解
\[
\boxed{\;
I\!\bigl(C_m[S]\bigr)=2^{s}\prod_{i=1}^{r}F_{\ell_i+3}\;} .
\]
从而对任意 \(q\ge 2\)，成立整性刚性
\[
2^m\Tr\!\bigl((T^{(q)})^m\bigr)\in\ZZ,
\qquad
2^mS_m(q)\in\ZZ.
\]
\end{conclusion}
\begin{proof}[证明要点]
当 \(S\subsetneq E(C_m)\) 时，\(C_m[S]\) 不可能包含长度为 \(m\) 的整环，故无环且最大度不超过 \(2\)，于是每个连通分量为路径或孤立点。
设 \(a_n:=I(P_n)\)（\(P_n\) 为 \(n\) 个顶点路径），按端点是否入选有递推
\(
a_n=a_{n-1}+a_{n-2}
\)
并由 \(a_1=2,a_2=3\) 得 \(a_n=F_{n+2}\)。
边数为 \(\ell\) 的路径分量有 \(\ell+1\) 个顶点，贡献 \(F_{\ell+3}\)；孤立点贡献 \(2\)。
分量独立相乘即得所示分解。
\par
再由结论 \ref{con:xi-terminal-replica-softcore-trace-moment-cycle-subgraph-independent-sets} 与
结论 \ref{con:xi-terminal-replica-softcore-exceptional-power-sum-cycle-sector-replacement}，
\[
2^m\Tr\!\bigl((T^{(q)})^m\bigr)=\sum_{S\subseteq E(C_m)}I(C_m[S])^{q},
\]
\[
2^mS_m(q)=\Omega_m(q)+\sum_{S\subsetneq E(C_m)}I(C_m[S])^{q},
\]
且 \(\Omega_m(q)\in\ZZ\)（结论 \ref{con:xi-terminal-replica-softcore-lucas-pell-omega-unified}），故整性成立。证毕。
\end{proof}

\begin{remark}[与二元 word/necklace 展开的同一性]\label{rem:xi-terminal-replica-softcore-trace-moment-cycle-subgraph-word-identification}
将边子集 \(S\subseteq E(C_m)\) 编码为二元词 \(\omega=\omega_1\cdots\omega_m\in\{0,1\}^m\)：
\[
\omega_t:=\mathbf{1}_{\{(t,t+1)\notin S\}}\quad(t\in\ZZ/m\ZZ).
\]
则结论 \ref{con:xi-terminal-replica-softcore-binary-necklace-trace-formula} 的词迹值 \(\tau(\omega)\) 满足
\(\tau(\omega)=I(C_m[S])\)，
从而结论 \ref{con:xi-terminal-replica-softcore-trace-moment-cycle-subgraph-independent-sets} 与其 word 求和式严格一致。
\end{remark}

\begin{conclusion}[例外幂和的循环子图展开与对称化替换项]\label{con:xi-terminal-replica-softcore-exceptional-power-sum-cycle-sector-replacement}
固定 \(q\ge 2\)。记
\[
S_m(q):=\sum_{i=0}^{q}\bigl(\lambda_i^{(q)}\bigr)^m,\qquad m\ge1.
\]
令
\[
M:=\begin{pmatrix}1&1\\[2pt]1&0\end{pmatrix},\qquad
\Omega_m(q):=\Tr\!\bigl(\mathrm{Sym}^q(M^m)\bigr).
\]
则对一切 \(m\ge 1\) 有
\[
\boxed{\;
2^mS_m(q)=\Omega_m(q)+\sum_{S\subsetneq E(C_m)} I\!\bigl(C_m[S]\bigr)^{\,q}\; } ,
\]
等价地，
\[
\boxed{\;
S_m(q)=\frac1{2^m}\left(\sum_{S\subsetneq E(C_m)} I\!\bigl(C_m[S]\bigr)^{\,q}+\Omega_m(q)\right)\; }.
\]
\end{conclusion}
\begin{proof}
设 \(\mathrm{Inv}\subset\RR^{U_q}\) 为 \(S_q\)-不变子空间，\(P\) 为到 \(\mathrm{Inv}\) 的正交投影。由于 \(T^{(q)}\) 与坐标置换作用可交换，
\[
S_m(q)=\Tr\!\Bigl(\bigl(T^{(q)}|_{\mathrm{Inv}}\bigr)^m\Bigr)=\Tr\!\bigl(P(T^{(q)})^m\bigr).
\]
把 \((T^{(q)})^m\) 按 \(T^{(q)}=\frac12(J+D^{(q)})\) 展开为 \(\varepsilon\in\{0,1\}^m\) 的和。
若 \(\varepsilon\neq(1,\dots,1)\)，对应字词 \(W_\varepsilon\) 至少含一个 \(J\)。
由 \(J=\mathbf 1\mathbf 1^{\mathsf T}\) 及 \(J,D^{(q)}\) 与置换作用可交换，任意含 \(J\) 的字词像空间包含于 \(\mathrm{Inv}\)，故
\[
\Tr(PW_\varepsilon)=\Tr(W_\varepsilon).
\]
这些项与全迹一致，按结论
\ref{con:xi-terminal-replica-softcore-trace-moment-cycle-subgraph-independent-sets}
给出
\[
\sum_{S\subsetneq E(C_m)}I\!\bigl(C_m[S]\bigr)^q.
\]
唯一剩余项为 \(\varepsilon=(1,\dots,1)\)，即 \((D^{(q)})^m\)。
在按位张量识别 \(\RR^{U_q}\cong(\RR^2)^{\otimes q}\) 下，
\[
D^{(q)}=M^{\otimes q},\qquad (D^{(q)})^m=(M^m)^{\otimes q}.
\]
投影 \(P\) 即对称化算子，其像同构于 \(\mathrm{Sym}^q(\RR^2)\)，从而
\[
\Tr\!\bigl(P(D^{(q)})^m\bigr)=\Tr\!\bigl(\mathrm{Sym}^q(M^m)\bigr)=\Omega_m(q).
\]
合并全部项即得结论。
\end{proof}

\begin{conclusion}[边伯努利形变下的随机环图独立集矩表示]\label{con:xi-terminal-replica-softcore-trace-moment-bernoulli-deformation-random-cycle}
取参数 \(p\in[0,1]\)，定义形变核
\[
T^{(q)}_{p}(u,v):=(1-p)+p\,\mathbf{1}_{\{u\wedge v=0\}},
\qquad u,v\in U_q.
\]
令随机子图 \(G_{m,p}\) 由 \(C_m\) 的每条边独立保留（概率 \(p\)）得到，并记
\(
Z(G_{m,p}):=I(G_{m,p})
\)。
则对任意整数 \(m\ge 1\) 与 \(q\ge 2\)，有严格恒等式
\[
\boxed{\;
\Tr\!\bigl((T^{(q)}_{p})^m\bigr)
=\mathbb E\!\bigl[Z(G_{m,p})^{q}\bigr]
=\sum_{S\subseteq E(C_m)}p^{|S|}(1-p)^{m-|S|}\,I(C_m[S])^{q}\; } .
\]
当 \(p=\tfrac12\) 时退化为原核 \(T^{(q)}\) 的迹矩公式。
\end{conclusion}
\begin{proof}[证明要点]
沿迹展开写
\[
\Tr\!\bigl((T^{(q)}_{p})^m\bigr)
=\sum_{u_1,\dots,u_m\in U_q}\prod_{t\in\ZZ/m\ZZ}\Bigl((1-p)+p\,\delta_t\Bigr),
\qquad
\delta_t:=\mathbf 1_{\{u_t\wedge u_{t+1}=0\}}.
\]
对每个因子按是否取 \(\delta_t\) 展开，即得到
\[
\prod_t\bigl((1-p)+p\,\delta_t\bigr)
=\sum_{S\subseteq E(C_m)}p^{|S|}(1-p)^{m-|S|}\prod_{(t,t+1)\in S}\delta_t.
\]
对固定 \(S\)，其内层 \((u_t)\)-求和与
结论 \ref{con:xi-terminal-replica-softcore-trace-moment-cycle-subgraph-independent-sets} 完全同型，
给出 \(I(C_m[S])^{q}\)；代回即得第一与第三式。
第二式仅是把权重 \(p^{|S|}(1-p)^{m-|S|}\) 解释为边独立保留模型的概率。证毕。
\end{proof}

\begin{conclusion}[固定 \(m\) 下的独立集谱反演与最小递推]\label{con:xi-terminal-replica-softcore-cycle-subgraph-spectrum-inversion-minimal-recurrence}
固定整数 \(m\ge 1\)，定义取值集合
\[
\mathcal A_m:=\{I(C_m[S]):S\subseteq E(C_m)\}\subset\ZZ_{>0}.
\]
设其不同元素为 \(a_1<\cdots<a_r\)，并记重数
\[
\mu_j:=\bigl|\{S\subseteq E(C_m):I(C_m[S])=a_j\}\bigr|.
\]
则对全部 \(q\ge 0\) 有
\[
\boxed{\;
2^m\Tr\!\bigl((T^{(q)})^m\bigr)=\sum_{j=1}^{r}\mu_j\,a_j^{q}\; } .
\]
因此 \(q\mapsto \Tr((T^{(q)})^m)\) 满足阶数 \(r\) 的齐次线性递推，其特征多项式可取
\[
P_m(x):=\prod_{j=1}^{r}(x-a_j)\in\ZZ[x].
\]
并且序列矩 \(\{2^m\Tr((T^{(q)})^m)\}_{q\ge 0}\) 对 \(\{(a_j,\mu_j)\}_{j=1}^{r}\) 具有限矩确定性：当 \(r\) 已知时，前 \(2r\) 项可唯一恢复全部 \((a_j,\mu_j)\)。
\end{conclusion}
\begin{proof}[证明要点]
首式由结论 \ref{con:xi-terminal-replica-softcore-trace-moment-cycle-subgraph-independent-sets}
对 \(\{S\subseteq E(C_m)\}\) 按同值 \(I(C_m[S])\) 分组即得。
有限指数和 \(\sum_j\mu_j a_j^q\) 被 \(\prod_j(x-a_j)\) 湮灭，故得递推。
\par
对反演性，记矩序列 \(M_q:=\sum_j\mu_j a_j^q\)。
其生成函数
\[
\sum_{q\ge 0}M_q z^q=\sum_{j=1}^{r}\frac{\mu_j}{1-a_j z}
\]
为分母次数 \(r\) 的有理函数。
由前 \(2r\) 项矩可唯一确定该有理函数（Pad\'e/Prony 型重构），
再由部分分式分解唯一性恢复 \((a_j,\mu_j)\)。证毕。
\end{proof}

\begin{conclusion}[低阶例外幂和闭式的统一回收：\(m=2\) 与 \(m=6\)]\label{con:xi-terminal-replica-softcore-exceptional-power-sum-m2-m6-recovered}
对任意 \(q\ge 2\)，有
\[
\boxed{\;
S_2(q)=\frac14\Bigl(4^{q}+2\cdot 3^{q}+\Omega_2(q)\Bigr)
=\frac14\Bigl(4^{q}+2\cdot 3^{q}+F_{2q+2}\Bigr)\; } ,
\]
以及
\[
\boxed{\;
\begin{aligned}
S_6(q)=\frac1{64}\Big(
&64^{q}+6\cdot48^{q}+6\cdot40^{q}+9\cdot36^{q}+6\cdot32^{q}+12\cdot30^{q}\\
&+2\cdot27^{q}+6\cdot26^{q}+3\cdot25^{q}+6\cdot24^{q}+6\cdot21^{q}+\Omega_6(q)
\Big)\; .
\end{aligned}
}
\]
\end{conclusion}
\begin{proof}[证明要点]
由结论 \ref{con:xi-terminal-replica-softcore-exceptional-power-sum-cycle-sector-replacement} 取 \(m=2\)，
真子边集仅有空集与两条单边：对应 \(I=4\) 与 \(I=3\)（重数 \(2\)），代入得
\(
S_2(q)=2^{-2}(4^q+2\cdot 3^q+\Omega_2(q))
\)；
再由 \(\Omega_2(q)=F_{2q+2}\)（结论 \ref{con:xi-terminal-replica-softcore-lucas-pell-omega-unified}）即得第一式。
\par
第二式与结论 \ref{con:xi-terminal-replica-softcore-exceptional-power-sum-m6} 一致，
亦可由 \(m=6\) 的子边集按连通型分类并代回同一替换公式得到。证毕。
\end{proof}

\begin{conclusion}[大复制阶的统一指数分层：次主比率 \(\tfrac34\)]\label{con:xi-terminal-replica-softcore-large-q-two-layer-asymptotics}
固定整数 \(m\ge 3\)。当 \(q\to\infty\) 时，
\[
\Tr\!\bigl((T^{(q)})^m\bigr)
=2^{m(q-1)}
\left(
1+m\Bigl(\frac34\Bigr)^q+O_m\!\Bigl(\Bigl(\frac58\Bigr)^q\Bigr)
\right),
\]
且同一展开对例外幂和成立：
\[
S_m(q)
=2^{m(q-1)}
\left(
1+m\Bigl(\frac34\Bigr)^q+O_m\!\Bigl(\Bigl(\frac58\Bigr)^q\Bigr)
\right).
\]
\end{conclusion}
\begin{proof}[证明要点]
由结论 \ref{con:xi-terminal-replica-softcore-trace-moment-cycle-subgraph-independent-sets}，
\[
2^m\Tr\!\bigl((T^{(q)})^m\bigr)=\sum_{S\subseteq E(C_m)}I(C_m[S])^{q}.
\]
其中 \(S=\varnothing\) 唯一给出最大值 \(I=2^m\)。
当 \(|S|=1\) 时 \(I=3\cdot 2^{m-2}\)，重数恰为 \(m\)。
若 \(|S|\ge 2\)，则要么含相邻两边（出现长度 \(2\) 路径分量），要么至少含两条不交边；
两种情形统一给出
\[
I(C_m[S])\le 5\cdot 2^{m-3}.
\]
故
\[
2^m\Tr\!\bigl((T^{(q)})^m\bigr)
=(2^m)^q+m(3\cdot 2^{m-2})^q+O_m\!\bigl((5\cdot 2^{m-3})^q\bigr),
\]
化简得第一式。
\par
再由结论 \ref{con:xi-terminal-replica-softcore-exceptional-power-sum-cycle-sector-replacement}，
\[
S_m(q)=\Tr\!\bigl((T^{(q)})^m\bigr)-2^{-m}\bigl(L_m^q-\Omega_m(q)\bigr).
\]
且 \(|L_m^q-\Omega_m(q)|\ll_m \varphi^{mq}\)，故差项为
\[
2^{-m}|L_m^q-\Omega_m(q)|
=2^{m(q-1)}\,O_m\!\Bigl(\Bigl(\frac{\varphi}{2}\Bigr)^{mq}\Bigr)
=2^{m(q-1)}\,O_m\!\Bigl(\Bigl(\frac58\Bigr)^q\Bigr),
\]
其中最后一步用到 \(m\ge 3\) 时 \((\varphi/2)^m<5/8\)。
于是第二式成立。证毕。
\end{proof}

\begin{conclusion}[时间纤维分裂的系数定律：Fibonacci 乘积基底与圆周重复指数]\label{con:xi-terminal-replica-softcore-cycle-subgraph-cyclic-composition-coeff-law}
固定整数 \(m\ge 1\)。
对任意真子图 \(C_m[S]\)（即 \(S\subsetneq E(C_m)\)），图 \(C_m[S]\) 必为若干条路径的离散并。
记其路径分量的顶点数依圆周顺序给出圆周分拆（cyclic composition）
\[
\boldsymbol n=(n_1,\dots,n_r),\qquad n_i\in\ZZ_{\ge 1},\ \sum_{i=1}^{r}n_i=m,
\]
并按圆周旋转等价。
令 Fibonacci 数列满足 \(F_0=0,F_1=1\) 与递推 \(F_{n+2}=F_{n+1}+F_n\)，并定义
\[
B(\boldsymbol n):=\prod_{i=1}^{r}F_{n_i+2}\in\ZZ_{>0}.
\]
再令 \(\mathrm{rep}(\boldsymbol n)\) 表示 \(\boldsymbol n\) 作为圆周序列的最大重复次数：存在且仅存在最小块 \(\boldsymbol n_0\) 使得（按圆周意义）
\(
\boldsymbol n=\boldsymbol n_0^{\ \mathrm{rep}(\boldsymbol n)}
\)。
则：
\begin{enumerate}
  \item \textbf{（指数基底）}若 \(S\subsetneq E(C_m)\) 的分量圆周分拆为 \(\boldsymbol n\)，则
  \[
  \boxed{\ I\!\bigl(C_m[S]\bigr)=B(\boldsymbol n)\ }.
  \]
  \item \textbf{（多重性）}对应于同一圆周分拆类 \(\boldsymbol n\) 的边子集个数严格等于
  \[
  \boxed{\ \kappa(\boldsymbol n)=\frac{m}{\mathrm{rep}(\boldsymbol n)}\ } .
  \]
  \item \textbf{（分组重写）}因此结论 \ref{con:xi-terminal-replica-softcore-exceptional-power-sum-cycle-sector-replacement} 的循环扇区替换可按圆周分拆分组为
  \[
  \boxed{\;
  2^m S_m(q)=\Omega_m(q)+\sum_{\boldsymbol n\in\mathrm{CycComp}(m)}\kappa(\boldsymbol n)\,B(\boldsymbol n)^{q}\; },
  \qquad q\ge 2,
  \]
  其中 \(\mathrm{CycComp}(m)\) 表示全部圆周分拆类（按圆周旋转取商）。
\end{enumerate}
\end{conclusion}
\begin{proof}[证明要点]
由于 \(S\subsetneq E(C_m)\)，图 \(C_m[S]\) 无环且每个顶点度数不超过 \(2\)，故其连通分量必为路径。
设某一分量为顶点数 \(n\) 的路径 \(P_n\)，记 \(a_n:=I(P_n)\)。
按是否选取端点拆分可得递推 \(a_n=a_{n-1}+a_{n-2}\)（\(n\ge 2\)），且 \(a_0=1,a_1=2\)，故 \(a_n=F_{n+2}\)。
从而 \(I(C_m[S])\) 由分量相乘得到 \(B(\boldsymbol n)\)。
\par
对多重性，令 \(Z:=E(C_m)\setminus S\neq\varnothing\) 为缺边集合，并按圆周顺序列出其元素。
相邻缺边之间的圆周距离（以顶点数计）给出 \(\boldsymbol n\)。
循环平移群 \(\ZZ/m\ZZ\) 作用于边子集并保持 \(\boldsymbol n\) 的圆周分拆类；
其稳定子由把缺边集合整体平移不变的旋转组成，而稳定子阶恰等于重复次数 \(\mathrm{rep}(\boldsymbol n)\)。
由轨道—稳定子定理得到轨道大小 \(m/\mathrm{rep}(\boldsymbol n)\)，即为出现次数 \(\kappa(\boldsymbol n)\)。
最后，将结论 \ref{con:xi-terminal-replica-softcore-exceptional-power-sum-cycle-sector-replacement} 中
\(\sum_{S\subsetneq E(C_m)}I(C_m[S])^{q}\)
按圆周分拆类分组并代入前两条即得重写式。
可复现核验脚本：\texttt{scripts/exp\_replica\_softcore\_cycle\_subgraph\_cyccomp\_audit.py}。证毕。
\end{proof}

\begin{conclusion}[谱泄漏局部化：非对称子空间的迹仅由全循环扇区贡献]\label{con:xi-terminal-replica-softcore-spectral-leakage-localization-full-cycle-sector}
令 \(\mathcal H_q:=(\CC^2)^{\otimes q}\)，并记 \(S_q\)-不变子空间 \(\mathcal H_q^{S_q}\cong\mathrm{Sym}^q(\CC^2)\)（结论 \ref{con:xi-terminal-replica-softcore-exceptional-spectrum-sympower-model}）。
则对任意整数 \(m\ge 1\) 与 \(q\ge 2\)，有
\[
\boxed{\;
\Tr\!\left(\left.(T^{(q)})^m\right|_{(\mathcal H_q^{S_q})^{\perp}}\right)
=2^{-m}\Bigl(I(C_m)^{q}-\Omega_m(q)\Bigr)\;}
\]
其中正交补取于 \(\mathcal H_q\) 的自然内积。
从而非对称方向的全部迹贡献被严格局部化到全循环扇区 \(S=E(C_m)\)（亦即字词 \(K^m\)）所触发的替换差项上。
\end{conclusion}
\begin{proof}[证明要点]
由 \(J=uu^{\top}\)（\(u=(1,1)^{\top}\)），任一含 \(J\) 的字词乘积 \(M(w)\) 皆为秩 \(1\) 矩阵；
因此 \(M(w)^{\otimes q}\) 的像空间包含于 \(\mathrm{span}\{a(w)^{\otimes q}\}\subset \mathcal H_q^{S_q}\)。
换言之，除 \(w=K^m\) 外的全部扇区项在 \((\mathcal H_q^{S_q})^{\perp}\) 上限制为零算子。
因此
\[
\Tr\!\left(\left.(T^{(q)})^m\right|_{(\mathcal H_q^{S_q})^{\perp}}\right)
=2^{-m}\left(\Tr\!\bigl((K^{\otimes q})^m\bigr)-\Tr\!\bigl(\mathrm{Sym}^q(K^m)\bigr)\right)
=2^{-m}\bigl(L_m^{q}-\Omega_m(q)\bigr),
\]
其中 \(\Tr((K^{\otimes q})^m)=\Tr((K^m)^{\otimes q})=(\Tr K^m)^q=L_m^q\)。
最后用 \(I(C_m)=L_m\) 得到所示形式。证毕。
\end{proof}

\endinput


\paragraph{环边积分分解与 Vieta 系数闭式：例外幂和的可枚举结构化重写}
沿用本节记号：\(q\ge 2\)、\(U_q=\{0,1\}^q\)、\(T^{(q)}\) 为软核矩阵，
\(\{\lambda_i^{(q)}\}_{i=0}^{q}\) 为 \(S_q\)-对称子空间上的 \(q+1\) 个单根例外特征值，
\[
S_m(q):=\sum_{i=0}^{q}\bigl(\lambda_i^{(q)}\bigr)^m,\qquad m\ge 1.
\]

\begin{conclusion}[例外幂和的环边路径积分分解与 Fibonacci 组件因子化]\label{con:xi-terminal-replica-softcore-cycle-edge-integral-decomposition-fibonacci-components}
固定 \(m\ge 1\)。记环的边位集合
\[
E_m:=\{(t,t+1):t\in\ZZ/m\ZZ\},
\]
并对任意 \(S\subseteq E_m\) 定义
\[
N_m(S):=\#\left\{x\in\{0,1\}^{\ZZ/m\ZZ}:\
x_t x_{t+1}=0\ \text{对全部 }(t,t+1)\in S\right\}.
\]
再定义
\[
\Omega_m(q):=\sum_{j=0}^{q}(-1)^{jm}\varphi^{m(q-2j)}.
\]
则对一切 \(q\ge 2\)、\(m\ge 1\)，有
\begin{equation}\label{eq:xi-cycle-edge-integral-trace}
\Tr\!\bigl((T^{(q)})^m\bigr)=2^{-m}\sum_{S\subseteq E_m}N_m(S)^q,
\end{equation}
\begin{equation}\label{eq:xi-cycle-edge-integral-exceptional}
S_m(q)=2^{-m}\!\left(\Omega_m(q)+\sum_{S\subsetneq E_m}N_m(S)^q\right).
\end{equation}
此外，当 \(S\subsetneq E_m\) 时，关联无向图分解为若干路径分量；若其分量顶点数为
\(
n_1,\dots,n_r
\)
（\(\sum_{\ell=1}^{r}n_\ell=m\)），则
\begin{equation}\label{eq:xi-cycle-edge-integral-fibonacci-factor}
N_m(S)=\prod_{\ell=1}^{r}F_{n_\ell+2}.
\end{equation}
\end{conclusion}

\begin{proof}
\textbf{第一步：迹矩展开与坐标解耦。}
由
\(
T^{(q)}_{u,v}=\frac12(1+\mathbf 1_{\{u\wedge v=0\}})
\)
及环路径迹展开，
\[
\Tr\!\bigl((T^{(q)})^m\bigr)
=\sum_{u_0,\dots,u_{m-1}\in U_q}
\prod_{t\in\ZZ/m\ZZ}T^{(q)}_{u_t,u_{t+1}},\qquad u_m:=u_0.
\]
对 \(m\) 个因子完全展开得
\[
\Tr\!\bigl((T^{(q)})^m\bigr)
=2^{-m}\sum_{S\subseteq E_m}
\#\left\{(u_t)_t\in U_q^m:\ u_t\wedge u_{t+1}=0,\ \forall (t,t+1)\in S\right\}.
\]
固定 \(S\) 后，约束按坐标 \(r=1,\dots,q\) 完全解耦；每个坐标都是同一二元环约束计数 \(N_m(S)\)。
故内层计数等于 \(N_m(S)^q\)，即得 \eqref{eq:xi-cycle-edge-integral-trace}。

\textbf{第二步：例外幂和替换。}
由全谱分解（例外谱与固定谱并置）：
\[
\Tr\!\bigl((T^{(q)})^m\bigr)
=S_m(q)+\sum_{j=0}^{q}\left(\binom{q}{j}-1\right)
\left(\frac{(-1)^j\varphi^{q-2j}}{2}\right)^m.
\]
化简第二项：
\[
2^{-m}\!\left(
\sum_{j=0}^{q}\binom{q}{j}(-1)^{jm}\varphi^{m(q-2j)}
-\sum_{j=0}^{q}(-1)^{jm}\varphi^{m(q-2j)}
\right)
=2^{-m}\!\left(L_m^q-\Omega_m(q)\right).
\]
其中第一和式由二项式定理给出
\[
\sum_{j=0}^{q}\binom{q}{j}(\varphi^m)^{q-j}\bigl((-1)^m\varphi^{-m}\bigr)^j
=\bigl(\varphi^m+(-1)^m\varphi^{-m}\bigr)^q=L_m^q.
\]
故
\[
\Tr\!\bigl((T^{(q)})^m\bigr)=S_m(q)+2^{-m}\bigl(L_m^q-\Omega_m(q)\bigr).
\]
再与 \eqref{eq:xi-cycle-edge-integral-trace} 联立，并用 \(N_m(E_m)=L_m\)（环 \(C_m\) 的独立集数），即得
\eqref{eq:xi-cycle-edge-integral-exceptional}。

\textbf{第三步：真子图的 Fibonacci 乘法分解。}
当 \(S\subsetneq E_m\) 时，去掉至少一个环边位，关联图无环且各分量为路径。
对 \(n\) 顶点路径 \(P_n\)，独立集数 \(a_n\) 满足
\(
a_n=a_{n-1}+a_{n-2}
\)、\(a_0=1,a_1=2\)，故 \(a_n=F_{n+2}\)。
各分量独立相乘得到 \eqref{eq:xi-cycle-edge-integral-fibonacci-factor}。证毕。
\end{proof}

\begin{corollary}[二次幂和的完全闭式与 \texorpdfstring{$F_{2q+2}$}{F\_(2q+2)} 出现]\label{cor:xi-terminal-replica-softcore-exceptional-power-sum-m2-fibonacci-closed}
\leanverified{paper\_xi\_terminal\_replica\_softcore\_exceptional\_power\_sum\_m2\_fibonacci\_closed}
对任意 \(q\ge 2\)，有
\begin{equation}\label{eq:xi-exceptional-s2-complete-closed-form}
S_2(q)=\frac14\left(4^q+2\cdot 3^q+F_{2q+2}\right).
\end{equation}
\end{corollary}

\begin{proof}
在结论 \ref{con:xi-terminal-replica-softcore-cycle-edge-integral-decomposition-fibonacci-components} 中取 \(m=2\)。
此时 \(E_2\) 含两个边位；真子集为 \(\varnothing\) 与两个单边位子集，故
\[
N_2(\varnothing)=4,\qquad N_2(S)=3\ \ (|S|=1).
\]
又
\[
\Omega_2(q)=\sum_{j=0}^{q}\varphi^{2(q-2j)}
=\frac{\varphi^{2(q+1)}-\varphi^{-2(q+1)}}{\sqrt5}
=F_{2q+2}.
\]
代入 \eqref{eq:xi-cycle-edge-integral-exceptional} 即得 \eqref{eq:xi-exceptional-s2-complete-closed-form}。证毕。
\end{proof}

\begin{conclusion}[例外特征多项式二阶 Vieta 系数闭式]\label{con:xi-terminal-replica-softcore-exceptional-vieta-e2-closed-form}
记
\[
\chi_{\mathrm{exc}}^{(q)}(x)
:=\prod_{i=0}^{q}\bigl(x-\lambda_i^{(q)}\bigr)
=x^{q+1}-e_1^{(q)}x^q+e_2^{(q)}x^{q-1}-\cdots+(-1)^{q+1}e_{q+1}^{(q)}.
\]
则对任意 \(q\ge 2\)，
\begin{equation}\label{eq:xi-exceptional-vieta-e2-closed}
e_2^{(q)}
=2^{q-2}F_{q+1}
-\frac14\Bigl(3^q+F_qF_{q+1}\Bigr).
\end{equation}
\end{conclusion}

\begin{proof}
由 Vieta 关系，
\[
e_1^{(q)}=\sum_{i=0}^{q}\lambda_i^{(q)}=S_1(q),\qquad
e_2^{(q)}=\sum_{0\le i<j\le q}\lambda_i^{(q)}\lambda_j^{(q)}
=\frac{S_1(q)^2-S_2(q)}{2}.
\]
由结论
\ref{con:xi-terminal-replica-softcore-exceptional-power-sum-m1-integer-collapse}
与推论 \ref{cor:xi-terminal-replica-softcore-exceptional-power-sum-m2-fibonacci-closed}，
\[
S_1(q)=2^{q-1}+\frac12F_{q+1},\qquad
S_2(q)=\frac14\bigl(4^q+2\cdot 3^q+F_{2q+2}\bigr).
\]
代入并用恒等式
\(
F_{2q+2}=F_{q+1}^2+2F_qF_{q+1}
\)
化简，即得 \eqref{eq:xi-exceptional-vieta-e2-closed}。证毕。
\end{proof}

\begin{conclusion}[六次幂和全闭式与 Fibonacci 组件谱系]\label{con:xi-terminal-replica-softcore-exceptional-power-sum-m6-fibonacci-component-spectrum}
对任意 \(q\ge 2\)，
\begin{equation}\label{eq:xi-exceptional-s6-fibonacci-spectrum-closed}
\begin{aligned}
S_6(q)=\frac1{64}\Big(
&64^q+6\cdot48^q+6\cdot40^q+9\cdot36^q+6\cdot32^q+12\cdot30^q\\
&+2\cdot27^q+6\cdot26^q+3\cdot25^q+6\cdot24^q+6\cdot21^q+\Omega_6(q)
\Big),
\end{aligned}
\end{equation}
且 \(\Omega_6\) 满足
\begin{equation}\label{eq:xi-omega6-second-order-recurrence}
\Omega_6(0)=1,\quad
\Omega_6(1)=18,\quad
\Omega_6(q)=18\,\Omega_6(q-1)-\Omega_6(q-2)\ \ (q\ge 2).
\end{equation}
\end{conclusion}

\begin{proof}
由 \eqref{eq:xi-cycle-edge-integral-exceptional}（\(m=6\)）：
\[
S_6(q)=2^{-6}\!\left(\Omega_6(q)+\sum_{S\subsetneq E_6}N_6(S)^q\right).
\]
按缺边位数 \(r:=|E_6\setminus S|\) 分类，结合
\eqref{eq:xi-cycle-edge-integral-fibonacci-factor} 逐类得到 \(N_6(S)\) 与重数：
\begin{enumerate}
  \item \(r=6\)：\(N=64\)，重数 \(1\)；
  \item \(r=5\)：\(N=48\)，重数 \(6\)；
  \item \(r=4\)：\(N=40\)（重数 \(6\)），\(N=36\)（重数 \(9\)）；
  \item \(r=3\)：\(N=32\)（重数 \(6\)），\(N=30\)（重数 \(12\)），\(N=27\)（重数 \(2\)）；
  \item \(r=2\)：\(N=26\)（重数 \(6\)），\(N=24\)（重数 \(6\)），\(N=25\)（重数 \(3\)）；
  \item \(r=1\)：\(N=21\)，重数 \(6\)。
\end{enumerate}
总重数 \(1+6+15+20+15+6=63=2^6-1\)。
代回即得 \eqref{eq:xi-exceptional-s6-fibonacci-spectrum-closed}。

又
\[
\Omega_6(q)=\sum_{j=0}^{q}\varphi^{6(q-2j)}
\]
是特征根 \(\varphi^6,\varphi^{-6}\) 的二指数线性组合，故满足
\[
\Omega_6(q)-(\varphi^6+\varphi^{-6})\Omega_6(q-1)+\Omega_6(q-2)=0.
\]
而 \(\varphi^6+\varphi^{-6}=L_6=18\)，并由定义得初值 \(1,18\)，即 \eqref{eq:xi-omega6-second-order-recurrence}。证毕。
\end{proof}

\begin{proposition}[固定 \texorpdfstring{$m$}{m} 的 \texorpdfstring{$q$}{q}-方向有限阶线性递推]\label{prop:xi-terminal-replica-softcore-fixed-m-q-recurrence-PAQ}
固定 \(m\ge 1\)。定义
\[
\mathcal A_m:=\{N_m(S):S\subsetneq E_m\}\subset\ZZ_{>0},
\]
\[
P_m(t):=\prod_{a\in\mathcal A_m}(t-a),
\qquad
Q_m(t):=t^2-L_m t+(-1)^m.
\]
记移位算子 \(E\) 作用为 \((Ef)(q)=f(q+1)\)。则
\begin{equation}\label{eq:xi-fixed-m-q-recurrence-PAQ}
P_m(E)\,Q_m(E)\,S_m(\cdot)\equiv 0.
\end{equation}
换言之，\(q\mapsto S_m(q)\) 满足一个由 \(\mathcal A_m\) 与 \(L_m\) 显式决定的整系数齐次递推；
其阶数为 \(|\mathcal A_m|+2\)。
\end{proposition}

\begin{proof}
由 \eqref{eq:xi-cycle-edge-integral-exceptional}，
\[
2^mS_m(q)=\Omega_m(q)+\sum_{S\subsetneq E_m}N_m(S)^q.
\]
将真子集项按同值 \(a\in\mathcal A_m\) 归并，写为
\[
\sum_{S\subsetneq E_m}N_m(S)^q=\sum_{a\in\mathcal A_m}c_m(a)\,a^q,\qquad c_m(a)\in\NN.
\]
每个指数序列 \(a^q\) 都被 \(E-a\) 湮灭，故总和被 \(P_m(E)\) 湮灭。

另一方面，\(\Omega_m(q)\) 的特征根为 \(\varphi^m\) 与 \((-1)^m\varphi^{-m}\)，因此满足
\[
\Omega_m(q+2)-L_m\Omega_m(q+1)+(-1)^m\Omega_m(q)=0,
\]
即 \(Q_m(E)\Omega_m\equiv 0\)。
对
\(
2^mS_m=\Omega_m+\sum_{a}c_m(a)a^q
\)
左作用 \(P_m(E)Q_m(E)\) 即得
\eqref{eq:xi-fixed-m-q-recurrence-PAQ}。证毕。
\end{proof}

\endinput

\paragraph{对称降维字词迹分解与指数通道层级}
本段将 \(S_m(q)=\Tr((A^{(q)})^m)\) 统一改写为 \(\{J,D\}\)-字词迹框架，其中
\(A^{(q)}=\frac12(J+D)\)、\(D=B_q^{\mathsf T}\) 且 \(J\) 为秩一算子。
该重写把 \(\Omega_m\)-通道与 Fibonacci 乘法通道在同一代数层级上分离。

\begin{conclusion}[Fibonacci--Wick 因子化：对称降维中任意混合字的迹闭式]\label{con:xi-terminal-replica-softcore-trace-free-semigroup-lift}
固定整数 \(q\ge 1\)。定义
\[
A^{(q)}_{k,l}:=\frac12\!\left(\binom{q}{l}+\binom{q-k}{l}\right),\qquad 0\le k,l\le q,
\]
\[
\mathbf 1:=(1,\dots,1)^{\mathsf T}\in\RR^{q+1},\qquad
b:=\bigl(\binom{q}{0},\binom{q}{1},\dots,\binom{q}{q}\bigr)^{\mathsf T},
\]
\[
B_q(i,j):=\binom{q-j}{i}\quad(0\le i,j\le q),\qquad
J:=\mathbf 1\,b^{\mathsf T},\qquad
D:=B_q^{\mathsf T}.
\]
则 \(A^{(q)}=\frac12(J+D)\)。对任意 \(m\ge 1\) 与循环字 \(w\in\{J,D\}^m\)：
\begin{enumerate}
  \item 若 \(w=D^m\)，则
  \[
  \Tr(w)=\Tr(B_q^m)=:\Omega_m(q).
  \]
  \item 若 \(w\) 至少含一个 \(J\)，则在循环同伦意义下存在唯一表示
  \[
  w\sim J D^{\ell_1}J D^{\ell_2}\cdots J D^{\ell_r},
  \qquad
  r\ge 1,\ \ell_i\ge 0,\ \sum_{i=1}^r(\ell_i+1)=m,
  \]
  并满足
  \[
  \Tr(w)=\prod_{i=1}^{r}F_{\ell_i+3}^{\,q}.
  \]
\end{enumerate}
\end{conclusion}
\begin{proof}
第一条由 \(D=B_q^{\mathsf T}\) 与 \(\Tr(X^{\mathsf T})=\Tr(X)\) 立即得到。
对第二条，因 \(J=\mathbf 1b^{\mathsf T}\) 为秩一算子，
\[
\Tr\!\bigl(JD^{\ell_1}\cdots JD^{\ell_r}\bigr)=\prod_{i=1}^{r}b^{\mathsf T}D^{\ell_i}\mathbf 1.
\]
又 \(D=B_q^{\mathsf T}\)，故
\[
b^{\mathsf T}D^\ell\mathbf 1=\mathbf 1^{\mathsf T}B_q^\ell b.
\]
利用 \(B_q\simeq\mathrm{Sym}^q(M)\)（\(M=\begin{psmallmatrix}1&1\\[1pt]1&0\end{psmallmatrix}\)）：
\(b\) 对应 \((x+y)^q\)，\(\mathbf 1^{\mathsf T}\) 对应在 \((1,1)\) 处取值。
由
\[
M^\ell:\ (x,y)\longmapsto\bigl(F_{\ell+1}x+F_\ell y,\ F_\ell x+F_{\ell-1}y\bigr),
\]
得
\[
(x+y)^q\xmapsto{\ \mathrm{Sym}^q(M^\ell)\ }\bigl(F_{\ell+2}x+F_{\ell+1}y\bigr)^q.
\]
在 \((1,1)\) 处取值即
\[
\mathbf 1^{\mathsf T}B_q^\ell b=(F_{\ell+2}+F_{\ell+1})^q=F_{\ell+3}^{\,q}.
\]
代回即可得到第二条。证毕。
\end{proof}

\begin{conclusion}[含秩一字母通道的纯乘法迹结构]\label{con:xi-terminal-replica-softcore-word-trace-fibonacci-factorization}
沿结论 \ref{con:xi-terminal-replica-softcore-trace-free-semigroup-lift} 的记号，若
\(w\in\{J,D\}^m\) 含至少一个 \(J\)，并在循环意义下写作
\[
w\sim J D^{\ell_1}J D^{\ell_2}\cdots J D^{\ell_r},
\]
则
\[
\Tr(w)=\prod_{i=1}^{r}F_{\ell_i+3}^{\,q}.
\]
唯一不含 \(J\) 的字词为 \(D^m\)，其迹为 \(\Omega_m(q)=\Tr(B_q^m)\)。
\end{conclusion}
\begin{proof}
直接由结论 \ref{con:xi-terminal-replica-softcore-trace-free-semigroup-lift} 的两条结论合并得到。证毕。
\end{proof}

\begin{conclusion}[例外幂和的全阶字词闭式与递推闭包]\label{con:xi-terminal-replica-softcore-exceptional-power-sum-word-closed-form}
令
\[
S_m(q):=\sum_{i=0}^{q}(\lambda_i^{(q)})^m=\Tr\!\bigl((A^{(q)})^m\bigr),
\]
并沿用 \(\Omega_m(q)\)（结论 \ref{con:xi-terminal-replica-softcore-lucas-pell-omega-unified}）。
则对任意 \(q\ge2,\ m\ge1\) 有
\[
\boxed{\;
S_m(q)=\frac1{2^m}\left(
\Omega_m(q)+
\sum_{\substack{w\in\{J,D\}^m\\ w\neq D^m}}
\Tr(w)
\right)\; }.
\]
并且：
\begin{enumerate}
  \item \(2^mS_m(q)\in\ZZ\)；
  \item 记
  \[
  \mathcal T_m:=\{\Tr(w):w\in\{J,D\}^m,\ w\neq D^m\}\subset\ZZ_{>0},
  \]
  则 \(q\mapsto S_m(q)\) 属于
  \[
  \mathrm{span}_{\QQ}\{\alpha^q:\alpha\in\mathcal T_m\}
  \ \oplus\
  \mathrm{span}_{\QQ}\{\Omega_m(q)\},
  \]
  从而满足首一整数系数线性递推；其湮灭多项式可取
  \[
  \boxed{\;
  \Bigl(x^2-L_mx+(-1)^m\Bigr)\prod_{\alpha\in\mathcal T_m}(x-\alpha)\; }.
  \]
\end{enumerate}
\end{conclusion}
\begin{proof}
由 \(A^{(q)}=\frac12(J+D)\)（结论 \ref{con:xi-terminal-replica-softcore-trace-free-semigroup-lift}），
\[
S_m(q)=2^{-m}\sum_{w\in\{J,D\}^m}\Tr(w).
\]
其中纯 \(D\)-字词贡献 \(\Tr(D^m)=\Omega_m(q)\)，其余字词由
结论 \ref{con:xi-terminal-replica-softcore-word-trace-fibonacci-factorization}
给出为正整数的 \(q\)-次幂，因此得到主公式。
整性由右端整数项直接推出。
再由 \(\Omega_m(q)=L_m\Omega_m(q-1)-(-1)^m\Omega_m(q-2)\)，
以及每个 \(\alpha^q\)（\(\alpha\in\mathcal T_m\)）满足一阶递推 \(x-\alpha\)，
故乘积多项式湮灭该序列。证毕。
\end{proof}

\begin{corollary}[例外谱二次幂和的全闭式]\label{cor:xi-terminal-replica-softcore-exceptional-power-sum-m2-closed-form}
对任意 \(q\ge 1\)，
\[
\boxed{\;
S_2(q)=\frac14\Bigl(4^q+2\cdot3^q+F_{2q+2}\Bigr)\; }.
\]
\leanverified{paper\_xi\_terminal\_replica\_softcore\_exceptional\_power\_sum\_m2\_closed\_form}
\end{corollary}
\begin{proof}
由定义
\(
S_2(q)=\Tr\bigl((A^{(q)})^2\bigr)
\)
且
\[
A^{(q)}_{k,\ell}
=\frac12\Bigl(\binom{q}{\ell}+\binom{q-k}{\ell}\Bigr),
\qquad
0\le k,\ell\le q,
\]
故
\[
S_2(q)=\sum_{k,\ell=0}^{q}A^{(q)}_{k,\ell}A^{(q)}_{\ell,k}
=\frac14\,(T_1+T_2+T_3+T_4),
\]
其中
\[
T_1:=\sum_{k,\ell}\binom{q}{\ell}\binom{q}{k}=4^q,
\]
\[
T_2:=\sum_{k,\ell}\binom{q}{\ell}\binom{q-\ell}{k}
=\sum_{\ell}\binom{q}{\ell}2^{q-\ell}=3^q,
\]
\[
T_3:=\sum_{k,\ell}\binom{q-k}{\ell}\binom{q}{k}
=\sum_{k}\binom{q}{k}2^{q-k}=3^q,
\]
以及
\[
T_4:=\sum_{k,\ell=0}^{q}\binom{q-k}{\ell}\binom{q-\ell}{k}
=\sum_{\substack{k,\ell\ge0\\k+\ell\le q}}
\binom{q-k}{\ell}\binom{q-\ell}{k}.
\]
令 \(c:=q-k-\ell\ge0\)，则
\[
T_4
=\sum_{\substack{k,\ell,c\ge0\\k+\ell+c=q}}
\binom{\ell+c}{\ell}\binom{k+c}{k}.
\]
取生成函数 \(G(z):=\sum_{q\ge0}T_4(q)\,z^q\)，交换求和次序并用
\[
\sum_{k\ge0}\binom{k+c}{k}z^k=(1-z)^{-(c+1)}
\]
得到
\[
G(z)
=\sum_{c\ge0}z^c
\Bigl(\sum_{k\ge0}\binom{k+c}{k}z^k\Bigr)
\Bigl(\sum_{\ell\ge0}\binom{\ell+c}{\ell}z^\ell\Bigr)
=\sum_{c\ge0}\frac{z^c}{(1-z)^{2c+2}}
=\frac1{1-3z+z^2}.
\]
而
\(
\sum_{q\ge0}F_{2q+2}z^q=\frac1{1-3z+z^2}
\)，故 \(T_4=F_{2q+2}\)。
代回即得
\[
S_2(q)=\frac14\Bigl(4^q+2\cdot3^q+F_{2q+2}\Bigr).
\]
证毕。
\end{proof}

\begin{conclusion}[六次幂和的全指数展开]\label{con:xi-terminal-replica-softcore-exceptional-power-sum-m6-word-expanded}
对任意 \(q\ge 2\)，
\[
\boxed{\;
\begin{aligned}
S_6(q)=\frac1{64}\Big(
&64^{q}
+6\cdot 48^{q}
+6\cdot 40^{q}
+9\cdot 36^{q}
+6\cdot 32^{q}
+12\cdot 30^{q}\\
&+2\cdot 27^{q}
+6\cdot 26^{q}
+3\cdot 25^{q}
+6\cdot 24^{q}
+6\cdot 21^{q}
+\Omega_6(q)
\Big)\; .
\end{aligned}
}
\]
\end{conclusion}
\begin{proof}
该式与结论 \ref{con:xi-terminal-replica-softcore-exceptional-power-sum-m6} 一致。
亦可由结论 \ref{con:xi-terminal-replica-softcore-word-trace-fibonacci-factorization}
对长度 \(6\) 的 \(\{J,D\}\)-字词按 \(D\)-段型分组计数，再代入
结论 \ref{con:xi-terminal-replica-softcore-exceptional-power-sum-word-closed-form} 得到。证毕。
\end{proof}

\begin{conclusion}[隧穿通道主导律：唯一极大指数与次极大指数]\label{con:xi-terminal-replica-softcore-tunneling-channel-dominance}
固定 \(m\ge 3\)，定义
\[
\Theta_m:=\max\{\Tr(w):w\in\{J,D\}^m,\ w\neq D^m\}.
\]
则 \(\Theta_m=2^m\)，且仅由 \(w=J^m\) 取得。
并且对所有 \(w\neq J^m\) 有
\[
\Tr(w)\le 3\cdot2^{m-2},
\]
等号当且仅当 \(w\) 含恰一段长度 \(1\) 的 \(D\)-块（等价于恰一枚 \(D\)，重数为 \(m\)）。
因此当 \(q\to\infty\) 时，
\[
\boxed{\;
S_m(q)=\frac1{2^m}\Big((2^m)^q+m(3\cdot2^{m-2})^q\Big)
+O_m\!\big((5\cdot2^{m-3})^q\big)\; }.
\]
\end{conclusion}
\begin{proof}
由结论 \ref{con:xi-terminal-replica-softcore-word-trace-fibonacci-factorization}，
任一含 \(J\) 字词可写为
\(
\Tr(w)=\prod_i F_{\ell_i+3}^{\,q}
\)，且 \(\sum_i(\ell_i+1)=m\)。
由 \(F_{\ell+3}\le 2^{\ell+1}\)（\(\ell\ge0\)，且仅 \(\ell=0\) 取等）得
\[
\Tr(w)\le 2^{mq},
\]
等价于指数底数上界 \(2^m\)，且等号当且仅当全部 \(\ell_i=0\)，即 \(w=J^m\)。
若 \(w\neq J^m\)，至少有一段 \(\ell_j\ge1\)。
在固定 \(\sum_i\ell_i\) 下，由
\[
F_{a+3}F_{b+3}=2F_{a+b+3}-F_aF_b<2F_{a+b+3}\qquad(a,b\ge1)
\]
可知将正段合并会严格增大乘积，故次极大由唯一正段给出。
若该段长度为 \(S\ge1\)，对应指数底数为
\[
F_{S+3}\,2^{m-S-1}.
\]
且 \(F_{S+3}/2^{S-1}\) 随 \(S\ge1\) 严格递减，故次极大在 \(S=1\) 取得
\(3\cdot2^{m-2}\)，重数为 \(m\)；下一层为 \(5\cdot2^{m-3}\)。
再结合 \(\Omega_m(q)=O(\varphi^{mq})\) 且 \(m\ge3\) 时 \(\varphi^m\le 5\cdot2^{m-3}\)，
代入结论 \ref{con:xi-terminal-replica-softcore-exceptional-power-sum-word-closed-form}
即得渐近式。证毕。
\end{proof}

\begin{conclusion}[二元投影词迹函数的二进--Fibonacci 正规形与二面体不变性]\label{con:xi-terminal-replica-softcore-binary-wordtrace-dihedral-normal-form}
设
\[
J:=\begin{pmatrix}1&1\\[2pt]1&1\end{pmatrix},
\qquad
M:=\begin{pmatrix}1&1\\[2pt]1&0\end{pmatrix}.
\]
对任意二元词 \(\omega=\omega_1\cdots \omega_m\in\{0,1\}^m\)，定义
\[
A(\omega):=\prod_{t=1}^{m}A_{\omega_t},
\qquad
A_0:=J,\ \ A_1:=M,
\qquad
\tau(\omega):=\Tr(A(\omega)).
\]
若 \(\omega=1^m\)，则 \(\tau(\omega)=L_m\)。若 \(\omega\neq 1^m\)，则存在整数
\(
r\ge 1,\ b_1,\dots,b_r\ge 1,\ a_1,\dots,a_r\ge 0
\)，使 \(\omega\) 在循环意义下可写为
\[
\omega\sim 0^{b_1}1^{a_1}0^{b_2}1^{a_2}\cdots 0^{b_r}1^{a_r}.
\]
记 \(z:=\sum_{i=1}^{r}b_i\)（零的总数），则有显式迹公式
\[
\tau(\omega)=2^{z-r}\prod_{i=1}^{r}F_{a_i+3}.
\]
进一步，\(\tau\) 在 \(\{0,1\}^m\) 上满足循环不变性与反转不变性：
\[
\tau(\omega)=\tau(\sigma\omega),\qquad
\tau(\omega)=\tau(\omega^{\mathrm{rev}}),
\]
其中 \(\sigma\) 为任意循环移位。因而 \(\tau\) 仅依赖于三项数据：零块个数 \(r\)、零总数 \(z\)、以及一块长度多重集 \(\{a_1,\dots,a_r\}\)。
\end{conclusion}
\begin{proof}[证明要点]
令 \(u:=(1,1)^{\top}\)，则 \(J=uu^{\top}\) 且 \(J^k=2^{k-1}J\)（\(k\ge 1\)）。
当 \(\omega\neq 1^m\) 时，\(A(\omega)\) 至少含一个 \(J\)，故 \(\det A(\omega)=0\) 且 \(\mathrm{rank}(A(\omega))\le 1\)。
按所示循环分解，
\[
A(\omega)=J^{b_1}M^{a_1}\cdots J^{b_r}M^{a_r}
=2^{z-r}\,J M^{a_1}J M^{a_2}\cdots J M^{a_r}.
\]
利用 \(J=uu^{\top}\) 有
\(
J M^a J=(u^{\top}M^a u)\,J
\)，迭代得
\[
\tau(\omega)=2^{z-r}\prod_{i=1}^{r}(u^{\top}M^{a_i}u).
\]
由 Fibonacci \(Q\)-矩阵恒等式
\[
M^n=
\begin{pmatrix}
F_{n+1}&F_n\\[2pt]
F_n&F_{n-1}
\end{pmatrix},
\]
得到 \(u^{\top}M^a u=F_{a+3}\)，从而结论成立。
循环不变性由迹的循环置换不变性给出；反转不变性可由
\(\Tr(X)=\Tr(X^{\top})\) 与上式仅依赖 \(\{a_i\}\) 的对称结构推出。证毕。
\end{proof}

\begin{conclusion}[\texorpdfstring{$\Tr((T^{(q)})^m)$}{Tr((T^(q))^m)} 的全阶词迹--张量幂闭式（互补编码）]\label{con:xi-terminal-replica-softcore-trace-moment-wordtrace-tensor-closed-complementary-coding}
令 \(V_q:=\{0,1\}^q\)，定义
\[
T^{(q)}_{u,v}:=2^{-\mathbf 1_{\{u\wedge v\neq 0\}}}\in\left\{1,\frac12\right\}.
\]
则对任意 \(q\ge 2\)、\(m\ge 1\)，有严格恒等式
\[
\Tr\!\bigl((T^{(q)})^m\bigr)
=2^{-m}\sum_{\omega\in\{0,1\}^m}\tau(\omega)^q,
\]
其中 \(\tau(\omega)\) 由结论 \ref{con:xi-terminal-replica-softcore-binary-wordtrace-dihedral-normal-form} 定义。
\end{conclusion}
\begin{proof}[证明要点]
由
\(
T^{(q)}_{u,v}=\frac12+\frac12\mathbf 1_{\{u\wedge v=0\}}
\)
可得
\[
T^{(q)}=\frac12\bigl(J^{\otimes q}+M^{\otimes q}\bigr).
\]
故
\[
\bigl(J^{\otimes q}+M^{\otimes q}\bigr)^m
=\sum_{\omega\in\{0,1\}^m}\bigl(A_{\omega_1}\cdots A_{\omega_m}\bigr)^{\otimes q}
=\sum_{\omega\in\{0,1\}^m}A(\omega)^{\otimes q}.
\]
取迹并用 \(\Tr(B^{\otimes q})=(\Tr B)^q\) 得
\[
\Tr\!\bigl((T^{(q)})^m\bigr)
=2^{-m}\sum_{\omega\in\{0,1\}^m}\Tr(A(\omega))^q
=2^{-m}\sum_{\omega\in\{0,1\}^m}\tau(\omega)^q.
\]
证毕。
\end{proof}

\begin{conclusion}[例外幂和的单词替换原理与对称幂角色分离]\label{con:xi-terminal-replica-softcore-exceptional-power-sum-word-substitution-principle}
记
\[
S_m(q):=\sum_{i=0}^{q}\bigl(\lambda_i^{(q)}\bigr)^m,
\]
其中 \(\{\lambda_i^{(q)}\}_{i=0}^{q}\) 为 \(T^{(q)}\) 在 \(S_q\)-不变子空间上的例外特征值。
则对任意 \(q\ge 2\)、\(m\ge 1\)，成立
\[
S_m(q)
=2^{-m}\!\left(
\Omega_m(q)+
\sum_{\omega\in\{0,1\}^m\setminus\{1^m\}}\tau(\omega)^q
\right),
\]
其中 \(\Omega_m(q)\) 的定义见结论
\ref{con:xi-terminal-replica-softcore-lucas-pell-omega-unified}。
等价地，
\[
S_m(q)-\Tr\!\bigl((T^{(q)})^m\bigr)
=2^{-m}\bigl(\Omega_m(q)-L_m^q\bigr).
\]
\end{conclusion}
\begin{proof}[证明要点]
由结论 \ref{con:xi-terminal-replica-softcore-exceptional-spectrum-sympower-model}，
\[
T^{(q)}\big|_{\mathrm{Sym}^q}
=\frac12\bigl(\mathrm{Sym}^q(J)+\mathrm{Sym}^q(M)\bigr),
\]
故
\[
S_m(q)=2^{-m}\sum_{\omega\in\{0,1\}^m}
\Tr\!\bigl(\mathrm{Sym}^q(A(\omega))\bigr).
\]
当 \(\omega\neq 1^m\) 时，\(A(\omega)\) 含 \(J\)，故 \(\det A(\omega)=0\) 且秩为 \(1\)，特征值为 \(\{\tau(\omega),0\}\)，于是
\[
\Tr\!\bigl(\mathrm{Sym}^q(A(\omega))\bigr)=\tau(\omega)^q.
\]
当 \(\omega=1^m\) 时，\(A(\omega)=M^m\)，其贡献即
\(
\Tr(\mathrm{Sym}^q(M^m))=\Omega_m(q)
\)。
故得第一式。再与结论
\ref{con:xi-terminal-replica-softcore-trace-moment-wordtrace-tensor-closed-complementary-coding}
相减，即得第二式。证毕。
\end{proof}

\begin{corollary}[Fibonacci 原始素因子导出的高指标刚性]\label{cor:xi-terminal-replica-softcore-fibonacci-primitive-divisor-high-index-rigidity}
\leanverified{paper\_xi\_terminal\_replica\_softcore\_fibonacci\_primitive\_divisor\_high\_index\_rigidity}
令 \(\omega,\omega'\in\{0,1\}^m\setminus\{1^m\}\)，并按结论
\ref{con:xi-terminal-replica-softcore-binary-wordtrace-dihedral-normal-form}
写成
\[
\tau(\omega)=2^{z-r}\prod_{i=1}^{r}F_{a_i+3},
\qquad
\tau(\omega')=2^{z'-r'}\prod_{j=1}^{r'}F_{a'_j+3}.
\]
若 \(\tau(\omega)=\tau(\omega')\)，则两侧下标多重集在截断区间 \((12,\infty)\) 上完全一致（含重数）：
\[
\{a_i+3:\ a_i+3>12\}_{i=1}^{r}
=
\{a'_j+3:\ a'_j+3>12\}_{j=1}^{r'}.
\]
等价地，\(\{a_i:\ a_i>9\}\) 与 \(\{a'_j:\ a'_j>9\}\) 作为多重集一致。
\end{corollary}
\begin{proof}[证明要点]
令
\(
N:=\max\!\bigl(\{a_i+3\}_{i=1}^{r}\cup\{a'_j+3\}_{j=1}^{r'}\bigr)
\)。
若 \(N>12\)，由 Fibonacci 原始素因子定理（除低阶例外 \(6,12\) 外，\(F_N\) 含不整除任一更小 \(F_k\) 的素因子）可取素数 \(p\mid F_N\)，且 \(p\nmid F_k\)（\(k<N\)）。
于是等式 \(\tau(\omega)=\tau(\omega')\) 两端中，\(p\) 只能来自下标 \(N\) 的 Fibonacci 因子，故 \(F_N\) 的出现次数在两端相同。
约去共同的 \(F_N\) 因子后重复该“最大下标下降”步骤，直至所有剩余下标 \(\le 12\)。
故结论成立。证毕。
\end{proof}

\begin{proposition}[迹矩与例外幂和的二面体轨道压缩表示]\label{prop:xi-terminal-replica-softcore-dihedral-orbit-compression-trace-moments}
记 \(\mathrm{Brace}^{(2)}_m\) 为长度 \(m\) 的二元手镯集合（即 \(\{0,1\}^m\) 在循环移位与反转生成的二面体作用下的等价类），
对 \(b\in \mathrm{Brace}^{(2)}_m\) 记其轨道大小为 \(|\mathrm{Orb}(b)|\)。
则对任意 \(q\ge 2\)、\(m\ge 1\)，有
\[
\Tr\!\bigl((T^{(q)})^m\bigr)
=2^{-m}\sum_{b\in \mathrm{Brace}^{(2)}_m}
|\mathrm{Orb}(b)|\,\tau(b)^q,
\]
以及
\[
S_m(q)
=2^{-m}\!\left(
\Omega_m(q)+
\sum_{b\in \mathrm{Brace}^{(2)}_m\setminus\{[1^m]\}}
|\mathrm{Orb}(b)|\,\tau(b)^q
\right).
\]
\end{proposition}
\begin{proof}
由结论 \ref{con:xi-terminal-replica-softcore-binary-wordtrace-dihedral-normal-form}，
\(\tau\) 在二面体轨道上常值。
将结论
\ref{con:xi-terminal-replica-softcore-trace-moment-wordtrace-tensor-closed-complementary-coding}
与结论
\ref{con:xi-terminal-replica-softcore-exceptional-power-sum-word-substitution-principle}
中的词求和按二面体轨道分组，即得两式。证毕。
\end{proof}

\endinput

\paragraph{分拆参数化的全空间迹分解、对称幂纠正项与 \texorpdfstring{$q$}{q}-方向动力学}
本段将结论 \ref{con:xi-terminal-replica-softcore-partition-trace-closed-form} 与
\ref{con:xi-terminal-replica-softcore-partition-exceptional-closed-form}
在统一分拆记号下重写为一组可直接用于生成函数、渐近层级与递推算子的闭式模板。

\begin{theorem}[全空间幂迹的分拆--Fibonacci 完备指数分解]\label{thm:xi-terminal-replica-softcore-fullspace-trace-partition-fib-complete}
令 \(q\ge 1\)、\(m\ge 1\)，并沿用
\(
T^{(q)}=\frac12(J^{\otimes q}+K^{\otimes q})
\)
（结论 \ref{con:xi-terminal-replica-softcore-binary-necklace-trace-formula}）。
对分拆 \(\pi\vdash m\) 记
\[
\pi=1^{m_1}2^{m_2}\cdots m^{m_m},\qquad
\ell(\pi):=\sum_{s=1}^{m}m_s,
\]
\[
\Theta(\pi):=\prod_{s=1}^{m}F_{s+2}^{\,m_s},
\qquad
c(\pi):=\frac{m}{\ell(\pi)}\cdot\frac{\ell(\pi)!}{\prod_{s=1}^{m}m_s!}\in\ZZ_{>0}.
\]
则
\begin{equation}\label{eq:xi-terminal-replica-softcore-fullspace-trace-partition-fib-complete}
\Tr\!\bigl((T^{(q)})^m\bigr)
=\frac1{2^m}\left(
L_m^{\,q}+\sum_{\pi\vdash m}c(\pi)\,\Theta(\pi)^{\,q}
\right).
\end{equation}
\end{theorem}
\begin{proof}
结论 \ref{con:xi-terminal-replica-softcore-partition-trace-closed-form} 已给出
\[
\Tr\!\bigl((T^{(q)})^m\bigr)
=\frac1{2^m}\left(L_m^{\,q}+\sum_{\lambda\vdash m}c_{m,\lambda}B_\lambda^{\,q}\right),
\]
其中 \(\lambda=1^{a_1}2^{a_2}\cdots\)。
取 \(m_s:=a_s\) 并记 \(\pi:=\lambda\)，则
\[
B_\lambda=\prod_{s\ge 1}F_{s+2}^{\,a_s}=\Theta(\pi),\qquad
c_{m,\lambda}=\frac{m}{\ell(\pi)}\frac{\ell(\pi)!}{\prod_s m_s!}=c(\pi),
\]
从而得到 \eqref{eq:xi-terminal-replica-softcore-fullspace-trace-partition-fib-complete}。证毕。
\end{proof}

\begin{theorem}[例外谱幂和的分拆闭式与对称幂表征]\label{thm:xi-terminal-replica-softcore-exceptional-power-partition-sympower}
令
\(
S_m(q):=\sum_{i=0}^{q}(\lambda_i^{(q)})^m
\)。
在定理 \ref{thm:xi-terminal-replica-softcore-fullspace-trace-partition-fib-complete} 的记号下，
有
\begin{equation}\label{eq:xi-terminal-replica-softcore-exceptional-power-partition-sympower}
S_m(q)=\frac1{2^m}\left(
\Omega_m(q)+\sum_{\pi\vdash m}c(\pi)\,\Theta(\pi)^{\,q}
\right),
\end{equation}
且
\begin{equation}\label{eq:xi-terminal-replica-softcore-omega-sympower-Km}
\Omega_m(q)=\Tr\!\bigl(\operatorname{Sym}^q(K^m)\bigr).
\end{equation}
\end{theorem}
\begin{proof}
公式 \eqref{eq:xi-terminal-replica-softcore-exceptional-power-partition-sympower}
是结论 \ref{con:xi-terminal-replica-softcore-partition-exceptional-closed-form}
在 \(\pi\)-记号下的等价改写；而
\eqref{eq:xi-terminal-replica-softcore-omega-sympower-Km}
由结论
\ref{con:xi-terminal-replica-softcore-exceptional-power-sum-word-trace-sympower}
直接给出。证毕。
\end{proof}

\begin{theorem}[固定 \texorpdfstring{$m$}{m} 时的 \texorpdfstring{$q$}{q}-方向有理生成函数]\label{thm:xi-terminal-replica-softcore-q-genfunc-rational-partition}
固定 \(m\ge 1\)，定义形式幂级数
\[
\mathcal T_m(z):=\sum_{q\ge 0}\Tr\!\bigl((T^{(q)})^m\bigr)z^q,
\qquad
\mathcal S_m(z):=\sum_{q\ge 0}S_m(q)z^q.
\]
则
\begin{equation}\label{eq:xi-terminal-replica-softcore-q-genfunc-trace}
\mathcal T_m(z)
=\frac1{2^m}\left(
\frac{1}{1-L_m z}
+\sum_{\pi\vdash m}\frac{c(\pi)}{1-\Theta(\pi)z}
\right),
\end{equation}
\begin{equation}\label{eq:xi-terminal-replica-softcore-q-genfunc-exceptional}
\mathcal S_m(z)
=\frac1{2^m}\left(
\frac{1}{1-L_m z+(-1)^m z^2}
+\sum_{\pi\vdash m}\frac{c(\pi)}{1-\Theta(\pi)z}
\right).
\end{equation}
\end{theorem}
\begin{proof}
由定理 \ref{thm:xi-terminal-replica-softcore-fullspace-trace-partition-fib-complete}，
\[
\Tr\!\bigl((T^{(q)})^m\bigr)
=2^{-m}\!\left(L_m^q+\sum_{\pi\vdash m}c(\pi)\Theta(\pi)^q\right),
\]
逐项用几何级数求和得 \eqref{eq:xi-terminal-replica-softcore-q-genfunc-trace}。

由定理 \ref{thm:xi-terminal-replica-softcore-exceptional-power-partition-sympower}，
\[
S_m(q)=2^{-m}\!\left(\Omega_m(q)+\sum_{\pi\vdash m}c(\pi)\Theta(\pi)^q\right).
\]
而 \(\Omega_m(q)\) 满足递推
\(
\Omega_m(q)=L_m\Omega_m(q-1)-(-1)^m\Omega_m(q-2)
\)
（结论 \ref{con:xi-terminal-replica-softcore-lucas-pell-omega-unified}），并有
\(\Omega_m(0)=1,\Omega_m(1)=L_m\)，故
\[
\sum_{q\ge 0}\Omega_m(q)z^q=\frac{1}{1-L_m z+(-1)^m z^2}.
\]
与分拆项求和叠加得到 \eqref{eq:xi-terminal-replica-softcore-q-genfunc-exceptional}。证毕。
\end{proof}
\leanverified{paper\_xi\_terminal\_replica\_softcore\_q\_genfunc\_rational\_partition}

\begin{theorem}[分拆权重层级与 \texorpdfstring{$q\to\infty$}{q->infty} 的前三阶指数渐近]\label{thm:xi-terminal-replica-softcore-partition-weight-hierarchy-asymptotics}
固定 \(m\ge 3\)。对 \(\pi\vdash m\) 定义 \(\Theta(\pi)\) 如上。则：
\begin{enumerate}
\item \(\max_{\pi\vdash m}\Theta(\pi)=2^m\)，且仅由 \(\pi=1^m\) 取得；
\item 在 \(\pi\neq 1^m\) 中，最大值为 \(3\cdot 2^{m-2}\)，且仅由 \(\pi=2,1^{m-2}\) 取得；
\item 排除前两者后，最大值为 \(5\cdot 2^{m-3}\)，且仅由 \(\pi=3,1^{m-3}\) 取得。
\end{enumerate}
从而
\begin{align}
S_m(q)
=&\ 2^{m(q-1)}
+\frac{m}{2^m}\bigl(3\cdot 2^{m-2}\bigr)^q
+\frac{m}{2^m}\bigl(5\cdot 2^{m-3}\bigr)^q\notag\\
&\ +O\!\left(\bigl(9\cdot 2^{m-4}\bigr)^q\right),
\qquad q\to\infty.\label{eq:xi-terminal-replica-softcore-top3-asymptotics}
\end{align}
\end{theorem}
\begin{proof}
设
\[
r_s:=\frac{F_{s+2}}{2^s}\qquad(s\ge 1).
\]
由 \(F_{s+3}=F_{s+2}+F_{s+1}<2F_{s+2}\) 得
\[
\frac{r_{s+1}}{r_s}=\frac{F_{s+3}}{2F_{s+2}}<1,
\]
故 \(r_s\) 严格递减，且 \(r_1=1\)、\(r_s<1\)（\(s\ge 2\)）。
对 \(\pi=1^{m_1}\cdots m^{m_m}\)：
\[
\Theta(\pi)
=\prod_{s=1}^{m}(2^s r_s)^{m_s}
=2^m\prod_{s=2}^{m}r_s^{m_s}.
\]
因此最大值仅在 \(m_s=0\ (s\ge 2)\) 时取得，即 \(\pi=1^m\)，得到第 1 条。

若 \(\pi\neq 1^m\)，为使 \(\prod_{s\ge2}r_s^{m_s}\) 最大，应只引入一次最小代价部件 \(s=2\)，即 \(\pi=2,1^{m-2}\)，给出第 2 条。
继续排除后，候选最大来自 \(\pi=3,1^{m-3}\) 与 \(\pi=2^2,1^{m-4}\)，对应
\[
5\cdot 2^{m-3}\quad\text{与}\quad 9\cdot 2^{m-4},
\]
且 \(5/8>9/16\)，故第 3 条成立。

由定理 \ref{thm:xi-terminal-replica-softcore-exceptional-power-partition-sympower}，
\[
S_m(q)=2^{-m}\!\left(\Omega_m(q)+\sum_{\pi\vdash m}c(\pi)\Theta(\pi)^q\right).
\]
前三项系数分别为 \(c(1^m)=1\)、\(c(2,1^{m-2})=m\)、\(c(3,1^{m-3})=m\)。
其余分拆项底数至多 \(9\cdot 2^{m-4}\)；且 \(\Omega_m(q)=O(\varphi^{mq})\)，并有
\(
\varphi^m<9\cdot 2^{m-4}
\)（\(m\ge3\)）。
故得 \eqref{eq:xi-terminal-replica-softcore-top3-asymptotics}。证毕。
\end{proof}

\begin{theorem}[固定 \texorpdfstring{$m$}{m} 的显式常系数递推（阶至多 \texorpdfstring{$p(m)+2$}{p(m)+2}）]\label{thm:xi-terminal-replica-softcore-q-recurrence-partition-annihilator}
固定 \(m\ge 1\)，记
\[
\mathscr B_m:=\{\Theta(\pi):\pi\vdash m\}\subset\ZZ_{>0},\qquad
a_q:=2^m S_m(q),\qquad b_q:=2^m\Tr\!\bigl((T^{(q)})^m\bigr).
\]
则
\begin{equation}\label{eq:xi-terminal-replica-softcore-q-recurrence-annihilator-a}
\Bigl(E^2-L_mE+(-1)^m\Bigr)\cdot\prod_{\theta\in\mathscr B_m}(E-\theta)\,a=0,
\end{equation}
\begin{equation}\label{eq:xi-terminal-replica-softcore-q-recurrence-annihilator-b}
(E-L_m)\cdot\prod_{\theta\in\mathscr B_m}(E-\theta)\,b=0.
\end{equation}
故 \(S_m(q)\) 满足阶数不超过 \(|\mathscr B_m|+2\le p(m)+2\) 的常系数线性递推。
\end{theorem}
\begin{proof}
由定理 \ref{thm:xi-terminal-replica-softcore-exceptional-power-partition-sympower}，
\[
a_q=\Omega_m(q)+\sum_{\pi\vdash m}c(\pi)\Theta(\pi)^q.
\]
每项 \(\Theta(\pi)^q\) 被 \(E-\Theta(\pi)\) 湮灭；\(\Omega_m\) 被
\(
E^2-L_mE+(-1)^m
\)
湮灭。
乘积算子因此湮灭 \(a\)，得 \eqref{eq:xi-terminal-replica-softcore-q-recurrence-annihilator-a}。

同理，由定理 \ref{thm:xi-terminal-replica-softcore-fullspace-trace-partition-fib-complete}，
\[
b_q=L_m^q+\sum_{\pi\vdash m}c(\pi)\Theta(\pi)^q,
\]
得 \eqref{eq:xi-terminal-replica-softcore-q-recurrence-annihilator-b}。证毕。
\end{proof}

\begin{corollary}[分拆系数总和恒等式]\label{cor:xi-terminal-replica-softcore-partition-coeff-sum-2m-minus-1}
对任意 \(m\ge 1\)，有
\begin{equation}\label{eq:xi-terminal-replica-softcore-partition-coeff-sum}
\sum_{\pi\vdash m}c(\pi)=2^m-1.
\end{equation}
\leanverified{paper\_xi\_terminal\_replica\_softcore\_partition\_coeff\_sum\_2m\_minus\_1}
\end{corollary}
\begin{proof}
由定理 \ref{thm:xi-terminal-replica-softcore-fullspace-trace-partition-fib-complete} 中的组合解释，\(c(\pi)\) 计数长度 \(m\) 的二元词中“含至少一个 \(J\)”且块长分拆为 \(\pi\) 的词数。
各分拆类互不相交并覆盖全部含 \(J\) 的词。总词数 \(2^m\)，唯一不含 \(J\) 的词是 \(K^m\)，故总数为 \(2^m-1\)。证毕。
\end{proof}

\begin{theorem}[分拆--Fibonacci 直积映射的同值核结构]\label{thm:xi-terminal-replica-softcore-partition-fib-collision-kernel}
固定 \(m\ge 1\)。设 \(\lambda,\mu\vdash m\)，并写
\[
\lambda=1^{a_1}2^{a_2}\cdots,\qquad
\mu=1^{b_1}2^{b_2}\cdots,\qquad
d_r:=a_r-b_r\in\ZZ.
\]
则
\[
B_\lambda=B_\mu
\iff
\exists\,t\in\ZZ:\ 
(d_1,d_2,d_4,d_{10})=t(2,-2,-2,1),\quad
d_r=0\ \ (r\notin\{1,2,4,10\}).
\]
\end{theorem}
\begin{proof}
由定义
\[
B_\lambda=B_\mu
\iff
\prod_{r\ge 1}F_{r+2}^{\,d_r}=1.
\]
又因 \(\lambda,\mu\) 同为 \(m\) 的分拆，必有
\[
\sum_{r\ge 1}r\,d_r=0.
\]
\textbf{第一步：高指标消去。}
对 Fibonacci 数列应用原始素因子定理（Carmichael--Zsigmondy 型结论）：
当 \(n\notin\{1,2,6,12\}\) 时，存在素数 \(p\mid F_n\) 且 \(p\nmid F_k\)（\(k<n\)）。
设 \(N:=\max\{r+2:d_r\neq 0\}\)。若 \(N\ge 13\)，则取 \(p\mid F_N\) 为上述原始素因子。
由于所有出现指标均不超过 \(N\)，乘积 \(\prod_rF_{r+2}^{d_r}\) 的 \(p\)-进赋值仅来自 \(F_N^{d_{N-2}}\)，故必有 \(d_{N-2}=0\)，与 \(N\) 的定义矛盾。
因此 \(d_r=0\) 对全部 \(r\ge 11\)。

\textbf{第二步：\(r\le 10\) 的素因子隔离。}
列出
\[
F_3=2,\ F_4=3,\ F_5=5,\ F_6=8,\ F_7=13,\ F_8=21,\ F_9=34,\ F_{10}=55,\ F_{11}=89,\ F_{12}=144.
\]
在该范围内，素因子 \(13,7,17,11,89\) 分别只出现于 \(F_7,F_8,F_9,F_{10},F_{11}\)，故
\[
d_5=d_6=d_7=d_8=d_9=0.
\]
再由 \(F_5=5\) 且 \(5\) 仅与 \(F_{10}\) 共现，而 \(d_8=0\)，得 \(d_3=0\)。
于是仅可能
\[
d_r\neq 0\Rightarrow r\in\{1,2,4,10\}.
\]

\textbf{第三步：线性求解。}
由
\[
F_3=2,\qquad F_4=3,\qquad F_6=2^3,\qquad F_{12}=2^4\cdot 3^2,
\]
等式 \(\prod_rF_{r+2}^{d_r}=1\) 化为
\[
d_2+2d_{10}=0,\qquad d_1+3d_4+4d_{10}=0.
\]
再与分拆权重约束
\[
d_1+2d_2+4d_4+10d_{10}=0
\]
联立，解得
\[
d_2=-2d_{10},\qquad d_4=-2d_{10},\qquad d_1=2d_{10}.
\]
设 \(t:=d_{10}\) 即得
\[
(d_1,d_2,d_4,d_{10})=t(2,-2,-2,1).
\]
反向蕴含由恒等式
\[
F_{12}\,F_3^2=F_6^2\,F_4^2
\]
立即成立。证毕。
\end{proof}

\begin{corollary}[最小同值分拆块与低阶单射性]\label{cor:xi-terminal-replica-softcore-partition-fib-minimal-collision}
设 \(\mathcal B_m:=\{B_\lambda:\lambda\vdash m\}\)。
则：
\begin{enumerate}
\item 当 \(m\le 11\) 时，映射 \(\lambda\mapsto B_\lambda\) 在 \(\{\lambda:\lambda\vdash m\}\) 上单射；
\item 最小非平凡同值出现在 \(m=12\)，并由唯一基本置换块给出
\[
(10,1,1)\sim(4,4,2,2),\qquad
B_{(10,1,1)}=B_{(4,4,2,2)}=576.
\]
\end{enumerate}
\end{corollary}
\begin{proof}
由定理 \ref{thm:xi-terminal-replica-softcore-partition-fib-collision-kernel}，
任意非零差向量均为 \(t(2,-2,-2,1)\)。
其正向质量与负向质量各自的加权和均为 \(12|t|\)。
故当 \(m\le 11\) 时不可能出现非平凡同值，单射成立。

当 \(m=12\) 时，仅有 \(|t|=1\) 可能；并且无法再叠加额外公共分拆部分。
因此唯一基本同值对即
\((10,1,1)\) 与 \((4,4,2,2)\)（互换对应 \(t=\pm1\)）。
数值等式由
\(
F_{12}F_3^2=F_6^2F_4^2=576
\)
给出。证毕。
\end{proof}

\begin{proposition}[同值等价类的整数参数化与类大小]\label{prop:xi-terminal-replica-softcore-partition-fib-equivalence-class-param}
固定 \(\lambda\vdash m\)，写 \(\lambda=1^{a_1}2^{a_2}3^{a_3}\cdots\)，并定义
\[
\mathcal E(\lambda):=\{\mu\vdash m:\ B_\mu=B_\lambda\}.
\]
则 \(\mu=1^{b_1}2^{b_2}\cdots\in\mathcal E(\lambda)\) 当且仅当：
\begin{enumerate}
\item 对全部 \(r\notin\{1,2,4,10\}\)，有 \(b_r=a_r\)；
\item 存在 \(s\in\ZZ\) 使
\[
(b_1,b_2,b_4,b_{10})
=(a_1,a_2,a_4,a_{10})+s(2,-2,-2,1),
\]
且右端各分量非负。
\end{enumerate}
等价地，\(s\) 的可行区间为
\[
s_{\min}:=\max\!\left\{-\left\lfloor\frac{a_1}{2}\right\rfloor,\,-a_{10}\right\},
\qquad
s_{\max}:=\min\!\left\{\left\lfloor\frac{a_2}{2}\right\rfloor,\ \left\lfloor\frac{a_4}{2}\right\rfloor\right\},
\]
并有
\[
\#\mathcal E(\lambda)=s_{\max}-s_{\min}+1.
\]
\end{proposition}
\begin{proof}
由定理 \ref{thm:xi-terminal-replica-softcore-partition-fib-collision-kernel}，
\(
d_r=a_r-b_r
\)
必须满足
\(
d=-s(2,-2,-2,1)
\)
（在 \(r=1,2,4,10\) 上）且其余分量为零。
于是得到充要条件。
非负整数约束等价于
\[
a_1+2s\ge 0,\quad a_{10}+s\ge 0,\quad a_2-2s\ge 0,\quad a_4-2s\ge 0,
\]
从而 \(s\) 构成闭整数区间 \([s_{\min},s_{\max}]\cap\ZZ\)，点数即上式。证毕。
\end{proof}

\begin{theorem}[分拆迹值去重计数的精确闭式]\label{thm:xi-terminal-replica-softcore-partition-fib-distinct-count-exact}
记 \(p(n)\) 为整数分拆函数（约定 \(p(n)=0\) 对 \(n<0\)），并定义
\[
N_m:=\bigl|\{B_\lambda:\lambda\vdash m\}\bigr|.
\]
则对任意 \(m\ge 0\) 有
\[
N_m=p(m)-p(m-12).
\]
\end{theorem}
\begin{proof}
在 \(\mathrm{Part}(m)\) 上定义无向图 \(G_m\)：
对每个 \(\nu\vdash(m-12)\)，连边
\[
\nu\sqcup(10,1,1)\longleftrightarrow \nu\sqcup(4,4,2,2).
\]
由恒等式 \(F_{12}F_3^2=F_6^2F_4^2\)，每条边两端具有相同 \(B\)-值。
反之，若两分拆同值，则由定理
\ref{thm:xi-terminal-replica-softcore-partition-fib-collision-kernel}
其差向量为 \(t(2,-2,-2,1)\)，即可由有限次上述基本块替换连通。
故 \(G_m\) 的连通分支与 \(B\)-值等价类一一对应，从而
\[
N_m=\#\mathrm{Comp}(G_m).
\]
又顶点数 \(\#V(G_m)=p(m)\)。
边与 \(\nu\vdash(m-12)\) 一一对应，故 \(\#E(G_m)=p(m-12)\)。
由命题 \ref{prop:xi-terminal-replica-softcore-partition-fib-equivalence-class-param}，
每个等价类由整数参数 \(s\) 形成区间，边对应 \(s\mapsto s\pm1\)，故每个连通分支是路径，\(G_m\) 为森林。
因此
\[
\#\mathrm{Comp}(G_m)=\#V(G_m)-\#E(G_m)=p(m)-p(m-12).
\]
证毕。
\end{proof}

\begin{proposition}[同值聚合系数的显式求和公式]\label{prop:xi-terminal-replica-softcore-partition-fib-aggregated-coeff}
固定 \(m\ge 1\)。对
\(
B\in\{B_\lambda:\lambda\vdash m\}
\)
定义聚合系数
\[
C_m(B):=\sum_{\substack{\lambda\vdash m\\ B_\lambda=B}}c_{m,\lambda}.
\]
取任意 \(\lambda=1^{a_1}2^{a_2}3^{a_3}\cdots\vdash m\)，并令 \(s_{\min},s_{\max}\) 如命题
\ref{prop:xi-terminal-replica-softcore-partition-fib-equivalence-class-param}。
对 \(s\in[s_{\min},s_{\max}]\cap\ZZ\)，定义 \(\lambda(s)\) 的部件重数
\[
a_1(s)=a_1+2s,\quad a_2(s)=a_2-2s,\quad a_4(s)=a_4-2s,\quad a_{10}(s)=a_{10}+s,
\]
\[
a_r(s)=a_r\ \ (r\notin\{1,2,4,10\}),\qquad
k(s):=\sum_{r\ge 1}a_r(s).
\]
则
\[
C_m(B_\lambda)=\sum_{s=s_{\min}}^{s_{\max}}
\frac{m\,(k(s)-1)!}{\prod_{r\ge 1}a_r(s)!}.
\]
\end{proposition}
\begin{proof}
由命题 \ref{prop:xi-terminal-replica-softcore-partition-fib-equivalence-class-param}，
\(
\{\mu\vdash m:B_\mu=B_\lambda\}
=\{\lambda(s):s_{\min}\le s\le s_{\max}\}
\)。
逐项代入结论
\ref{con:xi-terminal-replica-softcore-partition-trace-closed-form}
中的系数公式
\(
c_{m,\mu}=m\,(k_\mu-1)!/\prod_r a_r(\mu)!
\)
即得。证毕。
\end{proof}

\begin{corollary}[去重迹值计数的 Hardy--Ramanujan 渐近]\label{cor:xi-terminal-replica-softcore-partition-fib-distinct-count-asymptotic}
当 \(m\to\infty\) 时，
\[
N_m\sim \frac{\pi}{\sqrt2}\,m^{-3/2}
\exp\!\left(\pi\sqrt{\frac{2m}{3}}\right).
\]
\end{corollary}
\begin{proof}
由定理 \ref{thm:xi-terminal-replica-softcore-partition-fib-distinct-count-exact}，
\(
N_m=p(m)-p(m-12)
\)。
结合 Hardy--Ramanujan 渐近
\[
p(m)\sim \frac{1}{4m\sqrt3}\exp\!\left(\pi\sqrt{\frac{2m}{3}}\right),
\]
并用展开
\[
\sqrt{m-12}
=\sqrt m-\frac{6}{\sqrt m}+O(m^{-3/2}),
\]
得
\[
\frac{p(m-12)}{p(m)}
=\exp\!\left(-6\pi\sqrt{\frac{2}{3m}}+O(m^{-1})\right).
\]
于是
\[
N_m
=p(m)\left(1-\frac{p(m-12)}{p(m)}\right)
\sim p(m)\cdot 6\pi\sqrt{\frac{2}{3m}}.
\]
代入 \(p(m)\) 主项并化简常数得到
\[
N_m\sim \frac{\pi}{\sqrt2}\,m^{-3/2}
\exp\!\left(\pi\sqrt{\frac{2m}{3}}\right).
\]
证毕。
\end{proof}

\endinput


\begin{conclusion}[分拆--Fibonacci 迹像：\texorpdfstring{$\tau(\omega)$}{tau(omega)} 的分拆化刻画]\label{con:xi-terminal-replica-softcore-trace-partition-image}
固定整数 \(m\ge 1\)，沿用结论 \ref{con:xi-terminal-replica-softcore-binary-necklace-trace-formula} 的 \(\tau(\omega)\)。
对整数分拆 \(\lambda\vdash m\) 写作 \(\lambda=(\lambda_1,\dots,\lambda_k)\)（\(\lambda_i\ge 1\)、\(\sum_i\lambda_i=m\)），并定义
\[
B_\lambda:=\prod_{i=1}^{k}F_{\lambda_i+2}\in\ZZ_{>0}.
\]
则除全零词外的迹值集合满足
\[
\{\tau(\omega):\omega\in\{0,1\}^m,\ \omega\neq 0^m\}
=\{B_\lambda:\lambda\vdash m\},
\qquad
\tau(0^m)=L_m.
\]
因此 \(\Tr((T^{(q)})^m)\) 在 \(q\)-方向上为有限指数多项式，其指数底数集合由 \(\{L_m\}\cup\{B_\lambda:\lambda\vdash m\}\) 的数值取值所生成。
\end{conclusion}
\begin{proof}[证明要点]
由结论 \ref{con:xi-terminal-replica-softcore-binary-necklace-trace-formula}，
若 \(\eta\in\mathrm{Neck}_m^{(2)}\) 含至少一个 \(1\)，将其循环表示为
\(
\eta\sim 1\,0^{r_1}\,1\,0^{r_2}\cdots 1\,0^{r_k}
\)
（\(k\ge 1\)、\(r_i\ge 0\)、\(\sum_i r_i=m-k\)），则
\(
\tau(\eta)=\prod_{i=1}^{k}F_{r_i+3}
\)。
令 \(\lambda_i:=r_i+1\)，则 \((\lambda_1,\dots,\lambda_k)\) 为 \(m\) 的一个正整数组合，其分拆化（按降序重排）给出某个 \(\lambda\vdash m\)，并且
\(
\prod_i F_{r_i+3}=\prod_i F_{\lambda_i+2}=B_\lambda
\)。
迹的循环不变性给出 \(\tau(\omega)=\tau(\eta)\)，从而左侧包含于右侧；反向包含由对任意 \(\lambda=(\lambda_1,\dots,\lambda_k)\) 取 necklace 代表
\(
1\,0^{\lambda_1-1}\cdots 1\,0^{\lambda_k-1}
\)
直接构造得到。
全零词情形为 \(\tau(0^m)=\Tr(K^m)=L_m\)。证毕。
\end{proof}

\begin{conclusion}[全迹矩的分拆闭式：\texorpdfstring{$\Tr((T^{(q)})^m)$}{Tr((T^(q))^m)} 的分拆系数可计算展开]\label{con:xi-terminal-replica-softcore-partition-trace-closed-form}
固定整数 \(m\ge 1\) 与 \(q\ge 2\)，沿用 \(\{B_\lambda\}_{\lambda\vdash m}\)（结论 \ref{con:xi-terminal-replica-softcore-trace-partition-image}）。
对分拆 \(\lambda\vdash m\) 记其部件计数为 \(\lambda=1^{a_1}2^{a_2}\cdots\)，并记部件总数 \(k:=\sum_{r\ge1}a_r\)。
定义分拆系数
\[
c_{m,\lambda}:=\frac{m}{k}\cdot\frac{k!}{\prod_{r\ge1}a_r!}
= m\cdot\frac{(k-1)!}{\prod_{r\ge1}a_r!}\ \in\ \ZZ_{>0}.
\]
则对任意 \(q\ge 2\) 成立严格闭式
\[
\boxed{\;
\Tr\!\bigl((T^{(q)})^m\bigr)
=\frac1{2^m}\left(L_m^{q}+\sum_{\lambda\vdash m}c_{m,\lambda}\,B_\lambda^{q}\right)\; } ,
\qquad
B_\lambda=\prod_{r\ge1}F_{r+2}^{\,a_r}.
\]
并且系数满足总量守恒
\[
\sum_{\lambda\vdash m}c_{m,\lambda}=2^m-1.
\]
\end{conclusion}
\begin{proof}[证明要点]
由结论 \ref{con:xi-terminal-replica-softcore-binary-necklace-trace-formula}，
\(
2^{m}\Tr((T^{(q)})^m)=\sum_{\omega\in\{0,1\}^m}\tau(\omega)^q
\)，且全零词贡献为 \(\tau(0^m)^q=L_m^q\)。
对 \(\omega\neq 0^m\)，将其旋转使其以 \(1\) 起首，并以相邻 \(1\) 之间的循环块长得到一个由 \(\lambda\) 给出的正整数组合；其分拆化为某个 \(\lambda\vdash m\)，且按结论 \ref{con:xi-terminal-replica-softcore-trace-partition-image} 有 \(\tau(\omega)=B_\lambda\)。
\par
对固定 \(\lambda=1^{a_1}2^{a_2}\cdots\)（部件数 \(k\)），以 \(1\,0^{r-1}\) 作为长度为 \(r\) 的基本块，则以 \(1\) 起首且类型为 \(\lambda\) 的线性字词恰对应于对这 \(k\) 个基本块（含重数）作排列，故个数为 \(k!/\prod_r a_r!\)。
再对长度 \(m\) 的循环旋转作用分解为轨道：每个轨道中的字词均含恰 \(k\) 个 \(1\)，因此该轨道中以 \(1\) 起首的字词数与轨道总大小之比恒为 \(k/m\)。
将该比例对所有轨道求和得到类型 \(\lambda\) 的字词总数为 \((m/k)\cdot(k!/\prod_r a_r!)=c_{m,\lambda}\)，代回即得闭式。
最后一式由 \(\sum_{\lambda\vdash m}c_{m,\lambda}\) 计数所有非全零字词的总数 \(2^m-1\) 得出。证毕。
\end{proof}

\subsubsection{分拆闭式与剥离律：例外幂和的递推与生成函数证书}\label{subsubsec:xi-exceptional-partition-peeling-recurrence}

\begin{conclusion}[全迹矩的整底指数和：按词迹值聚类]\label{con:xi-terminal-replica-softcore-trace-moment-clustered-exponential-sum}
固定整数 \(m\ge 1\)，沿用结论 \ref{con:xi-terminal-replica-softcore-binary-necklace-trace-formula} 的 \(\tau(\omega)\)。
定义迹值集合与计数系数
\[
\mathcal A_m:=\{\tau(\omega):\omega\in\{0,1\}^m\}\subset\ZZ_{\ge 0},
\qquad
c_m(a):=\#\{\omega\in\{0,1\}^m:\tau(\omega)=a\}\in\ZZ_{\ge 0}.
\]
则对一切 \(q\ge 1\) 成立严格恒等式
\[
\boxed{\;
\Tr\!\bigl((T^{(q)})^{m}\bigr)=2^{-m}\sum_{a\in\mathcal A_m}c_m(a)\,a^q\; } .
\]
\end{conclusion}
\begin{proof}
由结论 \ref{con:xi-terminal-replica-softcore-binary-necklace-trace-formula} 的 word 求和式
\(
2^{m}\Tr((T^{(q)})^{m})=\sum_{\omega\in\{0,1\}^m}\tau(\omega)^q
\)，
按迹值 \(a=\tau(\omega)\) 分组即可得到所示聚类展开。证毕。
\end{proof}

\begin{conclusion}[奇素长度处的循环轨道整除律：循环平移作用下的系数刚性]\label{con:xi-terminal-replica-softcore-trace-moment-prime-orbit-divisibility}
令 \(m\ge 3\) 为奇素数，并沿用结论 \ref{con:xi-terminal-replica-softcore-trace-moment-clustered-exponential-sum} 的 \(\mathcal A_m\) 与 \(c_m(a)\)。
则恒有
\[
c_m(2^m)=1,\qquad c_m(L_m)=1,
\]
并且对一切 \(a\in\mathcal A_m\setminus\{2^m,L_m\}\) 成立整除律
\[
m\mid c_m(a).
\]
此外，在例外幂和的全阶分解中（结论 \ref{con:xi-terminal-replica-softcore-partition-exceptional-closed-form}），除首项 \(2^{mq}\) 与替换通道 \(\Omega_m(q)\) 外，其余纯指数项 \(a^q\) 的系数亦全部被 \(m\) 整除。
\end{conclusion}
\begin{proof}
令 \(\sigma\) 为长度 \(m\) 字词的循环平移（把首位移至末位）。
由迹的循环不变性 \(\Tr(X_1\cdots X_m)=\Tr(X_2\cdots X_mX_1)\)，对任意 \(\omega\in\{0,1\}^m\) 有
\(
\tau(\sigma(\omega))=\tau(\omega)
\)。
因此 \(c_m(a)\) 可按循环轨道分解为
\(
c_m(a)=\sum_{\mathcal O}\#\mathcal O
\)，其中求和遍历所有满足 \(\tau\equiv a\) 的 \(\langle\sigma\rangle\)-轨道。
\par
由于 \(m\) 为素数，轨道大小只能为 \(1\) 或 \(m\)，而轨道大小为 \(1\) 的情形仅对应恒定字词 \(0^m\) 与 \(1^m\)。
并且 \(\tau(1^m)=\Tr(J^m)=2^m\)。
若 \(\omega\neq 1^m\)，则 \(\omega\) 含至少一个 \(0\)；令 \(r\) 为其中 \(0\) 的个数，则 \(r\ge 1\)。
由结论 \ref{con:xi-terminal-replica-softcore-fixed-zero-count-trace-extremal}（取 \(r\ge 1\)）得
\(
\tau(\omega)\le F_{r+3}2^{m-r-1}\le F_4 2^{m-2}=3\cdot 2^{m-2}<2^m
\)，
故 \(2^m\) 仅由 \(\omega=1^m\) 取到，从而 \(c_m(2^m)=1\)。
\par
另一方面，\(\tau(0^m)=\Tr(K^m)=L_m\)。
若 \(\omega\neq 0^m\)，由结论 \ref{con:xi-terminal-replica-softcore-trace-partition-image} 知 \(\tau(\omega)=B_\lambda\)（某个 \(\lambda\vdash m\)），其中 \(B_\lambda=\prod_i F_{\lambda_i+2}\)。
对 \(a,b\ge 1\)，由 Fibonacci 加法公式与递推 \(F_{n+2}=F_{n+1}+F_n\) 可得
\[
F_{a+2}F_{b+2}=F_{a+b+2}+F_aF_b>F_{a+b+2},
\]
从而对任意含至少两个部件的分拆 \(\lambda\vdash m\) 有 \(B_\lambda>F_{m+2}\)，而单部件分拆 \((m)\) 给出 \(B_{(m)}=F_{m+2}\)。
故总有 \(\tau(\omega)=B_\lambda\ge F_{m+2}\)。
又
\(
F_{m+2}-L_m=(F_{m+1}+F_m)-(F_{m+1}+F_{m-1})=F_{m-2}>0
\)
（因 \(m\ge 3\)），于是 \(\tau(\omega)>L_m\)。
这表明 \(L_m\) 仅由 \(\omega=0^m\) 取到，从而 \(c_m(L_m)=1\)。
\par
综上，除 \(a\in\{2^m,L_m\}\) 外，其余迹值的计数皆为若干个大小为 \(m\) 的轨道之和，故被 \(m\) 整除。
\par
对例外幂和，结论 \ref{con:xi-terminal-replica-softcore-exceptional-power-sum-word-trace-sympower} 给出
\(
2^mS_m(q)=\Omega_m(q)+\sum_{w\neq K^m}(\Tr M(w))^q
\)；
其中 \(K^m\) 的词迹通道被 \(\Omega_m(q)\) 替换，而 \(J^m\) 的通道仍给出首项 \(2^{mq}\)。
其余非恒定字词在循环平移下形成大小为 \(m\) 的轨道并保持 \(\Tr M(w)\) 不变，因此对应指数项的系数必被 \(m\) 整除。证毕。
\end{proof}

\endinput



\begin{conclusion}[\texorpdfstring{$S_q$}{Sq}-不变子空间上的对称幂模型：例外谱等价于 \texorpdfstring{$\frac12(\mathrm{Sym}^q(J)+\mathrm{Sym}^q(K))$}{1/2(Sym^q(J)+Sym^q(K))} 的谱]\label{con:xi-terminal-replica-softcore-exceptional-spectrum-sympower-model}
令 \(\mathcal H_q:=(\CC^2)^{\otimes q}\)，并以标准基 \(e_0,e_1\in\CC^2\) 对 \(u\in U_q=\{0,1\}^q\) 记
\[
e_u:=e_{u_1}\otimes\cdots\otimes e_{u_q},
\]
从而得到自然同构 \(\CC^{U_q}\cong \mathcal H_q\)。
令 \(S_q\) 通过置换张量因子作用于 \(\mathcal H_q\)，记不变子空间为 \(\mathcal H_q^{S_q}\)。
则存在自然同构
\[
\mathcal H_q^{S_q}\cong \mathrm{Sym}^q(\CC^2),
\]
并且在该同构下有算子等价
\[
T^{(q)}\big|_{\mathcal H_q^{S_q}}
\ \cong\
\frac12\Bigl(\mathrm{Sym}^q(J)+\mathrm{Sym}^q(K)\Bigr).
\]
因此，例外谱 \(\{\lambda_i^{(q)}\}_{i=0}^{q}\) 精确等同于算子
\(\tfrac12(\mathrm{Sym}^q(J)+\mathrm{Sym}^q(K))\) 在 \(\mathrm{Sym}^q(\CC^2)\) 上的谱（含重数）。
\end{conclusion}
\begin{proof}[证明要点]
由对称化投影的通用性质，\(\mathcal H_q^{S_q}\) 与 \(\mathrm{Sym}^q(\CC^2)\) 存在自然同构。
沿用结论 \ref{con:xi-terminal-replica-softcore-binary-necklace-trace-formula} 的张量分解
\(
T^{(q)}=\tfrac12(J^{\otimes q}+K^{\otimes q})
\)；
由于对称幂表示与张量幂的 \(S_q\)-对称化操作相容，\(J^{\otimes q}\) 与 \(K^{\otimes q}\) 在
\(\mathcal H_q^{S_q}\) 上的诱导作用分别等价于 \(\mathrm{Sym}^q(J)\) 与 \(\mathrm{Sym}^q(K)\)。
谱等价由相似变换不改变特征值给出。证毕。
\end{proof}
\begin{conclusion}[非交换字词展开：\texorpdfstring{$2^m S_m(q)$}{2^m S_m(q)} 的对称幂迹公式与行列式二分]\label{con:xi-terminal-replica-softcore-exceptional-power-sum-word-trace-sympower}
对任意整数 \(m\ge 1\) 与 \(q\ge 2\)，记二字母自由半群 \(\{J,K\}^m\) 中的字词
\(w=A_1A_2\cdots A_m\)（\(A_t\in\{J,K\}\)），并记其对应的 \(2\times 2\) 乘积
\[
M(w):=A_1A_2\cdots A_m\in \Mat_{2}(\ZZ).
\]
则例外幂和满足无谱参数恒等式
\[
\boxed{\;
2^m S_m(q)=\sum_{w\in\{J,K\}^m}\Tr\!\left(\mathrm{Sym}^q(M(w))\right)\;}
\]
并且具有严格的行列式二分：
\begin{enumerate}
  \item 若 \(w\neq K^m\)，则 \(\det M(w)=0\)，从而 \(M(w)\) 为秩 \(1\) 矩阵，且
  \[
  \Tr\!\left(\mathrm{Sym}^q(M(w))\right)=\bigl(\Tr M(w)\bigr)^q.
  \]
  \item 若 \(w=K^m\)，则 \(\det(K^m)=(-1)^m\neq 0\)，并且
  \[
  \Tr\!\left(\mathrm{Sym}^q(K^m)\right)=\Omega_m(q),
  \]
  其中 \(\Omega_m(q)\) 定义见结论 \ref{con:xi-terminal-replica-softcore-lucas-pell-omega-unified}。
\end{enumerate}
进一步地，若 \(w\neq K^m\)，令 \(k\) 为 \(w\) 中字母 \(J\) 的出现次数；则在循环移位意义下存在分解
\[
w\sim JK^{r_1}JK^{r_2}\cdots JK^{r_k},
\qquad
r_i\ge 0,\ \ \sum_{i=1}^{k}r_i=m-k.
\]
记 \(\lambda(w)\vdash m\) 为 \((r_1+1,\dots,r_k+1)\) 的分拆化（按降序重排），则
\[
\Tr M(w)=B_{\lambda(w)},
\qquad
\mathrm{Spec}(M(w))=\{B_{\lambda(w)},0\},
\qquad
\Tr\!\left(\mathrm{Sym}^q(M(w))\right)=B_{\lambda(w)}^{\,q}\qquad(q\ge 0),
\]
其中 \(B_\lambda\) 由结论 \ref{con:xi-terminal-replica-softcore-trace-partition-image} 定义。
并且对固定分拆 \(\lambda\vdash m\)，满足 \(\lambda(w)=\lambda\) 的字词个数严格等于分拆系数 \(c_{m,\lambda}\)
（结论 \ref{con:xi-terminal-replica-softcore-partition-trace-closed-form}）。
因此
\[
2^m S_m(q)=\Omega_m(q)+\sum_{\substack{w\in\{J,K\}^m\\ w\neq K^m}}\bigl(\Tr M(w)\bigr)^q\in\ZZ,
\]
且当把 \(q\) 视为自由变量时，\(\Omega_m(q)\) 是右端唯一不属于纯指数项 \(a^q\) 类型的贡献通道。
\end{conclusion}
\begin{proof}[证明要点]
由结论 \ref{con:xi-terminal-replica-softcore-exceptional-spectrum-sympower-model}，
例外谱对应算子为 \(\tfrac12(\mathrm{Sym}^q(J)+\mathrm{Sym}^q(K))\) 作用在 \(\mathrm{Sym}^q(\CC^2)\) 上，
故
\(
2^m S_m(q)=\Tr\bigl((\mathrm{Sym}^q(J)+\mathrm{Sym}^q(K))^m\bigr)
\)。
展开非交换二项式并用 \(\mathrm{Sym}^q\) 的函子性
\(\mathrm{Sym}^q(A_1)\cdots \mathrm{Sym}^q(A_m)=\mathrm{Sym}^q(A_1\cdots A_m)\)
即得字词求和式。
行列式二分由 \(\det J=0\)、\(\det K=-1\) 与 \(\det\) 乘性给出。
秩 \(1\) 情形下 \(M(w)\) 的特征值为 \(\Tr M(w)\) 与 \(0\)，从而
\(\Tr(\mathrm{Sym}^q(M(w)))=(\Tr M(w))^q\)。
最后一式把 \(w=K^m\) 的贡献识别为 \(\Omega_m(q)\)（结论 \ref{con:xi-terminal-replica-softcore-lucas-pell-omega-unified}）。证毕。
\par
分拆参数化与计数结论由 \(\Tr M(w)\) 的分拆像刻画（结论 \ref{con:xi-terminal-replica-softcore-trace-partition-image}）与分拆系数展开（结论 \ref{con:xi-terminal-replica-softcore-partition-trace-closed-form}）给出。
\end{proof}
\begin{conclusion}[\texorpdfstring{$\{J,K\}$}{\{J,K\}}-字词迹的块化 Fibonacci 因子分解]\label{con:xi-terminal-replica-softcore-jk-word-trace-fib-block-factorization}
固定整数 \(m\ge 1\)。
取任意字词 \(w\in\{J,K\}^m\) 且含至少一个 \(J\)，并记 \(M(w)\) 如结论 \ref{con:xi-terminal-replica-softcore-exceptional-power-sum-word-trace-sympower}。
则 \(\det M(w)=0\)，且 \(M(w)\) 为秩 \(1\) 矩阵。
\par
更精确地，将相邻同字母合并后，\(w\) 具有唯一分块表示
\[
w=K^{a_0}J^{b_1}K^{a_1}J^{b_2}\cdots J^{b_r}K^{a_r},
\qquad
r\ge 1,\ a_0,\dots,a_r\in\ZZ_{\ge 0},\ b_1,\dots,b_r\in\ZZ_{\ge 1}.
\]
令 Fibonacci 数列满足 \(F_0=0,F_1=1\) 与递推 \(F_{n+2}=F_{n+1}+F_n\)。
则迹满足严格因子分解
\[
\boxed{\;
\Tr M(w)
=2^{(b_1+\cdots+b_r)-r}\,F_{a_0+a_r+3}\prod_{i=1}^{r-1}F_{a_i+3}\;} ,
\]
从而 \(M(w)\) 的特征值集合为 \(\mathrm{Spec}(M(w))=\{\Tr M(w),0\}\)。
\end{conclusion}
\begin{proof}
设 \(u=(1,1)^{\top}\)，则 \(J=uu^{\top}\)，从而 \(J^2=2J\) 并归纳得
\(
J^{b}=2^{b-1}J
\)（\(b\ge 1\)）。
因此
\[
M(w)=2^{(b_1+\cdots+b_r)-r}\cdot K^{a_0}J K^{a_1}J\cdots J K^{a_r}.
\]
由 \(K\) 的 Fibonacci 矩阵恒等式
\[
K^n=
\begin{pmatrix}
F_{n+1}&F_n\\[2pt]
F_n&F_{n-1}
\end{pmatrix}\quad(n\ge 1),
\qquad
K^0=I,
\]
可得 \(u^{\top}K^n u=F_{n+3}\)（\(n\ge 0\)）。
注意到 \(K^{a}J=(K^{a}u)u^{\top}\) 与 \(JK^{a}=u(u^{\top}K^{a})\) 皆为秩 \(1\) 矩阵，
从而
\[
K^{a_0}J K^{a_1}J\cdots J K^{a_r}
=(K^{a_0}u)\Bigl(\prod_{i=1}^{r-1}u^{\top}K^{a_i}u\Bigr)(u^{\top}K^{a_r}),
\]
其迹为
\[
\Tr\!\bigl(K^{a_0}J K^{a_1}J\cdots J K^{a_r}\bigr)
=(u^{\top}K^{a_r}K^{a_0}u)\prod_{i=1}^{r-1}(u^{\top}K^{a_i}u)
=F_{a_0+a_r+3}\prod_{i=1}^{r-1}F_{a_i+3}.
\]
乘回标量因子即得所示因子分解。
由于 \(\det J=0\)，且 \(w\) 含 \(J\)，故 \(\det M(w)=0\)；同时上式将 \(M(w)\) 表为一对列/行向量的外积，故 \(\mathrm{rank}(M(w))=1\)。证毕。
\end{proof}
\begin{conclusion}[例外幂和的线性字词全闭式：Fibonacci 乘法核与 dyadic 整性]\label{con:xi-terminal-replica-softcore-exceptional-power-sum-linear-word-closed-form}
对任意整数 \(m\ge 1\) 与 \(q\ge 2\)，沿用 \(\Omega_m(q)\)（结论 \ref{con:xi-terminal-replica-softcore-lucas-pell-omega-unified}）。
则
\[
\boxed{\;
S_m(q)
:=2^{-m}\Bigg(
\Omega_m(q)
\;+\;
\sum_{k=1}^{m}\ \sum_{\substack{a_0,\dots,a_k\ge 0\\ a_0+\cdots+a_k=m-k}}
\Bigl(F_{a_0+a_k+3}\prod_{i=1}^{k-1}F_{a_i+3}\Bigr)^q
\Bigg)\;}
\]
其中内层求和遍历所有线性分解字词
\[
w=K^{a_0}JK^{a_1}J\cdots JK^{a_k},
\qquad
\sum_{i=0}^{k}a_i=m-k.
\]
特别地，\(2^m S_m(q)\in\ZZ\)，且分子由上述 Fibonacci 乘法核的指数族与 \(\Omega_m(q)\) 显式给出。
\end{conclusion}
\begin{proof}[证明要点]
对 \(w\neq K^m\)，按首个 \(J\) 的位置唯一写成所示线性分解。
该情形可视为结论 \ref{con:xi-terminal-replica-softcore-jk-word-trace-fib-block-factorization} 在 \(b_1=\cdots=b_k=1\) 的特例；从而亦可直接读出
\[
\Tr M(w)=F_{a_0+a_k+3}\prod_{i=1}^{k-1}F_{a_i+3}.
\]
设 \(u=(1,1)^{\top}\)，则 \(J=uu^{\top}\)，并由结论 \ref{con:xi-terminal-replica-softcore-binary-necklace-trace-formula} 的计算有
\(u^{\top}K^{r}u=F_{r+3}\)。
因此
\[
\Tr M(w)=\Tr\!\bigl(K^{a_0}J K^{a_1}J\cdots J K^{a_k}\bigr)
=F_{a_0+a_k+3}\prod_{i=1}^{k-1}F_{a_i+3}.
\]
将其代入结论 \ref{con:xi-terminal-replica-softcore-exceptional-power-sum-word-trace-sympower} 的
\(
2^mS_m(q)=\Omega_m(q)+\sum_{w\neq K^m}(\Tr M(w))^q
\)
并按 \(k\) 与 \((a_0,\dots,a_k)\) 分组即得所示闭式。证毕。
\end{proof}

\begin{conclusion}[全迹矩的线性字词闭式：\texorpdfstring{$\Tr((T^{(q)})^m)$}{Tr((T^(q))^m)} 的纯指数分解]\label{con:xi-terminal-replica-softcore-trace-moment-linear-word-closed-form}
对任意整数 \(m\ge 1\) 与 \(q\ge 2\)，沿用结论 \ref{con:xi-terminal-replica-softcore-exceptional-power-sum-word-trace-sympower} 的字词与乘积 \(M(w)\)。
则全迹矩满足严格闭式
\[
\boxed{\;
\Tr\!\bigl((T^{(q)})^m\bigr)
=2^{-m}\Bigg(
L_m^{q}
\;+\;
\sum_{k=1}^{m}\ \sum_{\substack{a_0,\dots,a_k\ge 0\\ a_0+\cdots+a_k=m-k}}
\Bigl(F_{a_0+a_k+3}\prod_{i=1}^{k-1}F_{a_i+3}\Bigr)^{q}
\Bigg)\; } .
\]
\end{conclusion}
\begin{proof}[证明要点]
由结论 \ref{con:xi-terminal-replica-softcore-binary-necklace-trace-formula} 的张量分解
\(
T^{(q)}=\tfrac12(J^{\otimes q}+K^{\otimes q})
\)
与 \((AB)^{\otimes q}=A^{\otimes q}B^{\otimes q}\)，展开得
\(
(T^{(q)})^m=2^{-m}\sum_{w\in\{J,K\}^m}M(w)^{\otimes q}
\)；
再由 \(\Tr(M^{\otimes q})=(\Tr M)^q\) 得
\[
\Tr\!\bigl((T^{(q)})^m\bigr)=2^{-m}\sum_{w\in\{J,K\}^m}\bigl(\Tr M(w)\bigr)^q.
\]
其中 \(w=K^m\) 的贡献为 \(2^{-m}L_m^q\)。
若 \(w\neq K^m\)，则 \(w\) 含至少一个 \(J\)，并可唯一写为
\(
w=K^{a_0}JK^{a_1}J\cdots JK^{a_k}
\)
且 \(\sum_{i=0}^{k}a_i=m-k\)；
由结论 \ref{con:xi-terminal-replica-softcore-jk-word-trace-fib-block-factorization}（取 \(b_1=\cdots=b_k=1\)）得
\(
\Tr M(w)=F_{a_0+a_k+3}\prod_{i=1}^{k-1}F_{a_i+3}
\)。
将上述分解代回并按 \(k\) 与 \((a_0,\dots,a_k)\) 分组即得所示闭式。证毕。
\end{proof}
\begin{conclusion}[\texorpdfstring{$m$}{m}-方向生成函数：Fibonacci 指数核的单极点结构与 Lucas--Pell 极点簇]\label{con:xi-terminal-replica-softcore-exceptional-power-sum-m-direction-genfunc}
固定 \(q\ge 2\)。
定义形式幂级数
\[
\mathscr A_q(z):=\sum_{r\ge 0}F_{r+3}^{\,q} z^{r},
\qquad
\mathscr B_q(z):=\frac{\dd}{\dd z}\bigl(z\mathscr A_q(z)\bigr)
:=\sum_{r\ge 0}(r+1)F_{r+3}^{\,q} z^{r}.
\]
则秩一字词贡献的生成函数满足严格恒等式
\[
\sum_{m\ge 1}\left(\sum_{\substack{w\in\{J,K\}^m\\ w\neq K^m}}\bigl(\Tr M(w)\bigr)^q\right) z^{m}
:=\frac{z\,\mathscr B_q(z)}{1-z\mathscr A_q(z)}.
\]
等价地，令
\[
A_q(z):=z\mathscr A_q(z)=\sum_{j\ge 1}F_{j+2}^{\,q} z^{j},
\qquad
A_q'(z)=\mathscr B_q(z),
\]
则
\[
\sum_{m\ge 1}\bigl(2^m S_m(q)-\Omega_m(q)\bigr)\,z^m
=\frac{z\,A_q'(z)}{1-A_q(z)}.
\]
从而
\[
\sum_{m\ge 1}2^m S_m(q)\,z^m
:=\frac{z\,\mathscr B_q(z)}{1-z\mathscr A_q(z)}
\;+\;
\sum_{m\ge 1}\Omega_m(q)\,z^m.
\]
此外，若将 \(\Omega_0(q):=q+1\) 作为自然延拓，则
\[
\sum_{m\ge 0}\Omega_m(q)\,z^m
:=\sum_{j=0}^{q}\frac{1}{1-z\,(-1)^j\varphi^{\,q-2j}}.
\]
\end{conclusion}
\begin{proof}[证明要点]
由结论 \ref{con:xi-terminal-replica-softcore-exceptional-power-sum-linear-word-closed-form}，
对固定 \(k\ge 1\) 的线性分解 \(w=K^{a_0}J\cdots JK^{a_k}\) 有
\(
(\Tr M(w))^q=\bigl(F_{a_0+a_k+3}\bigr)^q\prod_{i=1}^{k-1}F_{a_i+3}^{\,q}
\)
且 \(m=k+\sum_{i=0}^{k}a_i\)。
对端点对 \((a_0,a_k)\) 以 \(r=a_0+a_k\) 重新参数化，并注意满足 \(a_0+a_k=r\) 的非负整对数为 \(r+1\)，
得到端点生成函数为 \(\mathscr B_q(z)\)，中间块贡献为 \(\mathscr A_q(z)^{k-1}\)，而每个 \(J\) 贡献一因子 \(z\)。
对所有 \(k\ge 1\) 求和得到
\(
z\mathscr B_q(z)\sum_{k\ge 1}(z\mathscr A_q(z))^{k-1}=\frac{z\mathscr B_q(z)}{1-z\mathscr A_q(z)}
\)。
第二条恒等式由结论 \ref{con:xi-terminal-replica-softcore-exceptional-power-sum-word-trace-sympower} 的
\(2^mS_m(q)=\Omega_m(q)+\sum_{w\neq K^m}(\Tr M(w))^q\) 逐项求和得到。
最后，\(\Omega_m(q)=\Tr(\mathrm{Sym}^q(K^m))\)，而 \(\mathrm{Sym}^q(K)\) 的谱为
\(\{(-1)^j\varphi^{q-2j}\}_{j=0}^{q}\)，故按几何级数求和得到部分分式表示。证毕。
\end{proof}
\begin{remark}[文献对照]\label{rem:xi-terminal-replica-softcore-word-trace-literature-contrast}
关于 Fibonacci/Lucas 矩阵恒等式与秩一因子分解 \(J=uu^{\top}\) 的标准推导，可参见 \cite{FQFibonacciLucasChap11Matrices,MSEFibonacciMatrices784710} 等资料。
本段把这些支点组织为一条可离线核验的指数展开链：由 \(\{J,K\}^m\) 的字词迹因子分解把 \(2^mS_m(q)\) 与 \(\Tr((T^{(q)})^m)\) 在 \(q\)-方向压缩为有限多底数指数族，并进一步给出素数长度处的循环轨道整除刚性与低阶 \(m\) 的显式聚类闭式。
\end{remark}
上述线性字词口径与 \(m\)-方向生成函数可视为对结论 \ref{con:xi-terminal-replica-softcore-binary-necklace-trace-formula} 的非循环展开；
进一步将字词按循环移位归并，可得到以下 necklace 剥离律。
\endinput


\begin{conclusion}[二元 necklace 剥离律：例外幂和 \texorpdfstring{$S_m(q)$}{S_m(q)} 的全阶闭式]\label{con:xi-terminal-replica-softcore-exceptional-power-sum-necklace-peeling}
对任意整数 \(m\ge 1\) 与 \(q\ge 2\)，沿用 \(\Omega_m(q)\)（结论 \ref{con:xi-terminal-replica-softcore-lucas-pell-omega-unified}）与
\(\mathrm{Neck}_m^{(2)}\)、\(\tau(\eta)\)（结论 \ref{con:xi-terminal-replica-softcore-binary-necklace-trace-formula}）的记号。
则存在严格恒等式
\[
\boxed{\;
S_m(q)
=2^{-m}\left(
\Omega_m(q)
\;+\!\!\!\sum_{\eta\in \mathrm{Neck}_m^{(2)}\setminus\{0^m\}}\#\mathrm{orb}(\eta)\,\tau(\eta)^{q}
\right)\; } .
\]
并且 \(\Omega_m(q)\) 可等价写为
\[
\Omega_m(q)=\Tr(B_q^{\,m}),
\qquad
B_q(i,j)=\binom{q-j}{i}\qquad(0\le i,j\le q),
\]
其中 \(B_q\) 为不交算子在 \(S_q\)-对称子空间上的压缩矩阵（结论 \ref{con:xi-terminal-disjointness-bq-sympower-unimodular-inverse}）。
该分解表明：例外幂和的全部“非 Lucas--Pell 部分”由含 \(1\) 的二元 necklace 族贡献，其权重完全由零段长度的 Fibonacci 乘法形 \(\tau(\eta)=\prod_i F_{r_i+3}\) 控制；而 \(\Omega_m\) 精确承担全零 necklace 通道的谱替换。
\end{conclusion}
\begin{proof}[证明要点]
由定理 \ref{thm:pom-replica-softcore-characteristic-factorization}，
\(
\Tr((T^{(q)})^m)=S_m(q)+\sum_{j=0}^{q}(\binom{q}{j}-1)\bigl(\frac{(-1)^j\varphi^{q-2j}}{2}\bigr)^m
\)。
对固定谱簇幂和作二项式求和与等比求和可得
\(
\sum_{j}(\binom{q}{j}-1)\bigl(\frac{(-1)^j\varphi^{q-2j}}{2}\bigr)^m
=2^{-m}\bigl(L_m^{q}-\Omega_m(q)\bigr)
\)
（其中 \(L_m=\varphi^m+(-1)^m\varphi^{-m}\)）。
将其代回并与结论 \ref{con:xi-terminal-replica-softcore-binary-necklace-trace-formula} 的
\(
2^{m}\Tr((T^{(q)})^m)=\sum_{\eta}\#\mathrm{orb}(\eta)\tau(\eta)^q
\)
合并，恰好消去全零 necklace 项 \(2^{-m}L_m^q\)，得到所示剥离律。
最后，\(\Omega_m(q)=\Tr(B_q^{\,m})\) 由 \(B_q=\mathrm{Sym}^q(K)\)（结论 \ref{con:xi-terminal-disjointness-bq-sympower-unimodular-inverse}）与 \(\mathrm{Sym}^q\) 的函子性给出：\(B_q^{\,m}=\mathrm{Sym}^q(K^m)\)，取迹即与结论 \ref{con:xi-terminal-replica-softcore-lucas-pell-omega-unified} 的闭式一致。
\end{proof}

\begin{conclusion}[对称化缺陷项的精确表达]\label{con:xi-terminal-replica-softcore-trace-exceptional-difference}
对任意整数 \(m\ge 1\) 与 \(q\ge 2\)，成立严格恒等式
\[
\boxed{\;
\Tr\!\bigl((T^{(q)})^m\bigr)-S_m(q)=2^{-m}\Bigl(L_m^{q}-\Omega_m(q)\Bigr)\; } .
\]
因此全迹矩与例外幂和的差值由单一缺陷项 \(L_m^{q}-\Omega_m(q)\) 完全刻画。
\end{conclusion}
\begin{proof}[证明要点]
由定理 \ref{thm:pom-replica-softcore-characteristic-factorization} 的谱分解，
\(
\Tr((T^{(q)})^m)=S_m(q)+\sum_{j=0}^{q}(\binom{q}{j}-1)\bigl(\frac{(-1)^j\varphi^{q-2j}}{2}\bigr)^m
\)。
并由上一个证明中的二项式与等比求和化简得到固定谱簇幂和为 \(2^{-m}(L_m^q-\Omega_m(q))\)，代回即得结论。证毕。
\end{proof}

\begin{conclusion}[例外幂和的分拆闭式：Fibonacci--partition 的全指数分解]\label{con:xi-terminal-replica-softcore-partition-exceptional-closed-form}
固定整数 \(m\ge 1\) 与 \(q\ge 2\)。
沿用结论 \ref{con:xi-terminal-replica-softcore-partition-trace-closed-form} 的 \(B_\lambda\) 与 \(c_{m,\lambda}\)。
则成立严格闭式
\[
\boxed{\;
S_m(q)
=\frac1{2^m}\left(\Omega_m(q)+\sum_{\lambda\vdash m}c_{m,\lambda}\,B_\lambda^{q}\right)\; } .
\]
该分解把 \(q\)-方向的全部复杂性分离为两类互不混淆的通道：由分拆数据产生的纯整数指数族 \(\{B_\lambda^q\}\)，以及唯一的二阶 Lucas--Pell 隧穿项 \(\Omega_m(q)=\Tr(\mathrm{Sym}^q(K^m))\)。
\end{conclusion}
\begin{proof}[证明要点]
由结论 \ref{con:xi-terminal-replica-softcore-partition-trace-closed-form} 与结论 \ref{con:xi-terminal-replica-softcore-trace-exceptional-difference}，
\[
S_m(q)
=\Tr\!\bigl((T^{(q)})^m\bigr)-2^{-m}\bigl(L_m^{q}-\Omega_m(q)\bigr)
=\frac1{2^m}\left(\Omega_m(q)+\sum_{\lambda\vdash m}c_{m,\lambda}B_\lambda^q\right).
\]
对 \(\Omega_m\) 的对称幂刻画来自结论 \ref{con:xi-terminal-replica-softcore-exceptional-power-sum-word-trace-sympower}。证毕。
\end{proof}

\begin{conclusion}[\texorpdfstring{$q$}{q}-方向最小线性递推与可取特征多项式]\label{con:xi-terminal-replica-softcore-q-recurrence-characteristic-upper-bound}
固定 \(m\ge 1\)。定义
\[
\mathcal A_m:=\bigl\{I(C_m[S])\mid S\subsetneq E(C_m)\bigr\}\subset\ZZ_{\ge0},
\qquad
b_q:=2^mS_m(q)\quad(q\ge0).
\]
则 \(b_q\) 满足常系数线性递推，其特征多项式可取为
\[
\boxed{\;
P_m(x):=\bigl(x^2-L_mx+(-1)^m\bigr)\prod_{a\in\mathcal A_m}(x-a)\; } .
\]
等价地，若 \(E\) 为移位算子 \(E:b_q\mapsto b_{q+1}\)，则 \(P_m(E)\,b\equiv0\)。
\end{conclusion}
\begin{proof}[证明要点]
由结论 \ref{con:xi-terminal-replica-softcore-exceptional-power-sum-cycle-sector-replacement}，
\[
b_q=\Omega_m(q)+\sum_{S\subsetneq E(C_m)}I(C_m[S])^q.
\]
对任意常数 \(a\)，序列 \(a^q\) 被一次算子 \(E-a\) 湮灭；
而 \(\Omega_m(q)\) 满足
\(
\Omega_m(q)=L_m\Omega_m(q-1)-(-1)^m\Omega_m(q-2)
\)，
故被 \(E^2-L_mE+(-1)^m\) 湮灭。
于是 \(b_q\) 被上述因子乘积湮灭，结论成立。
\end{proof}
\begin{conclusion}[\texorpdfstring{$q$}{q}-方向有理型生成函数]\label{con:xi-terminal-replica-softcore-q-ordinary-genfunc-partial-fractions}
固定 \(m\ge 1\)，记 \(b_q:=2^mS_m(q)\)（\(q\ge0\)）。
则普通生成函数满足
\[
\boxed{\;
\sum_{q\ge0}b_qz^q
=\frac{1}{1-L_m z+(-1)^m z^2}
+\sum_{S\subsetneq E(C_m)}\frac{1}{1-I(C_m[S])\,z}\; } .
\]
从而 \(\sum_{q\ge0}S_m(q)z^q\) 为有理函数，其极点包含于
\[
\left\{\frac1a:\ a\in\mathcal A_m\right\}
\cup
\left\{z:\ 1-L_m z+(-1)^m z^2=0\right\}.
\]
\end{conclusion}
\begin{proof}[证明要点]
由结论 \ref{con:xi-terminal-replica-softcore-exceptional-power-sum-cycle-sector-replacement}，
\[
b_q=\Omega_m(q)+\sum_{S\subsetneq E(C_m)}I(C_m[S])^q.
\]
逐项求和并用
\(
\sum_{q\ge0}a^qz^q=(1-az)^{-1}
\)
得到
\[
\sum_{q\ge0}\sum_{S\subsetneq E(C_m)}I(C_m[S])^qz^q
=\sum_{S\subsetneq E(C_m)}\frac{1}{1-I(C_m[S])\,z}.
\]
另一方面，\(\Omega_m(q)\) 满足二阶递推
\(
\Omega_m(q)=L_m\Omega_m(q-1)-(-1)^m\Omega_m(q-2)
\)
且初值 \(\Omega_m(0)=1,\Omega_m(1)=L_m\)，故
\[
\sum_{q\ge0}\Omega_m(q)z^q=\frac{1}{1-L_m z+(-1)^m z^2}.
\]
两式相加即得结论。
\end{proof}

\begin{conclusion}[代数整性：例外特征值的二倍是整数矩阵谱]\label{con:xi-terminal-replica-softcore-exceptional-spectrum-double-integer-matrix}
固定整数 \(q\ge 2\)，并沿用命题 \ref{prop:pom-multiplicity-composition-sq-symmetry-reduction} 的例外谱 \(\{\lambda_i^{(q)}\}_{i=0}^{q}\subset\RR\)（亦即 \(A^{(q)}\) 的全部特征值）。
定义缩放谱
\[
\mu_i^{(q)}:=2\lambda_i^{(q)}\qquad(0\le i\le q),
\]
并令整数矩阵 \(M_q\in\ZZ^{(q+1)\times(q+1)}\) 为
\[
\boxed{\;
(M_q)_{k,l}:=\binom{q}{l}+\binom{q-k}{l},
\qquad 0\le k,l\le q.\;}
\]
则 \(\{\mu_i^{(q)}\}_{i=0}^{q}\) 恰为 \(M_q\) 的全部特征值（含重数）。
从而 \(\mu_i^{(q)}\) 全为代数整数，并且其首一特征多项式
\[
\chi_q(x):=\prod_{i=0}^{q}(x-\mu_i^{(q)})=\det(xI_{q+1}-M_q)\in\ZZ[x].
\]
\end{conclusion}
\begin{proof}[证明要点]
由定义 \eqref{eq:pom-multiplicity-composition-Aq-def}，
\(
A^{(q)}_{k,l}=\frac12\bigl(\binom{q}{l}+\binom{q-k}{l}\bigr)
\)，
故 \(M_q=2A^{(q)}\) 且 \(M_q\) 逐项为整数。
命题 \ref{prop:pom-multiplicity-composition-sq-symmetry-reduction} 表明 \(\{\lambda_i^{(q)}\}\) 为 \(A^{(q)}\) 的全部特征值，故 \(\{\mu_i^{(q)}\}\) 为 \(2A^{(q)}=M_q\) 的全部特征值。
整数矩阵的特征值为代数整数，且其特征多项式属于 \(\ZZ[x]\)，从而结论成立。
\end{proof}

\input{subsubsec__xi-time-protocol-conclusions-part31d-exceptional-galois}

\begin{conclusion}[例外谱的有限可审计重构：由 \texorpdfstring{$2\times 2$}{2x2} 字迹枚举确定整系数特征多项式]\label{con:xi-terminal-replica-softcore-exceptional-spectrum-offline-newton-reconstruction}
固定整数 \(q\ge 2\)，沿用 \(J,K\)（结论 \ref{con:xi-terminal-replica-softcore-binary-necklace-trace-formula}）与 \(M(w)\)（结论 \ref{con:xi-terminal-replica-softcore-exceptional-power-sum-word-trace-sympower}）。
对每个整数 \(m\ge 1\)，定义整数数据
\[
P_m(q):=\sum_{w\in\{J,K\}^{m}}\bigl(\Tr M(w)\bigr)^{q}\in\ZZ_{\ge 0}.
\]
等价地，若记 \(B^{(q)}:=2T^{(q)}\)（从而 \(B^{(q)}_{u,v}=2\) 当 \(u\wedge v=0\)，否则 \(B^{(q)}_{u,v}=1\)），则 \(P_m(q)=\Tr\!\bigl((B^{(q)})^{m}\bigr)\)。
再定义例外缩放谱的幂和
\[
p_m(q):=\sum_{i=0}^{q}\bigl(\mu_i^{(q)}\bigr)^{m}=2^{m}S_m(q)\in\ZZ.
\]
则对一切 \(m\ge 1\) 成立严格恒等式
\[
\boxed{\;
p_m(q)=P_m(q)-L_m^{q}+\Omega_m(q)\; } ,
\]
其中 \(L_m=\Tr(K^m)\) 且 \(\Omega_m(q)=\Tr(\mathrm{Sym}^q(K^m))\)（结论 \ref{con:xi-terminal-replica-softcore-lucas-pell-omega-unified}）。
因此，\(\chi_q(x)\) 的全部整系数可由有限个整数 \(\{P_m(q)\}_{1\le m\le q+1}\) 离线重构：先由上式得到 \(\{p_m(q)\}_{1\le m\le q+1}\)，再由 Newton 恒等式把幂和 \(\{p_m\}\) 反演为首一多项式 \(\chi_q(x)=\prod_{i=0}^{q}(x-\mu_i^{(q)})\) 的系数。
\end{conclusion}
\begin{proof}[证明要点]
由结论 \ref{con:xi-terminal-replica-softcore-binary-necklace-trace-formula} 得
\(
P_m(q)=2^m\Tr((T^{(q)})^m)
\)，
而结论 \ref{con:xi-terminal-replica-softcore-trace-exceptional-difference} 给出
\(
2^m\Tr((T^{(q)})^m)=2^mS_m(q)+L_m^q-\Omega_m(q)
\)。
移项即得所示恒等式 \(p_m(q)=2^mS_m(q)=P_m(q)-L_m^q+\Omega_m(q)\)。
最后，\(\chi_q\) 为次数 \(q+1\) 的首一多项式，Newton 恒等式以有限步将幂和 \(\{p_m(q)\}_{1\le m\le q+1}\) 唯一确定为 \(\chi_q\) 的整系数；结合结论 \ref{con:xi-terminal-replica-softcore-exceptional-spectrum-double-integer-matrix} 的 \(\chi_q\in\ZZ[x]\) 即得离线重构结论。证毕。
\end{proof}

\paragraph{例外块整数模型的可逆性与整格缺陷}
在结论 \ref{con:xi-terminal-replica-softcore-exceptional-spectrum-double-integer-matrix} 已把例外缩放谱
\(\{\mu_i^{(q)}\}\) 识别为整数矩阵 \(M_q\) 的谱之后，
下述命题给出 \(M_q^{-1}\) 的完全闭式、相应的秩一分解、行列式刚性与 Smith 标准形，
从而把“例外块”的整性障碍压缩为一个唯一的 \(2\)-挠缺陷。

\begin{theorem}[例外块整数模型的逆矩阵闭式：反三角二项核与唯一半整数缺陷]\label{thm:xi-exceptional-integer-model-Mq-inverse-closed-form}
固定整数 \(q\ge 2\)，并沿用结论 \ref{con:xi-terminal-replica-softcore-exceptional-spectrum-double-integer-matrix} 的整数矩阵 \(M_q\)：
\[
(M_q)_{k\ell}:=\binom{q}{\ell}+\binom{q-k}{\ell},\qquad 0\le k,\ell\le q,
\]
并约定 \(\binom{n}{r}=0\)（当 \(r<0\) 或 \(r>n\)）。
则 \(M_q\) 可逆，且其逆矩阵条目具有闭式
\[
\boxed{\;
(M_q^{-1})_{k\ell}=
\begin{cases}
-\frac12, & (k,\ell)=(0,0),\\[4pt]
(-1)^{k+\ell-q}\binom{k}{k+\ell-q}, & k+\ell\ge q,\ (k,\ell)\neq(0,0),\\[4pt]
0, & k+\ell<q.
\end{cases}\;}
\]
\end{theorem}
\begin{proof}[证明要点]
令 \(N_q\) 为右端所给矩阵。
对固定 \((k,\ell)\)，由 \(N_{j\ell}=0\) 当 \(j<q-\ell\)，有
\[
(M_qN_q)_{k\ell}
=\sum_{j=q-\ell}^{q}\Bigl(\binom{q}{j}+\binom{q-k}{j}\Bigr)(-1)^{j+\ell-q}\binom{j}{j+\ell-q}
\;+\;\mathbf 1_{\{k=\ell=0\}}\Bigl(-\binom{q}{0}-\binom{q}{0}\Bigr)\frac12.
\]
利用恒等式 \(\binom{n}{j}\binom{j}{r}=\binom{n}{r}\binom{n-r}{j-r}\) 与二项反演
\(\sum_{t=0}^{n}(-1)^t\binom{n}{t}=\mathbf 1_{\{n=0\}}\)，
可将两段求和分别化为 \(\mathbf 1_{\{k=\ell\}}\) 与 \(-\mathbf 1_{\{k=\ell=0\}}\)，从而得到
\((M_qN_q)_{k\ell}=\delta_{k\ell}\)。
因此 \(N_q=M_q^{-1}\)。
\end{proof}
\leanverified{paper\_xi\_exceptional\_integer\_model\_Mq\_inverse\_closed\_form}

\begin{proposition}[结构分解：\texorpdfstring{$M_q$}{Mq} 与 \texorpdfstring{$M_q^{-1}$}{Mq^{-1}} 的 Pascal--反序变换秩一扰动]\label{prop:xi-exceptional-integer-model-Mq-rank1-decomposition}
令 \(P_q=[\binom{k}{\ell}]_{0\le k,\ell\le q}\) 为 Pascal 下三角矩阵，
令 \(R_q\) 为反序置换矩阵 \((R_q)_{k\ell}=\mathbf 1_{\{k+\ell=q\}}\)，
令 \(D_q=\mathrm{diag}(({-1})^{0},({-1})^{1},\dots,({-1})^{q})\)，
并记 \(E_{00}\) 为仅 \((0,0)\) 处为 \(1\) 的基矩阵。
再令 \(b_q:=(\binom{q}{0},\binom{q}{1},\dots,\binom{q}{q})^{\top}\) 与 \(\mathbf 1:=(1,\dots,1)^{\top}\in\ZZ^{q+1}\)。
则有完全分解
\[
\boxed{\;
M_q^{-1}=(-1)^q\,D_qP_qR_qD_q-\frac12\,E_{00}\;}
\]
以及等价的正向分解
\[
\boxed{\;
M_q=(-1)^q\,D_qR_qP_q^{-1}D_q\;+\;\mathbf 1\,b_q^{\top}.\;}
\]
\end{proposition}
\begin{proof}[证明要点]
由恒等式 \((P_qR_q)_{k\ell}=\binom{k}{q-\ell}\) 可知
\(\bigl(({-1})^qD_qP_qR_qD_q\bigr)_{k\ell}=(-1)^{k+\ell-q}\binom{k}{k+\ell-q}\)，
其恰与定理 \ref{thm:xi-exceptional-integer-model-Mq-inverse-closed-form} 在 \((k,\ell)\neq(0,0)\) 的条目一致，
而 \((0,0)\) 条目为 \(0\)，故第一式成立。
\par
对第二式，令 \(A:=(-1)^qD_qP_qR_qD_q\)。
由于 \(P_q\) 可逆且 \(R_q^2=I\)，有
\(
A^{-1}=(-1)^qD_qR_qP_q^{-1}D_q
\)。
再由 Sherman--Morrison 公式对秩一扰动 \(A-\frac12E_{00}\) 取逆，可得
\[
\bigl(A-\tfrac12E_{00}\bigr)^{-1}=A^{-1}+A^{-1}e_0e_0^{\top}A^{-1},
\]
其中 \(e_0\) 为第 \(0\) 个标准基向量。
直接计算可得 \(A^{-1}e_0=\mathbf 1\) 且 \(e_0^{\top}A^{-1}=b_q^{\top}\)，从而推出第二式。证毕。
\end{proof}

\begin{theorem}[行列式刚性与 Smith 标准形：余核恒为 \texorpdfstring{$\ZZ/2\ZZ$}{Z/2Z}]\label{thm:xi-exceptional-integer-model-Mq-det-snf}
\leanverified{paper\_xi\_exceptional\_integer\_model\_mq\_det\_snf}
对一切 \(q\ge 2\)，矩阵 \(M_q\) 的行列式满足
\[
\boxed{\;\det(M_q)=2\,(-1)^{\frac{q(q+1)}{2}}.\;}
\]
并且 \(M_q\) 的 Smith 标准形为
\[
\boxed{\;\mathrm{SNF}(M_q)=\mathrm{diag}(1,\dots,1,2),\;}
\]
从而余核恒同构于 \(\mathrm{coker}(M_q:\ZZ^{q+1}\to\ZZ^{q+1})\cong\ZZ/2\ZZ\)。
\end{theorem}
\begin{proof}[证明要点]
由命题 \ref{prop:xi-exceptional-integer-model-Mq-rank1-decomposition} 的正向分解，
令 \(B:=(-1)^qD_qR_qP_q^{-1}D_q\)，则 \(\det(B)=\det(R_q)=(-1)^{\frac{q(q+1)}{2}}\) 且 \(M_q=B+\mathbf 1 b_q^{\top}\)。
注意到 \(Be_0=\mathbf 1\)，故 \(B^{-1}\mathbf 1=e_0\)。
由矩阵行列式引理，
\[
\det(M_q)=\det(B)\bigl(1+b_q^{\top}B^{-1}\mathbf 1\bigr)
=\det(B)\bigl(1+b_q^{\top}e_0\bigr)
=2\det(B),
\]
即得行列式闭式。
\par
另一方面，由定理 \ref{thm:xi-exceptional-integer-model-Mq-inverse-closed-form}，
\(\mathrm{adj}(M_q)=\det(M_q)\,M_q^{-1}\) 为整数矩阵，并且其 \((0,0)\) 元为
\(\det(M_q)\cdot(-\tfrac12)=\pm1\)，
从而所有 \(1\times 1\) 余子式的最大公因子为 \(1\)。
结合 \(|\det(M_q)|=2\)，Smith 不变量只能为 \((1,\dots,1,2)\)，结论成立。证毕。
\end{proof}

\begin{corollary}[丢番图可解性的唯一模障碍：仅取决于首坐标奇偶]\label{cor:xi-exceptional-integer-model-Mq-diophantine-parity}
设 \(q\ge 2\) 且 \(b=(b_0,\dots,b_q)^{\top}\in\ZZ^{q+1}\)。
线性丢番图方程
\[
M_qx=b,\qquad x\in\ZZ^{q+1},
\]
有解当且仅当 \(b_0\) 为偶数；并且在可解时，解集为满秩格 \(M_q^{-1}\ZZ^{q+1}\) 的一个陪集。
\end{corollary}
\begin{proof}
由定理 \ref{thm:xi-exceptional-integer-model-Mq-inverse-closed-form}，解唯一地给为 \(x=M_q^{-1}b\)。
其中第 \(0\) 坐标满足
\[
x_0=(M_q^{-1})_{00}b_0+(M_q^{-1})_{0q}b_q=-\frac12 b_0+b_q,
\]
而其余坐标仅为整数线性组合。
故 \(x\in\ZZ^{q+1}\) 当且仅当 \(b_0\equiv 0\pmod 2\)。
陪集性质由定理 \ref{thm:xi-exceptional-integer-model-Mq-det-snf} 的满秩与 \(\mathrm{coker}(M_q)\cong\ZZ/2\ZZ\) 立得。证毕。
\end{proof}

\begin{corollary}[例外谱的整格压缩：非平凡分母集中于零层且仅发生一次]\label{cor:xi-exceptional-integer-model-Mq-single-denominator-defect}
设 \(q\ge 2\)，并记 \(\chi_q(x)=\det(xI_{q+1}-M_q)\in\ZZ[x]\) 的根为 \(\{\mu_i^{(q)}\}_{i=0}^{q}\)。
则：
\begin{enumerate}
  \item \(\prod_{i=0}^{q}\mu_i^{(q)}=\det(M_q)=2(-1)^{\frac{q(q+1)}{2}}\)，从而例外谱的乘积仅含一个素因子 \(2\)；
  \item \(2M_q^{-1}\in\Mat_{q+1}(\ZZ)\)，并且 \(2M_q^{-1}\) 的唯一奇数条目恰为 \((0,0)\) 元（等于 \(\pm1\)）。
\end{enumerate}
\end{corollary}
\begin{proof}
第 (1) 条由 \(\prod_i\mu_i^{(q)}=\chi_q(0)=\det(-M_q)\) 与定理 \ref{thm:xi-exceptional-integer-model-Mq-det-snf} 立即得到。
第 (2) 条由定理 \ref{thm:xi-exceptional-integer-model-Mq-inverse-closed-form} 的条目式直接读出：除 \((0,0)\) 外皆为整数，而 \((0,0)=-\tfrac12\)。证毕。
\end{proof}

\begin{conclusion}[加权互补核的闭式逆与分母支撑]\label{con:xi-exceptional-integer-model-weighted-complementary-kernel-inverse-denominator-support}
取整数 \(q\ge 2\) 与 \(s\in\NN_{\ge 1}\)，并定义整数矩阵 \(M_q^{(s)}\in\ZZ^{(q+1)\times(q+1)}\) 为
\[
\boxed{\;
\bigl(M_q^{(s)}\bigr)_{k\ell}:=\binom{q}{\ell}+s\,\binom{q-k}{\ell},
\qquad 0\le k,\ell\le q.\;}
\]
则 \(M_q^{(s)}\) 在 \(\QQ\) 上可逆，并满足闭式分解
\[
\boxed{\;
\bigl(M_q^{(s)}\bigr)^{-1}
\;=\;\frac1s\,P_q^{-1}R_q\;-\;\frac{1}{s(s+1)}\,E_{00}\; }.
\]
特别地，\(s\bigl(M_q^{(s)}\bigr)^{-1}\) 的全部条目均为整数，唯一可能引入额外分母的条目为 \((0,0)\) 元，并且该额外分母仅来自因子 \((s+1)\)。
\end{conclusion}
\begin{proof}
令 \(b_q:=(\binom{q}{0},\binom{q}{1},\dots,\binom{q}{q})^{\top}\) 与 \(\mathbf 1:=(1,\dots,1)^{\top}\in\ZZ^{q+1}\)。
由 \((R_qP_q)_{k\ell}=\binom{q-k}{\ell}\) 可写
\[
M_q^{(s)}=\mathbf 1\,b_q^{\top}+s\,R_qP_q.
\]
记 \(N_q:=R_qP_q\)，则 \(N_q\) 幺模且 \(N_q^{-1}=P_q^{-1}R_q\)。
左乘 \(N_q^{-1}\) 得到秩一扰动的标准形
\[
N_q^{-1}M_q^{(s)}=sI_{q+1}+u\,v^{\top},
\qquad
u:=N_q^{-1}\mathbf 1,\quad v:=b_q.
\]
注意到 \(e_q^{\top}P_q=b_q^{\top}\) 与 \(e_q^{\top}R_q=e_0^{\top}\)，故
\(
v^{\top}N_q^{-1}=b_q^{\top}P_q^{-1}R_q=e_0^{\top}
\)。
另一方面，\(N_qe_0=\mathbf 1\) 蕴含 \(u=N_q^{-1}\mathbf 1=e_0\)。
因此 \(v^{\top}u=e_0^{\top}e_0=1\)，由 Sherman--Morrison 公式得到
\[
\bigl(sI_{q+1}+uv^{\top}\bigr)^{-1}
\;=\;\frac1s I_{q+1}-\frac{1}{s(s+1)}\,u v^{\top}.
\]
两侧右乘 \(N_q^{-1}\) 并代入 \(u=e_0\) 与 \(v^{\top}N_q^{-1}=e_0^{\top}\) 即得
\[
\bigl(M_q^{(s)}\bigr)^{-1}
\;=\;\bigl(sI_{q+1}+uv^{\top}\bigr)^{-1}N_q^{-1}
\;=\;\frac1s\,P_q^{-1}R_q-\frac{1}{s(s+1)}\,E_{00}.
\]
最后的整性断言由 \(P_q^{-1}R_q\in\Mat_{q+1}(\ZZ)\) 与 \(E_{00}\) 的支撑性质直接推出。证毕。
\end{proof}

\begin{conclusion}[Smith 标准形与余核的完全分解]\label{con:xi-exceptional-integer-model-weighted-complementary-kernel-snf-cokernel}
在结论 \ref{con:xi-exceptional-integer-model-weighted-complementary-kernel-inverse-denominator-support} 的设定下，对任意 \(q\ge 2\)、\(s\ge 1\) 有
\[
\boxed{\;
\mathrm{SNF}\bigl(M_q^{(s)}\bigr)
\;=\;
\mathrm{diag}\Bigl(1,\underbrace{s,\dots,s}_{q-1\ \mathrm{times}},\,s(s+1)\Bigr).\;}
\]
等价地，
\[
\boxed{\;
\mathrm{coker}\!\bigl(M_q^{(s)}:\ZZ^{q+1}\to\ZZ^{q+1}\bigr)
\ \cong\ (\ZZ/s\ZZ)^{\oplus(q-1)}\ \oplus\ \ZZ/(s(s+1))\ZZ.\;}
\]
\end{conclusion}
\begin{proof}
由上一个证明中的幺模矩阵 \(N_q=R_qP_q\) 有
\[
N_q^{-1}M_q^{(s)}=sI_{q+1}+e_0\,b_q^{\top}.
\]
右乘幺模矩阵不改变余核同构型，故
\(
\mathrm{coker}(M_q^{(s)})\cong \mathrm{coker}(sI_{q+1}+e_0b_q^{\top})
\)。
进一步注意到 \(b_q^{\top}N_q^{-1}=e_0^{\top}\)，从而
\[
\bigl(sI_{q+1}+e_0b_q^{\top}\bigr)N_q^{-1}=sN_q^{-1}+e_0e_0^{\top}.
\]
因此 \(\mathrm{coker}(M_q^{(s)})\cong \mathrm{coker}(A)\)，其中
\[
A:=sI_{q+1}+\mathbf 1\,e_0^{\top}.
\]
设 \(\{e_i\}_{0\le i\le q}\) 为 \(\ZZ^{q+1}\) 的标准基并记 \(G:=\ZZ^{q+1}/A\ZZ^{q+1}\)。
由
\[
Ae_0=(s+1)e_0+\sum_{i=1}^{q}e_i,\qquad
Ae_i=s e_i\ (1\le i\le q),
\]
在 \(G\) 中得到关系
\[
s\,\bar e_i=0\ (1\le i\le q),\qquad (s+1)\bar e_0+\sum_{i=1}^{q}\bar e_i=0,
\]
其中 \(\bar e_i\) 表示 \(e_i\) 在 \(G\) 中的像。
令 \(H:=\langle \bar e_1,\dots,\bar e_q\rangle\)，则 \(H\cong(\ZZ/s\ZZ)^{\oplus q}\)。
记 \(t:=\sum_{i=1}^{q}\bar e_i\in H\)，则 \(t=-(s+1)\bar e_0\)，从而
\(
s(s+1)\bar e_0=-s t=0
\)，故 \(\langle \bar e_0\rangle\cong \ZZ/(s(s+1))\ZZ\)。
另一方面，\(H/\langle t\rangle\cong(\ZZ/s\ZZ)^{\oplus(q-1)}\)，并且 \(t\) 已由 \(\bar e_0\) 的倍数控制，从而
\[
G\cong \bigl(H/\langle t\rangle\bigr)\oplus \langle \bar e_0\rangle
\cong(\ZZ/s\ZZ)^{\oplus(q-1)}\oplus \ZZ/(s(s+1))\ZZ.
\]
由有限生成交换群结构定理读出 Smith 对角不变因子即得结论。证毕。
\end{proof}

\begin{corollary}[整数可解性的显式同余判据]\label{cor:xi-exceptional-integer-model-weighted-complementary-kernel-diophantine-congruence}
令 \(q\ge 2\)、\(s\ge 1\)，并取 \(b=(b_0,\dots,b_q)^{\top}\in\ZZ^{q+1}\)。
线性丢番图方程
\[
M_q^{(s)}x=b,\qquad x\in\ZZ^{q+1},
\]
存在整数解当且仅当同余条件
\[
\boxed{\;
b_0\equiv 0\pmod{s+1}
\quad\text{且}\quad
b_i\equiv \frac{b_0}{s+1}\pmod{s}\ \ (1\le i\le q)\;}
\]
同时成立。
\end{corollary}
\begin{proof}
由 \(M_q^{(s)}=A\,R_qP_q\)（其中 \(A=sI_{q+1}+\mathbf 1 e_0^{\top}\)）可知
\(
M_q^{(s)}x=b
\)
与
\(
Ay=b
\)
等价，其中 \(y:=R_qP_qx\in\ZZ^{q+1}\)（因为 \(R_qP_q\) 幺模）。
令 \(y=(y_0,\dots,y_q)^{\top}\)，则 \(Ay=b\) 逐坐标为
\[
(s+1)y_0=b_0,\qquad sy_i+y_0=b_i\ (1\le i\le q).
\]
因此可解当且仅当 \(b_0\equiv 0\pmod{s+1}\) 使 \(y_0=b_0/(s+1)\in\ZZ\)，并且对每个 \(i\ge 1\) 有
\(
b_i-y_0\equiv 0\pmod{s}
\)，即 \(b_i\equiv b_0/(s+1)\pmod{s}\)。
充分性由上述等式直接构造 \(y\) 得到。证毕。
\end{proof}

\begin{corollary}[\texorpdfstring{$p$}{p}-局部可逆性与坏素支撑]\label{cor:xi-exceptional-integer-model-weighted-complementary-kernel-localization-bad-primes}
令 \(p\) 为素数并取 \(q\ge 2\)、\(s\ge 1\)。
若 \(p\nmid s(s+1)\)，则 \(M_q^{(s)}\) 在 \(\ZZ_p\) 上可逆。
若 \(p\mid s\)，则 \(\mathrm{coker}(M_q^{(s)})\) 的 \(p\)-初等分量至少含有秩 \(q-1\) 的 \((\ZZ/p^{v_p(s)}\ZZ)\)-挠直和因子；
若 \(p\mid(s+1)\)，则除来自 \(s\) 的挠分量外，额外出现一个循环 \(p\)-挠因子，其阶为 \(p^{v_p(s+1)}\)（嵌入于 \(\ZZ/(s(s+1))\ZZ\) 的 \(p\)-分量）。
\end{corollary}
\begin{proof}
由结论 \ref{con:xi-exceptional-integer-model-weighted-complementary-kernel-snf-cokernel} 的余核分解，
\(
\mathrm{coker}(M_q^{(s)})\cong(\ZZ/s\ZZ)^{\oplus(q-1)}\oplus \ZZ/(s(s+1))\ZZ
\)，其挠素因子集合恰为 \(\pi(s(s+1))\)。
若 \(p\nmid s(s+1)\)，则在 \(\ZZ_p\) 上余核消失，从而矩阵可逆。
其余断言由直和分解与 \(p\)-初等分量的标准分解直接读出。证毕。
\end{proof}

\begin{remark}[审计脚本]\label{rem:xi-exceptional-integer-model-Mq-audit-script}
定理 \ref{thm:xi-exceptional-integer-model-Mq-inverse-closed-form}--推论 \ref{cor:xi-exceptional-integer-model-Mq-single-denominator-defect} 的矩阵恒等式与 Smith 标准形可由脚本 \texttt{scripts/exp\_xi\_exceptional\_integer\_model\_Mq\_audit.py} 在有限 \(q\) 范围内进行一致性核验。
\end{remark}
\begin{conclusion}[两参数 binomial 矩阵的条目级逆矩阵与行列式闭式]\label{con:xi-terminal-replica-softcore-mq-ab-inverse-det}
固定整数 \(q\ge 2\)，取参数 \(a,b\in\ZZ\) 并定义 \((q+1)\times(q+1)\) 整数矩阵
\[
M_q(a,b):=\bigl(m_{k\ell}\bigr)_{0\le k,\ell\le q},
\qquad
m_{k\ell}:=a\binom{q}{\ell}+b\binom{q-k}{\ell}.
\]
记 Pascal 矩阵 \(P_q=(\binom{k}{\ell})_{0\le k,\ell\le q}\)，反序置换矩阵 \(R_q\) 满足 \((R_qx)_k=x_{q-k}\)，并记 \(E_{00}\) 为仅 \((0,0)\) 元为 \(1\) 的基矩阵。
则当 \(b\neq 0\) 且 \(a+b\neq 0\) 时，\(M_q(a,b)\) 可逆，并且
\[
\boxed{\;
\det M_q(a,b)=(-1)^{\frac{q(q+1)}{2}}\,b^{q}(a+b)\;} .
\]
其逆矩阵满足完全条目级闭式
\[
\boxed{\;
\bigl(M_q(a,b)^{-1}\bigr)_{k\ell}
\;=\;
\frac{(-1)^{k+\ell-q}}{b}\binom{k}{k+\ell-q}\,\mathbf 1_{\{k+\ell\ge q\}}
\;-\;\frac{a}{b(a+b)}\,\mathbf 1_{\{k=\ell=0\}}\;} .
\]
特别地，当 \((a,b)=(1,1)\) 时有 \(M_q(1,1)=M_q\)（结论 \ref{con:xi-terminal-replica-softcore-exceptional-spectrum-double-integer-matrix}）。
\end{conclusion}
\begin{proof}
令 \(\beta:=(\binom{q}{0},\binom{q}{1},\dots,\binom{q}{q})^{\top}\) 与 \(\mathbf{1}:=(1,\dots,1)^{\top}\in\ZZ^{q+1}\)，并设
\(
N_q:=R_qP_q
\)。
则 \((N_q)_{k\ell}=\binom{q-k}{\ell}\) 且 \(N_qe_0=\mathbf{1}\)；又因 \(P_q\) 为下三角幺模矩阵，故 \(N_q\) 幺模且
\(
\det N_q=\det R_q=(-1)^{\frac{q(q+1)}{2}}
\)。
利用 \(m_{k\ell}\) 的分离结构，可写
\[
M_q(a,b)=a\,\mathbf{1}\beta^{\top}+b\,N_q.
\]
左乘 \(N_q^{-1}=P_q^{-1}R_q\) 得到秩一扰动的标准形
\[
N_q^{-1}M_q(a,b)=bI_{q+1}+a\,e_0\beta^{\top},
\qquad
\beta^{\top}e_0=\binom{q}{0}=1.
\]
因此由 Sherman--Morrison 公式，
\[
\bigl(bI_{q+1}+a\,e_0\beta^{\top}\bigr)^{-1}
\;=\;
\frac1bI_{q+1}-\frac{a}{b(a+b)}\,e_0\beta^{\top},
\qquad
\det\!\bigl(bI_{q+1}+a\,e_0\beta^{\top}\bigr)=b^{q}(a+b).
\]
结合 \(\det M_q(a,b)=\det N_q\cdot \det(N_q^{-1}M_q(a,b))\) 即得行列式闭式。
\par
为得到条目级逆矩阵，注意 \(P_q^{-1}\) 的经典条目为
\(
(P_q^{-1})_{k\ell}=(-1)^{k-\ell}\binom{k}{\ell}\mathbf 1_{\{k\ge \ell\}}
\)，
故
\[
(N_q^{-1})_{k\ell}=(P_q^{-1}R_q)_{k\ell}
=(-1)^{k+\ell-q}\binom{k}{k+\ell-q}\mathbf 1_{\{k+\ell\ge q\}}.
\]
又有 \(\beta^{\top}N_q^{-1}=e_0^{\top}\)，从而
\[
M_q(a,b)^{-1}
=\bigl(bI_{q+1}+a\,e_0\beta^{\top}\bigr)^{-1}N_q^{-1}
=\frac1bN_q^{-1}-\frac{a}{b(a+b)}\,e_0e_0^{\top},
\]
即得到所示条目公式。证毕。
\end{proof}

\begin{conclusion}[全阶子式最大公因子与 Smith 标准形的刚性]\label{con:xi-terminal-replica-softcore-mq-ab-smith-snf}
固定 \(q\ge 2\) 并取 \(a,b\in\ZZ\)。
记 \(g:=\gcd(a,b)\)，并令 \(d_k\) 表示矩阵 \(M_q(a,b)\) 的所有 \(k\times k\) 子式的最大公因子（约定 \(d_0=1\)）。
则
\[
d_k=g\,|b|^{k-1}\quad(1\le k\le q),
\qquad
d_{q+1}=|b|^{q}\,|a+b|.
\]
从而 \(M_q(a,b)\) 的 Smith 标准形（取对角元为正并满足整除链）为
\[
\mathrm{SNF}\big(M_q(a,b)\big)
=\mathrm{diag}\Bigl(g,\underbrace{|b|,\dots,|b|}_{q-1\ \mathrm{times}},\frac{|b(a+b)|}{g}\Bigr).
\]
\end{conclusion}
\begin{proof}
先写 \(M_q(a,b)=g\,M_q(a/g,b/g)\)，故只需证明 \(\gcd(a,b)=1\) 的情形；一般情形由整体乘子 \(g\) 同步放大 Smith 不变量得到。
\par
设 \(\gcd(a,b)=1\)，并沿用上一个证明中的幺模矩阵 \(N_q\) 与
\(
A_q:=N_q^{-1}M_q(a,b)=bI_{q+1}+a\,e_0\beta^{\top}
\)。
由于 \(N_q\) 幺模，\(M_q(a,b)\) 与 \(A_q\) 幺模等价，因此它们具有相同的 \(d_k\) 与 Smith 不变量。
注意 \(A_q\) 具有显式上三角形：
\[
A_q=
\begin{pmatrix}
a+b & a\binom{q}{1} & \cdots & a\binom{q}{q-1} & a\\
0& b\\
& & \ddots\\
0&&& b\\
0&&&& b
\end{pmatrix}.
\]
对 \(1\le k\le q\)，任取 \(k\times k\) 子矩阵：若其不含第 \(0\) 行，则它为对角子块，子式为 \(\pm b^{k}\)；
若含第 \(0\) 行，则为使子式非零，必须同时包含至少 \(k-1\) 个不同的对角行，因此每个非零子式必含因子 \(b^{k-1}\)。
另一方面，取行集 \(\{0,1,\dots,k-1\}\) 与列集 \(\{q,1,\dots,k-1\}\)，对应子式恰为 \(\pm a\,b^{k-1}\)；
由 \(\gcd(a,b)=1\) 得该子式在 \(b^{k-1}\) 层面上不可再共同整除，故 \(d_k=|b|^{k-1}\)。
最后 \(|\det A_q|=|b|^{q}|a+b|\) 给出 \(d_{q+1}\)。
\par
由 Smith 标准形理论 \(d_k=\prod_{i=1}^{k}s_i\)（\(s_i\) 为 Smith 对角不变量），即可解出所示 \(\mathrm{SNF}\)。
证毕。
\end{proof}

\begin{conclusion}[原始情形的余核分解：一个 \texorpdfstring{$|b(a+b)|$}{|b(a+b)|} 循环因子与 \texorpdfstring{$(q-1)$}{(q-1)} 个 \texorpdfstring{$|b|$}{|b|} 挠因子]\label{con:xi-terminal-replica-softcore-mq-ab-cokernel-splitting}
在 \(\gcd(a,b)=1\) 且 \(b\neq 0,a+b\neq 0\) 下，令
\(
A_q=bI_{q+1}+a\,e_0\beta^{\top}
\)
如上，并令 \(\bar f_0,\dots,\bar f_q\) 表示标准基在余核
\(
\mathrm{coker}(A_q)=\ZZ^{q+1}/A_q\ZZ^{q+1}
\)
中的像。
则：
\begin{enumerate}
  \item \(\bar f_q\) 的阶恰为 \(|b(a+b)|\)；
  \item 对每个 \(1\le k\le q-1\)，元素 \(\bar f_k-\binom{q}{k}\bar f_q\) 的阶恰为 \(|b|\)，且它们彼此独立；
  \item 存在同构
  \[
  \mathrm{coker}(A_q)\cong (\ZZ/|b|\ZZ)^{q-1}\ \oplus\ \ZZ/|b(a+b)|\ZZ.
  \]
\end{enumerate}
\end{conclusion}
\begin{proof}
由列关系
\[
A_q e_0=(a+b)f_0,\qquad
A_q e_q=af_0+bf_q,\qquad
A_q e_k=a\binom{q}{k}f_0+bf_k\ (1\le k\le q-1),
\]
在余核中得到
\[
(a+b)\bar f_0=0,\quad
b\bar f_q=-a\bar f_0,\quad
b\bar f_k=-a\binom{q}{k}\bar f_0.
\]
由此对 \(1\le k\le q-1\) 有
\(
b(\bar f_k-\binom{q}{k}\bar f_q)=0
\)，
故这些元素均为 \(|b|\)-挠元。
考虑投影同态
\(
\pi:\ZZ^{q+1}\to (\ZZ/b\ZZ)^{q-1},
\ \pi(y):=(y_1,\dots,y_{q-1})\bmod b
\)：
由于 \((A_qx)_k=bx_k\)（\(1\le k\le q-1\)），\(\pi\) 在像子群上恒为零，从而诱导良定同态
\(
\bar\pi:\mathrm{coker}(A_q)\to(\ZZ/|b|\ZZ)^{q-1}
\)。
显然 \(\bar\pi(\bar f_q)=0\) 且 \(\bar\pi(\bar f_k-\binom{q}{k}\bar f_q)=e_k\)，因此这些 \(|b|\)-挠元阶为 \(|b|\) 且相互独立。
\par
另一方面，由 \((a+b)b\bar f_q=-a(a+b)\bar f_0=0\) 得 \(|b(a+b)|\,\bar f_q=0\)，故 \(\bar f_q\) 的阶整除 \(|b(a+b)|\)。
取 \(s\in\ZZ\) 满足 \(as\equiv -1\pmod{a+b}\)（可取因为 \(\gcd(a,a+b)=1\)），并定义同态
\[
\psi:\ZZ^{q+1}\to \ZZ/(b(a+b))\ZZ,\qquad
\psi(f_0):=bs,\ \ \psi(f_q):=1,\ \ \psi(f_k):=\binom{q}{k}\ (1\le k\le q-1).
\]
直接计算可得 \(\psi(A_qe_0)=\psi(A_qe_k)=0\)（\(1\le k\le q\)），故 \(\psi\) 诱导同态
\(
\bar\psi:\mathrm{coker}(A_q)\twoheadrightarrow \ZZ/(b(a+b))\ZZ
\)
且 \(\bar\psi(\bar f_q)=1\)。
因此 \(\bar f_q\) 的阶至少为 \(|b(a+b)|\)，结合上界即知其阶恰为 \(|b(a+b)|\)。
\par
由此可知由 \(\bar f_q\) 与 \(\{\bar f_k-\binom{q}{k}\bar f_q\}_{1\le k\le q-1}\) 生成的子群大小为
\(
|b(a+b)|\cdot |b|^{q-1}=|b|^{q}|a+b|=|\det A_q|
\)，
而 \(|\mathrm{coker}(A_q)|=|\det A_q|\)，故该子群即为全体余核并给出所示直和分解。证毕。
\end{proof}

\begin{conclusion}[整数可解性判别的反序差分闭式]\label{con:xi-terminal-replica-softcore-mq-ab-integer-solvability}
令 \(q\ge 2\)，并设 \(\gcd(a,b)=1\)、\(b\neq 0\)。
对任意 \(y\in\ZZ^{q+1}\)，线性方程组
\[
M_q(a,b)\,x=y,\qquad x\in\ZZ^{q+1},
\]
可解当且仅当下列同余同时成立。
令
\[
y':=P_q^{-1}R_q\,y,
\qquad
y'_k=\sum_{j=0}^{k}(-1)^{k-j}\binom{k}{j}\,y_{q-j}.
\]
则可解性等价于
\[
b\mid y'_k\ (1\le k\le q),\qquad
a+b\mid\Bigl(y'_0-\sum_{k=1}^{q}\binom{q}{k}\frac{y'_k}{b}\Bigr).
\]
在此条件下，整数解可取为
\[
x_k=\frac{y'_k}{b}\ (1\le k\le q),\qquad
x_0=\frac{y'_0-\sum_{k=1}^{q}\binom{q}{k}\frac{y'_k}{b}}{a+b}.
\]
\end{conclusion}
\begin{proof}
由结论 \ref{con:xi-terminal-replica-softcore-mq-ab-inverse-det} 的幺模等价 \(A_q=N_q^{-1}M_q(a,b)\)（其中 \(N_q=R_qP_q\)），
\(
M_q(a,b)x=y
\)
等价于
\(
A_qx=y'
\)，其中 \(y'=N_q^{-1}y=P_q^{-1}R_qy\)。
对 \(k\ge 1\) 行有 \(bx_k=y'_k\)，故必须且只需 \(b\mid y'_k\) 并令 \(x_k=y'_k/b\)。
代回第 \(0\) 行即得
\[
(a+b)x_0+\sum_{k=1}^{q}a\binom{q}{k}\frac{y'_k}{b}=y'_0,
\]
从而得到第二个整除条件与所示 \(x_0\) 的闭式表达。证毕。
\end{proof}

\endinput


\paragraph{反向二项核的加权正交闭式：对角自伴随权的唯一性、范数恒等式与谱投影条目化}

固定整数 \(q\ge 2\)。
记反向二项核 \(B_q\in \Mat_{q+1}(\ZZ)\) 为
\[
(B_q)_{k\ell}:=\binom{q-\ell}{k}\qquad(0\le k,\ell\le q),
\]
并令二项权矩阵
\[
W_q:=\mathrm{diag}\Bigl(\binom{q}{0}^{-1},\binom{q}{1}^{-1},\dots,\binom{q}{q}^{-1}\Bigr),
\qquad
\langle u,v\rangle_{W_q}:=u^{\mathsf T}W_qv.
\]

\begin{conclusion}[二项权自伴随化的唯一性：对角权仅差整体比例]\label{con:xi-bq-unique-diagonal-selfadjoint}
设 \(D=\mathrm{diag}(d_0,\dots,d_q)\succ 0\) 为对角矩阵并满足
\[
B_q^{\mathsf T}D=DB_q.
\]
则 \(D\)（在正比例意义下）唯一确定，并且必为
\[
\boxed{\;
D\ \propto\ W_q
\;=\;
\mathrm{diag}\Bigl(\binom{q}{0}^{-1},\binom{q}{1}^{-1},\dots,\binom{q}{q}^{-1}\Bigr).\;}
\]
\end{conclusion}
\begin{proof}
将 \(B_q^{\mathsf T}D=DB_q\) 取 \((k,0)\) 元，得到
\[
(B_q^{\mathsf T}D)_{k0}=(B_q)_{0k}\,d_0=\binom{q-k}{0}\,d_0=d_0,
\qquad
(DB_q)_{k0}=d_k(B_q)_{k0}=d_k\binom{q}{k}.
\]
故对一切 \(0\le k\le q\) 有 \(d_k=d_0/\binom{q}{k}\)，从而 \(D\) 由 \(d_0\) 唯一确定并与 \(W_q\) 成正比例。证毕。
\end{proof}

\begin{conclusion}[黄金谱的正交范数闭式：线性型本征向量的 \texorpdfstring{$W_q$}{Wq}-规范化]\label{con:xi-bq-golden-spectrum-orthogonality-norm}
令 \(\varphi=\frac{1+\sqrt5}{2}\)。
对每个 \(j\in\{0,1,\dots,q\}\)，定义齐次多项式
\[
P_j(x,y):=(\varphi x+y)^{q-j}(x-\varphi y)^{j}
\,=\,\sum_{k=0}^{q}c^{(j)}_k\,x^{q-k}y^{k},
\qquad
c^{(j)}:=(c^{(j)}_0,\dots,c^{(j)}_q)^{\mathsf T}\in\RR^{q+1}.
\]
则 \(c^{(j)}\) 为 \(B_q\) 的本征向量且本征值为
\[
\boxed{\;\mu_j=(-1)^j\varphi^{\,q-2j}\;}
\qquad(0\le j\le q).
\]
并且在 \(\langle\cdot,\cdot\rangle_{W_q}\) 下满足严格正交与范数闭式
\[
\langle c^{(j)},c^{(j')}\rangle_{W_q}=0\quad(j\neq j'),
\qquad
\boxed{\;\langle c^{(j)},c^{(j)}\rangle_{W_q}=\frac{(\varphi+2)^{q}}{\binom{q}{j}}.\;}
\]
因此归一化向量
\[
u^{(j)}:=\sqrt{\frac{\binom{q}{j}}{(\varphi+2)^q}}\;c^{(j)}
\]
构成一组 \(W_q\)-正交规范基：\(u^{(j)\mathsf T}W_qu^{(j')}=\delta_{jj'}\)。
\end{conclusion}
\begin{proof}[证明要点]
\par\noindent\textbf{(i) 代入算子与本征关系}
对任意系数向量 \(c=(c_0,\dots,c_q)^{\mathsf T}\)，令 \(F_c(x,y):=\sum_{k=0}^{q}c_kx^{q-k}y^k\)。
由二项式展开可验证 \(B_q\) 对应代入作用 \((x,y)\mapsto(x+y,x)\)：
\[
F_{B_qc}(x,y)=F_c(x+y,x).
\]
对线性型 \(L_+(x,y):=\varphi x+y\)、\(L_-(x,y):=x-\varphi y\) 有
\[
L_+(x+y,x)=\varphi(x+y)+x=\varphi^2x+\varphi y=\varphi\,L_+(x,y),
\]
\[
L_-(x+y,x)=(x+y)-\varphi x=y-\varphi^{-1}x=-\varphi^{-1}L_-(x,y),
\]
其中使用 \(\varphi^2=\varphi+1\)。
故
\[
P_j(x+y,x)=\varphi^{q-j}(-\varphi^{-1})^{j}P_j(x,y)=(-1)^j\varphi^{q-2j}P_j(x,y),
\]
从而 \(B_qc^{(j)}=\mu_jc^{(j)}\)。

\par\noindent\textbf{(ii) 正交性}
由命题 \ref{prop:pom-bq-weighted-selfadjoint-binomial}，\(B_q\) 在 \(\langle\cdot,\cdot\rangle_{W_q}\) 下自伴随；
又由命题 \ref{prop:pom-bq-is-symq-and-spectrum}，本征值 \(\{\mu_j\}_{j=0}^{q}\) 两两不同。
因此不同本征值对应的本征向量两两正交，得到第一式。

\par\noindent\textbf{(iii) 范数闭式}
定义齐次次数 \(q\) 的 Fischer--apolar 配对
\[
\langle F,G\rangle_{\mathrm{ap}}
:=\frac{1}{q!}\,F(\partial_x,\partial_y)\,G(x,y)\big|_{(0,0)}.
\]
对基 \(x^{q-k}y^k\) 的直接计算给出
\(
\langle F_c,F_{c'}\rangle_{\mathrm{ap}}=c^{\mathsf T}W_qc'
\)，
故只需计算 \(\langle P_j,P_j\rangle_{\mathrm{ap}}\)。
令方向导数算子
\[
D_+:=\varphi\,\partial_x+\partial_y,\qquad D_-:=\partial_x-\varphi\,\partial_y.
\]
则 \(D_+(\varphi x+y)=\varphi^2+1=\varphi+2\)、\(D_+(x-\varphi y)=0\)，以及 \(D_-(\varphi x+y)=0\)、\(D_-(x-\varphi y)=\varphi+2\)。
因此 \(D_+\) 仅作用于第一因子而 \(D_-\) 仅作用于第二因子，得到
\[
D_+^{\,q-j}D_-^{\,j}P_j
=(q-j)!\,j!\,(\varphi+2)^{q}.
\]
再由定义
\(
\langle P_j,P_j\rangle_{\mathrm{ap}}
=\frac{1}{q!}D_+^{\,q-j}D_-^{\,j}P_j
\)
得到
\(
\langle c^{(j)},c^{(j)}\rangle_{W_q}
=\frac{(q-j)!\,j!}{q!}(\varphi+2)^q=(\varphi+2)^q/\binom{q}{j}
\)。
证毕。
\end{proof}

\begin{conclusion}[谱投影与逆本征矩阵闭式：无需求逆即可条目级重构 \texorpdfstring{$B_q^{\,m}$}{Bq^m}]\label{con:xi-bq-spectral-projector-entrywise}
令
\[
C_q:=\bigl[c^{(0)},c^{(1)},\dots,c^{(q)}\bigr]\in\RR^{(q+1)\times(q+1)},
\qquad
\Lambda_q:=\mathrm{diag}(\mu_0,\dots,\mu_q),
\]
并记 \(W_q\) 如上。
则 \(C_q\) 可逆且其逆矩阵满足闭式
\[
\boxed{\;
C_q^{-1}=(\varphi+2)^{-q}\,W_q^{-1}\,C_q^{\mathsf T}\,W_q\;}
\]
并由此对任意整数 \(m\ge 0\) 有谱投影展开
\[
\boxed{\;
B_q^{\,m}
=(\varphi+2)^{-q}\,C_q\,\Lambda_q^{\,m}\,W_q^{-1}\,C_q^{\mathsf T}\,W_q\; } .
\]
特别地，每个条目 \((k,\ell)\) 具有完全显式有限和表示
\[
\boxed{\;
\bigl(B_q^{\,m}\bigr)_{k\ell}
=\frac{1}{(\varphi+2)^{q}\binom{q}{\ell}}
\sum_{j=0}^{q}\binom{q}{j}\,\mu_j^{m}\,c^{(j)}_{k}\,c^{(j)}_{\ell}\; } .
\]
\end{conclusion}
\begin{proof}
由结论 \ref{con:xi-bq-golden-spectrum-orthogonality-norm} 的正交范数闭式，矩阵式可写为
\[
C_q^{\mathsf T}W_qC_q=(\varphi+2)^q\,W_q.
\]
左乘 \(W_q^{-1}\) 得 \(W_q^{-1}C_q^{\mathsf T}W_qC_q=(\varphi+2)^q I\)，从而
\(
C_q^{-1}=(\varphi+2)^{-q}W_q^{-1}C_q^{\mathsf T}W_q
\)。
由对角化 \(B_q=C_q\Lambda_qC_q^{-1}\) 得到所示 \(B_q^{\,m}\) 的谱投影展开；条目公式由右端逐项展开并注意 \(W_q\) 的对角性读出。证毕。
\end{proof}

\begin{conclusion}[完备性核恒等式：一条“黄金--二项”的离散分辨公式]\label{con:xi-bq-resolution-identity}
沿结论 \ref{con:xi-bq-spectral-projector-entrywise} 的记号，取 \(m=0\) 得对任意 \(0\le k,\ell\le q\) 成立恒等式
\[
\boxed{\;
\sum_{j=0}^{q}\binom{q}{j}\,c^{(j)}_{k}\,c^{(j)}_{\ell}
=(\varphi+2)^{q}\binom{q}{\ell}\,\delta_{k\ell}\; } .
\]
\end{conclusion}
\begin{proof}
将结论 \ref{con:xi-bq-spectral-projector-entrywise} 的条目公式取 \(m=0\)，并用 \((B_q^0)_{k\ell}=\delta_{k\ell}\) 即得。证毕。
\end{proof}

\begin{proposition}[反向二项核幂的 Fibonacci 条目闭式]\label{prop:xi-bq-power-entrywise-fibonacci-binomial}
取 Fibonacci 数列 \(F_0=0,F_1=1,F_{n+1}=F_n+F_{n-1}\)，并记
\[
K:=\begin{pmatrix}1&1\\[2pt]1&0\end{pmatrix},
\qquad
K^m=\begin{pmatrix}F_{m+1}&F_m\\[2pt]F_m&F_{m-1}\end{pmatrix}\qquad(m\ge 1).
\]
则对任意整数 \(q\ge 0\)、\(m\ge 1\) 及 \(0\le k,\ell\le q\)，成立条目闭式
\begin{equation}\label{eq:xi-bq-power-entrywise-fibonacci-binomial}
\bigl(B_q\bigr)^{m}_{k,\ell}
=
\sum_{t=\max(0,k+\ell-q)}^{\min(k,\ell)}
\binom{\ell}{t}\binom{q-\ell}{k-t}\,
F_{m+1}^{\,q-k-\ell+t}\,
F_{m}^{\,k+\ell-2t}\,
F_{m-1}^{\,t}.
\end{equation}
等价地，令
\[
C^{(q)}:=B_q^{\mathsf T},\qquad
C^{(q)}_{i,j}=\binom{q-i}{j}\qquad(0\le i,j\le q),
\]
则对任意 \(0\le i,j\le q\) 有
\begin{equation}\label{eq:xi-Cq-power-entrywise-fibonacci-binomial}
\bigl(C^{(q)}\bigr)^{m}_{i,j}
=
\sum_{t=\max(0,i+j-q)}^{\min(i,j)}
\binom{i}{t}\binom{q-i}{j-t}\,
F_{m+1}^{\,q-i-j+t}\,
F_{m}^{\,i+j-2t}\,
F_{m-1}^{\,t}.
\end{equation}
\end{proposition}
\begin{proof}
由 \((x,y)\mapsto(x+y,x)\) 的迭代对应 \(K^m\) 的代入作用，
\(
F_{B_q^{\,m}c}(x,y)=F_c\bigl((K^m)_{11}x+(K^m)_{12}y,\ (K^m)_{21}x+(K^m)_{22}y\bigr)
\)。
因此 \(\bigl(B_q\bigr)^m_{k,\ell}\) 等于齐次多项式
\(
(F_{m+1}x+F_my)^{q-\ell}(F_mx+F_{m-1}y)^{\ell}
\)
中 \(x^{q-k}y^k\) 的系数。
对两因子分别作二项式展开并令第二因子贡献的 \(y\)-次数为 \(t\)，则第一因子贡献 \(y\)-次数为 \(k-t\)；
求和范围由次数约束截断为
\(
t\in[\max(0,k+\ell-q),\min(k,\ell)]
\)。
收集系数与幂次即得 \eqref{eq:xi-bq-power-entrywise-fibonacci-binomial}。
将 \(\bigl(C^{(q)}\bigr)^m=(B_q^m)^{\mathsf T}\) 代入并交换指标得到 \eqref{eq:xi-Cq-power-entrywise-fibonacci-binomial}。
\end{proof}
\leanverified{paper\_xi\_bq\_power\_entrywise\_fibonacci\_binomial}

\begin{corollary}[边界条目与行和闭式]\label{cor:xi-Cq-power-boundary-and-row-sum}
在 \eqref{eq:xi-Cq-power-entrywise-fibonacci-binomial} 的记号下，对任意 \(q\ge 0,m\ge 1\) 与 \(0\le i,j\le q\)，成立：
\[
\bigl(C^{(q)}\bigr)^m_{0,j}=\binom{q}{j}\,F_{m+1}^{\,q-j}F_m^{\,j},
\qquad
\bigl(C^{(q)}\bigr)^m_{i,0}=F_{m+1}^{\,q-i}F_m^{\,i},
\qquad
\bigl(C^{(q)}\bigr)^m_{i,q}=F_m^{\,q-i}F_{m-1}^{\,i},
\]
以及
\[
\sum_{j=0}^{q}\bigl(C^{(q)}\bigr)^m_{i,j}=F_{m+2}^{\,q-i}F_{m+1}^{\,i}.
\]
\end{corollary}
\begin{proof}
前三式由 \eqref{eq:xi-Cq-power-entrywise-fibonacci-binomial} 在 \(i=0\) 或 \(j\in\{0,q\}\) 时求和退化为单项直接得到。
对行和恒等式，取
\(
P_i(x,y):=(F_{m+1}x+F_my)^{q-i}(F_mx+F_{m-1}y)^i
\)。
则 \(\bigl(C^{(q)}\bigr)^m_{i,j}\) 为 \(P_i\) 中 \(x^{q-j}y^j\) 的系数，故
\[
\sum_{j=0}^{q}\bigl(C^{(q)}\bigr)^m_{i,j}=P_i(1,1)
=(F_{m+1}+F_m)^{q-i}(F_m+F_{m-1})^i
=F_{m+2}^{\,q-i}F_{m+1}^{\,i}.
\]
\end{proof}
\leanverified{paper\_xi\_cq\_power\_boundary\_and\_row\_sum}

\begin{theorem}[严格符号正则性：幂矩阵各阶子式的统一符号律]\label{thm:xi-Cq-power-strict-sign-regular}
\leanverified{paper\_xi\_cq\_power\_strict\_sign\_regular}
对任意整数 \(q\ge 0\) 与 \(m\ge 1\)，矩阵 \(\bigl(C^{(q)}\bigr)^m\) 为严格符号正则矩阵。
更精确地，对每个阶数 \(1\le k\le q+1\)，任意非零 \(k\times k\) 子式的符号恒等于
\[
\varepsilon_k(m):=(-1)^{m\binom{k}{2}}.
\]
特别地，当 \(m\) 为偶数时，\(\bigl(C^{(q)}\bigr)^m\) 为全阶完全非负；当 \(m\) 为奇数时，其非零 \(k\)-阶子式符号恒为 \((-1)^{\binom{k}{2}}\)。
\end{theorem}
\begin{proof}
\par\noindent\textbf{(i) 基准符号律.}
令 \(P^{(q)}\) 为 Pascal 矩阵 \(P^{(q)}_{i,j}=\binom{i}{j}\mathbf 1_{\{i\ge j\}}\)，并令 \(R^{(q)}\) 为反序置换矩阵 \(R^{(q)}_{i,j}=\mathbf 1_{\{i+j=q\}}\)。
则 \(C^{(q)}=R^{(q)}P^{(q)}\)。
由于 \(P^{(q)}\) 的任意非零子式为正整数且 \(R^{(q)}\) 的任意非零 \(k\)-阶子式符号为 \((-1)^{\binom{k}{2}}\)，由 Cauchy--Binet 公式可得：\(C^{(q)}\) 的任意非零 \(k\)-阶子式符号恒为 \((-1)^{\binom{k}{2}}\)。
\par
\noindent\textbf{(ii) 偶幂的全阶完全非负性.}
由命题 \ref{prop:xi-bq-power-entrywise-fibonacci-binomial}，有 \(\bigl(C^{(q)}\bigr)^2=\mathrm{Sym}^q(K^2)^{\mathsf T}\)。
其中 \(K^2=\bigl(\begin{smallmatrix}2&1\\ 1&1\end{smallmatrix}\bigr)\) 各项为正且 \(\det(K^2)=1\)。
取 Gauss 分解 \(K^2=LDU\)（\(L,U\) 为单位下/上三角且含正参数，\(D\) 为正对角），由表示同态性
\(
\mathrm{Sym}^q(K^2)=\mathrm{Sym}^q(L)\mathrm{Sym}^q(D)\mathrm{Sym}^q(U)
\)。
三者均为全阶完全非负矩阵，且乘积保持全阶完全非负；从而 \(\bigl(C^{(q)}\bigr)^2\) 全阶完全非负，并由乘法封闭性得到任意偶数 \(m=2r\) 时 \(\bigl(C^{(q)}\bigr)^m\) 全阶完全非负。
\par
\noindent\textbf{(iii) 奇幂的符号传递.}
令 \(m=2r+1\)。对任意固定的行集 \(I\)、列集 \(J\)（\(|I|=|J|=k\)），由 Cauchy--Binet 公式
\[
\det\!\Bigl(\bigl(C^{(q)}\bigr)^{m}[I,J]\Bigr)
=\sum_{S:\ |S|=k}\det\!\Bigl(\bigl(C^{(q)}\bigr)^{2r}[I,S]\Bigr)\,\det\!\bigl(C^{(q)}[S,J]\bigr).
\]
由第 (ii) 步，\(\det((C^{(q)})^{2r}[I,S])\ge 0\)；
由第 (i) 步，\(\det(C^{(q)}[S,J])\) 或为 \(0\)，或符号为 \((-1)^{\binom{k}{2}}\)。
因此所有非零求和项同号，不发生抵消，故左端或为零，或符号恒为 \((-1)^{\binom{k}{2}}\)。
合并偶幂情形即得 \(\varepsilon_k(m)=(-1)^{m\binom{k}{2}}\)。
\end{proof}

\endinput


\paragraph{黄金对合矩阵的对称幂：缩放对合、二项互易与两点谱塌缩}
固定整数 \(q\ge 2\) 并令黄金比例 \(\varphi=\frac{1+\sqrt5}{2}\)。
定义
\[
G:=\begin{pmatrix}\varphi&1\\[2pt]1&-\varphi\end{pmatrix}\in\Mat_2(\ZZ[\varphi]),
\qquad
G^2=(\varphi^2+1)I_2=(\varphi+2)I_2.
\]
令 \(\RR[x,y]_q\) 为齐次次数 \(q\) 的二元多项式空间，并令线性代换算子
\[
(\mathcal{G}_q f)(x,y):=f(\varphi x+y,\ x-\varphi y)
\]
作用其上。取单项式基 \(\{x^{q-j}y^j\}_{j=0}^{q}\)，记 \(\mathcal{G}_q\) 的矩阵为 \(C_q=(c_{k j})_{0\le k,j\le q}\)，其中
\[
c_{k j}:=[x^{q-k}y^k]\bigl(\varphi x+y\bigr)^{q-j}\bigl(x-\varphi y\bigr)^j.
\]
由构造 \(C_q=\mathrm{Sym}^q(G)\)。

\begin{conclusion}[缩放对合与二项互易：\texorpdfstring{$C_q^2=(\varphi+2)^qI$}{Cq^2=(phi+2)^q I}]\label{con:xi-symq-golden-involution-reciprocity}
令二项权矩阵
\[
W_q:=\mathrm{diag}\Bigl(\binom{q}{0}^{-1},\binom{q}{1}^{-1},\dots,\binom{q}{q}^{-1}\Bigr).
\]
则成立矩阵恒等式
\[
\boxed{\;
C_q^2=(\varphi+2)^{q}\,I_{q+1},
\qquad
C_q^{\mathsf T}W_q=W_qC_q.\;}
\]
等价地，对一切 \(0\le k,j\le q\) 有精确互易律
\[
\boxed{\;
\binom{q}{j}\,c_{k j}=\binom{q}{k}\,c_{j k}\;}
\]
并且逆矩阵满足闭式
\[
\boxed{\;
C_q^{-1}=(\varphi+2)^{-q}\,C_q.\;}
\]
\end{conclusion}
\begin{proof}
由于 \(C_q=\mathrm{Sym}^q(G)\) 且 \(G^2=(\varphi+2)I_2\)，由对称幂函子性得到
\[
C_q^2=\mathrm{Sym}^q(G)^2=\mathrm{Sym}^q(G^2)=\mathrm{Sym}^q\!\bigl((\varphi+2)I_2\bigr)=(\varphi+2)^q I_{q+1},
\]
从而 \(C_q^{-1}=(\varphi+2)^{-q}C_q\)。
\par
对加权自伴随性，沿用 Fischer--apolar 配对
\(
\langle F,G\rangle_{\mathrm{ap}}=\frac{1}{q!}F(\partial_x,\partial_y)G(x,y)\big|_{(0,0)}
\)。
对基 \(\{x^{q-k}y^k\}_{k=0}^{q}\) 的直接计算给出其 Gram 矩阵恰为 \(W_q\)。
线性代换 \(f\mapsto f\circ G\) 在该配对下的伴随等于 \(f\mapsto f\circ G^{\mathsf T}\)；由于 \(G\) 为实对称矩阵，得到 \(\mathcal{G}_q\) 的自伴随性，故 \(C_q^{\mathsf T}W_q=W_qC_q\)。
逐项展开即得所示互易律。证毕。
\end{proof}

\begin{conclusion}[谱的两点塌缩：\texorpdfstring{$\mathrm{Spec}(C_q)\subset\{\pm(\varphi+2)^{q/2}\}$}{Spec(Cq) in {±(phi+2)^{q/2}}}]\label{con:xi-symq-golden-spectrum-two-point}
\(G\) 的特征值为 \(\pm\sqrt{\varphi+2}\)。
因此 \(C_q=\mathrm{Sym}^q(G)\) 的全部特征值仅为两点
\[
\boxed{\;
\lambda_{+}=(\varphi+2)^{q/2},
\qquad
\lambda_{-}=-(\varphi+2)^{q/2}\;}
\]
其重数分别为
\[
m_{+}(q)=\#\{0\le j\le q:\ j\equiv 0\ (\mathrm{mod}\ 2)\}=\Bigl\lfloor\frac q2\Bigr\rfloor+1,
\qquad
m_{-}(q)=\#\{0\le j\le q:\ j\equiv 1\ (\mathrm{mod}\ 2)\}=\Bigl\lceil\frac q2\Bigr\rceil.
\]
\end{conclusion}
\begin{proof}
由于 \(G^2=(\varphi+2)I_2\) 且 \(\Tr(G)=0\)，\(G\) 的特征值为 \(\pm\sqrt{\varphi+2}\)。
设其为 \((\alpha,\beta)=(\sqrt{\varphi+2},-\sqrt{\varphi+2})\)。
对称幂表示的谱函子性给出：\(\mathrm{Sym}^q(G)\) 的特征值为 \(\alpha^{q-j}\beta^j\)（\(0\le j\le q\)），故
\(
\alpha^{q-j}\beta^j=(-1)^j(\varphi+2)^{q/2}
\)
仅取两值。
奇偶计数给出重数。证毕。
\end{proof}

\begin{conclusion}[迹与行列式闭式：奇偶律与代数范数坍缩]\label{con:xi-symq-golden-trace-det-norm}
成立迹的奇偶律
\[
\boxed{\;
\tr(C_q)=
\begin{cases}
(\varphi+2)^{q/2},& q\ \text{为偶},\\[2pt]
0,& q\ \text{为奇},
\end{cases}}
\]
以及行列式闭式
\[
\boxed{\;
\det(C_q)=(-1)^{\frac{q(q+1)}{2}}\,(\varphi+2)^{\frac{q(q+1)}{2}}.\;}
\]
并且在二次数域 \(\QQ(\sqrt5)\) 上的代数范数满足
\[
\boxed{\;
\mathrm{N}_{\QQ(\sqrt5)/\QQ}\!\bigl(\det(C_q)\bigr)=5^{\frac{q(q+1)}{2}}.\;}
\]
\end{conclusion}
\begin{proof}
由结论 \ref{con:xi-symq-golden-spectrum-two-point} 的两点谱与重数，
\[
\tr(C_q)=(m_+(q)-m_-(q))(\varphi+2)^{q/2},
\qquad
\det(C_q)=(\varphi+2)^{\frac{q(q+1)}{2}}(-1)^{m_-(q)}.
\]
当 \(q\) 为偶时 \(m_+(q)-m_-(q)=1\)，当 \(q\) 为奇时差为 \(0\)，得到迹公式。
又由 \(m_-(q)\equiv \frac{q(q+1)}{2}\ (\mathrm{mod}\ 2)\) 得行列式符号闭式。
\par
记 \(\widehat\varphi=\frac{1-\sqrt5}{2}\) 为 \(\varphi\) 的共轭，则
\(
\mathrm{N}_{\QQ(\sqrt5)/\QQ}(\varphi+2)=(\varphi+2)(\widehat\varphi+2)=5
\)。
由于 \(\mathrm{N}((-1)^{\frac{q(q+1)}{2}})=1\)，范数结论随之成立。证毕。
\end{proof}

\begin{conclusion}[函数演算的二维闭包：\texorpdfstring{$\QQ(\sqrt5)[C_q]$}{Q(sqrt5)[Cq]} 由两点插值控制]\label{con:xi-symq-golden-functional-calculus}
设 \(p\in\QQ(\sqrt5)[X]\) 为任意多项式，则存在唯一 \((a_p,b_p)\in\QQ(\sqrt5)^2\) 使
\[
\boxed{\;
p(C_q)=a_p\,I_{q+1}+b_p\,C_q.\;}
\]
若记 \(s=(\varphi+2)^{q/2}>0\)，则 \((a_p,b_p)\) 可由对称插值给出：
\[
a_p=\frac{p(s)+p(-s)}{2},
\qquad
b_p=\frac{p(s)-p(-s)}{2s}.
\]
\end{conclusion}
\begin{proof}
由结论 \ref{con:xi-symq-golden-involution-reciprocity}，\(C_q\) 满足 \(C_q^2=(\varphi+2)^q I_{q+1}\)，故最小多项式整除 \(X^2-(\varphi+2)^q\)，从而 \(\QQ(\sqrt5)[C_q]\) 为至多二维代数。
将 \(p\) 在 \(\{\pm s\}\) 上作偶/奇分解并代回 \(C_q\) 即得所示表达与唯一性。证毕。
\end{proof}

\paragraph{整数环上的阶梯 Smith 结构：\texorpdfstring{$\ZZ[\varphi]$}{Z[phi]} 与 \texorpdfstring{$(\varphi+2)$}{phi+2}-进骨架}
在二次数域的整数环 \(\mathcal O:=\ZZ[\varphi]\) 中，元素
\[
\pi_\varphi:=\varphi+2
\]
满足 \(\mathrm N_{\QQ(\sqrt5)/\QQ}(\pi_\varphi)=5\)（见结论 \ref{con:xi-symq-golden-trace-det-norm}），因此 \((\pi_\varphi)\) 给出 \(5\) 上方的唯一素理想。
下述结论表明：\(C_q\) 在 \(\mathcal O\) 上不仅满足缩放对合 \(C_q^2=\pi_\varphi^q I\)，其整格层面还具有严格的阶梯 Smith 结构。

\begin{conclusion}[幺模阶梯分解与 Smith 正规形]\label{con:xi-symq-golden-cq-smith-staircase}
令
\[
\ell:=\varphi x+y,\qquad m:=(\varphi-1)x+y,
\qquad
x-\varphi y=2\ell-\pi_\varphi m.
\]
对 \(0\le r\le q\) 定义
\[
Q_r(x,y):=\ell^{\,q-r}m^{\,r},
\qquad
P_j(x,y):=\ell^{\,q-j}(x-\varphi y)^{j}\qquad(0\le j\le q),
\]
并令 \(Q_q\in\Mat_{q+1}(\mathcal O)\) 为列向量 \((Q_0,\dots,Q_q)\) 相对于基 \(\{x^{q-k}y^k\}_{k=0}^{q}\) 的系数矩阵。
再令上三角矩阵 \(U_q=(u_{rj})_{0\le r,j\le q}\in\Mat_{q+1}(\mathcal O)\) 由
\[
u_{rj}:=
\begin{cases}
(-1)^r\,2^{\,j-r}\binom{j}{r},& r\le j,\\
0,& r>j.
\end{cases}
\]
则在 \(\mathcal O\) 上成立完全恒等式
\[
\boxed{\;
C_q
=
Q_q\cdot \mathrm{diag}(1,\pi_\varphi,\pi_\varphi^2,\dots,\pi_\varphi^{q})\cdot U_q\; }.
\]
并且 \(Q_q,U_q\in \mathrm{GL}_{q+1}(\mathcal O)\) 为幺模矩阵。
因此 \(C_q\) 的 Smith 正规形（按整除链排列）为
\[
\boxed{\;
\mathrm{SNF}_{\mathcal O}(C_q)=\mathrm{diag}(1,\pi_\varphi,\pi_\varphi^2,\dots,\pi_\varphi^{q})\; }.
\]
\end{conclusion}
\begin{proof}
由
\(
x-\varphi y=2\ell-\pi_\varphi m
\)
对每个 \(j\) 作二项式展开并合并 \(\ell^{q-j}\) 得
\[
P_j=\sum_{r=0}^{j}(-1)^r\,2^{\,j-r}\binom{j}{r}\,\pi_\varphi^{\,r}\,Q_r,
\]
从而以“系数列向量逐列相等”的意义得到矩阵恒等式
\(
C_q=Q_qD_qU_q
\)，其中 \(D_q=\mathrm{diag}(\pi_\varphi^{\,r})_{r=0}^{q}\)。
\par
对幺模性：线性型 \((\ell,m)\) 的系数列向量分别为 \((\varphi,1)^{\top}\) 与 \((\varphi-1,1)^{\top}\)，故其变换矩阵
\[
H:=
\begin{pmatrix}\varphi&\varphi-1\\[2pt]1&1\end{pmatrix}\in\mathrm{SL}_2(\mathcal O)
\]
满足 \((\ell,m)=(x,y)H\) 且 \(\det H=1\)。
因此 \(Q_q=\mathrm{Sym}^q(H)\) 且
\(
\det(Q_q)=\det(H)^{q(q+1)/2}=1
\)，从而 \(Q_q\) 幺模。
另一方面，\(U_q\) 为上三角矩阵且对角元为 \(u_{jj}=(-1)^j\in\mathcal O^{\times}\)，故 \(U_q\) 幺模。
最后，Smith 正规形由幺模等价唯一性与对角阶梯指数 \(0,1,\dots,q\) 立得。证毕。
\end{proof}

\paragraph{\texorpdfstring{$\pi_\varphi^q$}{pi^q}-自逆方程与 Smith 不变量的互补对称性}
在 \(\mathcal O\) 上，方程 \(AB=BA=\pi_\varphi^q I\) 将 Smith 不变因子严格限制为 \(\pi_\varphi\)-幂，并在同一组左右幺模变换下诱导互补对角化。
自平方根情形 \(A^2=\pi_\varphi^q I\) 进一步强迫指数多重集的对称性；在“互异指数”极限下由此导出阶梯刚性。

\begin{proposition}[互补 Smith 指数：\texorpdfstring{$\pi_\varphi^q$}{pi^q}-分解的同步对角化]\label{prop:xi-pi-power-smith-complementarity}
\leanverified{paper\_xi\_pi\_power\_smith\_complementarity}
令 \(n\ge 1\) 并取 \(A,B\in\Mat_n(\mathcal O)\) 满足
\[
AB=BA=\pi_\varphi^{q}I_n.
\]
则存在整数 \(0\le a_1\le\cdots\le a_n\le q\) 与幺模矩阵 \(P,Q\in\mathrm{GL}_n(\mathcal O)\) 使得
\[
A=P\,\mathrm{diag}(\pi_\varphi^{a_1},\dots,\pi_\varphi^{a_n})\,Q,
\qquad
B=Q^{-1}\,\mathrm{diag}(\pi_\varphi^{q-a_1},\dots,\pi_\varphi^{q-a_n})\,P^{-1}.
\]
特别地，\(A\) 的 Smith 不变因子序列唯一决定 \(B\) 的 Smith 不变因子序列（逐项互补于 \(q\)）。
\end{proposition}
\begin{proof}
由 \(\mathcal O\) 为 PID，取 Smith 正规形使
\(
UAV=\mathrm{diag}(d_1,\dots,d_n)
\)，其中 \(U,V\in\mathrm{GL}_n(\mathcal O)\) 幺模且 \(d_i\mid d_{i+1}\)。
令 \(B':=V^{-1}BU^{-1}\)，则由 \(AB=\pi_\varphi^qI\) 得
\[
\mathrm{diag}(d_1,\dots,d_n)\,B'=\pi_\varphi^q I_n.
\]
比较非对角元可知 \(B'\) 为对角矩阵，且对每个 \(i\) 有 \(d_i\,B'_{ii}=\pi_\varphi^q\)，从而 \(d_i\mid \pi_\varphi^q\)。
由于 \((\pi_\varphi)\) 为素理想且 \(\mathcal O\) 为 UFD，\(d_i\) 与某个 \(\pi_\varphi^{a_i}\) 相关联，且可吸收单位元使 \(d_i=\pi_\varphi^{a_i}\)（并满足 \(0\le a_1\le\cdots\le a_n\le q\)）。
此时 \(B'=\mathrm{diag}(\pi_\varphi^{q-a_1},\dots,\pi_\varphi^{q-a_n})\)，并令 \(P:=U^{-1}\)、\(Q:=V^{-1}\) 即得所示分解。证毕。
\end{proof}

\begin{corollary}[自平方根的 Smith 指数对称性]\label{cor:xi-smith-spectrum-self-square-root-symmetry}
\leanverified{paper\_xi\_smith\_spectrum\_self\_square\_root\_symmetry}
设 \(A\in\Mat_n(\mathcal O)\) 满足 \(A^2=\pi_\varphi^q I_n\)，并令其 Smith 正规形的对角不变因子为
\(
\pi_\varphi^{a_1}\mid\cdots\mid \pi_\varphi^{a_n}
\)。
则多重集 \(\{a_1,\dots,a_n\}\) 在映射 \(r\mapsto q-r\) 下封闭；等价地，按升序排列后恒有
\[
a_i+a_{n+1-i}=q\qquad(1\le i\le n).
\]
\end{corollary}
\begin{proof}
将命题 \ref{prop:xi-pi-power-smith-complementarity} 作用于 \((A,B)=(A,A)\) 即得 \(\{a_i\}=\{q-a_i\}\)（按多重度计）；升序重排得到配对恒等式。证毕。
\end{proof}

\begin{corollary}[互异指数的阶梯刚性]\label{cor:xi-smith-staircase-rigidity-distinct}
\leanverified{paper\_xi\_smith\_staircase\_rigidity\_distinct}
设 \(A\in\Mat_{q+1}(\mathcal O)\) 满足 \(A^2=\pi_\varphi^q I_{q+1}\)，并令其 Smith 指数为 \(0\le a_1\le\cdots\le a_{q+1}\le q\)。
若 \(\{a_1,\dots,a_{q+1}\}\) 两两不同，则必有
\[
\{a_1,\dots,a_{q+1}\}=\{0,1,\dots,q\},
\]
从而（按整除链排列，确定到单位元）
\[
\mathrm{SNF}_{\mathcal O}(A)=\mathrm{diag}(1,\pi_\varphi,\pi_\varphi^2,\dots,\pi_\varphi^{q}).
\]
\end{corollary}
\begin{proof}
由推论 \ref{cor:xi-smith-spectrum-self-square-root-symmetry}，集合 \(\{a_i\}\) 在 \(r\mapsto q-r\) 下封闭。
由于其含 \(q+1\) 个互异元素且每个 \(a_i\in\{0,1,\dots,q\}\)，只能等于整个集合 \(\{0,1,\dots,q\}\)。
代回 Smith 正规形即可。证毕。
\end{proof}

\begin{corollary}[子式理想的 \texorpdfstring{$\pi_\varphi$}{pi}-阶]\label{cor:xi-symq-golden-cq-determinantal-divisors}
令 \(d_k(C_q)\subset\mathcal O\) 表示由 \(C_q\) 的全部 \(k\times k\) 子式生成的理想（\(1\le k\le q+1\)）。
则在单位元意义下有
\[
\boxed{\;
d_k(C_q)\ \sim\ \pi_\varphi^{\,\frac{k(k-1)}{2}}\; }.
\]
\end{corollary}
\begin{proof}
由结论 \ref{con:xi-symq-golden-cq-smith-staircase} 的 Smith 正规形，
\(d_k(C_q)\) 等于前 \(k\) 个不变因子的乘积所生成的理想，即
\(
\prod_{r=0}^{k-1}\pi_\varphi^{\,r}=\pi_\varphi^{k(k-1)/2}
\)（确定到单位元）。证毕。
\end{proof}
\leanverified{paper\_xi\_symq\_golden\_cq\_determinantal\_divisors}

\begin{corollary}[模 \texorpdfstring{$\pi_\varphi^t$}{pi^t} 的阶梯约化与生成秩]\label{cor:xi-symq-golden-cq-mod-pi-power-staircase}
对整数 \(t\) 满足 \(1\le t\le q+1\)，在商环 \(\mathcal O/(\pi_\varphi^t)\) 上有等价变换
\[
\boxed{\;
C_q\bmod \pi_\varphi^t
\ \sim\
\mathrm{diag}(1,\pi_\varphi,\dots,\pi_\varphi^{t-1},0,\dots,0)\;}
\]
其中 “\(\sim\)” 表示左右乘以 \(\mathrm{GL}_{q+1}(\mathcal O/(\pi_\varphi^t))\) 的可逆矩阵。
因此像模 \(\mathrm{Im}(C_q\bmod \pi_\varphi^t)\) 的最小生成元数为 \(t\)；
等价地，
\[
\dim_{\mathcal O/(\pi_\varphi)}\Bigl(\mathrm{Im}(C_q\bmod \pi_\varphi^t)\big/\pi_\varphi\,\mathrm{Im}(C_q\bmod \pi_\varphi^t)\Bigr)=t.
\]
\end{corollary}
\begin{proof}
将结论 \ref{con:xi-symq-golden-cq-smith-staircase} 的幺模分解在 \(\mathcal O/(\pi_\varphi^t)\) 下约化即可；此时 \(\pi_\varphi^{r}\equiv 0\) 当且仅当 \(r\ge t\)，从而得到对角形态。
最小生成元数等于上述商空间的维数，为 Artin 局部环上的标准等价刻画。证毕。
\end{proof}

\begin{corollary}[余核的完全分解与 \texorpdfstring{$5$}{5}-规模]\label{cor:xi-symq-golden-cq-cokernel-structure}
\leanverified{paper\_xi\_symq\_golden\_cq\_cokernel\_structure}
有 \(\mathcal O\)-模同构
\[
\boxed{\;
\mathrm{coker}_{\mathcal O}(C_q)\ \cong\ \bigoplus_{r=1}^{q}\mathcal O/(\pi_\varphi^{r})\;}
\]
从而其 Jordan--H\"older 长度为 \(\sum_{r=1}^{q}r=\frac{q(q+1)}{2}\)，并且其有限基数为
\[
\boxed{\;
\#\,\mathrm{coker}_{\mathcal O}(C_q)=\prod_{r=1}^{q}\mathrm N(\pi_\varphi^{r})=5^{\frac{q(q+1)}{2}}.\;}
\]
\end{corollary}
\begin{proof}
由 Smith 正规形 \(\mathrm{diag}(1,\pi_\varphi,\dots,\pi_\varphi^{q})\)，余核同构为除去单位元因子后的直和 \(\bigoplus_{r=1}^q\mathcal O/(\pi_\varphi^r)\)。
长度与基数由加性/乘法性以及 \(\mathrm N(\pi_\varphi)=5\) 得到。证毕。
\end{proof}

\begin{corollary}[模 \texorpdfstring{$\pi_\varphi$}{pi} 的秩一外积与列比律]\label{cor:xi-symq-golden-cq-mod-pi-rankone-outer}
在剩余域 \(\mathcal O/(\pi_\varphi)\cong\FF_5\) 中有同余
\[
x-\varphi y\equiv 2(\varphi x+y)\pmod{\pi_\varphi},
\]
从而对每个 \(0\le j\le q\) 有
\[
(\varphi x+y)^{q-j}(x-\varphi y)^j\equiv 2^{j}(\varphi x+y)^q\pmod{\pi_\varphi}.
\]
令 \(u\in(\mathcal O/\pi_\varphi)^{q+1}\) 为 \((\varphi x+y)^q\) 在基 \(\{x^{q-k}y^k\}\) 下的系数列向量，令
\(
v:=(2^{0},2^{1},\dots,2^{q})^{\top}\in(\mathcal O/\pi_\varphi)^{q+1}
\)。
则
\[
\boxed{\;
C_q\bmod \pi_\varphi\;=\;u\,v^{\mathsf T}\;}
\]
因而 \(\mathrm{rank}_{\FF_5}(C_q\bmod \pi_\varphi)=1\) 且每一列满足严格列比
\[
\boxed{\;
(C_q\bmod \pi_\varphi)_{\bullet,j}=2^{j}(C_q\bmod \pi_\varphi)_{\bullet,0}\; }.
\]
并且由结论 \ref{con:xi-symq-golden-involution-reciprocity} 的 \(C_q^2=\pi_\varphi^q I_{q+1}\) 得
\[
\boxed{\;(C_q\bmod \pi_\varphi)^2=0\;},
\]
故 \(\mathrm{Im}(C_q\bmod \pi_\varphi)\subseteq \ker(C_q\bmod \pi_\varphi)\)。
\end{corollary}
\begin{proof}
由结论 \ref{con:xi-symq-golden-cq-smith-staircase} 的恒等式 \(x-\varphi y=2\ell-\pi_\varphi m\) 取模 \(\pi_\varphi\) 即得所示线性型同余；其余为逐列系数同余转译。证毕。
\end{proof}

\paragraph{预像塔、模幂截断与限制标量：阶梯 Smith 结构的伴随分解}
固定整数 \(q\ge 2\)，并沿用上文关于 \(\mathcal O=\ZZ[\varphi]\)、\(\pi_\varphi=\varphi+2\) 与 \(C_q\) 的记号。

\begin{conclusion}[预像子晶格塔的指数闭式与 \texorpdfstring{$\FF_5$}{F5}-分级维数]\label{con:xi-symq-golden-cq-preimage-tower-index}
令 \(L:=\mathcal O^{q+1}\)，并把 \(C_q\) 视为 \(L\to L\) 的 \(\mathcal O\)-线性同态。
对每个整数 \(s\ge 0\) 定义预像子晶格
\[
L(s):=\{x\in L:\ C_qx\in \pi_\varphi^{\,s}L\}.
\]
则其指数满足闭式
\[
\boxed{\;
[L:L(s)]
=5^{\sum_{r=0}^{q}\max\{0,s-r\}}
=
\begin{cases}
5^{\frac{s(s+1)}{2}},& 0\le s\le q,\\[4pt]
5^{(q+1)s-\frac{q(q+1)}{2}},& s\ge q.
\end{cases}\;}
\]
并且每一层商 \(L(s)/L(s+1)\) 都是 \(\mathcal O/(\pi_\varphi)\cong\FF_5\) 上的向量空间，且维数满足
\[
\boxed{\;
\dim_{\FF_5}\bigl(L(s)/L(s+1)\bigr)=\min\{s+1,q+1\}\; } .
\]
\end{conclusion}
\begin{proof}
由结论 \ref{con:xi-symq-golden-cq-smith-staircase} 的 Smith 正规形，存在幺模矩阵 \(P,Q\in\mathrm{GL}_{q+1}(\mathcal O)\) 使
\(
PC_qQ=\mathrm{diag}(1,\pi_\varphi,\dots,\pi_\varphi^q)
 \)。
令 \(y:=Q^{-1}x\)，则 \(x\in L(s)\) 等价于 \(\mathrm{diag}(\pi_\varphi^r)_{r=0}^q\,y\in \pi_\varphi^{\,s}L\)，亦即对每个 \(0\le r\le q\) 有
\(
y_r\in \pi_\varphi^{\max\{0,s-r\}}\mathcal O
\)。
因此
\[
[L:L(s)]
=\prod_{r=0}^{q}\bigl[\mathcal O:\pi_\varphi^{\max\{0,s-r\}}\mathcal O\bigr]
=\prod_{r=0}^{q}5^{\max\{0,s-r\}}
=5^{\sum_{r=0}^q\max\{0,s-r\}},
\]
其中使用 \(\#(\mathcal O/(\pi_\varphi^k))=\mathrm N(\pi_\varphi^k)=5^k\)。
对和式作分段求和得到所示闭式。
同理，
\(
L(s)/L(s+1)\cong\bigoplus_{r=0}^{q}\pi_\varphi^{\max\{0,s-r\}}\mathcal O\big/\pi_\varphi^{\max\{0,s+1-r\}}\mathcal O
\)，
每个直和分量要么为 \(0\)，要么同构于 \(\mathcal O/(\pi_\varphi)\cong\FF_5\)；
非零分量恰对应于 \(r\le s\)，其个数为 \(\min\{s+1,q+1\}\)。
证毕。
\end{proof}

\begin{conclusion}[模 \texorpdfstring{$\pi_\varphi^t$}{pi^t} 的外积截断：显式秩一算子和与生成系]\label{con:xi-symq-golden-cq-mod-pi-power-outer-truncation}
取整数 \(t\) 满足 \(1\le t\le q+1\)，并令 \(R_t:=\mathcal O/(\pi_\varphi^t)\)。
沿用结论 \ref{con:xi-symq-golden-cq-smith-staircase} 的幺模分解
\[
C_q=Q_q\cdot \mathrm{diag}(1,\pi_\varphi,\dots,\pi_\varphi^q)\cdot U_q.
\]
记 \(q_r\in\mathcal O^{q+1}\) 为 \(Q_q\) 的第 \(r\) 列向量、记 \(u_r^{\mathsf T}\in\mathcal O^{1\times(q+1)}\) 为 \(U_q\) 的第 \(r\) 行向量（\(0\le r\le q\)）。
并以 \(\overline q_r\) 表示 \(q_r\) 在 \(R_t^{q+1}\) 中的像。
则成立精确外积展开
\[
\boxed{\;
C_q=\sum_{r=0}^{q}\pi_\varphi^{\,r}\,q_r\,u_r^{\mathsf T}\;}
\]
并在 \(R_t\) 上满足截断同余
\[
\boxed{\;
C_q\equiv \sum_{r=0}^{t-1}\pi_\varphi^{\,r}\,q_r\,u_r^{\mathsf T}\pmod{\pi_\varphi^{t}}\;}
\]
从而 \(\mathrm{Im}(C_q\bmod \pi_\varphi^t)\) 可由 \(\{\,\pi_\varphi^{\,r}\,\overline q_r\,\}_{r=0}^{t-1}\) 生成，并且其最小生成元数为 \(t\)（见推论 \ref{cor:xi-symq-golden-cq-mod-pi-power-staircase}）。
\end{conclusion}
\begin{proof}
将 \(Q_q\mathrm{diag}(\pi_\varphi^r)U_q\) 按对角分量展开得到
\(
C_q=\sum_{r=0}^{q}\pi_\varphi^r(Q_qE_{rr}U_q)
\)，
其中 \(E_{rr}\) 为第 \((r,r)\) 处为 \(1\) 的标准矩阵单位。
而 \(Q_qE_{rr}\) 仅保留第 \(r\) 列、\(E_{rr}U_q\) 仅保留第 \(r\) 行，故 \(Q_qE_{rr}U_q=q_ru_r^{\mathsf T}\)，得到外积展开。
将等式在 \(R_t\) 下约化并注意 \(\pi_\varphi^r\equiv 0\pmod{\pi_\varphi^t}\) 当且仅当 \(r\ge t\)，即得截断同余。
由外积形态，右端每一项的像均包含于 \(R_t\cdot(\pi_\varphi^{\,r}\overline q_r)\)，故像模由所列向量生成；
最小性由推论 \ref{cor:xi-symq-golden-cq-mod-pi-power-staircase} 给出。证毕。
\end{proof}

\begin{conclusion}[模 \texorpdfstring{$\pi_\varphi^t$}{pi^t} 下核、像、余核的完全同构型；\texorpdfstring{$\ker\cong\mathrm{coker}$}{ker≅coker} 为严格结构恒等]\label{con:xi-symq-golden-cq-mod-pi-power-kernel-image-cokernel}
令 \(R_t:=\mathcal O/(\pi_\varphi^t)\)，并记 \(C_{q,t}:R_t^{q+1}\to R_t^{q+1}\) 为由 \(C_q\) 诱导的 \(R_t\)-线性同态。
则存在 \(R_t\)-模同构
\[
\boxed{\;
\mathrm{Im}(C_{q,t})\ \cong\ \bigoplus_{s=1}^{t}\mathcal O/(\pi_\varphi^{s})\;}
\]
\[
\boxed{\;
\ker(C_{q,t})\ \cong\ \Bigl(\bigoplus_{s=1}^{t-1}\mathcal O/(\pi_\varphi^{s})\Bigr)\ \oplus\ \bigl(\mathcal O/(\pi_\varphi^{t})\bigr)^{q+1-t}\;}
\]
\[
\boxed{\;
\mathrm{coker}(C_{q,t})\ \cong\ \Bigl(\bigoplus_{s=1}^{t-1}\mathcal O/(\pi_\varphi^{s})\Bigr)\ \oplus\ \bigl(\mathcal O/(\pi_\varphi^{t})\bigr)^{q+1-t}\;}
\]
从而对每个 \(t\) 都有严格同构 \(\ker(C_{q,t})\cong \mathrm{coker}(C_{q,t})\)。
\end{conclusion}
\begin{proof}
由 Smith 正规形，\(C_{q,t}\) 在 \(R_t\) 上与对角算子
\(
(x_r)_{r=0}^{q}\mapsto (\pi_\varphi^{\,r}x_r)_{r=0}^{q}
\)
幺模等价。
因此像在第 \(r\) 分量为 \(\pi_\varphi^{\,r}R_t\)；当 \(r<t\) 时乘法 \(R_t\to\pi_\varphi^{\,r}R_t\) 的核为 \(\pi_\varphi^{t-r}R_t\)，故 \(\pi_\varphi^{\,r}R_t\cong R_t/(\pi_\varphi^{t-r})\cong \mathcal O/(\pi_\varphi^{t-r})\)，而当 \(r\ge t\) 时像分量为 \(0\)。
将 \(r=0,\dots,t-1\) 重标为 \(s=t-r\) 得到像的直和分解。
\par
核在第 \(r\) 分量为 \(\{x\in R_t:\ \pi_\varphi^{\,r}x=0\}\)。
当 \(r<t\) 时该集合为 \(\pi_\varphi^{t-r}R_t\cong R_t/(\pi_\varphi^{r})\cong \mathcal O/(\pi_\varphi^{r})\)；
当 \(r\ge t\) 时 \(\pi_\varphi^{\,r}=0\) 于 \(R_t\)，故核分量为 \(R_t\cong \mathcal O/(\pi_\varphi^{t})\)。
汇总得到核的分解。
\par
余核在第 \(r\) 分量为 \(R_t/\pi_\varphi^{\,r}R_t\)。
当 \(r<t\) 时同构于 \(\mathcal O/(\pi_\varphi^{r})\)，而当 \(r\ge t\) 时同构于 \(R_t\cong \mathcal O/(\pi_\varphi^{t})\)，故余核与核同构。证毕。
\end{proof}

\begin{conclusion}[限制标量后的逐级块化分解：\texorpdfstring{$\widetilde C_q$}{Ctilde} 的整数幺模对角骨架可显式写出]\label{con:xi-symq-golden-cq-restriction-of-scalars-block-diagonal}
定义乘法表示 \(\iota:\mathcal O\to\Mat_2(\ZZ)\) 为：对 \(a+b\varphi\in\mathcal O\),
\[
\iota(a+b\varphi):=
\begin{pmatrix}a&b\\[2pt]b&a+b\end{pmatrix}.
\]
令 \(A:=\iota(\pi_\varphi)=\bigl(\begin{smallmatrix}2&1\\[2pt]1&3\end{smallmatrix}\bigr)\)，则 \(\iota(\pi_\varphi^{\,r})=A^{r}\)。
对任意 \(\mathcal O\)-矩阵 \(M=[m_{ij}]\)，定义其限制标量矩阵 \(\widetilde M\in\Mat_{2n}(\ZZ)\) 为以 \(2\times 2\) 块 \(\iota(m_{ij})\) 组成的整数矩阵（其中 \(M\in\Mat_n(\mathcal O)\)）。
\par
则沿结论 \ref{con:xi-symq-golden-cq-smith-staircase} 的幺模分解，有严格恒等式
\[
\boxed{\;
\widetilde C_q
=\widetilde Q_q\cdot \mathrm{diag}(A^{0},A^{1},\dots,A^{q})\cdot \widetilde U_q\;}
\]
其中 \(\widetilde Q_q,\widetilde U_q\in\mathrm{GL}_{2(q+1)}(\ZZ)\) 为整数幺模矩阵。
\end{conclusion}
\begin{proof}
\(\iota\) 为环同态，故对块化构造满足 \(\widetilde{MN}=\widetilde M\,\widetilde N\)。
对角阶梯阵的块化为
\[
\widetilde{\mathrm{diag}(1,\pi_\varphi,\dots,\pi_\varphi^q)}
=\mathrm{diag}\bigl(\iota(1),\iota(\pi_\varphi),\dots,\iota(\pi_\varphi^q)\bigr)
=\mathrm{diag}(A^0,A^1,\dots,A^q),
\]
从而由结论 \ref{con:xi-symq-golden-cq-smith-staircase} 得到所示分解。
又由于 \(Q_q,U_q\in\mathrm{GL}_{q+1}(\mathcal O)\) 作为 \(\mathcal O\)-线性自同构在限制标量后仍为 \(\ZZ\)-线性自同构，故 \(\widetilde Q_q,\widetilde U_q\) 为整数幺模矩阵。证毕。
\end{proof}

\begin{conclusion}[\texorpdfstring{$\varpi^m$}{pi^m}-自逆矩阵的 Smith 对称律：离散赋值环上的互补指数约束]\label{con:xi-dvr-selfinverse-smith-symmetry}
令 \(R\) 为离散赋值环，\(\varpi\) 为其均匀元，\(K\) 为分式域。
设 \(M\in\Mat_n(R)\cap\mathrm{GL}_n(K)\) 满足
\[
M^2=\varpi^{m}I_n
\qquad(\text{等价地 }M^{-1}=\varpi^{-m}M).
\]
设 \(M\) 的 Smith 正规形不变指数为 \(0\le s_1\le\cdots\le s_n\)，即 \(M\) 在 \(R\) 上幺模等价于
\(\mathrm{diag}(\varpi^{s_1},\dots,\varpi^{s_n})\)。
则必有对称律
\[
\boxed{\;
s_i+s_{n+1-i}=m,\qquad 1\le i\le n.\;}
\]
\end{conclusion}
\begin{proof}
取 Smith 分解 \(M=U\,\mathrm{diag}(\varpi^{s_1},\dots,\varpi^{s_n})\,V\)（\(U,V\in\mathrm{GL}_n(R)\)）。
则
\[
\varpi^{m}M^{-1}
=\varpi^{m}V^{-1}\,\mathrm{diag}(\varpi^{-s_1},\dots,\varpi^{-s_n})\,U^{-1}
=V^{-1}\,\mathrm{diag}(\varpi^{m-s_1},\dots,\varpi^{m-s_n})\,U^{-1}.
\]
由假设 \(\varpi^{m}M^{-1}=M\) 且 Smith 指数在 \(R\) 上除重排外唯一，指数多重集 \(\{s_1,\dots,s_n\}\) 与 \(\{m-s_1,\dots,m-s_n\}\) 相同。
将二者按升序排列即得 \(s_i=m-s_{n+1-i}\)。证毕。
\end{proof}

\begin{proposition}[阶梯 Smith 模板：对称幂与线性型同余导出的幺模分解]\label{prop:xi-smith-staircase-template-symq}
设 \(R\) 为交换整环，\(\pi\in R\) 为非零非单位元。
取
\(
H=\begin{pmatrix}a&b\\ c&d\end{pmatrix}\in\mathrm{SL}_2(R)
\)
并令 \(\ell=ax+cy\)、\(m=bx+dy\)。
若存在单位元 \(u\in R^\times\) 与另一线性型 \(\ell_2\) 满足
\(
\ell_2=u\,\ell-\pi\,m
\)，
则对任意 \(q\ge 1\)，由
\(
P_j:=\ell^{\,q-j}\ell_2^{\,j}
\)
生成的系数矩阵 \(C_q(R)\) 具有双侧幺模分解
\[
\boxed{\;
C_q(R)=\mathrm{Sym}^q(H)\cdot \mathrm{diag}(1,\pi,\pi^2,\dots,\pi^q)\cdot U_q(R)\;}
\]
其中 \(U_q(R)\) 为上三角矩阵并满足
\[
U_q(R)_{rj}=
\begin{cases}
(-1)^r\,u^{\,j-r}\binom{j}{r},& r\le j,\\
0,& r>j.
\end{cases}
\]
若进一步 \(R\) 为 PID，则 \(C_q(R)\) 的 Smith 正规形为 \(\mathrm{diag}(1,\pi,\dots,\pi^q)\)。
\end{proposition}
\begin{proof}
由二项式展开
\(
\ell_2^{j}=(u\ell-\pi m)^{j}
\)
得到 \(P_j=\sum_{r\le j}(-1)^r u^{j-r}\binom{j}{r}\pi^r\,Q_r\)（其中 \(Q_r=\ell^{q-r}m^r\)），从而给出所示矩阵分解。
由于 \(\mathrm{Sym}^q(H)\in\mathrm{GL}_{q+1}(R)\) 且 \(U_q(R)\) 对角元为 \(((-1)^j)_{0\le j\le q}\subset R^\times\)，二者均幺模。
PID 情形下 Smith 形由幺模等价唯一性推出。证毕。
\end{proof}

\input{subsubsec__xi-time-protocol-conclusions-part31i-exterior-power-smith-schubert}

\input{subsubsec__xi-time-protocol-conclusions-part31i-integer-blockization}

\endinput


\begin{conclusion}[固定零计数的迹极值律：Fibonacci--二幂标度的尖锐上确界]\label{con:xi-terminal-replica-softcore-fixed-zero-count-trace-extremal}
固定整数 \(m\ge 2\)。
对任意字词 \(w=A_1\cdots A_m\in\{J,K\}^m\) 且含至少一个 \(J\)，记
\[
r:=\#\{t:A_t=K\}\in\{0,1,\dots,m-1\},
\qquad
M(w):=A_1\cdots A_m.
\]
则有尖锐上界
\[
\Tr M(w)\le F_{r+3}\,2^{m-r-1},
\]
且等号成立当且仅当 \(w\) 在循环移位意义下等价于 \(J^{m-r}K^{r}\)（亦即所有 \(K\) 在某个代表元中汇聚为单一连续簇）。
因此，在结论 \ref{con:xi-terminal-replica-softcore-exceptional-power-sum-linear-word-closed-form} 的指数族中，
底数按 \(r\) 分层，其第 \(r\) 层的最大底数精确为 \(F_{r+3}2^{m-r-1}\)。
\end{conclusion}
\begin{proof}[证明要点]
由于迹对循环移位不变，可将 \(w\) 视为二元 necklace 并应用结论 \ref{con:xi-terminal-replica-softcore-binary-necklace-trace-formula} 的分解：
若 \(w\) 含 \(m-r\) 个 \(J\)，则存在 \(r_1,\dots,r_{m-r}\ge 0\) 使 \(\sum_i r_i=r\) 且
\(
\Tr M(w)=\prod_{i=1}^{m-r}F_{r_i+3}.
\)
利用 Fibonacci 的乘法上界
\(
F_{a+3}F_{b+3}\le 2F_{a+b+3}
\)
（\(a,b\ge 0\)），迭代合并零簇得到
\(
\prod_{i}F_{r_i+3}\le 2^{m-r-1}F_{r+3}
\)。
当且仅当除一项外其余 \(r_i=0\) 时取等，从而对应于单簇形 \(J^{m-r}K^{r}\)。证毕。
\end{proof}
\begin{conclusion}[含 \(J\) 字词迹的全局极值与等号刻画]\label{con:xi-terminal-replica-softcore-jk-word-trace-global-extrema}
固定整数 \(m\ge 1\)。
令
\[
\mathcal T_m:=\{\Tr M(w): w\in\{J,K\}^m,\ w\text{ 含 }J\}\subset\ZZ_{\ge 0},
\]
其中 \(M(w)\) 定义见结论 \ref{con:xi-terminal-replica-softcore-exceptional-power-sum-word-trace-sympower}。
则有严格极值律
\[
\boxed{\;\min\mathcal T_m=F_{m+2},\qquad \max\mathcal T_m=2^m.\;}
\]
并且 \(\Tr M(w)=F_{m+2}\) 当且仅当 \(w\) 恰含一个 \(J\)（位置任意），亦即 \(w=K^{a}JK^{m-1-a}\)（\(0\le a\le m-1\)）；
而 \(\Tr M(w)=2^m\) 当且仅当 \(w=J^m\)。
\end{conclusion}
\begin{proof}
当 \(m=1\) 时结论由 \(\Tr(J)=2=F_3\) 直接成立；以下设 \(m\ge 2\)。
\par
\textbf{(i) 最大值与等号刻画.}
字词 \(w=J^m\) 满足 \(\Tr M(w)=\Tr(J^m)=2^m\)。
若 \(w\neq J^m\) 且含 \(J\)，则其 \(K\) 计数满足 \(r\ge 1\)；
由结论 \ref{con:xi-terminal-replica-softcore-fixed-zero-count-trace-extremal} 得
\[
\Tr M(w)\le F_{r+3}\,2^{m-r-1}\le F_4\,2^{m-2}=3\cdot 2^{m-2}<2^m,
\]
故最大值唯一且仅由 \(w=J^m\) 取到。
\par
\textbf{(ii) 最小值与等号刻画.}
由结论 \ref{con:xi-terminal-replica-softcore-trace-partition-image}，
任意含 \(J\) 的迹值均可写为 \(B_\lambda=\prod_{i}F_{\lambda_i+2}\)，其中 \(\lambda\vdash m\)。
若 \(\lambda=(m)\)，则 \(B_\lambda=F_{m+2}\)；
若 \(\lambda\) 含至少两个部件，则对任意 \(a,b\ge 1\) 有
\[
F_{a+2}F_{b+2}=F_{a+b+2}+F_aF_b>F_{a+b+2},
\]
从而反复合并部件得到 \(B_\lambda>F_{m+2}\)。
因此 \(\min\mathcal T_m=F_{m+2}\)，且仅由 \(\lambda=(m)\) 取到。
该分拆对应于字词恰含一个 \(J\) 的情形；由迹的循环不变性可将其写成 \(K^{a}JK^{m-1-a}\)（\(0\le a\le m-1\)）。
证毕。
\end{proof}
\endinput


\subsubsection{例外谱的秩一交错、Perron 根的首项夹逼与 \texorpdfstring{$\mathrm{rank}\,1$}{rank-1} 谱塌缩：固定张量谱的可核验谱护栏}
沿用命题 \ref{prop:pom-multiplicity-composition-sq-symmetry-reduction} 的例外块 \(A^{(q)}\)（定义 \eqref{eq:pom-multiplicity-composition-Aq-def}）及其谱
\(\{\lambda_i^{(q)}\}_{i=0}^{q}\)（按降序排列）。
本段把例外谱与固定张量谱 \(\{(-1)^k\varphi^{q-2k}/2\}_{k=0}^{q}\) 的关系压缩为一条严格交错链，并由此导出符号惯性、Perron 根的可核验首项夹逼，以及非 Perron 谱相对于主特征值的严格塌缩。

\begin{theorem}[秩一扰动的严格交错：例外谱夹在固定张量谱之间]\label{thm:xi-replica-softcore-exceptional-interlacing}
\leanverified{paper\_xi\_replica\_softcore\_exceptional\_interlacing}
固定整数 \(q\ge 1\)，记 \(b_k:=\binom{q}{k}\) 并令 \(\mathbf{1}:=(1,\dots,1)^{\mathsf T}\in\RR^{q+1}\)、\(b:=(b_0,\dots,b_q)^{\mathsf T}\)。
定义
\[
C^{(q)}_{k,\ell}:=\binom{q-k}{\ell}\qquad(0\le k,\ell\le q),
\qquad
A^{(q)}=\frac12\bigl(\mathbf{1}b^{\mathsf T}+C^{(q)}\bigr),
\]
并令
\[
D:=\mathrm{diag}\bigl(\sqrt{b_0},\dots,\sqrt{b_q}\bigr),\qquad
u:=D\mathbf{1}\in\RR^{q+1},
\qquad
\widetilde C^{(q)}:=D C^{(q)}D^{-1},\qquad
\widetilde A^{(q)}:=D A^{(q)}D^{-1}.
\]
则 \(\widetilde C^{(q)}\) 为实对称矩阵，且
\[
\widetilde A^{(q)}=\frac12\Bigl(\widetilde C^{(q)}+uu^{\mathsf T}\Bigr).
\]
记 \(\nu_0>\nu_1>\cdots>\nu_q\) 为 \(\frac12\widetilde C^{(q)}\) 的全部特征值（按降序），
并记 \(\lambda_0^{(q)}>\lambda_1^{(q)}>\cdots>\lambda_q^{(q)}\) 为 \(A^{(q)}\) 的全部特征值（亦即 \(\widetilde A^{(q)}\) 的特征值）。
则成立严格交错链
\[
\lambda_0^{(q)}>\nu_0>\lambda_1^{(q)}>\nu_1>\cdots>\lambda_q^{(q)}>\nu_q.
\]
并且谱集合满足
\[
\{\nu_j\}_{j=0}^{q}=\left\{\frac{(-1)^k\varphi^{\,q-2k}}{2}:0\le k\le q\right\}
\qquad(\text{仅差排列}).
\]
\end{theorem}
\begin{proof}
由 \(C^{(q)}=B_q^{\mathsf T}\)（其中 \(B_q(i,j)=\binom{q-j}{i}\)）与命题 \ref{prop:pom-bq-weighted-selfadjoint-binomial} 的加权自伴随性
\(
W_qB_q=B_q^{\mathsf T}W_q
\),
其中 \(W_q=\mathrm{diag}(b_0^{-1},\dots,b_q^{-1})=D^{-2}\)，可得
\[
\widetilde C^{(q)}
=D B_q^{\mathsf T}D^{-1}
=\bigl(D^{-1}B_qD\bigr)^{\mathsf T}
=D^{-1}B_qD
=\bigl(\widetilde C^{(q)}\bigr)^{\mathsf T},
\]
故 \(\widetilde C^{(q)}\) 对称。
又因
\(
\widetilde A^{(q)}=D\cdot \frac12(\mathbf{1}b^{\mathsf T}+C^{(q)})\cdot D^{-1}
=\frac12(\widetilde C^{(q)}+uu^{\mathsf T})
\)
且 \(u=D\mathbf{1}\)，故所示分解成立。
\par
由命题 \ref{prop:pom-bq-is-symq-and-spectrum}，
\(
\mathrm{spec}(B_q)=\{(-1)^k\varphi^{q-2k}\}_{k=0}^{q}
\)
且全为单根；因此
\(
\mathrm{spec}(C^{(q)})=\mathrm{spec}(B_q^{\mathsf T})=\mathrm{spec}(B_q)
\)
而 \(\widetilde C^{(q)}\) 与 \(C^{(q)}\) 相似，故
\(
\mathrm{spec}(\tfrac12\widetilde C^{(q)})=\{(-1)^k\varphi^{q-2k}/2\}_{k=0}^{q}
\),
得到 \(\{\nu_j\}\) 的谱集合表达。
\par
令 \(B:=\frac12\widetilde C^{(q)}\)，则 \(B\) 为实对称且单谱。
由经典秩一扰动谱理论，\(B+\frac12 uu^{\mathsf T}\) 的特征值与 \(B\) 交错；
并且若 \(u\) 与 \(B\) 的每个本征向量都不正交，则交错为严格。
下证该非正交性。
\par
令 \(v^{(k)}\) 为 \(B_q\) 对应本征值 \(\mu_k=(-1)^k\varphi^{q-2k}\) 的本征向量，并按命题 \ref{prop:pom-bq-explicit-diagonalization-linear-forms} 取其为多项式本征向量
\(
P_{q,k}(x,y)=(\varphi x+y)^{q-k}(\psi x+y)^k
\)
（\(\psi=-\varphi^{-1}\)）在基 \(\{x^{q-j}y^j\}_{j=0}^{q}\) 下的坐标向量。
由 \(W_qB_q=B_q^{\mathsf T}W_q\) 可知
\(
w^{(k)}:=W_qv^{(k)}
\)
为 \(C^{(q)}=B_q^{\mathsf T}\) 的本征向量且本征值仍为 \(\mu_k\)；
因此 \(Dw^{(k)}\) 为 \(\widetilde C^{(q)}=DC^{(q)}D^{-1}\) 的本征向量。
于是
\[
\bigl\langle u, Dw^{(k)}\bigr\rangle
=(D\mathbf{1})^{\mathsf T}(Dw^{(k)})
=\mathbf{1}^{\mathsf T}D^{2}w^{(k)}
=\mathbf{1}^{\mathsf T}\mathrm{diag}(b)\,W_qv^{(k)}
=\mathbf{1}^{\mathsf T}v^{(k)}.
\]
而 \(\mathbf{1}^{\mathsf T}v^{(k)}\) 等于 \(P_{q,k}(1,1)\)（系数求和即在 \((1,1)\) 处取值），从而
\[
P_{q,k}(1,1)=(\varphi+1)^{q-k}(\psi+1)^{k}=(\varphi^2)^{q-k}(\varphi^{-2})^{k}=\varphi^{2(q-2k)}\neq 0.
\]
故 \(\langle u, Dw^{(k)}\rangle\neq 0\) 对所有 \(k\) 成立，严格交错成立。证毕。
\end{proof}

\begin{corollary}[符号惯性与零点排除]\label{cor:xi-replica-softcore-exceptional-inertia}
对任意整数 \(q\ge 1\)，例外谱 \(\{\lambda_i^{(q)}\}_{i=0}^{q}\) 的符号分布为：
\[
\#\{i:\lambda_i^{(q)}>0\}=\left\lfloor\frac{q}{2}\right\rfloor+1,
\qquad
\#\{i:\lambda_i^{(q)}<0\}=\left\lceil\frac{q}{2}\right\rceil.
\]
尤其地，\(0\notin \mathrm{spec}(A^{(q)})\)。
\end{corollary}
\begin{proof}
定理 \ref{thm:xi-replica-softcore-exceptional-interlacing} 已给出
\(
\mathrm{spec}(\frac12\widetilde C^{(q)})=\{(-1)^k\varphi^{q-2k}/2\}
\)
且全为非零单根，因此其负特征值个数恰为 \(\lceil q/2\rceil\)。
又因 \(\widetilde A^{(q)}=\frac12(\widetilde C^{(q)}+uu^{\mathsf T})\) 是对称矩阵加上半正定秩一扰动，
由惯性定理（或等价地：秩一扰动至多改变一个符号）可知 \(\widetilde A^{(q)}\) 的负特征值个数只能取 \(\lceil q/2\rceil\) 或 \(\lceil q/2\rceil-1\)。
\par
另一方面，由定理 \ref{thm:pom-replica-softcore-exceptional-spectrum-product}，
\(
\det(A^{(q)})=(-1)^{q(q+1)/2}2^{-q}\neq 0
\)，
故 \(\widetilde A^{(q)}\) 亦可逆且无零特征值；
其行列式符号决定负特征值个数的奇偶性。
而 \(\lceil q/2\rceil\) 与 \(\lceil q/2\rceil-1\) 奇偶性相反，
因此只能取 \(\lceil q/2\rceil\)。
由维数为 \(q+1\) 得正特征值个数为 \(\lfloor q/2\rfloor+1\)，并立即推出零点排除。证毕。
\end{proof}

\begin{theorem}[Perron 根的首项夹逼与极限]\label{thm:xi-replica-softcore-perron-leading-term-bounds}
记 Perron 根（最大特征值）为 \(\rho_q:=\lambda_0^{(q)}=\rho(A^{(q)})\)。
则对任意整数 \(q\ge 1\) 有双侧界
\[
2^{q-1}+\frac12\Bigl(\frac32\Bigr)^{q}
\ \le\ \rho_q\ \le\
\frac12\sqrt{4^{q}+2\cdot 3^{q}+F_{2q+2}},
\]
其中 \(F_n\) 为 Fibonacci 数列（\(F_0=0,F_1=1\)）。
并且首项渐近满足
\[
\lim_{q\to\infty}\frac{\rho_q-2^{q-1}}{(3/2)^q}=\frac12.
\]
\end{theorem}
\begin{proof}
取定理 \ref{thm:xi-replica-softcore-exceptional-interlacing} 的相似对称表示
\(
\widetilde A^{(q)}=\frac12(\widetilde C^{(q)}+uu^{\mathsf T})
\)；
其谱与 \(A^{(q)}\) 相同，故 \(\rho_q=\lambda_{\max}(\widetilde A^{(q)})\)。
由 Rayleigh 商得
\[
\rho_q\ge \frac{u^{\mathsf T}\widetilde A^{(q)}u}{u^{\mathsf T}u}
=\frac12\cdot\frac{u^{\mathsf T}\widetilde C^{(q)}u+(u^{\mathsf T}u)^2}{u^{\mathsf T}u}.
\]
注意 \(u^{\mathsf T}u=\sum_{k=0}^{q}\binom{q}{k}=2^{q}\)。
又
\[
u^{\mathsf T}\widetilde C^{(q)}u
=(D\mathbf{1})^{\mathsf T}(DC^{(q)}D^{-1})(D\mathbf{1})
=\mathbf{1}^{\mathsf T}D^{2}C^{(q)}\mathbf{1}
=b^{\mathsf T}C^{(q)}\mathbf{1}.
\]
而 \(\bigl(C^{(q)}\mathbf{1}\bigr)_k=\sum_{\ell=0}^{q}\binom{q-k}{\ell}=2^{q-k}\)，故
\[
b^{\mathsf T}C^{(q)}\mathbf{1}
=\sum_{k=0}^{q}\binom{q}{k}2^{q-k}=(2+1)^{q}=3^{q}.
\]
代回即得下界
\(
\rho_q\ge \frac12(3^{q}+4^{q})/2^{q}=2^{q-1}+\frac12(3/2)^q
\)。
\par
上界由二次矩给出：\(\rho_q^2\le \sum_{i=0}^{q}(\lambda_i^{(q)})^2=S_2(q)\)。
而结论 \ref{con:xi-terminal-replica-softcore-exceptional-spectrum-square-sum} 给出
\(
S_2(q)=4^{q-1}+\frac12\,3^{q}+\frac14F_{2q+2}
=\frac14(4^{q}+2\cdot 3^{q}+F_{2q+2})
\)，
故得所示上界。
\par
最后的极限：将上界写为
\[
\rho_q\le 2^{q-1}\sqrt{1+\varepsilon_q},
\qquad
\varepsilon_q:=2\Bigl(\frac34\Bigr)^{q}+\frac{F_{2q+2}}{4^{q}}.
\]
由不等式 \(\sqrt{1+\varepsilon}-1\le \varepsilon/2\)（\(\varepsilon\ge 0\)）得
\[
0\le \rho_q-2^{q-1}
\le
2^{q-1}\cdot\frac{\varepsilon_q}{2}
=\frac12\Bigl(\frac32\Bigr)^{q}+\frac{F_{2q+2}}{2^{q+2}}.
\]
结合下界
\(
\rho_q-2^{q-1}\ge \frac12(3/2)^q
\)
得到
\[
\frac12
\le
\frac{\rho_q-2^{q-1}}{(3/2)^q}
\le
\frac12+\frac{F_{2q+2}}{4\cdot 3^q}.
\]
而 \(F_{2q+2}=O(\varphi^{2q})\) 且 \(\varphi^2/3<1\)，故右端余项趋于 \(0\)；
夹逼即得极限为 \(1/2\)。证毕。
\end{proof}

\begin{theorem}[例外块的 \texorpdfstring{$\mathrm{rank}\,1$}{rank-1} 谱塌缩与迹主项]\label{thm:xi-replica-softcore-rankone-spectral-collapse}
沿用定理 \ref{thm:xi-replica-softcore-exceptional-interlacing} 与
\ref{thm:xi-replica-softcore-perron-leading-term-bounds} 的记号，并定义
\[
\Lambda_q:=\max_{1\le i\le q}|\lambda_i^{(q)}|.
\]
则对一切 \(q\ge 1\) 都有
\[
\boxed{\;
\Lambda_q\le \frac{\varphi^q}{2}.\;}
\]
从而
\[
\boxed{\;
\frac{\Lambda_q}{\rho_q}
\le
\left(\frac{\varphi}{2}\right)^q(1+o(1))
\longrightarrow 0.\;}
\]
另一方面，Perron 根满足精确归一化
\[
\boxed{\;
\frac{\rho_q}{2^{q-1}}
=
1+\left(\frac34\right)^q
+o\!\left(\left(\frac34\right)^q\right).\;}
\]
进一步，对任意固定整数 \(m\ge 1\)，有
\[
\boxed{\;
\Tr\!\bigl((A^{(q)})^m\bigr)
=
\rho_q^{\,m}
+O\!\left(q\left(\frac{\varphi^q}{2}\right)^m\right).\;}
\]
特别地，
\[
\boxed{\;
\frac{\Tr\!\bigl((A^{(q)})^m\bigr)}{2^{m(q-1)}}\longrightarrow 1
\qquad(q\to\infty).\;}
\]
\end{theorem}
\begin{proof}
由定理 \ref{thm:xi-replica-softcore-exceptional-interlacing}，对每个 \(1\le i\le q\) 都有
\[
\nu_i<\lambda_i^{(q)}<\nu_{i-1},
\]
其中
\[
\{\nu_j\}_{j=0}^{q}
=
\left\{\frac{(-1)^k\varphi^{q-2k}}{2}:0\le k\le q\right\}.
\]
因此
\[
|\lambda_i^{(q)}|
\le
\max\{|\nu_{i-1}|,|\nu_i|\}
\le
\frac{\varphi^q}{2},
\]
从而得到 \(\Lambda_q\le \varphi^q/2\)。

再由定理 \ref{thm:xi-replica-softcore-perron-leading-term-bounds}，
\[
\rho_q
=
2^{q-1}
+\frac12\left(\frac32\right)^q
+o\!\left(\left(\frac32\right)^q\right),
\]
故
\[
\frac{\rho_q}{2^{q-1}}
=
1+\frac{\frac12(3/2)^q}{2^{q-1}}
+o\!\left(\frac{(3/2)^q}{2^{q-1}}\right)
=
1+\left(\frac34\right)^q+o\!\left(\left(\frac34\right)^q\right).
\]
又因为 \(\rho_q=2^{q-1}(1+o(1))\)，便有
\[
\frac{\Lambda_q}{\rho_q}
\le
\frac{\varphi^q/2}{2^{q-1}(1+o(1))}
=
\left(\frac{\varphi}{2}\right)^q(1+o(1))
\longrightarrow 0.
\]

最后，由谱分解
\[
\Tr\!\bigl((A^{(q)})^m\bigr)
=
\sum_{i=0}^{q}\lambda_i^{(q)\,m}
=
\rho_q^{\,m}+\sum_{i=1}^{q}\lambda_i^{(q)\,m}.
\]
利用上界 \(|\lambda_i^{(q)}|\le \varphi^q/2\) 得
\[
\left|\sum_{i=1}^{q}\lambda_i^{(q)\,m}\right|
\le
q\left(\frac{\varphi^q}{2}\right)^m,
\]
于是
\[
\Tr\!\bigl((A^{(q)})^m\bigr)
=
\rho_q^{\,m}
+O\!\left(q\left(\frac{\varphi^q}{2}\right)^m\right).
\]
两边再除以 \(2^{m(q-1)}\)。由上式的 Perron 归一化可知
\(
\rho_q^{\,m}/2^{m(q-1)}\to 1
\)，并且
\[
\frac{q(\varphi^q/2)^m}{2^{m(q-1)}}
=
q\left(\frac{\varphi}{2}\right)^{mq}
\longrightarrow 0.
\]
故所示极限成立。证毕。
\end{proof}

\endinput


\begin{conclusion}[\texorpdfstring{$\Omega_m(q)$}{Omega\_m(q)} 的表示论闭式与二阶整系数递推]\label{con:xi-replica-softcore-exceptional-symq-evaluation-coupling}
令
\[
M:=\begin{pmatrix}1&1\\[2pt]1&0\end{pmatrix},\qquad
\mathcal V_q:=\mathrm{Sym}^q(\RR^2),
\]
并定义
\[
(\mathcal B_q f)(x,y):=f(x+y,x),\qquad
(\mathcal E_q f)(x,y):=f(1,1)\,(x+y)^q.
\]
在基 \(\{x^{q-j}y^j\}_{j=0}^{q}\) 下，\(\mathcal B_q\) 的矩阵为
\[
B_q=(\binom{q-j}{i})_{0\le i,j\le q},
\]
且
\[
\frac12(\mathcal B_q+\mathcal E_q)\ \sim\ A^{(q)}.
\]
特别地，对任意整数 \(m\ge1,\ q\ge0\)，若 \(\alpha,\beta\) 为 \(M^m\) 的特征值（即
\(\alpha=\varphi^m,\ \beta=(-1)^m\varphi^{-m}\)），则
\[
\Omega_m(q):=\Tr(B_q^m)
=\Tr\!\bigl(\mathrm{Sym}^q(M^m)\bigr)
=\sum_{j=0}^{q}\alpha^{\,q-j}\beta^{\,j}.
\]
等价地，
\[
\Omega_m(q)=\sum_{j=0}^{q}(-1)^{mj}\varphi^{m(q-2j)}.
\]
因此 \(\Omega_m(q)\in\ZZ\)，并满足
\[
\Omega_m(0)=1,\qquad \Omega_m(1)=L_m,
\]
\[
\Omega_m(q)=L_m\,\Omega_m(q-1)-(-1)^m\,\Omega_m(q-2)\qquad(q\ge2).
\]
\end{conclusion}
\begin{proof}
取同构
\[
\Phi:\RR^{q+1}\to\mathcal V_q,\qquad
\Phi(c_0,\dots,c_q)(x,y):=\sum_{j=0}^{q}c_jx^{q-j}y^j.
\]
直接计算得
\[
\mathcal B_q(x^{q-j}y^j)=(x+y)^{q-j}x^j
=\sum_{i=0}^{q-j}\binom{q-j}{i}x^{q-i}y^i,
\]
故 \(\mathcal B_q\) 的系数矩阵是 \(B_q\)。
同理 \(\mathcal E_q\) 的矩阵为 \(b\mathbf1^{\mathsf T}\)（\(b_i=\binom qi\)），从而
\[
[\tfrac12(\mathcal B_q+\mathcal E_q)]_{\Phi}
=\tfrac12(B_q+b\mathbf1^{\mathsf T})
=(A^{(q)})^{\mathsf T},
\]
因此与 \(A^{(q)}\) 同谱。
又因 \(B_q=\mathrm{Sym}^q(M)\)，函子性给出
\[
B_q^m=\mathrm{Sym}^q(M)^m=\mathrm{Sym}^q(M^m).
\]
设 \(M^m\) 的特征值为 \(\alpha,\beta\)，则 \(\mathrm{Sym}^q(M^m)\) 的特征值为
\(\{\alpha^{q-j}\beta^j\}_{j=0}^{q}\)，其和即 \(\Omega_m(q)\) 的闭式。
最后由
\(\alpha+\beta=L_m\)、\(\alpha\beta=(-1)^m\)
及 \(h_q\) 的标准递推即得二阶关系。证毕。
\end{proof}

\begin{theorem}[负迹矩规范化与 Lucas--Pell 校正项的同一性]\label{thm:xi-exceptional-theta-omega-identification}
\leanverified{paper\_xi\_exceptional\_theta\_omega\_identification}
对任意整数 \(m\ge 1,\ q\ge 0\)，定义
\[
C^{(q)}_{k,\ell}:=\binom{q-k}{\ell}= (B_q^{\mathsf T})_{k,\ell},
\qquad
\Theta_m(q):=
\begin{cases}
\operatorname{tr}\!\bigl((C^{(q)})^{-m}\bigr),& m\ \text{为偶数},\\[1mm]
(-1)^q\operatorname{tr}\!\bigl((C^{(q)})^{-m}\bigr),& m\ \text{为奇数}.
\end{cases}
\]
则 \(\Theta_m(q)\in\ZZ\)，并满足
\[
\Theta_m(0)=1,\qquad
\Theta_m(1)=L_m,
\]
\[
\Theta_m(q)=L_m\,\Theta_m(q-1)-(-1)^m\,\Theta_m(q-2)\qquad(q\ge 2).
\]
此外，对所有 \(m,q\) 均有
\[
\Theta_m(q)=\Omega_m(q),
\]
其中 \(\Omega_m(q)\) 见结论 \ref{con:xi-replica-softcore-exceptional-symq-evaluation-coupling}。
\end{theorem}
\begin{proof}
由 \(C^{(q)}=B_q^{\mathsf T}\) 与转置不变性可得
\[
\operatorname{tr}\!\bigl((C^{(q)})^{-m}\bigr)
=\operatorname{tr}\!\bigl((B_q)^{-m}\bigr)
=\sum_{j=0}^{q}(-1)^{jm}\varphi^{-m(q-2j)}.
\]
若 \(m\) 为偶数，则
\[
\Theta_m(q)=\sum_{j=0}^{q}\varphi^{-m(q-2j)}
=\sum_{r=0}^{q}\varphi^{m(q-2r)}
=\Omega_m(q),
\]
其中第二步取换元 \(r=q-j\)。
若 \(m\) 为奇数，则
\[
\Theta_m(q)
=(-1)^q\sum_{j=0}^{q}(-1)^j\varphi^{-m(q-2j)}
=\sum_{r=0}^{q}(-1)^r\varphi^{m(q-2r)}
=\Omega_m(q),
\]
同样取 \(r=q-j\)。
于是 \(\Theta_m(q)=\Omega_m(q)\) 对全部 \(m,q\) 成立。
其整性、初值与二阶递推立即由结论 \ref{con:xi-replica-softcore-exceptional-symq-evaluation-coupling} 转移得到。证毕。
\end{proof}

\begin{conclusion}[例外谱的显式世俗方程与闭式本征向量]\label{con:xi-replica-softcore-exceptional-secular-explicit-weights}
令
\[
\varphi:=\frac{1+\sqrt5}{2},
\qquad
L_+(x,y):=\varphi x+y,
\qquad
L_-(x,y):=y-\varphi^{-1}x,
\]
并定义
\[
f_k(x,y):=L_+(x,y)^{q-k}L_-(x,y)^k
\qquad(k=0,1,\dots,q).
\]
则 \(\{f_k\}_{k=0}^{q}\) 构成 \(\mathcal V_q\) 的一组基，且
\[
\mathcal B_qf_k=\mu_kf_k,
\qquad
\mu_k:=(-1)^k\varphi^{q-2k}.
\]
再定义权重
\[
\beta_{q,k}:=\binom{q}{k}\frac{\varphi^{q-2k}}{(\sqrt5)^q},
\qquad
\omega_{q,k}:=\binom{q}{k}\frac{\varphi^{3q-6k}}{(\sqrt5)^q}>0.
\]
则 \(\lambda\in\RR\) 属于例外谱当且仅当满足
\[
\sum_{k=0}^{q}\frac{\omega_{q,k}}{2\lambda-\mu_k}=1,
\]
且任意此类 \(\lambda\) 的本征向量可写为
\[
v_\lambda=\sum_{k=0}^{q}c_kf_k,
\qquad
c_k\propto\frac{\beta_{q,k}}{2\lambda-\mu_k}.
\]
\end{conclusion}
\begin{proof}
由
\[
L_+(x+y,x)=\varphi(x+y)+x=\varphi(\varphi x+y)=\varphi L_+(x,y),
\]
\[
L_-(x+y,x)=x-\varphi^{-1}(x+y)=-\varphi^{-1}\bigl(y-\varphi^{-1}x\bigr)=-\varphi^{-1}L_-(x,y),
\]
得
\[
\mathcal B_qf_k=\varphi^{q-k}(-\varphi^{-1})^kf_k=\mu_kf_k.
\]
又由线性方程
\(
x+y=aL_++bL_-
\)
解得
\(
a=\varphi/\sqrt5,\ b=\varphi^{-1}/\sqrt5
\)，故
\[
(x+y)^q=\sum_{k=0}^{q}\binom{q}{k}a^{q-k}b^k\,L_+^{q-k}L_-^k
=\sum_{k=0}^{q}\beta_{q,k}f_k.
\]
并且
\[
f_k(1,1)=L_+(1,1)^{q-k}L_-(1,1)^k
=(\varphi^2)^{q-k}(\varphi^{-2})^k
=\varphi^{2q-4k},
\]
故
\[
\omega_{q,k}=f_k(1,1)\beta_{q,k}.
\]
若 \(v=\sum_kc_kf_k\) 且 \(\mathcal A_q^{\top}v=\lambda v\)，则
\[
\frac12\bigl(\mathcal B_qv+v(1,1)(x+y)^q\bigr)=\lambda v.
\]
比较 \(f_k\)-系数得
\[
(\mu_k-2\lambda)c_k+\beta_{q,k}\,v(1,1)=0,
\]
从而
\[
c_k=\frac{\beta_{q,k}\,v(1,1)}{2\lambda-\mu_k}.
\]
代回
\(
v(1,1)=\sum_kc_kf_k(1,1)
\)
并消去 \(v(1,1)\neq 0\) 即得世俗方程；向量系数同时得到。证毕。
\end{proof}

\begin{conclusion}[例外谱与 \texorpdfstring{$\{\mu_k/2\}$}{\{mu_k/2\}} 的严格交错与单根性]\label{con:xi-replica-softcore-exceptional-secular-strict-interlacing}
将极点集合 \(\{\mu_k/2\}_{k=0}^{q}\) 按降序排列为
\[
\pi_0>\pi_1>\cdots>\pi_q.
\]
则例外谱可按降序记为 \(\lambda_0>\lambda_1>\cdots>\lambda_q\)，并满足
\[
\lambda_0\in(\pi_0,\infty),
\qquad
\lambda_j\in(\pi_{j-1},\pi_j)\quad(j=1,\dots,q).
\]
因此例外特征值全部为实且互异，并且无根落在任一 \(\mu_k/2\) 上。
\end{conclusion}
\begin{proof}
记
\[
g_q(\lambda):=\sum_{k=0}^{q}\frac{\omega_{q,k}}{2\lambda-\mu_k}.
\]
因 \(\omega_{q,k}>0\)，在 \(\lambda\notin\{\mu_k/2\}\) 上有
\[
g_q'(\lambda)=-2\sum_{k=0}^{q}\frac{\omega_{q,k}}{(2\lambda-\mu_k)^2}<0.
\]
故 \(g_q\) 在每个连通分支上严格递减。
又在任一极点两侧 \(g_q\) 取 \(\pm\infty\) 跳变，且 \(\lambda\to+\infty\) 时 \(g_q(\lambda)\to 0\)。
因此方程 \(g_q(\lambda)=1\) 在 \((\pi_0,\infty)\) 上恰一根，在每个 \((\pi_{j-1},\pi_j)\) 上亦恰一根。
由 \(g_q'\neq 0\) 这些根全为单根。
结合结论 \ref{con:xi-replica-softcore-exceptional-secular-explicit-weights}，根集即例外谱。证毕。
\end{proof}

\begin{conclusion}[例外特征多项式的拉格朗日型闭式]\label{con:xi-replica-softcore-exceptional-characteristic-lagrange-form}
定义
\[
P_q(z):=\prod_{k=0}^{q}(z-\mu_k),
\qquad
Q_q(z):=\sum_{k=0}^{q}\omega_{q,k}\prod_{\substack{0\le j\le q\\ j\ne k}}(z-\mu_j).
\]
则例外特征多项式满足
\[
\chi_{\mathrm{exc}}^{(q)}(\lambda)
=2^{-(q+1)}\bigl(P_q(2\lambda)-Q_q(2\lambda)\bigr).
\]
等价地，\(P_q\) 给出纯张量谱节点 \(\{\mu_k\}\) 的主项，\(Q_q\) 给出评价耦合诱导的秩一修正项。
\end{conclusion}
\begin{proof}
在基 \(\{f_k\}\) 下，
\[
[\mathcal A_q^{\top}]
=\frac12\bigl(\mathrm{diag}(\mu_0,\dots,\mu_q)+\beta\,\ell^{\mathsf T}\bigr),
\]
其中 \(\beta_k=\beta_{q,k}\)、\(\ell_k=f_k(1,1)\)，且 \(\omega_{q,k}=\beta_k\ell_k\)。
故
\[
\chi_{\mathrm{exc}}^{(q)}(\lambda)
=\det(\lambda I-\mathcal A_q^{\top})
=2^{-(q+1)}\det\!\bigl(2\lambda I-\mathrm{diag}(\mu)-\beta\ell^{\mathsf T}\bigr).
\]
对秩一扰动应用行列式引理：
\[
\det\!\bigl(2\lambda I-\mathrm{diag}(\mu)-\beta\ell^{\mathsf T}\bigr)
=\det(2\lambda I-\mathrm{diag}(\mu))
\Bigl(1-\ell^{\mathsf T}(2\lambda I-\mathrm{diag}(\mu))^{-1}\beta\Bigr).
\]
第一因子即 \(P_q(2\lambda)\)，第二因子展开后恰给出
\[
P_q(2\lambda)-Q_q(2\lambda).
\]
代回即得。证毕。
\end{proof}

\begin{conclusion}[例外谱惯性定理：负特征值精确计数]\label{con:xi-replica-softcore-exceptional-secular-origin-rigidity}
对任意整数 \(q\ge2\)，\(A^{(q)}\) 无零特征值，且其负特征值个数满足
\[
\#\{0\le i\le q:\ \lambda_i^{(q)}<0\}=\left\lfloor\frac{q+1}{2}\right\rfloor.
\]
等价地，
\[
\#\{\lambda_i^{(q)}>0\}=\left\lceil\frac{q+1}{2}\right\rceil,\qquad
0\notin\mathrm{spec}(A^{(q)}).
\]
\end{conclusion}
\begin{proof}
由结论 \ref{con:xi-replica-softcore-exceptional-secular-strict-interlacing}，
方程 \(g_q(\lambda)=1\) 的根在极点
\(\pi_k=\mu_k/2\) 之间严格交错，其中
\[
\mu_k=(-1)^k\varphi^{q-2k}\neq0,\qquad k=0,\dots,q.
\]
因此例外谱全实、单根，且每个相邻极点区间恰含一根。
再由
\[
g_q(\lambda)=\sum_{k=0}^{q}\frac{\omega_{q,k}}{2\lambda-\mu_k},\qquad
\omega_{q,k}=\binom{q}{k}\frac{\varphi^{3q-6k}}{(\sqrt5)^q}>0,
\]
计算
\[
g_q(0)
=-\frac1{(\sqrt5)^q}\sum_{k=0}^{q}\binom{q}{k}(-1)^k\varphi^{2q-4k}
=-\frac{\varphi^{2q}}{(\sqrt5)^q}(1-\varphi^{-4})^q
=-1.
\]
在包含 \(0\) 的交错区间上，\(g_q\) 严格递减且 \(g_q(0)=-1<1\)，
故该区间内的根位于 \(0\) 的负侧。
其余区间中根号与极点号一致，故正根数等于正极点数。
又 \(\mu_k>0\iff k\) 为偶数，偶指标个数为 \(\lceil(q+1)/2\rceil\)，
于是负根数为 \((q+1)-\lceil(q+1)/2\rceil=\lfloor(q+1)/2\rfloor\)。
从而 \(0\) 非特征值；并与
\(\det(A^{(q)})=(-1)^{q(q+1)/2}2^{-q}\neq0\)
（定理 \ref{thm:pom-replica-softcore-exceptional-spectrum-product}）一致。证毕。
\end{proof}

\begin{conclusion}[张量谱节点处的显式避让值]\label{con:xi-replica-softcore-exceptional-node-avoidance-value}
对任意 \(k\in\{0,\dots,q\}\)，有
\[
\chi_{\mathrm{exc}}^{(q)}\!\left(\frac{\mu_k}{2}\right)
=-2^{-(q+1)}\,\omega_{q,k}\!\!\prod_{\substack{0\le j\le q\\ j\ne k}}\!(\mu_k-\mu_j)\neq 0.
\]
因此例外谱与 \(\{\mu_k/2\}_{k=0}^{q}\) 完全分离，且分离量由上式给出。
\end{conclusion}
\begin{proof}
把结论 \ref{con:xi-replica-softcore-exceptional-characteristic-lagrange-form} 的闭式在 \(2\lambda=\mu_k\) 处取值。
此时 \(P_q(\mu_k)=0\)，而 \(Q_q(\mu_k)\) 仅保留 \(k\)-项，故
\[
\chi_{\mathrm{exc}}^{(q)}(\mu_k/2)
=2^{-(q+1)}\bigl(0-Q_q(\mu_k)\bigr)
=-2^{-(q+1)}\,\omega_{q,k}\!\!\prod_{j\ne k}(\mu_k-\mu_j).
\]
由于 \(\omega_{q,k}>0\) 且 \(\{\mu_j\}\) 互异，右端非零。证毕。
\end{proof}

\input{subsubsec__xi-time-protocol-conclusions-part31q}
\input{subsubsec__xi-time-protocol-conclusions-part31v-bq-auditable-spectral-certificates}
\input{subsubsec__xi-time-protocol-conclusions-part31r}
\input{subsubsec__xi-time-protocol-conclusions-part31t}

\endinput


\begin{conclusion}[定边数约束下循环子图独立集极值与例外幂和主导渐近]\label{con:xi-terminal-replica-softcore-exceptional-moments-exponential-hierarchy}
令 \(m\ge3\)。对每个 \(0\le k\le m\)，有
\[
\boxed{\;
\max_{\substack{S\subseteq E(C_m)\\ |S|=k}} I\!\bigl(C_m[S]\bigr)=
\begin{cases}
2^m, & k=0,\\[3pt]
F_{k+3}\,2^{m-k-1}, & 1\le k\le m-1,\\[3pt]
L_m, & k=m.
\end{cases}
\;}
\]
当 \(1\le k\le m-1\) 时，取到最大值当且仅当 \(S\) 的 \(k\) 条边在 \(C_m\) 上构成单一连续路径（其余 \(m-k-1\) 个顶点孤立）。

并且，当 \(q\to\infty\) 时，
\[
\boxed{\;
S_m(q)=2^{m(q-1)}\left(1+m\left(\frac34\right)^q+O_m\!\left(\left(\frac58\right)^q\right)\right)\; } .
\]
\end{conclusion}
\begin{proof}
当 \(k=0\) 时 \(S=\varnothing\)，显然 \(I(C_m[S])=2^m\)。
当 \(k=m\) 时 \(S=E(C_m)\) 唯一，故 \(I(C_m[S])=I(C_m)=L_m\)。

下设 \(1\le k\le m-1\)。
此时 \(C_m[S]\) 无环，且每个连通分量为路径或孤立点。
设路径分量条数为 \(r\ge1\)，其边数为 \(\ell_1,\dots,\ell_r\ge1\)，满足
\(
\sum_{i=1}^{r}\ell_i=k
\)；
孤立点数为 \(m-(k+r)\)。
由路径独立集计数 \(I(P_{\ell+1})=F_{\ell+3}\) 得
\[
I\!\bigl(C_m[S]\bigr)=2^{m-(k+r)}\prod_{i=1}^{r}F_{\ell_i+3}.
\]
对 \(a,b\ge1\)，定义
\[
g_b:=2F_{a+b+3}-F_{a+3}F_{b+3}.
\]
由 Fibonacci 递推可得 \(g_{b+1}=g_b+g_{b-1}\)，且
\(
g_1=F_a\ge0,\ g_2=F_{a+1}\ge0
\)，故 \(g_b\ge0\)，即
\[
F_{a+3}F_{b+3}\le 2F_{a+b+3}.
\]
迭代合并 \(\ell_1,\dots,\ell_r\) 得
\[
\prod_{i=1}^{r}F_{\ell_i+3}\le 2^{r-1}F_{k+3}.
\]
代回即
\[
I\!\bigl(C_m[S]\bigr)\le 2^{m-(k+r)}\cdot 2^{r-1}F_{k+3}
=F_{k+3}\,2^{m-k-1}.
\]
当且仅当 \(r=1\) 且 \(\ell_1=k\) 取等，即 \(S\) 为连续 \(k\) 边路径。

再证渐近式。由结论
\ref{con:xi-terminal-replica-softcore-exceptional-power-sum-cycle-sector-replacement}，
\[
2^mS_m(q)=\Omega_m(q)+\sum_{S\subsetneq E(C_m)}I\!\bigl(C_m[S]\bigr)^q.
\]
其中 \(S=\varnothing\) 给出主项 \((2^m)^q\)；
\(|S|=1\) 给出次主项 \(m(3\cdot 2^{m-2})^q\)。
当 \(|S|=k\ge2\) 时，上述极值公式给出
\[
I\!\bigl(C_m[S]\bigr)\le F_{k+3}\,2^{m-k-1}\le 5\cdot 2^{m-3},
\]
其中第二个不等式由 \(F_{k+3}\le 5\cdot 2^{k-2}\)（\(k\ge2\)，可由 \(F_{n+1}\le2F_n\) 归纳）得到。
故所有 \(k\ge2\) 的贡献统一为 \(O_m((5\cdot 2^{m-3})^q)\)。
同时 \(|\Omega_m(q)|\ll_m \varphi^{mq}\)，且 \(m\ge3\) 时
\(
(\varphi/2)^m<(5/8)
\)，故 \(\Omega_m(q)\) 亦可并入同一余项尺度。
两边同除 \(2^m\) 即得
\[
S_m(q)=2^{m(q-1)}\left(1+m\left(\frac34\right)^q+O_m\!\left(\left(\frac58\right)^q\right)\right).
\]
\end{proof}

\begin{conclusion}[全迹矩的同阶主导渐近]\label{con:xi-terminal-replica-softcore-trace-moments-exponential-hierarchy}
固定 \(m\ge3\)。当 \(q\to\infty\) 时，
\[
\Tr\!\bigl((T^{(q)})^m\bigr)
=2^{m(q-1)}\left(1+m\left(\frac34\right)^q+O_m\!\left(\left(\frac58\right)^q\right)\right).
\]
\end{conclusion}
\begin{proof}[证明要点]
由结论 \ref{con:xi-terminal-replica-softcore-trace-exceptional-difference}，
\[
\Tr\!\bigl((T^{(q)})^m\bigr)
=S_m(q)+2^{-m}\bigl(L_m^q-\Omega_m(q)\bigr).
\]
又 \(L_m\ll \varphi^m\)、\(|\Omega_m(q)|\ll_m\varphi^{mq}\)，故差项满足
\[
2^{-m}\bigl|L_m^q-\Omega_m(q)\bigr|
=2^{m(q-1)}O_m\!\left(\left(\frac{\varphi}{2}\right)^{mq}\right)
=2^{m(q-1)}O_m\!\left(\left(\frac58\right)^q\right),
\]
并入上一结论的余项即可。
\end{proof}

\endinput


\begin{conclusion}[例外谱 \texorpdfstring{$\zeta$}{zeta}-函数的 Euler 乘积：原始二元 necklace 分解]\label{con:xi-terminal-replica-softcore-exceptional-spectrum-zeta-euler-product}
定义例外谱 \(\zeta\)-函数（形式幂级数意义）
\[
\mathcal Z_{\mathrm{exc}}^{(q)}(z)
:=\exp\!\left(\sum_{m\ge 1}\frac{S_m(q)}{m}z^m\right)
=\prod_{i=0}^{q}\frac{1}{1-\lambda_i^{(q)}z}
=\det\!\bigl(I_{q+1}-zA^{(q)}\bigr)^{-1}.
\]
令 \(\mathcal P^{(2)}\) 为所有含至少一个 \(1\) 的原始二元 necklace 之集合；
对 \(p\in\mathcal P^{(2)}\) 记其长度为 \(|p|\)，并以结论 \ref{con:xi-terminal-replica-softcore-binary-necklace-trace-formula} 定义 \(\tau(p)\)。
则存在严格的 Euler 乘积分解（形式恒等式）
\[
\boxed{\;
\mathcal Z_{\mathrm{exc}}^{(q)}(z)
=\det\!\Bigl(I_{q+1}-\frac z2\,B_q\Bigr)^{-1}
\prod_{p\in\mathcal P^{(2)}}\left(1-\tau(p)^q\left(\frac z2\right)^{|p|}\right)^{-1}\; } .
\]
该分解把例外谱的闭路矩信息拆解为两条互不混淆的通道：有限维二项式动力学通道 \(\det(I-\tfrac z2 B_q)^{-1}\)，以及由原始二元 necklace 索引的无限 Euler 乘积通道，其局部因子以 Fibonacci 乘法权重 \(\tau(p)^q\) 精确编码闭路增益。
\end{conclusion}
\begin{proof}[证明要点]
由 \(\log\det(I-zM)=-\sum_{m\ge 1}\Tr(M^m)z^m/m\) 得
\(
\det(I-\tfrac z2B_q)^{-1}=\exp(\sum_{m\ge 1}\Omega_m(q)(z/2)^m/m)
\)。
另一方面，结论 \ref{con:xi-terminal-replica-softcore-exceptional-power-sum-necklace-peeling} 把
\(
\sum_{m\ge 1}S_m(q)z^m/m
\)
分解为上述 \(\Omega_m\) 通道与一项
\(
\sum_{m\ge 1}\frac{1}{m}\sum_{\eta\ne 0^m}\#\mathrm{orb}(\eta)\bigl(\tau(\eta)^q(z/2)^m\bigr)
\)。
对后一项应用经典的 necklace--primitive 分解恒等式（Witt Euler 乘积；将所有 necklace 按其本原周期分解），即得 \(\prod_{p\in\mathcal P^{(2)}}(1-\tau(p)^q(z/2)^{|p|})^{-1}\)。
\end{proof}

\endinput

\paragraph{倒数矩通道的诱导表示刚性、分拆闭式与逐项卷积公式}
沿用结论链 \ref{subsubsec:xi-exceptional-spectrum-necklace-chain} 的记号。固定 \(q\ge 2\)，记
\[
C^{(q)}_{i,j}:=\binom{q-i}{j}\qquad(0\le i,j\le q),\qquad
A^{(q)}=\frac12\Bigl(\mathbf 1(b^{(q)})^{\mathsf T}+C^{(q)}\Bigr),
\]
并定义
\[
U^{(q)}:=2\bigl(C^{(q)}\bigr)^{-1},\qquad
S_{-m}(q):=\Tr\!\bigl((A^{(q)})^{-m}\bigr)\quad(m\ge1).
\]
采用延拓 Fibonacci 约定
\[
F_{-1}:=1,\qquad F_0:=0,\qquad F_1:=1.
\]
此外定义倒数 \(\Omega\)-通道
\[
\Omega_m^{(-)}(q):=\Tr\!\bigl((C^{(q)})^{-m}\bigr),
\]
则由 \(\mathrm{spec}(C^{(q)})=\{(-1)^j\varphi^{q-2j}\}_{j=0}^{q}\) 可得
\[
\Omega_m^{(-)}(q)=(-1)^{mq}\Omega_m(q),
\]
其中 \(\Omega_m(q)\) 见结论 \ref{con:xi-terminal-replica-softcore-lucas-pell-omega-unified}。

\begin{theorem}[诱导表示角元刚性与 Fibonacci 指数律]\label{thm:xi-replica-softcore-reciprocal-corner-fibonacci-power-law}
对任意整数 \(m\ge0\)，成立严格恒等式
\[
\boxed{\;
\bigl(U^{(q)}\bigr)^m_{0,0}
=2^m(-1)^{mq}F_{m-1}^{\,q}\; }.
\]
\end{theorem}
\leanverified{paper\_xi\_replica\_softcore\_reciprocal\_corner\_fibonacci\_power\_law}
\begin{proof}
令
\[
Q:=\begin{pmatrix}1&1\\[2pt]1&0\end{pmatrix}.
\]
在齐次多项式空间 \(\Sym^q(\CC^2)\) 的基
\(
\{x^{q-j}y^j\}_{j=0}^{q}
\)
下，二项展开给出
\[
\Sym^q(Q)=\bigl(C^{(q)}\bigr)^{\mathsf T},
\]
从而
\[
\bigl(C^{(q)}\bigr)^{-1}
=\bigl(\Sym^q(Q^{-1})\bigr)^{\mathsf T}.
\]
于是
\[
\bigl(U^{(q)}\bigr)^m
=2^m\Bigl(\Sym^q(Q^{-1})^m\Bigr)^{\mathsf T}
=2^m\Bigl(\Sym^q(Q^{-m})\Bigr)^{\mathsf T}.
\]
另一方面，Fibonacci 矩阵恒等式给出
\[
Q^m=
\begin{pmatrix}
F_{m+1}&F_m\\
F_m&F_{m-1}
\end{pmatrix},
\qquad
Q^{-m}=(-1)^m
\begin{pmatrix}
F_{m-1}&-F_m\\
-F_m&F_{m+1}
\end{pmatrix}.
\]
设 \(Q^{-m}=\begin{psmallmatrix}\alpha&\beta\\ \gamma&\delta\end{psmallmatrix}\)，则
\(
\bigl(\Sym^q(Q^{-m})\bigr)_{0,0}
\)
等于 \((\alpha x+\beta y)^q\) 的 \(x^q\) 系数，即 \(\alpha^q\)。
由 \(\alpha=(-1)^mF_{m-1}\)，得到
\[
\bigl(U^{(q)}\bigr)^m_{0,0}
=2^m\alpha^q
=2^m(-1)^{mq}F_{m-1}^{\,q}.
\]
证毕。
\end{proof}

\begin{theorem}[例外谱倒数矩的 Fibonacci--partition 分拆闭式]\label{thm:xi-replica-softcore-reciprocal-moment-partition-closed-form}
\leanverified{paper\_xi\_replica\_softcore\_reciprocal\_moment\_partition\_closed\_form}
固定 \(m\ge1\)。对分拆
\[
\pi=1^{m_1}2^{m_2}\cdots m^{m_m}\vdash m,
\]
记
\[
\ell(\pi):=\sum_{s=1}^{m}m_s,\qquad
c(\pi):=\frac{m}{\ell(\pi)}\cdot\frac{\ell(\pi)!}{\prod_{s=1}^{m}m_s!},
\]
\[
\Psi(\pi):=\prod_{s=1}^{m}F_{s-2}^{\,m_s}.
\]
则对任意 \(q\ge2\) 成立
\[
\boxed{\;
S_{-m}(q)
=2^m\Omega_m^{(-)}(q)
+\sum_{\substack{\pi\vdash m\\ \ell(\pi)\ge1}}
(-1)^{\ell(\pi)+q(m-\ell(\pi))}
c(\pi)\,2^{m-\ell(\pi)}\,\Psi(\pi)^q\; }.
\]
\end{theorem}
\begin{proof}
由定理 \ref{thm:xi-replica-softcore-exceptional-block-inverse-integral-closed}，
\[
\bigl(A^{(q)}\bigr)^{-1}=U^{(q)}-E,\qquad E:=E_{00}=e_0e_0^{\mathsf T},
\]
故
\[
\bigl(A^{(q)}\bigr)^{-m}=(U^{(q)}-E)^m.
\]
对该式作非交换词展开并取迹。

不含 \(E\) 的唯一词为 \(U^{(q)m}\)，贡献
\[
\Tr\!\bigl(U^{(q)m}\bigr)
=2^m\Tr\!\bigl((C^{(q)})^{-m}\bigr)
=2^m\Omega_m^{(-)}(q).
\]

设词中 \(E\) 出现 \(\ell\ge1\) 次。利用迹的循环不变性，可把每个循环类唯一写成块长
\(
(s_1,\dots,s_\ell)
\)
（\(s_r\ge1\)、\(\sum_r s_r=m\)），其意义是每块为 \(E(U^{(q)})^{s_r-1}\)。
记对应分拆为 \(\pi\vdash m\)。
由
\[
E(U^{(q)})^aE=\bigl((U^{(q)})^a\bigr)_{0,0}E
\]
和定理 \ref{thm:xi-replica-softcore-reciprocal-corner-fibonacci-power-law}，
\[
\Tr\!\bigl(E(U^{(q)})^{s_1-1}\cdots E(U^{(q)})^{s_\ell-1}\bigr)
=\prod_{r=1}^{\ell}2^{s_r-1}(-1)^{q(s_r-1)}F_{s_r-2}^{\,q}
=2^{m-\ell}(-1)^{q(m-\ell)}\Psi(\pi)^q.
\]
再乘上 \((U^{(q)}-E)^m\) 中 \(\ell\) 次选取 \((-E)\) 的符号 \((-1)^\ell\)，得到每个循环类贡献
\[
(-1)^\ell\,2^{m-\ell}(-1)^{q(m-\ell)}\Psi(\pi)^q.
\]
固定分拆 \(\pi\) 时，其循环类计数为
\[
c(\pi)=\frac{m}{\ell(\pi)}\cdot\frac{\ell(\pi)!}{\prod_s m_s!}.
\]
对全部分拆求和即得结论。
\end{proof}

\begin{theorem}[例外谱倒数幂和的环分解公式]\label{thm:xi-replica-softcore-reciprocal-moment-ring-decomposition}
\leanverified{paper\_xi\_replica\_softcore\_reciprocal\_moment\_ring\_decomposition}
固定 \(m\ge 1\)。令
\[
s_n:=(-1)^{nq}F_{n-1}^{\,q}\qquad(n\ge 0),
\]
并约定 \(F_{-1}=1\)。
则对任意 \(q\ge 2\) 成立等价闭式
\[
\boxed{\;
S_{-m}(q)
=2^m\Omega_m^{(-)}(q)
\sum_{k=1}^{m}(-1)^k\,2^{m-k}\,\frac{m}{k}
\sum_{\substack{n_1+\cdots+n_k=m-k\\ n_j\ge 0}}
\ \prod_{j=1}^{k}s_{n_j}\; }.
\]
特别地，\(S_{-m}(q)\in\ZZ\) 对一切 \(m\ge 1\)、\(q\ge 2\) 成立。
\end{theorem}
\begin{proof}
沿用定理 \ref{thm:xi-replica-softcore-reciprocal-moment-partition-closed-form} 证明中的展开
\(
(A^{(q)})^{-m}=(U^{(q)}-E)^m
\)
与循环类归并，但不再按分拆 \(\pi\) 进一步合并。

\par
设在某一项中 \(E\) 出现 \(k\ge 1\) 次。
由迹的循环不变性，可把该循环类唯一编码为一个 \(k\)-元非负整数组合 \((n_1,\dots,n_k)\) 满足 \(\sum_{j=1}^{k}n_j=m-k\)，其意义为相邻两次 \(E\) 之间出现 \(U^{(q)}\) 的次数。
每个循环类在长度为 \(m\) 的线性展开中出现次数为 \(m/k\)。
又由
\[
E\bigl(U^{(q)}\bigr)^{n}E=\bigl(U^{(q)}\bigr)^{n}_{0,0}\,E
=2^n(-1)^{nq}F_{n-1}^{\,q}\,E
\]
（定理 \ref{thm:xi-replica-softcore-reciprocal-corner-fibonacci-power-law}，并用延拓约定使 \(n=0\) 时一致），得到该循环类对迹的贡献为
\[
(-1)^k\,2^{m-k}\prod_{j=1}^{k}\bigl((-1)^{n_jq}F_{n_j-1}^{\,q}\bigr),
\]
其中 \((-1)^k\) 来自选取 \((-E)\) 的符号。
对所有 \(k\) 与 \((n_1,\dots,n_k)\) 求和并加上无 \(E\) 的项 \(2^m\Omega_m^{(-)}(q)\)，即得所述公式。
整数性由 \( (A^{(q)})^{-1}\in\ZZ^{(q+1)\times(q+1)} \) 推出：
\(
S_{-m}(q)=\Tr(((A^{(q)})^{-1})^m)\in\ZZ
\)。
\end{proof}

\begin{corollary}[倒数三至六次矩的显式闭式]\label{cor:xi-replica-softcore-reciprocal-moment-m3-m6-closed}
\leanverified{paper\_xi\_replica\_softcore\_reciprocal\_moment\_m3\_m6\_closed}
对任意 \(q\ge2\)，成立
\[
\boxed{\;
S_{-3}(q)=8\,\Omega_3^{(-)}(q)-13\;},
\]
\[
\boxed{\;
S_{-4}(q)=16\,\Omega_4^{(-)}(q)-32(-1)^q+17\;},
\]
\[
\boxed{\;
S_{-5}(q)=32\,\Omega_5^{(-)}(q)-80\cdot 2^q+40(-1)^q-21\;},
\]
\[
\boxed{\;
S_{-6}(q)=64\,\Omega_6^{(-)}(q)+96\cdot 2^q-192(-1)^q3^q-48(-1)^q+73\; }.
\]
\end{corollary}
\begin{proof}
在定理 \ref{thm:xi-replica-softcore-reciprocal-moment-partition-closed-form} 中分别取 \(m=3,4,5,6\)。
含部件 \(2\) 的分拆项因 \(F_0=0\) 全部消失；其余有限项直接化简即得。
\end{proof}

\begin{corollary}[倒数一次至八次矩的完全闭式]\label{cor:xi-replica-softcore-reciprocal-moment-m1-m8-closed}
对任意 \(q\ge 2\)，成立
\[
\boxed{\;
S_{-1}(q)=2\,\Omega_1^{(-)}(q)-1=2(-1)^qF_{q+1}-1\;},
\]
\[
\boxed{\;
S_{-2}(q)=4\,\Omega_2^{(-)}(q)+1=4F_{2q+2}+1\;},
\]
以及推论 \ref{cor:xi-replica-softcore-reciprocal-moment-m3-m6-closed} 的 \(m=3,4,5,6\) 公式；并且
\[
\boxed{\;
S_{-7}(q)=128\,\Omega_7^{(-)}(q)-448\cdot 5^{q}+224(-1)^q 3^{q}-112\cdot 2^{q}+280(-1)^q-141\;},
\]
\[
\boxed{\;
S_{-8}(q)=256\,\Omega_8^{(-)}(q)-1024(-1)^q 8^{q}+512\cdot 5^{q}-256(-1)^q 3^{q}+640\cdot 2^{q}-576(-1)^q+481\; }.
\]
\end{corollary}
\begin{proof}
在定理 \ref{thm:xi-replica-softcore-reciprocal-moment-partition-closed-form} 中分别取 \(m=1,2,7,8\)。
当 \(m=1\) 时仅保留 \(\ell(\pi)=1\) 的常数项，得到 \(S_{-1}(q)=2\Omega_1^{(-)}(q)-1\)；
当 \(m=2\) 时仅出现 \(\Psi(\pi)\in\{F_{-1},F_0\}\) 且 \(F_0=0\) 使非平凡分拆项消失，得到 \(S_{-2}(q)=4\Omega_2^{(-)}(q)+1\)。
当 \(m=7,8\) 时同理由有限分拆项直接化简得到所示闭式。
\end{proof}

\begin{theorem}[全矩阵逆幂迹的 Lucas 主项与有限修正]\label{thm:xi-replica-softcore-full-trace-negative-power-lucas-finite-correction}
令 \(q\ge 2\)，并记 Lucas 数列 \(L_0=2\)、\(L_1=1\)、\(L_{m+1}=L_m+L_{m-1}\)。
则对任意整数 \(m\ge 1\) 成立
\[
\boxed{\;
\Tr\!\bigl((T^{(q)})^{-m}\bigr)
=2^m(-1)^{mq}L_m^{\,q}
\sum_{k=1}^{m}(-1)^k\,2^{m-k}\,\frac{m}{k}
\sum_{\substack{n_1+\cdots+n_k=m-k\\ n_j\ge 0}}
\ \prod_{j=1}^{k}\Bigl((-1)^{n_jq}F_{n_j-1}^{\,q}\Bigr)\; }.
\]
从而对 \(m=1,\dots,8\)（其中 \(L_1=1,L_2=3,L_3=4,L_4=7,L_5=11,L_6=18,L_7=29,L_8=47\)）有完全闭式
\[
\begin{aligned}
\Tr\!\bigl((T^{(q)})^{-1}\bigr)&=2(-1)^q-1,\\
\Tr\!\bigl((T^{(q)})^{-2}\bigr)&=4\cdot 3^q+1,\\
\Tr\!\bigl((T^{(q)})^{-3}\bigr)&=8(-1)^q\cdot 4^q-13,\\
\Tr\!\bigl((T^{(q)})^{-4}\bigr)&=16\cdot 7^q+17-32(-1)^q,\\
\Tr\!\bigl((T^{(q)})^{-5}\bigr)&=32(-1)^q\cdot 11^q-80\cdot 2^q-21+40(-1)^q,\\
\Tr\!\bigl((T^{(q)})^{-6}\bigr)&=64\cdot 18^q+96\cdot 2^q-192(-1)^q 3^q+73-48(-1)^q,\\
\Tr\!\bigl((T^{(q)})^{-7}\bigr)&=128(-1)^q\cdot 29^q-448\cdot 5^q+224(-1)^q 3^q-112\cdot 2^q+280(-1)^q-141,\\
\Tr\!\bigl((T^{(q)})^{-8}\bigr)&=256\cdot 47^q-1024(-1)^q 8^q+512\cdot 5^q-256(-1)^q 3^q+640\cdot 2^q-576(-1)^q+481.
\end{aligned}
\]
\end{theorem}
\begin{proof}
由定理 \ref{thm:xi-replica-softcore-full-spectrum-direct-sum} 的例外--体谱直和分解，
\[
\Tr\!\bigl((T^{(q)})^{-m}\bigr)
=S_{-m}(q)+\sum_{i=0}^{q}\Bigl(\binom{q}{i}-1\Bigr)\Bigl(\frac{2}{\mu_i^{(q)}}\Bigr)^{m},
\qquad
\mu_i^{(q)}=\varphi^{\,q-i}\psi^{\,i}.
\]
体谱部分满足
\[
\sum_{i=0}^{q}\binom{q}{i}\bigl(\mu_i^{(q)}\bigr)^{-m}
=(\varphi^{-m}+\psi^{-m})^q
=\bigl((-1)^mL_m\bigr)^q
=(-1)^{mq}L_m^{\,q},
\]
而
\(
\sum_{i=0}^{q}\bigl(\mu_i^{(q)}\bigr)^{-m}
=\Tr\!\bigl((C^{(q)})^{-m}\bigr)=\Omega_m^{(-)}(q)
\)
（见本段定义）。
因此体谱贡献为
\(
2^m\bigl((-1)^{mq}L_m^{\,q}-\Omega_m^{(-)}(q)\bigr)
\)。
将其与定理 \ref{thm:xi-replica-softcore-reciprocal-moment-ring-decomposition} 中的 \(S_{-m}(q)\) 相加，\(\Omega_m^{(-)}(q)\) 项严格相消，得到一般公式；将 \(m=1,\dots,8\) 的有限和逐一化简即可得到列出的闭式。证毕。
\end{proof}

\begin{theorem}[固定阶倒数矩的 \(q\)-方向递推因子]\label{thm:xi-replica-softcore-reciprocal-moment-q-recurrence}
\leanverified{paper\_xi\_replica\_softcore\_reciprocal\_moment\_q\_recurrence}
固定 \(m\ge1\)。定义
\[
\mathcal A_m^{-}:=
\Bigl\{
(-1)^{m-\ell(\pi)}\Psi(\pi)\,:\,
\pi\vdash m,\ \ell(\pi)\ge1,\ \Psi(\pi)\neq0
\Bigr\}\subset\ZZ.
\]
记 Lucas 数 \(L_m=\varphi^m+(-\varphi^{-1})^m\)，并置
\[
\widetilde L_m:=(-1)^mL_m.
\]
则序列 \(q\mapsto S_{-m}(q)\) 满足常系数线性递推，其特征多项式可取为
\[
\boxed{\;
P_m^{-}(X)=\bigl(X^2-\widetilde L_m X+(-1)^m\bigr)
\prod_{\alpha\in\mathcal A_m^{-}}(X-\alpha)\; }.
\]
将重复根合并并去除零系数通道后，得到该序列的最小特征多项式。
\end{theorem}
\begin{proof}
由定理 \ref{thm:xi-replica-softcore-reciprocal-moment-partition-closed-form}，
\[
S_{-m}(q)=2^m\Omega_m^{(-)}(q)+\sum_{\alpha\in\mathcal A_m^{-}}B_\alpha\,\alpha^q,
\]
其中 \(B_\alpha\in\QQ\) 与 \(q\) 无关。
又
\[
\Omega_m^{(-)}(q)
=u_m(\varphi^{-m})^q+v_m\bigl(((-1)^m\varphi^{m})\bigr)^q
\]
（\(u_m,v_m\neq0\)），故其特征二次因子为
\[
(X-\varphi^{-m})(X-(-1)^m\varphi^m)
=X^2-\widetilde L_m X+(-1)^m.
\]
每个指数项 \(\alpha^q\) 被一次因子 \(X-\alpha\) 湮灭，
因此 \(P_m^{-}\) 湮灭整个序列。合并重复根并删除零系数通道后即得最小特征多项式。
\end{proof}

\begin{theorem}[逆二项矩阵负幂的逐项卷积闭式]\label{thm:xi-replica-softcore-binomial-inverse-negative-power-entrywise}
对任意 \(q\ge2\)、\(m\ge0\) 与 \(0\le i,j\le q\)，成立
\[
\boxed{
\begin{aligned}
\bigl(C^{(q)}\bigr)^{-m}_{i,j}
=(-1)^{mq+i+j}
\sum_{t=\max(0,i+j-q)}^{\min(i,j)}
&\binom{i}{t}\binom{q-i}{j-t}\\
&\cdot F_{m-1}^{\,q-i-j+t}
F_m^{\,i+j-2t}
F_{m+1}^{\,t}.
\end{aligned}
}
\]
等价地，
\[
\bigl(U^{(q)}\bigr)^m_{i,j}=2^m\bigl(C^{(q)}\bigr)^{-m}_{i,j}.
\]
\end{theorem}
\leanverified{paper\_xi\_replica\_softcore\_binomial\_inverse\_negative\_power\_entrywise}
\begin{proof}
仍记
\(
Q=\begin{psmallmatrix}1&1\\1&0\end{psmallmatrix}
\)。
由 \(\Sym^q(Q)=\bigl(C^{(q)}\bigr)^{\mathsf T}\)，得
\[
\bigl(C^{(q)}\bigr)^{-m}
=\Bigl(\Sym^q(Q^{-m})\Bigr)^{\mathsf T}.
\]
设
\[
Q^{-m}=
\begin{pmatrix}
a&b\\ c&d
\end{pmatrix},
\]
则 \(\bigl(\Sym^q(Q^{-m})\bigr)_{j,i}\) 等于
\(
(ax+by)^{q-i}(cx+dy)^i
\)
中 \(x^{q-j}y^j\) 的系数，故
\[
\bigl(C^{(q)}\bigr)^{-m}_{i,j}
=\sum_{t}
\binom{i}{t}\binom{q-i}{j-t}
a^{\,q-i-(j-t)}b^{\,j-t}c^{\,i-t}d^{\,t},
\]
其中求和区间由指数非负性给出
\(
t\in[\max(0,i+j-q),\min(i,j)].
\)
代入
\[
Q^{-m}=(-1)^m
\begin{pmatrix}
F_{m-1}&-F_m\\ -F_m&F_{m+1}
\end{pmatrix},
\]
并整理符号指数，即得
\(
(-1)^{mq+i+j}
\)
与所示三参数卷积因子。证毕。
\end{proof}

\endinput


\subsubsection{跨支隧穿作用量与分裂纤维的首项可见性：\texorpdfstring{$\ZZ_2$}{Z/2} 扭曲 Ihara-\texorpdfstring{$\zeta$}{zeta}、谱双重态裂分与端点角窗}
沿用定义 \ref{def:ihara-tunneling-action} 对隧穿作用量 \(A_K\) 的约定（无回溯 Ihara prime 环的最小正能量），并沿用引理 \ref{lem:pom-one-bit} 的隐藏位二支分解
\(
d_m(x)=d_{m,0}(x)+d_{m,1}(x)
\)；
同时沿用结论 \ref{con:xi-timefiber-poisson-busemann-identity} 的端点 \(-1\) Poisson--Busemann 时间纤维坐标
\(
B^{-1}(w)=\frac{1-|w|^2}{|1+w|^2}
\)。
在该固定口径下，隐藏位二支分裂在协议图/无回溯 prime 侧诱导一个 \(\ZZ_2\) 扭曲通道；
其首个可见阶数由一纯整数不变量 \(A_{\mathrm{flip}}\) 完全刻画，并可在 Euler 乘积比值、Perron 双重态裂分与端点读出分辨率三条接口上同时读出。

\begin{conclusion}[分裂隧穿作用量的 \texorpdfstring{$\zeta$}{zeta}-检测：扭曲/未扭曲比值的 \texorpdfstring{$u$}{u}-首阶指数]\label{con:xi-split-tunneling-zeta-detect-aflip}
设 \(\Gamma\) 为有限有向图，并在其 Hashimoto 无回溯边空间上取 Ihara prime 环（tailless、primitive、non-backtracking 闭环）为 \(p\)。
令 \(\kappa:E(\Gamma)\to\ZZ_{\ge 0}\) 为能量权重，\(\ell:E(\Gamma)\to\ZZ_{>0}\) 为时间权重，并对闭环取总权重 \(\kappa(p),\ell(p)\)。
设 \(\beta:E(\Gamma)\to\ZZ_2\) 为由隐藏位二支分裂诱导的 \(\ZZ_2\)-余圈，并记
\[
\chi_\beta(p):=(-1)^{\sum_{e\in p}\beta(e)}\in\{\pm 1\}.
\]
定义二变量带权 Ihara \(\zeta\) 与其 \(\ZZ_2\)-扭曲通道（形式幂级数意义）
\[
Z_K(u,t):=\prod_{p}\bigl(1-u^{\kappa(p)}t^{\ell(p)}\bigr)^{-1},
\qquad
Z_K(u,t;\beta):=\prod_{p}\bigl(1-\chi_\beta(p)\,u^{\kappa(p)}t^{\ell(p)}\bigr)^{-1}.
\]
若零能量骨架不发生 sheet-交换（即对所有 \(p\) 有 \(\kappa(p)=0\Rightarrow \chi_\beta(p)=+1\)），则可定义扩展取值的不变量
\[
A_{\mathrm{flip}}
:=\min\{\kappa(p):\ p\ \text{为 Ihara prime 环且}\ \chi_\beta(p)=-1\}\in\ZZ_{>0}\cup\{+\infty\},
\]
其中约定 \(\min\varnothing:=+\infty\)。
当 \(A_{\mathrm{flip}}=+\infty\) 时，所有 prime 环均满足 \(\chi_\beta(p)=+1\)，从而
\[
Z_K(u,t;\beta)\equiv Z_K(u,t),
\qquad
\frac{Z_K(u,t;\beta)}{Z_K(u,t)}\equiv 1.
\]
一般地，比值的 \(u\)-首阶指数由 \(A_{\mathrm{flip}}\) 给出：
\[
[u^k]\left(\frac{Z_K(u,t;\beta)}{Z_K(u,t)}-1\right)=0\quad(0\le k<A_{\mathrm{flip}}),
\]
且当 \(A_{\mathrm{flip}}<+\infty\) 时，其首项系数具有闭式
\[
\frac{Z_K(u,t;\beta)}{Z_K(u,t)}
=1-2\Bigl(\sum_{\substack{p\ \mathrm{Ihara\ prime}\\ \kappa(p)=A_{\mathrm{flip}},\ \chi_\beta(p)=-1}} t^{\ell(p)}\Bigr)u^{A_{\mathrm{flip}}}
+O\bigl(u^{A_{\mathrm{flip}}+1}\bigr).
\]
因此，隐藏位二支分裂对“零温隧穿”的首个可见阶数并不由 \(A_K\) 决定，而由跨支隧穿作用量 \(A_{\mathrm{flip}}\) 完全刻画。
\end{conclusion}
\begin{proof}[证明要点]
把 \(x_p:=u^{\kappa(p)}t^{\ell(p)}\) 代入，则逐 prime 环有严格因子化
\[
\frac{Z_K(u,t;\beta)}{Z_K(u,t)}=\prod_{p}\frac{1-x_p}{1-\chi_\beta(p)x_p}.
\]
若对所有 prime 环均有 \(\chi_\beta(p)=+1\)，则乘积中每一因子恒为 \(1\)，从而比值恒为 \(1\)，并按约定 \(A_{\mathrm{flip}}=\min\varnothing=+\infty\)。
下设存在某条 prime 环满足 \(\chi_\beta(p)=-1\)。
当 \(\chi_\beta(p)=+1\) 时该因子恒为 \(1\)；当 \(\chi_\beta(p)=-1\) 时
\(
\frac{1-x_p}{1+x_p}=1-2x_p+O(x_p^2)
\)。
零能量假设排除了 \(\kappa(p)=0\) 且 \(\chi_\beta(p)=-1\) 的情形，从而最小能量阶由
\(
\min\{\kappa(p):\chi_\beta(p)=-1\}=A_{\mathrm{flip}}
\)
给出，首项系数由所有满足 \(\kappa(p)=A_{\mathrm{flip}}\) 的 prime 环线性叠加得到。
\end{proof}

\begin{conclusion}[跨支轨道的奇次重复选择律：扭曲比值的 Euler 乘积仅含奇次幂]\label{con:xi-split-tunneling-odd-power-selection-law}
在结论 \ref{con:xi-split-tunneling-zeta-detect-aflip} 的设定下，将 Ihara prime 环按 \(\chi_\beta(p)=\pm 1\) 分割，则有恒等式
\[
\frac{Z_K(u,t;\beta)}{Z_K(u,t)}
=
\prod_{\substack{p\ \mathrm{Ihara\ prime}\\ \chi_\beta(p)=-1}}
\frac{1-u^{\kappa(p)}t^{\ell(p)}}{1+u^{\kappa(p)}t^{\ell(p)}}.
\]
从而其对数展开满足严格的奇次选择律（形式幂级数意义）
\[
\log\frac{Z_K(u,t;\beta)}{Z_K(u,t)}
=
-2\sum_{\substack{p\ \mathrm{Ihara\ prime}\\ \chi_\beta(p)=-1}}
\sum_{j\ge 0}\frac{1}{2j+1}\Bigl(u^{\kappa(p)}t^{\ell(p)}\Bigr)^{2j+1},
\]
即每一条 \(\chi_\beta(p)=-1\) 的 Ihara prime 环只以奇次重复 \((2j+1)\) 进入 \(\log(Z_K(u,t;\beta)/Z_K(u,t))\)，所有偶次重复在比值中完全消隐。
\end{conclusion}
\begin{proof}[证明要点]
结论 \ref{con:xi-split-tunneling-zeta-detect-aflip} 的证明已给出逐 prime 环因子化
\[
\frac{Z_K(u,t;\beta)}{Z_K(u,t)}=\prod_{p}\frac{1-x_p}{1-\chi_\beta(p)x_p},\qquad x_p:=u^{\kappa(p)}t^{\ell(p)}.
\]
当 \(\chi_\beta(p)=+1\) 时因子恒为 \(1\)；当 \(\chi_\beta(p)=-1\) 时因子为 \((1-x_p)/(1+x_p)\)，从而得到乘积式。
对各因子取对数并用恒等式
\(
\log\frac{1-x}{1+x}=-2\sum_{j\ge 0}\frac{x^{2j+1}}{2j+1}
\)
（\(|x|<1\)，或按形式幂级数理解）即得对数展开。
\end{proof}

\begin{conclusion}[\texorpdfstring{$\ZZ_2$}{Z/2} 表示分解的行列式比：扭曲比值等价于两套 Fredholm 行列式之商]\label{con:xi-split-tunneling-z2-determinant-ratio}
设由 \(\beta\) 定义的两层覆盖之提升转移算子 \(\widetilde M(u,t)\) 与覆盖交换对称 \(J\) 可对易，并在 \(J\) 的 \(\pm1\) 特征子空间上诱导两套同维算子 \(M^{\pm}(u,t)\)。
则存在（由覆盖理论与字符分解固定的）同一性
\[
Z_K(u,t)=\det\bigl(I-M^{+}(u,t)\bigr)^{-1},
\qquad
Z_K(u,t;\beta)=\det\bigl(I-M^{-}(u,t)\bigr)^{-1},
\]
从而
\[
\frac{Z_K(u,t;\beta)}{Z_K(u,t)}
=
\frac{\det\bigl(I-M^{+}(u,t)\bigr)}{\det\bigl(I-M^{-}(u,t)\bigr)}.
\]
进一步，若在 sheet 基下 \(\widetilde M(u,t)\) 可写成对称块矩阵
\[
\widetilde M(u,t)=
\begin{pmatrix}
A(u,t) & B(u,t)\\
B(u,t) & A(u,t)
\end{pmatrix},
\]
则 \(M^{+}=A+B\)、\(M^{-}=A-B\)；
并且令
\[
S(u,t):=\bigl(I-A(u,t)\bigr)^{-1}B(u,t),
\]
则有等式
\[
\frac{Z_K(u,t;\beta)}{Z_K(u,t)}
=
\frac{\det\bigl(I-S(u,t)\bigr)}{\det\bigl(I+S(u,t)\bigr)}.
\]
该表示把跨支效应压缩为一对同阶 Fredholm 行列式之比；其对数展开仅含 \(S\) 的奇次幂，且 \(S(u,t)^2\) 对应单 sheet 上的 \(0\to1\to0\) 二步隧穿复合。
\end{conclusion}
\begin{proof}[证明要点]
第一段为 \(J\)-等变分解下的字符分量同一性：平凡表示与符号表示分别对应未扭曲与 \(\ZZ_2\)-扭曲通道。
第二段在基变换 \((v,v)\)、\((v,-v)\) 下把对称块矩阵对角化得到 \(A\pm B\)。
最后一式由恒等式
\[
I-(A\pm B)=\bigl(I-A\bigr)\bigl(I\mp (I-A)^{-1}B\bigr)
\]
取行列式并相除即得。
\end{proof}

\begin{conclusion}[分裂余圈的置换论刻画：符号角色与取向单值化类]\label{con:xi-split-beta-sign-orientation-w1}
设一套有限状态实现使每条边 \(e\in E(\Gamma)\) 在某一固定的隐藏态纤维（大小 \(n\ge 2\)）上诱导一个重标号置换 \(\sigma_e\in \mathfrak S_n\)；
从而对任意闭环 \(p\)（按边序列复合）可定义其置换 holonomy
\(
\sigma(p):=\prod_{e\in p}\sigma_e\in\mathfrak S_n
\)。
定义 \(\ZZ_2\)-余圈
\[
\beta_{\sigma}:E(\Gamma)\to\ZZ_2,\qquad
\beta_{\sigma}(e)\equiv \operatorname{sgn}(\sigma_e)\ (\mathrm{mod}\ 2),
\]
并记 \(\chi_{\beta_\sigma}(p):=(-1)^{\sum_{e\in p}\beta_\sigma(e)}\)。
则对任意闭环 \(p\) 有严格恒等式
\[
\chi_{\beta_\sigma}(p)=\operatorname{sgn}\!\bigl(\sigma(p)\bigr).
\]
因此 \([\beta_\sigma]\in H^1(\Gamma;\ZZ_2)\) 正是单值化置换表示经由符号角色诱导的取向类（等价地：与 \(\mathrm{Or}([n])\) 相关联的取向 \(\ZZ_2\)-局部系统的第一 Stiefel--Whitney 类）；
特别地，\([\beta_\sigma]=0\) 当且仅当 \(\sigma(p)\in\mathfrak A_n\) 对每条闭环 \(p\) 成立（即 monodromy 可约化到交错群）。
\par
在本节的 \(\ZZ_2\) 扭曲 Ihara-\(\zeta\) 口径下，当 \([\beta_\sigma]=0\) 时必有 \(A_{\mathrm{flip}}=+\infty\)；
而当 \([\beta_\sigma]\neq 0\) 时，\(A_{\mathrm{flip}}\) 等于 \([\beta_\sigma]\) 的 \(\ZZ_2\)-奇闭链能量极小值（结论 \ref{con:xi-split-tunneling-aflip-homological-minimization}）。
\end{conclusion}
\begin{proof}[证明要点]
由 \(\operatorname{sgn}:\mathfrak S_n\to\{\pm1\}\) 为群同态，
\[
\operatorname{sgn}\!\bigl(\sigma(p)\bigr)
=\prod_{e\in p}\operatorname{sgn}(\sigma_e)
=(-1)^{\sum_{e\in p}(\operatorname{sgn}(\sigma_e)\bmod 2)}
=(-1)^{\sum_{e\in p}\beta_\sigma(e)}
=\chi_{\beta_\sigma}(p),
\]
得到第一式。
其余断言为上述恒等式在同调层面的翻译：\([\beta_\sigma]\) 即为 \(\operatorname{sgn}\circ\sigma\) 的同调类，故 \([\beta_\sigma]=0\) 等价于符号单值化平凡，亦即闭环 holonomy 全落入 \(\ker(\operatorname{sgn})=\mathfrak A_n\)。
最后两句分别由 \(A_{\mathrm{flip}}\) 的定义与结论 \ref{con:xi-split-tunneling-aflip-homological-minimization} 直接给出。证毕。
\end{proof}

\paragraph{有限状态骨架与闭环探针：从机制可解释性到可审计谱--同调指纹}
机制可解释性研究中，已有工作把 Transformer 的某些状态跟踪与推理行为抽取为有限状态骨架，从而把“电路”提升为有限有向图上的动力系统与加权算子并允许离线核验与可重复审计 \cite{SomvanshiIslamRafeTustiChakrabortyBaitullahChowdhuryAlnawmasiDuttaDas2026BridgingBlackBoxSurveyMechInterp,zhang-etal-2025-finite}。
在本文的协议接口中，这一抽取对应于“有限状态实现 \(\Rightarrow\) 行列式 \(\zeta\) 指纹 \(\Rightarrow\) 有限部/谱指纹 \(\Rightarrow\) 同调不变量”的闭环；
其中 \(\ZZ_2\)-余圈扭曲及其跨支激活指数 \(A_{\mathrm{flip}}\)（结论 \ref{con:xi-split-tunneling-zeta-detect-aflip}）给出一个可推广的“分支耦合首阶能量”原型。
下述结论把该原型从 \(\ZZ_2\) 推进到一般有限群的表示分解，并给出把有限签名进一步压缩为有限条迹数据的完备化准则。

\begin{conclusion}[有限群扭曲 Ihara-\texorpdfstring{$\zeta$}{zeta} 的首激活指数：表示通道的低能不可见性]\label{con:xi-ihara-finite-group-twisted-activation-exponent}
设 \(\Gamma\) 为有限有向图，并取其 Hashimoto 无回溯边空间上的 Ihara prime 环（tailless、primitive、non-backtracking 闭环）为 \(p\)。
令 \(\kappa:E(\Gamma)\to\ZZ_{\ge 0}\) 为能量权重，\(\ell:E(\Gamma)\to\ZZ_{>0}\) 为时间权重，并对闭环记总权重 \(\kappa(p),\ell(p)\)。

令 \(H\) 为有限群，并给定一个边上的 holonomy 余圈 \(\beta:E(\Gamma)\to H\)；对闭环定义
\[
\beta(p):=\prod_{e\in p}\beta(e)\in H.
\]
令 \(\rho:H\to \GL(V_\rho)\) 为复表示，记 \(d_\rho:=\dim V_\rho\)。
定义表示扭曲的二变量 Ihara-\(\zeta\)（形式幂级数意义）
\[
Z_K(u,t;\beta,\rho):=\prod_{p}\det\!\Bigl(I-u^{\kappa(p)}t^{\ell(p)}\,\rho(\beta(p))\Bigr)^{-1},
\qquad
Z_K(u,t):=\prod_{p}\bigl(1-u^{\kappa(p)}t^{\ell(p)}\bigr)^{-1}.
\]
假设零能量骨架无 holonomy：
\[
\forall p,\ \kappa(p)=0\ \Longrightarrow\ \beta(p)=e_H.
\]
定义 \(\rho\)-激活指数（约定 \(\min\varnothing:=+\infty\)）
\[
A(\rho):=\min\bigl\{\kappa(p):\ p\ \text{为 Ihara prime 环且}\ \rho(\beta(p))\neq I\bigr\}\in\ZZ_{>0}\cup\{+\infty\}.
\]
则：
\begin{enumerate}
  \item 若 \(A(\rho)=+\infty\)，则 \(Z_K(u,t;\beta,\rho)\equiv Z_K(u,t)^{d_\rho}\)，从而比值
  \[
  \frac{Z_K(u,t;\beta,\rho)}{Z_K(u,t)^{d_\rho}}\equiv 1.
  \]
  \item 若 \(A(\rho)<+\infty\)，则存在 \(F_\rho(t)\in\CC[[t]]\) 使得
  \[
  \frac{Z_K(u,t;\beta,\rho)}{Z_K(u,t)^{d_\rho}}
  =1+u^{A(\rho)}F_\rho(t)+O\!\bigl(u^{A(\rho)+1}\bigr),
  \]
  且可取首项系数的显式表达
  \[
  F_\rho(t)=\sum_{\substack{p\ \mathrm{Ihara\ prime}\\ \kappa(p)=A(\rho)}}\bigl(\Tr\rho(\beta(p))-d_\rho\bigr)\,t^{\ell(p)}.
  \]
\end{enumerate}
\end{conclusion}
\begin{proof}[证明要点]
对任意矩阵 \(M\) 与形式变量 \(x\) 有恒等式
\[
\det(I-xM)^{-1}
=\exp\!\Bigl(\sum_{m\ge 1}\frac{x^m}{m}\Tr(M^m)\Bigr).
\]
对每条 prime 环 \(p\) 令 \(x_p:=u^{\kappa(p)}t^{\ell(p)}\)。
则
\[
\log Z_K(u,t;\beta,\rho)
=\sum_{p}\sum_{m\ge 1}\frac{x_p^m}{m}\Tr\!\bigl(\rho(\beta(p))^m\bigr),
\qquad
\log Z_K(u,t)^{d_\rho}
=\sum_{p}\sum_{m\ge 1}\frac{x_p^m}{m}\,d_\rho.
\]
相减得到
\[
\log\frac{Z_K(u,t;\beta,\rho)}{Z_K(u,t)^{d_\rho}}
=\sum_{p}\sum_{m\ge 1}\frac{u^{m\kappa(p)}t^{m\ell(p)}}{m}\Bigl(\Tr(\rho(\beta(p))^m)-d_\rho\Bigr).
\]
由假设，\(\kappa(p)=0\Rightarrow \beta(p)=e_H\Rightarrow \rho(\beta(p))^m=I\)，故零能量项严格消失。
对任意满足 \(\rho(\beta(p))\neq I\) 的 prime 环，最小的 \(u\)-阶由 \(m=1\) 给出；
因此对数比值的首个非零 \(u\)-阶为 \(A(\rho)\)，其系数为 \(\sum_{\kappa(p)=A(\rho)}(\Tr\rho(\beta(p))-d_\rho)t^{\ell(p)}\)。
最后用 \(\exp(u^{A}g+O(u^{A+1}))=1+u^{A}g+O(u^{A+1})\)（\(A\ge 1\)）得到所示展开。证毕。
\end{proof}

\begin{conclusion}[表示分层的超选结构：低能展开在临界阶前对所有非平凡通道不可见]\label{con:xi-ihara-finite-group-superselection-before-amin}
沿用结论 \ref{con:xi-ihara-finite-group-twisted-activation-exponent} 的设定。
定义
\[
A_{\min}:=\min\bigl\{A(\rho):\rho\in\widehat H,\ \rho\ \text{为非平凡不可约表示}\bigr\}\in\ZZ_{>0}\cup\{+\infty\}.
\]
则对任意整数 \(k<A_{\min}\) 与任意不可约表示 \(\rho\in\widehat H\) 都有
\[
[u^k]\left(\frac{Z_K(u,t;\beta,\rho)}{Z_K(u,t)^{d_\rho}}-1\right)=0\qquad\text{于}\ \CC[[t]].
\]
换言之，所有非平凡表象通道在 \(u^{A_{\min}}\) 之前严格不可见：低能展开在阶 \(<A_{\min}\) 上形成表示分层的超选结构。
\end{conclusion}
\begin{proof}[证明要点]
由结论 \ref{con:xi-ihara-finite-group-twisted-activation-exponent}，任意 \(\rho\) 的首个非零 \(u\)-阶为 \(A(\rho)\)。
当 \(k<A_{\min}\) 时对所有非平凡不可约 \(\rho\) 均有 \(k<A(\rho)\)，故对应系数恒为零。证毕。
\end{proof}

\begin{remark}[\texorpdfstring{$\ZZ_2$}{Z/2} 情形的退化]
取 \(H=\ZZ_2\) 且 \(\rho\) 为其非平凡字符，则 \(d_\rho=1\)，并且
\[
\frac{Z_K(u,t;\beta,\rho)}{Z_K(u,t)^{d_\rho}}
=\frac{Z_K(u,t;\beta)}{Z_K(u,t)}.
\]
此时 \(A(\rho)\) 退化为结论 \ref{con:xi-split-tunneling-zeta-detect-aflip} 的 \(A_{\mathrm{flip}}\)，而上述首激活展开与奇次选择律分别对应结论 \ref{con:xi-split-tunneling-zeta-detect-aflip} 与 \ref{con:xi-split-tunneling-odd-power-selection-law}。
\end{remark}

\begin{proposition}[阿贝尔 holonomy 的严格反演：角色族确定 holonomy 分解族]\label{prop:xi-abelian-holonomy-fourier-inversion}
\leanverified{paper\_xi\_abelian\_holonomy\_fourier\_inversion}
在结论 \ref{con:xi-ihara-finite-group-twisted-activation-exponent} 的设定下，进一步假设 \(H\) 为有限阿贝尔群。
取一类闭走廊集合 \(\mathcal{W}\)（例如无回溯闭走廊或 Hashimoto 边空间上的闭游走），并对闭走廊 \(w\) 定义总权重 \(\kappa(w),\ell(w)\) 与 holonomy \(\beta(w)\)（沿边乘积）。
对任意字符 \(\chi\in\widehat H\) 定义扭曲生成函数
\[
W_\chi(u,t):=\sum_{w\in\mathcal{W}}u^{\kappa(w)}t^{\ell(w)}\chi(\beta(w)).
\]
对任意 \(g\in H\) 定义 holonomy 分解生成函数
\[
W_g(u,t):=\sum_{\substack{w\in\mathcal{W}\\ \beta(w)=g}}u^{\kappa(w)}t^{\ell(w)}.
\]
则成立严格 Fourier 反演公式
\[
\boxed{\;
W_g(u,t)=\frac{1}{|H|}\sum_{\chi\in\widehat H}\chi(g)^{-1}W_\chi(u,t)\qquad(g\in H).\;}
\]
因此字符族 \(\{W_\chi\}_{\chi\in\widehat H}\) 在系数层面完整决定 holonomy 分解族 \(\{W_g\}_{g\in H}\)。
\end{proposition}
\begin{proof}
按 holonomy 分解，
\(
W_\chi(u,t)=\sum_{g\in H}\chi(g)W_g(u,t)
\)。
对有限阿贝尔群，字符满足正交关系
\(
\frac1{|H|}\sum_{\chi\in\widehat H}\chi(g)\chi(h)^{-1}=\mathbf 1_{\{g=h\}}
\)，
代入即得反演公式。证毕。
\end{proof}

\begin{conclusion}[迹数据的 Newton 完备性：行列式 \texorpdfstring{$\zeta$}{zeta} 指纹由有限条闭环迹确定]\label{con:xi-newton-trace-completeness-determinant-zeta}
设 \(R\) 为含 \(\QQ\) 的交换环，\(B\in M_N(R)\)。
令
\[
p_n:=\Tr(B^n)\in R\qquad (1\le n\le N).
\]
则首一特征多项式 \(\chi_B(\lambda)=\det(\lambda I-B)\in R[\lambda]\) 的系数由 \(\{p_n\}_{n=1}^{N}\) 唯一确定。
等价地，
\[
\det(I-tB)\in R[t]
\]
由 \(\{p_n\}_{n=1}^{N}\) 唯一确定，从而矩阵 \(\zeta\) 指纹
\[
\zeta_B(t):=\det(I-tB)^{-1}\in R[[t]]
\]
亦由 \(\{p_n\}_{n=1}^{N}\) 唯一确定。
\end{conclusion}
\begin{proof}[证明要点]
Newton 恒等式以有限步将幂和 \(\{p_n\}_{n=1}^{N}\) 反演为基本对称多项式 \(\{e_k\}_{k=1}^{N}\)，而
\(
\det(I-tB)=\sum_{k=0}^{N}(-1)^ke_kt^k
\)。
故结论成立。证毕。
\end{proof}

\begin{conclusion}[有限签名的迹化压缩：以有限条迹替代 Perron 根坐标仍保持切片判等完备性]\label{con:xi-trace-compressed-signature-complete}
在有界复杂度切片 \(\mathbf{PW}^{\ast}_{\le}\) 的设定下，沿用有限签名 \(\Sigma_m(w)\) 的定义（注 \ref{rem:pom-finite-signature-sigmam}）与保守性定理 \ref{thm:pom-sem-conservative}。
对每个整数 \(q\ge 2\) 与温度参数 \(u\)（落在切片网格 \(\mathcal{G}_{\mathrm{grid}}\) 上），令 \(\widetilde A_q(u)\) 表示该坐标所用的有限维 moment-kernel realization（例如 \S\ref{subsec:pom-mom-rewrite-normal} 的对称幂商核），并记 \(N_q:=\dim \widetilde A_q(u)\)（在切片口径下仅依赖 \(q\)）。
定义迹化签名
\[
\Sigma_m^{\mathrm{tr}}(w):=
\Bigl(
G^{\mathrm{BC}}_m(w),\
\{\mathsf{Anom}_{G}(K;\theta)\}_{\theta\in\Theta_{\mathrm{grid}}},\
\{\Tr(\widetilde A_q(u)^n)\}_{(q,u)\in\mathcal{G}_{\mathrm{grid}},\ 1\le n\le N_q}
\Bigr).
\]
则在 \(\mathbf{PW}^{\ast}_{\le}\) 上成立严格蕴含：
\[
\Sigma_m^{\mathrm{tr}}(w_1)=\Sigma_m^{\mathrm{tr}}(w_2)\quad\Longrightarrow\quad w_1\simeq w_2.
\]
\end{conclusion}
\begin{proof}[证明要点]
若 \(\Sigma_m^{\mathrm{tr}}(w_1)=\Sigma_m^{\mathrm{tr}}(w_2)\)，则对每个 \((q,u)\in\mathcal{G}_{\mathrm{grid}}\) 有
\(
\Tr(\widetilde A_q^{(1)}(u)^n)=\Tr(\widetilde A_q^{(2)}(u)^n)
\)（\(1\le n\le N_q\)）。
由结论 \ref{con:xi-newton-trace-completeness-determinant-zeta}（取 \(B=\widetilde A_q(u)\)）可推出
\(
\det(I-t\widetilde A_q^{(1)}(u))=\det(I-t\widetilde A_q^{(2)}(u))
\)，
从而谱半径（Perron 根）一致：\(r_q^{(1)}(u)=r_q^{(2)}(u)\)。
因此 \(\Sigma_m(w_1)=\Sigma_m(w_2)\)，再由定理 \ref{thm:pom-sem-conservative} 得 \(w_1\simeq w_2\)。证毕。
\end{proof}

\endinput


\begin{conclusion}[同调极小化刻画：\(A_{\mathrm{flip}}\) 是 \(\lbrack\beta\rbrack\) 的 \texorpdfstring{$\ZZ_2$}{Z/2}-奇闭链能量极小值]\label{con:xi-split-tunneling-aflip-homological-minimization}
令 \(\Gamma\) 为有限有向图，\(\beta:E(\Gamma)\to\ZZ_2\) 为 \(1\)-余圈，\(\kappa:E(\Gamma)\to\ZZ_{\ge 0}\) 为能量权重。
对任意闭链 \(c\)（\(\ZZ_2\)-系数的 \(1\)-循环）记其能量
\(
\kappa(c):=\sum_{e\in c}\kappa(e)
\)，并记配对
\(
\langle[\beta],c\rangle:=\sum_{e\in c}\beta(e)\in\ZZ_2
\)。
则在约定 \(\min\varnothing:=+\infty\) 下有
\[
A_{\mathrm{flip}}
=
\min\bigl\{\kappa(c):\ c\ \text{为}\ \ZZ_2\text{-}1\text{-循环且}\ \langle[\beta],c\rangle=1\bigr\}\in\ZZ_{>0}\cup\{+\infty\}.
\]
特别地，\(A_{\mathrm{flip}}=+\infty\) 当且仅当 \([\beta]=0\in H^{1}(\Gamma;\ZZ_2)\)，亦即 \(\beta\) 为余边界并对每条闭链满足 \(\langle[\beta],c\rangle=0\)。
特别地，\(A_{\mathrm{flip}}\) 仅依赖同调类 \([\beta]\in H^1(\Gamma;\ZZ_2)\) 与权重 \(\kappa\) 的 \(\ZZ_2\)-循环极小化问题，而与“primitive/无回溯”的语法选择无关；primitive 条件只决定实现该极小值的代表是否需要进一步去除幂重复。
\end{conclusion}
\begin{proof}[证明要点]
任一 Ihara prime 环 \(p\) 诱导一个 \(\ZZ_2\)-循环类 \(c_p\)，且 \(\chi_\beta(p)=-1\) 当且仅当 \(\langle[\beta],c_p\rangle=1\)。
因此 \(A_{\mathrm{flip}}\) 的定义等价于在满足奇配对约束的闭链中极小化能量。
反向方向：对任一 \(\langle[\beta],c\rangle=1\) 的 \(\ZZ_2\)-循环，取其在图中分解为有限条闭环之并（\(\ZZ_2\) 意义下的回路分解），至少有一条分量闭环仍满足奇配对；其能量不超过 \(\kappa(c)\)。
据此极小值由某条闭环实现，并与上式极小化问题一致。
\end{proof}

\begin{conclusion}[最短路证书：\(A_{\mathrm{flip}}\) 可化为两层覆盖图上的一次最短路距离]\label{con:xi-split-tunneling-aflip-shortest-path-certificate}
构造 \(\beta\) 的两层覆盖（等价的“奇偶扩展图”）\(\Gamma_\beta\)：顶点为 \(V(\Gamma)\times\ZZ_2\)，对每条边 \(e:v\to v'\) 引入边
\[
(v,\epsilon)\longrightarrow(v',\epsilon+\beta(e)),
\]
并赋权 \(\kappa(e)\)。
则有严格等式
\[
A_{\mathrm{flip}}
=
\min_{v\in V(\Gamma)}\ \operatorname{dist}_{\Gamma_\beta}\bigl((v,0),(v,1)\bigr),
\]
其中右端为带权最短路距离（允许重复点与边；若两点不连通则约定该距离为 \(+\infty\)）。
因此 \(A_{\mathrm{flip}}\) 的数值可由一次标准最短路算法给出，并且最短路本身即为跨支隧穿作用量的可审计见证。
\end{conclusion}
\begin{proof}[证明要点]
在 \(\Gamma_\beta\) 中，从 \((v,0)\) 到 \((v,1)\) 的任意路径投影到 \(\Gamma\) 给出一条从 \(v\) 到 \(v\) 的闭路，其 \(\beta\)-总和为 \(1\)（奇），能量恰为路径权重之和。
反之，\(\Gamma\) 中任一 \(\beta\)-奇闭路提升后把 \((v,0)\) 送到 \((v,1)\)。
因此两侧极小化问题严格等价。
\end{proof}

\begin{conclusion}[零温谱双重态的裂分律：两层提升 Perron 根的离散瞬子系数]\label{con:xi-split-tunneling-perron-doublet-splitting-law}
在结论 \ref{con:xi-split-tunneling-zeta-detect-aflip} 的设定下，令 \(\widetilde B(u,t)\) 为由 \(\beta\) 定义之两层覆盖上的加权无回溯转移算子（Hashimoto 边空间上的提升算子），并按 sheet 分解为 \(2\times2\) 块矩阵
\[
\widetilde B(u,t)=
\begin{pmatrix}
B_{00}(u,t) & B_{01}(u,t)\\
B_{10}(u,t) & B_{11}(u,t)
\end{pmatrix}.
\]
假设零能量骨架不发生 sheet-交换，从而 \(B_{01}(0,t)=B_{10}(0,t)=0\)；
并假设共同对角块
\(
B_{00}(0,t)=B_{11}(0,t)=:B_{0}(t)
\)
为 primitive，且其 Perron--Frobenius 特征值 \(\lambda_{0}(t)\) 单。
取左右 Perron 向量 \(r(t),\ell(t)\) 归一化为 \(\ell(t)^{\!\top}r(t)=1\)。
令
\[
C(t):=[u^{A_{\mathrm{flip}}}]\,B_{01}(u,t),
\qquad
D(t):=[u^{A_{\mathrm{flip}}}]\,B_{10}(u,t)
\]
为跨支耦合的首个非零系数块（方括号表示取 \(u^{A_{\mathrm{flip}}}\) 系数）。
并记
\[
\gamma(t):=\sqrt{\bigl(\ell(t)^{\!\top}C(t)r(t)\bigr)\bigl(\ell(t)^{\!\top}D(t)r(t)\bigr)}.
\]
则 \(\widetilde B(u,t)\) 的两条顶端特征值 \(\lambda_{\pm}(u,t)\) 满足严格裂分渐近式
\[
\lambda_{+}(u,t)-\lambda_{-}(u,t)
=2\gamma(t)\,u^{A_{\mathrm{flip}}}+O\bigl(u^{A_{\mathrm{flip}}+1}\bigr),
\qquad (u\downarrow 0),
\]
其中 \(r=r(t)\)、\(\ell=\ell(t)\)。
该裂分律给出一条离散“瞬子公式”：跨支 prime 环的最小能量 \(A_{\mathrm{flip}}\) 决定谱双重态的首个可见裂分阶数，而首项系数由首阶跨支耦合的 Perron 投影给出。
\end{conclusion}
\begin{proof}[证明要点]
由假设，\(\widetilde B(0,t)=\mathrm{diag}(B_0(t),B_0(t))\) 在 \(\lambda_0(t)\) 处为二重简并。
把 \(\widetilde B(u,t)\) 在 \(u=0\) 处作幂级数展开，\(\beta\)-扭曲首次引入跨支耦合即出现在 \(u^{A_{\mathrm{flip}}}\) 阶，
从而在简并子空间 \(\Span\{(r,0),(0,r)\}\) 上得到有效 \(2\times2\) 微扰矩阵
\(
\bigl(\begin{smallmatrix} a & \ell^{\!\top}Cr\\ \ell^{\!\top}Dr & a\end{smallmatrix}\bigr)
\)
（对角修正 \(a\) 仅平移两条特征值的平均值）。
对该矩阵取特征值即得两根之差的首项系数为
\(
2\sqrt{(\ell^{\!\top}Cr)(\ell^{\!\top}Dr)}
\)，
其余项吸收至 \(O(u^{A_{\mathrm{flip}}+1})\)。
\end{proof}

\begin{conclusion}[余圈同调与滑动块共轭下的分裂隧穿指纹不变性]\label{con:xi-split-tunneling-cohomology-conjugacy-invariance}
设两套有限状态实现 \((\Gamma,\ell)\)、\((\Gamma',\ell')\) 表示同一协议，并存在滑动块共轭 \(\Phi\) 使时间余圈同调等价（\(\ell'\circ\Phi=\ell+\partial b\)）。
若进一步假设能量余圈与分裂余圈在该共轭下亦仅差余边界（即 \(\kappa'\circ\Phi=\kappa+\partial\varphi\)、\(\beta'\circ\Phi=\beta+\partial\psi\)），则对任一闭环 \(p\) 有
\[
\kappa'\!\bigl(\Phi(p)\bigr)=\kappa(p),
\qquad
\chi_{\beta'}\!\bigl(\Phi(p)\bigr)=\chi_{\beta}(p).
\]
从而：
\begin{enumerate}
  \item 隧穿作用量 \(A_K\)（定义 \ref{def:ihara-tunneling-action}）与跨支隧穿作用量 \(A_{\mathrm{flip}}\) 均为协议语义不变量；
  \item 结论 \ref{con:xi-split-tunneling-zeta-detect-aflip} 的 \(u\)-首阶指数与首项系数，以及结论 \ref{con:xi-split-tunneling-perron-doublet-splitting-law} 的谱裂分首项系数，均仅依赖于同调类 \([\kappa]\in H^{1}(\Gamma;\ZZ)\)、\([\beta]\in H^{1}(\Gamma;\ZZ_2)\) 的协议级像，而不依赖于具体编码/recoding。
\end{enumerate}
\end{conclusion}
\begin{proof}[证明要点]
对闭环 \(p\) 而言，余边界项在闭合求和中消失：\(\sum_{e\in p}\partial f(e)=0\)。
因此闭环的能量和与 \(\ZZ_2\) holonomy 在共轭下保持不变，随之 \(A_K\)、\(A_{\mathrm{flip}}\) 及其由 prime 环枚举/谱投影构造出的首项系数均保持不变。
\end{proof}

\begin{conclusion}[超选结构与两尺度弛豫：\texorpdfstring{$A_{\mathrm{flip}}$}{A_flip} 给出跨支混合的 Arrhenius 指数]\label{con:xi-split-tunneling-superselection-arrhenius-exponent}
令 \(\widetilde P(u)\) 为由 \(\widetilde B(u,1)\) 逐行归一化得到的随机转移算子，并假设在 \(u=0\) 时两张 sheet 各自不可约且非周期。
则在 \(u\downarrow 0\) 的低温极限中：
\begin{enumerate}
  \item 跨 sheet 的最慢弛豫时间尺度具有显式主项：
  \[
  1-\lambda_2(u)=\frac{2\gamma(1)}{\lambda_0(1)}u^{A_{\mathrm{flip}}}+O\bigl(u^{A_{\mathrm{flip}}+1}\bigr),
  \]
  从而跨支混合时间满足
  \[
  \tau_{\mathrm{mix}}(u)=\frac{\lambda_0(1)}{2\gamma(1)}u^{-A_{\mathrm{flip}}}\bigl(1+o(1)\bigr).
  \]
  \item 若 \(A_{\mathrm{flip}}>A_K\)，则所有能量阶 \(<A_{\mathrm{flip}}\) 的贡献均在每张 sheet 内闭合；
  因此在低温展开的 \(u^{A_K},u^{A_K+1},\dots,u^{A_{\mathrm{flip}}-1}\) 各阶上形成严格“超选结构”：跨支耦合在 \(\Theta(u^{A_{\mathrm{flip}}})\) 之前不可见。
\end{enumerate}
\end{conclusion}
\begin{proof}[证明要点]
在 \(u=0\) 时 \(\widetilde P(0)\) 为两块互不连通的不可约链，因而存在一条简并特征值 \(1\) 的双重态。
跨支跃迁的最小能量阶为 \(A_{\mathrm{flip}}\) 意味着跨块耦合概率规模为 \(\Theta(u^{A_{\mathrm{flip}}})\)；结合结论 \ref{con:xi-split-tunneling-perron-doublet-splitting-law} 的裂分系数 \(\gamma(1)\) 并注意行归一化对谱的比例缩放，得到
\(
1-\lambda_2(u)=\frac{2\gamma(1)}{\lambda_0(1)}u^{A_{\mathrm{flip}}}+O(u^{A_{\mathrm{flip}}+1})
\)，从而 \(\tau_{\mathrm{mix}}(u)\asymp(1-\lambda_2(u))^{-1}\) 给出所示主项。
若 \(A_{\mathrm{flip}}>A_K\)，则在 \(u^{A_{\mathrm{flip}}}\) 阶之前跨块项严格为零，得到逐阶超选结构。
\end{proof}

\begin{conclusion}[端点压缩的锐常数：Poisson--Busemann 高度阈值的角窗长度主项]\label{con:xi-split-tunneling-endpoint-window-sqrt-law}
在圆盘模型中令 \(w=\rho e^{\mathrm{i}\theta}\)，并取端点 \(-1\) 的 Poisson--Busemann 权重（结论 \ref{con:xi-timefiber-poisson-busemann-identity}）
\[
B^{-1}(w)=\frac{1-\rho^2}{|1+\rho e^{\mathrm{i}\theta}|^2}.
\]
固定阈值 \(y_0>0\)，定义端点角窗
\[
W(\rho;y_0):=\bigl\{\theta\in[-\pi,\pi]:\ B^{-1}(\rho e^{\mathrm{i}\theta})\ge y_0\bigr\}.
\]
则当 \(\rho\uparrow 1\) 时，\(W(\rho;y_0)\) 为以 \(\theta=\pi\) 为中心的对称区间，且其 Lebesgue 长度满足
\[
|W(\rho;y_0)|
=
\frac{2}{\sqrt{y_0}}\sqrt{1-\rho^2}+O(1-\rho^2).
\]
特别地，取 \(\rho(u):=\sqrt{1-u^{A}}\)（\(A\in\ZZ_{>0}\)）则
\[
|W(\rho(u);y_0)|
=
\frac{2}{\sqrt{y_0}}u^{A/2}+O(u^{A}).
\]
结合“等半径纤维双端触及并把高高度贡献压缩到端点邻域”的端点压缩机制（定理 \ref{thm:dephys-timefiber-two-ended}），得到分辨率律：
若某类效应的低温阶数为 \(u^{A}\)，则其在端点压缩读出中必局域于角宽 \(\asymp u^{A/2}\) 的端点窗；取 \(A=A_{\mathrm{flip}}\) 即得跨支隧穿首项的端点读出分辨率下界。
\end{conclusion}
\begin{proof}[证明要点]
由恒等式 \(|1+\rho e^{\mathrm{i}\theta}|^2=1+\rho^2+2\rho\cos\theta\)，不等式 \(B^{-1}(\rho e^{\mathrm{i}\theta})\ge y_0\) 等价于
\[
1+\rho^2+2\rho\cos\theta\le \frac{1-\rho^2}{y_0}.
\]
令 \(\theta=\pi\pm\delta\)（\(\delta\ge 0\)），则 \(\cos\theta=-\cos\delta\)，从而边界 \(\delta=\delta(\rho)\) 由一元方程确定且 \(\delta(\rho)\to 0\) 当 \(\rho\uparrow 1\)。
对小 \(\delta\) 用 \(\cos\delta=1-\delta^2/2+O(\delta^4)\) 展开，并代入 \(\rho\uparrow 1\) 的尺度 \(1-\rho^2\to 0\)，得到
\[
\delta(\rho)=\frac{1}{\sqrt{y_0}}\sqrt{1-\rho^2}+O(1-\rho^2),
\]
从而 \(|W(\rho;y_0)|=2\delta(\rho)\) 给出所示主项与误差阶。
\end{proof}

\begin{conclusion}[\texorpdfstring{$A_{\mathrm{flip}}=1$}{A_flip=1} 的可审计充要证书：一条见证边与零能量闭合]\label{con:xi-split-tunneling-aflip1-auditable-certificate}
沿用结论 \ref{con:xi-split-tunneling-zeta-detect-aflip}，并进一步假设零能量骨架在两张 sheet 上各自闭合（等价地：\(\kappa(e)=0\Rightarrow \beta(e)=0\)）。
则以下两条件等价：
\begin{enumerate}
  \item \(A_{\mathrm{flip}}=1\)；
  \item 存在一条无回溯边 \(e\) 满足 \(\kappa(e)=1\) 且 \(\beta(e)=1\)，并且存在一条仅由零能量边组成的无回溯路径 \(\gamma\) 使 \(p:=e\gamma\) 构成 tailless 的 Ihara prime 环。
\end{enumerate}
该证书在“见证边 \(+\) 零能量无回溯闭合”的意义上是严格有限可核验的，从而把跨支隧穿首激发态的存在性判别归约为一次有限状态搜索。
\end{conclusion}
\begin{proof}[证明要点]
\((1)\Rightarrow(2)\)：由 \(A_{\mathrm{flip}}=1\) 存在 Ihara prime 环 \(p\) 满足 \(\kappa(p)=1\) 且 \(\chi_\beta(p)=-1\)。
由于 \(\kappa\ge 0\) 且取值为整数，必存在且仅存在一条边 \(e\in p\) 满足 \(\kappa(e)=1\)，其余边能量为零。
零能量闭合假设给出这些边皆满足 \(\beta=0\)，故 \(\chi_\beta(p)=-1\) 强制 \(\beta(e)=1\)。
取 \(\gamma\) 为 \(p\) 中删去该边后的剩余路径即可。
\par
\((2)\Rightarrow(1)\)：由 \(p=e\gamma\) 为 Ihara prime 环且 \(\kappa(e)=1\)、\(\kappa(\gamma)=0\)，得 \(\kappa(p)=1\)；
又由 \(\beta(e)=1\) 且 \(\beta(\gamma)=0\) 得 \(\chi_\beta(p)=-1\)，从而 \(A_{\mathrm{flip}}\le 1\)。
结合 \(A_{\mathrm{flip}}\in\ZZ_{>0}\cup\{+\infty\}\) 且 \(A_{\mathrm{flip}}\le 1\) 得 \(A_{\mathrm{flip}}=1\)。
\end{proof}

\begin{conclusion}[有限观测反演：由有限阶迹数据抽取 \texorpdfstring{$A_{\mathrm{flip}}$}{A_flip} 与裂分系数]\label{con:xi-split-tunneling-finite-trace-inversion}
令 \(d:=\dim(\widetilde B(u,t))\) 为两层提升后的 Hashimoto 边空间维数，并假设 \(\widetilde B(u,t)\) 的条目在 \(u\) 上为多项式（能量权重有限）。
则：
\begin{enumerate}
  \item \(\det(I-\widetilde B(u,t))\) 的系数在 \(u\) 上为多项式，次数存在上界 \(O(d)\)；
  因此 \(Z_K(u,t;\beta)/Z_K(u,t)\) 作为 \(u\)-形式幂级数的首阶指数 \(A_{\mathrm{flip}}\) 与首项系数由有限多个 \(u\)-系数决定。
  \item 通过恒等式
  \[
  \log\det(I-\widetilde B)
  =-\sum_{n\ge 1}\frac{1}{n}\Tr(\widetilde B^{\,n}),
  \]
  结合 Newton 恒等式，可由有限族迹数据 \(\{\Tr(\widetilde B(u,t)^n)\}_{1\le n\le d}\) 的 \(u\)-展开系数有限地重建 \(\det(I-\widetilde B(u,t))\) 的 \(u\)-首项结构，进而抽取：
  \[
  A_{\mathrm{flip}},\qquad
  \sum_{\kappa(p)=A_{\mathrm{flip}},\ \chi_\beta(p)=-1}t^{\ell(p)},
  \qquad
  \sqrt{(\ell^{\!\top}Cr)(\ell^{\!\top}Dr)}.
  \]
\end{enumerate}
从而“时间纤维分裂 \(\times\) 隧穿”的首个可见效应可被封闭为一套有限阶矩/迹的可审计证书体系，而不依赖极限论证或不可控拟合。
\end{conclusion}
\begin{proof}[证明要点]
第一条由行列式展开与条目次数上界直接给出。
第二条由 \(\log\det\) 的迹级数恒等式给出 ghost（迹）坐标，并由 Newton 恒等式把有限个迹幂和反演为特征多项式系数；对这些系数再作 \(u\)-展开即可定位首个非零阶与其系数。
\end{proof}

\endinput



\subsubsection{素数前缀的 Euler 因子对齐：由分裂乘法约化实现的 \texorpdfstring{$\zeta/L$}{zeta/L} 桥与半平面一致逼近}
前述 Bernoulli\((p)\) 倾斜核 \(Q_p(u)=\bigl(P(p)_{ij}u^{g(i,j)}\bigr)\) 的 Perron 根 \(\lambda(u,p)=\rho(Q_p(u))\) 受四次谱方程刚性约束（定理 \ref{thm:fold-gauge-anomaly-bernoulli-p-laws}，式 \eqref{eq:fold-gauge-anomaly-bernoulli-p-pressure-quartic}），从而把局部代数数据与 Euler 型乘积接口锁定为可审计对象。
本段给出一条算术几何侧的对应：通过在有限素数前缀上强制分裂乘法约化，可把椭圆曲线的局部 Euler 因子与 \(\zeta\)（或二次 Dirichlet \(L\)）的前缀因子严格对齐，并在 \(\Re(s)>\tfrac32\) 上获得统一逼近与全阶解析稳定。

\begin{conclusion}[primorial 标度下的 \texorpdfstring{$\zeta$}{zeta}-一致逼近：分裂乘法约化的前缀对齐与尾和误差界]\label{con:xi-primorial-elliptic-zeta-uniform-approx}
令 \(p_1<p_2<\cdots\) 为奇素数序列，并记
\[
M_n:=\prod_{j=1}^{n}p_j.
\]
取整数 \(d_n\) 满足
\[
\gcd\!\bigl(d_n,2M_n(M_n-1)\bigr)=1,
\qquad
\left(\frac{d_n}{p}\right)=\left(\frac{-1}{p}\right)\ \ \ \forall\,p\mid M_n,
\]
并定义二次扭曲的 Legendre 曲线
\[
E_n:\quad d_n y^2=x(x-1)(x-M_n)\qquad(\text{定义于 }\QQ).
\]
令去除 \(p=2\) 的部分欧拉积
\[
L^{(2)}(E_n,s):=\prod_{p\neq 2} L_p(E_n,s)^{-1},
\qquad
\zeta^{(2)}(s):=\prod_{p\neq 2}(1-p^{-s})^{-1}=\zeta(s)\,(1-2^{-s}).
\]
则对任意 \(\sigma>\tfrac32\)，存在常数 \(C_\sigma>0\) 使得对一切满足 \(\Re(s)\ge\sigma\) 的复数 \(s\)，有一致估计
\[
\left|\log\frac{L^{(2)}(E_n,s)}{\zeta^{(2)}(s)}\right|
\ \le\
C_\sigma\sum_{p>p_n}\Bigl(p^{\frac12-\sigma}+p^{1-2\sigma}\Bigr),
\]
特别地，
\[
\frac{L^{(2)}(E_n,s)}{\zeta^{(2)}(s)}\xrightarrow[n\to\infty]{}1
\quad\text{在半平面 }\Re(s)\ge\sigma\text{ 上局部一致收敛。}
\]
\end{conclusion}
\begin{proof}[证明要点]
对任意奇素数 \(p\le p_n\)，有 \(p\mid M_n\)，故在 \(\FF_p\) 上
\[
E_n:\ d_n y^2=x(x-1)(x-M_n)\ \equiv\ d_n y^2=x^2(x-1).
\]
特殊纤维在 \((x,y)=(0,0)\) 处为结点；其切锥由最低次项给出为 \(d_n y^2+x^2=0\)。
该切锥在 \(\FF_p\) 上可分解当且仅当 \(\bigl(\frac{-d_n}{p}\bigr)=1\)。
由假设 \(\bigl(\frac{d_n}{p}\bigr)=\bigl(\frac{-1}{p}\bigr)\) 得 \(\bigl(\frac{-d_n}{p}\bigr)=1\)，故在所有 \(p\le p_n\)（\(p\neq 2\)）处均为分裂乘法约化，并且局部 Euler 因子满足
\[
L_p(E_n,s)=1-p^{-s}\qquad(p\le p_n,\ p\neq 2).
\]
因此
\[
\frac{L^{(2)}(E_n,s)}{\zeta^{(2)}(s)}
=\prod_{p>p_n,\ p\neq 2}\frac{1-p^{-s}}{L_p(E_n,s)}.
\]
\par
固定 \(\Re(s)=\sigma>\tfrac32\)，令 \(t=p^{-s}\)（于是 \(|t|=p^{-\sigma}\)）。
对每个 \(p>p_n\)，存在 \(a_p\in\ZZ\) 与 \(b_p\in\{0,1\}\) 使
\[
L_p(E_n,s)=1-a_p t+b_p pt^2,
\]
其中 \(b_p=1\) 当且仅当 \(E_n\) 在 \(p\) 处良好约化；而在坏约化处可取 \(b_p=0\) 且 \(a_p\in\{0,\pm1\}\)。
并且在良好约化处有 Hasse 界 \(|a_p|\le 2\sqrt p\)，故总可统一写作 \(|a_p|\le 2\sqrt p\)。
因此对充分大的 \(p\)（依赖于 \(\sigma\) 但与 \(n\) 无关）可使 \(|a_p t|+|b_p pt^2|\le\tfrac12\)，从而
\[
\left|\frac{1-t}{L_p(E_n,s)}-1\right|
=\left|\frac{(a_p-1)t-b_p pt^2}{1-a_pt+b_p pt^2}\right|
\ \ll_\sigma\ (|a_p|+1)p^{-\sigma}+p^{1-2\sigma}
\ \ll\ p^{\frac12-\sigma}+p^{1-2\sigma}.
\]
对有限个小素数把常数吸收入 \(C_\sigma\)，得到
\(
\sum_{p>p_n}\bigl|\frac{1-p^{-s}}{L_p(E_n,s)}-1\bigr|
\ll_\sigma
\sum_{p>p_n}(p^{1/2-\sigma}+p^{1-2\sigma})
\)；
右端在 \(\sigma>\tfrac32\) 时收敛且其尾和随 \(n\to\infty\) 趋于 \(0\)。
由无穷乘积的一般判别（\(\sum|u_p-1|<\infty\Rightarrow\prod u_p\) 收敛且尾积趋于 \(1\)）并结合 \(|\log u|\ll|u-1|\)（在 \(|u-1|\le\tfrac12\) 区域内），得到所示对数估计与局部一致收敛。证毕。
\end{proof}

\begin{conclusion}[二次 Dirichlet--椭圆桥：二次特征的前缀 Euler 因子可由乘法退化精确编译]\label{con:xi-primorial-quadratic-dirichlet-elliptic-approx}
令 \(\chi\) 为任意二次 Dirichlet 特征（模 \(q\)，取值于 \(\{0,\pm1\}\)）。
在结论 \ref{con:xi-primorial-elliptic-zeta-uniform-approx} 的 \(M_n\) 基础上，取整数 \(d_{n,\chi}\) 满足
\[
\gcd\!\bigl(d_{n,\chi},2qM_n(M_n-1)\bigr)=1,
\qquad
\left(\frac{d_{n,\chi}}{p}\right)=\chi(p)\left(\frac{-1}{p}\right)\ \ \ \forall\,p\le p_n,\ p\nmid q,
\]
并定义
\[
E_{n,\chi}:\quad d_{n,\chi}y^2=x(x-1)(x-M_n).
\]
记去除 \(2q\) 的部分欧拉积
\[
L^{(2q)}(\chi,s):=\prod_{p\nmid 2q}(1-\chi(p)p^{-s})^{-1},
\qquad
L^{(2q)}(E_{n,\chi},s):=\prod_{p\nmid 2q} L_p(E_{n,\chi},s)^{-1}.
\]
则对任意 \(\sigma>\tfrac32\)，成立局部一致收敛
\[
\frac{L^{(2q)}(E_{n,\chi},s)}{L^{(2q)}(\chi,s)}\xrightarrow[n\to\infty]{}1
\quad\text{在 }\Re(s)\ge\sigma\text{ 上，}
\]
并且存在常数 \(C_\sigma>0\) 使得对 \(\Re(s)\ge\sigma\) 有同型尾和界
\[
\left|\log\frac{L^{(2q)}(E_{n,\chi},s)}{L^{(2q)}(\chi,s)}\right|
\ \le\
C_\sigma\sum_{p>p_n}\Bigl(p^{\frac12-\sigma}+p^{1-2\sigma}\Bigr).
\]
\end{conclusion}
\begin{proof}[证明要点]
对任意素数 \(p\le p_n\) 且 \(p\nmid q\)，有 \(p\mid M_n\)，故 \(E_{n,\chi}\) 在 \(p\) 处为乘法约化并在模 \(p\) 下退化为结点曲线。
与结论 \ref{con:xi-primorial-elliptic-zeta-uniform-approx} 相同，分裂型由切锥 \(d_{n,\chi}y^2+x^2=0\) 的平方类决定：
\[
a_p(E_{n,\chi})=
\left(\frac{-d_{n,\chi}}{p}\right)
\in\{\pm1\},
\qquad
L_p(E_{n,\chi},s)=1-a_p(E_{n,\chi})p^{-s}.
\]
由构造 \(\bigl(\frac{d_{n,\chi}}{p}\bigr)=\chi(p)\bigl(\frac{-1}{p}\bigr)\) 推得
\(
\bigl(\frac{-d_{n,\chi}}{p}\bigr)=\chi(p)
\)，
从而
\[
L_p(E_{n,\chi},s)=1-\chi(p)p^{-s}\qquad(p\le p_n,\ p\nmid 2q).
\]
于是
\[
\frac{L^{(2q)}(E_{n,\chi},s)}{L^{(2q)}(\chi,s)}
=\prod_{p>p_n,\ p\nmid 2q}\frac{1-\chi(p)p^{-s}}{L_p(E_{n,\chi},s)}.
\]
对 \(p>p_n\) 的解析收敛与对数尾和估计与结论 \ref{con:xi-primorial-elliptic-zeta-uniform-approx} 完全同型（只需注意 \(|\chi(p)|\le 1\) 且加法/乘法/良约化三类局部因子均满足同一阶的 \(p^{1/2-\sigma}+p^{1-2\sigma}\) 控制）。
此处 \(d_{n,\chi}\) 的存在性由有限素数集合上指定 Legendre 符号的中国剩余构造保证。证毕。
\end{proof}

\begin{conclusion}[比值 Dirichlet 级数的粗支撑与系数稳定：前缀对齐 \texorpdfstring{$\Rightarrow$}{=>} 有限支撑坍缩]\label{con:xi-primorial-ratio-dirichlet-support-stability}
在结论 \ref{con:xi-primorial-elliptic-zeta-uniform-approx} 的设定下，令
\[
R_n(s):=\frac{L^{(2)}(E_n,s)}{\zeta^{(2)}(s)}
=\prod_{p>p_n,\ p\neq 2}\frac{1-p^{-s}}{L_p(E_n,s)}
=\sum_{m\ge 1} b_{n}(m)\,m^{-s},
\qquad (\Re(s)>\tfrac32),
\]
其中最后一式为 \(\Re(s)>\tfrac32\) 上的绝对收敛 Dirichlet 展开。
则：
\begin{enumerate}
  \item 若 \(b_n(m)\neq 0\)，则 \(m\) 的全部素因子均严格大于 \(p_n\)；
  \item 对任意固定整数 \(m\ge 2\)，存在 \(N(m)\) 使得对所有 \(n\ge N(m)\) 恒有 \(b_n(m)=0\)。
\end{enumerate}
因此 \(R_n(s)\) 的 Dirichlet 系数在逐点意义下最终稳定为常数项 \(b_n(1)=1\)。
\end{conclusion}
\begin{proof}[证明要点]
由 Euler 乘积表达式可见，\(R_n\) 的全部局部因子仅来自 \(p>p_n\) 的素数。
在 \(\Re(s)>\tfrac32\) 上展开为 Dirichlet 级数时，每一项 \(m^{-s}\) 仅能由这些素数的幂与互素乘法闭包生成，故若 \(b_n(m)\neq 0\) 则 \(m\) 的每个素因子都满足 \(>p_n\)，得第 (1) 条。
\par
对固定 \(m\ge 2\)，取 \(N(m)\) 使 \(p_{N(m)}>\max\{p:\ p\mid m\}\)。
则当 \(n\ge N(m)\) 时，\(m\) 含有某个素因子 \(\le p_n\)，与第 (1) 条矛盾，故 \(b_n(m)=0\)，得第 (2) 条。证毕。
\end{proof}

\begin{conclusion}[乘法退化处 Euler 因子的线性化：局部 Weil--Deligne 表示的单 Jordan 块强制机制]\label{con:xi-primorial-wd-linear-euler-factor}
固定素数 \(\ell\)。
对结论 \ref{con:xi-primorial-quadratic-dirichlet-elliptic-approx} 的 \(E_{n,\chi}\)，对任意素数 \(p\le p_n\) 且 \(p\nmid 2q\ell\)，其 \(\ell\)-进 Tate 模 \(V_\ell(E_{n,\chi})\) 在 \(G_{\QQ_p}\) 上诱导的 Weil--Deligne 表示 \((\rho_p,N_p)\) 满足：
\begin{enumerate}
  \item \(N_p\neq 0\) 且 \(\dim\ker(N_p)=1\)；
  \item Frobenius 在 \(\ker(N_p)\) 上的特征值为 \(\chi(p)\)；
  \item 因而局部 \(L\)-因子满足
  \[
  L_p(E_{n,\chi},s)=\det\!\left(1-\mathrm{Fr}_p\,p^{-s}\mid \ker(N_p)\right)=1-\chi(p)p^{-s}.
  \]
\end{enumerate}
该陈述给出“良约化处二次 Euler 因子在乘法退化处坍缩为一次因子”的结构原因：线性化由 \(N_p\) 的非零单 Jordan 块（Tate 曲线型半稳定退化）所强制。
\end{conclusion}
\begin{proof}[证明要点]
由 \(p\le p_n\Rightarrow p\mid M_n\) 可知 \(E_{n,\chi}\) 在 \(p\) 处为乘法约化。
对乘法约化椭圆曲线，局部 \(\ell\)-进表示为 Tate 曲线型半稳定表示，其惯性上出现秩 \(1\) 的幂零算子 \(N_p\) 且 \(\ker(N_p)\) 一维。
分裂/非分裂型对应 Frobenius 在该一维商（等价的一维不变子空间）上的特征值 \(\pm1\)；
由结论 \ref{con:xi-primorial-quadratic-dirichlet-elliptic-approx} 的构造该符号正是 \(\chi(p)\)。
Weil--Deligne 口径下局部因子由 \(\ker(N_p)\) 上的 Frobenius 作用给出，从而得到 \(L_p(E_{n,\chi},s)=1-\chi(p)p^{-s}\)。证毕。
\end{proof}

\begin{conclusion}[全阶解析稳定：对数导数与 Taylor 系数在 \texorpdfstring{$\Re(s)>\tfrac32$}{Re(s)>3/2} 上的局部一致收敛]\label{con:xi-primorial-higher-analytic-stability}
在结论 \ref{con:xi-primorial-elliptic-zeta-uniform-approx} 的设定下，固定 \(\sigma>\tfrac32\) 与整数 \(k\ge 0\)。
则在闭半平面 \(\Re(s)\ge\sigma\) 上存在常数 \(C_{\sigma,k}>0\) 使得
\[
\sup_{\Re(s)\ge\sigma}\left|
\frac{d^k}{ds^k}\log L^{(2)}(E_n,s)
\;-\;
\frac{d^k}{ds^k}\log \zeta^{(2)}(s)
\right|
\ \le\
C_{\sigma,k}\sum_{p>p_n}(\log p)^k\Bigl(p^{\frac12-\sigma}+p^{1-2\sigma}\Bigr),
\]
从而
\[
\frac{d^k}{ds^k}\log L^{(2)}(E_n,s)\xrightarrow[n\to\infty]{\text{局部一致}}
\frac{d^k}{ds^k}\log \zeta^{(2)}(s)
\quad(\Re(s)>\tfrac32).
\]
特别地，对任意 \(s_0\) 满足 \(\Re(s_0)>\tfrac32\)，函数 \(L^{(2)}(E_n,s)\) 在 \(s_0\) 的 Taylor 展开系数随 \(n\to\infty\) 逐项收敛到 \(\zeta^{(2)}(s)\) 的相应系数。
\end{conclusion}
\begin{proof}[证明要点]
由结论 \ref{con:xi-primorial-elliptic-zeta-uniform-approx} 的乘积分解，写
\[
\log\frac{L^{(2)}(E_n,s)}{\zeta^{(2)}(s)}
=\sum_{p>p_n}\log\!\left(\frac{1-p^{-s}}{L_p(E_n,s)}\right).
\]
在 \(\Re(s)\ge\sigma>\tfrac32\) 上，上式由结论 \ref{con:xi-primorial-elliptic-zeta-uniform-approx} 的逐素数估计给出正常收敛，并且其尾和随 \(n\to\infty\) 归零。
由于每一项均为 \(t=p^{-s}\) 的有界解析函数复合 \(t=t(s)\)，对 \(s\) 求导至阶 \(k\) 仅引入至多 \((\log p)^k\) 的因子；
结合 \(|a_p|\le 2\sqrt p\)（良约化处）与坏约化处 \(a_p\in\{0,\pm1\}\)，以及 \(|t|=p^{-\sigma}\)，得到逐项一致界
\[
\left|\frac{d^k}{ds^k}\log\!\left(\frac{1-p^{-s}}{L_p(E_n,s)}\right)\right|
\ll_{\sigma,k}(\log p)^k\Bigl(p^{\frac12-\sigma}+p^{1-2\sigma}\Bigr).
\]
对 \(p>p_n\) 求和即得所示整体上界。
右端和在 \(\sigma>\tfrac32\) 下收敛且其尾和随 \(n\to\infty\) 趋于 \(0\)，于是导数序列在 \(\Re(s)\ge\sigma\) 上局部一致收敛。
Taylor 系数的逐项收敛由局部一致收敛与 Weierstrass 定理得到。证毕。
\end{proof}

\endinput



\subsubsection{隧穿作用量与算术几何接口的末端汇总}\label{subsubsec:xi-tunneling-arithmetic-geometry-terminal-summary}

\paragraph{预解式椭圆曲线与 \texorpdfstring{$S_3$}{S3} 闭包 Jacobian 的算术分解}
沿用结论 \ref{con:xi-terminal-zm-resolvent-cubic-plane-cubic-geometry}、结论 \ref{con:xi-terminal-zm-discriminant-ridge-etale-c3-explicit-model} 与定理 \ref{thm:xi-terminal-zm-s3-closure-jacobian-splitting} 的记号。
本段给出预解式平面三次的极小 Weierstrass 模型、Lee--Yang 判别式剥离恒等式、与导子 \(37\) 椭圆曲线的显式 \(2\)-同余，以及由此导出的 \(C_6\) 与 \(X_e\) 的 Jacobian 椭圆化分解。

\begin{theorem}[预解式椭圆曲线的极小 Weierstrass 模型、判别式与半稳定性]\label{con:xi-resolvent-cubic-short-weierstrass-12096}
设平面三次曲线 \(\mathcal R\subset\mathbb A^2_{y,z}\) 由
\[
\mathscr R(z,y):=z^3+(1+2y)z^2-(1+4y+4y^2)z-(1+5y+13y^2+8y^3)=0
\]
给出。则 \(\mathcal R\) 在 \(\QQ\) 上与下述短 Weierstrass 曲线双有理同构：
\[
E_R:\quad Y^2=X^3-\frac{28}{3}X-\frac{667}{108}.
\]
其不变量为
\[
c_4(E_R)=448,\qquad c_6(E_R)=5336,\qquad \Delta(E_R)=35557=31^2\cdot 37,
\]
\[
j(E_R)=\frac{c_4(E_R)^3}{\Delta(E_R)}=\frac{89915392}{35557}
=\frac{2^{18}\cdot 7^3}{31^2\cdot 37}.
\]
并且 \(E_R\) 在 \(p=31\) 处为非分裂乘法约化型 \(I_2\)，在 \(p=37\) 处为非分裂乘法约化型 \(I_1\)；
因此 \(E_R\) 为半稳定椭圆曲线，导子为
\[
N(E_R)=31\cdot 37=1147.
\]
等价地，可取 \(\QQ\)-同构整系数模型
\[
E_R^{\mathrm{int}}:\quad y^2=x^3-12096x+288144,
\]
其判别式满足 \(\Delta(E_R^{\mathrm{int}})=2^{12}3^{12}\Delta(E_R)\)。
\end{theorem}
\begin{proof}
由结论 \ref{con:xi-terminal-zm-resolvent-cubic-plane-cubic-geometry}，\(\mathcal R\) 为非奇异平面三次，故其规范化为亏格 \(1\) 曲线；
取有理点 \(P=(y,z)=(0,1)\)，以过 \(P\) 的直线族参数化可得与 \(\mathcal R\) 双有理等价的二重覆盖模型，再经标准变量代换化为上式短 Weierstrass 方程。
\par
对 \(A=-28/3\)、\(B=-667/108\) 应用短 Weierstrass 公式
\[
c_4=-48A,\qquad c_6=-864B,\qquad \Delta=-16(4A^3+27B^2),
\]
得到 \(c_4=448\)、\(c_6=5336\)、\(\Delta=35557\)，并由 \(j=c_4^3/\Delta\) 得所示 \(j\)-不变量。
\par
对 \(p\in\{31,37\}\)，有 \(v_p(\Delta)>0\) 且 \(v_p(c_4)=0\)，故均为乘法约化，且型别分别为 \(I_{v_p(\Delta)}=I_2,I_1\)。
再由乘法约化分裂判据 \(\left(\frac{-c_6}{p}\right)\)：
\[
-c_6\equiv 27\pmod{31},\qquad -c_6\equiv 29\pmod{37},
\]
两者均为非二次剩余，故两处均非分裂。
半稳定性与乘法约化情形下的导子指数给出 \(N(E_R)=31\cdot 37\)。
整系数模型由 \((x,y)=(36X,216Y)\) 直接得到。证毕。
\end{proof}

\begin{proposition}[Lee--Yang 判别式的 \texorpdfstring{$3$}{3}-幂剥离与预解式椭圆判别式恒等式]\label{prop:xi-resolvent-disc-leyang-3adic-peeling}
\leanverified{paper\_xi\_resolvent\_disc\_leyang\_3adic\_peeling}
设
\[
P_{\mathrm{LY}}(y):=256y^3+411y^2+165y+32.
\]
则有严格恒等式
\[
\Delta(E_R)=-\frac{\mathrm{Disc}(P_{\mathrm{LY}})}{3^9}=31^2\cdot 37.
\]
\end{proposition}
\begin{proof}
由结论 \ref{con:xi-terminal-zm-leyang-cubic-arithmetic} 已知
\(
\mathrm{Disc}(P_{\mathrm{LY}})=-3^9\cdot 31^2\cdot 37
\)；
由定理 \ref{con:xi-resolvent-cubic-short-weierstrass-12096} 有
\(
\Delta(E_R)=31^2\cdot 37
\)。
两式相除即得结论。证毕。
\end{proof}

\begin{theorem}[导子 \texorpdfstring{$37$}{37} 曲线与导子 \texorpdfstring{$1147$}{1147} 预解式曲线的显式 \texorpdfstring{$2$}{2}-同余]\label{con:xi-er-badprimes-31-37-and-2division-qsqrt37}
令
\[
E_{37}:\quad y^2+y=x^3-x
\qquad\Bigl(\text{等价于 }Y^2=X^3-X+\frac14\Bigr),
\]
并令 \(E_R\) 如定理 \ref{con:xi-resolvent-cubic-short-weierstrass-12096}。
记两者的 \(2\)-除法三次分别为
\[
f_{37}(u)=u^3-u+\frac14,\qquad
f_R(v)=v^3-\frac{28}{3}v-\frac{667}{108}.
\]
则存在显式 Tschirnhaus 变换
\[
v=4u^2-u-\frac{8}{3},
\qquad
u=\frac{224-39v-36v^2}{279},
\]
使得 \(u\) 为 \(f_{37}\) 的根当且仅当 \(v\) 为 \(f_R\) 的根。
因此诱导同构
\[
\QQ(E_{37}[2])\cong\QQ(E_R[2]),
\qquad
E_{37}[2]\simeq E_R[2]\ \ (\text{作为 }G_\QQ\text{-模}).
\]
\end{theorem}
\begin{proof}
直接代入并对 \(f_{37},f_R\) 作带余除法可得
\[
f_R\!\left(4u^2-u-\frac83\right)\equiv 0\pmod{f_{37}(u)},
\qquad
f_{37}\!\left(\frac{224-39v-36v^2}{279}\right)\equiv 0\pmod{f_R(v)},
\]
从而给出三次子域同构 \(\QQ(u)\cong\QQ(v)\)。
\par
另一方面，
\[
\mathrm{Disc}(f_{37})=\frac{37}{16},\qquad
\mathrm{Disc}(f_R)=\frac{\Delta(E_R)}{16}=\frac{35557}{16},
\]
两者平方自由部分同为 \(37\)。
对不可约三次，\(2\)-除法分裂域等于“任一根生成的三次子域”与“判别式平方根域”的复合；
故两者均为相应三次子域与 \(\QQ(\sqrt{37})\) 的复合。
由三次子域同构与二次 resolvent 一致，即得 \(\QQ(E_{37}[2])\cong\QQ(E_R[2])\)。
利用 \(\mathrm{GL}_2(\FF_2)\cong S_3\) 的标准识别，得到 \(G_\QQ\)-模同构 \(E_{37}[2]\simeq E_R[2]\)。证毕。
\end{proof}

\begin{corollary}[Frobenius 迹的模 \texorpdfstring{$2$}{2} 同余与点数奇偶同步]\label{cor:xi-er-frobenius-mod2-parity-sync}
对任意素数 \(p\nmid 2\cdot 31\cdot 37\)，有
\[
a_p(E_{37})\equiv a_p(E_R)\pmod 2.
\]
等价地，
\[
\#E_{37}(\FF_p)\equiv \#E_R(\FF_p)\pmod 2.
\]
\end{corollary}
\begin{proof}
由定理 \ref{con:xi-er-badprimes-31-37-and-2division-qsqrt37}，\(E_{37}[2]\simeq E_R[2]\) 为 \(G_\QQ\)-模同构。
对好素数 \(p\nmid 2N(E_{37})N(E_R)\)，有
\(
\mathrm{tr}(\mathrm{Frob}_p\mid E[2])\equiv a_p(E)\pmod 2
\)。
故两曲线在模 \(2\) 下 Frobenius 迹一致，点数奇偶同余式随之成立。证毕。
\end{proof}

\begin{theorem}[六次 \texorpdfstring{$S_3$}{S3}-闭包曲线的双椭圆性与 Prym 因子平方分解]\label{con:xi-s3-closure-jacobian-er2-refinement}
设
\[
C:\ w^2=-y(y-1)P_{\mathrm{LY}}(y),
\qquad
P_{\mathrm{LY}}(y)=256y^3+411y^2+165y+32,
\]
并令 \(C_6\) 为系统
\[
w^2=-y(y-1)P_{\mathrm{LY}}(y),\qquad \mathscr R(z,y)=0
\]
确定的归一化光滑射影模型。
记 \(\pi:C_6\to C\) 为 \((y,w,z)\mapsto (y,w)\)；
由既有结果，\(\pi\) 为无分歧循环三次覆盖，且 \(C_6\to\PP^1_y\) 为次数 \(6\) 的 \(S_3\)-伽罗瓦覆盖，并有
\[
J(C_6)\sim_{\QQ}J(C)\times A_2,\qquad J(C)\ \text{在 }\QQ\text{ 上简单}.
\]
则：
\begin{enumerate}
  \item 自然投影
  \[
  \rho:C_6\longrightarrow \overline{\mathcal R},\qquad (y,w,z)\longmapsto (y,z)
  \]
  为定义于 \(\QQ\) 的次数 \(2\) 态射，故 \(C_6\) 为双椭圆曲线。
  \item \(A_2\sim_{\QQ}E_R^2\)，从而
  \[
  J(C_6)\sim_{\QQ}J(C)\times E_R^2.
  \]
\end{enumerate}
\end{theorem}
\begin{proof}
第一条直接来自方程结构：固定 \((y,z)\in\overline{\mathcal R}\) 后，\(w\) 由二次方程
\(
w^2=-y(y-1)P_{\mathrm{LY}}(y)
\)
确定，故一般纤维有 \(2\) 个点；归一化后得到定义于 \(\QQ\) 的有限二次态射 \(\rho\)。
\par
第二条中，\(\rho\) 诱导非零同态
\[
\rho_*:J(C_6)\longrightarrow J(\overline{\mathcal R})\cong E_R.
\]
其像为椭圆曲线 \(E_R\)。
结合 \(J(C_6)\sim_{\QQ}J(C)\times A_2\) 且 \(J(C)\) 在 \(\QQ\) 上简单并无椭圆因子，\(E_R\) 必为 \(A_2\) 的 \(\QQ\)-同源因子。
\par
又因 \(C_6\to\PP^1_y\) 为 \(S_3\)-伽罗瓦覆盖，\(\QQ[S_3]\) 作用嵌入 \(\mathrm{End}^0_{\QQ}(J(C_6))\)。
在标准二维表示分量上得到
\(
M_2(\QQ)\subseteq \mathrm{End}^0_{\QQ}(A_2)
\)。
因此 \(A_2\) 在 \(\QQ\)-同源意义下必为某椭圆曲线平方 \(E^2\)。
由前述 \(E_R\) 已是 \(A_2\) 的同源因子，得 \(E\sim_{\QQ}E_R\)，故 \(A_2\sim_{\QQ}E_R^2\)。
代回即得 \(J(C_6)\sim_{\QQ}J(C)\times E_R^2\)。证毕。
\end{proof}

\begin{corollary}[\texorpdfstring{$C_6$}{C6} 的局部 Hasse--Weil 因子分解]\label{con:xi-s3-closure-local-lfactor-square-er}
对任意好素数 \(p\nmid 2\cdot 3\cdot 31\cdot 37\)，记
\[
P_{C_6,p}(T):=\det\!\left(1-T\mathrm{Frob}_p\mid H^1_{\acute et}(C_{6,\overline{\FF}_p},\QQ_\ell)\right).
\]
则有
\[
P_{C_6,p}(T)=P_{C,p}(T)\cdot\bigl(1-a_p(E_R)T+pT^2\bigr)^2.
\]
\end{corollary}
\begin{proof}
由定理 \ref{con:xi-s3-closure-jacobian-er2-refinement} 的 \(\QQ\)-同源分解
\(
J(C_6)\sim_{\QQ}J(C)\times E_R^2
\)，
在好素数处 \(\ell\)-进 Tate 模与 \(H^1\) 的 Frobenius 作用均按直和分解。
故局部特征多项式按因子乘积分解，且椭圆因子出现两次。证毕。
\end{proof}

\begin{corollary}[\texorpdfstring{$S_4$}{S4}-闭包曲线 \texorpdfstring{$X_e$}{Xe} 的 Jacobian 分解进一步椭圆化]\label{cor:xi-s4-jacobian-er2-elliptification}
沿用 \(S_4\)-闭包曲线 \(X_e\) 的同源分解
\[
J(X_e)\sim_{\QQ}A_{\mathrm{sgn}}\times A_2\times A_3\times A_3',
\qquad
A_{\mathrm{sgn}}\sim_{\QQ}J(C),
\]
并结合 \(J(C_6)\sim_{\QQ}A_{\mathrm{sgn}}\times A_2\)。
则由定理 \ref{con:xi-s3-closure-jacobian-er2-refinement} 可得
\[
J(X_e)\sim_{\QQ}J(C)\times E_R^2\times A_3\times A_3'.
\]
\end{corollary}
\begin{proof}
将 \(A_{\mathrm{sgn}}\sim_{\QQ}J(C)\) 与定理 \ref{con:xi-s3-closure-jacobian-er2-refinement} 给出的
\(
A_2\sim_{\QQ}E_R^2
\)
代入 \(J(X_e)\) 的既有分解式，直接得到结论。证毕。
\end{proof}

\paragraph{\texorpdfstring{$S_3$}{S3} 闭包的双覆盖分歧与等变 Jacobian 精化}
沿用本段既定记号：\(\mathscr R(z,y)=0\) 为预解式平面三次，\(\overline{\mathcal R}\) 为其射影闭包，
\[
C:\ w^2=-y(y-1)P_{\mathrm{LY}}(y),\qquad
P_{\mathrm{LY}}(y)=256y^3+411y^2+165y+32,
\]
且 \(C_6\) 为纤维积 \(C\times_{\PP^1_y}\overline{\mathcal R}\) 的归一化。

\begin{theorem}[从 \texorpdfstring{$C_6$}{C6} 到 \texorpdfstring{$\overline{\mathcal R}$}{Rbar} 的双覆盖与显式分歧除子]\label{thm:xi-s3-closure-double-cover-explicit-branch-divisor}
定义自然投影
\[
\pi_{\overline{\mathcal R}}:C_6\longrightarrow \overline{\mathcal R},\qquad
(y,w,z)\longmapsto (y,z).
\]
则 \(\pi_{\overline{\mathcal R}}\) 为次数 \(2\) 的有限态射。其分歧除子的约化支撑恰为六点
\[
\bigl\{Q_0,\ Q_1,\ Q_\infty,\ Q_{\alpha_1},\ Q_{\alpha_2},\ Q_{\alpha_3}\bigr\},
\]
其中：
\[
Q_0=(y,z)=(0,1),\qquad
Q_1=(y,z)=(1,3),\qquad
Q_\infty=[2:1:0]\in\overline{\mathcal R},
\]
\[
P_{\mathrm{LY}}(\alpha_i)=0,\qquad
Q_{\alpha_i}=(y,z)=\bigl(\alpha_i,s(\alpha_i)\bigr)\ (i=1,2,3),
\]
并且
\[
A(y):=1+2y,\qquad
s(y):=A(y)+\frac{9y(y-1)}{4A(y)^2}.
\]
等价地，
\[
\operatorname{Ram}(\pi_{\overline{\mathcal R}})_{\mathrm{red}}
=Q_0+Q_1+Q_\infty+\sum_{i=1}^{3}Q_{\alpha_i}.
\]
\end{theorem}

\begin{proof}
次数 \(2\) 来自模型
\(
w^2=-y(y-1)P_{\mathrm{LY}}(y)
\)
在 \((y,z)\) 固定后对 \(w\) 给出二次方程。

记 \(\overline{\mathcal R}\to\PP^1_y\) 为 \(y\)-投影。由
\[
\operatorname{Disc}_z\bigl(\mathscr R(z,y)\bigr)=-y(y-1)P_{\mathrm{LY}}(y)
\]
可知其分歧值集合为
\(
B=\{0,1,\infty\}\cup\{P_{\mathrm{LY}}=0\}
\)。
在每个 \(y_0\in B\) 上，\(\mathscr R(\cdot,y_0)\) 纤维型为 \((2,1)\)：存在唯一分歧点 \(Q_{\mathrm{ram}}\)（局部分歧指数 \(2\)）与唯一未分歧点 \(Q_{\mathrm{unr}}\)（局部分歧指数 \(1\)）。

将 \(C\to\PP^1_y\) 的二次分歧沿 \(\overline{\mathcal R}\to\PP^1_y\) 基变换，并在 \(Q\in\overline{\mathcal R}\) 处应用驯分歧 Abhyankar 引理，得
\[
e\!\left(\pi_{\overline{\mathcal R}}\ \text{at}\ Q\right)
=\frac{2}{\gcd\!\left(2,e\!\left(\overline{\mathcal R}\to\PP^1_y\ \text{at}\ Q\right)\right)}.
\]
故在 \(Q_{\mathrm{ram}}\) 处 \(\gcd(2,2)=2\)，不分歧；在 \(Q_{\mathrm{unr}}\) 处 \(\gcd(2,1)=1\)，分歧。于是每个 \(y_0\in B\) 恰贡献一个分歧点，共 \(6\) 点。

显式坐标如下。先对有限 \(y_0\in B\)，把
\[
\mathscr R(z,y_0)=z^3+bz^2+cz+d
\]
写为 \((z-r)^2(z-s)\)。
对“判别式为零且非三重根”的单首三次，简单根满足
\[
s=\frac{-b^3+4bc-9d}{b^2-3c}.
\]
在本题 \(b=A(y)\)、\(c=-A(y)^2\)、\(d=-(A(y)^3+y(y-1))\) 下化简得
\[
s(y)=A(y)+\frac{9y(y-1)}{4A(y)^2}.
\]
代入 \(y_0=0,1\) 得 \(s(0)=1,\ s(1)=3\)，即 \(Q_0,Q_1\)。
对 \(y_0=\alpha_i\)（\(P_{\mathrm{LY}}(\alpha_i)=0\)）得 \(Q_{\alpha_i}\)。

在无穷远处，将 \(\mathscr R\) 齐次化后取 \(W=0\)，得到
\[
(Z-2Y)(Z+2Y)^2=0,
\]
故 \([2:1:0]\) 为未分歧点、\([-2:1:0]\) 为分歧点；
按上述判据，\(\pi_{\overline{\mathcal R}}\) 在 \([2:1:0]\) 处分歧。证毕。
\end{proof}

\begin{theorem}[拉回核的三循环刻画与 \texorpdfstring{$J(C_6)$}{J(C6)} 的等变精细分解]\label{thm:xi-s3-closure-pullback-kernel-cyclic3-equivariant-splitting}
令 \(K=\QQ(\sqrt{-3})\)。
取结论 \ref{con:xi-terminal-zm-discriminant-ridge-kummer-divisor-3torsion} 的 \(3\)-挠点
\[
\tau\in J(C)(K)[3]\setminus\{0\},\qquad
G:=\langle\tau\rangle,
\]
并记 \(\pi:C_6\to C\) 为对应 Kummer 数据 \(t^3=u\) 所定义的无分歧循环三重覆盖。
则：
\begin{enumerate}
  \item \(\ker(\pi^\ast:J(C)\to J(C_6))=G\)（作为 \(\bar\QQ\)-点子群，且对 \(\mathrm{Gal}(\bar\QQ/\QQ)\) 稳定）；
  \item \(\pi^\ast J(C)\sim_{\QQ}J(C)/G\)；
  \item 设 \(\iota\) 为任一换位提升，\(C_6/\langle\iota\rangle\simeq\overline{\mathcal R}\)，则
  \[
  J(C_6)\sim_{\QQ}(J(C)/G)\times J(\overline{\mathcal R})^2.
  \]
\end{enumerate}
\end{theorem}

\begin{proof}
由 Kummer 覆盖结构，存在阶 \(3\) 线丛 \(L\in\operatorname{Pic}^0(C)\)（其类对应 \(\tau\)）使
\[
\pi_*\mathcal O_{C_6}\cong \mathcal O_C\oplus L^{-1}\oplus L^{-2},
\]
且在 \(C_6\) 上 \(L\) 被立方根平凡化，故 \(G=\langle L\rangle\subseteq\ker(\pi^\ast)\)。

反向包含由循环覆盖下降理论给出：满足 \(\pi^\ast M\simeq\mathcal O_{C_6}\) 的 \(M\in\operatorname{Pic}^0(C)\)
与 \(C_3\)-特征标同构，
\[
\operatorname{Hom}(C_3,\mathbf G_m)\cong\ZZ/3\ZZ,
\]
其生成元正是 \(L\)。故 \(\ker(\pi^\ast)=\langle L\rangle=G\)，即得 (1) 与 (2)。

再由结论 \ref{con:xi-s3-closure-jacobian-er2-refinement}
\[
J(C_6)\sim_{\QQ}J(C)\times J(\overline{\mathcal R})^2.
\]
将第一因子以同源 \(\varphi:J(C)\to J(C)/G\)（结论 \ref{con:xi-terminal-zm-discriminant-ridge-3isogeny-descent}）替换，即得
\[
J(C_6)\sim_{\QQ}(J(C)/G)\times J(\overline{\mathcal R})^2.
\]
证毕。
\end{proof}

\begin{corollary}[局部 \texorpdfstring{$\zeta$}{zeta} 因子分解与点数恒等式的等价重写]\label{cor:xi-s3-closure-local-zeta-factorization-c-r-square}
对任意同时使 \(C,\overline{\mathcal R},C_6\) 良约化的素数 \(p\)，定义
\[
P_{X,p}(T):=\det\!\left(1-T\mathrm{Frob}_p\mid H^1_{\acute et}(X_{\overline{\FF}_p},\QQ_\ell)\right),
\qquad X\in\{C,\overline{\mathcal R},C_6\}.
\]
则
\begin{equation}\label{eq:xi-c6-local-zeta-factorization-cr2}
P_{C_6,p}(T)=P_{C,p}(T)\,P_{\overline{\mathcal R},p}(T)^2.
\end{equation}
等价地，对任意 \(n\ge 1\)：
\begin{equation}\label{eq:xi-c6-point-count-identity-cr2}
\#C_6(\FF_{p^n})
=\#C(\FF_{p^n})+2\,\#\overline{\mathcal R}(\FF_{p^n})-2(p^n+1).
\end{equation}
\end{corollary}

\begin{proof}
由定理 \ref{thm:xi-s3-closure-pullback-kernel-cyclic3-equivariant-splitting} 的同源分解，
在好素数处 \(H^1\) 的 \(\ell\)-进表示按直和分解，得到
\eqref{eq:xi-c6-local-zeta-factorization-cr2}。
再用
\[
\#X(\FF_{p^n})=p^n+1-\operatorname{Tr}\!\left(\mathrm{Frob}_p^n\mid H^1_{\acute et}(X_{\overline{\FF}_p},\QQ_\ell)\right)
\]
逐项取迹即可得到 \eqref{eq:xi-c6-point-count-identity-cr2}。证毕。
\end{proof}

\begin{proposition}[标准等型因子的显式三坐标实现与 \texorpdfstring{$S_3$}{S3} 矩阵表示]\label{prop:xi-s3-closure-standard-isotypic-explicit-r3-realization}
\leanverified{paper\_xi\_s3\_closure\_standard\_isotypic\_explicit\_r3\_realization}
令 \(H_1,H_2,H_3\subset S_3\) 为三个位形互换的换位子群，并固定商映射
\[
f_i:C_6\to C_6/H_i\simeq \overline{\mathcal R},\qquad i=1,2,3.
\]
定义
\[
\Psi:J(C_6)\longrightarrow J(\overline{\mathcal R})^3,\qquad
D\longmapsto \bigl((f_1)_*D,\ (f_2)_*D,\ (f_3)_*D\bigr).
\]
则：
\begin{enumerate}
  \item \(\operatorname{Im}(\Psi)\subseteq W\)，其中
  \[
  W:=\{(P_1,P_2,P_3)\in J(\overline{\mathcal R})^3:\ P_1+P_2+P_3=0\};
  \]
  并且 \(\Psi\) 诱导
  \[
  J(C_6)/\pi^\ast J(C)\ \sim_{\QQ}\ W\ \sim_{\QQ}\ J(\overline{\mathcal R})^2.
  \]
  \item \(S_3\) 通过三坐标置换作用于 \(J(\overline{\mathcal R})^3\) 并稳定 \(W\)；
  在基
  \[
  v_1=(1,-1,0),\qquad v_2=(0,1,-1)
  \]
  下，若约定
  \(
  \sigma\cdot(P_1,P_2,P_3)=(P_{\sigma^{-1}(1)},P_{\sigma^{-1}(2)},P_{\sigma^{-1}(3)})
  \)，则
  \[
  (123)\longleftrightarrow
  \begin{pmatrix}
  0 & -1\\
  1 & -1
  \end{pmatrix},
  \qquad
  (12)\longleftrightarrow
  \begin{pmatrix}
  -1 & 1\\
  0 & 1
  \end{pmatrix}.
  \]
  从而 \((123)\) 的最小多项式为 \(x^2+x+1\)，且在 \(W\) 上无不动向量。
\end{enumerate}
\end{proposition}

\begin{proof}
每个 \(f_i\) 对应换位子群平均投影，\(\Psi\) 因而 \(S_3\)-等变。
三者求和对应平凡表示方向；而 \(C_6/S_3\simeq\PP^1\)（亏格 \(0\)）意味着 Jacobian 的平凡等型为零，
故
\[
(f_1)_*D+(f_2)_*D+(f_3)_*D=0,
\]
即 \(\operatorname{Im}(\Psi)\subseteq W\)。

由定理 \ref{thm:xi-s3-closure-pullback-kernel-cyclic3-equivariant-splitting}，
\(
J(C_6)/\pi^\ast J(C)
\)
是标准等型二维阿贝尔部分；而 \(W\) 正是置换表示中的标准子空间，
两者维数同为 \(2\) 且 \(\Psi\) 非零、等变，故诱导 \(\QQ\)-同源。

矩阵由坐标置换直接计算：
\[
(123)\cdot v_1=v_2,\qquad
(123)\cdot v_2=-v_1-v_2,
\]
\[
(12)\cdot v_1=-v_1,\qquad
(12)\cdot v_2=v_1+v_2,
\]
即得所示矩阵。
其中 \((123)\) 的特征多项式 \(x^2+x+1\) 无根 \(1\)，故无不动向量。证毕。
\end{proof}

\input{thm__xi-terminal-zm-s3-endoscopic-motive-polarization-endomorphism-rigidity}

\endinput

\paragraph{二面体商椭圆化、纤维积分歧奇偶律与 Prym 分解的闭式化}
沿用结论 \ref{con:xi-terminal-zm-resolvent-cubic-plane-cubic-geometry}、
定理 \ref{con:xi-resolvent-cubic-short-weierstrass-12096}、
定理 \ref{thm:xi-s3-closure-double-cover-explicit-branch-divisor}
与定理 \ref{thm:xi-terminal-zm-s3-closure-jacobian-splitting} 的记号。

\begin{theorem}[二面体商预解式曲线的显式 Weierstrass 化与双有理对应]\label{thm:xi-resolvent-dihedral-elliptic-birational}
设
\[
R(y,z):=z^{3}+(2y+1)z^{2}-(4y^{2}+4y+1)z-(8y^{3}+13y^{2}+5y+1)=0,
\]
并记其射影光滑模型为 \(C_{\mathrm{res}}\)。
取有理点 \(P=(y,z)=(0,1)\in C_{\mathrm{res}}(\QQ)\)。
定义
\[
m:=\frac{z-1}{y},\qquad
A(m):=m^{3}+2m^{2}-4m-8,\qquad
B(m):=4m^{2}-17,
\]
\[
u:=2A(m)\,y+B(m).
\]
则在 \(C_{\mathrm{res}}\) 的函数域中恒有
\[
u^{2}=-4m^{3}-16m^{2}+16m+65.
\]
从而得到支配有理映射
\[
C_{\mathrm{res}}\dashrightarrow E_{\mathrm D},\qquad
E_{\mathrm D}:\ u^{2}=-4m^{3}-16m^{2}+16m+65.
\]
其逆映射在开集 \(A(m)\neq 0\) 上可写为
\[
y=\frac{u-B(m)}{2A(m)},\qquad z=1+my.
\]
因此 \(C_{\mathrm{res}}\) 与 \(E_{\mathrm D}\) 在 \(\QQ\) 上双有理等价。
\end{theorem}
\begin{proof}
过 \(P\) 的直线族写作 \(z=1+my\)。
代入 \(R(y,z)=0\) 得
\[
R(y,1+my)=y\Bigl(A(m)y^{2}+B(m)y+C(m)\Bigr),\qquad C(m):=4m-7.
\]
去除因子 \(y\) 后得到二次方程
\[
A(m)y^{2}+B(m)y+C(m)=0.
\]
其判别式为
\[
\Delta_m:=B(m)^2-4A(m)C(m)=-4m^{3}-16m^{2}+16m+65.
\]
由根公式立得
\[
\bigl(2A(m)y+B(m)\bigr)^2=\Delta_m.
\]
按 \(u:=2A(m)y+B(m)\) 即得
\[
u^{2}=-4m^{3}-16m^{2}+16m+65.
\]
反向若给定 \((m,u)\in E_{\mathrm D}\) 且 \(A(m)\neq 0\)，取
\(
y=(u-B(m))/(2A(m))
\) 并置 \(z=1+my\)，可直接代回验证 \(R(y,z)=0\)。
两映射在共同定义域上互为逆，故为双有理等价。证毕。
\end{proof}

\input{subsubsec__xi-ed-cubic-y-branchfield-leyang}

\begin{corollary}[\texorpdfstring{$E_{\mathrm D}$}{ED} 的判别式与 \texorpdfstring{$j$}{j}-不变量闭式]\label{cor:xi-ed-discriminant-j-closed-form}
\leanverified{paper\_xi\_ed\_discriminant\_j\_closed\_form}
在定理 \ref{thm:xi-resolvent-dihedral-elliptic-birational} 的记号下，
令 \(m=x-\frac43\)，则 \(E_{\mathrm D}\) 化为
\[
u^{2}=4x^{3}-\frac{112}{3}x-\frac{667}{27}
\]
（即 \(\wp\)-形式 \(u^2=4x^3-g_2x-g_3\)）。
因此
\[
g_2=\frac{112}{3},\qquad g_3=\frac{667}{27},
\]
\[
\Delta(E_{\mathrm D})=g_2^3-27g_3^2=31^2\cdot 37=35557,
\]
\[
j(E_{\mathrm D})
=1728\,\frac{g_2^3}{\Delta(E_{\mathrm D})}
=\frac{2^{18}\cdot 7^3}{31^2\cdot 37}
=\frac{89915392}{35557}.
\]
\end{corollary}
\begin{proof}
代换 \(m=x-\frac43\) 后直接展开并配方即可得到 \(g_2,g_3\)。
其余为 Weierstrass 不变量标准公式。
该闭式与定理 \ref{con:xi-resolvent-cubic-short-weierstrass-12096} 完全一致。证毕。
\end{proof}

\begin{theorem}[Lee--Yang 三次判别式与 \texorpdfstring{$E_{\mathrm D}$}{ED} 判别式同步]\label{thm:xi-leyang-discriminant-sync-ed}
\leanverified{paper\_xi\_leyang\_discriminant\_sync\_ed}
令
\[
g(y):=256y^{3}+411y^{2}+165y+32.
\]
则有恒等式
\[
\operatorname{Disc}(g)=-3^{9}\cdot \Delta(E_{\mathrm D})
=-3^{9}\cdot 31^{2}\cdot 37.
\]
\end{theorem}
\begin{proof}
由结论 \ref{con:xi-terminal-zm-leyang-cubic-arithmetic}，
\(
\operatorname{Disc}(g)=-3^9\cdot 31^2\cdot 37
\)。
由推论 \ref{cor:xi-ed-discriminant-j-closed-form}，
\(
\Delta(E_{\mathrm D})=31^2\cdot 37
\)。
合并即得结论。证毕。
\end{proof}

\begin{theorem}[\texorpdfstring{$E_{\mathrm D}(\QQ)$}{ED(Q)} 挠子群平凡且秩为正]\label{thm:xi-ed-mordell-weil-torsion-trivial-rank-positive}
设
\[
E_{\mathrm D}':\ v^{2}=4x^{3}-16x^{2}-16x+65.
\]
则 \(E_{\mathrm D}'\) 与 \(E_{\mathrm D}\) 在 \(\QQ\) 上同构，且
\[
E_{\mathrm D}(\QQ)_{\mathrm{tors}}=0,\qquad
\operatorname{rank}E_{\mathrm D}(\QQ)\ge 1.
\]
\end{theorem}
\begin{proof}
对好素数 \(p\notin\{31,37\}\)，约化同态给出单射
\[
E_{\mathrm D}(\QQ)_{\mathrm{tors}}\hookrightarrow E_{\mathrm D}(\FF_p).
\]
取 \(p=5,7,11\)，有
\[
\#E_{\mathrm D}(\FF_5)=8,\qquad
\#E_{\mathrm D}(\FF_7)=13,\qquad
\#E_{\mathrm D}(\FF_{11})=15.
\]
故
\[
\gcd(8,13,15)=1,
\]
从而 \(E_{\mathrm D}(\QQ)_{\mathrm{tors}}\) 只能为平凡群。
另一方面，
\[
P:=(-2,1)\in E_{\mathrm D}'(\QQ)
\]
满足方程，因挠子群平凡，\(P\) 为无穷阶点，故秩至少为 \(1\)。证毕。
\end{proof}

\begin{theorem}[\texorpdfstring{$S_3$}{S3} 闭包的纤维积刻画与二次分支奇偶律]\label{thm:xi-s3-fiberproduct-branch-parity-xa-xv-ed}
记
\[
\Delta(y):=-y(y-1)g(y),\qquad
X_A:\ t^2=\Delta(y),
\]
并令 \(E_{\mathrm D}\to \PP^1_y\) 为由 \(R(y,z)=0\) 诱导的三次覆盖。
定义
\[
X_V:=\operatorname{Norm}\!\bigl(X_A\times_{\PP^1_y}E_{\mathrm D}\bigr).
\]
则：
\begin{enumerate}
  \item \(X_V\to X_A\) 为无分歧三次循环覆盖；
  \item 自然投影
  \[
  q:\ X_V\longrightarrow E_{\mathrm D}
  \]
  为二次覆盖，并且其分支点恰为满足
  \[
  \operatorname{ord}_Q\!\bigl(\Delta(y)\bigr)\equiv 1\pmod 2
  \]
  的 \(Q\in E_{\mathrm D}\)；
  \item 等价地，分支支撑正是六个分歧值
  \(
  y\in\{0,1,\infty\}\cup\{g(y)=0\}
  \)
  的各自未分歧原像点，故
  \[
  \deg \operatorname{Branch}(q)=6.
  \]
\end{enumerate}
\end{theorem}
\begin{proof}
函数域写作
\[
\QQ(X_A)=\QQ(y,t),\ t^2=\Delta(y),\qquad
\QQ(E_{\mathrm D})=\QQ(y,z),\ R(y,z)=0,
\]
\[
\QQ(X_V)=\QQ(y,z,t).
\]
由结论 \ref{con:xi-terminal-zm-s4-normal-subgroup-quotients-explicit-models-genus} 可知
\(\QQ(X_V)\) 正是该纤维积规范化的函数域。
在 \(\QQ(X_A)\) 上，三次预解式的判别式 \(\Delta(y)\) 已成为平方，故三次扩张的闭包群降为 \(C_3\)；
结合驯分歧基变换判据，得到 \(X_V\to X_A\) 为无分歧循环三次覆盖。
\par
对 \(q\) 而言，
\[
\QQ(E_{\mathrm D})\subset \QQ(E_{\mathrm D})\bigl(\sqrt{\Delta(y)}\bigr)
\]
为二次扩张，分歧判据即
\(
\operatorname{ord}_Q(\Delta(y))
\)
奇偶性。
又 \(E_{\mathrm D}\to\PP^1_y\) 在每个分歧值纤维型为 \((2,1)\)。
若 \(Q_{\mathrm{ram}}\) 为三次覆盖的分歧点，则 \(\operatorname{ord}_{Q_{\mathrm{ram}}}(y-y_0)=2\)；
若 \(Q_{\mathrm{unr}}\) 为未分歧点，则 \(\operatorname{ord}_{Q_{\mathrm{unr}}}(y-y_0)=1\)。
因 \(\Delta(y)\) 在每个有限分歧值处均是一阶零（在 \(y=\infty\) 处为五阶极点），得到
\[
\operatorname{ord}_{Q_{\mathrm{ram}}}\Delta(y)\equiv 0\pmod2,\qquad
\operatorname{ord}_{Q_{\mathrm{unr}}}\Delta(y)\equiv 1\pmod2.
\]
于是每个分歧值恰贡献一个分支点，共六点。该结论与定理
\ref{thm:xi-s3-closure-double-cover-explicit-branch-divisor}
中的显式分支除子一致。证毕。
\end{proof}

\begin{theorem}[分支二次覆盖的 Prym 分解]\label{thm:xi-prym-q-decomposition-jacxa-ed}
在定理 \ref{thm:xi-s3-fiberproduct-branch-parity-xa-xv-ed} 的设定下，
记 \(q:X_V\to E_{\mathrm D}\)。
则存在 \(\QQ\)-同源
\[
\operatorname{Prym}(q)\ \sim_{\QQ}\ \operatorname{Jac}(X_A)\times E_{\mathrm D}.
\]
\end{theorem}
\begin{proof}
对二次覆盖 \(q:C\to D\) 有标准同源分解
\[
\operatorname{Jac}(C)\sim_{\QQ}\operatorname{Jac}(D)\times \operatorname{Prym}(q).
\]
取 \(C=X_V,\ D=E_{\mathrm D}\)，得
\[
\operatorname{Jac}(X_V)\sim_{\QQ}E_{\mathrm D}\times \operatorname{Prym}(q).
\]
另一方面，由定理 \ref{thm:xi-terminal-zm-s3-closure-jacobian-splitting}
在同一对象识别 \(X_V=C_6,\ X_A=C,\ E_{\mathrm D}\sim \overline{\mathcal R}\) 下可写为
\[
\operatorname{Jac}(X_V)\sim_{\QQ}\operatorname{Jac}(X_A)\times E_{\mathrm D}^{2}.
\]
比较两式并消去一个 \(E_{\mathrm D}\) 因子即得结论。证毕。
\end{proof}

\begin{theorem}[\texorpdfstring{$\operatorname{Jac}(X_A)$}{Jac(XA)} 在 \(\QQ\) 上不可分解为椭圆曲线乘积]\label{thm:xi-jac-xa-q-indecomposable-via-p13}
不存在椭圆曲线 \(E_1,E_2/\QQ\) 使
\[
\operatorname{Jac}(X_A)\sim_{\QQ}E_1\times E_2.
\]
\end{theorem}
\begin{proof}
取好素数 \(p=13\) 并记
\[
N_1:=\#X_A(\FF_{13})=14,\qquad N_2:=\#X_A(\FF_{13^2})=166.
\]
设局部因子
\[
P_{13}(T)=1-a_1T+a_2T^2-13a_1T^3+13^2T^4.
\]
由标准关系
\[
a_1=13+1-N_1=0,\qquad
S_2:=13^2+1-N_2=4,\qquad
a_2=\frac{a_1^2-S_2}{2}=-2,
\]
得到
\[
P_{13}(T)=169T^4-2T^2+1.
\]
若其在 \(\QQ[T]\) 可约，仅可能分解为两二次因子。
偶分解情形对应 \(169U^2-2U+1\)（\(U=T^2\)）在 \(\QQ[U]\) 可约；但其判别式
\[
(-2)^2-4\cdot169\cdot 1=-672
\]
非平方，故不可能。
含奇次项相消的情形只能写作
\[
(13T^2+aT+1)(13T^2-aT+1),
\]
比较 \(T^2\) 系数得 \(26-a^2=-2\)，即 \(a^2=28\notin\QQ^{\times2}\)，亦不可能。
故 \(P_{13}(T)\) 在 \(\QQ[T]\) 不可约。
\par
若 \(\operatorname{Jac}(X_A)\sim_{\QQ}E_1\times E_2\)，则
\[
P_{13}(T)=\prod_{i=1}^{2}\bigl(1-a_{13}(E_i)T+13T^2\bigr)
\in \ZZ[T]
\]
必可约，矛盾。
因此 \(\operatorname{Jac}(X_A)\) 在 \(\QQ\)-同源意义下不可分解为椭圆曲线乘积。证毕。
\end{proof}

\endinput

\paragraph{Lee--Yang 判别式的 Cardano--Kummer 规范形与三次线性化}
在本段中固定
\[
H_{\mathrm{LY}}:\quad u^2=-y(y-1)P(y),\qquad P(y)=256y^3+411y^2+165y+32,
\]
并记
\[
A(y):=128y^3+219y^2+69y+16\in\ZZ[y].
\]

\begin{theorem}[Lee--Yang 判别式的六次幂塌缩恒等式]\label{thm:xi-leyang-sextic-perfect-power-collapse}
\leanverified{paper\_xi\_leyang\_sextic\_perfect\_power\_collapse}
在 \(\ZZ[y]\) 中成立整式恒等式
\[
A(y)^2+27\bigl(-y(y-1)P(y)\bigr)=256(2y+1)^6.
\]
等价地，在 \(\QQ(H_{\mathrm{LY}})\) 中有
\[
A(y)^2+27u^2=256(2y+1)^6.
\]
\end{theorem}
\begin{proof}
直接展开并比较系数即可得到第一式。
将 \(u^2=-y(y-1)P(y)\) 代入即得第二式。证毕。
\end{proof}

\begin{proposition}[Cardano 分子的主除子与有限零点集中]\label{prop:xi-leyang-cardano-numerator-divisor-visibility}
令 \(K:=\QQ(\sqrt{-3})\)，记 \(s:=\sqrt{-3}\)。
在 \(K(H_{\mathrm{LY}})\) 中定义
\[
N_-(y,u):=A(y)-3su,\qquad N_+(y,u):=A(y)+3su.
\]
则
\[
N_-(y,u)\,N_+(y,u)=256(2y+1)^6.
\]
记 \(\infty\) 为 \(H_{\mathrm{LY}}\) 的唯一无穷远点，并定义
\[
Q_-:=\left(y=-\frac12,\ u=-\frac94 s\right),\qquad
Q_+:=\left(y=-\frac12,\ u=\frac94 s\right)\in H_{\mathrm{LY}}(K).
\]
则主除子满足
\[
\operatorname{div}(N_-)=6Q_--6\infty,\qquad
\operatorname{div}(N_+)=6Q_+-6\infty.
\]
特别地，\(N_-\)（分别 \(N_+\)）在有限处仅于 \(Q_-\)（分别 \(Q_+\)）有零点，且零阶均为 \(6\)。
\end{proposition}
\begin{proof}
由定理 \ref{thm:xi-leyang-sextic-perfect-power-collapse} 立即得到乘积恒等式。
在 \(y=-\frac12\) 处，
\[
P\!\left(-\frac12\right)=\frac{81}{4},\qquad
u^2=-y(y-1)P(y)\big|_{y=-1/2}=-\frac{243}{16}=\left(\frac94\right)^2(-3),
\]
故该纤维恰为两点 \(Q_\pm\)。
又
\(
A\!\left(-\frac12\right)=\frac{81}{4}
\)，
代入即得
\[
N_-(Q_-)=0,\ N_+(Q_-)\neq 0;\qquad
N_+(Q_+)=0,\ N_-(Q_+)\neq 0.
\]
由于 \(u(Q_\pm)\neq 0\)，两点均为非分歧点，局部参数可取 \(t=y+\frac12\)，从而
\[
\operatorname{ord}_{Q_\pm}(2y+1)=1,\qquad
\operatorname{ord}_{Q_\pm}\bigl((2y+1)^6\bigr)=6.
\]
由乘积恒等式与非零因子可得
\[
\operatorname{ord}_{Q_-}(N_-)=6,\ \operatorname{ord}_{Q_-}(N_+)=0,\qquad
\operatorname{ord}_{Q_+}(N_+)=6,\ \operatorname{ord}_{Q_+}(N_-)=0.
\]
且 \((2y+1)^6\) 在其他有限点无零，故 \(N_\pm\) 无其他有限零点。
在无穷远点，奇数次数模型给出
\[
\operatorname{ord}_\infty(y)=-2,\qquad \operatorname{ord}_\infty(u)=-5.
\]
故
\(
\operatorname{ord}_\infty(A(y))=-6
\) 且
\(
\operatorname{ord}_\infty(3su)=-5
\)，从而
\(
\operatorname{ord}_\infty(N_\pm)=-6
\)。
综上得到所示主除子。证毕。
\end{proof}

\begin{theorem}[范数一单位与显式 \texorpdfstring{$3$}{3}-挠点证书]\label{thm:xi-leyang-norm-one-unit-explicit-3torsion}
在 \(K(H_{\mathrm{LY}})\) 中定义
\[
\eta:=\frac{A(y)-3su}{16(2y+1)^3}.
\]
记超椭圆对合 \(\iota:(y,u)\mapsto (y,-u)\)。
则
\[
\iota(\eta)=\eta^{-1},\qquad \eta\,\iota(\eta)=1.
\]
并且
\[
\operatorname{div}(\eta)=3(Q_--Q_+).
\]
因此
\[
\tau:=[Q_--Q_+]\in J(H_{\mathrm{LY}})(K)
\]
为显式非平凡 \(3\)-挠点，即 \(3\tau=0\) 且 \(\tau\neq 0\)。
\end{theorem}
\begin{proof}
由命题 \ref{prop:xi-leyang-cardano-numerator-divisor-visibility} 的乘积恒等式，
\[
(A-3su)(A+3su)=256(2y+1)^6,
\]
故
\[
\eta\,\iota(\eta)=
\frac{(A-3su)(A+3su)}{256(2y+1)^6}=1.
\]
再由命题 \ref{prop:xi-leyang-cardano-numerator-divisor-visibility}
\[
\operatorname{div}(A-3su)=6Q_--6\infty,
\]
以及
\[
\operatorname{div}(2y+1)=Q_-+Q_+-2\infty
\]
可得
\[
\operatorname{div}\!\bigl((2y+1)^3\bigr)=3Q_-+3Q_+-6\infty.
\]
因此
\[
\operatorname{div}(\eta)
=\bigl(6Q_--6\infty\bigr)-\bigl(3Q_-+3Q_+-6\infty\bigr)
=3(Q_--Q_+).
\]
故 \(3[Q_--Q_+]=0\)。
若 \([Q_--Q_+]=0\)，则存在除子为 \(Q_--Q_+\) 的函数，给出 \(H_{\mathrm{LY}}\to\PP^1\) 的次数 \(1\) 映射，
迫使曲线有亏格 \(0\)，与 \(g(H_{\mathrm{LY}})=2\) 矛盾。
故 \(\tau\neq 0\)，其阶即为 \(3\)。证毕。
\end{proof}

\begin{theorem}[循环三次无分歧覆盖的 Kummer 规范形]\label{thm:xi-leyang-kummer-c3-etale-normal-form}
在 \(K=\QQ(\sqrt{-3})\) 上考虑扩张
\[
K(H_{\mathrm{LY}})\subset K(H_{\mathrm{LY}})(\xi),\qquad \xi^3=\eta,
\]
其中 \(\eta\) 见定理 \ref{thm:xi-leyang-norm-one-unit-explicit-3torsion}。
则该扩张对应一个 \(K\)-上循环 \(C_3\) 无分歧覆盖
\[
\pi_\eta:\ C_\eta\longrightarrow H_{\mathrm{LY}}.
\]
进一步，设
\[
T:=\operatorname{Res}^1_{K/\QQ}\mathbf G_m,
\]
则 \(\eta\,\iota(\eta)=1\) 给出态射
\[
\varphi_\eta:\ H_{\mathrm{LY}}\longrightarrow T,
\]
且覆盖 \(\pi_\eta\) 等同于将 \(T\) 上乘三同态
\[
[3]:T\to T,\qquad t\mapsto t^3
\]
沿 \(\varphi_\eta\) 拉回所得的 \(\mu_3\)-扭量。
\end{theorem}
\begin{proof}
由定理 \ref{thm:xi-leyang-norm-one-unit-explicit-3torsion}，
\(
\operatorname{div}(\eta)=3(Q_--Q_+)
\)，
故任意离散赋值 \(v\) 满足 \(v(\eta)\equiv 0\pmod 3\)。
且 \(K\) 含 \(\mu_3\)，故 Kummer 理论适用：
\(
\xi^3=\eta
\)
定义的扩张为循环三次，且由赋值可整除性知处处无分歧。
\par
又因 \(\eta\,\iota(\eta)=1\)，\(\eta\) 落在范数一子群，因而定义到
\(
T=\operatorname{Res}^1_{K/\QQ}\mathbf G_m
\)
的态射 \(\varphi_\eta\)。
在函数域层面，沿 \([3]\) 的拉回正对应“加入 \(\xi\) 且满足 \(\xi^3=\eta\)”；
故覆盖描述与扭量描述一致。证毕。
\end{proof}

\begin{theorem}[Cardano--Chebyshev 线性化与三次分解根的迹表达]\label{thm:xi-leyang-cardano-chebyshev-linearization-z}
沿定理 \ref{thm:xi-leyang-kummer-c3-etale-normal-form}，令
\[
s:=\xi+\xi^{-1}\in K(C_\eta).
\]
则有恒等式
\[
\eta+\eta^{-1}=s^3-3s.
\]
并且在 \(K(C_\eta)\) 中，变量 \(z\) 由
\[
R(z,y)=z^3+(1+2y)z^2-(1+4y+4y^2)z-(1+5y+13y^2+8y^3)=0
\]
刻画时，满足显式表达
\[
z=\frac{2y+1}{3}\,(2s-1)=\frac{2y+1}{3}\Bigl(2(\xi+\xi^{-1})-1\Bigr).
\]
\end{theorem}
\begin{proof}
由 \(\eta=\xi^3\) 得
\[
\eta+\eta^{-1}=\xi^3+\xi^{-3}=(\xi+\xi^{-1})^3-3(\xi+\xi^{-1})=s^3-3s.
\]
另一方面，令
\[
s=\frac12\left(1+\frac{3z}{2y+1}\right)
\quad\Longleftrightarrow\quad
z=\frac{2y+1}{3}(2s-1).
\]
直接代入可得恒等式
\[
s^3-3s-\frac{A(y)}{8(2y+1)^3}
=\frac{27}{8(2y+1)^3}\,R(z,y).
\]
故 \(R(z,y)=0\) 当且仅当
\[
s^3-3s=\frac{A(y)}{8(2y+1)^3}.
\]
又
\[
\eta+\eta^{-1}
=\frac{A(y)-3su}{16(2y+1)^3}+\frac{A(y)+3su}{16(2y+1)^3}
=\frac{A(y)}{8(2y+1)^3},
\]
于是得到所述 \(z\) 的表达式。证毕。
\end{proof}

\begin{corollary}[分解三次曲线的通用 Chebyshev 规范形与判别式回收]\label{cor:xi-leyang-chebyshev-normal-form-and-discriminant-recovery}
定义
\[
t(y):=\frac{A(y)}{8(2y+1)^3}\in\QQ(y).
\]
则在双有理变换
\[
s=\frac12\left(1+\frac{3z}{2y+1}\right)
\]
下，三次分解曲线 \(R(z,y)=0\) 等价于
\[
s^3-3s=t(y).
\]
其关于 \(s\) 的判别式为
\[
\operatorname{Disc}_s\bigl(s^3-3s-t(y)\bigr)=27\bigl(4-t(y)^2\bigr)
=\frac{3^6}{2^6}\cdot\frac{u^2}{(2y+1)^6}.
\]
从而二次分辨率子域由
\[
\sqrt{4-t(y)^2}\sim \frac{u}{(2y+1)^3}
\]
生成，等同于 \(\QQ(y,u)=\QQ(H_{\mathrm{LY}})\)。
\end{corollary}
\begin{proof}
第一式由定理 \ref{thm:xi-leyang-cardano-chebyshev-linearization-z} 直接得到。
对三次 \(x^3+px+q\) 的判别式公式
\(
\operatorname{Disc}= -4p^3-27q^2
\)
代入 \(p=-3,q=-t(y)\) 得
\[
\operatorname{Disc}_s=27(4-t(y)^2).
\]
由定理 \ref{thm:xi-leyang-sextic-perfect-power-collapse}
\[
A(y)^2+27u^2=256(2y+1)^6
\]
可得
\[
4-t(y)^2
=\frac{256(2y+1)^6-A(y)^2}{64(2y+1)^6}
=\frac{27u^2}{64(2y+1)^6},
\]
从而得到所述判别式与二次分辨率表达。证毕。
\end{proof}

\begin{theorem}[\texorpdfstring{$t(y)$}{t(y)} 的三分支分歧型与显式临界值]\label{thm:xi-leyang-ty-three-branch-profile}
将
\[
t(y)=\frac{A(y)}{8(2y+1)^3}
\]
视为映射 \(\PP^1_y\to\PP^1_t\)。
则其次数为 \(3\)，且分歧结构完全显式如下：
\begin{enumerate}
  \item 有限临界点由
  \[
  2y^2-6y+1=0
  \quad\Longleftrightarrow\quad
  y=\frac{3\pm\sqrt7}{2},
  \]
  两点均为简单临界点，分歧指数为 \(2\)；
  \item \(y=-\frac12\) 为三重极点，映到 \(t=\infty\)，其分歧指数为 \(3\)；
  \item 两个有限分歧值为
  \[
  t_\pm=t\!\left(\frac{3\pm\sqrt7}{2}\right)=\frac{139}{72}\pm\frac{7\sqrt7}{144}.
  \]
\end{enumerate}
因此分歧值集合恰为
\[
\{\infty,t_+,t_-\},
\]
分歧型为
\[
\infty:(3),\qquad t_+:(2,1),\qquad t_-:(2,1),
\]
并满足 Riemann--Hurwitz 总分歧量
\[
(3-1)+(2-1)+(2-1)=4=2\cdot 3-2.
\]
\end{theorem}
\begin{proof}
直接求导得
\[
\frac{dt}{dy}
=-\frac{27}{8}\cdot\frac{2y^2-6y+1}{(2y+1)^4}.
\]
故有限临界点如上，且因其为单根，对应分歧指数 \(2\)。
点 \(y=-\frac12\) 为三重极点，故其像为 \(\infty\)，且局部次数为 \(3\)。
代入临界点即得 \(t_\pm\)。
分歧型与总分歧量由局部次数直接读出。证毕。
\end{proof}

\input{subsubsec__xi-time-protocol-conclusions-part34d-leyang-cardano-kummer-branching}
\input{subsubsec__xi-time-protocol-conclusions-part34d-kummer-s3-quotients-branch-spectrum}
\input{subsubsec__xi-time-protocol-conclusions-part34d-explicit-torsor-localfactor-cubicresidue}
\endinput

\paragraph{\texorpdfstring{$V_4$}{V4} 子覆盖的三点分支刚性、双双覆盖数据与 Prym 重构}
沿用结论 \ref{con:xi-terminal-zm-s4-hurwitz-subgroup-spectrum-cycletypes} 的记号，
记
\[
X:=\widetilde X,\qquad V_4:=V_4^{\mathrm{norm}}\triangleleft S_4,\qquad
C_6:=X/V_4,
\]
并令
\[
q:X\to C_6,\qquad \pi_6:C_6\to \PP^1_y
\]
分别为自然商映射。

\begin{theorem}[\texorpdfstring{$V_4$}{V4} 子覆盖的分支刚性与惯性全异性]\label{thm:xi-v4-subcover-branch-rigidity-inertia-distinct}
\leanverified{paper\_xi\_v4\_subcover\_branch\_rigidity\_inertia\_distinct}
设 \(Q_1,Q_2,Q_3\in C_6\) 为 \(\pi_6^{-1}(\infty)\) 的三点。
则：
\begin{enumerate}
  \item \(q\) 的分支点恰为 \(\{Q_1,Q_2,Q_3\}\)；
  \item 在三点处的惯性子群
  \[
  I_i:=\mathrm{Inertia}_{Q_i}(q)\le V_4
  \]
  均为阶 \(2\) 且两两不同，并满足
  \[
  \{I_1,I_2,I_3\}=\{H\le V_4:\ |H|=2\}.
  \]
\end{enumerate}
\end{theorem}
\begin{proof}
设 \(\varphi:X\to \PP^1_y\) 为 \(S_4\)-正则覆盖。
由惯性交叉判据
\[
I_P(q)=I_P(\varphi)\cap V_4
\]
可知：在五个换位型分歧值上，\(I_P(\varphi)\) 由换位生成，与 \(V_4\) 交平凡，故 \(q\) 在这些纤维上不分歧。
在 \(y=\infty\) 处，\(I_P(\varphi)\cong C_4=\langle\rho\rangle\)，且 \(\rho^2\in V_4\) 为双换位，故
\[
I_P(\varphi)\cap V_4=\langle \rho^2\rangle\cong C_2\neq 1.
\]
因此 \(q\) 仅在 \(\pi_6^{-1}(\infty)\) 上可能分歧。
\par
由结论 \ref{con:xi-terminal-zm-s4-hurwitz-subgroup-spectrum-cycletypes} 中 \(H=V_4\) 的行可知 \(\rho\) 在 \(S_4/V_4\) 上循环型为 \(2^3\)，故 \(\pi_6^{-1}(\infty)\) 恰有三点，且每点在 \(\pi_6\) 中分歧指数为 \(2\)。
再对 \(q\) 应用 Riemann--Hurwitz：
\[
2g(X)-2=|V_4|\,(2g(C_6)-2)+\sum_{P\in X}(e_P(q)-1).
\]
代入 \(g(X)=16,\ g(C_6)=4,\ |V_4|=4\)，得总分歧贡献为 \(6\)。
若基点 \(Q\in C_6\) 的惯性阶为 \(2\)，则其上方有 \(4/2=2\) 个点，每点贡献 \(1\)，故 \(Q\) 对总分歧贡献为 \(2\)。
于是恰需 \(3\) 个分支基点，故分支集正是 \(\{Q_1,Q_2,Q_3\}\)。
\par
记三点惯性为 \(I_i\)。
取任意阶 \(2\) 子群 \(H\le V_4\)，考虑双覆盖
\[
p_H:X/H\to C_6.
\]
在 \(Q_i\) 处分歧当且仅当 \(I_i\not\subseteq H\)，对阶 \(2\) 情形等价于 \(I_i\neq H\)。
故分歧点数
\[
r_H=\#\{i:\ I_i\neq H\}
\]
必须为偶数（双覆盖分歧点数为偶数），从而
\[
n_H:=\#\{i:\ I_i=H\}=3-r_H
\]
为奇数。
若某个 \(H\) 有 \(n_H=3\)，则其余两个阶 \(2\) 子群的 \(n_H=0\)，与“奇数”矛盾。
故每个 \(H\) 皆有 \(n_H=1\)，即三者恰各出现一次。证毕。
\end{proof}

\begin{corollary}[三条属 \texorpdfstring{$8$}{8} 的二次中间曲线与纤维积回收]\label{cor:xi-v4-intermediate-double-covers-genus8-fiberproduct}
\leanverified{paper\_xi\_v4\_intermediate\_double\_covers\_genus8\_fiberproduct}
令 \(H_1,H_2,H_3\le V_4\) 为三个非平凡阶 \(2\) 子群，并置
\[
D_i:=X/H_i,\qquad i=1,2,3.
\]
则：
\begin{enumerate}
  \item \(D_i\to C_6\) 为次数 \(2\) 覆盖，且恰在
  \[
  \{Q_1,Q_2,Q_3\}\setminus\{Q_i\}
  \]
  上分歧；并有 \(g(D_i)=8\)；
  \item 对任意 \(i\neq j\)，有
  \[
  \operatorname{Norm}\!\bigl(D_i\times_{C_6}D_j\bigr)\cong X.
  \]
\end{enumerate}
\end{corollary}
\begin{proof}
由定理 \ref{thm:xi-v4-subcover-branch-rigidity-inertia-distinct}，
\(\mathrm{Inertia}_{Q_i}(q)=H_i\)（重标号后可取此形）。
于是 \(D_i\to C_6\) 在 \(Q_\ell\) 处分歧当且仅当 \(H_\ell\neq H_i\)，即恰在另外两点分歧。
Riemann--Hurwitz 给出
\[
2g(D_i)-2=2(2g(C_6)-2)+2=2\cdot 6+2=14,
\]
故 \(g(D_i)=8\)。
\par
又对 \(i\neq j\)，有 \(\langle H_i,H_j\rangle=V_4\) 且 \(H_i\cap H_j=\{1\}\)。
故函数域满足
\[
\QQ(X)=\QQ(D_i)\,\QQ(D_j)\quad\text{于 }\QQ(C_6)\text{ 上的合成},
\]
等价于纤维积归一化同构为 \(X\)。证毕。
\end{proof}

\begin{theorem}[\texorpdfstring{$X\to C_6$}{X to C6} 的显式双双覆盖结构]\label{thm:xi-v4-bidouble-cover-explicit-linebundles}
在定理 \ref{thm:xi-v4-subcover-branch-rigidity-inertia-distinct} 的标号下，存在线丛
\[
L_1,L_2,L_3\in \mathrm{Pic}^1(C_6)
\]
使
\[
L_1^{\otimes2}\simeq \mathcal O_{C_6}(Q_2+Q_3),\quad
L_2^{\otimes2}\simeq \mathcal O_{C_6}(Q_1+Q_3),\quad
L_3^{\otimes2}\simeq \mathcal O_{C_6}(Q_1+Q_2),
\]
且
\[
q_*\mathcal O_X\simeq
\mathcal O_{C_6}\oplus L_1^{-1}\oplus L_2^{-1}\oplus L_3^{-1}.
\]
三条中间双覆盖 \(D_i\to C_6\) 分别由 \(L_i\) 给出，其分歧除子即 \(Q_j+Q_k\)（\(\{i,j,k\}=\{1,2,3\}\)）。
\end{theorem}
\begin{proof}
\((\ZZ/2\ZZ)^2\)-Galois 覆盖的标准分类给出：
给定三条角色线丛 \(L_i\) 与兼容关系
\(
2L_i\sim D_j+D_k
\)
可唯一组装双双覆盖，并满足推前分解
\(
q_*\mathcal O_X
\)
的角色直和表达。
由推论 \ref{cor:xi-v4-intermediate-double-covers-genus8-fiberproduct}，三条中间双覆盖分歧分别为 \(Q_j+Q_k\)，故可取上述 \(L_i\)；
其次数由 \(\deg(Q_j+Q_k)=2\) 立得 \(\deg L_i=1\)。证毕。
\end{proof}

\begin{theorem}[Jacobian 四因子分解与 Prym 三重对称]\label{thm:xi-v4-jacobian-prym-triple-decomposition}
令
\[
P_i:=\operatorname{Prym}(D_i/C_6),\qquad i=1,2,3.
\]
则：
\begin{enumerate}
  \item \(\dim P_i=4\)，并由双覆盖 \(D_i\to C_6\) 的两点分歧诱导标准 Prym 极化；
  \item 存在自然 \(\QQ\)-同种分解
  \[
  J(X)\sim_{\QQ}J(C_6)\times P_1\times P_2\times P_3;
  \]
  \item \(S_4/V_4\simeq S_3\) 对 \(\{H_1,H_2,H_3\}\) 的传递置换诱导
  \[
  (P_1,P_2,P_3)
  \]
  的传递置换，从而三者两两同构（因而两两同种）。
\end{enumerate}
\end{theorem}
\begin{proof}
由推论 \ref{cor:xi-v4-intermediate-double-covers-genus8-fiberproduct}，\(g(D_i)=8\) 与 \(g(C_6)=4\)，故
\(
\dim P_i=g(D_i)-g(C_6)=4
\)。
\par
对群 \(V_4\) 的群代数幂等元分解给出 \(J(X)\) 的 isotypic 分裂：平凡角色对应 \(J(C_6)\)，三个非平凡角色分别对应三个中间双覆盖的 Prym 因子，故得
\[
J(X)\sim_{\QQ}J(C_6)\times P_1\times P_2\times P_3.
\]
\par
又 \(V_4\triangleleft S_4\)，商群 \(S_4/V_4\simeq S_3\) 传递作用于三个非平凡阶 \(2\) 子群，故对相应三条双覆盖与 Prym 因子诱导传递置换。
于是 \(P_1,P_2,P_3\) 两两同构。证毕。
\end{proof}

\begin{theorem}[\texorpdfstring{$V_4$}{V4} 特征分量的正则微分均匀律]\label{thm:xi-v4-character-eigenspace-uniform-dimension-four}
\leanverified{paper\_xi\_v4\_character\_eigenspace\_uniform\_dimension\_four}
记
\[
V_4^\vee:=\mathrm{Hom}(V_4,\{\pm1\}).
\]
则
\[
H^0(X,\omega_X)=\bigoplus_{\chi\in V_4^\vee}H^0(X,\omega_X)_\chi,
\]
并满足
\[
\dim H^0(X,\omega_X)_\chi=4\qquad(\forall\,\chi\in V_4^\vee).
\]
\end{theorem}
\begin{proof}
由定理 \ref{thm:xi-v4-bidouble-cover-explicit-linebundles} 的角色推前分解，
\[
q_*\omega_X
\simeq
\omega_{C_6}\oplus\bigoplus_{\chi\neq 1}\bigl(\omega_{C_6}\otimes L_\chi\bigr),
\]
其中 \(L_\chi\in\mathrm{Pic}^1(C_6)\)。
故平凡角色分量
\[
h^0(C_6,\omega_{C_6})=g(C_6)=4.
\]
对任意非平凡 \(\chi\)，\(\deg(\omega_{C_6}\otimes L_\chi)=6+1=7\)。
由 Riemann--Roch，
\[
h^0(C_6,\omega_{C_6}\otimes L_\chi)-h^0(C_6,L_\chi^{-1})=7-4+1=4.
\]
又 \(\deg L_\chi^{-1}=-1\)，故 \(h^0(C_6,L_\chi^{-1})=0\)。
于是非平凡分量维数亦为 \(4\)。证毕。
\end{proof}

\begin{theorem}[由 \texorpdfstring{$(C_6,Q_1+Q_2+Q_3)$}{(C6,B)} 的函子性重构]\label{thm:xi-v4-reconstruction-from-c6-threepoints-prym-triple}
记
\[
B:=Q_1+Q_2+Q_3\in \mathrm{Div}(C_6).
\]
则可由 \((C_6,B)\) 进行函子性重构：
\begin{enumerate}
  \item 取三条双覆盖
  \[
  D_i\to C_6,\qquad \operatorname{Br}(D_i/C_6)=B-Q_i;
  \]
  \item 由这三条双覆盖的双双覆盖组装恢复 \(q:X\to C_6\)（同构意义下唯一）；
  \item 记 \(P:=P_i\)（由定理 \ref{thm:xi-v4-jacobian-prym-triple-decomposition}，不依赖 \(i\)），则
  \[
  J(X)\sim_{\QQ}J(C_6)\times P^{\,3},
  \]
  且该分解与 \(S_3\simeq S_4/V_4\) 的置换对称相容。
\end{enumerate}
\end{theorem}
\begin{proof}
第 1 步与第 2 步由双双覆盖分类与定理
\ref{thm:xi-v4-bidouble-cover-explicit-linebundles}
直接给出。
第 3 步由定理 \ref{thm:xi-v4-jacobian-prym-triple-decomposition} 立即得到。证毕。
\end{proof}

\endinput

\paragraph{倍乘映射诱导的重量模对应：四次生成元、判别式核与几何预算的统一刚性}
在固定椭圆曲线
\[
E:\quad Z^2=X^3-X+\tfrac14
\]
上定义
\[
y:=Z+X^2-\tfrac12\in K:=\QQ(E),\qquad L:=\QQ(y).
\]
并记
\[
\Pi(X,y):=X^4-X^3-(2y+1)X^2+X+y(y+1)=0.
\]

\begin{theorem}[二倍乘与三倍乘的四次模方程及首项扭点刚性]\label{thm:xi-elliptic-weight-modular-f2f3-leading-rigidity-closure}
\leanverified{paper\_xi\_elliptic\_weight\_modular\_f2f3\_leading\_rigidity\_closure}
记 \(y_n:=y\circ[n]\in K\)（\(n=2,3\)）。
则存在本原且在 \(\ZZ[y,Y_n]\) 中不可约的四次多项式
\[
F_n(y,Y_n)\in \ZZ[y,Y_n],\qquad F_n(y,y_n)=0.
\]
其中：
\begin{enumerate}
  \item \(F_2(y,Y_2)=\sum_{j=0}^4 f_j(y)Y_2^j\) 与 \(F_3(y,Y_3)=\sum_{j=0}^4 g_j(y)Y_3^j\) 的完整整系数展开可由显式证书直接给出；
  \item 最高次系数满足
  \[
  f_4(y)=P_2(y)^4,\qquad g_4(y)=P_3(y)^4,
  \]
  其中
  \[
  P_2(y)=16y^3-8y^2-4y+1,
  \]
  \[
  P_3(y)=81y^8-324y^7+297y^6+729y^5-216y^4-1395y^3-480y^2-120y+7.
  \]
\end{enumerate}
并且 \(P_2,P_3\) 分别刻画 \(E[2]\setminus\{O\}\)、\(E[3]\setminus\{O\}\) 在 \(y\)-投影下的像。
\end{theorem}
\begin{proof}
四次生成元的存在性、不可约性与本原性由
\(\Pi(X,y)\) 与倍乘有理表达的消元直接给出，见
定理 \ref{thm:xi-elliptic-weight-modular-f2f3-leading-rigidity-closure}。
\par
首项四次幂结构来自“扭点像消元 + 极点阶 \(4\)”机制：对 \(n=2\) 与 \(n=3\)，清分母后首项分别由
\(
P_2^4
\) 与
\(
P_3^4
\)
控制，见定理
\ref{thm:xi-elliptic-weight-modular-f2f3-leading-rigidity-closure}。
证毕。
\end{proof}

\begin{theorem}[判别式平方自由核的统一性与索引平方分解]\label{thm:xi-elliptic-weight-modular-discriminant-kernel-unified-closure}
\leanverified{paper\_xi\_elliptic\_weight\_modular\_discriminant\_kernel\_unified\_closure}
设
\[
C(y):=256y^3+411y^2+165y+32,\qquad \Delta(y):=-y(y-1)C(y).
\]
则关于主变量的判别式满足
\[
\mathrm{Disc}_{Y_2}(F_2)=\Delta(y)\,\Delta_{12}(y)^2\,\Delta_{26}(y)^2,
\]
\[
\mathrm{Disc}_{Y_3}(F_3)=\Delta(y)\,\Delta_{32}(y)^2\,\Delta_{66}(y)^2,
\]
其中 \(\Delta_{12},\Delta_{26},\Delta_{32},\Delta_{66}\in\ZZ[y]\)，且前三级多项式可显式写出。
换言之，\(F_2\) 与 \(F_3\) 的判别式具有同一平方自由分支核 \(\Delta(y)\)。
\end{theorem}
\begin{proof}
由
\[
\mathrm{Disc}_X\!\bigl(\Pi(X,y)\bigr)=\Delta(y)
\]
可得扩张 \(K/L\) 的平方自由分支核。
另一方面，\(y_2,y_3\) 均为 \(K/L\) 的原始元，故其最小多项式判别式与域判别式之比为索引平方；
于是平方自由部分必一致并等于 \(\Delta(y)\)。
将余因子分解即得上述平方项结构。
详式与分解结果见定理
\ref{thm:xi-elliptic-weight-modular-discriminant-kernel-unified-closure}。
证毕。
\end{proof}

\begin{theorem}[泛型伽罗瓦群与对应曲线奇点预算的联合刚性]\label{thm:xi-elliptic-weight-modular-galois-delta-joint-rigidity}
\leanverified{paper\_xi\_elliptic\_weight\_modular\_galois\_delta\_joint\_rigidity}
对 \(n=2,3\)，\(F_n(y,Y_n)\) 作为 \(\QQ(y)\) 上四次不可约多项式满足
\[
\mathrm{Gal}\!\bigl(F_n/\QQ(y)\bigr)\cong S_4.
\]
令 \(\mathcal D_n\subset \PP^1_y\times\PP^1_{Y_n}\) 由 \(F_n(y,Y_n)=0\) 定义，则
\[
\deg_y(F_n)=4n^2,\qquad \deg_{Y_n}(F_n)=4,
\]
\[
g_a(\mathcal D_n)=(4n^2-1)(4-1)=12n^2-3.
\]
又由 \(\QQ(\mathcal D_n)\cong \QQ(E)\) 可知其正规化亏格为 \(1\)，从而
\[
\sum_{P\in\mathrm{Sing}(\mathcal D_n)}\delta_P
=g_a(\mathcal D_n)-1
=12n^2-4.
\]
特别地，
\[
\sum\delta_P=
\begin{cases}
44,& n=2,\\
104,& n=3.
\end{cases}
\]
\end{theorem}
\begin{proof}
\(S_4\) 泛型伽罗瓦群结论由特化判别法与 resolvent 不可约性给出，见定理
\ref{thm:xi-elliptic-weight-modular-galois-delta-joint-rigidity}。
\par
双次数由 \(\deg(y)=4\) 与 \(\deg([n])=n^2\) 导出；函数域同构
\(
\QQ(\mathcal D_n)\cong \QQ(E)
\)
给出正规化亏格 \(1\)。
结合 \(\PP^1\times\PP^1\) 的算术亏格公式与亏格差恒等式，得到 \(\delta\)-总预算。
详见定理
\ref{thm:xi-elliptic-weight-modular-galois-delta-joint-rigidity}。
证毕。
\end{proof}

\endinput

\paragraph{边界单参数退化下的极限混合 Hodge 结构与 \texorpdfstring{$S_4$}{S4} 等型退化谱}
沿用结论 \ref{con:xi-terminal-zm-s4-hurwitz-subgroup-spectrum-cycletypes}、\ref{con:xi-terminal-zm-s4-chevalley-weil-explicit-decomposition} 与 \ref{con:xi-terminal-zm-s4-jacobian-group-algebra-prym-e3-b3} 的记号。
设
\[
X:=\widetilde X,\qquad
V:=H^{1,0}(X)\cong 2\varepsilon\oplus \rho_2\oplus \rho_3\oplus 3\rho_3'.
\]
考虑固定局部分歧型 \((4,2,2,2,2,2)\) 的标记 \(S_4\)-Galois 覆盖的连通 Hurwitz 分支 \(\mathcal H\)，取一条横截边界 \(\overline{\mathcal H}\setminus\mathcal H\) 的单参数半稳定退化
\[
\pi:\mathcal X\to \Delta,\qquad X_t:=\pi^{-1}(t)\ (t\neq 0),
\]
并记局部单调与其对数为
\[
T\in \mathrm{Sp}\!\bigl(H^1(X_t,\QQ)\bigr),\qquad N:=\log T,\qquad N^2=0.
\]

\begin{theorem}[边界单参数退化的消失圈 \texorpdfstring{$S_4$}{S4}-模为诱导表示]\label{thm:xi-s4-boundary-vanishing-cycle-induced-module}
\leanverified{paper\_xi\_s4\_boundary\_vanishing\_cycle\_induced\_module}
记
\[
W_0(N):=\operatorname{Im}(N)\subset H^1(X_t,\QQ).
\]
设两条发生碰撞的换位局部单值为 \(g_i,g_{i+1}\in S_4\)，并置
\[
H:=\langle g_i,g_{i+1}\rangle,\qquad
p:=g_ig_{i+1},\qquad
r:=\operatorname{ord}(p)\in\{1,2,3\}.
\]
则存在 \(H\) 的唯一非平凡一维特征 \(\xi\) 满足 \(\xi|_{\langle p\rangle}\equiv 1\)（等价于 \(\xi\) 经由 \(H/\langle p\rangle\simeq C_2\) 的符号特征给出），并且
\[
W_0(N)\cong \operatorname{Ind}_H^{S_4}(\xi)
\qquad\text{（作为 }\QQ[S_4]\text{-模）}.
\]
其不可约分解按 \(r\) 分类为
\[
\begin{array}{c|l}
r=1\ (H\simeq C_2) &
W_0(N)\cong \varepsilon\oplus \rho_2\oplus \rho_3\oplus 2\rho_3'\\
\hline
r=2\ (H\simeq V_4) &
W_0(N)\cong \rho_3\oplus \rho_3'\\
\hline
r=3\ (H\simeq S_3) &
W_0(N)\cong \varepsilon\oplus \rho_3'
\end{array}
\]
\end{theorem}
\begin{proof}[证明要点]
半稳定退化的极限混合 Hodge 结构理论给出
\(
W_0(N)\cong H_1(\Gamma,\QQ)
\)，其中 \(\Gamma\) 为中心纤维对偶图。
在“两条换位分歧点碰撞”的 admissible-cover 局部模型中，中心纤维由一个主分量与 \([S_4:H]\) 个有理尾分量组成；每个尾分量与主分量之间形成一条图圈。
群 \(H\) 在该图圈上的作用由 \(H/\langle p\rangle\simeq C_2\) 的非平凡元给出负号，因而对应特征 \(\xi\)。
整体 \(S_4\)-作用通过左陪集 \(S_4/H\) 的传递置换提升，得到
\(
W_0(N)\cong \operatorname{Ind}_H^{S_4}(\xi)
\)。
对 \(H\simeq C_2,V_4,S_3\) 三类情形，用 Frobenius 互反律
\[
m_W=\dim\operatorname{Hom}_{S_4}\!\bigl(W,\operatorname{Ind}_H^{S_4}\xi\bigr)
=\dim\operatorname{Hom}_H(W,\xi)
\]
逐一计算重数，即得表中分解。证毕。
\end{proof}

\begin{corollary}[三类边界处主分量的 \texorpdfstring{$H^{1,0}$}{H10} 表示完全可判定]\label{cor:xi-s4-boundary-main-component-h10-decomposition}
记中心纤维正规化为 \(\widetilde X_0\)，其主分量记为 \(X_{\mathrm{main}}\)。
由于尾分量均为有理曲线，
\[
H^1(\widetilde X_0,\QQ)\cong H^1(X_{\mathrm{main}},\QQ).
\]
在 Grothendieck 群意义下有
\[
[H^1(X_{\mathrm{main}},\QQ)]
=[H^1(X_t,\QQ)]-2[W_0(N)].
\]
从而
\[
[H^{1,0}(X_{\mathrm{main}})]=[V]-[W_0(N)].
\]
因此三类边界上
\[
\begin{array}{c|l}
r=1 & H^{1,0}(X_{\mathrm{main}})\cong \varepsilon\oplus \rho_3'\\
\hline
r=2 & H^{1,0}(X_{\mathrm{main}})\cong 2\varepsilon\oplus \rho_2\oplus 2\rho_3'\\
\hline
r=3 & H^{1,0}(X_{\mathrm{main}})\cong \varepsilon\oplus \rho_2\oplus \rho_3\oplus 2\rho_3'
\end{array}
\]
\end{corollary}
\begin{proof}[证明要点]
半稳定退化下
\[
\operatorname{Gr}_1^W H^1_{\lim}\cong H^1(\widetilde X_0,\QQ),\qquad
\operatorname{Gr}_0^W H^1_{\lim}\cong W_0(N),\qquad
\operatorname{Gr}_2^W H^1_{\lim}\cong W_0(N)(-1),
\]
故半单化意义下
\[
H^1(X_t,\QQ)^{\mathrm{ss}}
\simeq W_0(N)\oplus H^1(\widetilde X_0,\QQ)\oplus W_0(N).
\]
于是得到 \(H^1(X_{\mathrm{main}})\) 的 Grothendieck 等式；再取 \((1,0)\)-部分并代入定理 \ref{thm:xi-s4-boundary-vanishing-cycle-induced-module} 的三类分解即可。证毕。
\end{proof}

\begin{theorem}[Jacobian 半稳定约化中各 \texorpdfstring{$\QQ[S_4]$}{Q[S4]}-等型因子的环面秩]\label{thm:xi-s4-isotypic-jacobian-toric-rank-by-boundary-type}
\leanverified{paper\_xi\_s4\_isotypic\_jacobian\_toric\_rank\_by\_boundary\_type}
记 \(J_t:=\operatorname{Jac}(X_t)\)，对不可约表示 \(W\in\{\varepsilon,\rho_2,\rho_3,\rho_3'\}\) 定义等型因子
\[
A_W:=e_W(J_t).
\]
设半稳定极限中 \(A_W\) 的环面秩为 \(t_W\)，则
\[
t_W=\dim_{\QQ}\!\bigl(W_0(N)\cap H^1(A_W,\QQ)\bigr).
\]
由定理 \ref{thm:xi-s4-boundary-vanishing-cycle-induced-module} 可得
\[
\begin{array}{c|c}
r & (t_{\varepsilon},t_{\rho_2},t_{\rho_3},t_{\rho_3'})\\
\hline
1 & (1,2,3,6)\\
2 & (0,0,3,3)\\
3 & (1,0,0,3)
\end{array}
\]
\end{theorem}
\begin{proof}[证明要点]
对结点曲线 Jacobian 的半稳定极限，环面部分维数等于权重零部分维数；并且 \(\QQ[S_4]\)-作用与极限混合 Hodge 结构相容。
故每个等型因子的环面秩由 \(W_0(N)\) 在该等型中的分量维数给出。
将定理 \ref{thm:xi-s4-boundary-vanishing-cycle-induced-module} 的分解逐项读取即可得到表中数值。证毕。
\end{proof}

\begin{corollary}[三类边界的等型零化模式内禀判别]\label{cor:xi-s4-boundary-type-intrinsic-by-isotypic-nilpotent-pattern}
记 \(N_W\) 为 \(N\) 在 \(H^1(A_W,\QQ)\) 上的限制。
则三类边界可由下列模式唯一刻画：
\[
\begin{array}{c|l}
r=1 & N_{\rho_2}\neq 0\ \ (\Leftrightarrow\ t_{\rho_2}=2)\\
\hline
r=2 & N_{\rho_3}\neq 0,\ N_{\varepsilon}=0\ \ (\Leftrightarrow\ (t_{\rho_3},t_{\varepsilon})=(3,0))\\
\hline
r=3 & N_{\varepsilon}\neq 0,\ N_{\rho_2}=N_{\rho_3}=0\ \ (\Leftrightarrow\ (t_{\varepsilon},t_{\rho_2},t_{\rho_3})=(1,0,0))
\end{array}
\]
\end{corollary}
\begin{proof}[证明要点]
由定理 \ref{thm:xi-s4-isotypic-jacobian-toric-rank-by-boundary-type}，
\[
N_W\neq 0\iff t_W>0.
\]
三类边界对应的 \((t_W)\) 模式两两不同，故可由等型零化/非零化模式内禀区分。证毕。
\end{proof}

\begin{theorem}[主分量 Jacobian 的等型分解与表示缺失]\label{thm:xi-s4-main-component-jacobian-isotypic-missing-pattern}
\leanverified{paper\_xi\_s4\_main\_component\_jacobian\_isotypic\_missing\_pattern}
设 \(J_{\mathrm{main}}:=\operatorname{Jac}(X_{\mathrm{main}})\)。
则其 \(\QQ[S_4]\)-等型由推论 \ref{cor:xi-s4-boundary-main-component-h10-decomposition} 完全决定，并满足：
\[
\begin{array}{c|l}
r=1 & e_{\rho_2}(J_{\mathrm{main}})=0,\ \ e_{\rho_3}(J_{\mathrm{main}})=0\\
\hline
r=2 & e_{\rho_3}(J_{\mathrm{main}})=0\\
\hline
r=3 & \varepsilon,\rho_2,\rho_3,\rho_3'\ \text{四类等型均出现}
\end{array}
\]
\end{theorem}
\begin{proof}[证明要点]
平滑曲线 Jacobian 的等型出现与否等价于对应不可约表示在 \(H^{1,0}\) 中出现与否。
将推论 \ref{cor:xi-s4-boundary-main-component-h10-decomposition} 的三类分解逐项读取，即得上述缺失/出现判别。证毕。
\end{proof}

\begin{proposition}[等型周期映射的最大退化判据]\label{prop:xi-s4-isotypic-period-map-maximal-degeneration}
\leanverified{paper\_xi\_s4\_isotypic\_period\_map\_maximal\_degeneration}
设 \(\mathcal P_W\) 为参数化 \(A_W\) 的等型周期映射。
则：
\begin{enumerate}
  \item 当 \(r=1\) 时，\(\rho_2\) 与 \(\rho_3\) 两个等型满足 \(t_{\rho_2}=2,\ t_{\rho_3}=3\)，对应因子在边界处达到最大幂零退化；
  \item 当 \(r=2\) 时，\(\rho_3\) 等型达到最大幂零退化，而 \(\varepsilon,\rho_2\) 等型具有潜在良好约化；
  \item 当 \(r=3\) 时，\(\rho_2,\rho_3\) 等型具有潜在良好约化，而 \(\varepsilon\) 等型出现一维环面退化。
\end{enumerate}
\end{proposition}
\begin{proof}[证明要点]
对半稳定退化阿贝尔簇，等型环面秩 \(t_W\) 等于对应权重零维数；当 \(t_W=\dim A_W\) 时即为最大幂零退化。
将定理 \ref{thm:xi-s4-isotypic-jacobian-toric-rank-by-boundary-type} 的环面秩表逐项代入即可得到三条断言。证毕。
\end{proof}


\endinput


\paragraph{判别式消去、Widom 边缘与 Lee--Yang 简单边缘的体积量子化标度}
在前述 Lee--Yang 判别式脊线与分支几何框架下，可将“最近上半平面边缘”的临界条件统一改写为实代数消去问题，并在简单边缘情形下导出体积指标的二次量子化律。
下述命题链把判别式分支、有限阶递推模型与局部 Puiseux 参数化接入同一可审计接口。

\begin{theorem}[判别式--Widom 消去原理]\label{thm:xi-discriminant-widom-elimination}
设
\[
\Pi(\lambda,y;\theta)\in \ZZ[\lambda,y,\theta_1,\dots,\theta_m]
\]
在 \(\lambda\) 上首一、次数 \(d\ge 2\)，且其系数对 \((y,\theta)\) 为实系数多项式。
记
\[
D(y;\theta):=\operatorname{Disc}_{\lambda}\!\bigl(\Pi(\lambda,y;\theta)\bigr)\in \ZZ[y,\theta_1,\dots,\theta_m].
\]
设存在开集 \(U\subset\RR^m\) 与解析映射 \(y_\star:U\to\CC\) 满足：
\begin{enumerate}
  \item \(D(y_\star(\theta);\theta)=0\)，且 \(\partial_y D(y_\star(\theta);\theta)\neq 0\)；
  \item \(\Im y_\star(\theta)>0\)。
\end{enumerate}
定义边缘高度
\[
\delta(\theta):=\Im y_\star(\theta),
\qquad
\mathrm{Crit}(\delta):=\{\theta\in U:\nabla_\theta\delta(\theta)=0\}.
\]
写 \(y=u+iv\)（\(u,v\in\RR,\ v>0\)）并分解
\[
D(u+iv;\theta)=D_R(u,v;\theta)+i\,D_I(u,v;\theta),
\]
其中 \(D_R,D_I\in\ZZ[u,v,\theta]\)。
对 \(j=1,\dots,m\) 定义
\[
E_j(u,v;\theta):=\Im\!\Bigl(\partial_{\theta_j}D(u+iv;\theta)\,
\overline{\partial_y D(u+iv;\theta)}\Bigr),
\]
则 \(E_j\in\ZZ[u,v,\theta]\)。
令
\[
\mathcal V:=
\Bigl\{
(u,v,\theta):
D_R=0,\ D_I=0,\ E_j=0\ \forall j
\Bigr\}.
\]
则有
\[
\mathrm{Crit}(\delta)\subset \pi_\theta(\mathcal V),
\]
其中 \(\pi_\theta\) 为到 \(\theta\)-空间的投影。

当 \(m=1\) 且 \(\delta'\not\equiv 0\) 时，\(\mathrm{Crit}(\delta)\) 在任意紧区间内为有限集；并且其元素均满足某个整系数消去多项式方程，因而均为代数数。
\end{theorem}
\begin{proof}
由 \(\partial_y D(y_\star(\theta);\theta)\neq 0\) 与隐函数定理，\(y_\star\) 在 \(U\) 上解析，且
\[
\partial_{\theta_j}y_\star(\theta)=
-\frac{\partial_{\theta_j}D(y_\star(\theta);\theta)}
{\partial_y D(y_\star(\theta);\theta)}.
\]
故
\[
\partial_{\theta_j}\delta(\theta)=
\Im\bigl(\partial_{\theta_j}y_\star(\theta)\bigr)=0
\iff
\Im\!\left(
\frac{\partial_{\theta_j}D}{\partial_y D}(y_\star(\theta);\theta)
\right)=0.
\]
清分母得
\[
\Im\!\Bigl(
\partial_{\theta_j}D(y_\star(\theta);\theta)\,
\overline{\partial_y D(y_\star(\theta);\theta)}
\Bigr)=0.
\]
写 \(y_\star(\theta)=u(\theta)+iv(\theta)\) 并利用 \(D\) 对 \(y\) 为实系数多项式，复共轭对应 \(v\mapsto -v\)，故上式可写为多项式条件
\[
E_j\bigl(u(\theta),v(\theta);\theta\bigr)=0.
\]
连同 \(D_R=D_I=0\)，得到
\(
(u(\theta),v(\theta),\theta)\in\mathcal V
\)，从而
\(
\theta\in\pi_\theta(\mathcal V)
\)。

若 \(m=1\) 且 \(\delta'\not\equiv 0\)，则 \(\delta'\) 为非常零解析函数，其零点在区间内离散；紧区间内仅有有限个。
另一方面，\(\pi_\theta(\mathcal V)\) 由消去 \(u,v\) 得到整系数多项式关系，故其点为代数数。证毕。
\end{proof}

\begin{corollary}[固定阶递推/有限状态 Hankel 模型的 Widom 线半代数闭包]\label{cor:xi-widom-semialgebraic-for-fixed-order-recurrence}
设 \(\{Z_n(y;\theta)\}_{n\ge 0}\) 满足固定阶线性递推
\[
Z_{n+d}(y;\theta)=\sum_{k=0}^{d-1} a_k(y;\theta)\,Z_{n+k}(y;\theta),
\qquad
a_k\in \ZZ[y,\theta_1,\dots,\theta_m].
\]
设伴随矩阵 \(M(y;\theta)\) 的特征多项式为 \(\Pi(\lambda,y;\theta)\)，并记
\[
D(y;\theta)=\operatorname{Disc}_{\lambda}\!\bigl(\Pi(\lambda,y;\theta)\bigr).
\]
若在开集 \(U\subset\RR^m\) 内存在唯一最近上半平面简单分支
\[
y_\star(\theta)\in\{y:\ D(y;\theta)=0,\ \Im y>0\},
\qquad
\partial_y D(y_\star(\theta);\theta)\neq 0,
\]
则：
\begin{enumerate}
  \item 由 \(\delta(\theta)=\Im y_\star(\theta)\) 定义的 Widom 临界集是实代数投影，因而为半代数集合；
  \item 当 \(m=1\) 且 \(\delta'\not\equiv 0\) 时，Widom 临界点在任意紧区间内有限且皆为代数数；
  \item 该临界集可由 \((D_R,D_I,E_1,\dots,E_m)\) 的 Gröbner 消去（或等价消去算法）构造。
\end{enumerate}
\end{corollary}
\begin{proof}
递推给出表示 \(Z_n=\mathbf u^\top M^n\mathbf v\)，分支碰撞由 \(\Pi\) 的 \(\lambda\)-判别式控制。
在简单分支假设下，结论直接由定理 \ref{thm:xi-discriminant-widom-elimination} 与实代数投影封闭性得到。证毕。
\end{proof}

\begin{theorem}[简单 Lee--Yang 边缘的 \texorpdfstring{$n^{-2}$}{n^-2} 量子化律与局部零点参数化]\label{thm:xi-simple-leyang-edge-n2-quantization}
\leanverified{paper\_xi\_simple\_leyang\_edge\_n2\_quantization}
设 \(\Pi(\lambda,y)\in\CC[\lambda,y]\) 在 \(\lambda\) 上首一、次数 \(d\ge 2\)。
取与之同特征多项式的一族 \(d\times d\) 矩阵 \(M(y)\)，并定义
\[
Z_n(y):=\tr\!\bigl(M(y)^n\bigr)=\sum_{j=1}^{d}\lambda_j(y)^n.
\]
设存在 \((\lambda_0,y_0)\in\CC^2\) 使
\[
\Pi(\lambda_0,y_0)=0,\quad
\partial_\lambda\Pi(\lambda_0,y_0)=0,\quad
\partial_{\lambda\lambda}\Pi(\lambda_0,y_0)\neq 0,\quad
\partial_y\Pi(\lambda_0,y_0)\neq 0,
\]
并存在邻域 \(V\ni y_0\) 满足：
\begin{enumerate}
  \item 两个分支在 \(y_0\) 合并，
  \[
  \lambda_\pm(y)=\lambda_0\pm a\,(y-y_0)^{1/2}+O(y-y_0),
  \qquad
  a^2=-\frac{2\,\partial_y\Pi(\lambda_0,y_0)}{\partial_{\lambda\lambda}\Pi(\lambda_0,y_0)};
  \]
  \item 存在 \(q\in(0,1)\) 使对 \(y\in V\setminus\{y_0\}\) 有
  \[
  \max_{j\ge 3}|\lambda_j(y)|\le q\,|\lambda_0|.
  \]
\end{enumerate}
则对每个固定 \(k\in\NN\cup\{0\}\)，当 \(n\to\infty\) 时，\(Z_n\) 在 \(V\) 内存在零点 \(y_{n,k}\) 满足
\[
y_{n,k}
=
y_0
-\frac{\pi^2(2k+1)^2}{4n^2}\cdot\frac{\lambda_0^{\,2}}{a^2}
+O(n^{-3}).
\]
并且这些零点在 \(V\) 内最终均为单零点，其离开 \(y_0\) 的主方向由
\(\arg(-\lambda_0^2/a^2)\) 决定（共轭成对出现）。
\end{theorem}
\begin{proof}
由泰勒展开
\[
\Pi(\lambda,y)
=
\frac12\Pi_{\lambda\lambda}(\lambda_0,y_0)(\lambda-\lambda_0)^2
+\Pi_y(\lambda_0,y_0)(y-y_0)
+O\!\bigl(|\lambda-\lambda_0|^3+|y-y_0||\lambda-\lambda_0|+|y-y_0|^2\bigr),
\]
解 \(\Pi(\lambda,y)=0\) 得
\[
(\lambda-\lambda_0)^2
=
-\frac{2\Pi_y}{\Pi_{\lambda\lambda}}(y-y_0)+O(|y-y_0|^{3/2}),
\]
即得到两条 Puiseux 分支。

由谱支配性写作
\[
Z_n(y)=\lambda_+(y)^n+\lambda_-(y)^n+R_n(y),
\qquad
|R_n(y)|\le C\,(q|\lambda_0|)^n.
\]
定义
\[
r(y):=\frac{\lambda_+(y)}{\lambda_-(y)}
=1+\frac{2a}{\lambda_0}\sqrt{y-y_0}+O(y-y_0),
\]
故
\[
\log r(y)=\frac{2a}{\lambda_0}\sqrt{y-y_0}+O(y-y_0).
\]
主方程 \(\lambda_+^n+\lambda_-^n=0\) 等价于 \(r(y)^n=-1\)，即
\[
\log r(y)=\frac{(2k+1)\pi i}{n}
\]
（取局部一致分支）。
代入并解得
\[
\sqrt{y-y_0}
=
\frac{(2k+1)\pi i}{n}\cdot\frac{\lambda_0}{2a}+O(n^{-2}),
\]
平方即得
\[
y-y_0
=
-\frac{\pi^2(2k+1)^2}{4n^2}\cdot\frac{\lambda_0^2}{a^2}
+O(n^{-3}).
\]
最后在以该近似点为中心、半径 \(O(n^{-2})\) 的小圆上，利用
\(
|R_n|=o(|\lambda_+^n+\lambda_-^n|)
\)
与 Rouché 定理，得到每个近似根附近恰有一个真根，且位移为 \(O(n^{-3})\)；
导数主项非退化给出最终单零性。方向由主项相位给出。证毕。
\end{proof}

\begin{proposition}[最近零点位移常数的导数闭式与二点体积外推]\label{prop:xi-nearest-root-shift-constant-and-richardson}
\leanverified{paper\_xi\_nearest\_root\_shift\_constant\_and\_richardson}
在定理 \ref{thm:xi-simple-leyang-edge-n2-quantization} 条件下，令 \(y_{n,0}\) 为 \(V\) 内最靠近 \(y_0\) 的零点支。
则存在常数
\[
C_\star
:=
-\frac{\pi^2}{4}\cdot\frac{\lambda_0^2}{a^2}
=
\frac{\pi^2\,\lambda_0^{\,2}\,\partial_{\lambda\lambda}\Pi(\lambda_0,y_0)}
{8\,\partial_y\Pi(\lambda_0,y_0)},
\]
使
\[
y_{n,0}=y_0+\frac{C_\star}{n^2}+O(n^{-3}).
\]
定义二点外推
\[
\widehat y_0^{(n)}
:=
\frac{(n+1)^2y_{n+1,0}-n^2y_{n,0}}{(n+1)^2-n^2},
\]
则
\[
\widehat y_0^{(n)}=y_0+O(n^{-3}).
\]
\end{proposition}
\begin{proof}
第一式为定理 \ref{thm:xi-simple-leyang-edge-n2-quantization} 在 \(k=0\) 的直接特例；
再由
\(
a^2=-2\Pi_y/\Pi_{\lambda\lambda}
\)
得到导数闭式。
外推式是对
\(
y_{n,0}=y_0+C_\star n^{-2}+O(n^{-3})
\)
的 Richardson 一次消元，\(n^{-2}\) 主项被抵消，故误差提升至 \(O(n^{-3})\)。证毕。
\end{proof}

\begin{conjecture}[有限状态有效配分函数的 Widom 代数闭包与二次体积律]\label{conj:xi-effective-partition-widom-algebraic-closure}
设 \(\mathcal Z_V(\mu,T)\) 为有限体积 \(V\) 的有效模型配分函数，取 \(y=e^\mu\)。
对每个固定截断阶 \(d\)，设 \(\mathcal Z_{V,d}(y;T)\) 在 \(y\) 上为多项式，且随体积指标满足固定阶线性递推（等价于有限维伴随矩阵表示）。
令 \(\Pi_d(\lambda,y;T)\) 为对应特征多项式，\(D_d(y;T)=\operatorname{Disc}_\lambda(\Pi_d)\)。
猜想：
\begin{enumerate}
  \item 对每个固定 \(d\)，由最近上半平面分支 \(y_{\star,d}(T)\) 定义的 Widom 临界集
  \(
  \partial_T\Im y_{\star,d}(T)=0
  \)
  为代数投影，因而是半代数集合；
  \item 固定 \(d\) 且 \(V\to\infty\) 时，最近零点轨迹满足定理 \ref{thm:xi-simple-leyang-edge-n2-quantization} 型 \(n^{-2}\) 主标度，并可由命题 \ref{prop:xi-nearest-root-shift-constant-and-richardson} 的外推策略系统提纯；
  \item 当 \(d\to\infty\) 时，上述代数 Widom 集合在适当拓扑（例如 Hausdorff 局部收敛）下逼近真实 Widom 线。
\end{enumerate}
该猜想将超临界交叉的边缘定位问题组织为判别式分支与代数消去的可计算链条。
\end{conjecture}

\endinput


\paragraph{Fold 纤维复杂度、Fibonacci--Lee--Yang 边界与算术共振接口}
在既有 Fold--certificate 语义与 Lee--Yang 单位圆框架下，可把纤维多重度、算法复杂度、零点输运和算术分裂行为统一写入同一审计账本：纤维层给出微态复杂度增益，Fibonacci 乘积给出自然边界与等分布极限，Joukowsky 输运把单位圆零点闭包到椭圆层，而同余读出侧的相干模数与 Frobenius 类数据给出可有限验证的算术接口。

\begin{theorem}[Fold 的算法全息账本]\label{thm:xi-fold-algorithmic-holographic-ledger}
设 \(m\ge 1\)，\(\Omega_m:=\{0,1\}^m\)，稳定类型空间为 \(X_m\)，并令
\(
\Fold_m:\Omega_m\to X_m
\)
为规范折叠映射。
对 \(x\in X_m\) 定义纤维多重度
\[
d_m(x):=\bigl|\Fold_m^{-1}(x)\bigr|,
\qquad
w_m(x):=\frac{d_m(x)}{2^m}.
\]
记 \(K(\cdot\mid m)\) 为给定 \(m\) 的前缀型 Kolmogorov 复杂度。
若 \(\omega\sim\mathrm{Unif}(\Omega_m)\)，并置 \(x:=\Fold_m(\omega)\)，则存在与 \(m,c\) 无关的绝对常数 \(C>0\)，使任意 \(c\in\mathbb N,\ c\ge 1\) 满足
\[
\mathbb P\!\left(
K(\omega\mid m)\ge K(x\mid m)+\lfloor\log_2 d_m(x)\rfloor-C(\log m+c)
\right)\ge 1-2^{-c},
\]
并且对所有 \(\omega\in\Omega_m\) 恒有
\[
K(\omega\mid m)\le K(x\mid m)+\lceil\log_2 d_m(x)\rceil+C\log m.
\]
\end{theorem}
\begin{proof}
对每个纤维 \(\Fold_m^{-1}(x)\) 取字典序枚举，定义可计算编码
\[
E_m:\Omega_m\to\bigsqcup_{x\in X_m}\{x\}\times\{0,1,\dots,d_m(x)-1\},
\qquad
E_m(\omega)=(x,i),
\]
其中 \(x=\Fold_m(\omega)\)，\(i\) 为 \(\omega\) 在该纤维中的秩。
由 \(\Fold_m\) 可计算且纤维有限可枚举，\(E_m\) 及其逆均可计算。
因此
\[
K(\omega\mid m)\le K(x,i\mid m)+O(1)
\le K(x\mid m)+K(i\mid x,m)+O(\log m).
\]
又 \(0\le i<d_m(x)\)，可取自描述编码使
\[
K(i\mid x,m)\le \lceil\log_2 d_m(x)\rceil+O(\log m),
\]
得到上界。

对下界，固定 \(x\) 后，\(\omega\) 在纤维 \(\Fold_m^{-1}(x)\) 上均匀。
长度小于 \(\lfloor\log_2 d_m(x)\rfloor-c\) 的前缀描述至多产生
\(
2^{\lfloor\log_2 d_m(x)\rfloor-c}\le 2^{-c}d_m(x)
\)
个对象，故条件概率至少 \(1-2^{-c}\) 的 \(\omega\) 满足
\[
K(\omega\mid x,m)\ge \lfloor\log_2 d_m(x)\rfloor-c.
\]
再由链式下界与 \(x\) 可由 \(\omega\) 可计算恢复，
\[
K(\omega\mid m)\ge K(x\mid m)+K(\omega\mid x,m)-O(\log m),
\]
合并即得概率下界。证毕。
\end{proof}

\begin{corollary}[稳定类型可极简而微态仍可近随机]\label{cor:xi-fold-simple-type-complex-microstate}
在定理 \ref{thm:xi-fold-algorithmic-holographic-ledger} 的记号下，
对任意 \(m\) 与任意 \(x\in X_m\)，存在
\(
\omega\in\Fold_m^{-1}(x)
\)
使
\[
K(\omega\mid m)\ge
K(x\mid m)+\lfloor\log_2 d_m(x)\rfloor-O(\log m).
\]
因此可取一列 \(m\to\infty\) 的对象 \((\omega_m)\)，其稳定类型
\(
x_m:=\Fold_m(\omega_m)
\)
满足 \(K(x_m\mid m)=O(\log m)\)，而
\(
K(\omega_m\mid m)=\Omega(\log d_m(x_m))
\)。
当 \(d_m(x_m)\) 指数增长时，微态复杂度可线性增长而稳定类型仍保持低描述复杂度。
\end{corollary}
\begin{proof}
固定 \(x\)。
若纤维内所有 \(\omega\) 都满足
\(
K(\omega\mid m)<K(x\mid m)+\lfloor\log_2 d_m(x)\rfloor-C\log m
\)，
则可由不足 \(d_m(x)\) 个程序覆盖整纤维，矛盾。
故至少存在一个 \(\omega\) 达到所示下界。
令 \(x_m\) 取为低描述复杂度的规范形并使用该下界，即得结论。证毕。
\end{proof}

\begin{theorem}[Fibonacci--Lee--Yang 乘积的单位圆自然边界]\label{thm:xi-fibonacci-leyang-natural-boundary}
令 \(F_1=F_2=1\) 为 Fibonacci 数列，并定义
\[
\mathcal Z(z):=\prod_{j=1}^{\infty}\bigl(1+z^{F_{j+1}}\bigr),
\qquad |z|<1.
\]
则无穷乘积在 \(|z|<1\) 绝对收敛并定义解析函数。
此外，单位圆 \(\{|z|=1\}\) 是 \(\mathcal Z\) 的自然边界：\(\mathcal Z\) 不能在单位圆任一点作解析延拓。
并且其单位圆零点集合稠密。
\end{theorem}
\begin{proof}
当 \(|z|<1\) 时，\(\sum_{j\ge 1}|z|^{F_{j+1}}<\infty\)，故
\(
\sum_j\log(1+z^{F_{j+1}})
\)
绝对收敛，乘积解析。

每个因子 \(1+z^{F_{j+1}}\) 的零点为
\[
z=\exp\!\left(\frac{(2k+1)\pi i}{F_{j+1}}\right),\qquad
k=0,\dots,F_{j+1}-1,
\]
全部位于单位圆。
由于 \(F_{j+1}\to\infty\)，这些点在单位圆上稠密。

反设 \(\mathcal Z\) 可在某个 \(\zeta\in\{|z|=1\}\) 的邻域 \(U\) 解析延拓为 \(\widetilde{\mathcal Z}\)。
由零点稠密性，\(U\) 内存在以 \(\zeta\) 为聚点的无穷零点列，故 \(\widetilde{\mathcal Z}\) 在包含 \(\zeta\) 的连通分支上恒为零。
特别地，\(\widetilde{\mathcal Z}\equiv 0\) 于 \(U\cap\{|z|<1\}\)。
由解析延拓唯一性，\(\mathcal Z\equiv 0\) 于单位圆内，和 \(\mathcal Z(0)=1\) 矛盾。
故单位圆为自然边界。证毕。
\end{proof}

\begin{theorem}[Fibonacci--Lee--Yang 零点的 Haar 等分布]\label{thm:xi-fibonacci-leyang-haar-equidistribution}
令
\[
P_m(z):=\prod_{j=1}^{m}\bigl(1+z^{F_{j+1}}\bigr),
\]
并令 \(\mu_m\) 为其零点（按重数）的归一化计数测度。
则
\[
\mu_m\Rightarrow \mu_{\mathrm{Haar}}
\qquad (m\to\infty),
\]
其中 \(\mu_{\mathrm{Haar}}\) 为单位圆上的 Haar 概率测度。
\end{theorem}
\begin{proof}
记
\(
D_m:=\deg P_m=\sum_{j=1}^{m}F_{j+1}=F_{m+3}-2
\)。
对 \(n\ge 1\)，设 \(\nu_n\) 为方程 \(z^n=-1\) 的根在单位圆上的均匀计数测度，则
\[
\mu_m=\frac{1}{D_m}\sum_{j=1}^{m}F_{j+1}\,\nu_{F_{j+1}}.
\]
对整数 \(k\)，其 Fourier 系数
\[
\widehat{\nu_n}(k)
:=\int z^k\,d\nu_n(z)
=\frac1n\sum_{\ell=0}^{n-1}
\exp\!\left(\frac{(2\ell+1)\pi i k}{n}\right)
=
\begin{cases}
(-1)^{k/n}, & n\mid k,\\
0, & n\nmid k.
\end{cases}
\]
故
\[
\widehat{\mu_m}(k)
=\frac{1}{D_m}\sum_{j=1}^{m}F_{j+1}\,\widehat{\nu_{F_{j+1}}}(k).
\]
当 \(k\neq 0\) 时，仅有限多个 \(j\) 满足 \(F_{j+1}\mid k\)，其加权和有界，而 \(D_m\to\infty\)，故
\(
\widehat{\mu_m}(k)\to 0
\)。
并且 \(\widehat{\mu_m}(0)=1\)。
这与 Haar 测度 Fourier 系数一致，故 \(\mu_m\Rightarrow\mu_{\mathrm{Haar}}\)。证毕。
\end{proof}

\begin{corollary}[Joukowsky--G\"odel 椭圆输运下的零点等分布]\label{cor:xi-joukowsky-godel-ellipse-transport-equidistribution}
对任意 \(r>0\)，定义
\[
J_r(e^{i\theta})=re^{i\theta}+r^{-1}e^{-i\theta},
\]
其像记为椭圆 \(E_r\)，并记
\(
\mu_r:=(J_r)_{\#}\mu_{\mathrm{Haar}}
\)。
令
\[
Q_{m,r}(w):=
\operatorname{Res}_z\!\bigl(P_m(z),\,rz^2-wz+r^{-1}\bigr),
\]
并令 \(\nu_{m,r}\) 为 \(Q_{m,r}\) 的零点归一化计数测度（按重数）。
则
\[
\nu_{m,r}=(J_r)_{\#}\mu_m,
\qquad
\nu_{m,r}\Rightarrow \mu_r\quad(m\to\infty).
\]
特别地，\(Q_{m,r}\) 的全部零点落在 \(E_r\) 上。
\end{corollary}
\begin{proof}
结式定义保证：\(w\) 是 \(Q_{m,r}\) 的根，当且仅当存在 \(z\) 使
\(
P_m(z)=0
\)
且
\(
rz^2-wz+r^{-1}=0
\)。
对 \(|z|=1\) 的零点，后式等价于 \(w=J_r(z)\)，因而零点多重集正是单位圆零点在 \(J_r\) 下的推前，故 \(\nu_{m,r}=(J_r)_{\#}\mu_m\)。
再由定理 \ref{thm:xi-fibonacci-leyang-haar-equidistribution} 与连续映射定理即得弱收敛。证毕。
\end{proof}

\begin{theorem}[Fibonacci 同余扩散的完全可解性与逆扩散闭式]\label{thm:xi-fib-congruence-diffusion-inverse-closed}
\leanverified{paper\_xi\_fib\_congruence\_diffusion\_inverse\_closed}
取整数 \(m\ge 1\)，令 Fibonacci 数列满足 \(F_1=F_2=1\)，并记
\[
M:=F_{m+2},\qquad G_m:=\ZZ/(M\ZZ),\qquad \widehat G_m:=\operatorname{Hom}(G_m,\mathbb T).
\]
对 \(\Omega_t:=\{0,1\}^t\) 中的微态 \(\omega=(\omega_1,\dots,\omega_t)\) 定义
\[
N(\omega):=\sum_{j=1}^{t}\omega_jF_{j+1}\in\ZZ,
\qquad
r(\omega):=N(\omega)\bmod M\in G_m,
\]
并令
\[
d_t(r):=\bigl|\{\omega\in\Omega_t:\ r(\omega)=r\}\bigr|,
\qquad
p_t(r):=\frac{d_t(r)}{2^t}\qquad(r\in G_m).
\]
考虑随机过程 \((R_t)_{t=0}^{m}\)：
\[
R_0=0,\qquad
R_t\equiv R_{t-1}+B_tF_{t+1}\pmod M,
\qquad
B_t\stackrel{\mathrm{i.i.d.}}{\sim}\mathrm{Bernoulli}\!\left(\tfrac12\right).
\]
则对每个 \(t\in\{0,\dots,m\}\)：
\begin{enumerate}
  \item \(\mathcal L(R_t)=p_t\)。
  \item 对任意角色 \(\chi\in\widehat G_m\)，Fourier 系数满足完全因子化
  \[
  \widehat p_t(\chi):=\EE[\chi(R_t)]
  =2^{-t}\prod_{j=1}^{t}\bigl(1+\chi(F_{j+1})\bigr).
  \]
  \item 令 \(u(r):=\frac{1}{M}\) 为 \(G_m\) 上均匀分布，则有 Parseval 恒等式
  \[
  \|p_t-u\|_2^2=\frac1M\sum_{\chi\neq 1}\bigl|\widehat p_t(\chi)\bigr|^2,
  \qquad
  \mathrm{TV}(p_t,u)\le \frac12\sqrt{M}\,\|p_t-u\|_2.
  \]
  并且由 Rényi--\(2\) 熵下界推出
  \[
  D_{\mathrm{KL}}(p_t\mid u)\le \log\!\left(1+\sum_{\chi\neq 1}\bigl|\widehat p_t(\chi)\bigr|^2\right).
  \]
  \item 逆扩散的一步后验具有闭式 logistic 比值：对任意 \(r\in G_m\)，
  \[
  \PP(B_t=1\mid R_t=r)
  =\frac{p_{t-1}(r-F_{t+1})}{p_{t-1}(r)+p_{t-1}(r-F_{t+1})}
  =\sigma\!\Bigl(\log p_{t-1}(r-F_{t+1})-\log p_{t-1}(r)\Bigr),
  \]
  其中 \(\sigma(x)=\frac{1}{1+e^{-x}}\)。
\end{enumerate}
\end{theorem}
\begin{proof}
由一步递推 \(R_t\equiv R_{t-1}\) 或 \(R_t\equiv R_{t-1}+F_{t+1}\ (\mathrm{mod}\ M)\)（对应 \(B_t=0,1\)）得到
\[
p_t(r)=\frac12p_{t-1}(r)+\frac12p_{t-1}(r-F_{t+1}),
\qquad p_0=\delta_0,
\]
其与 \(d_t/2^t\) 的组合递推一致，故 \(\mathcal L(R_t)=p_t\)。
对 \(\chi\in\widehat G_m\)，由独立性与 \(\chi(B_tF_{t+1})\in\{1,\chi(F_{t+1})\}\) 得
\[
\widehat p_t(\chi)=\EE[\chi(R_{t-1})]\cdot \EE[\chi(B_tF_{t+1})]
=\widehat p_{t-1}(\chi)\cdot \frac12\bigl(1+\chi(F_{t+1})\bigr),
\]
归纳得乘积公式。
\par
令 Fourier 变换约定为 \(\widehat f(\chi)=\sum_{r\in G_m}f(r)\chi(r)\)，则 \(\widehat p_t(\chi)=\EE[\chi(R_t)]\) 且 \(\widehat u(1)=1,\widehat u(\chi)=0(\chi\neq 1)\)。
由 Parseval 得
\(
\|p_t-u\|_2^2=\frac1M\sum_{\chi\neq 1}|\widehat p_t(\chi)|^2
\)。
全变差界由 \(\|f\|_1\le \sqrt{M}\|f\|_2\) 立即推出。
又由 \(H(p)\ge H_2(p)=-\log\sum_r p(r)^2\) 得
\[
D_{\mathrm{KL}}(p_t\mid u)=\log M-H(p_t)\le \log\!\bigl(M\sum_r p_t(r)^2\bigr)
=\log\!\bigl(1+M\|p_t-u\|_2^2\bigr),
\]
代入 Parseval 恒等式即得 KL 上界。
\par
最后由 Bayes 公式
\[
\PP(B_t=1\mid R_t=r)
=\frac{\PP(R_{t-1}=r-F_{t+1},B_t=1)}{\PP(R_t=r)}
=\frac{\frac12p_{t-1}(r-F_{t+1})}{\frac12(p_{t-1}(r)+p_{t-1}(r-F_{t+1}))},
\]
化简得到所示比值与 logistic 表达。
\end{proof}

\begin{proposition}[角色截断投影的最优秩逼近与误差尾能量闭式]\label{prop:xi-fib-congruence-fourier-rankk-optimal}
\leanverified{paper\_xi\_fib\_congruence\_fourier\_rankk\_optimal}
沿用定理 \ref{thm:xi-fib-congruence-diffusion-inverse-closed} 的 \(G_m,\widehat G_m\) 与分布 \(p_m\)。
对任意子集 \(S\subset\widehat G_m\)（并约定 \(1\in S\)），定义 Fourier 截断
\[
p_m^{(S)}(r):=\frac1M\sum_{\chi\in S}\widehat p_m(\chi)\,\overline{\chi(r)}.
\]
则 \(p_m^{(S)}\) 是 \(p_m\) 在 Hilbert 空间 \(L^2(G_m)\) 上到子空间 \(\operatorname{span}S\) 的正交投影，因而在所有形如
\(
f(r)=\frac1M\sum_{\chi\in S}c_\chi\overline{\chi(r)}
\)
的函数中唯一最小化 \(\|p_m-f\|_2\)。
并且最小误差满足闭式
\[
\|p_m-p_m^{(S)}\|_2^2=\frac1M\sum_{\chi\notin S}\bigl|\widehat p_m(\chi)\bigr|^2.
\]
特别地，在约束 \(|S|=k\) 下，最优秩-\(k\) 逼近由选取 \(|\widehat p_m(\chi)|\) 最大的 \(k\) 个角色实现，其最小误差等于对应 Fourier 能量尾和。
\end{proposition}
\begin{proof}
有限阿贝尔群上角色构成 \(L^2(G_m)\) 的正交基。
故 \(p_m^{(S)}\) 为正交投影并由投影定理给出最优性；误差平方由 Parseval 对被截去系数求和得到。
最后在 \(|S|=k\) 约束下，误差尾和显然由保留最大的 \(k\) 项最小化。
\end{proof}

\begin{remark}[与可解析扩散、线性状态空间与有限证书的外部镜像]\label{rem:xi-fib-congruence-external-mirror-diffusion-attention}
定理 \ref{thm:xi-fib-congruence-diffusion-inverse-closed} 给出一个可完全解析的离散扩散--逆扩散对偶模型：逆扩散的“score”退化为对数密度差，收敛误差由非平凡角色能量闭式控制。
命题 \ref{prop:xi-fib-congruence-fourier-rankk-optimal} 则把低秩截断的最优误差精确压缩为 Fourier 尾能量。
这些结构与近期对线性注意力、掩码扩散与可验证机制证书的解析研究在形式上同型，可作为本文同余口径的一个可比参照系。\cite{LeeSohnLee2026LinearAttentionICLNoForgetting,LiangTanShroffLiang2026MaskedDiffusionSharpTV,HadadKatzBassan2026AutomatedCircuitDiscoveryProvable}
\end{remark}

\begin{theorem}[黄金均值稳定类型上的掩码后验计数与转移矩阵闭式]\label{thm:xi-masked-posterior-golden-mean-transfer-matrix}
\leanverified{paper\_xi\_masked\_posterior\_golden\_mean\_transfer\_matrix}
令 \(m\ge 1\)，并令
\(
X_m\subset\{0,1\}^m
\)
为黄金均值稳定类型集合（禁止相邻子词 \(11\)，从而 \(|X_m|=F_{m+2}\)，见附录 \ref{subsec:sync-kernel-counting}）。
取 \(U_m\) 为 \(X_m\) 上的均匀分布。
对任意掩码观测 \(y\in\{0,1,\bot\}^m\)（\(\bot\) 表示未知位），定义一致集合
\[
X_m(y):=\{x\in X_m:\ \forall i,\ y_i\in\{0,1\}\Rightarrow x_i=y_i\},
\]
则后验 \(U_m(\cdot\mid y)\) 为 \(X_m(y)\) 上的均匀分布。
令黄金均值约束的转移矩阵为
\[
A:=
\begin{pmatrix}
1&1\\
1&0
\end{pmatrix},
\]
其行/列分别以状态 \(0,1\) 编号（允许转移 \(0\to 0,1\) 与 \(1\to 0\)）。
设 \(I\subseteq\{1,\dots,m\}\) 为一段极大掩码块（连续 \(\bot\)）的指标区间，长度为 \(L:=|I|\)；
令该块左右相邻的已知边界位为 \(a,b\in\{0,1\}\)（若块触及端点，则约定端点之外的边界位为 \(0\)）。
则：
\begin{enumerate}
  \item 该掩码块的可补全数为
  \[
  \#\{\text{补全}\}=(A^{L+1})_{a,b}.
  \]
  \item 对块内第 \(j\) 个位置（\(1\le j\le L\)），其后验边缘概率满足
  \[
  \boxed{\;
  \PP(x_{I_0+j-1}=1\mid y)=
  \frac{(A^{j})_{a,1}\,(A^{L+1-j})_{1,b}}{(A^{L+1})_{a,b}}\;},
  \]
  其中 \(I_0\) 为该块的起始指标（即 \(I=\{I_0,\dots,I_0+L-1\}\)）。
\end{enumerate}
\end{theorem}
\begin{proof}
黄金均值约束等价于在二状态有向图 \(0\rightleftarrows 1\)（删除 \(1\to 1\)）上的路径约束。
因此长度为 \(L\) 的补全过程 \((z_1,\dots,z_L)\) 连同边界 \(z_0=a,z_{L+1}=b\) 与约束 \(z_i z_{i+1}=0\) 等价于从顶点 \(a\) 到 \(b\) 的长度 \(L+1\) 路径计数，给出 (1)。
若进一步固定块内第 \(j\) 位为 \(1\)，则路径必须在第 \(j\) 步到达顶点 \(1\)，其计数由路径拼接分解为 \((A^{j})_{a,1}(A^{L+1-j})_{1,b}\)。
在均匀后验下，边缘概率即为“固定为 \(1\) 的计数 / 总计数”，得到 (2)。证毕。
\end{proof}

\begin{remark}[与逆扩散 logistic 比值的一致形式]\label{rem:xi-masked-posterior-logistic-parallel}
定理 \ref{thm:xi-masked-posterior-golden-mean-transfer-matrix} 的边缘闭式可视为“按端点条件计数的局部桥分解”，其分子/分母均为有限步转移计数。
与定理 \ref{thm:xi-fib-congruence-diffusion-inverse-closed} 的逆扩散后验
\(\PP(B_t=1\mid R_t=r)=\sigma(\log p_{t-1}(r-F_{t+1})-\log p_{t-1}(r))\)
同型：两者均把 Bayes 去噪器压缩为对数计数差（或对数密度差）的 logistic 形式，只是一个发生在稳定语法 \(X_m\) 的局部桥约束上，另一个发生在同余群 \(G_m\) 的加性扩散上。
\end{remark}

\endinput


\begin{theorem}[投影词相干模数与有限可证性]\label{thm:xi-coherence-modulus-finite-decidability}
设 \(w\) 为有限投影词，\(\mathrm{Mod}(w)\) 为其调用的模集合，且
\[
m^\ast(w):=\mathrm{lcm}\bigl(\mathrm{Mod}(w)\bigr).
\]
记 \(\Phi_w\) 为 \(w\) 在逆极限同余对象上的解释，并定义相干模数
\[
\kappa(w):=
\min\Bigl\{
n\mid m^\ast(w):
\Phi_w\ \text{可通过}\ \mathbb Z/n\mathbb Z\ \text{因子化}
\Bigr\}.
\]
在本文既定的投影词语义（串接对应态射复合、模依赖按并集合并）下，成立：
\begin{enumerate}
  \item \(\kappa(w)\) 良定义且 \(\kappa(w)\mid m^\ast(w)\)；
  \item 对任意词 \(u,v\) 与 \(k\ge 1\)，
  \[
  \kappa(uv)=\mathrm{lcm}\bigl(\kappa(u),\kappa(v)\bigr),
  \qquad
  \kappa(u^k)=\kappa(u);
  \]
  \item 令 \(n:=\mathrm{lcm}(\kappa(u),\kappa(v))\)，则
  \[
  \Phi_u=\Phi_v\ \text{在逆极限对象上成立}
  \Longleftrightarrow
  \Phi_u=\Phi_v\ \text{在}\ \mathbb Z/n\mathbb Z\ \text{上逐点成立}.
  \]
\end{enumerate}
\end{theorem}
\leanverified{paper\_xi\_coherence\_modulus\_finite\_decidability}
\begin{proof}
由有限词稳定性，\(\Phi_w\) 至少可通过 \(m^\ast(w)\) 因子化，故定义集合非空；
该集合属于 \(m^\ast(w)\) 的因子集，有限且按整除有最小元，故 (1) 成立。

若 \(\Phi_u,\Phi_v\) 分别可通过 \(n_u,n_v\) 因子化，则复合 \(\Phi_{uv}\) 可通过 \(\mathrm{lcm}(n_u,n_v)\) 因子化，得
\(
\kappa(uv)\mid\mathrm{lcm}(\kappa(u),\kappa(v))
\)。
反向由语义中的模并集闭包：\(uv\) 的可行模必须同时覆盖 \(u,v\) 的最小相干模，故
\(
\mathrm{lcm}(\kappa(u),\kappa(v))\mid\kappa(uv)
\)。
从而得等式；幂等式同理。

对 (3)，由 (2) 可知 \(\Phi_u,\Phi_v\) 均可通过 \(\mathbb Z/n\mathbb Z\) 因子化。
设 \(\pi_n\) 为逆极限到 \(\mathbb Z/n\mathbb Z\) 的投影，则
\(
\Phi_u=\bar\Phi_u\circ\pi_n,\ \Phi_v=\bar\Phi_v\circ\pi_n
\)。
于是逆极限上逐点相等当且仅当 \(\bar\Phi_u=\bar\Phi_v\)，即在 \(\mathbb Z/n\mathbb Z\) 上逐点相等。证毕。
\end{proof}

\begin{theorem}[非交换审计的线性外置证书尺度]\label{thm:xi-sl2-linear-external-certificate}
\leanverified{paper\_xi\_sl2\_linear\_external\_certificate}
令
\[
R=
\begin{pmatrix}
1&1\\
0&1
\end{pmatrix},
\qquad
L=
\begin{pmatrix}
1&0\\
1&1
\end{pmatrix}
\in \mathrm{SL}_2(\mathbb Z_{\ge 0}),
\]
并对有限正整序列 \(a=(a_1,\dots,a_T)\) 定义
\[
M(a):=R^{a_1}L\,R^{a_2}L\cdots R^{a_T}L.
\]
若存在 \(A_{\max}\) 使 \(1\le a_t\le A_{\max}\)（\(1\le t\le T\)），
则存在仅依赖于 \(A_{\max}\) 的常数 \(0<c_1<c_2\) 使
\[
2^{c_1T}\le \|M(a)\|_\infty\le 2^{c_2T},
\]
其中 \(\|\cdot\|_\infty\) 为矩阵条目最大值。
从而二进制位长
\[
\ell_2\!\bigl(\|M(a)\|_\infty\bigr):=
\lfloor\log_2\|M(a)\|_\infty\rfloor+1
\]
满足 \(\ell_2(\|M(a)\|_\infty)=\Theta(T)\)。
\end{theorem}
\begin{proof}
记 \(B_t:=R^{a_t}L\)。
由
\[
R^a=
\begin{pmatrix}
1&a\\
0&1
\end{pmatrix},
\qquad
R^aL=
\begin{pmatrix}
1+a&a\\
1&1
\end{pmatrix},
\]
得 \(\|B_t\|_\infty\le A_{\max}+1=:M_0\)。
对任意非负 \(2\times2\) 矩阵 \(U,V\) 有
\(
\|UV\|_\infty\le 2\|U\|_\infty\|V\|_\infty
\)，
迭代即得
\[
\|M(a)\|_\infty\le (2M_0)^T.
\]

又因 \(a_t\ge 1\)，逐项比较得
\[
B_t\succeq RL=
\begin{pmatrix}
2&1\\
1&1
\end{pmatrix},
\qquad
M(a)\succeq (RL)^T.
\]
而
\[
(RL)^T=
\begin{pmatrix}
F_{2T+1}&F_{2T}\\
F_{2T}&F_{2T-1}
\end{pmatrix},
\]
故
\[
\|M(a)\|_\infty\ge F_{2T+1}
\ge \frac{\varphi^{2T-1}}{\sqrt5}
\ge 2^{c_1T}
\]
（对某个仅依赖常数的 \(c_1>0\)）。
与上界合并即得结论。证毕。
\end{proof}

\begin{theorem}[区间谱的 Lee--Yang 提升与 6D 超立方体椭圆闭包]\label{thm:xi-interval-spectrum-leyang-lift-6d-hypercube}
设 \(P(\lambda)\in\mathbb R[\lambda]\) 的全部零点属于闭区间 \([-2,2]\)，定义
\[
\mathcal L[P](z):=z^{\deg P}\,P\!\bigl(z+z^{-1}\bigr).
\]
则 \(\mathcal L[P]\) 的全部零点都位于单位圆 \(\{|z|=1\}\)。
进一步，对任意 \(r>0\)，经 Joukowsky 输运
\(
w=J_r(z)=rz+r^{-1}z^{-1}
\)
所得结式多项式零点全部落在椭圆 \(E_r\) 上。

特别地，6D 超立方体邻接多项式 \(\chi_{H_6}\) 具有实谱 \(\{6-2k:0\le k\le 6\}\subset[-6,6]\)；
令 \(P_6(t):=\chi_{H_6}(3t)\)，则 \(P_6\) 的零点位于 \([-2,2]\)，故
\(
\mathcal L[P_6]
\)
自动满足 Lee--Yang 单位圆零点性质，并在 \(J_r\) 输运下形成椭圆零点族。
\end{theorem}
\begin{proof}
对任意零点 \(\lambda_0\in[-2,2]\)，存在 \(\theta\in\mathbb R\) 使
\(
\lambda_0=2\cos\theta=e^{i\theta}+e^{-i\theta}
\)。
于是 \(P(\lambda_0)=0\) 当且仅当
\(
P(z+z^{-1})=0
\)
在 \(z=e^{\pm i\theta}\) 处成立，且 \(|z|=1\)。
乘子 \(z^{\deg P}\) 仅清除负幂，不改变有限点零点集合，故 \(\mathcal L[P]\) 零点均在单位圆。
再由结式定义，单位圆零点经 \(w=J_r(z)\) 推前后即得椭圆零点闭包。证毕。
\end{proof}

\begin{theorem}[Lee--Yang--Fold 特征和的 Pisano 周期分解与衰减指数闭式]\label{thm:xi-leyang-fold-pisano-factorization-decay}
令 Fibonacci 数列满足 \(F_1=F_2=1\) 与递推 \(F_{n+2}=F_{n+1}+F_n\)。
取奇素数 \(p\)，并记 \(\zeta_p:=\exp(2\pi i/p)\)。
对 \(m\ge 1\) 定义特征和
\begin{equation}\label{eq:xi-leyang-fold-Sm-def}
S_m(p):=
2^{-m}\sum_{\omega\in\{0,1\}^m}\zeta_p^{\sum_{j=1}^{m}\omega_jF_{j+1}}
=
\prod_{j=1}^{m}\frac{1+\zeta_p^{F_{j+1}}}{2}.
\end{equation}
令 \(T:=\pi(p)\) 为 Fibonacci 数列在模 \(p\) 下的 Pisano 周期（见 \eqref{eq:xi-leyang-fold-Sm-def} 的指数序列），则对任意分解 \(m=qT+r\)（\(q\in\NN\)，\(0\le r<T\)）有严格因子化
\begin{equation}\label{eq:xi-leyang-fold-pisano-block-factorization}
S_m(p)=\bigl(S_T(p)\bigr)^q\,S_r(p).
\end{equation}
并且极限
\begin{equation}\label{eq:xi-leyang-fold-gamma-def}
\gamma(p):=-\lim_{m\to\infty}\frac{1}{m}\log|S_m(p)|
\end{equation}
存在，且满足闭式
\begin{equation}\label{eq:xi-leyang-fold-gamma-closed}
\gamma(p)= -\frac{1}{T}\log\abs{S_T(p)}
 = -\frac{1}{T}\sum_{j=1}^{T}\log\abs{\cos\!\Bigl(\frac{\pi F_{j+1}}{p}\Bigr)} \;>\;0.
\end{equation}
\end{theorem}
\begin{proof}
对 \eqref{eq:xi-leyang-fold-Sm-def} 的和式逐坐标分离得
\[
S_m(p)=2^{-m}\prod_{j=1}^{m}\Bigl(\sum_{\omega_j\in\{0,1\}}\zeta_p^{\omega_jF_{j+1}}\Bigr)
=\prod_{j=1}^{m}\frac{1+\zeta_p^{F_{j+1}}}{2}.
\]
由于 \(F_{n+T}\equiv F_n\ (\bmod\ p)\)（\(T=\pi(p)\)），因子 \(\frac{1+\zeta_p^{F_{j+1}}}{2}\) 在 \(j\) 上以周期 \(T\) 重复，从而得到分解 \eqref{eq:xi-leyang-fold-pisano-block-factorization}。
对 \eqref{eq:xi-leyang-fold-pisano-block-factorization} 取绝对值对数并除以 \(m=qT+r\)，注意到 \(r\) 有界而 \(q/m\to 1/T\)，得到极限 \eqref{eq:xi-leyang-fold-gamma-def} 存在且等于 \(-\frac1T\log|S_T(p)|\)。
\par
最后，对任意整数 \(k\) 有恒等式
\[
\left|\frac{1+\zeta_p^{k}}{2}\right|
=\left|\cos\!\Bigl(\frac{\pi k}{p}\Bigr)\right|,
\]
代入 \(k=F_{j+1}\) 即得 \eqref{eq:xi-leyang-fold-gamma-closed} 的第二个表达式。
当 \(p\) 为奇素数时，\(\cos(\pi k/p)\neq 0\) 对一切 \(k\in\ZZ\) 成立；并且在积 \(\prod_{j=1}^{T}\frac{1+\zeta_p^{F_{j+1}}}{2}\) 中 \(j=1\) 给出因子 \(\abs{\cos(\pi/p)}<1\)，从而 \(|S_T(p)|\in(0,1)\) 且 \(\gamma(p)>0\)。
\end{proof}

\begin{proposition}[Pisano 块的 cyclotomic 范数恒等式与共轭平均衰减指数]\label{prop:xi-leyang-fold-cyclotomic-norm-conjugate-mean}
在定理 \ref{thm:xi-leyang-fold-pisano-factorization-decay} 的设定下，取 \(g\in(\ZZ/p\ZZ)^\times\)，并定义扭曲特征和
\begin{equation}\label{eq:xi-leyang-fold-Sm-twist}
S_m(p;g):=
2^{-m}\sum_{\omega\in\{0,1\}^m}\zeta_p^{g\sum_{j=1}^{m}\omega_jF_{j+1}}
=\prod_{j=1}^{m}\frac{1+\zeta_p^{gF_{j+1}}}{2}.
\end{equation}
令 \(T:=\pi(p)\)，并定义 Pisano 块
\[
A_{p,g}:=S_T(p;g)\in \QQ(\zeta_p).
\]
设
\[
\rho(p):=\min\{n\ge 1:\ p\mid F_n\}
\]
为 Fibonacci 数列的出现阶，则 \(\rho(p)\mid T\) 且 \(p\mid F_n\iff \rho(p)\mid n\)（例如见 \cite{Wall1960FibonacciModm}）。
在此记号下，\(A_{p,1}\) 的范数满足闭式
\begin{equation}\label{eq:xi-leyang-fold-norm-Ap1}
\mathrm N_{\QQ(\zeta_p)/\QQ}\!\bigl(A_{p,1}\bigr)
=\prod_{g=1}^{p-1}A_{p,g}
=2^{-(p-1)\left(T-\frac{T}{\rho(p)}\right)}\in\QQ_{>0},
\end{equation}
从而
\begin{equation}\label{eq:xi-leyang-fold-absprod-Apg}
\prod_{g=1}^{p-1}\abs{A_{p,g}}
=2^{-(p-1)\left(T-\frac{T}{\rho(p)}\right)}.
\end{equation}
定义共轭衰减指数
\[
\gamma_g(p):=-\frac{1}{T}\log\abs{A_{p,g}},
\]
则其共轭平均值满足
\begin{equation}\label{eq:xi-leyang-fold-mean-gamma}
\frac{1}{p-1}\sum_{g=1}^{p-1}\gamma_g(p)=\left(1-\frac{1}{\rho(p)}\right)\log 2.
\end{equation}
\end{proposition}
\begin{proof}
\(\Gal(\QQ(\zeta_p)/\QQ)\cong(\ZZ/p\ZZ)^\times\) 由 \(\sigma_g(\zeta_p)=\zeta_p^g\) 给出；由 \eqref{eq:xi-leyang-fold-Sm-twist} 得 \(\sigma_g(A_{p,1})=A_{p,g}\)，因此
\[
\mathrm N(A_{p,1})
=\prod_{g=1}^{p-1}A_{p,g}
=\prod_{g=1}^{p-1}\prod_{j=1}^{T}\frac{1+\zeta_p^{gF_{j+1}}}{2}
=\prod_{j=1}^{T}\prod_{g=1}^{p-1}\frac{1+\zeta_p^{gF_{j+1}}}{2}.
\]
对固定 \(j\) 分两类：
若 \(p\mid F_{j+1}\)，则 \(\zeta_p^{gF_{j+1}}=1\) 且内积等于 \(1\)。
若 \(p\nmid F_{j+1}\)，则映射 \(g\mapsto gF_{j+1}\) 在 \((\ZZ/p\ZZ)^\times\) 上为置换，从而
\[
\prod_{g=1}^{p-1}(1+\zeta_p^{gF_{j+1}})=\prod_{k=1}^{p-1}(1+\zeta_p^k).
\]
由 \(\prod_{k=1}^{p-1}(x-\zeta_p^k)=\frac{x^p-1}{x-1}\) 取 \(x=-1\)（\(p\) 奇使 \(p-1\) 为偶）得
\(\prod_{k=1}^{p-1}(1+\zeta_p^k)=1\)，因此该类 \(j\) 的内积为 \(2^{-(p-1)}\)。
\par
由出现阶性质 \(p\mid F_n\iff \rho(p)\mid n\) 且 \(\rho(p)\mid T\) 可知，在 \(j\in\{1,\dots,T\}\) 中满足 \(p\mid F_{j+1}\) 的指标恰有 \(T/\rho(p)\) 个。
故满足 \(p\nmid F_{j+1}\) 的指标数为 \(T-T/\rho(p)\)，从而得到 \eqref{eq:xi-leyang-fold-norm-Ap1}。
对绝对值并利用范数与共轭乘积一致性得到 \eqref{eq:xi-leyang-fold-absprod-Apg}，再取对数即得 \eqref{eq:xi-leyang-fold-mean-gamma}。
\end{proof}

\begin{corollary}[Fibonacci--cyclotomic \(2\)-单位性：Pisano 块乘积仅在 \((2)\) 处有算术支撑]\label{cor:xi-leyang-fold-pisano-block-2unit}
取奇素数 \(p\)，令 \(T=\pi(p)\) 并以 \(\rho(p)\) 表示出现阶。
在 \(\ZZ[\zeta_p]\) 中定义
\[
U_p:=\prod_{j=1}^{T}\bigl(1+\zeta_p^{F_{j+1}}\bigr)\in\ZZ[\zeta_p].
\]
则其绝对范数满足闭式
\begin{equation}\label{eq:xi-leyang-fold-norm-Up}
\mathrm N_{\QQ(\zeta_p)/\QQ}(U_p)=2^{(p-1)\,T/\rho(p)}.
\end{equation}
特别地，主理想 \((U_p)\subset\ZZ[\zeta_p]\) 的素理想分解中不出现任何位于奇素数 \(\ell\neq 2\) 之上的素理想；
换言之，\(U_p\) 是一个 \(2\)-单位（其理想支撑仅位于 \((2)\) 上）。
\end{corollary}
\begin{proof}
由定义 \(U_p=2^T A_{p,1}\) 且 \(\mathrm N(2^T)=2^{T(p-1)}\)，再结合命题 \ref{prop:xi-leyang-fold-cyclotomic-norm-conjugate-mean} 的 \eqref{eq:xi-leyang-fold-norm-Ap1} 得到
\[
\mathrm N(U_p)=2^{T(p-1)}\cdot 2^{-(p-1)\left(T-\frac{T}{\rho(p)}\right)}
=2^{(p-1)\,T/\rho(p)}.
\]
若某素理想 \(\mathfrak p\subset\ZZ[\zeta_p]\) 位于奇素数 \(\ell\neq 2\) 之上且 \(\mathfrak p\mid (U_p)\)，则理想范数 \(\mathrm N(\mathfrak p)=\ell^f\) 为奇数，从而 \(\ell\mid \mathrm N(U_p)\)，与 \eqref{eq:xi-leyang-fold-norm-Up} 右端为纯 \(2\) 次幂矛盾。
故 \((U_p)\) 的支撑仅能位于 \((2)\) 上。
\end{proof}

\begin{corollary}[由原始素因子定理导出的无穷共轭平均衰减平台]\label{cor:xi-leyang-fold-infinite-mean-gamma-platform}
对每个整数 \(n>12\)，存在素数 \(p\) 使 \(\rho(p)=n\)（例如由 Lucas 序列的原始素因子定理推出，见 \cite{BiluHanrotVoutier2001PrimitiveDivisorsLucas}）。
从而对每个 \(n>12\) 都存在素数 \(p\) 满足
\[
\frac{1}{p-1}\sum_{g=1}^{p-1}\gamma_g(p)=\left(1-\frac{1}{n}\right)\log 2,
\]
并且这些值随 \(n\to\infty\) 单调趋于 \(\log 2\)。
\end{corollary}
\begin{proof}
取 \(n>12\)。原始素因子定理给出：\(F_n\) 存在一个原始素因子 \(p\)，即 \(p\mid F_n\) 且对一切 \(1\le k<n\) 有 \(p\nmid F_k\)。
由出现阶定义即得 \(\rho(p)=n\)。
代入命题 \ref{prop:xi-leyang-fold-cyclotomic-norm-conjugate-mean} 的平均公式 \eqref{eq:xi-leyang-fold-mean-gamma} 得到所示恒等式；并且右端随 \(n\) 增大严格递增且以 \(\log 2\) 为极限。
\end{proof}

\begin{proposition}[Pisano 块的实值性与相位锁定：整周期上无净相位]\label{prop:xi-leyang-fold-pisano-block-real}
取奇素数 \(p\)，令 \(T=\pi(p)\) 且 \(g\in(\ZZ/p\ZZ)^\times\)。
则 Pisano 块 \(A_{p,g}=S_T(p;g)\) 为实数。
更强地，令
\[
W_T:=\sum_{j=1}^{T}F_{j+1},
\]
则 \(W_T\equiv 0\ (\bmod\ p)\)，并且对补映射 \(\omega\mapsto \omega^c\)（\(\omega_j^c:=1-\omega_j\)）有严格共轭配对，从而 \(A_{p,g}\in\RR\)。
\end{proposition}
\begin{proof}
由恒等式 \(\sum_{k=0}^{n}F_k=F_{n+2}-1\) 得
\(
W_T=\sum_{k=2}^{T+1}F_k=F_{T+3}-2
\)。
由 Pisano 周期定义 \((F_T,F_{T+1})\equiv(0,1)\ (\bmod\ p)\)，从而 \(F_{T+2}\equiv 1\) 且 \(F_{T+3}\equiv 2\)，于是 \(W_T\equiv 0\ (\bmod\ p)\)。
\par
将 \eqref{eq:xi-leyang-fold-Sm-twist} 在 \(m=T\) 写为
\[
A_{p,g}
=2^{-T}\sum_{\omega\in\{0,1\}^T}\zeta_p^{g\sum_{j=1}^{T}\omega_jF_{j+1}}.
\]
对任意 \(\omega\)，记 \(\omega^c\) 为补，则
\[
\sum_{j=1}^{T}\omega^c_jF_{j+1}=W_T-\sum_{j=1}^{T}\omega_jF_{j+1}\equiv -\sum_{j=1}^{T}\omega_jF_{j+1}\pmod p,
\]
从而
\[
\zeta_p^{g\sum \omega^c_jF_{j+1}}
=\zeta_p^{-g\sum \omega_jF_{j+1}}
=\overline{\zeta_p^{g\sum \omega_jF_{j+1}}}.
\]
故和式在复共轭下不变，得到 \(A_{p,g}\in\RR\)。
\end{proof}

\begin{conjecture}[6D--Lee--Yang--Fold 的 Chebotarev 共振律]\label{conj:xi-6d-leyang-fold-chebotarev-resonance}
沿用文内记号：令 \(L_2\) 为 6D 超立方体相关多项式 \(P_2\) 的分裂域，\(L_{\mathrm{LY}}\) 为 Lee--Yang 三次多项式的分裂域，记复合域
\[
K:=L_2L_{\mathrm{LY}}.
\]
对 \(K/\mathbb Q\) 的不分歧素数 \(p\)，记 Frobenius 共轭类为
\(
\operatorname{Frob}_p\subset\operatorname{Gal}(K/\mathbb Q)
\)。
定义
\[
S_m(p):=
2^{-m}\sum_{\omega\in\{0,1\}^m}
\exp\!\left(
\frac{2\pi i}{p}\sum_{j=1}^{m}\omega_jF_{j+1}
\right)
=
\prod_{j=1}^{m}\frac{1+\exp(2\pi iF_{j+1}/p)}{2},
\]
并令
\[
\gamma(p):=
-\lim_{m\to\infty}\frac{1}{m}\log|S_m(p)|
\]
（当 \(p\) 为奇素数时极限由定理 \ref{thm:xi-leyang-fold-pisano-factorization-decay} 保证存在并具有闭式）。

猜想：\(\gamma(p)\) 对每个不分歧 \(p\) 存在，且仅依赖于 \(\operatorname{Frob}_p\) 的共轭类。
此外，\(\operatorname{Frob}_p\) 为恒等类（即 \(p\) 在 \(K\) 中完全分裂）时，\(\gamma(p)\) 严格低于 Chebotarev 意义下的典型值，从而在模 \(p\) 的 Fold 纤维多重度 Fourier 指纹中产生可检测的算术偏振窗口。
\end{conjecture}

\endinput

\paragraph{外置可逆记录轴的条件复杂度下界与 Lee--Yang 零点定量误差界}
本段把两类“不可压缩预算”与“可审计等分布速度”并置到同一时间账本链条中：前者把纤维多重度转化为外置记录在条件 Kolmogorov 复杂度上的硬下界；后者把 Fibonacci--Lee--Yang 零点的 Fourier 系数显式化为可求和的误差尾，从而得到椭圆推前下的同型定量界。

\begin{proposition}[可逆外置记录的条件复杂度障碍]\label{prop:xi-external-record-conditional-K-lb}
\leanverified{paper\_xi\_external\_record\_conditional\_k\_lb}
沿用定理 \ref{thm:xi-fold-algorithmic-holographic-ledger} 的 \(\Fold_m,\Omega_m,X_m,d_m,w_m\) 记号。
设 \(r:\Omega_m\to\{0,1\}^\ast\) 为可计算映射，且
\[
\omega\ \longmapsto\ (\Fold_m(\omega),r(\omega))
\]
为单射。
令 \(U\sim\mathrm{Unif}(\Omega_m)\)，并记 \(X:=\Fold_m(U)\)。
则对任意整数 \(c\ge 1\)，存在与 \(m,c\) 无关的绝对常数 \(C>0\) 使
\[
\PP\!\left(
K\bigl(r(U)\mid X,m\bigr)
\ge
\bigl\lfloor \log_2 d_m(X)\bigr\rfloor -c-C
\right)\ \ge\ 1-2^{1-c}.
\]
从而
\[
\EE\Bigl[K\bigl(r(U)\mid X,m\bigr)\Bigr]
\ge
\EE\bigl[\log_2 d_m(X)\bigr]-O(1).
\]
\end{proposition}

\begin{proof}
固定 \(x\in X_m\)，令纤维 \(S_x:=\Fold_m^{-1}(x)\)（\(|S_x|=d_m(x)\)），并记像集
\(
R_x:=\{r(\omega):\omega\in S_x\}
\)。
单射假设给出 \(|R_x|=d_m(x)\)。
前缀 Kolmogorov 复杂度的计数界断言：满足 \(K(y\mid x,m)\le k\) 的二进制串至多 \(2^{k+1}\) 个。
取
\(
k=\lfloor\log_2 d_m(x)\rfloor-c-C
\)
得到
\[
\#\Bigl\{y\in R_x:\ K(y\mid x,m)\le \log_2 d_m(x)-c-C\Bigr\}
\le 2^{1-c}\,d_m(x).
\]
在条件分布 \(U\mid(X=x)\) 的均匀性下，上式等价于
\[
\PP\Bigl(K(r(U)\mid x,m)\le \log_2 d_m(x)-c-C\mid X=x\Bigr)\le 2^{1-c}.
\]
对 \(X\sim w_m\) 去条件得到概率结论。
期望下界由对阈值事件取期望或用 \(\EE Z\ge a\cdot \PP(Z\ge a)\) 得到。证毕。
\end{proof}

\begin{theorem}[Fibonacci--Lee--Yang 零点的可求和 Fourier 误差界与椭圆推前稳定性]\label{thm:xi-fibonacci-leyang-summable-fourier-error}
沿用定理 \ref{thm:xi-fibonacci-leyang-haar-equidistribution} 的 \(P_m,\mu_m\) 记号，并记 \(D_m=\deg P_m\)。
设可积测试函数 \(f\) 在单位圆上的 Fourier 系数为
\[
f(e^{i\theta})=\sum_{k\in\ZZ} c_k e^{ik\theta},
\qquad
\sum_{k\neq 0}|c_k|\,|k|\log(1+|k|)<\infty.
\]
则存在绝对常数 \(C>0\) 使对所有 \(m\ge 1\) 有
\[
\left|\int f\,d\mu_m-\int f\,d\mu_{\mathrm{Haar}}\right|
\ \le\
\frac{C}{D_m}\sum_{k\neq 0}|c_k|\,|k|\log(1+|k|).
\]
进一步，对任意 \(r>0\) 与推论 \ref{cor:xi-joukowsky-godel-ellipse-transport-equidistribution} 的椭圆零点测度 \(\nu_{m,r}\)，若 \(F\) 为椭圆 \(E_r\) 上的测试函数且 \(f:=F\circ J_r\) 满足上述可求和条件，则同样成立
\[
\left|\int F\,d\nu_{m,r}-\int F\,d\mu_r\right|
\ \le\
\frac{C}{D_m}\sum_{k\neq 0}|c_k|\,|k|\log(1+|k|).
\]
\end{theorem}

\begin{proof}
由定理 \ref{thm:xi-fibonacci-leyang-haar-equidistribution} 的分解
\(
\mu_m=\frac{1}{D_m}\sum_{j=1}^m F_{j+1}\nu_{F_{j+1}}
\)
与 \(\widehat{\nu_n}(k)=0\)（\(n\nmid k\)）可得
\[
\widehat\mu_m(k)
=\frac{1}{D_m}\sum_{\substack{1\le j\le m\\ F_{j+1}\mid k}}F_{j+1}\,\widehat{\nu_{F_{j+1}}}(k),
\qquad k\neq 0,
\]
从而
\[
|\widehat\mu_m(k)|
\le
\frac{1}{D_m}\sum_{\substack{1\le j\le m\\ F_{j+1}\mid k}}F_{j+1}
\le
\frac{|k|}{D_m}\cdot \#\{j:\ F_{j+1}\le |k|\}.
\]
Fibonacci 指数增长给出 \(\#\{j:\ F_{j+1}\le |k|\}\le C_0\log(1+|k|)\)，故
\[
|\widehat\mu_m(k)|\le \frac{C_1}{D_m}\,|k|\log(1+|k|).
\]
于是
\[
\int f\,d\mu_m-\int f\,d\mu_{\mathrm{Haar}}
=\sum_{k\neq 0} c_k\,\widehat\mu_m(k),
\]
由三角不等式与可求和条件即得第一式。
\par
椭圆情形由 \(\nu_{m,r}=(J_r)_{\#}\mu_m\)、\(\mu_r=(J_r)_{\#}\mu_{\mathrm{Haar}}\) 得
\(
\int F\,d\nu_{m,r}-\int F\,d\mu_r=\int (F\circ J_r)\,d\mu_m-\int (F\circ J_r)\,d\mu_{\mathrm{Haar}}
\)，
代入第一式即得。证毕。
\end{proof}

\endinput


\paragraph{Gödel 纤维索引的模均匀化、profinite 极限与 Koenigs--Joukowsky 参数化}
在满射折叠与外置寄存器语义下，纤维内的均匀索引可在有限模层面给出统一的均匀化与去耦误差界；该误差在 \(\Fold_m\) 的 Fibonacci 规模律下呈指数衰减，并提升为 \(\widehat{\mathbb Z}\) 上的 Haar 极限。与此同时，Witt 伸缩在 Koenigs 坐标中的平移化与素数寄存器分解共同给出一条从“算术生成元”到“Joukowsky 椭圆族”的可审计同构通道。

\begin{theorem}[Gödel 纤维索引的模均匀性与输出独立性定量界]\label{thm:xi-godel-fiber-index-mod-uniformity-independence}
设 \(\Omega,X\) 为有限集合，\(f:\Omega\twoheadrightarrow X\) 为满射。
记
\[
d(x):=\bigl|f^{-1}(x)\bigr|,\qquad x\in X.
\]
对每个 \(x\in X\) 固定一个双射
\[
\iota_x:f^{-1}(x)\stackrel{\sim}{\longrightarrow}\{0,1,\dots,d(x)-1\}.
\]
令 \(\omega\sim\mathrm{Unif}(\Omega)\)，设 \(X_\ast:=f(\omega)\)，并定义
\[
I:=\iota_{X_\ast}(\omega),\qquad
R_q:=I\bmod q\in\mathbb Z/q\mathbb Z
\]
（\(q\ge 2\)）。
记 \(U_q:=\mathrm{Unif}(\mathbb Z/q\mathbb Z)\)。
则总变差距离满足
\begin{equation}\label{eq:xi-godel-mod-uniform-tv-bound}
\bigl\|\mathcal L(R_q)-U_q\bigr\|_{\mathrm{TV}}
\le
\min\!\left\{1,\frac{q\,|X|}{4\,|\Omega|}\right\},
\end{equation}
以及
\begin{equation}\label{eq:xi-godel-joint-independence-tv-bound}
\bigl\|\mathcal L(X_\ast,R_q)-\mathcal L(X_\ast)\otimes U_q\bigr\|_{\mathrm{TV}}
\le
\min\!\left\{1,\frac{q\,|X|}{4\,|\Omega|}\right\}.
\end{equation}
特别地，误差仅由规模比 \(|X|/|\Omega|\) 控制，与纤维形状 \(\{d(x)\}_{x\in X}\) 的细节无关。
\leanverified{paper\_xi\_godel\_fiber\_index\_mod\_uniformity\_independence}
\end{theorem}
\begin{proof}
固定 \(x\in X\)，写 \(d(x)=a_xq+b_x\)（\(0\le b_x<q\)）。
在条件 \(X_\ast=x\) 下，\(I\) 在 \(\{0,\dots,d(x)-1\}\) 上均匀，故
\(
R_q
\)
在 \(b_x\) 个剩余类上概率为 \((a_x+1)/d(x)\)，在其余 \(q-b_x\) 个类上概率为 \(a_x/d(x)\)。
于是
\[
\left\|\mathcal L(R_q\mid X_\ast=x)-U_q\right\|_{\mathrm{TV}}
=\frac{b_x(q-b_x)}{q\,d(x)}
\le \frac{q}{4\,d(x)}.
\]
又 \(\mathbb P(X_\ast=x)=d(x)/|\Omega|\)。
利用总变差对混合的凸性可得
\[
\|\mathcal L(R_q)-U_q\|_{\mathrm{TV}}
\le
\sum_{x\in X}\frac{d(x)}{|\Omega|}\cdot\frac{q}{4d(x)}
=\frac{q|X|}{4|\Omega|},
\]
并与平凡上界 \(1\) 取最小得到
\eqref{eq:xi-godel-mod-uniform-tv-bound}。

对联合分布，按 \(X_\ast\) 分层有精确分解
\[
\left\|\mathcal L(X_\ast,R_q)-\mathcal L(X_\ast)\otimes U_q\right\|_{\mathrm{TV}}
=
\sum_{x\in X}\mathbb P(X_\ast=x)
\left\|\mathcal L(R_q\mid X_\ast=x)-U_q\right\|_{\mathrm{TV}},
\]
代入同一单纤维界即得
\eqref{eq:xi-godel-joint-independence-tv-bound}。证毕。
\end{proof}

\begin{corollary}[\(\Fold_m\) 的 profinite Haar 极限与指数收敛率]\label{cor:xi-fold-profinite-haar-exponential-rate}
令 \(\Fold_m:\Omega_m\twoheadrightarrow X_m\) 为折叠算子，\(\Omega_m=\{0,1\}^m\)，并令 \(\omega\sim\mathrm{Unif}(\Omega_m)\)。
在每条纤维上任选双射索引，定义 \(I_m\) 与 \(R_{m,q}:=I_m\bmod q\)（\(q\ge 2\)）。
则
\begin{equation}\label{eq:xi-fold-mod-tv-rate-general}
\bigl\|\mathcal L(R_{m,q})-U_q\bigr\|_{\mathrm{TV}}
\le
\min\!\left\{1,\frac{q\,|X_m|}{4\cdot 2^m}\right\},
\end{equation}
且
\begin{equation}\label{eq:xi-fold-joint-mod-tv-rate-general}
\bigl\|\mathcal L(\Fold_m(\omega),R_{m,q})-\mathcal L(\Fold_m(\omega))\otimes U_q\bigr\|_{\mathrm{TV}}
\le
\min\!\left\{1,\frac{q\,|X_m|}{4\cdot 2^m}\right\}.
\end{equation}
若 \(|X_m|=F_{m+2}\)（无相邻 \(1\) 的二进制词计数），则
\begin{equation}\label{eq:xi-fold-mod-tv-rate-fibonacci}
\bigl\|\mathcal L(R_{m,q})-U_q\bigr\|_{\mathrm{TV}}
=
O\!\left(q\left(\frac{\varphi}{2}\right)^m\right),
\qquad
\varphi=\frac{1+\sqrt5}{2}.
\end{equation}
并且把 \(I_m\) 经自然嵌入 \(\mathbb Z\hookrightarrow\widehat{\mathbb Z}\) 记为 \(\widehat I_m\)，则
\begin{equation}\label{eq:xi-fold-profinite-haar-limit}
\widehat I_m \Longrightarrow \mathrm{Haar}_{\widehat{\mathbb Z}},
\qquad m\to\infty.
\end{equation}
\end{corollary}
\begin{proof}
\eqref{eq:xi-fold-mod-tv-rate-general} 与
\eqref{eq:xi-fold-joint-mod-tv-rate-general}
由定理 \ref{thm:xi-godel-fiber-index-mod-uniformity-independence}
直接代入 \(|\Omega_m|=2^m\) 得到。
\eqref{eq:xi-fold-mod-tv-rate-fibonacci}
由 \(|X_m|=F_{m+2}\sim \varphi^{m+2}/\sqrt5\) 立得。

对
\eqref{eq:xi-fold-profinite-haar-limit}，
只需检验每个有限商 \(\pi_q:\widehat{\mathbb Z}\to\mathbb Z/q\mathbb Z\) 的推前分布收敛到 \(U_q\)。
而 \(\pi_q(\widehat I_m)=R_{m,q}\)，其收敛已由
\eqref{eq:xi-fold-mod-tv-rate-general} 给出。
由于 \(\widehat{\mathbb Z}\) 的柱集合由有限模投影生成，故得到弱收敛到 Haar 测度。证毕。
\end{proof}

\begin{theorem}[外置寄存器的信息屏蔽上界]\label{thm:xi-godel-register-mutual-information-bound}
\leanverified{paper\_xi\_godel\_register\_mutual\_information\_bound}
在定理 \ref{thm:xi-godel-fiber-index-mod-uniformity-independence} 的设定下，令
\[
P_x:=\mathcal L(R_q\mid X_\ast=x),\qquad
\eta_x:=\|P_x-U_q\|_{\mathrm{TV}},
\qquad
h(t):=-t\log t-(1-t)\log(1-t).
\]
则
\begin{equation}\label{eq:xi-godel-mutual-information-bound-general}
\mathbf I(X_\ast;R_q)
\le
\varepsilon_q\log(q-1)+h(\varepsilon_q),
\qquad
\varepsilon_q:=\min\!\left\{1-\frac1q,\frac{q|X|}{4|\Omega|}\right\}.
\end{equation}
尤其在
\(
\frac{q|X|}{4|\Omega|}\le \frac12
\)
的制度下，
\[
\mathbf I(X_\ast;R_q)
\le
\frac{q|X|}{4|\Omega|}\log(q-1)
+h\!\left(\frac{q|X|}{4|\Omega|}\right).
\]
对 \(\Fold_m\) 且 \(|X_m|=F_{m+2}\) 的情形，任意固定 \(q\) 有
\begin{equation}\label{eq:xi-fold-register-mutual-information-rate}
\mathbf I\!\bigl(\Fold_m(\omega);\,I_m\bmod q\bigr)
=
O\!\left(q\left(\frac{\varphi}{2}\right)^m\log q\right).
\end{equation}
\end{theorem}
\begin{proof}
由 \(H(R_q)\le\log q\) 得
\[
\mathbf I(X_\ast;R_q)
=H(R_q)-H(R_q\mid X_\ast)
\le \log q-\sum_{x}\mathbb P(X_\ast=x)H(P_x).
\]
对每个 \(x\)，Fannes--Audenaert 不等式给出
\[
\log q-H(P_x)\le \eta_x\log(q-1)+h(\eta_x),
\qquad
\eta_x\in\left[0,1-\frac1q\right].
\]
故
\[
\mathbf I(X_\ast;R_q)
\le
\sum_x\mathbb P(X_\ast=x)\bigl[\eta_x\log(q-1)+h(\eta_x)\bigr].
\]
利用 \(h\) 的凹性与 Jensen，不等式右侧至多为
\[
\bar\eta\,\log(q-1)+h(\bar\eta),
\qquad
\bar\eta:=\sum_x\mathbb P(X_\ast=x)\eta_x.
\]
由定理 \ref{thm:xi-godel-fiber-index-mod-uniformity-independence} 的单纤维界，
\(
\bar\eta\le q|X|/(4|\Omega|)
\)，并且 \(\bar\eta\le 1-1/q\) 总成立。
函数 \(t\mapsto t\log(q-1)+h(t)\) 在 \([0,1-1/q]\) 单调递增，故可把 \(\bar\eta\) 以 \(\varepsilon_q\) 替换，得到
\eqref{eq:xi-godel-mutual-information-bound-general}。
\eqref{eq:xi-fold-register-mutual-information-rate}
由 \(|\Omega_m|=2^m\)、\(|X_m|=F_{m+2}\) 与
\eqref{eq:xi-godel-mutual-information-bound-general}
直接推出。证毕。
\end{proof}

\begin{proposition}[Koenigs 平移化与素数生成元的谱分解]\label{prop:xi-koenigs-prime-generator-spectral-factorization}
\leanverified{paper\_xi\_koenigs\_prime\_generator\_spectral\_factorization}
设 Koenigs 深度坐标满足 Witt 伸缩平移律
\begin{equation}\label{eq:xi-koenigs-witt-translation-law}
y(ms)=y(s)+\log m,\qquad m\ge 1.
\end{equation}
设素数寄存器半群为
\[
\mathsf P:=\bigoplus_{p\ \mathrm{prime}}\mathbb Z_{\ge 0},
\]
并定义值映射
\begin{equation}\label{eq:xi-prime-register-value-map}
N:\mathsf P\to\mathbb N_{\ge 1},
\qquad
N(r):=\prod_{p}p^{r(p)}.
\end{equation}
则对任意 \(r\in\mathsf P\) 有
\begin{equation}\label{eq:xi-koenigs-prime-register-translation}
y\!\bigl(N(r)s\bigr)
=
y(s)+\sum_{p}r(p)\log p.
\end{equation}
进一步，在 \(L^2(\mathbb R,dy)\) 上定义 Koopman 表示
\[
(T_m f)(y):=f(y+\log m).
\]
其傅里叶像满足
\begin{equation}\label{eq:xi-koopman-fourier-diagonalization}
\widehat{T_m f}(t)=e^{it\log m}\,\widehat f(t)
=m^{it}\widehat f(t)
=\left(\prod_{p}p^{it\,r_m(p)}\right)\widehat f(t),
\end{equation}
其中 \(m=\prod_p p^{r_m(p)}\)。
故伸缩半群在频域上完全对角化，且相位谱值位于单位圆 \(\mathbb S^1\) 并按素数生成元乘法分离。
\end{proposition}
\begin{proof}
\eqref{eq:xi-koenigs-prime-register-translation}
由
\eqref{eq:xi-koenigs-witt-translation-law}
与
\eqref{eq:xi-prime-register-value-map}
及
\(
\log N(r)=\sum_p r(p)\log p
\)
直接得到。
对 \(T_m\) 取傅里叶变换即得
\eqref{eq:xi-koopman-fourier-diagonalization}，
最后代入 \(m\) 的素因子分解完成证明。证毕。
\end{proof}

\begin{corollary}[Joukowsky 椭圆族的素数寄存器参数化]\label{cor:xi-joukowsky-prime-register-ellipse-parameterization}
令
\[
S=y+y^{-1},\qquad
|y|=\sqrt L,\quad L>0.
\]
写 \(y=\sqrt L\,e^{i\theta}\)（\(\theta\in[0,2\pi)\)），则
\begin{equation}\label{eq:xi-joukowsky-ellipse-param}
S(\theta)
=
\left(\sqrt L+\frac{1}{\sqrt L}\right)\cos\theta
+i\left(\sqrt L-\frac{1}{\sqrt L}\right)\sin\theta,
\end{equation}
对应椭圆半轴
\[
a(L)=\sqrt L+\frac{1}{\sqrt L},
\qquad
b(L)=\sqrt L-\frac{1}{\sqrt L}.
\]
若 \(L=N(r)\)（\(r\in\mathsf P\)），则
\begin{equation}\label{eq:xi-joukowsky-prime-register-semiaxes}
a(r)=\sqrt{N(r)}+\frac{1}{\sqrt{N(r)}},
\qquad
b(r)=\sqrt{N(r)}-\frac{1}{\sqrt{N(r)}}.
\end{equation}
并且对应 Koenigs 平移长度为
\begin{equation}\label{eq:xi-joukowsky-koenigs-prime-shift}
\Delta y(r)=\log\sqrt{N(r)}
=\frac12\sum_{p}r(p)\log p.
\end{equation}
因此，选择寄存器状态 \(r\) 与选择椭圆尺度参数 \(L\)（等价地 \(\Delta y\)）一一对应；椭圆族由素数生成元坐标唯一标定。
\end{corollary}
\begin{proof}
\eqref{eq:xi-joukowsky-ellipse-param}
由 \(y=\sqrt L\,e^{i\theta}\) 直接代入 \(S=y+y^{-1}\) 得到。
再以 \(L=N(r)\) 代入即得
\eqref{eq:xi-joukowsky-prime-register-semiaxes}。
式
\eqref{eq:xi-joukowsky-koenigs-prime-shift}
由
\(
\log\sqrt{N(r)}=\frac12\log N(r)=\frac12\sum_p r(p)\log p
\)
与命题 \ref{prop:xi-koenigs-prime-generator-spectral-factorization}
一致。证毕。
\end{proof}

\paragraph{Joukowsky--Gödel 椭圆输运的结式闭式、二次分支核与可逆性}
本段把椭圆输运的结式消元写成二次扩张上的范数，并据此把“不可见相位”精确压缩为二次分支生成的二幂核；同时给出单个非退化椭圆零点对尺度参数的代数可逆性。

\begin{theorem}[Joukowsky--Gödel 输运结式的二次扩张闭式]\label{thm:xi-joukowsky-godel-resultant-quadratic-norm-closed}
设 \(K\subset\mathbb C\) 为一子域且 \(\mathrm{char}(K)\neq 2\)。
取 \(P(z)\in K[z]\) 为次数 \(n\) 的多项式，并取 \(r\in K^{\times}\)。
定义
\[
Q_r(w):=\operatorname{Res}_z\!\bigl(P(z),\,rz^2-wz+r^{-1}\bigr)\in K[w].
\]
则在二次扩张 \(K\bigl(w,\sqrt{w^2-4}\bigr)\) 中恒有恒等式
\[
Q_r(w)
=
r^n\,
P\!\Bigl(\frac{w+\sqrt{w^2-4}}{2r}\Bigr)\,
P\!\Bigl(\frac{w-\sqrt{w^2-4}}{2r}\Bigr).
\]
等价地，若记
\[
\zeta_{\pm}(w):=\frac{w\pm\sqrt{w^2-4}}{2},
\]
则
\[
Q_r(w)
=
r^n\,
\operatorname{Norm}_{K(w,\sqrt{w^2-4})/K(w)}\!\Bigl(P\bigl(\zeta_+(w)/r\bigr)\Bigr).
\]
\end{theorem}
\begin{proof}
记
\[
B(z):=rz^2-wz+r^{-1}\in K[w][z].
\]
在 \(K\bigl(w,\sqrt{w^2-4}\bigr)\) 中有分解
\[
B(z)=r\,(z-\beta_+)(z-\beta_-),
\qquad
\beta_{\pm}=\frac{w\pm\sqrt{w^2-4}}{2r}.
\]
由结式的基本恒等式（以 \(B\) 的首项系数表达），得到
\[
\operatorname{Res}_z(P,B)
=
\bigl(\mathrm{lc}_z(B)\bigr)^n\prod_{B(\beta)=0}P(\beta)
=
r^n\,P(\beta_+)\,P(\beta_-).
\]
这给出第一条闭式。
由于右端在 \(\sqrt{w^2-4}\mapsto-\sqrt{w^2-4}\) 下不变，故其系数落回 \(K(w)\)，并因其作为结式已属于 \(K[w]\)，结论成立。
范数形式由二次扩张的范数定义
\(
\operatorname{Norm}(\xi)=\xi\cdot\xi^{\sigma}
\)
（\(\sigma\) 为非平凡自同构）直接改写得到。证毕。
\end{proof}

\begin{corollary}[输运的二次分支核与二幂扩张]\label{cor:xi-transport-quadratic-branch-2power-kernel}
沿用定理 \ref{thm:xi-joukowsky-godel-resultant-quadratic-norm-closed} 的记号，并假设 \(\mathrm{char}(K)=0\) 且 \(P\) 可分。
把 \(r\) 视为超越元并在基域 \(K(r)\) 上工作。
令 \(L_P\) 为 \(P\) 在 \(K(r)\) 上的分裂域，\(L_Q\) 为 \(Q_r\) 在 \(K(r)\) 上的分裂域。
则有包含 \(L_Q\subseteq L_P\)，并存在有限集合 \(\mathcal W\subset L_Q\)（可取为 \(Q_r\) 的全体根）使
\[
L_P
=
L_Q\Bigl(\,\sqrt{w^2-4}\;:\; w\in\mathcal W\,\Bigr),
\qquad
[L_P:L_Q]=2^m\ \text{对某个 }m\ge 0.
\]
\end{corollary}
\begin{proof}
任取 \(P\) 的根 \(z\in L_P\)。
定义
\[
w:=rz+r^{-1}z^{-1}.
\]
由结式零化的定义，存在同一 \(z\) 使 \(P(z)=0\) 且 \(rz^2-wz+r^{-1}=0\) 同时成立，故 \(w\) 是 \(Q_r\) 的根，因而 \(w\in L_Q\subseteq L_P\)。
又由二次方程
\(
rz^2-wz+r^{-1}=0
\)
可知 \(z\) 落在 \(L_Q(\sqrt{w^2-4})\) 中。
对全部根重复该过程，得到
\[
L_P\subseteq L_Q\Bigl(\,\sqrt{w^2-4}\;:\; w\in\mathcal W\,\Bigr).
\]
反向包含显然，故等号成立。
由于每一步均为至多二次扩张，故 \([L_P:L_Q]\) 为二的幂。证毕。
\end{proof}

\begin{theorem}[超立方体 Lee--Yang 提升的完全二次分解与尺度不增分裂域]\label{thm:xi-hypercube-leyang-lift-complete-quadratic-factorization-scale-invariant-splitting}
令 \(H_d\) 为 \(d\)-维超立方体图，\(A_d\) 为其邻接矩阵，并令
\[
\chi_d(\lambda):=\det(\lambda I-A_d),
\qquad
P_d(u):=\chi_d\!\Bigl(\frac d2\,u\Bigr).
\]
定义 Lee--Yang 提升多项式
\[
L_d(z):=z^{2^d}\,P_d\!\bigl(z+z^{-1}\bigr)\in\mathbb Z[z],
\]
以及其椭圆输运结式
\[
Q_{d,r}(w):=\operatorname{Res}_z\!\bigl(L_d(z),\,rz^2-wz+r^{-1}\bigr).
\]
记
\[
u_k:=2-\frac{4k}{d}\qquad(k=0,1,\dots,d).
\]
则成立：
\begin{enumerate}
  \item 在单位圆侧有完全二次分解
  \[
  L_d(z)=\prod_{k=0}^{d}\bigl(z^2-u_k z+1\bigr)^{\binom{d}{k}}.
  \]
  \item 对任意 \(r\in\mathbb C^{\times}\)，在椭圆侧有完全二次分解
  \[
  Q_{d,r}(w)=\prod_{k=0}^{d}
  \Bigl(w^2-(r+r^{-1})u_k\,w+\bigl(u_k^2+(r-r^{-1})^2\bigr)\Bigr)^{\binom{d}{k}}.
  \]
  \item 若 \(r\in\mathbb Q_{>0}\)，则 \(Q_{d,r}\) 的全部根均落在固定的多二次扩张
  \[
  K_d:=\mathbb Q\bigl(\sqrt{u_k^2-4}\;:\;1\le k\le d-1\bigr),
  \]
  从而 \(Q_{d,r}\) 在 \(K_d\) 上分裂，且该分裂域不依赖于 \(r\) 的取值。
\end{enumerate}
并且当 \(d=6\) 时有
\[
K_6=\mathbb Q(i,\sqrt2,\sqrt5).
\]
\end{theorem}
\begin{proof}
超立方体 \(H_d\) 的谱为 \(\lambda_k=d-2k\)（重数 \(\binom{d}{k}\)），故
\[
\chi_d(\lambda)=\prod_{k=0}^d(\lambda-(d-2k))^{\binom{d}{k}},
\qquad
P_d(u)=\prod_{k=0}^d(u-u_k)^{\binom{d}{k}}.
\]
于是
\[
L_d(z)=z^{2^d}\prod_{k=0}^d\bigl(z+z^{-1}-u_k\bigr)^{\binom{d}{k}}
=\prod_{k=0}^d\bigl(z^2-u_k z+1\bigr)^{\binom{d}{k}},
\]
得到第一条。

由结式对乘积的乘法性，第二条化归为计算
\(
\operatorname{Res}_z(z^2-uz+1,\,rz^2-wz+r^{-1})
\)。
若 \(z^2-uz+1=0\)，则 \(z^{-1}=u-z\)。
代入 \(w=rz+r^{-1}z^{-1}\) 得
\[
w=(r-r^{-1})z+r^{-1}u,\qquad
z=\frac{w-r^{-1}u}{r-r^{-1}}.
\]
将该式代回 \(z^2-uz+1=0\) 并整理，得到关于 \(w\) 的二次方程
\[
w^2-(r+r^{-1})u\,w+\bigl(u^2+(r-r^{-1})^2\bigr)=0.
\]
故所求结式即为该二次多项式；取 \(u=u_k\) 并按重数相乘即得第二条。

对第三条，二次因子的判别式为
\[
\Delta_k=(r+r^{-1})^2u_k^2-4\bigl(u_k^2+(r-r^{-1})^2\bigr)
=(r-r^{-1})^2(u_k^2-4).
\]
若 \(r\in\mathbb Q\)，则 \(r-r^{-1}\in\mathbb Q\)，故每个根均属于 \(\mathbb Q(\sqrt{u_k^2-4})\)，从而全体根属于 \(K_d\)。
当 \(d=6\) 时，
\(
u_k\in\{\pm\frac43,\pm\frac23,0,\pm2\}
\)，
故 \(\sqrt{u_k^2-4}\) 生成的多二次数域等价于 \(\mathbb Q(i,\sqrt2,\sqrt5)\)。
证毕。
\end{proof}

\begin{proposition}[诱导子图类的 Lee--Yang 闭包在椭圆输运下保持]\label{prop:xi-induced-subgraph-leyang-to-elliptic-closure}
令 \(G\) 为 \(H_d\) 的任一诱导子图，\(A_G\) 为其邻接矩阵。
定义缩放特征多项式
\[
P_G(u):=\det\!\Bigl(\frac d2\,u\,I-A_G\Bigr),
\qquad
L_G(z):=z^{|V(G)|}\,P_G\!\bigl(z+z^{-1}\bigr).
\]
则：
\begin{enumerate}
  \item \(P_G\) 的全部零点均落在实区间 \([-2,2]\)；
  \item \(L_G\) 的全部非零零点均落在单位圆 \(\mathbb T\)；
  \item 对任意 \(r>0\)，若定义
  \[
  Q_{G,r}(w):=\operatorname{Res}_z\!\bigl(L_G(z),\,rz^2-wz+r^{-1}\bigr),
  \]
  则 \(Q_{G,r}\) 的全部零点均落在椭圆 \(E_r=J_r(\mathbb T)\) 上。
\end{enumerate}
\end{proposition}
\begin{proof}
因 \(G\) 为诱导子图，\(A_G\) 为 \(A_d\) 的主子矩阵。
由 Cauchy 交错定理与 \(\sigma(A_d)\subset[-d,d]\)，得 \(\sigma(A_G)\subset[-d,d]\)。
于是 \(P_G(u)=0\) 的根为 \(u=2\lambda/d\)（\(\lambda\in\sigma(A_G)\)），故落在 \([-2,2]\)，得到第一条。
第二条由恒等式
\(
z+z^{-1}\in[-2,2]\iff |z|=1
\)
与 \(L_G(z)=0\iff P_G(z+z^{-1})=0\)（忽略平凡因子 \(z^{|V(G)|}\)）推出。
第三条由结式零化定义与第二条：若 \(w\) 为 \(Q_{G,r}\) 的根，则存在 \(z\in\mathbb T\) 使 \(w=J_r(z)=rz+r^{-1}z^{-1}\)，故 \(w\in E_r\)。证毕。
\end{proof}

\begin{theorem}[单个非退化椭圆零点对尺度参数的代数可逆性]\label{thm:xi-single-elliptic-zero-reconstructs-scale}
取 \(u\in K\) 满足 \(u^2\neq 4\)，并令未知 \(r\in K^{\times}\) 与未知 \(w\in\overline K\) 满足
\begin{equation}\label{eq:xi-elliptic-transport-quadratic-relation}
w^2-(r+r^{-1})u\,w+\bigl(u^2+(r-r^{-1})^2\bigr)=0.
\end{equation}
令
\[
\Delta_u:=u^2-4,\qquad \Delta_w:=w^2-4.
\]
则 \(a:=r+r^{-1}\) 为二次方程
\[
a^2-(uw)a+\bigl(w^2+u^2-4\bigr)=0
\]
的根，从而
\[
a=\frac{uw\pm\sqrt{\Delta_u\Delta_w}}{2}.
\]
进而 \(r\) 为
\(
t^2-a t+1=0
\)
的根，即
\[
r=\frac{a\pm\sqrt{a^2-4}}{2}.
\]
\end{theorem}
\begin{proof}
把 \eqref{eq:xi-elliptic-transport-quadratic-relation} 视为关于
\(
a=r+r^{-1}
\)
与
\(
b=r-r^{-1}
\)
的约束，并注意 \(a^2-b^2=4\)。
由 \eqref{eq:xi-elliptic-transport-quadratic-relation} 得
\[
w^2-auw+u^2+b^2=0
\quad\Rightarrow\quad
b^2=auw-w^2-u^2.
\]
代入 \(a^2-b^2=4\) 得
\[
a^2-(uw)a+(w^2+u^2-4)=0,
\]
从而
\[
a=\frac{uw\pm\sqrt{u^2w^2-4(w^2+u^2-4)}}{2}
=\frac{uw\pm\sqrt{(u^2-4)(w^2-4)}}{2}.
\]
给定 \(a\) 后，由 \(r+r^{-1}=a\) 可得 \(r\) 为 \(t^2-a t+1=0\) 的根，公式成立。证毕。
\end{proof}

\begin{conjecture}[输运二次分支的典型满秩独立性]\label{conj:xi-transport-quadratic-branches-typical-full-rank}
设 \(K\) 为特征 \(0\) 的数域，\(r\) 为超越元，\(P(z)\in K[z]\) 为次数 \(n\) 的可分多项式，并满足 Lee--Yang 单位圆零点条件且无额外对称性退化。
令
\[
Q_r(w):=\operatorname{Res}_z\!\bigl(P(z),\,rz^2-wz+r^{-1}\bigr).
\]
记 \(L_P,L_Q\) 分别为 \(P,Q_r\) 在 \(K(r)\) 上的分裂域。
猜想：由推论 \ref{cor:xi-transport-quadratic-branch-2power-kernel} 给出的二次分支元
\(\sqrt{w^2-4}\)（\(w\) 遍历 \(Q_r\) 的根）在 \(L_P\) 中达到最大独立秩，从而在自然作用下
\[
\operatorname{Gal}(L_P/K(r))
\cong
(\mathbb Z/2\mathbb Z)^n\rtimes \operatorname{Gal}(L_Q/K(r)),
\]
且该半直积中的 \(2\)-核秩达到理论上界 \(n\)。
\end{conjecture}

\endinput


\paragraph{分布式格、素数账本与椭圆度量的统一外置}
在纤维侧，定理 \ref{thm:pom-fiber-birkhoff-fence-ideal-lattice} 给出 Fold 纤维的分布式格完备表示；在账本侧，定义 \ref{def:cdim-median-godel-prime-code}--定理 \ref{thm:cdim-squarefree-median-metric-ellipse-realization} 给出平方自由素数寄存器与中值度量的闭式表达及椭圆实现。
本段在一般有限偏序的理想格层面给出同一外置链条，并给出可直接在纤维格上套用的等距与中值闭式。

\begin{theorem}[Gödel--Birkhoff 算术外置：理想格的平方自由素数实现]\label{thm:xi-godel-birkhoff-ideal-lattice-squarefree}
\leanverified{paper\_xi\_godel\_birkhoff\_ideal\_lattice\_squarefree}
设 \(P\) 为有限偏序集，\(J(P)\) 为其序理想（下闭集）按包含关系组成的分布式格。
取单射标号
\[
q:P\hookrightarrow \mathbb P
\]
（\(\mathbb P\) 为素数集合），并记
\[
\NN_{\square}:=\{n\in\NN:\ \nu_p(n)\in\{0,1\}\ \forall p\},
\]
其中 \(\nu_p\) 为 \(p\)-进赋值。
定义 Gödel 乘积码
\[
\Phi_q:J(P)\longrightarrow \NN_{\square},\qquad
\Phi_q(I):=\prod_{p\in I}q(p).
\]
则：
\begin{enumerate}
  \item \(\Phi_q\) 为单射且为格同态；对任意 \(I,J\in J(P)\)，有
  \[
  \Phi_q(I\wedge J)=\gcd(\Phi_q(I),\Phi_q(J)),\qquad
  \Phi_q(I\vee J)=\operatorname{lcm}(\Phi_q(I),\Phi_q(J)).
  \]
  \item 设 \(n\in\NN_{\square}\) 的素因子均属于 \(q(P)\)，则
  \[
  n\in\Phi_q(J(P))
  \quad\Longleftrightarrow\quad
  \forall\, a\le_P b,\ \ q(b)\mid n\Rightarrow q(a)\mid n.
  \]
  换言之，像集由偏序诱导的“素因子下闭性”完全刻画。
\end{enumerate}
\end{theorem}
\begin{proof}
在 \(J(P)\) 中有 \(I\wedge J=I\cap J\) 与 \(I\vee J=I\cup J\)。
由于 \(q\) 单射，\(\Phi_q\) 将集合交/并分别送为平方自由整数的 \(\gcd/\operatorname{lcm}\)，从而得同态关系。
单射性由素因子唯一分解定理立得：\(\Phi_q(I)=\Phi_q(J)\) 蕴含两边素因子集合相同，从而 \(I=J\)。

对第二条，若 \(n=\Phi_q(I)\)，则 \(q(b)\mid n\) 等价于 \(b\in I\)，又 \(I\) 下闭给出 \(a\le_P b\Rightarrow a\in I\)，故 \(q(a)\mid n\)。
反之，若 \(n\) 满足所示蕴含，定义
\[
I(n):=\{p\in P:\ q(p)\mid n\}.
\]
若 \(b\in I(n)\) 且 \(a\le_P b\)，则由蕴含得 \(q(a)\mid n\)，故 \(a\in I(n)\)，从而 \(I(n)\in J(P)\)，且 \(\Phi_q(I(n))=n\)。
证毕。
\end{proof}

\begin{corollary}[Möbius 函数的算术化：区间 Möbius 与数论 Möbius 的同一闭式]\label{cor:xi-ideal-lattice-mobius-arithmeticization}
在定理 \ref{thm:xi-godel-birkhoff-ideal-lattice-squarefree} 的设定下，记 \(\mu_{J(P)}\) 为格 \(J(P)\) 的 Möbius 函数，\(\mu\) 为数论 Möbius 函数。
对任意 \(I\subseteq J\)（序理想包含），记
\[
D:=J\setminus I\subseteq P,\qquad
r:=\frac{\Phi_q(J)}{\Phi_q(I)}=\prod_{p\in D}q(p)\in\NN_{\square}.
\]
则
\[
\mu_{J(P)}(I,J)=
\begin{cases}
\mu(r),& D\ \text{在 \(P\) 中为反链},\\
0,& \text{否则}.
\end{cases}
\]
\end{corollary}
\begin{proof}
映射 \(K\mapsto K\setminus I\) 给出区间同构 \([I,J]\cong J(D)\)（\(D\) 取诱导偏序），故
\(\mu_{J(P)}(I,J)=\mu_{J(D)}(\hat 0,\hat 1)\)。
若 \(D\) 为反链，则 \(J(D)\) 为布尔格，故 \(\mu_{J(D)}(\hat 0,\hat 1)=(-1)^{|D|}\)。
若 \(D\) 含可比对，则 \(J(D)\) 非布尔且顶端 Möbius 数为 \(0\)（见定理 \ref{thm:pom-fence-mobius-rigidity} 的同型论证）。
另一方面，\(r\) 为平方自由，故 \(\mu(r)=(-1)^{|D|}\)。
合并即得结论。证毕。
\end{proof}

\begin{proposition}[多链谱的偏序分层分解：zeta 多链计数与素因子分层计数一致]\label{prop:xi-ideal-lattice-multichain-prime-stratification}
沿用定理 \ref{thm:xi-godel-birkhoff-ideal-lattice-squarefree} 的记号。
对 \(k\in\NN\)，记
\[
\mathcal C_k(P):=\{I_0\subseteq I_1\subseteq\cdots\subseteq I_k\subseteq P:\ I_t\in J(P)\},
\qquad
Z_{J(P)}(k):=|\mathcal C_k(P)|.
\]
令顶元 Gödel 码
\(
n_{\max}:=\Phi_q(P)=\prod_{p\in P}q(p)
\)。
则 \(Z_{J(P)}(k)\) 等于满足下述条件的 \((k+1)\)-元组 \((a_0,\dots,a_k)\) 的数目：
\begin{enumerate}
  \item \(a_0a_1\cdots a_k=n_{\max}\) 且两两互素（等价于把素因子集合 \(q(P)\) 划分为 \(k+1\) 个层）；
  \item 由 \(q(p)\mid a_{\lambda(p)}\) 唯一确定的 \(\lambda:P\to\{0,\dots,k\}\) 为保序映射：
  \[
  a\le_P b\ \Longrightarrow\ \lambda(a)\le \lambda(b).
  \]
\end{enumerate}
\end{proposition}
\begin{proof}
给定多链 \(I_0\subseteq\cdots\subseteq I_k\)，令 \(S_0:=I_0\) 且 \(S_t:=I_t\setminus I_{t-1}\)（\(t\ge 1\)），则 \(P=\bigsqcup_{t=0}^k S_t\)。
置 \(a_t:=\prod_{p\in S_t}q(p)\)，得 \(a_0\cdots a_k=n_{\max}\) 且互素。
由每个 \(I_t\) 的下闭性，若 \(a\le_P b\) 且 \(b\in S_j\)，则 \(b\in I_j\Rightarrow a\in I_j\)，从而 \(a\in S_i\) 且 \(i\le j\)，即 \(\lambda\) 保序。

反向地，若给定 \((a_0,\dots,a_k)\) 满足两条，则由唯一分解得到划分 \(P=\bigsqcup S_t\)，并令 \(I_t:=\bigsqcup_{j\le t}S_j\)。
保序性保证每个 \(I_t\) 下闭，从而得到一条多链。
两构造互逆，故计数相等。证毕。
\end{proof}

\begin{theorem}[素数权 Hasse 度量与 \texorpdfstring{$\delta$}{delta} 的等距性：测地线与线性扩张]\label{thm:xi-ideal-lattice-hasse-metric-delta-geodesics}
在 \(J(P)\) 的 Hasse 覆盖图上赋权：若 \(I\prec I\cup\{p\}\) 为覆盖关系，则令边长
\[
\ell(I,I\cup\{p\}):=\log q(p).
\]
记诱导的最短路度量为 \(d_q\)。
则对任意 \(I,J\in J(P)\) 有严格恒等式
\[
d_q(I,J)
=
\log\!\frac{\operatorname{lcm}(\Phi_q(I),\Phi_q(J))}{\gcd(\Phi_q(I),\Phi_q(J))}
=
\sum_{p\in I\triangle J}\log q(p).
\]
特别地，从 \(\varnothing\) 到 \(P\) 的 \(d_q\)-测地线与 \(P\) 的线性扩张一一对应，从而测地线条数等于 \(e(P)\)，并且
\[
\log\#\{\text{从 }\varnothing\text{ 到 }P\text{ 的 }d_q\text{-测地线}\}=\log e(P).
\]
\end{theorem}
\begin{proof}
任一边对应于添加（或删除）一个元素 \(p\)，并贡献长度 \(\log q(p)\)。
因此任意路径从 \(I\) 到 \(J\) 至少需要对每个 \(I\triangle J\) 中元素执行一次变更，故长度不小于 \(\sum_{p\in I\triangle J}\log q(p)\)。
另一方面，总可先删去 \(I\setminus J\) 中的极大元，再加入 \(J\setminus I\) 中的极小元，且每步保持下闭性；该路径长度恰为该下界，从而得到距离等式。
将集合对称差在 \(\Phi_q\) 下写为平方自由整数的 \(\operatorname{lcm}/\gcd\) 比值，得到第一条恒等式。

当 \(I=\varnothing,J=P\) 时，最短路必须恰添加每个元素一次；下闭性强制添加顺序与偏序兼容，即线性扩张。
反向由线性扩张取前缀下闭包得到测地线，故二者一一对应。证毕。
\end{proof}

\begin{theorem}[平方自由账本到椭圆族的等距中值实现]\label{thm:xi-ideal-lattice-ellipse-median-realization}
定义椭圆族
\[
\mathcal E:=\{E_n:\ n\in\NN_{\square}\},
\qquad
E_n:=\left\{(x,y)\in\RR^2:\ \frac{x^2}{n^2}+y^2n^2\le 1\right\}.
\]
并令 \(\delta_{\mathcal E}(E_n,E_m):=\log(\operatorname{lcm}(n,m)/\gcd(n,m))\)。
定义复合映射
\[
\Psi_q:J(P)\to\mathcal E,\qquad \Psi_q(I):=E_{\Phi_q(I)}.
\]
则 \(\Psi_q\) 为等距嵌入：
\[
\delta_{\mathcal E}(\Psi_q(I),\Psi_q(J))=d_q(I,J)\qquad(\forall I,J\in J(P)).
\]
并且 \(\Psi_q\) 保持三元中值；若记
\[
\mathrm{med}_{J(P)}(I_1,I_2,I_3):=(I_1\cap I_2)\cup(I_2\cap I_3)\cup(I_3\cap I_1),
\qquad n_i:=\Phi_q(I_i),
\]
则
\[
\Phi_q(\mathrm{med}_{J(P)}(I_1,I_2,I_3))
=
\frac{\gcd(n_1,n_2)\gcd(n_2,n_3)\gcd(n_3,n_1)}{\gcd(n_1,n_2,n_3)^2}.
\]
\end{theorem}
\begin{proof}
等距性由定理 \ref{thm:xi-ideal-lattice-hasse-metric-delta-geodesics} 与 \(\delta_{\mathcal E}\) 的定义直接给出。
中值闭式由定理 \ref{thm:cdim-squarefree-median-metric-ellipse-realization} 的平方自由中值公式与定理 \ref{thm:xi-godel-birkhoff-ideal-lattice-squarefree} 的 \(\gcd/\operatorname{lcm}\) 同态性得到；集合公式即分布式格的中值恒等式
\(
\mathrm{med}(a,b,c)=(a\wedge b)\vee(b\wedge c)\vee(c\wedge a)
\)
在 \(J(P)\) 上的具体化。证毕。
\end{proof}

\begin{corollary}[Fold 纤维格的素数账本外置与椭圆等距实现]\label{cor:xi-fold-fiber-ideal-lattice-prime-ledger-ellipse-isometry}
固定 \(x\in X_m\)，令纤维分布式格 \(\mathcal L_x\) 由定理 \ref{thm:pom-fiber-birkhoff-fence-ideal-lattice} 给出，并取其 Birkhoff 基底偏序
\(
P_x\simeq \bigsqcup_i Z_{\ell_i}
\)
使 \(\mathcal L_x\simeq J(P_x)\)。
则任选单射标号 \(q:P_x\hookrightarrow\mathbb P\)，由定理 \ref{thm:xi-godel-birkhoff-ideal-lattice-squarefree}--定理 \ref{thm:xi-ideal-lattice-ellipse-median-realization} 得到纤维的平方自由账本嵌入与椭圆等距实现：
\[
\mathcal L_x
\xhookrightarrow{\ \Phi_q\ }
\NN_{\square}
\xhookrightarrow{\ n\mapsto E_n\ }
\mathcal E,
\]
并且纤维上的素数权 Hasse 距离与算术度量
\(
\delta(n,m)=\log(\operatorname{lcm}(n,m)/\gcd(n,m))
\)
严格一致，三元中值由 \(\gcd\) 闭式唯一确定。
\end{corollary}

\endinput


\paragraph{极大纤维诱导的循环审计轴与素数解封机制}
最大纤维多重度给出一条在偶数分辨率层上不可回避的整除约束，从而在 profinite 读出轴上强制出现一个极小的循环审计子系统。
该子系统在算术完备化意义下与 \(\widehat{\ZZ}\) 等价，并由 Fibonacci 原始素因子定理保证其素数支撑在层间无限次严格扩张。

\begin{definition}[循环审计链与极小链]\label{def:xi-cyclic-audit-chain-minimal}
在分辨率 \(m\) 上，令 \(\Fold_m:\Omega_m\to X_m\) 为折叠算子，并记纤维多重度
\(
d_m(x):=|\Fold_m^{-1}(x)|
\)
及其极大值
\[
D_m:=\max_{x\in X_m} d_m(x).
\]
称一列正整数 \((M_k)_{k\ge 1}\) 为\textbf{循环审计链}，若满足
\[
M_k\mid M_{k+1},
\qquad
D_{2k}\mid M_k
\qquad(k\ge 1).
\]
并定义由极大纤维诱导的\textbf{极小链}
\[
N_k:=\operatorname{lcm}(D_2,D_4,\dots,D_{2k})
\qquad(k\ge 1).
\]
\end{definition}

\begin{theorem}[极小循环审计链与其泛性质]\label{thm:xi-minimal-cyclic-audit-chain-universal}
在定义 \ref{def:xi-cyclic-audit-chain-minimal} 的设定下：
\begin{enumerate}
  \item \((N_k)_{k\ge 1}\) 为循环审计链；
  \item 对任意循环审计链 \((M_k)_{k\ge 1}\)，逐层整除关系 \(N_k\mid M_k\) 对一切 \(k\) 成立；
  \item 令
  \[
  \mathfrak R_{\min}:=\varprojlim_{k}\ZZ/N_k\ZZ,
  \qquad
  \mathfrak R:=\varprojlim_{k}\ZZ/M_k\ZZ,
  \]
  则存在自然的连续满同态 \(\mathfrak R\twoheadrightarrow\mathfrak R_{\min}\)。
  因而 \(\mathfrak R_{\min}\) 在“由循环审计链生成的逆极限寄存器”范畴中满足初始性（任意对象到其存在唯一经由分量投影给出的满同态）。
\end{enumerate}
\end{theorem}
\begin{proof}
由最小公倍数定义立得 \(N_k\mid N_{k+1}\)，且每个 \(D_{2j}\mid N_k\)（\(j\le k\)），故 \((N_k)\) 为循环审计链。

若 \((M_k)\) 为循环审计链，则对任意 \(j\le k\) 有 \(D_{2j}\mid M_j\) 且 \(M_j\mid M_k\)，从而 \(D_{2j}\mid M_k\)。
于是 \(N_k=\operatorname{lcm}(D_2,\dots,D_{2k})\mid M_k\)，得到第二条。

由 \(N_k\mid M_k\) 给出自然满射 \(\ZZ/M_k\ZZ\twoheadrightarrow\ZZ/N_k\ZZ\)，并与转移映射相容；取逆极限即得连续满同态 \(\mathfrak R\twoheadrightarrow\mathfrak R_{\min}\)。
初始性表述即该构造的泛性质重写。证毕。
\end{proof}

\begin{theorem}[极小循环审计轴的算术完备性：\texorpdfstring{$\mathfrak R_{\min}\cong\widehat{\ZZ}$}{Rmin ≅ Zhat}]\label{thm:xi-rmin-is-zhat}
\leanverified{paper\_xi\_rmin\_is\_zhat}
沿用定义 \ref{def:xi-cyclic-audit-chain-minimal} 的记号。
则有自然同构
\[
\mathfrak R_{\min}\ \cong\ \widehat{\ZZ}.
\]
\end{theorem}
\begin{proof}
由定理 \ref{thm:pom-max-fiber}（亦见 \S\ref{subsubsec:window6-discrete-results} 的汇总口径），偶数层极大纤维满足
\[
D_{2k}=F_{k+2},
\]
其中 \(F_n\) 为 Fibonacci 数列（取 \(F_0=0,F_1=1\)）。
因此 \(N_k=\operatorname{lcm}(F_3,F_4,\dots,F_{k+2})\)。

对任意 \(n\ge 2\)，令 \(\pi(n)\) 为 Fibonacci 数列在模 \(n\) 下的 Pisano 周期。
由周期定义，\((F_0,F_1)\equiv(0,1)\pmod n\) 在第 \(\pi(n)\) 步回归，故 \(F_{\pi(n)}\equiv 0\pmod n\)，即 \(n\mid F_{\pi(n)}\)。
令 \(k:=\pi(n)-2\)，则 \(n\mid F_{k+2}=D_{2k}\mid N_k\)。
因而 \((N_k)\) 在整除偏序中共终：对每个 \(n\ge 1\) 存在 \(k\) 使 \(n\mid N_k\)（\(n=1\) 平凡）。

由 profinite 完备化的标准性质，共终子系统上的逆极限与全系统一致：
\[
\varprojlim_{k}\ZZ/N_k\ZZ\ \cong\ \varprojlim_{n}\ZZ/n\ZZ\ =\ \widehat{\ZZ}.
\]
证毕。
\end{proof}

\begin{theorem}[原始素因子强迫律：极小链的素数支撑无限解封]\label{thm:xi-primitive-prime-divisor-forces-unbounded-prime-support}
\leanverified{paper\_xi\_primitive\_prime\_divisor\_forces\_unbounded\_prime\_support}
令 \(N_k=\operatorname{lcm}(D_2,D_4,\dots,D_{2k})\)。
则存在 \(k_0\) 使得对所有 \(k\ge k_0\)，存在素数 \(p_k\) 满足
\[
p_k\mid N_k,
\qquad
p_k\nmid N_{k-1}.
\]
从而集合 \(\{p:\ p\mid N_k\}\) 随 \(k\) 增长而严格扩张无穷次。
\end{theorem}
\begin{proof}
由 \(N_k=\operatorname{lcm}(N_{k-1},D_{2k})\)，若存在素数 \(p\mid D_{2k}\) 且 \(p\nmid N_{k-1}\)，则 \(p\mid N_k\) 且 \(p\nmid N_{k-1}\)。
又由 \(D_{2k}=F_{k+2}\)，问题化归为 Fibonacci 数列的原始素因子存在性。

原始素因子定理断言：除有限例外外，每个 \(F_n\) 具有原始素因子，即存在素数 \(p\mid F_n\) 且对所有 \(m<n\) 有 \(p\nmid F_m\)。
对 Fibonacci 情形，可取例外阈值 \(n\ge 13\)（等价于 \(k\ge 11\)）。
于是对每个 \(k\ge 11\) 可取 \(p_k\mid F_{k+2}\) 且 \(p_k\nmid F_j\)（\(j<k+2\)），从而 \(p_k\nmid \operatorname{lcm}(F_3,\dots,F_{k+1})=N_{k-1}\)，并且 \(p_k\mid N_k\)。
取 \(k_0:=11\) 得结论。证毕。
\end{proof}

\begin{corollary}[任意循环审计链不可能由有限素数集支撑]\label{cor:xi-any-cyclic-audit-chain-needs-infinitely-many-primes}
对任意循环审计链 \((M_k)_{k\ge 1}\)，其素因子集合必为无限集。
并且当 \(k\ge 11\) 时，存在素数 \(p_k\) 使
\[
p_k\mid M_k,\qquad p_k\nmid M_{k-1}.
\]
\end{corollary}
\begin{proof}
由定理 \ref{thm:xi-minimal-cyclic-audit-chain-universal}，对一切 \(k\) 有 \(N_k\mid M_k\)。
再由定理 \ref{thm:xi-primitive-prime-divisor-forces-unbounded-prime-support}，当 \(k\ge 11\) 时可取 \(p_k\mid N_k\) 且 \(p_k\nmid N_{k-1}\)。
于是 \(p_k\mid M_k\)；若 \(p_k\mid M_{k-1}\)，则 \(p_k\mid N_{k-1}\)，矛盾。
因此新素数在 \(M_k\) 中无限次解封，素因子集合必无限。证毕。
\end{proof}

\begin{conjecture}[非循环情形的强泛性质：\texorpdfstring{$\widehat{\ZZ}$}{Zhat} 作为不可约素数审计核]\label{conj:xi-zhat-as-irreducible-prime-audit-kernel}
设 \(\mathfrak R\) 为一个 profinite 审计寄存器，且在每个偶数分辨率 \(2k\) 上，其可读有限商足以对某条达到 \(D_{2k}\) 的纤维给出可核验的满索引证书。
则存在连续满同态
\[
\mathfrak R\twoheadrightarrow \widehat{\ZZ}.
\]
该猜想宣称：即便不先验限定为循环型，任何能够逐层核验区分极大纤维的审计寄存器都必然包含一个不可再约化的 \(\widehat{\ZZ}\) 因子；循环子范畴的情形由定理 \ref{thm:xi-minimal-cyclic-audit-chain-universal}--定理 \ref{thm:xi-rmin-is-zhat} 与推论 \ref{cor:xi-any-cyclic-audit-chain-needs-infinitely-many-primes} 说明。
\end{conjecture}

\endinput


\paragraph{互素模族的外置寄存器去相关、最小熵下界与碰撞稀释}
定理 \ref{thm:xi-godel-fiber-index-mod-uniformity-independence} 给出单一模数下的均匀化与输出独立性定量界。
当外置残差由一族互素模共同给出时，中国剩余同构与总变差的 data-processing 性质允许将单模结果提升为向量残差的同时去相关。
在同一误差界下，可进一步得到条件最小熵下界与二阶碰撞指纹的稀释律。

\begin{theorem}[有限互素模族的外置残差去相关]\label{thm:xi-godel-coprime-mod-family-decoupling}
沿用定理 \ref{thm:xi-godel-fiber-index-mod-uniformity-independence} 的设定。
取两两互素整数 \(q_1,\dots,q_k\ge 2\)，并令
\[
q:=\prod_{j=1}^k q_j,
\qquad
\mathbf R_{\mathbf q}:=\bigl(R_{q_1},\dots,R_{q_k}\bigr)\in \prod_{j=1}^k\mathbb Z/q_j\mathbb Z.
\]
记 \(U_{\mathbf q}:=\bigotimes_{j=1}^k U_{q_j}\) 为直积均匀分布。
则有总变差界
\[
\bigl\|\mathcal L(\mathbf R_{\mathbf q})-U_{\mathbf q}\bigr\|_{\mathrm{TV}}
\le
\min\!\left\{1,\frac{q\,|X|}{4|\Omega|}\right\},
\]
以及联合去相关界
\[
\bigl\|\mathcal L(X_\ast,\mathbf R_{\mathbf q})-\mathcal L(X_\ast)\otimes U_{\mathbf q}\bigr\|_{\mathrm{TV}}
\le
\min\!\left\{1,\frac{q\,|X|}{4|\Omega|}\right\}.
\]
\leanverified{paper\_xi\_godel\_coprime\_mod\_family\_decoupling}
\end{theorem}
\begin{proof}
由中国剩余定理存在群同构
\(
\psi:\mathbb Z/q\mathbb Z\stackrel{\sim}{\to}\prod_{j=1}^k\mathbb Z/q_j\mathbb Z
\)
使得 \(\mathbf R_{\mathbf q}=\psi(R_q)\)，且 \(U_q\) 在 \(\psi\) 下推前为 \(U_{\mathbf q}\)。
由于总变差距离在可测映射推前下收缩，
\[
\|\mathcal L(\mathbf R_{\mathbf q})-U_{\mathbf q}\|_{\mathrm{TV}}
\le
\|\mathcal L(R_q)-U_q\|_{\mathrm{TV}},
\qquad
\|\mathcal L(X_\ast,\mathbf R_{\mathbf q})-\mathcal L(X_\ast)\otimes U_{\mathbf q}\|_{\mathrm{TV}}
\le
\|\mathcal L(X_\ast,R_q)-\mathcal L(X_\ast)\otimes U_q\|_{\mathrm{TV}}.
\]
将定理 \ref{thm:xi-godel-fiber-index-mod-uniformity-independence} 代入模数 \(q\) 即得结论。证毕。
\end{proof}

\begin{corollary}[条件最小熵的几乎处处下界]\label{cor:xi-godel-conditional-minentropy-lower-bound}
\leanverified{paper\_xi\_godel\_conditional\_minentropy\_lower\_bound}
在定理 \ref{thm:xi-godel-coprime-mod-family-decoupling} 的设定下，对每个 \(x\in X\) 定义条件偏差
\[
\eta_x:=\left\|\mathcal L(\mathbf R_{\mathbf q}\mid X_\ast=x)-U_{\mathbf q}\right\|_{\mathrm{TV}}.
\]
则成立精确分解
\[
\left\|\mathcal L(X_\ast,\mathbf R_{\mathbf q})-\mathcal L(X_\ast)\otimes U_{\mathbf q}\right\|_{\mathrm{TV}}
=\sum_{x\in X}\PP(X_\ast=x)\,\eta_x,
\]
从而 \(\EE[\eta_{X_\ast}]\le \varepsilon\)，其中 \(\varepsilon:=\min\{1,q|X|/(4|\Omega|)\}\)。
因此对任意 \(\delta\in(0,1]\)，令
\(
G_\delta:=\{x\in X:\ \eta_x\le \delta\}
\)，则
\[
\PP(X_\ast\notin G_\delta)\le \varepsilon/\delta.
\]
并且对每个 \(x\in G_\delta\) 有条件最小熵下界
\[
H_\infty(\mathbf R_{\mathbf q}\mid X_\ast=x)\ \ge\ \log q-\log(1+\delta q),
\]
其中 \(H_\infty(Y):=-\log\max_y\PP(Y=y)\)。
\end{corollary}
\begin{proof}
前两式由总变差对分层混合的精确分解与定理 \ref{thm:xi-godel-coprime-mod-family-decoupling} 直接得到。
由 Markov 不等式，
\(
\PP(\eta_{X_\ast}>\delta)\le \EE[\eta_{X_\ast}]/\delta\le \varepsilon/\delta
\)。
若 \(x\in G_\delta\)，则对任意 \(\mathbf r\) 由 \(\|\mathcal L(\mathbf R_{\mathbf q}\mid X_\ast=x)-U_{\mathbf q}\|_{\mathrm{TV}}\le \delta\) 得
\(
\PP(\mathbf R_{\mathbf q}=\mathbf r\mid X_\ast=x)\le q^{-1}+\delta
\)，
从而
\(
H_\infty(\mathbf R_{\mathbf q}\mid X_\ast=x)\ge -\log(q^{-1}+\delta)=\log q-\log(1+\delta q)
\)。
证毕。
\end{proof}

\begin{corollary}[素数指数寄存器的均匀性与稳定类型独立性]\label{cor:xi-godel-prime-exponent-register-uniform-independence}
\leanverified{paper\_xi\_godel\_prime\_exponent\_register\_uniform\_independence}
在折叠设定 \(X_\ast=\Fold_m(\omega)\)（\(\omega\sim\mathrm{Unif}(\Omega_m)\)）下，沿用推论 \ref{cor:xi-fold-profinite-haar-exponential-rate} 的纤维索引 \(I_m\)。
取互异素数 \(p_1,\dots,p_k\)，并令
\[
E_{m,j}:=I_m\bmod p_j\in\{0,1,\dots,p_j-1\},
\qquad
\mathbf E_m:=(E_{m,1},\dots,E_{m,k}),
\qquad
q:=\prod_{j=1}^k p_j.
\]
则有
\[
\bigl\|\mathcal L(\mathbf E_m)-\textstyle\bigotimes_{j=1}^k U_{p_j}\bigr\|_{\mathrm{TV}}
\le
\min\!\left\{1,\frac{q\,|X_m|}{4\cdot 2^m}\right\},
\]
以及
\[
\bigl\|\mathcal L(\Fold_m(\omega),\mathbf E_m)-\mathcal L(\Fold_m(\omega))\otimes \textstyle\bigotimes_{j=1}^k U_{p_j}\bigr\|_{\mathrm{TV}}
\le
\min\!\left\{1,\frac{q\,|X_m|}{4\cdot 2^m}\right\}.
\]
进一步令
\[
N_m:=\prod_{j=1}^k p_j^{E_{m,j}},
\qquad
\mathcal N:=\Bigl\{\prod_{j=1}^k p_j^{e_j}:\ 0\le e_j<p_j\Bigr\},
\]
则 \(N_m\) 的分布与 \(\mathcal N\) 上均匀分布的总变差距离满足同一上界，并且 \((\Fold_m(\omega),N_m)\) 与乘积分布的总变差距离满足同一上界。
\end{corollary}
\begin{proof}
对 \(\mathbf E_m\) 的两条界由定理 \ref{thm:xi-godel-coprime-mod-family-decoupling} 取 \(q_j=p_j\) 并代入 \(|\Omega_m|=2^m\)、\(|X|=|X_m|\) 得到。
映射 \((e_1,\dots,e_k)\mapsto \prod_j p_j^{e_j}\) 在 \(\prod_j\{0,\dots,p_j-1\}\) 与 \(\mathcal N\) 之间为双射，故推前保持总变差距离，从而 \(N_m\) 的结论成立；联合分布同理。证毕。
\end{proof}

\begin{proposition}[外置模寄存器的碰撞稀释律]\label{prop:xi-godel-mod-register-collision-dilution}
\leanverified{paper\_xi\_godel\_mod\_register\_collision\_dilution}
沿用推论 \ref{cor:xi-fold-profinite-haar-exponential-rate} 的折叠设定，令 \(q\ge 2\) 并记
\(
P_{m,q}:=\mathcal L(\Fold_m(\omega),R_{m,q})
\)、
\(
w_m:=\mathcal L(\Fold_m(\omega))
\)。
对任意分布 \(\nu\) 定义碰撞概率
\(
\mathrm{Col}(\nu):=\sum_z \nu(z)^2
\)。
则有不等式
\[
\bigl|\mathrm{Col}(P_{m,q})-\tfrac{1}{q}\mathrm{Col}(w_m)\bigr|
\le
2\varepsilon_{m,q},
\qquad
\varepsilon_{m,q}:=\min\!\left\{1,\frac{q\,|X_m|}{4\cdot 2^m}\right\}.
\]
\end{proposition}
\begin{proof}
由 \eqref{eq:xi-fold-joint-mod-tv-rate-general} 得
\(
\|P_{m,q}-w_m\otimes U_q\|_{\mathrm{TV}}\le \varepsilon_{m,q}
\)。
对任意两分布 \(\nu,\nu'\) 有
\[
\bigl|\mathrm{Col}(\nu)-\mathrm{Col}(\nu')\bigr|
=\Bigl|\sum_z(\nu(z)-\nu'(z))(\nu(z)+\nu'(z))\Bigr|
\le \sum_z|\nu(z)-\nu'(z)|\cdot(\nu(z)+\nu'(z))
\le 2\|\nu-\nu'\|_{\mathrm{TV}}.
\]
故
\(
|\mathrm{Col}(P_{m,q})-\mathrm{Col}(w_m\otimes U_q)|\le 2\varepsilon_{m,q}
\)。
另一方面
\(
\mathrm{Col}(w_m\otimes U_q)=\sum_{x,r}w_m(x)^2q^{-2}=\frac{1}{q}\mathrm{Col}(w_m)
\)。
合并即得结论。证毕。
\end{proof}

\begin{remark}[秩索引外置的 Kolmogorov 不可压缩性]\label{rem:xi-godel-rank-index-kolmogorov-incompressible}
对折叠映射 \(\Fold_m\) 的任意可计算纤维索引 \(I_m\)，
将定理 \ref{thm:pom-reversible-external-residual-kolmogorov-lower-bound} 应用于残差 \(R(\omega)=I_m(\omega)\) 得到：
在给定 \(X=\Fold_m(\omega)\) 与 \(m\) 的条件下，\(I_m\) 的前缀 Kolmogorov 复杂度以概率 \(1-2^{-\Omega(1)}\) 达到 \(\log_2 d_m(X)\) 的最优阶，并且存在相应的确定性上界（推论 \ref{cor:pom-reversible-external-residual-kolmogorov-typical-equality}）。
\end{remark}

\endinput


\paragraph{Lee--Yang 扭点像的消元判别式与 \texorpdfstring{$\mathrm{GL}_2(\FF_p)$}{GL2(Fp)} 符号机制}
设
\[
E:\quad Y^2=4X^3-4X+1,
\]
并定义度 \(4\) 的有理函数
\[
y:=X^2+\frac{Y-1}{2}.
\]
消去 \(Y\) 得到
\[
\Pi(X,y)=X^4-X^3-(2y+1)X^2+X+y(y+1)=0.
\]
以下在此固定输入下讨论 \(p\)-扭点像的消元多项式、不可约性、判别式平方类及其群论来源。

\begin{theorem}[五扭点 Lee--Yang 像的消元多项式、不可约性与判别式平方类]\label{thm:xi-leyang-image-five-torsion-elimination-irreducibility-discriminant}
令 \(\psi_5(X)\) 为曲线 \(E\) 的五分裂多项式对应的 \(X\)-版本，可取
\[
\begin{aligned}
\psi_5(X)=\;&5X^{12}-62X^{10}+95X^9-105X^8-60X^7+285X^6\\
&-174X^5-5X^4-5X^3+35X^2-15X+2.
\end{aligned}
\]
定义
\[
L_5(y):=\operatorname{Res}_X\!\bigl(\psi_5(X),\Pi(X,y)\bigr)\in\ZZ[y].
\]
则 \(L_5\) 的 \(y\)-次数为 \(24\)，首项系数为 \(5^4\)，并显式等于
\[
\begin{aligned}
L_5(y)=\;&625y^{24}-23500y^{23}+244475y^{22}+855915y^{21}-11037859y^{20}\\
&-89302755y^{19}-22555640y^{18}+761997855y^{17}+522337946y^{16}\\
&-7536023160y^{15}+2387307475y^{14}+22058496400y^{13}-18687397720y^{12}\\
&-33571140480y^{11}+30722417580y^{10}+11215028400y^9-17547790815y^8\\
&+800735580y^7+3666848562y^6-1687774224y^5-262565904y^4\\
&+264107184y^3-48718224y^2-3276800y+10048.
\end{aligned}
\]
并满足：
\begin{enumerate}
  \item \(L_5\) 在 \(\QQ[y]\) 上不可约；
  \item
  \[
  \operatorname{Disc}(L_5)=5\cdot D_5^2,\qquad D_5\in\ZZ.
  \]
\end{enumerate}
其中
\[
\operatorname{Disc}(L_5)
=
2^{32}\cdot 5^{35}\cdot 19^2\cdot 37^{94}\cdot 359^2\cdot 1181^2\cdot 1399^2\cdot 2389^2\cdot 94201^2\cdot P_5^2,
\]
\[
P_5=
8470418073546514738455998677740719775605918902449324446515157612831208141799,
\]
可取
\[
D_5=
2^{16}\cdot 5^{17}\cdot 19\cdot 37^{47}\cdot 359\cdot 1181\cdot 1399\cdot 2389\cdot 94201\cdot P_5.
\]
\end{theorem}
\begin{proof}
由 \(y=X^2+\frac{Y-1}{2}\) 与曲线方程消去 \(Y\) 得 \(\Pi(X,y)=0\)。
非零五扭点的 \(X\)-坐标由 \(\psi_5(X)=0\) 刻画，故
\[
y\bigl(E[5]\setminus\{O\}\bigr)
\]
在代数闭包中的消元约束由
\(
\operatorname{Res}_X(\psi_5,\Pi)
\)
给出。由符号消元计算得到上述显式多项式与次数。

将 \(L_5\) 模 \(3\) 约化，可得其在 \(\FF_3[y]\) 上不可约（次数保持为 \(24\)）。
因 \(3\nmid 5^4\)，由 Gauss 引理与模素数不可约判别，\(L_5\) 在 \(\QQ[y]\) 上不可约。

最后，判别式及其素因子分解由整数判别式计算与分解直接得到，故
\(
\operatorname{Disc}(L_5)=5\cdot D_5^2
\)。
证毕。
\end{proof}

\begin{proposition}[五扭点像的分离性与二次分解域]\label{prop:xi-leyang-image-five-torsion-separation-quadratic-field}
在定理 \ref{thm:xi-leyang-image-five-torsion-elimination-irreducibility-discriminant} 设定下，映射
\[
y:E\longrightarrow \PP^1
\]
在 \(E[5]\setminus\{O\}\) 上单射。
此外，设 \(K_5\) 为 \(L_5\) 的分裂域，则其由判别式平方类诱导的二次子域为
\[
\QQ(\sqrt{5}).
\]
\end{proposition}
\begin{proof}
\[
|E[5]\setminus\{O\}|=5^2-1=24.
\]
由定理 \ref{thm:xi-leyang-image-five-torsion-elimination-irreducibility-discriminant}，\(\deg L_5=24\) 且 \(\operatorname{Disc}(L_5)\neq 0\)，故 \(L_5\) 在特征 \(0\) 下可分，具有 \(24\) 个互异根。
这些根覆盖全部 \(y(P)\)（\(P\in E[5]\setminus\{O\}\)），故像集基数亦为 \(24\)，从而单射。

又
\[
\operatorname{Disc}(L_5)=5\cdot(\text{平方}),
\]
故判别式平方类对应二次域为 \(\QQ(\sqrt{5})\)。证毕。
\end{proof}

\begin{theorem}[黄金二次数域与 Lee--Yang 主三次闭包的线性互不相交]\label{thm:xi-sqrt5-axis-leyang-cubic-linear-disjoint}
令黄金二次数域
\[
K_\varphi:=\QQ(\sqrt{5}),
\]
并令
\[
P_{\mathrm{LY}}(y):=256y^{3}+411y^{2}+165y+32\in\ZZ[y],
\qquad
L_{\mathrm{LY}}:=\mathrm{Spl}\bigl(P_{\mathrm{LY}}\bigr).
\]
则有
\[
K_\varphi\cap L_{\mathrm{LY}}=\QQ,
\]
从而 \(K_\varphi\) 与 \(L_{\mathrm{LY}}\) 在 \(\QQ\) 上线性互不相交，并且
\[
[K_\varphi L_{\mathrm{LY}}:\QQ]=[K_\varphi:\QQ]\cdot[L_{\mathrm{LY}}:\QQ]=2\cdot 6=12,
\qquad
\Gal(K_\varphi L_{\mathrm{LY}}/\QQ)\cong \Gal(L_{\mathrm{LY}}/\QQ)\times C_2.
\]
\end{theorem}

\begin{proof}
二次扩张 \(K_\varphi/\QQ\) 的判别式为 \(5\)，故仅在素数 \(5\) 处分歧。
另一方面，由结论 \ref{con:xi-terminal-37-squarefree-coupling}，
\[
\Disc\bigl(P_{\mathrm{LY}}\bigr)=-3^9\cdot 31^2\cdot 37,
\]
从而 \(L_{\mathrm{LY}}/\QQ\) 的分歧素数集合包含于 \(\{3,31,37\}\)，尤其在 \(5\) 处不分歧。
若交域 \(K_\varphi\cap L_{\mathrm{LY}}\) 为非平凡二次扩张，则它在 \(5\) 处必分歧，矛盾。
故交域只能为 \(\QQ\)，线性互不相交成立。
次数与直积型伽罗瓦群结论由线性互不相交的标准性质给出（并用 \([K_\varphi:\QQ]=2\)）。证毕。
\end{proof}

\begin{corollary}[双轴完全分裂的 Chebotarev 密度]\label{cor:xi-sqrt5-leyang-double-split-density-1-12}
对任意素数 \(p\nmid 3\cdot 31\cdot 37\cdot 5\)，命题
\[
p\ \text{在}\ K_\varphi\ \text{与}\ L_{\mathrm{LY}}\ \text{中同时完全分裂}
\]
等价于 \(p\) 在合成扩张 \(K_\varphi L_{\mathrm{LY}}/\QQ\) 中完全分裂；其自然密度为
\[
\boxed{\;\delta=\frac{1}{12}\;}
\]
\end{corollary}

\begin{proof}
由定理 \ref{thm:xi-sqrt5-axis-leyang-cubic-linear-disjoint}，合成扩张 \(K_\varphi L_{\mathrm{LY}}/\QQ\) 为次数 \(12\) 的伽罗瓦扩张。
对 \(p\nmid 3\cdot 31\cdot 37\cdot 5\)，上述扩张在 \(p\) 处不分歧，且“在两边同时完全分裂”等价于在合成域中完全分裂。
由 Chebotarev 密度定理，完全分裂对应 Frobenius 共轭类为恒等元，密度为 \(|\{1\}|/|\Gal|=1/12\)。证毕。
\end{proof}

\begin{theorem}[七扭点 Lee--Yang 像的消元多项式、不可约性与判别式平方类]\label{thm:xi-leyang-image-seven-torsion-elimination-irreducibility-discriminant}
沿用 \((E,y,\Pi)\)。
令 \(\psi_7(X)\) 为曲线 \(E\) 的七分裂多项式对应的 \(X\)-版本，可取
\[
\begin{aligned}
\psi_7(X)=\;&7X^{24}-392X^{22}+1148X^{21}+5929X^{20}-40040X^{19}-16072X^{18}\\
&+351792X^{17}-282282X^{16}-230104X^{15}+1369368X^{14}-1302812X^{13}\\
&-192942X^{12}+829752X^{11}-527450X^{10}-310170X^9+369012X^8\\
&-99576X^7-31892X^6+42406X^5-23235X^4-3556X^3+2394X^2-784X+80.
\end{aligned}
\]
定义
\[
L_7(y):=\operatorname{Res}_X\!\bigl(\psi_7(X),\Pi(X,y)\bigr)\in\ZZ[y].
\]
则 \(L_7\) 的 \(y\)-次数为 \(48\)，首项系数为 \(7^4\)，并显式等于
\[
\begin{aligned}
L_7(y)=\;&2401y^{48}-216090y^{47}+8508661y^{46}-195748830y^{45}+2930004106y^{44}\\
&-30316346864y^{43}+222334655189y^{42}-1157342551860y^{41}+3991074225594y^{40}\\
&-7202709575204y^{39}-10657811407452y^{38}+112205816041752y^{37}-378960173753194y^{36}\\
&+813792173949916y^{35}-1272579912391434y^{34}+1701235626444784y^{33}\\
&-1636138513560978y^{32}-197998457901880y^{31}+4420149463274830y^{30}\\
&-10825321480542788y^{29}+19180870490835505y^{28}-24357792947158800y^{27}\\
&+20556374176128478y^{26}+2129138162785994y^{25}-33507974213689042y^{24}\\
&+75751558581142088y^{23}-119383507152698410y^{22}+143124429682746056y^{21}\\
&-131987006124578241y^{20}+91192136369810754y^{19}-46742852057176406y^{18}\\
&+16707500387178988y^{17}-3486011714189714y^{16}-99972999135036y^{15}\\
&+521443046392150y^{14}-802990216949360y^{13}+728350156733752y^{12}\\
&-514747019925328y^{11}+284639553390553y^{10}-120638941704504y^9\\
&+37113475496646y^8-6945498996040y^7-140075424201y^6+262233631328y^5\\
&-153469178224y^4+46882788672y^3-7564012320y^2+578009600y-16056320.
\end{aligned}
\]
并满足：
\begin{enumerate}
  \item \(L_7\) 在 \(\QQ[y]\) 上不可约；
  \item
  \[
  \operatorname{Disc}(L_7)=-7\cdot D_7^2,\qquad D_7\in\ZZ.
  \]
\end{enumerate}
其中
\[
\operatorname{Disc}(L_7)=
-3^{12}\cdot 7^{63}\cdot 37^{380}\cdot 43^2\cdot 191^2\cdot 143977^2\cdot 2730979^2\cdot 2473661231^2\cdot P_7^2,
\]
\[
\begin{aligned}
P_7=\;&896605137920355289318322421516871363636708810373820861334045825064975794195448006252805195318424980131241138817965808845409424354511493219373761004696208382587053432155947941,
\end{aligned}
\]
可取
\[
D_7=
3^{6}\cdot 7^{31}\cdot 37^{190}\cdot 43\cdot 191\cdot 143977\cdot 2730979\cdot 2473661231\cdot P_7.
\]
\end{theorem}
\begin{proof}
与五扭点情形同理，
\(
L_7
\)
由 \(\psi_7\) 与 \(\Pi\) 的消元定义，显式与次数来自符号消元。

模 \(5\) 约化后，\(L_7\) 在 \(\FF_5[y]\) 上不可约，且 \(5\nmid 7^4\)。
由 Gauss 引理，\(L_7\) 在 \(\QQ[y]\) 上不可约。

判别式及其整数分解由计算直接给出，从而
\(
\operatorname{Disc}(L_7)=-7\cdot D_7^2
\)。
证毕。
\end{proof}

\begin{proposition}[七扭点像的分离性与二次分解域]\label{prop:xi-leyang-image-seven-torsion-separation-quadratic-field}
在定理 \ref{thm:xi-leyang-image-seven-torsion-elimination-irreducibility-discriminant} 设定下，映射 \(y\) 在 \(E[7]\setminus\{O\}\) 上单射。
此外，\(L_7\) 分裂域由判别式平方类诱导的二次子域为
\[
\QQ(\sqrt{-7}).
\]
\end{proposition}
\begin{proof}
\[
|E[7]\setminus\{O\}|=7^2-1=48.
\]
由定理 \ref{thm:xi-leyang-image-seven-torsion-elimination-irreducibility-discriminant}，\(L_7\) 为次数 \(48\) 的可分多项式，故有 \(48\) 个互异根，恰对应全部七扭点像值，故单射。

又
\[
\operatorname{Disc}(L_7)=-7\cdot(\text{平方}),
\]
故二次子域为 \(\QQ(\sqrt{-7})\)。证毕。
\end{proof}

\begin{theorem}[\texorpdfstring{$\mathrm{GL}_2(\FF_p)$}{GL2(Fp)} 在 \texorpdfstring{$\FF_p^2\setminus\{0\}$}{Fp^2\\{0\\}} 上的符号特征]\label{thm:xi-gl2-fp-action-sign-character-legendre-det}
\leanverified{paper\_xi\_gl2\_fp\_action\_sign\_character\_legendre\_det}
令 \(p\) 为奇素数，\(V:=\FF_p^2\)，\(V^\times:=V\setminus\{0\}\)，并令
\[
G:=\mathrm{GL}_2(\FF_p)
\]
以自然方式作用于 \(V^\times\)。
记该置换表示的符号特征为
\[
\operatorname{sgn}_{V^\times}:G\to\{\pm1\}.
\]
则对任意 \(g\in G\) 成立
\[
\operatorname{sgn}_{V^\times}(g)=\left(\frac{\det(g)}{p}\right),
\]
其中 \(\bigl(\frac{\cdot}{p}\bigr)\) 为 Legendre 符号。

因而，若某不可约多项式 \(F\in\QQ[y]\) 的根上伽罗瓦作用与
\(
G\curvearrowright V^\times
\)
同构，则其判别式平方类满足
\[
\QQ\!\bigl(\sqrt{\operatorname{Disc}(F)}\bigr)
=
\QQ\!\bigl(\sqrt{(-1)^{(p-1)/2}p}\bigr),
\]
等价地
\[
\operatorname{Disc}(F)\equiv (-1)^{(p-1)/2}p
\quad\text{于}\quad
\QQ^\times/(\QQ^\times)^2.
\]
\end{theorem}
\begin{proof}
\(\operatorname{sgn}_{V^\times}\) 是群同态。
由于
\[
G^{\mathrm{ab}}\simeq \FF_p^\times
\]
并由 \(\det\) 实现，任意 \(G\to\{\pm1\}\) 同态必经由 \(\det\) 因子化，故
\(
\operatorname{sgn}_{V^\times}
\)
或平凡，或等于
\(
\left(\frac{\det}{p}\right)
\)。

取
\[
g_a:=\operatorname{diag}(a,1),\qquad a\in\FF_p^\times.
\]
其作用为 \((x,y)\mapsto(ax,y)\)。
分解
\[
V^\times=V_0\sqcup W,\qquad
V_0:=\{(0,y):y\neq0\},\ 
W:=\{(x,y):x\neq0\}.
\]
在 \(V_0\) 上 \(g_a\) 恒等。
对每个固定 \(y\in\FF_p\)，纤维
\[
W_y:=\{(x,y):x\in\FF_p^\times\}
\]
上作用同构于 \(\FF_p^\times\) 上的乘法 \(x\mapsto ax\)。

若 \(a\) 在 \(\FF_p^\times\) 中阶为 \(d\)，则每个 \(W_y\) 分解为 \((p-1)/d\) 个长度 \(d\) 的循环，故 \(W\) 上循环总数为 \(p(p-1)/d\)，从而
\[
\operatorname{sgn}_{V^\times}(g_a)
=(-1)^{\frac{p(p-1)}{d}(d-1)}
=(-1)^{\frac{p-1}{d}}.
\]
当 \(a\) 为二次剩余时，\(\frac{p-1}{d}\) 为偶，符号为 \(+1\)；
当 \(a\) 为非二次剩余时，\(\frac{p-1}{d}\) 为奇，符号为 \(-1\)。
故
\[
\operatorname{sgn}_{V^\times}(g_a)=\left(\frac{a}{p}\right)=\left(\frac{\det(g_a)}{p}\right).
\]
于是排除平凡特征并得全体 \(g\) 上恒等式
\[
\operatorname{sgn}_{V^\times}(g)=\left(\frac{\det(g)}{p}\right).
\]

后半段是标准判别式--符号角色对应：判别式平方类由根置换符号特征切出的二次扩张确定，代入上式即得所述二次域与平方类。证毕。
\end{proof}

\begin{conjecture}[奇素数 \texorpdfstring{$p$}{p} 的统一不可约性与判别式平方类模式]\label{conj:xi-leyang-torsion-image-uniform-pattern-all-odd-primes}
在固定 \((E,y,\Pi)\) 下，对每个奇素数 \(p\) 定义
\[
L_p(y):=\operatorname{Res}_X\!\bigl(\psi_p(X),\Pi(X,y)\bigr)\in\ZZ[y],
\]
其中 \(\psi_p\) 为 \(p\)-分裂多项式的 \(X\)-版本。

猜想成立：
\begin{enumerate}
  \item \(L_p\) 在 \(\QQ[y]\) 上不可约；
  \item
  \[
  \deg L_p=p^2-1;
  \]
  \item
  \[
  \operatorname{Disc}(L_p)=(-1)^{(p-1)/2}p\cdot(\text{整数平方});
  \]
  \item \(L_p\) 根上的伽罗瓦作用同构于
  \[
  \mathrm{GL}_2(\FF_p)\curvearrowright \FF_p^2\setminus\{0\}.
  \]
\end{enumerate}
该模式与本节 \(p=5,7\) 结果及既有低阶情形兼容，并将判别式二次子域统一锁定为
\[
\QQ\!\bigl(\sqrt{(-1)^{(p-1)/2}p}\bigr).
\]
\end{conjecture}

\paragraph{Lee--Yang 判别式证书与对数坐标下的数论分层}
沿用 \(2026\)-\(02\)-\(22\) 文档的谱核记号，记 \(\Pi(\lambda,y)\in\overline{\mathbb Q}[\lambda,y]\) 为主谱核的最小多项式，并满足
\[
\mathrm{Disc}_{\lambda}\bigl(\Pi(\lambda,y)\bigr)=-y(y-1)P_{\mathrm{LY}}(y).
\]
在该可编译判别式口径下，分歧值、诱导奇点与零点凝聚曲线均由显式代数对象给出；进一步在 \(s=\log y\) 坐标可得到代数量与对数超越量之间的严格分层。

\begin{theorem}[简单特征值解析分支的代数喷射封闭性]\label{thm:xi-leyang-simple-root-algebraic-jet-closure}
\leanverified{paper\_xi\_leyang\_simple\_root\_algebraic\_jet\_closure}
设 \(K\subset\mathbb C\) 为数域，\(M(t)\in\mathrm{Mat}_{n}(K[t])\)，并记
\[
\chi(\lambda,t):=\det(\lambda I_n-M(t))\in K[\lambda,t].
\]
若存在 \((\lambda_0,t_0)\in\overline K\times K\) 使
\[
\chi(\lambda_0,t_0)=0,\qquad
\partial_{\lambda}\chi(\lambda_0,t_0)\neq 0,
\]
则在 \(t_0\) 的某邻域内存在唯一解析函数 \(\lambda(t)\) 满足
\[
\lambda(t_0)=\lambda_0,\qquad
\chi(\lambda(t),t)\equiv 0.
\]
并且对任意整数 \(m\ge 0\)，皆有
\[
\lambda^{(m)}(t_0)\in K(\lambda_0)\subset\overline K.
\]
\end{theorem}
\begin{proof}
由隐函数定理，\(\partial_{\lambda}\chi(\lambda_0,t_0)\neq0\) 蕴含上述解析支存在且唯一。
对恒等式 \(\chi(\lambda(t),t)\equiv0\) 求导并在 \(t=t_0\) 处代入，得
\[
\chi_{\lambda}(\lambda_0,t_0)\lambda'(t_0)+\chi_t(\lambda_0,t_0)=0,
\]
从而
\[
\lambda'(t_0)
=-\frac{\chi_t(\lambda_0,t_0)}{\chi_{\lambda}(\lambda_0,t_0)}
\in K(\lambda_0).
\]
对 \(m\ge2\) 重复求导，\(\lambda^{(m)}(t_0)\) 的线性系数始终为 \(\chi_{\lambda}(\lambda_0,t_0)\neq0\)，右端仅含 \(\chi\) 在 \((\lambda_0,t_0)\) 的偏导值与低阶导数 \(\lambda^{(j)}(t_0)\)（\(j<m\)）的有理组合。
以 \(m\) 归纳即得 \(\lambda^{(m)}(t_0)\in K(\lambda_0)\)。证毕。
\end{proof}

\begin{corollary}[有限核热力学的导数代数性与熵项超越性]\label{cor:xi-leyang-finite-kernel-algebraic-vs-transcendental-layering}
设 \(\lambda_+(y)\) 为某有限核转移矩阵的主特征值解析支，且满足代数方程 \(\Pi(\lambda_+(y),y)=0\)。
若在 \(y=1\) 处
\[
\partial_{\lambda}\Pi\!\bigl(\lambda_+(1),1\bigr)\neq0,
\]
并定义
\[
\psi(s):=\log\lambda_+(e^s)-\log2,
\]
则对任意 \(m\ge1\)，有 \(\psi^{(m)}(0)\in\overline{\mathbb Q}\)。
此外，若
\[
\frac{\lambda_+(1)}{2}\in\overline{\mathbb Q}_{>0}\setminus\{1\},
\]
则
\[
\psi(0)=\log\!\frac{\lambda_+(1)}{2}
\]
为超越数。
特别地，记
\[
\kappa_{\infty}:=\frac{6\sqrt5-5}{250},
\qquad
c_1:=\log\!\frac{2}{\varphi},
\qquad
\varphi:=\frac{1+\sqrt5}{2},
\]
则 \(\kappa_{\infty}\in\overline{\mathbb Q}\)、\(c_1\) 为超越数，且对任意 \(r\in\overline{\mathbb Q}^{\times}\) 均有
\[
\kappa_{\infty}\neq r\,c_1^{2},
\qquad
\frac{\kappa_{\infty}}{c_1^{2}}\ \text{为超越数}.
\]
\end{corollary}
\begin{proof}
把定理 \ref{thm:xi-leyang-simple-root-algebraic-jet-closure} 应用于 \(\chi(\lambda,y)=\Pi(\lambda,y)\) 在 \(y=1\) 的简单根分支，得 \(\lambda_+^{(m)}(1)\in\overline{\mathbb Q}\)（\(m\ge0\)）。
与 \(y=e^s\) 复合后，\(\lambda_+(e^s)\) 在 \(s=0\) 的各阶导数仍由代数数经有理运算给出，故 \(\psi^{(m)}(0)\in\overline{\mathbb Q}\)（\(m\ge1\)）。
若 \(\lambda_+(1)/2\in\overline{\mathbb Q}_{>0}\setminus\{1\}\)，则 \(\log(\lambda_+(1)/2)\) 由 Lindemann--Weierstrass 定理为超越数 \cite{Baker1975TranscendentalNumberTheory}。
又 \(\kappa_{\infty}\) 显然代数，而 \(2/\varphi\in\overline{\mathbb Q}_{>0}\setminus\{1\}\) 使 \(c_1\) 超越，故 \(c_1^2\) 亦超越。
于是 \(r c_1^2\)（\(r\in\overline{\mathbb Q}^{\times}\)）超越，不可能等于代数数 \(\kappa_{\infty}\)。
若 \(\kappa_{\infty}/c_1^2\) 代数，则 \(c_1^2=\kappa_{\infty}/(\kappa_{\infty}/c_1^2)\) 代数，矛盾。证毕。
\end{proof}

\begin{theorem}[Lee--Yang 分歧值在对数坐标下的必然超越性]\label{thm:xi-leyang-log-singularity-transcendence}
设 \(y_{\star}\in\overline{\mathbb Q}\setminus\{0,1\}\) 为 \(\lambda_+(y)\) 的任一分歧值（等价地，\(\mathrm{Disc}_{\lambda}(\Pi)\) 的零点之一）。
取任意对数分支并记
\[
s_{\star}\in\mathbb C,\qquad e^{s_{\star}}=y_{\star}.
\]
则 \(s_{\star}\) 必为超越数。
特别地，在该文档给出的最近复奇点
\[
y_{\mathrm c}^{\pm}=\frac{3\pm i\sqrt{111}}{96}
\]
处，任一
\[
s_{\mathrm c}^{\pm}=\log y_{\mathrm c}^{\pm}
\]
均为超越数，且
\[
\Re(s_{\mathrm c}^{+})
=\log|y_{\mathrm c}^{+}|
=\log\sqrt{\frac{5}{384}}
\]
为超越实数。
\end{theorem}
\begin{proof}
反设 \(s_{\star}\) 为代数数。
若 \(s_{\star}\neq0\)，则由 Lindemann--Weierstrass 定理，\(e^{s_{\star}}\) 必为超越数，与 \(e^{s_{\star}}=y_{\star}\in\overline{\mathbb Q}\) 矛盾。
若 \(s_{\star}=0\)，则 \(y_{\star}=1\)，与 \(y_{\star}\neq1\) 矛盾。
故 \(s_{\star}\) 超越。
对 \(y_{\mathrm c}^{\pm}\) 代入同理。
又 \(|y_{\mathrm c}^{+}|=\sqrt{5/384}\in\overline{\mathbb Q}_{>0}\setminus\{1\}\)，故其对数由同一定理为超越实数。证毕。
\end{proof}

\begin{proposition}[Lee--Yang 振荡角频率与 \texorpdfstring{$\pi$}{pi} 的不可公度性]\label{prop:xi-leyang-angle-over-pi-irrational}
令
\[
y_{\mathrm c}^{+}:=\frac{3+i\sqrt{111}}{96},
\qquad
\theta:=\arg(y_{\mathrm c}^{+})\in(0,\pi).
\]
则
\[
\frac{\theta}{\pi}\notin\mathbb Q.
\]
\end{proposition}
\begin{proof}
反设 \(\theta/\pi\in\mathbb Q\)。
则 \(\zeta:=e^{i\theta}\) 为单位根，且
\[
\zeta=\frac{y_{\mathrm c}^{+}}{|y_{\mathrm c}^{+}|}
=\frac{3+i\sqrt{111}}{2\sqrt{30}}
\in F:=\mathbb Q(\sqrt{-111},\sqrt{30}).
\]
域 \(F/\mathbb Q\) 为双二次扩张，\(\mathrm{Gal}(F/\mathbb Q)\cong V_4\)。
因此若 \(\zeta\) 的阶为 \(n\)，则 \(\mathbb Q(\zeta_n)\subseteq F\)，其伽罗瓦群 \((\mathbb Z/n\mathbb Z)^{\times}\) 必是 \(V_4\) 的子商，故指数至多为 \(2\) 且次数 \(\varphi(n)\le4\)。
可得 \(n\in\{1,2,3,4,6,8,12\}\)。
然而 \(F\) 的三个二次子域仅为
\[
\mathbb Q(\sqrt{-111}),\qquad
\mathbb Q(\sqrt{30}),\qquad
\mathbb Q(\sqrt{-3330}),
\]
其中不含 \(\mathbb Q(i)\) 与 \(\mathbb Q(\sqrt{-3})\)。
故 \(n\notin\{3,4,6,8,12\}\)，只能 \(n\in\{1,2\}\)，即 \(\zeta=\pm1\)。
但 \(y_{\mathrm c}^{+}\) 非实，故 \(\zeta\neq\pm1\)，矛盾。
因此 \(\theta/\pi\notin\mathbb Q\)。证毕。
\end{proof}

\begin{theorem}[平方耦合关系的代数--超越障碍]\label{thm:xi-leyang-kappa-vs-c1-square-obstruction}
令
\[
\kappa_{\infty}=\frac{6\sqrt5-5}{250},
\qquad
c_1=\log\!\frac{2}{\varphi},
\qquad
\varphi=\frac{1+\sqrt5}{2}.
\]
则对任意 \(r\in\overline{\mathbb Q}^{\times}\) 均有
\[
\kappa_{\infty}\neq r\,c_1^2.
\]
并且
\[
\frac{\kappa_{\infty}}{c_1^2}
\]
为超越数。
\end{theorem}
\begin{proof}
\(\kappa_{\infty}\) 为代数数。
由 \(2/\varphi\in\overline{\mathbb Q}_{>0}\setminus\{1\}\)，\(c_1=\log(2/\varphi)\) 为超越数，故 \(c_1^2\) 超越。
于是任意 \(r\in\overline{\mathbb Q}^{\times}\) 的乘积 \(r c_1^2\) 仍超越，不能等于代数数 \(\kappa_{\infty}\)。
若 \(\kappa_{\infty}/c_1^2\) 代数，则 \(c_1^2=\kappa_{\infty}/(\kappa_{\infty}/c_1^2)\) 代数，矛盾。证毕。
\end{proof}

\begin{corollary}[代数输运常数与对数熵常数之间不存在闭合平方律]\label{cor:xi-leyang-algebraic-transport-vs-log-entropy-no-square-law}
在有限核谱多项式框架中，凡可写为“简单特征支局部喷射经有理运算与有限代数消元”所得的输运常数，均为代数数。
另一方面，对任意
\[
\alpha\in\overline{\mathbb Q}_{>0}\setminus\{1\},
\qquad
c:=\log\alpha,
\]
有 \(c^2\) 超越。
因此对任意代数常数 \(T\in\overline{\mathbb Q}\) 与任意 \(r\in\overline{\mathbb Q}^{\times}\)，恒有
\[
T\neq r\,c^2.
\]
\end{corollary}
\begin{proof}
第一部分由定理 \ref{thm:xi-leyang-simple-root-algebraic-jet-closure} 的导数闭合性与代数消元稳定性给出。
第二部分由 Lindemann--Weierstrass 定理给出 \(c\) 超越，进而 \(c^2\) 超越。
代数数不可能等于上述超越数，结论成立。证毕。
\end{proof}

\begin{conjecture}[对数奇点常数的超越独立性]\label{conj:xi-leyang-log-singularity-transcendence-independence}
在同一归一化下，三实常数
\[
\pi,\qquad
\log\!\frac{2}{\varphi},\qquad
\log\sqrt{\frac{5}{384}}
\]
在 \(\mathbb Q\) 上代数独立。
该猜想若成立，则 Lee--Yang 端点诱导的振荡频率与熵率常数在超越层面不存在任何非平凡代数耦合。
\end{conjecture}

\paragraph{与深度/素数接口的兼容性}
定理 \ref{thm:conclusion-rate-cdim-budget-inequality} 与定理 \ref{thm:cdim-budget-completeness}
给出的预算与剖面下界，与命题 \ref{prop:xi-koenigs-prime-generator-spectral-factorization} 及推论 \ref{cor:xi-joukowsky-prime-register-ellipse-parameterization} 相容：相位圆因子、素数寄存器与 Koenigs 深度轴在结构层分别对应自由秩、外置账本秩与平移尺度，不发生维度重复计费。
\endinput


\paragraph{Jensen 缺陷的递归帽核表象与误差鲁棒证书链}
本段在附录命题 \ref{prop:app-jensen-formula-defect} 的 Jensen 求和公式口径下，
把圆盘解析函数的零点半径分布改写为可递归细化的二阶差分反演问题。
核心对象是对数半径坐标中的 Jensen 缺陷序列，其局部二阶差分给出窗口化“帽核”积分，从而可在噪声下形成可执行的存在性与重数证书。

\begin{theorem}[Jensen 缺陷的帽核二阶差分表示]\label{thm:xi-jensen-defect-hat-kernel-second-difference}
设 \(F\) 在单位圆盘 \(\mathbb D\) 内解析，且 \(F(0)\neq 0\)。
定义 Jensen 缺陷
\[
D_F(r):=\frac{1}{2\pi}\int_0^{2\pi}\log|F(re^{i\theta})|\,\dd\theta-\log|F(0)|,
\qquad 0<r<1.
\]
令 \(x:=\log r\in(-\infty,0)\)，并记
\[
J_F(x):=D_F(e^x).
\]
设 \(F\) 在 \(\mathbb D\) 内的零点为 \(\{a_k\}\)（按重数 \(m_k\) 计），定义
\[
t_k:=\log|a_k|\in(-\infty,0),
\qquad
\mu_F:=\sum_k m_k\,\delta_{t_k}.
\]
则对任意 \(h>0\) 与 \(x\pm h<0\)，二阶差分
\[
\Delta_h^2J_F(x):=J_F(x+h)-2J_F(x)+J_F(x-h)
\]
满足精确恒等式
\[
\boxed{\;
\Delta_h^2J_F(x)=\int_{\mathbb R}(h-|x-t|)_+\,\dd\mu_F(t)
=\sum_k m_k\,(h-|x-t_k|)_+\; }.
\]
其中 \((u)_+:=\max\{u,0\}\)。
特别地，
\[
\Delta_h^2J_F(x)=0
\quad\Longleftrightarrow\quad
\mu_F\bigl((x-h,x+h)\bigr)=0.
\]
\leanverified{paper\_xi\_jensen\_defect\_hat\_kernel\_second\_difference}
\end{theorem}
\begin{proof}
由命题 \ref{prop:app-jensen-formula-defect}，
\[
D_F(r)=\sum_{|a_k|<r}m_k\log\frac{r}{|a_k|}.
\]
代入 \(r=e^x\) 得
\[
J_F(x)=\sum_k m_k\,(x-t_k)_+.
\]
对固定 \(x\pm h<0\)，仅有有限个 \(|a_k|<e^{x+h}\) 贡献非零项，故可逐项作二阶差分：
\[
\Delta_h^2J_F(x)
=\sum_k m_k\Bigl[(x+h-t_k)_+-2(x-t_k)_++(x-h-t_k)_+\Bigr].
\]
记 \(\phi_t(x):=(x-t)_+\)。
分段计算可得
\[
\phi_t(x+h)-2\phi_t(x)+\phi_t(x-h)=(h-|x-t|)_+.
\]
故得主恒等式。
又 \((h-|x-t|)_+\ge 0\)，且其为零当且仅当 \(|x-t|\ge h\)。
于是 \(\Delta_h^2J_F(x)=0\) 当且仅当窗口 \((x-h,x+h)\) 内无原子质量，即
\(\mu_F((x-h,x+h))=0\)。证毕。
\end{proof}

\begin{corollary}[带噪声三点采样下的存在性证书与重数下界]\label{cor:xi-jensen-defect-noisy-threepoint-certificate}
在定理 \ref{thm:xi-jensen-defect-hat-kernel-second-difference} 设定下，
若仅能观测到近似值
\[
\widetilde J(x-h),\ \widetilde J(x),\ \widetilde J(x+h),
\]
并满足
\[
|\widetilde J(\xi)-J_F(\xi)|\le \varepsilon,
\qquad
\xi\in\{x-h,x,x+h\},
\]
定义可计算量
\[
\widetilde W_h(x):=\widetilde J(x+h)-2\widetilde J(x)+\widetilde J(x-h).
\]
则
\[
\bigl|\widetilde W_h(x)-\Delta_h^2J_F(x)\bigr|\le 4\varepsilon.
\]
进一步成立：
\begin{enumerate}
  \item 若 \(\widetilde W_h(x)>4\varepsilon\)，则
  \[
  \mu_F((x-h,x+h))\ge 1,
  \]
  即窗口对应半径带内必有盘内零点；
  \item 记窗口内总重数
  \[
  M_{x,h}:=\mu_F((x-h,x+h)),
  \]
  则
  \[
  M_{x,h}\ge
  \max\!\left\{0,\left\lceil\frac{\widetilde W_h(x)-4\varepsilon}{h}\right\rceil\right\}.
  \]
\end{enumerate}
\end{corollary}
\begin{proof}
误差估计由三角不等式：
\[
\begin{aligned}
|\widetilde W_h-\Delta_h^2J_F|
&\le |\widetilde J(x+h)-J_F(x+h)|
+2|\widetilde J(x)-J_F(x)|\\
&\quad +|\widetilde J(x-h)-J_F(x-h)|
\le 4\varepsilon.
\end{aligned}
\]
若 \(\widetilde W_h>4\varepsilon\)，则 \(\Delta_h^2J_F\ge \widetilde W_h-4\varepsilon>0\)。
由定理 \ref{thm:xi-jensen-defect-hat-kernel-second-difference}，窗口内必有原子，故第一条成立。
第二条由
\[
\Delta_h^2J_F
=\sum_{t_k\in(x-h,x+h)}m_k(h-|x-t_k|)\le h\,M_{x,h}
\]
与 \(\Delta_h^2J_F\ge \widetilde W_h-4\varepsilon\) 合并得到。证毕。
\end{proof}

\begin{theorem}[递归细化中的误差--分辨率耦合阈值]\label{thm:xi-jensen-recursive-resolution-noise-coupling}
在定理 \ref{thm:xi-jensen-defect-hat-kernel-second-difference} 框架下，考虑以
\[
\widetilde W_h(x)=\widetilde J(x+h)-2\widetilde J(x)+\widetilde J(x-h)
\]
作为递归触发量的中心化判据（以 \(\widetilde W_h(x)>0\) 作为下一层细化触发条件）。
若目标情形为“窗口中心 \(x\) 处存在单重零点”（即 \(\mu_F\) 在 \(t=x\) 处有一个单位原子），则：
\begin{enumerate}
  \item 触发充分条件为
  \[
  \varepsilon<\frac{h}{4}.
  \]
  在该条件下，对所有满足误差界的噪声实现，均有 \(\widetilde W_h(x)>0\)；
  \item 该门槛在上述触发规则下是尖锐的：若 \(\varepsilon\ge h/4\)，则存在合法样本与合法噪声使 \(\widetilde W_h(x)\le 0\)，从而本层触发可失败；
  \item 因此在递归细化 \(h\mapsto h/2\) 下，若保持同一触发机制，则噪声预算须同步线性收缩
  \[
  \varepsilon\mapsto \varepsilon/2.
  \]
\end{enumerate}
\end{theorem}
\begin{proof}
由定理 \ref{thm:xi-jensen-defect-hat-kernel-second-difference}，中心单重零点在尺度 \(h\) 的理想信号为
\[
\Delta_h^2J_F(x)=(h-|x-x|)_+=h.
\]
由推论 \ref{cor:xi-jensen-defect-noisy-threepoint-certificate}，
\[
\widetilde W_h(x)=\Delta_h^2J_F(x)+\eta_h,\qquad |\eta_h|\le 4\varepsilon.
\]
故最坏情形下
\[
\widetilde W_h(x)\ge h-4\varepsilon.
\]
若 \(\varepsilon<h/4\)，则 \(h-4\varepsilon>0\)，得到第一条。

为证尖锐性，取 \(F_1(w):=1-w/a\) 且 \(a=e^x\in(0,1)\)，则 \(F_1(0)=1\)，并在 \(w=a\) 处有单重零点。
由 Jensen 公式
\[
J_{F_1}(x-h)=0,\quad J_{F_1}(x)=0,\quad J_{F_1}(x+h)=h,
\]
因此 \(\Delta_h^2J_{F_1}(x)=h\)。
当 \(\varepsilon\ge h/4\) 时，取允许噪声 \(\eta_h=-4\varepsilon\)，得到
\[
\widetilde W_h(x)=h-4\varepsilon\le 0,
\]
触发失败，证得第二条。

第三条由第一条在尺度 \(h_n=h_0\,2^{-n}\) 上逐层套用：
\(
\varepsilon_n<h_n/4
\)
等价于 \(\varepsilon_{n+1}<\varepsilon_n/2\)。证毕。
\end{proof}

\begin{theorem}[全域误差下的边界层不可判别性]\label{thm:xi-jensen-boundary-layer-indistinguishability}
设
\[
\mathcal A:=\{F:\mathbb D\to\mathbb C\ \text{解析},\ F(0)=1\}.
\]
给定任意 \(\varepsilon>0\) 与任意 \(\eta\in(0,\varepsilon]\)，存在 \(F_0,F_1\in\mathcal A\) 满足：
\begin{enumerate}
  \item \(F_0\) 在 \(\mathbb D\) 内无零点；
  \item \(F_1\) 在 \(\mathbb D\) 内有单重零点 \(a\)，且 \(|a|=e^{-\eta}\)；
  \item 两者 Jensen 缺陷在全半径上不可分辨：
  \[
  \sup_{0<r<1}|D_{F_1}(r)-D_{F_0}(r)|=\eta\le\varepsilon.
  \]
\end{enumerate}
因此在统一加性误差 \(\varepsilon\) 下，任何仅基于 \(D_F(r)\) 全域观测的确定性核验流程都不能排除 \(\log\)-深度不超过 \(\varepsilon\) 的边界层内部零点。
\end{theorem}
\begin{proof}
取
\[
F_0(w)\equiv 1,\qquad
F_1(w):=1-\frac{w}{a},
\qquad |a|=e^{-\eta}.
\]
则 \(F_0,F_1\in\mathcal A\)，且 \(F_1\) 在 \(w=a\) 处有单重零点。
由 Jensen 公式，
\[
D_{F_1}(r)=
\begin{cases}
0, & 0<r\le |a|,\\[1mm]
\log\!\frac{r}{|a|}, & |a|<r<1.
\end{cases}
\]
于是
\[
\sup_{0<r<1}D_{F_1}(r)=\log\frac{1}{|a|}=\eta.
\]
又 \(D_{F_0}(r)\equiv 0\)，故
\[
\sup_{0<r<1}|D_{F_1}(r)-D_{F_0}(r)|=\eta\le\varepsilon.
\]
结论成立。证毕。
\end{proof}

\paragraph{与圆盘化 RH 核验链的对齐}
由附录定理 \ref{thm:app-rh-iff-disk-zero-free}，完成后的 \(\xi\)-函数经右半平面到圆盘的双全纯变换后，
“右半平面离线零点存在性”与“圆盘内零点存在性”等价。
定理 \ref{thm:xi-jensen-defect-hat-kernel-second-difference}--\ref{thm:xi-jensen-boundary-layer-indistinguishability}
因此把该判定问题转化为：
\[
\text{零点计数测度}\ \mu_F
\quad\Longleftrightarrow\quad
\text{Jensen 缺陷的帽核二阶差分多尺度反演}
\]
并给出显式噪声门槛与边界层硬下界，从而把“零点递归—误差控制”写成可审计的尺度耦合律。

\endinput


\paragraph{混合位置--尺度 Poisson 核粗粒化下 \texorpdfstring{$f$}{f}-散度的二阶矩主导律}
设缺陷测度为
\[
\nu=\sum_{j=1}^{M}m_j\,\delta_{(\gamma_j,\delta_j)},
\qquad
m_j>0,\ \delta_j>0,
\]
并记 Poisson 核
\[
P_s(x):=\frac{1}{\pi}\frac{s}{x^2+s^2},
\qquad s>0.
\]
定义 Poisson 场
\[
G_t(x):=\int_{\mathbb R\times(0,\infty)}P_{t+\delta}(x-\gamma)\,\dd\nu(\gamma,\delta)
=\sum_{j=1}^{M}m_j\,P_{t+\delta_j}(x-\gamma_j).
\]
令总质量
\[
M_0:=\nu\bigl(\mathbb R\times(0,\infty)\bigr)=\sum_{j=1}^{M}m_j,
\]
并定义归一化密度
\[
h_t(x):=\frac{G_t(x)}{M_0}
=\int P_{t+\delta}(x-\gamma)\,\dd\mu(\gamma,\delta),
\qquad
\mu:=\frac{\nu}{M_0}.
\]
记
\[
\bar\gamma:=\int \gamma\,\dd\mu,\qquad
\bar\delta:=\int \delta\,\dd\mu,\qquad
s_t:=t+\bar\delta,
\]
及基线核
\[
g_t(x):=P_{s_t}(x-\bar\gamma).
\]
定义复二阶中心矩
\[
Z:=(\gamma-\bar\gamma)+i(\delta-\bar\delta),
\qquad
m_2:=\mathbb E_\mu[Z^2]\in\mathbb C.
\]
对满足 \(f(1)=0\)、在 \(1\) 邻域 \(C^2\) 且凸的函数 \(f\)，记
\[
D_f(p\mid q):=\int_{\mathbb R}q(x)\,f\!\left(\frac{p(x)}{q(x)}\right)\dd x.
\]

\begin{theorem}[混合位置--尺度 Poisson 粗粒化的 \texorpdfstring{$f$}{f}-散度四次锐渐近]\label{thm:xi-poisson-cauchy-f-divergence-sharp-asymptotic}
若
\[
\int_{\mathbb R\times(0,\infty)}\bigl(|\gamma|^3+|\delta|^3\bigr)\,\dd\mu(\gamma,\delta)<\infty,
\]
则当 \(t\to\infty\) 时，
\[
D_f(h_t\mid g_t)
=\frac{f''(1)}{8}\frac{|m_2|^2}{s_t^4}
+o\!\left(s_t^{-4}\right).
\]
等价地，
\[
\lim_{t\to\infty}\frac{s_t^4}{|m_2|^2}D_f(h_t\mid g_t)=\frac{f''(1)}{8}.
\]
\end{theorem}
\begin{proof}
记
\[
P(x,s):=\frac{1}{\pi}\frac{s}{x^2+s^2}.
\]
固定 \(x\)，考虑
\[
\Phi_{t,x}(\gamma,\delta):=P(x-\gamma,t+\delta).
\]
在点 \((\bar\gamma,\bar\delta)\) 对 \((\gamma,\delta)\) 作二阶 Taylor 展开并对 \(\mu\) 取期望。
由中心化条件
\(
\mathbb E_\mu(\gamma-\bar\gamma)=\mathbb E_\mu(\delta-\bar\delta)=0
\)
可知一阶项消失，从而
\[
h_t(x)=g_t(x)
+\frac{1}{2}\Var_\mu(\gamma)\,\partial_{xx}P
-\Cov_\mu(\gamma,\delta)\,\partial_{xs}P
+\frac{1}{2}\Var_\mu(\delta)\,\partial_{ss}P
+R_t(x),
\]
其中各偏导均在 \((x-\bar\gamma,s_t)\) 处取值。
由 Poisson 核三阶导数尺度 \(\partial^{(3)}P=O(s_t^{-4})\) 与三阶矩有界，得到
\[
\left\|\frac{R_t}{g_t}\right\|_{L^\infty(\mathbb R)}=O(s_t^{-3}).
\]

Poisson 核在上半平面满足调和关系
\[
\partial_{xx}P+\partial_{ss}P=0,
\]
故二阶主项化简为
\[
h_t-g_t
=\frac{1}{2}\bigl(\Var_\mu(\gamma)-\Var_\mu(\delta)\bigr)\partial_{xx}P
-\Cov_\mu(\gamma,\delta)\partial_{xs}P
+R_t.
\]
记
\[
\alpha:=\Var_\mu(\gamma)-\Var_\mu(\delta),\qquad
\beta:=-\Cov_\mu(\gamma,\delta),\qquad
y:=\frac{x-\bar\gamma}{s_t}.
\]
则
\[
g_t(x)=\frac{1}{\pi s_t}\frac{1}{1+y^2},
\]
并且直接计算可得
\[
\frac{\partial_{xx}P}{g_t}
=\frac{2}{s_t^2}\frac{3y^2-1}{(1+y^2)^2},
\qquad
\frac{\partial_{xs}P}{g_t}
=\frac{1}{s_t^2}\frac{2y(3-y^2)}{(1+y^2)^2}.
\]
定义
\[
u_1(y):=\frac{3y^2-1}{(1+y^2)^2}\quad(\text{偶函数}),\qquad
u_2(y):=\frac{2y(3-y^2)}{(1+y^2)^2}\quad(\text{奇函数}),
\]
以及相对扰动
\[
\varepsilon_t(x):=\frac{h_t(x)}{g_t(x)}-1.
\]
于是
\[
\varepsilon_t(x)
=\frac{1}{s_t^2}\bigl(\alpha u_1(y)+\beta u_2(y)\bigr)
+O(s_t^{-3})
\quad\text{于 }L^\infty(\mathbb R),
\]
故 \(\varepsilon_t\to0\) 一致成立。

由 \(f\in C^2\)（在 \(1\) 邻域）有一致展开
\[
f(1+u)=f(1)+f'(1)u+\frac{f''(1)}{2}u^2+o(u^2)\qquad(u\to0).
\]
又
\[
\int_{\mathbb R}g_t\varepsilon_t\,\dd x
=\int_{\mathbb R}(h_t-g_t)\,\dd x=0,
\]
故线性项积分为零，从而
\[
D_f(h_t\mid g_t)
=\frac{f''(1)}{2}\int_{\mathbb R}g_t(x)\varepsilon_t(x)^2\,\dd x
+o\!\left(\int g_t\varepsilon_t^2\right).
\]
令
\[
G(y):=\frac{1}{\pi(1+y^2)},
\qquad
x=\bar\gamma+s_t y,
\]
则 \(g_t(x)\dd x=G(y)\dd y\)，并且
\[
\int g_t\varepsilon_t^2
=\frac{1}{s_t^4}\int_{\mathbb R}G(y)\bigl(\alpha u_1(y)+\beta u_2(y)\bigr)^2\dd y
+o(s_t^{-4}).
\]
由奇偶性，交叉项积分为零。
再由 Beta 积分计算可得
\[
\int_{\mathbb R}G(y)u_1(y)^2\,\dd y=\frac14,
\qquad
\int_{\mathbb R}G(y)u_2(y)^2\,\dd y=1.
\]
故
\[
\int g_t\varepsilon_t^2
=\frac{1}{s_t^4}\left(\frac{\alpha^2}{4}+\beta^2\right)
+o(s_t^{-4}),
\]
进而
\[
D_f(h_t\mid g_t)
=\frac{f''(1)}{8}\frac{\alpha^2+4\beta^2}{s_t^4}
+o(s_t^{-4}).
\]
另一方面，
\[
\begin{aligned}
m_2
&=\mathbb E_\mu\!\left[\bigl((\gamma-\bar\gamma)+i(\delta-\bar\delta)\bigr)^2\right]\\
&=\Var_\mu(\gamma)-\Var_\mu(\delta)+2i\,\Cov_\mu(\gamma,\delta)
=\alpha-2i\beta,
\end{aligned}
\]
故
\[
|m_2|^2=\alpha^2+4\beta^2.
\]
代回即得结论。证毕。
\end{proof}

\begin{corollary}[KL 耗散五次律与分数 Fisher 信息首项]\label{cor:xi-poisson-cauchy-kl-fractional-fisher-fifth-law}
在定理 \ref{thm:xi-poisson-cauchy-f-divergence-sharp-asymptotic} 的设定下，取
\[
f(u)=u\log u,
\]
则
\[
D_{\mathrm{KL}}(h_t\mid g_t)
=\frac{|m_2|^2}{8s_t^4}
+o(s_t^{-4}).
\]
若记沿 Poisson 半群轨道的 KL 耗散率为分数 Fisher 信息 \(I_1\)，即
\[
\frac{\dd}{\dd t}D_{\mathrm{KL}}(p_t\mid q_t)=-I_1(p_t\mid q_t),
\qquad p_t=P_t*p,\ q_t=P_t*q,
\]
则本设定下
\[
I_1(h_t\mid g_t)
=\frac{|m_2|^2}{2s_t^5}
+o(s_t^{-5}).
\]
\end{corollary}
\begin{proof}
由 \(f''(1)=1\) 即得 KL 的四次渐近式。
又由卷积可加性 \(P_{t+\delta}=P_t*P_\delta\)，有
\[
h_t=P_t*h_0,\qquad
g_t=P_t*g_0,
\]
其中
\[
h_0(x)=\int P_\delta(x-\gamma)\,\dd\mu(\gamma,\delta),
\qquad
g_0(x)=P_{\bar\delta}(x-\bar\gamma).
\]
故 \((h_t,g_t)\) 同属一条 Poisson 半群推进轨道，耗散恒等式可用。
记
\[
D_{\mathrm{KL}}(h_t\mid g_t)=\frac{|m_2|^2}{8s_t^4}+r_t,\qquad r_t=o(s_t^{-4}),
\]
并利用 Poisson 平滑给出的可微性，可得
\[
\frac{\dd}{\dd t}D_{\mathrm{KL}}(h_t\mid g_t)
=-\frac{|m_2|^2}{2s_t^5}+o(s_t^{-5}).
\]
代入
\(
\frac{\dd}{\dd t}D_{\mathrm{KL}}=-I_1
\)
即得结论。证毕。
\end{proof}

\begin{proposition}[超四次衰减的伪协方差刚性]\label{prop:xi-poisson-cauchy-superquartic-rigidity}
在定理 \ref{thm:xi-poisson-cauchy-f-divergence-sharp-asymptotic} 的设定下，下列命题等价：
\begin{enumerate}
  \item
  \[
  D_{\mathrm{KL}}(h_t\mid g_t)=o(s_t^{-4})\qquad (t\to\infty);
  \]
  \item
  \[
  m_2=\mathbb E_\mu\!\left[\bigl((\gamma-\bar\gamma)+i(\delta-\bar\delta)\bigr)^2\right]=0.
  \]
\end{enumerate}
进一步地，若 \(m_2=0\) 且三阶矩有限，则
\[
D_{\mathrm{KL}}(h_t\mid g_t)=O(s_t^{-6}).
\]
\end{proposition}
\begin{proof}
由推论 \ref{cor:xi-poisson-cauchy-kl-fractional-fisher-fifth-law}，
\[
D_{\mathrm{KL}}(h_t\mid g_t)
=\frac{|m_2|^2}{8s_t^4}+o(s_t^{-4}),
\]
故两条命题等价。

当 \(m_2=0\) 时，二阶主项整体消失。
由定理 \ref{thm:xi-poisson-cauchy-f-divergence-sharp-asymptotic} 的展开结构与余项估计，
\[
\left\|\frac{h_t-g_t}{g_t}\right\|_{L^\infty}=O(s_t^{-3}).
\]
记 \(\varepsilon_t=h_t/g_t-1\)，则
\[
\|\varepsilon_t\|_{L^\infty}=O(s_t^{-3}).
\]
利用
\[
(1+u)\log(1+u)-u=O(u^2)\qquad(u\to0),
\]
得到
\[
\begin{aligned}
D_{\mathrm{KL}}(h_t\mid g_t)
&=\int g_t\Bigl((1+\varepsilon_t)\log(1+\varepsilon_t)-\varepsilon_t\Bigr)\dd x\\
&=O\!\left(\int g_t\varepsilon_t^2\,\dd x\right)
=O(s_t^{-6}).
\end{aligned}
\]
证毕。
\end{proof}

\begin{theorem}[端点 \texorpdfstring{$-1$}{-1} 的 Busemann--时间纤维替换下的边界层四次律]\label{thm:xi-poisson-cauchy-busemann-boundary-fourth-law}
设
\[
B^{-1}(w)=\frac{1-|w|^2}{|1+w|^2}=\Im\!\bigl(C^{-1}(w)\bigr),
\]
并令 \(w\to-1\) 非切向逼近边界。
在定理 \ref{thm:xi-poisson-cauchy-f-divergence-sharp-asymptotic} 的假设下，对任意满足条件的凸函数 \(f\)，有
\[
D_f\!\left(h_{B^{-1}(w)}\mid g_{B^{-1}(w)}\right)
=\frac{f''(1)}{8}\frac{|m_2|^2}{\bigl(B^{-1}(w)+\bar\delta\bigr)^4}
+o\!\left(\bigl(B^{-1}(w)\bigr)^{-4}\right).
\]
\end{theorem}
\begin{proof}
由边界几何恒等式可知非切向 \(w\to-1\) 蕴含 \(B^{-1}(w)\to\infty\)。
将定理 \ref{thm:xi-poisson-cauchy-f-divergence-sharp-asymptotic} 的时间参数 \(t\) 以
\[
t=B^{-1}(w)
\]
替换，即得结论。证毕。
\end{proof}

\begin{theorem}[双缺陷构型在强粗粒化下的可分辨首项]\label{thm:xi-poisson-cauchy-two-mixtures-distinguishability}
设 \(\mu,\mu'\) 为 \(\mathbb R\times(0,\infty)\) 上概率测度，且满足三阶矩有限，并具有一致的一阶矩
\[
\int \gamma\,\dd\mu=\int \gamma\,\dd\mu'=\bar\gamma,
\qquad
\int \delta\,\dd\mu=\int \delta\,\dd\mu'=\bar\delta.
\]
定义
\[
h_t(x):=\int P_{t+\delta}(x-\gamma)\,\dd\mu(\gamma,\delta),
\qquad
h_t'(x):=\int P_{t+\delta}(x-\gamma)\,\dd\mu'(\gamma,\delta),
\]
以及
\[
m_2:=\int\bigl((\gamma-\bar\gamma)+i(\delta-\bar\delta)\bigr)^2\,\dd\mu,
\qquad
m_2':=\int\bigl((\gamma-\bar\gamma)+i(\delta-\bar\delta)\bigr)^2\,\dd\mu'.
\]
则对任意满足定理 \ref{thm:xi-poisson-cauchy-f-divergence-sharp-asymptotic} 条件的凸函数 \(f\)，当 \(t\to\infty\) 时
\[
D_f(h_t\mid h_t')
=\frac{f''(1)}{8}\frac{|m_2-m_2'|^2}{(t+\bar\delta)^4}
+o(t^{-4}).
\]
特别地，
\[
D_{\mathrm{KL}}(h_t\mid h_t')
=\frac{1}{8}\frac{|m_2-m_2'|^2}{(t+\bar\delta)^4}
+o(t^{-4}).
\]
\end{theorem}
\begin{proof}
设共同基线
\[
g_t(x):=P_{t+\bar\delta}(x-\bar\gamma).
\]
与定理 \ref{thm:xi-poisson-cauchy-f-divergence-sharp-asymptotic} 同样的二阶展开给出
\[
\frac{h_t}{g_t}-1
=\frac{1}{(t+\bar\delta)^2}\bigl(\alpha u_1+\beta u_2\bigr)
+O\bigl((t+\bar\delta)^{-3}\bigr),
\]
\[
\frac{h_t'}{g_t}-1
=\frac{1}{(t+\bar\delta)^2}\bigl(\alpha' u_1+\beta' u_2\bigr)
+O\bigl((t+\bar\delta)^{-3}\bigr),
\]
其中
\[
\alpha=\Var_\mu(\gamma)-\Var_\mu(\delta),\quad
\beta=-\Cov_\mu(\gamma,\delta),
\]
\[
\alpha'=\Var_{\mu'}(\gamma)-\Var_{\mu'}(\delta),\quad
\beta'=-\Cov_{\mu'}(\gamma,\delta).
\]
令
\[
\varepsilon_t:=\frac{h_t}{g_t}-1,\qquad
\varepsilon_t':=\frac{h_t'}{g_t}-1.
\]
则
\[
\frac{h_t}{h_t'}
=\frac{1+\varepsilon_t}{1+\varepsilon_t'}
=1+(\varepsilon_t-\varepsilon_t')+O\bigl((t+\bar\delta)^{-4}\bigr)
\]
于 \(L^\infty\) 成立，且
\[
\varepsilon_t-\varepsilon_t'
=\frac{1}{(t+\bar\delta)^2}
\Bigl((\alpha-\alpha')u_1+(\beta-\beta')u_2\Bigr)
+O\bigl((t+\bar\delta)^{-3}\bigr).
\]
对
\[
D_f(h_t\mid h_t')=\int h_t'(x)\,f\!\left(\frac{h_t(x)}{h_t'(x)}\right)\dd x
\]
在 \(1\) 处作二阶展开，并利用
\[
\int h_t'\!\left(\frac{h_t}{h_t'}-1\right)\dd x=0
\]
消去线性项；再以
\[
h_t'=g_t+O\bigl((t+\bar\delta)^{-3}\bigr)
\]
替换权重，仅产生低阶误差，主项化为
\[
\frac{f''(1)}{2}\int g_t(\varepsilon_t-\varepsilon_t')^2\dd x.
\]
按与定理 \ref{thm:xi-poisson-cauchy-f-divergence-sharp-asymptotic} 相同的奇偶分解与积分恒等式
\[
\int G u_1^2=\frac14,\qquad \int G u_2^2=1
\]
可得
\[
D_f(h_t\mid h_t')
=\frac{f''(1)}{8}\frac{(\alpha-\alpha')^2+4(\beta-\beta')^2}{(t+\bar\delta)^4}
+o(t^{-4}).
\]
又
\[
m_2-m_2'=(\alpha-\alpha')-2i(\beta-\beta'),
\]
故
\[
|m_2-m_2'|^2=(\alpha-\alpha')^2+4(\beta-\beta')^2.
\]
代回即得结论。证毕。
\end{proof}

\endinput


\paragraph{二尺度反混叠与离散曲率刚性的统一误差预算}
设有限缺陷指数谱像写为
\[
u_n(\Delta)=4\pi\sum_{j=1}^{\kappa}\frac{m_j\delta_j}{1+\delta_j}
\exp\!\bigl(-(1+\delta_j)n\Delta\bigr)\exp\!\bigl(-i\gamma_j n\Delta\bigr),
\qquad n\ge 0,
\]
其中 \(\kappa\ge 1\)、\(\delta_j>0\)、\(\gamma_j\in\mathbb R\)、\(m_j>0\)。
记单模态谱像
\[
z_j(\Delta):=\exp\!\bigl(-(1+\delta_j)\Delta\bigr)\exp(-i\gamma_j\Delta)\in\mathbb D.
\]

\begin{theorem}[二尺度共动指数谱像的无混叠可识别性与 Diophantine 相位门槛]\label{thm:xi-two-scale-comoving-spectrum-diophantine-threshold}
\leanverified{paper\_xi\_two\_scale\_comoving\_spectrum\_diophantine\_threshold}
取两步长 \(\Delta_1,\Delta_2>0\)，并设
\[
\alpha:=\frac{\Delta_1}{\Delta_2}\notin\mathbb Q.
\]
则成立：
\begin{enumerate}
  \item 映射
  \[
  (\gamma,\delta)\longmapsto \bigl(z(\Delta_1),z(\Delta_2)\bigr)
  \]
  在 \(\mathbb R\times(0,\infty)\) 上单射，故 \(\bigl(z(\Delta_1),z(\Delta_2)\bigr)\) 唯一决定 \((\gamma,\delta)\)；
  \item 进一步设 \(\alpha\) 满足 \((c,\nu)\)-Diophantine 条件：存在 \(c>0\)、\(\nu\ge2\)，对任意既约有理数 \(p/q\)（\(q\ge1\)）有
  \[
  \left|\alpha-\frac{p}{q}\right|\ge\frac{c}{q^\nu}.
  \]
  不失一般性取 \(\Delta_1\le\Delta_2\)。若 \(|\gamma|\le T\)，并定义观测相位误差
  \[
  \eta:=\max_{r=1,2}\operatorname{dist}\!\bigl(\arg z^{\mathrm{obs}}(\Delta_r)-\arg z(\Delta_r),\,2\pi\mathbb Z\bigr),
  \]
  当
  \[
  \eta<
  \frac{\pi c}{
  \bigl(\lceil T\Delta_2/(2\pi)\rceil+1\bigr)^{\nu-1}
  }
  \]
  时，\(\gamma\) 的整数提升在区间 \([-T,T]\) 内唯一，且任意一致重构 \(\gamma^{\mathrm{rec}}\) 满足
  \[
  |\gamma^{\mathrm{rec}}-\gamma|
  \le
  \frac{\eta}{\Delta_1}+\frac{\eta}{\Delta_2}.
  \]
  同时
  \[
  |\delta^{\mathrm{rec}}-\delta|
  \le
  \min_{r=1,2}
  \frac{1}{\Delta_r}
  \left|
  \log|z^{\mathrm{obs}}(\Delta_r)|-\log|z(\Delta_r)|
  \right|.
  \]
\end{enumerate}
\end{theorem}
\begin{proof}
由
\[
|z(\Delta_r)|=\exp\!\bigl(-(1+\delta)\Delta_r\bigr)
\]
可得
\[
\delta=-\frac{1}{\Delta_r}\log|z(\Delta_r)|-1,
\]
故 \(\delta\) 由任一通道模长唯一确定。
若两组参数 \((\gamma,\delta)\)、\((\gamma',\delta)\) 诱导同一双通道相位，则
\[
e^{-i(\gamma-\gamma')\Delta_r}=1,\qquad r=1,2,
\]
即
\[
(\gamma-\gamma')\Delta_r\in2\pi\mathbb Z.
\]
从而存在整数 \(k_1,k_2\) 使
\[
\gamma-\gamma'=\frac{2\pi k_1}{\Delta_1}=\frac{2\pi k_2}{\Delta_2},
\]
故 \(\alpha=\Delta_1/\Delta_2=k_1/k_2\)。由 \(\alpha\notin\mathbb Q\) 得 \(k_1=k_2=0\)，即 \(\gamma=\gamma'\)。单射性成立。

对误差门槛部分，取相位代表元 \(\theta_r\in\mathbb R\) 满足
\[
e^{-i\theta_r}=\frac{z(\Delta_r)}{|z(\Delta_r)|},
\qquad
\theta_r=-\gamma\Delta_r.
\]
观测相位写为
\[
\arg z^{\mathrm{obs}}(\Delta_r)=\theta_r+e_r\pmod{2\pi},
\qquad |e_r|\le\eta.
\]
任意重构提升整数 \(k_r\in\mathbb Z\) 给出
\[
\gamma^{\mathrm{rec}}\Delta_r=-\theta_r-e_r+2\pi k_r.
\]
于是
\[
\Delta_r(\gamma^{\mathrm{rec}}-\gamma)=2\pi k_r-e_r.
\]
若发生错误提升，则 \((k_1,k_2)\neq(0,0)\)。
将上式分别除以 \(\Delta_1,\Delta_2\) 并相减：
\[
\left|
\frac{k_1}{\Delta_1}-\frac{k_2}{\Delta_2}
\right|
\le
\frac{\eta}{2\pi}\left(\frac1{\Delta_1}+\frac1{\Delta_2}\right).
\]
若 \(k_2=0\)，则由上式与 \(\Delta_1\le\Delta_2\) 得
\[
\frac{|k_1|}{\Delta_1}
\le
\frac{\eta}{2\pi}\left(\frac1{\Delta_1}+\frac1{\Delta_2}\right)
\le
\frac{\eta}{\pi\Delta_1}
<
\frac1{\Delta_1},
\]
故 \(k_1=0\)，与错误提升矛盾，因此必有 \(k_2\neq0\)。
于是
\[
\left|\frac{k_1}{k_2}-\alpha\right|
\le
\frac{\eta}{2\pi|k_2|}(1+\alpha).
\]
由 \(\Delta_1\le\Delta_2\) 得 \(\alpha\le1\)，故右侧不超过 \(\eta/(\pi|k_2|)\)。
另一方面，若候选频率 \(\gamma^{\mathrm{rec}}\in[-T,T]\)，则
\[
|k_2|\le \left\lceil\frac{T\Delta_2}{2\pi}\right\rceil+1=:Q_T.
\]
Diophantine 下界给出
\[
\left|\frac{k_1}{k_2}-\alpha\right|
\ge
\frac{c}{|k_2|^\nu}
\ge
\frac{c}{Q_T^\nu}.
\]
若
\[
\eta<\frac{\pi c}{Q_T^{\nu-1}},
\]
则与上界矛盾，故不存在非零 \((k_1,k_2)\)，提升唯一。

在唯一提升情形 \(k_r=0\) 下，
\[
|\gamma^{\mathrm{rec}}-\gamma|
=\frac{|e_r|}{\Delta_r}\le\frac{\eta}{\Delta_r},\qquad r=1,2,
\]
从而
\[
|\gamma^{\mathrm{rec}}-\gamma|
\le
\frac{\eta}{\Delta_1}+\frac{\eta}{\Delta_2}.
\]
\(\delta\) 的误差界由
\[
\delta=-\frac1{\Delta_r}\log|z(\Delta_r)|-1
\]
逐通道直接传播并取最小值得证。证毕。
\end{proof}

\begin{corollary}[黄金比例步长比值的统一门槛常数]\label{cor:xi-golden-ratio-stepsize-unaliasing-threshold}
若
\[
\alpha=\varphi:=\frac{1+\sqrt5}{2},
\]
则由 Hurwitz--Markov 极值性质，对任意既约有理数 \(p/q\) 有
\[
\left|\varphi-\frac{p}{q}\right|\ge\frac{1}{\sqrt5\,q^2}.
\]
因此定理 \ref{thm:xi-two-scale-comoving-spectrum-diophantine-threshold} 中可取
\[
(c,\nu)=\left(\frac1{\sqrt5},\,2\right),
\]
并得到统一相位容忍阈值
\[
\eta<
\frac{\pi}{
\sqrt5\left(\left\lceil T\Delta_2/(2\pi)\right\rceil+1\right)
}
\asymp T^{-1}.
\]
\end{corollary}
\begin{proof}
将
\(
c=1/\sqrt5,\ \nu=2
\)
代入定理 \ref{thm:xi-two-scale-comoving-spectrum-diophantine-threshold} 的门槛式即可。证毕。
\end{proof}

\begin{theorem}[离散曲率控制下的近单模态刚性]\label{thm:xi-discrete-curvature-near-unimodal-rigidity}
设有限谱
\[
\{a_j\}_{j=1}^{\kappa}\subset[1,\infty),\qquad w_j>0,
\]
并定义
\[
Z(u):=\sum_{j=1}^{\kappa}w_j e^{-ua_j},\qquad u\ge0.
\]
令倾斜分布
\[
\mathbb P_u(A=a_j):=\frac{w_j e^{-ua_j}}{Z(u)}.
\]
则
\[
\partial_u\log Z(u)=-\mathbb E_u[A],\qquad
\partial_u^2\log Z(u)=\Var_u(A)\ge0.
\]
给定 \(\Delta>0\)，定义离散曲率
\[
\kappa_Z(u;\Delta):=\log Z(u+\Delta)-2\log Z(u)+\log Z(u-\Delta),
\qquad u\ge\Delta.
\]
则成立：
\begin{enumerate}
  \item 精确积分表示
  \[
  \kappa_Z(u;\Delta)=
  \int_{-\Delta}^{\Delta}(\Delta-|s|)\Var_{u+s}(A)\,\dd s;
  \]
  \item 若 \(\kappa_Z(u;\Delta)\le\varepsilon\)，则存在 \(u^\ast\in[u-\Delta,u+\Delta]\) 使
  \[
  \Var_{u^\ast}(A)\le\frac{2\varepsilon}{\Delta^2}.
  \]
  因而对任意 \(r>0\) 有
  \[
  \mathbb P_{u^\ast}\!\bigl(|A-\mathbb E_{u^\ast}A|\ge r\bigr)
  \le
  \frac{2\varepsilon}{\Delta^2r^2}.
  \]
\end{enumerate}
\end{theorem}
\begin{proof}
对任意 \(f\in C^2\) 有恒等式
\[
f(u+\Delta)-2f(u)+f(u-\Delta)
=
\int_{-\Delta}^{\Delta}(\Delta-|s|)f''(u+s)\,\dd s.
\]
取 \(f=\log Z\) 即得第一条。

若对一切 \(s\in[-\Delta,\Delta]\) 都有
\(
\Var_{u+s}(A)>2\varepsilon/\Delta^2
\)，
则由第一条与
\[
\int_{-\Delta}^{\Delta}(\Delta-|s|)\,\dd s=\Delta^2
\]
推出 \(\kappa_Z(u;\Delta)>2\varepsilon\)，与假设矛盾，故存在 \(u^\ast\) 满足所述方差上界。
尾界由 Chebyshev 不等式直接给出。证毕。
\end{proof}

\begin{theorem}[二尺度反混叠与离散曲率刚性的全链路停止证书]\label{thm:xi-two-scale-unaliasing-discrete-curvature-stop-certificate}
在定理 \ref{thm:xi-two-scale-comoving-spectrum-diophantine-threshold} 的模型下，设两通道步长比值 \(\alpha=\Delta_1/\Delta_2\) 满足 \((c,\nu)\)-Diophantine 条件，并设同一缺陷谱诱导配分函数 \(Z\)。
若在某尺度参数 \(u\) 上有
\[
\kappa_Z(u;\Delta_r)\le\varepsilon_r,\qquad r=1,2,
\]
则存在 \(u_r^\ast\in[u-\Delta_r,u+\Delta_r]\) 使
\[
\Var_{u_r^\ast}(A)\le\frac{2\varepsilon_r}{\Delta_r^2}.
\]
进一步假设重构流程给出残差 \(\mathcal R_r\) 与相位误差 \(\eta_r\)，并满足可审计放大模型
\[
\eta_r\le \Lambda_r\!\bigl(\Var_{u_r^\ast}(A)\bigr)\,\mathcal R_r,
\]
其中 \(\Lambda_r\) 对其自变量单调不减。
则有
\[
\eta_r\le
\Lambda_r\!\left(\frac{2\varepsilon_r}{\Delta_r^2}\right)\mathcal R_r.
\]
记 \(\eta:=\max\{\eta_1,\eta_2\}\)。若
\[
\eta<
\frac{\pi c}{
\bigl(\lceil T\Delta_2/(2\pi)\rceil+1\bigr)^{\nu-1}
},
\]
则所有 \(|\gamma_j|\le T\) 模态的整数提升唯一，且
\[
|\gamma_j^{\mathrm{rec}}-\gamma_j|
\le
\frac{\eta_1}{\Delta_1}+\frac{\eta_2}{\Delta_2},
\]
\[
|\delta_j^{\mathrm{rec}}-\delta_j|
\le
\min_{r=1,2}
\frac{1}{\Delta_r}
\left|
\log|z_j^{\mathrm{obs}}(\Delta_r)|-\log|z_j(\Delta_r)|
\right|.
\]
因此，一旦曲率阈值与 Diophantine 相位阈值同时满足，即得到可审计停止条件：后续递归仅改变常数项，不再改变无混叠可识别性与误差阶。
\end{theorem}
\begin{proof}
前半部分由定理 \ref{thm:xi-discrete-curvature-near-unimodal-rigidity} 逐通道应用得到。
单调性假设给出
\[
\eta_r\le
\Lambda_r\!\left(\frac{2\varepsilon_r}{\Delta_r^2}\right)\mathcal R_r.
\]
再将 \(\eta=\max\{\eta_1,\eta_2\}\) 代入定理 \ref{thm:xi-two-scale-comoving-spectrum-diophantine-threshold} 即得提升唯一性与 \((\gamma_j,\delta_j)\) 的显式误差界。
由于该界的主阶由阈值结构决定，阈值满足后继续递归不改变识别性与误差阶，仅可能改善前常数，故停止判据成立。证毕。
\end{proof}

\endinput


\paragraph{判别式层级分解、三重分歧并合与 Galois 对称闭合}
沿用前述四次特征多项式 \(F(\lambda;u,p)\) 与其 \(\lambda\)-判别式分解记号
\[
\operatorname{Disc}_{\lambda}F(\lambda;u,p)
=-p^3u(1-p)\,D(u,p),
\qquad
\deg_u D=11.
\]

\begin{theorem}[判别式之判别式的完全分解与唯一临界偏置]\label{thm:xi-discriminant-of-discriminant-complete-factorization}
多项式 \(D(u,p)\) 关于变量 \(u\) 的判别式在 \(\mathbb Z[p]\) 上满足
\[
\operatorname{Disc}_{u}D(u,p)
=-20639121408\,
p^{48}(p-1)^{34}(3p+1)^2
E_8(p)^2
F_8(p)
G_{16}(p)^3,
\]
其中
\[
\begin{aligned}
E_8(p)=\;&243p^{8}-324p^{7}-216p^{6}+4124p^{5}-8938p^{4}+6548p^{3}-1712p^{2}+20p-1,\\
F_8(p)=\;&84375p^{8}-421875p^{7}+843750p^{6}-847374p^{5}+428395p^{4}\\
&-83215p^{3}-6312p^{2}+1296p+1024,\\
G_{16}(p)=\;&8p^{16}-68p^{15}+246p^{14}-14979p^{13}+135362p^{12}-459861p^{11}+19453803p^{10}\\
&-74533695p^{9}+96659367p^{8}-26757131p^{7}-34511365p^{6}+16694550p^{5}\\
&+6485075p^{4}-2181640p^{3}-916704p^{2}-63336p+6272.
\end{aligned}
\]
并且在 \(p\in(0,1)\) 上成立：
\begin{enumerate}
  \item \(E_8\) 与 \(F_8\) 在 \((0,1)\) 内无实根；
  \item \(G_{16}\) 在 \((0,1)\) 内恰有一个实根 \(p_\star\)，数值为
  \[
  p_\star\approx0.053415794573903219\ldots.
  \]
\end{enumerate}
因此在 \(p\in(0,1)\) 内，\(D(\cdot,p)\) 出现 \(u\)-重根的唯一参数为 \(p=p_\star\)。
\end{theorem}
\begin{proof}
对 \(D(\cdot,p)\) 直接计算 \(\operatorname{Disc}_{u}\) 并在 \(\mathbb Z[p]\) 上完全因子分解，即得所示恒等式。
该等式属于代数恒等式，可通过展开系数逐项校核。
对 \(E_8,F_8,G_{16}\) 在区间 \((0,1)\) 上分别构造 Sturm 序列并计数端点符号变号差，可得实根数依次为 \(0,0,1\)。
故 \(\operatorname{Disc}_{u}D\) 在 \((0,1)\) 内仅在 \(p_\star\) 处为零。证毕。
\end{proof}

\begin{theorem}[唯一三阶分歧点与三重谱碰撞]\label{thm:xi-unique-cubic-branchpoint-triple-collision}
令 \(p_\star\) 为定理 \ref{thm:xi-discriminant-of-discriminant-complete-factorization} 中
\(G_{16}\) 在 \((0,1)\) 的唯一实根。
则存在唯一 \(u_\star<0\) 使
\[
D(u_\star,p_\star)=0,
\qquad
\partial_u D(u_\star,p_\star)=0,
\]
且 \(u_\star\) 为 \(D(\cdot,p_\star)\) 的二重根，其余 \(u\)-根均为单根。数值上
\[
u_\star\approx-4.67345092038166580786413589157\ldots.
\]

进一步，在谱曲线
\[
\mathcal C_{p_\star}:\quad F(\lambda;u,p_\star)=0
\]
上，投影
\(
\pi:\mathcal C_{p_\star}\to\mathbb P_u^1
\)
在 \(u=u_\star\) 处出现唯一三阶分歧点：存在唯一实数 \(\lambda_\star\) 满足
\[
F(\lambda_\star;u_\star,p_\star)
=\partial_\lambda F(\lambda_\star;u_\star,p_\star)
=\partial_\lambda^2F(\lambda_\star;u_\star,p_\star)=0,
\]
\[
\partial_\lambda^3F(\lambda_\star;u_\star,p_\star)\neq0,
\qquad
\lambda_\star\approx0.387716992134592608093404839266\ldots.
\]
等价地，\(F(\cdot;u_\star,p_\star)\) 在 \(\lambda=\lambda_\star\) 处为三重根，而非 \((2,2)\) 型双二重点。

负实轴上的分裂型为：
\[
\#\{u<0:D(u,p)=0\}
=
\begin{cases}
2, & 0<p<p_\star,\\
1\ (\text{按二重根计}), & p=p_\star,\\
0, & p_\star<p<1.
\end{cases}
\]
\end{theorem}
\begin{proof}
由定理 \ref{thm:xi-discriminant-of-discriminant-complete-factorization}，\(p\in(0,1)\) 内发生 \(u\)-重根的参数仅有 \(p_\star\)，从而得到唯一并合参数。
对 \(F\) 与 \(\partial_\lambda F\) 作 subresultant 链，二次 subresultant 可写为
\[
S_2(\lambda;u,p)=A(u,p)\lambda^2+B(u,p)\lambda+C(u,p),
\]
其中
\[
\begin{aligned}
A(u,p)&=5p^{2}-8pu-2p-3,\\
B(u,p)&=12p^{3}u^{3}-2p^{3}-12p^{2}u^{3}-10p^{2}u+4p^{2}+10pu-2p,\\
C(u,p)&=-p^{4}u^{3}+2p^{3}u^{3}-15p^{3}u-p^{2}u^{3}+14p^{2}u+pu.
\end{aligned}
\]
并有
\[
\operatorname{Disc}_{\lambda}S_2(\lambda;u,p)=4p(p-1)\,H(u,p),
\]
其中 \(H(u,p)\) 为显式六次多项式。
在 \((u_\star,p_\star)\) 上，\(D=\partial_uD=0\) 与 \(H=0\) 联合给出
\(
\gcd(F,\partial_\lambda F)=(\lambda-\lambda_\star)^2
\)；
故 \(F\) 在 \(\lambda_\star\) 处的重数为 \(3\)，并由三阶导非零排除更高重数。

负实根分裂由参数连续性与“仅在 \(p_\star\) 处发生并合”得到：在
\((0,p_\star)\) 与 \((p_\star,1)\) 两连通分支内，负实根数为局部常数。
取代表点作 Sturm 计数即可确定分别为 \(2\) 与 \(0\)。证毕。
\end{proof}

\begin{theorem}[分支多项式 \(D(u,p)\) 的泛型 Galois 群为 \(S_{11}\)]\label{thm:xi-branch-polynomial-generic-galois-s11}
将 \(D(u,p)\) 视为 \(u\) 的十一次多项式，系数域为 \(\mathbb Q(p)\)。
则
\[
\operatorname{Gal}\bigl(D(u,p)/\mathbb Q(p)\bigr)\cong S_{11}.
\]
\end{theorem}
\begin{proof}
取专化 \(p=\frac23\)，清分母后得到
\[
\begin{aligned}
d(u)=\;&216u^{11}-1080u^{9}-352u^{8}+6678u^{7}+2544u^{6}-11358u^{5}\\
&-6534u^{4}+15168u^{3}-8019u^{2}+1620u-108.
\end{aligned}
\]
其模 \(11\) 约化在 \(\mathbb F_{11}[u]\) 上不可约，故 Frobenius 给出一个 \(11\)-循环；
模 \(5\) 约化的次数分解型为 \((7,2,1,1)\)，故群内含 \(7\)-循环与奇置换。
由 Jordan 定理，素数次数本原群且含长度 \(7\) 的循环时包含 \(A_{11}\)；
又因存在奇置换，故该专化群为 \(S_{11}\)。

专化群是泛型群的子群，既然专化已达最大可能 \(S_{11}\)，泛型群只能是 \(S_{11}\)。
证毕。
\end{proof}

\begin{theorem}[四次谱覆盖 \(F(\lambda;u,p)\) 的泛型 Galois 群为 \(S_4\)]\label{thm:xi-quartic-spectrum-generic-galois-s4}
将 \(F(\lambda;u,p)\) 视为 \(\lambda\) 的四次多项式，系数域为 \(\mathbb Q(u,p)\)。
则
\[
\operatorname{Gal}\bigl(F(\lambda;u,p)/\mathbb Q(u,p)\bigr)\cong S_4.
\]
\end{theorem}
\begin{proof}
取专化 \(p=\frac23,\ u=2\)，得到
\[
F\!\left(\lambda;2,\frac23\right)
=\lambda^4-\frac13\lambda^3-\frac{14}{9}\lambda^2-\frac{20}{27}\lambda+\frac{8}{27}.
\]
清分母后
\[
\widetilde F(\lambda)=27\lambda^4-9\lambda^3-42\lambda^2-20\lambda+8.
\]
对 \(\widetilde F\) 作分裂域计算（例如借助 resolvent 与判别式平方性判据）可得群阶为 \(24\)，即专化群为 \(S_4\)。
专化群为泛型群子群，而 \(S_4\) 已是四次情形最大可能，故泛型群即 \(S_4\)。证毕。
\end{proof}

\begin{theorem}[实分支值的全局拓扑与两相结构]\label{thm:xi-real-branch-topology-two-phase-structure}
定义分支值集合
\[
\mathcal B_p:=\{0\}\cup\{u\in\mathbb C:\ D(u,p)=0\},
\]
其中 \(u=0\) 来自 \(\operatorname{Disc}_{\lambda}F\) 的显式因子 \(u\)。
对任意 \(p\in(0,1)\)，成立：
\begin{enumerate}
  \item
  \[
  D(0,p)=4p^2(p-1)^4(3p+1)>0,
  \]
  因而分支值在实轴上不穿越 \(u=0\)；
  \item 当 \(p\neq\frac12\) 时，\(D(u,p)=0\) 在正实轴上恰有三条实根：两条位于 \((0,1)\)，一条位于 \((1,\infty)\)；
  当 \(p=\frac12\) 时额外出现单根 \(u=1\)；
  \item 令 \(p_\star\) 如定理 \ref{thm:xi-discriminant-of-discriminant-complete-factorization}，则负实轴上存在唯一相变：
  \[
  \begin{cases}
  0<p<p_\star:\ \#\{u<0:D(u,p)=0\}=2,\\
  p=p_\star:\ \text{两根并合为 }u_\star<0\text{ 的二重根},\\
  p_\star<p<1:\ \#\{u<0:D(u,p)=0\}=0.
  \end{cases}
  \]
\end{enumerate}
\end{theorem}
\begin{proof}
第一条由 \(D(0,p)\) 的闭式直接得到。

若开区间 \(I\subset(0,1)\) 上恒有 \(\operatorname{Disc}_uD\neq0\)，则 \(D(\cdot,p)\) 的实根数及其在各分区
\((-\infty,0),(0,1),(1,\infty)\) 的分布在 \(I\) 内保持不变。
由定理 \ref{thm:xi-discriminant-of-discriminant-complete-factorization}，\((0,1)\) 内仅在 \(p_\star\) 处可发生重根事件，故根分布仅可能在该点改变。
在两侧任取代表点进行 Sturm 计数，即得负实轴根数分别为 \(2\) 与 \(0\)。

对 \(p=\frac12\)，由 \(D(1,p)\) 的直接代入可验证 \(u=1\) 为单根，且对应于边界点事件，不改变正轴总根数结构。证毕。
\end{proof}

\begin{corollary}[分支值延拓的满对称单群作用]\label{cor:xi-full-symmetric-geometric-monodromy-s11}
分支多项式 \(D(u,p)\) 的几何 Galois 群亦为 \(S_{11}\)。
等价地，参数 \(p\) 在去除 \(\operatorname{Disc}_uD=0\) 的参数空间内沿闭路解析延拓时，
十一个分支值 \(\{u:D(u,p)=0\}\) 的单值化排列给出满对称群作用。
\end{corollary}
\begin{proof}
由定理 \ref{thm:xi-branch-polynomial-generic-galois-s11}，算术 Galois 群为 \(S_{11}\)。
几何群只能为 \(S_{11}\) 或 \(A_{11}\)。
若几何群为 \(A_{11}\)，则 \(\operatorname{Disc}_uD\) 在 \(\mathbb Q(p)\) 中应为平方。
然而定理 \ref{thm:xi-discriminant-of-discriminant-complete-factorization} 的分解中
\(
F_8(p)^1
\)
与
\(
G_{16}(p)^3
\)
均为奇指数因子，故判别式非平方，排除 \(A_{11}\)。
因此几何群为 \(S_{11}\)。证毕。
\end{proof}

\endinput


\paragraph{软核张量算子的全谱分解、割线方程与惯性平衡}
设 \(q\ge2\)，\(U_q:=\{0,1\}^q\)，并按标准张量识别
\[
\mathbb R^{U_q}\cong(\mathbb R^2)^{\otimes q}.
\]
记
\[
K:=\begin{pmatrix}1&1\\[2pt]1&0\end{pmatrix},
\qquad
J:=\begin{pmatrix}1&1\\[2pt]1&1\end{pmatrix},
\]
\[
D^{(q)}:=K^{\otimes q},
\qquad
J_q:=J^{\otimes q},
\qquad
T^{(q)}:=\frac12\bigl(J_q+D^{(q)}\bigr).
\]
令
\[
\varphi:=\frac{1+\sqrt5}{2},
\qquad
\psi:=\frac{1-\sqrt5}{2}=-\varphi^{-1},
\]
\[
\mu_i^{(q)}:=\varphi^{\,q-i}\psi^{\,i}=(-1)^i\varphi^{\,q-2i},
\qquad 0\le i\le q.
\]
并记 \(\mathcal H_q:=\mathrm{Sym}^q(\mathbb R^2)\subset\mathbb R^{U_q}\) 为 \(S_q\)-对称子空间。

\begin{theorem}[全谱的例外--体谱直和分解与精确重数]\label{thm:xi-replica-softcore-full-spectrum-direct-sum}
对每个 \(0\le i\le q\)，\(\mu_i^{(q)}/2\) 是 \(T^{(q)}\) 的特征值，其代数重数恰为
\[
\binom{q}{i}-1.
\]
其余恰好 \(q+1\) 个特征值为 \(\left.T^{(q)}\right|_{\mathcal H_q}\) 的谱
\(\{\lambda_0^{(q)},\dots,\lambda_q^{(q)}\}\)，并与体谱不相交。

等价地，存在正交直和分解
\[
\mathbb R^{U_q}
=
\mathcal H_q\ \oplus\
\bigoplus_{i=0}^{q}\bigl(E_i\cap\mathbf 1^\perp\bigr),
\]
其中 \(E_i:=\ker(D^{(q)}-\mu_i^{(q)}I)\)、\(\mathbf1\) 为全 \(1\) 向量，且
\[
\left.T^{(q)}\right|_{E_i\cap\mathbf1^\perp}
=
\frac{\mu_i^{(q)}}{2}I.
\]
\end{theorem}
\begin{proof}
矩阵 \(K\) 为实对称矩阵，特征值为 \(\varphi,\psi\)，故 \(D^{(q)}=K^{\otimes q}\) 的特征值为
\(\mu_i^{(q)}\)，且
\[
\dim E_i=\binom{q}{i}.
\]
又 \(J_q\) 为秩一算子且
\[
\operatorname{Ran}(J_q)=\operatorname{span}\{\mathbf1\},
\qquad
J_qx=0\iff x\perp\mathbf1.
\]
因此对任意 \(x\in E_i\cap\mathbf1^\perp\)，
\[
T^{(q)}x=\frac12\bigl(D^{(q)}x+J_qx\bigr)=\frac{\mu_i^{(q)}}2x.
\]

记 \(x_i:=\operatorname{Proj}_{E_i}\mathbf1\)。
由于 \((1,1)\) 在 \(K\) 的两条特征方向上投影均非零，故 \(x_i\neq0\) 对全部 \(i\) 成立，从而
\[
\dim(E_i\cap\mathbf1^\perp)=\dim E_i-1=\binom{q}{i}-1.
\]
体谱总维数为
\[
\sum_{i=0}^{q}\left(\binom{q}{i}-1\right)=2^q-(q+1).
\]
余下 \(q+1\) 维空间由 \(\{x_i\}_{i=0}^{q}\) 张成；该空间 \(S_q\)-不变且维数与 \(\mathcal H_q\) 相同，故即为 \(\mathcal H_q\)。
由结论 \ref{con:xi-replica-softcore-exceptional-secular-strict-interlacing}，\(\mathcal H_q\) 上的例外谱与 \(\{\mu_i^{(q)}/2\}\) 严格交错，因而与体谱不相交。证毕。
\end{proof}

\begin{corollary}[特征多项式分解与体谱因子闭式]\label{cor:xi-replica-softcore-charpoly-bulk-factorization}
记
\[
\chi_q(x):=\det(xI-T^{(q)}),
\qquad
\chi_q^{\mathrm{exc}}(x):=\det\!\left(xI-\left.T^{(q)}\right|_{\mathcal H_q}\right).
\]
则
\[
\chi_q(x)
=
\chi_q^{\mathrm{exc}}(x)\,
\prod_{i=0}^{q}\left(x-\frac{\mu_i^{(q)}}2\right)^{\binom{q}{i}-1}.
\]
若记 \(m_q^{\mathrm{exc}}\) 为 \(\left.T^{(q)}\right|_{\mathcal H_q}\) 的最小多项式，则
\[
m_q(x)=m_q^{\mathrm{exc}}(x)\prod_{i=0}^{q}\left(x-\frac{\mu_i^{(q)}}2\right).
\]
\end{corollary}
\begin{proof}
由定理 \ref{thm:xi-replica-softcore-full-spectrum-direct-sum} 的不变直和分解，特征多项式按分块相乘。
又体谱点在对应分块上为纯标量，指数即其重数。
最小多项式由各不变分量最小多项式的最小公倍式给出，并由谱分离得到所示因子。证毕。
\end{proof}

\begin{theorem}[例外谱的加权割线方程]\label{thm:xi-replica-softcore-exceptional-weighted-secular}
定义
\[
\alpha:=1+\frac{2}{\sqrt5},
\qquad
\beta:=1-\frac{2}{\sqrt5},
\]
\[
w_i^{(q)}:=\binom{q}{i}\alpha^{\,q-i}\beta^{\,i}>0,
\qquad 0\le i\le q.
\]
则 \(\left.T^{(q)}\right|_{\mathcal H_q}\) 的全部特征值恰为方程
\[
\sum_{i=0}^{q}\frac{w_i^{(q)}}{\lambda-\mu_i^{(q)}/2}=2
\]
的全部根。
并且
\[
\sum_{i=0}^{q}w_i^{(q)}=2^q,
\]
所有根均为单根。
\end{theorem}
\begin{proof}
取 \(K\) 的单位正交本征向量 \(e_+,e_-\)（对应 \(\varphi,\psi\)），并记
\[
c_+:=\langle(1,1),e_+\rangle,\qquad c_-:=\langle(1,1),e_-\rangle.
\]
直接计算得
\[
c_+^2=1+\frac{2}{\sqrt5}=\alpha,\qquad
c_-^2=1-\frac{2}{\sqrt5}=\beta.
\]
在 \(\mathcal H_q\) 的标准对称本征基 \(\{s_i\}_{i=0}^{q}\) 下，
\[
\left.D^{(q)}\right|_{\mathcal H_q}=\operatorname{diag}\bigl(\mu_0^{(q)},\dots,\mu_q^{(q)}\bigr),
\]
且 \(\mathbf1\) 的坐标向量 \(g\) 满足
\[
g_i^2=\binom{q}{i}c_+^{2(q-i)}c_-^{2i}=w_i^{(q)}.
\]
故
\[
\left.T^{(q)}\right|_{\mathcal H_q}
=
\operatorname{diag}\!\left(\frac{\mu_i^{(q)}}2\right)_{i=0}^{q}
+\frac12\,gg^{\mathsf T}.
\]
由矩阵行列式引理，
\[
0=
\det\!\left(
\lambda I-\operatorname{diag}\!\left(\frac{\mu_i^{(q)}}2\right)-\frac12gg^{\mathsf T}
\right)
\]
\[
=
\prod_{i=0}^{q}\left(\lambda-\frac{\mu_i^{(q)}}2\right)
\left(
1-\frac12\sum_{i=0}^{q}\frac{g_i^2}{\lambda-\mu_i^{(q)}/2}
\right).
\]
又 \(g_i\neq0\)，故极点 \(\lambda=\mu_i^{(q)}/2\) 不属例外谱，从而得到割线方程。
\(\sum_i w_i^{(q)}=(\alpha+\beta)^q=2^q\) 由二项式定理即得。
根的单性由
\(
\frac{\mathrm d}{\mathrm d\lambda}\sum_i\frac{w_i^{(q)}}{\lambda-\mu_i^{(q)}/2}<0
\)
推出。证毕。
\end{proof}

\begin{proposition}[秩一割线方程的耦合单调性与导数--范数同一律]\label{prop:xi-rankone-secular-monotonicity-derivative-norm}
\leanverified{paper\_xi\_rankone\_secular\_monotonicity\_derivative\_norm}
取互异实数 \(t_0,\dots,t_q\) 与权重 \(w_i>0\)，并令
\[
D:=\mathrm{diag}(t_0,\dots,t_q),\qquad g:=(\sqrt{w_0},\dots,\sqrt{w_q})^{\mathsf T}.
\]
对任意 \(b>0\)，考虑秩一扰动
\[
A(b):=D+b\,gg^{\mathsf T}.
\]
则 \(A(b)\) 的谱由割线方程
\begin{equation}\label{eq:xi-rankone-secular-general}
\Phi(\lambda):=\sum_{i=0}^{q}\frac{w_i}{\lambda-t_i}=\frac1b,
\qquad
\lambda\notin\{t_0,\dots,t_q\},
\end{equation}
给出，且任意根分支 \(\lambda=\lambda(b)\) 在其定义域内满足严格单调与导数闭式
\begin{equation}\label{eq:xi-rankone-secular-derivative}
\frac{\mathrm d\lambda}{\mathrm db}
=\frac{1}{b^2\displaystyle\sum_{i=0}^{q}\frac{w_i}{(\lambda-t_i)^2}}
\;>\;0.
\end{equation}
并且与 \(\lambda(b)\) 对应的特征向量可取为
\[
v(b):=(\lambda I-D)^{-1}g,\qquad
v_i(b)=\frac{\sqrt{w_i}}{\lambda(b)-t_i},
\]
从而耦合量与范数满足恒等式
\[
g^{\mathsf T}v(b)=\frac1b,
\qquad
\|v(b)\|_2^2=\sum_{i=0}^{q}\frac{w_i}{(\lambda(b)-t_i)^2}=\frac{1}{b^2\,\lambda'(b)}.
\]
\end{proposition}
\begin{proof}
由矩阵行列式引理，
\[
\det(\lambda I-A(b))
=\det(\lambda I-D)\Bigl(1-b\,g^{\mathsf T}(\lambda I-D)^{-1}g\Bigr),
\]
其中
\(
g^{\mathsf T}(\lambda I-D)^{-1}g=\sum_i\frac{w_i}{\lambda-t_i}=\Phi(\lambda)
\)。
故 \(\lambda\) 为特征值当且仅当 \eqref{eq:xi-rankone-secular-general} 成立；并且 \(\Phi'(\lambda)=-\sum_i\frac{w_i}{(\lambda-t_i)^2}<0\)，根均为单根。
对 \(\Phi(\lambda(b))=1/b\) 作隐式微分即得 \eqref{eq:xi-rankone-secular-derivative}，从而 \(\lambda'(b)>0\)。
\par
取 \(v(b)=(\lambda I-D)^{-1}g\)，则
\[
(A(b)-\lambda I)v
=(D-\lambda I)v+b\,g(g^{\mathsf T}v)
=-g+b\,g\,\Phi(\lambda)
=0,
\]
其中最后一步使用割线方程；故 \(v(b)\) 为特征向量。
其坐标式与 \(g^{\mathsf T}v=1/b\) 由定义直接给出；
范数恒等式由坐标式与 \eqref{eq:xi-rankone-secular-derivative} 联立得到。
\end{proof}

\begin{theorem}[秩一耦合的谱纤维分裂：小耦合与大耦合极限]\label{thm:xi-rankone-secular-small-large-coupling}
沿用命题 \ref{prop:xi-rankone-secular-monotonicity-derivative-norm} 的记号。
令极点集合 \(t_i\) 两两不同，并记
\[
G(\lambda):=\prod_{i=0}^{q}(\lambda-t_i),
\qquad
N(\lambda):=\sum_{i=0}^{q} w_i\!\!\prod_{\substack{0\le j\le q\\ j\ne i}}\!\!(\lambda-t_j),
\qquad
\Phi(\lambda)=\frac{N(\lambda)}{G(\lambda)}.
\]
则对 \(b>0\) 的全部特征值（即 \(\Phi(\lambda)=1/b\) 的全部根）有：
\begin{enumerate}
  \item 当 \(b\to 0^+\) 时，对每个 \(i\in\{0,\dots,q\}\) 存在唯一根分支 \(\lambda_i(b)\) 满足
  \[
  \lambda_i(b)=t_i+bw_i+O(b^2),
  \qquad
  \lambda_i(b)\to t_i.
  \]
  \item 当 \(b\to+\infty\) 时，存在唯一根分支 \(\lambda_\infty(b)\to+\infty\) 且
  \[
  \lambda_\infty(b)=b\sum_{i=0}^{q}w_i+\frac{\sum_{i=0}^{q}w_it_i}{\sum_{i=0}^{q}w_i}+O(b^{-1}).
  \]
  其余 \(q\) 条根分支在 \(b\to+\infty\) 时分别收敛到多项式 \(N(\lambda)\) 的 \(q\) 个实零点；这些零点均为单根并严格夹在相邻极点之间。
\end{enumerate}
\end{theorem}
\begin{proof}
\par\noindent\textbf{(i) 小耦合.}
固定 \(i\)。在 \(\lambda=t_i\) 的穿孔邻域内，
\[
\Phi(\lambda)=\frac{w_i}{\lambda-t_i}+H_i(\lambda),
\qquad
H_i(\lambda):=\sum_{j\ne i}\frac{w_j}{\lambda-t_j},
\]
其中 \(H_i\) 在 \(\lambda=t_i\) 解析。
令 \(\lambda=t_i+\delta\)，方程 \(\Phi(\lambda)=1/b\) 等价于
\(
w_i/\delta+H_i(t_i)+O(\delta)=1/b
\)。
当 \(b\to 0^+\) 时右端趋于 \(+\infty\)，迫使 \(\delta\to 0\)，并由上式解得
\(\delta=bw_i+O(b^2)\)。
由于 \(\Phi'(\lambda)<0\)，在足够小的邻域内根唯一，从而得到唯一根分支。
\par
\noindent\textbf{(ii) 大耦合.}
当 \(|\lambda|\to\infty\)，有渐近展开
\[
\Phi(\lambda)=\sum_{i=0}^{q}\frac{w_i}{\lambda-t_i}
=\frac{W}{\lambda}+\frac{M_1}{\lambda^2}+O(\lambda^{-3}),
\qquad
W:=\sum_{i=0}^{q}w_i,\quad M_1:=\sum_{i=0}^{q}w_it_i.
\]
令 \(\Phi(\lambda)=1/b\) 并令 \(b\to+\infty\) 得到唯一逃逸根满足 \(\lambda\sim bW\)；
进一步迭代代回即得所示渐近展开。
\par
其余根的收敛：注意 \(N(\lambda)=0\) 当且仅当 \(\Phi(\lambda)=0\)。
在每个相邻极点之间的开区间上，\(\Phi\) 连续且严格单调（\(\Phi'<0\)），并在端点趋于 \(\pm\infty\)，因此 \(N\) 在每个此类区间内恰有一个实零点且为单根。
同理，\(\Phi(\lambda)=1/b\) 在每个区间内亦恰有一个实根；令 \(b\to+\infty\) 使右端趋于 \(0\)，由单调性与连续性得该根收敛到区间内的零点。
\end{proof}

\endinput


\begin{theorem}[二项--几何谱测度的 Fibonacci 幂矩塌缩]\label{thm:xi-replica-softcore-fibonacci-moment-collapse}
沿用定理 \ref{thm:xi-replica-softcore-exceptional-weighted-secular} 的权重 \(w_i^{(q)}\) 与节点
\(
d_i:=\mu_i^{(q)}/2=\frac12\varphi^{\,q-i}\psi^{\,i}=\frac12(-1)^i\varphi^{\,q-2i}
\)（\(0\le i\le q\)）。
则对任意整数 \(m\ge 0\)，成立严格恒等式
\begin{equation}\label{eq:xi-replica-softcore-fibonacci-moment-collapse}
\sum_{i=0}^{q} w_i^{(q)}\,d_i^{\,m}
=2^{-m}\,F_{m+3}^{\,q}.
\end{equation}
因此所有幂矩均在 \(\QQ\) 内闭合；右端不含任何 \(\sqrt5\) 的显式项。
\end{theorem}
\begin{proof}
由 \(\psi=-\varphi^{-1}\)，有
\(
d_i^{\,m}=2^{-m}\varphi^{m(q-i)}\psi^{mi}
\)。
因此
\[
\sum_{i=0}^{q} w_i^{(q)}\,d_i^{\,m}
=2^{-m}\sum_{i=0}^{q}\binom{q}{i}\alpha^{q-i}\beta^{i}\varphi^{m(q-i)}\psi^{mi}
=2^{-m}\sum_{i=0}^{q}\binom{q}{i}(\alpha\varphi^m)^{q-i}(\beta\psi^m)^{i}.
\]
二项式定理给出
\[
\sum_{i=0}^{q} w_i^{(q)}\,d_i^{\,m}
=2^{-m}\bigl(\alpha\varphi^m+\beta\psi^m\bigr)^{q}.
\]
再由
\(
\alpha=\varphi^3/\sqrt5
\)、\(
\beta=-\psi^3/\sqrt5
\)，得到
\[
\alpha\varphi^m+\beta\psi^m=\frac{\varphi^{m+3}-\psi^{m+3}}{\sqrt5}=F_{m+3},
\]
代回即得 \eqref{eq:xi-replica-softcore-fibonacci-moment-collapse}。证毕。
\end{proof}
\leanverified{paper\_xi\_replica\_softcore\_fibonacci\_moment\_collapse}

\paragraph{对称幂谱的 Vandermonde 判别式与秩一耦合结果式}
在本段，Fibonacci 幂矩塌缩 \eqref{eq:xi-replica-softcore-fibonacci-moment-collapse} 与秩一割线方程链（命题 \ref{prop:xi-rankone-secular-monotonicity-derivative-norm}）联结为一条纯代数的判别式/结果式闭式，并进一步给出 Hankel 行列式的同型桥接。

\begin{theorem}[Fibonacci--Vandermonde 判别式分解]\label{thm:xi-symq-k-discriminant-fibonacci-vandermonde}
\leanverified{paper\_xi\_symq\_k\_discriminant\_fibonacci\_vandermonde}
令
\[
K=\begin{pmatrix}1&1\\[2pt]1&0\end{pmatrix},
\qquad
B_q:=\mathrm{Sym}^q(K)\in M_{q+1}(\ZZ),
\qquad
g_q(z):=\det(zI_{q+1}-B_q).
\]
则对一切 \(q\ge 2\)，成立严格恒等式
\[
\Disc(g_q)
=5^{\frac{q(q+1)}2}\prod_{k=1}^{q}F_k^{\,2(q+1-k)}.
\]
\end{theorem}
\begin{proof}
由于 \(K\) 为实对称矩阵，其特征值为 \(\varphi\) 与 \(\psi=-\varphi^{-1}\)；于是
\(\mathrm{Sym}^q(K)\) 的谱为
\[
\mu_i=\varphi^{\,q-i}\psi^{\,i}=\varphi^{\,q}\,r^i,
\qquad
r:=\psi/\varphi=-\varphi^{-2},
\qquad
0\le i\le q.
\]
从而 \(g_q(z)=\prod_{i=0}^{q}(z-\mu_i)\) 为首一多项式，并有
\[
\Disc(g_q)=\prod_{0\le i<j\le q}(\mu_j-\mu_i)^2.
\]
注意
\(
\mu_j-\mu_i=\varphi^{\,q}r^i(r^{j-i}-1)
\)，故 Vandermonde 乘积可分解为
\[
\prod_{0\le i<j\le q}(\mu_j-\mu_i)
=\varphi^{\,q\sum_{i=0}^{q}(q-i)}\,
r^{\sum_{i=0}^{q}i(q-i)}\,
\prod_{k=1}^{q}(r^k-1)^{q+1-k}.
\]
平方后得
\[
\Disc(g_q)
=\varphi^{\,2q\sum_{i=0}^{q}(q-i)}\,
r^{2\sum_{i=0}^{q}i(q-i)}\,
\prod_{k=1}^{q}(r^k-1)^{2(q+1-k)}.
\]
其中
\[
\sum_{i=0}^{q}(q-i)=\frac{q(q+1)}2,
\qquad
\sum_{i=0}^{q}i(q-i)=\frac{q(q+1)(q-1)}6.
\]
又因 \(r=-\varphi^{-2}\) 且 \(\psi^k=(-1)^k\varphi^{-k}\)，Binet 恒等式给出
\[
1-(-\varphi^{-2})^k
=\frac{\varphi^k-\psi^k}{\varphi^k}
=\frac{\sqrt5\,F_k}{\varphi^k},
\]
从而
\[
(r^k-1)^2=(1-(-\varphi^{-2})^k)^2=\frac{5F_k^2}{\varphi^{2k}}.
\]
代入并合并指数可见所有 \(\varphi\)-幂严格消去，于是得到所示闭式。证毕。
\end{proof}

\begin{theorem}[秩一耦合的 Sylvester 结果式：谱差积的 Fibonacci 分解]\label{thm:xi-symq-kj-resultant-fibonacci}\leanverified{paper\_xi\_symq\_kj\_resultant\_fibonacci}
令
\[
J=\begin{pmatrix}1&1\\[2pt]1&1\end{pmatrix},
\qquad
E_q:=\mathrm{Sym}^q(J),
\qquad
A_q:=B_q+E_q,
\qquad
p_q(z):=\det(zI_{q+1}-A_q).
\]
则对一切 \(q\ge 2\)，成立严格恒等式
\[
\operatorname{Res}_z(g_q,p_q)
=(-1)^{\frac{(q+1)(q+2)}2}
\left(\prod_{i=0}^{q}\binom{q}{i}\right)
\left(\prod_{k=1}^{q}F_k^{\,2(q+1-k)}\right).
\]
\end{theorem}
\begin{proof}
记 \(s:=(1,1)^{\mathsf T}\)。由于 \(J=ss^{\mathsf T}\)，有
\[
E_q=\mathrm{Sym}^q(J)=u\,u^{\mathsf T},
\qquad
u:=s^{\otimes q}\in\mathrm{Sym}^q(\RR^2).
\]
取 \(B_q\) 的一组标准正交特征基 \(\{e_i\}_{i=0}^{q}\) 使 \(B_q e_i=\mu_i e_i\)。
将 \(u\) 展开为 \(u=\sum_{i=0}^{q}u_i e_i\)，则在该特征基下 \(A_q=\mathrm{diag}(\mu_i)+uu^{\mathsf T}\)。
由矩阵行列式引理得对任意 \(z\) 有
\[
p_q(z)=g_q(z)\Bigl(1-\sum_{i=0}^{q}\frac{u_i^2}{z-\mu_i}\Bigr).
\]
因此对每个根 \(\mu_i\) 有
\[
p_q(\mu_i)=-u_i^2\,g_q'(\mu_i).
\]
由结果式的根值表述，
\[
\operatorname{Res}_z(g_q,p_q)=\prod_{i=0}^{q}p_q(\mu_i)
=(-1)^{q+1}\Bigl(\prod_{i=0}^{q}u_i^2\Bigr)\Bigl(\prod_{i=0}^{q}g_q'(\mu_i)\Bigr).
\]
又因 \(g_q(z)=\prod_{i=0}^{q}(z-\mu_i)\) 首一，
\[
\prod_{i=0}^{q}g_q'(\mu_i)
=\prod_{i=0}^{q}\prod_{j\ne i}(\mu_i-\mu_j)
=(-1)^{\frac{q(q+1)}2}\Disc(g_q).
\]
为计算 \(\prod u_i^2\)，取 \(K\) 的单位正交本征向量 \(e_+,e_-\)（对应 \(\varphi,\psi\)），并记
\[
c_+:=\langle s,e_+\rangle,\qquad c_-:=\langle s,e_-\rangle.
\]
则由对称幂的二项展开，在 \(\mathrm{Sym}^q(\RR^2)\) 的标准正交对称本征基下有
\[
u_i=\sqrt{\binom{q}{i}}\,c_+^{\,q-i}c_-^{\,i},
\qquad 0\le i\le q.
\]
故
\[
\prod_{i=0}^{q}u_i^2
=\left(\prod_{i=0}^{q}\binom{q}{i}\right)
c_+^{\,2\sum_{i=0}^{q}(q-i)}c_-^{\,2\sum_{i=0}^{q}i}
=\left(\prod_{i=0}^{q}\binom{q}{i}\right)(c_+c_-)^{q(q+1)}.
\]
而 \(c_+^2=1+\frac{2}{\sqrt5}\)、\(c_-^2=1-\frac{2}{\sqrt5}\)，从而
\((c_+c_-)^2=\frac15\)，即 \(c_+c_-=5^{-1/2}\)。
因此
\[
\prod_{i=0}^{q}u_i^2
=\left(\prod_{i=0}^{q}\binom{q}{i}\right)5^{-\frac{q(q+1)}2}.
\]
将其与定理 \ref{thm:xi-symq-k-discriminant-fibonacci-vandermonde} 的 \(\Disc(g_q)\) 代回，并合并符号指数
\[
(-1)^{q+1+\frac{q(q+1)}2}=(-1)^{\frac{(q+1)(q+2)}2},
\]
即得所示结果式恒等式。证毕。
\end{proof}

\begin{proposition}[Hankel--Sylvester 桥接：谱差结果式的幂矩行列式模型]\label{prop:xi-symq-hankel-sylvester-bridge}\leanverified{paper\_xi\_symq\_hankel\_sylvester\_bridge}
令
\[
H_q:=\bigl(F_{i+j+3}^{\,q}\bigr)_{0\le i,j\le q}.
\]
则对一切 \(q\ge 2\)，成立
\[
\operatorname{Res}_z(g_q,p_q)
=(-1)^{\frac{(q+1)(q+2)}2}\det(H_q).
\]
\end{proposition}
\begin{proof}
由幂矩塌缩 \eqref{eq:xi-replica-softcore-fibonacci-moment-collapse}，对任意 \(n\ge 0\) 有
\[
F_{n+3}^{\,q}=2^{n}\sum_{i=0}^{q}w_i^{(q)}\,d_i^{\,n},
\qquad
d_i=\mu_i/2,
\]
其中 \(w_i^{(q)}>0\) 与 \(d_i\) 见定理 \ref{thm:xi-replica-softcore-fibonacci-moment-collapse}。
因此
\[
H_q
=D\cdot \Bigl(\sum_{i=0}^{q}w_i^{(q)}\,v(d_i)v(d_i)^{\mathsf T}\Bigr)\cdot D,
\qquad
D:=\mathrm{diag}(2^0,2^1,\dots,2^{q}),
\]
其中 \(v(d):=(1,d,\dots,d^{q})^{\mathsf T}\)。
由 Gram--Vandermonde 行列式公式得到
\[
\det(H_q)
=2^{q(q+1)}\Bigl(\prod_{i=0}^{q}w_i^{(q)}\Bigr)\prod_{0\le i<j\le q}(d_j-d_i)^2.
\]
将 \(d_i=\mu_i/2\) 代入并整理得
\[
\det(H_q)
=\Bigl(\prod_{i=0}^{q}w_i^{(q)}\Bigr)\prod_{0\le i<j\le q}(\mu_j-\mu_i)^2
=\Bigl(\prod_{i=0}^{q}w_i^{(q)}\Bigr)\Disc(g_q).
\]
而 \(\prod_{i=0}^{q}w_i^{(q)}=\bigl(\prod_{i=0}^{q}\binom{q}{i}\bigr)\alpha^{\frac{q(q+1)}2}\beta^{\frac{q(q+1)}2}\)
且 \(\alpha\beta=\frac15\)，从而
\(\prod w_i^{(q)}=\bigl(\prod_{i=0}^{q}\binom{q}{i}\bigr)5^{-\frac{q(q+1)}2}\)。
结合定理 \ref{thm:xi-symq-k-discriminant-fibonacci-vandermonde} 即得
\[
\det(H_q)
=\left(\prod_{i=0}^{q}\binom{q}{i}\right)\left(\prod_{k=1}^{q}F_k^{\,2(q+1-k)}\right).
\]
与定理 \ref{thm:xi-symq-kj-resultant-fibonacci} 对比即得所示桥接恒等式。证毕。
\end{proof}

\begin{theorem}[Hankel 秩与线性递推：最小阶的可检验刻画]\label{thm:xi-hankel-rank-linear-recurrence}
令 \((a_n)_{n\ge 0}\) 为复数列，定义无限 Hankel 矩阵
\[
H=\bigl(a_{i+j}\bigr)_{i,j\ge 0},
\]
以及主子式
\[
\Delta_m:=\det\bigl(a_{i+j}\bigr)_{0\le i,j\le m}\qquad(m\ge 0).
\]
则下列条件等价：
\begin{enumerate}
  \item 存在整数 \(k\ge 1\) 与常数 \(c_0,\dots,c_{k-1}\in\CC\)，使得对一切 \(n\ge 0\) 有
  \[
  a_{n+k}=\sum_{j=0}^{k-1}c_j\,a_{n+j};
  \]
  \item \(\mathrm{rank}(H)\le k\)；
  \item 对所有 \(m\ge k\) 有 \(\Delta_m=0\)。
\end{enumerate}
并且最小可行递推阶数恰等于 \(\mathrm{rank}(H)\)（等价地：恰为使得 \(\Delta_m=0\) 对所有 \(m\ge k\) 成立的最小 \(k\)）。
\leanverified{paper\_xi\_hankel\_rank\_linear\_recurrence}
\end{theorem}
\begin{proof}
\((1)\Rightarrow(2)\)：递推给出对每个 \(j\ge k\)，Hankel 的第 \(j\) 列是前 \(k\) 列的线性组合，故秩至多 \(k\)。
\((2)\Rightarrow(3)\)：秩至多 \(k\) 等价于任意 \((k+1)\times(k+1)\) 子式为零，尤其所有主子式 \(\Delta_m\) 在 \(m\ge k\) 时为零。
\((3)\Rightarrow(1)\)：若所有 \((k+1)\times(k+1)\) 主子式为零，则 Hankel 前 \(k+1\) 列线性相关，存在不全为零的 \((\gamma_0,\dots,\gamma_k)\) 使
\[
\sum_{j=0}^k \gamma_j a_{n+j}=0\qquad(\forall n\ge 0),
\]
从而得到阶至多 \(k\) 的递推。最小阶数等于 Hankel 秩由列空间的最小生成元长度与秩的定义直接推出。证毕。
\end{proof}

\begin{proposition}[耦合参数的主因子定律]\label{prop:xi-symq-rankone-coupling-t-power}
\leanverified{paper\_xi\_symq\_rankone\_coupling\_t\_power}
对任意 \(t\in\CC\)，定义
\[
p_{q,t}(z):=\det\!\bigl(zI_{q+1}-(B_q+tE_q)\bigr).
\]
则对一切 \(q\ge 2\) 有
\[
\operatorname{Res}_z(g_q,p_{q,t})
=t^{\,q+1}\operatorname{Res}_z(g_q,p_{q,1}).
\]
\end{proposition}
\begin{proof}
沿用定理 \ref{thm:xi-symq-kj-resultant-fibonacci} 的特征基与 \(u\) 的记号，有
\(
B_q+tE_q=\mathrm{diag}(\mu_i)+tuu^{\mathsf T}
\)。
故对每个 \(\mu_i\) 成立 \(p_{q,t}(\mu_i)=-t\,u_i^2g_q'(\mu_i)\)，从而
\[
\operatorname{Res}_z(g_q,p_{q,t})
=\prod_{i=0}^{q}p_{q,t}(\mu_i)
=t^{q+1}\prod_{i=0}^{q}p_{q,1}(\mu_i)
=t^{q+1}\operatorname{Res}_z(g_q,p_{q,1}).
\]
证毕。
\end{proof}

\begin{proposition}[谱差积的 \(p\)-进估值公式：数字和项与 Fibonacci 赋值项]\label{prop:xi-symq-resultant-padic-valuation}
\leanverified{paper\_xi\_symq\_resultant\_padic\_valuation}
令
\[
R_q:=\bigl|\operatorname{Res}_z(g_q,p_q)\bigr|
=\left(\prod_{i=0}^{q}\binom{q}{i}\right)\left(\prod_{k=1}^{q}F_k^{\,2(q+1-k)}\right).
\]
则对任意素数 \(p\) 有
\[
\nu_p(R_q)=\sum_{i=0}^{q}\nu_p\!\binom{q}{i}
\;+\;\sum_{k=1}^{q}2(q+1-k)\,\nu_p(F_k),
\]
并且
\[
\sum_{i=0}^{q}\nu_p\!\binom{q}{i}
=\frac{2\sum_{n=0}^{q}s_p(n)-(q+1)s_p(q)}{p-1},
\]
其中 \(s_p(n)\) 表示 \(n\) 的 \(p\)-进展开的数字和。
\end{proposition}
\begin{proof}
第一式由 \(R_q\) 的乘积结构与 \(p\)-进估值的可加性直接得到。
第二式由 Legendre 公式
\(
\nu_p(n!)=\frac{n-s_p(n)}{p-1}
\)
推出
\[
\nu_p\binom{q}{i}
=\nu_p(q!)-\nu_p(i!)-\nu_p((q-i)!)
=\frac{s_p(i)+s_p(q-i)-s_p(q)}{p-1},
\]
对 \(i=0,\dots,q\) 求和并利用 \(\sum_{i=0}^{q}s_p(q-i)=\sum_{i=0}^{q}s_p(i)\) 即得。证毕。
\end{proof}

\begin{proposition}[判别式/结果式比值的纯 \(5\)-幂缺陷]\label{prop:xi-symq-disc-over-resultant-pure-5-defect}
对一切 \(q\ge 2\) 成立
\[
\frac{\Disc(g_q)}{\operatorname{Res}_z(g_q,p_q)}
=(-1)^{\frac{(q+1)(q+2)}2}\cdot
\frac{5^{\frac{q(q+1)}2}}{\displaystyle\prod_{i=0}^{q}\binom{q}{i}}.
\]
\end{proposition}
\begin{proof}
由定理 \ref{thm:xi-symq-k-discriminant-fibonacci-vandermonde} 与定理 \ref{thm:xi-symq-kj-resultant-fibonacci} 直接相除即可；
其中 \(\prod_{k=1}^{q}F_k^{\,2(q+1-k)}\) 在该比值中严格消去。证毕。
\end{proof}

\endinput


\begin{theorem}[Perron 根的纯 Fibonacci 固定点方程]\label{thm:xi-replica-softcore-perron-fibonacci-fixed-point}
记 \(\rho_q:=\rho(A^{(q)})\) 为 \(\left.T^{(q)}\right|_{\mathcal H_q}\) 的 Perron 特征值。
则存在唯一实数 \(x_q>1\) 使
\[
\rho_q=2^{q-1}x_q,
\]
并且 \(x_q\) 满足严格收敛的固定点方程
\begin{equation}\label{eq:xi-replica-softcore-perron-fixed-point}
x_q
=1+\sum_{m=1}^{\infty}\left(\frac{F_{m+3}}{2^{m+1}}\right)^{q}\,x_q^{-m}.
\end{equation}
该刻画在每个 \(q\ge 2\) 下给出唯一正根。
\end{theorem}
\begin{proof}
由定理 \ref{thm:xi-replica-softcore-exceptional-weighted-secular}，Perron 根 \(\rho_q\) 为割线方程
\begin{equation}\label{eq:xi-replica-softcore-perron-secular}
\sum_{i=0}^{q}\frac{w_i^{(q)}}{\rho_q-d_i}=2
\end{equation}
在顶端区间的唯一解，其中 \(d_i=\mu_i^{(q)}/2\)。
并且由 Rayleigh 商下界（取向量 \(u=D\mathbf 1\) 于相似对称表示 \(\widetilde A^{(q)}\) 上计算）可得
\(
\rho_q\ge (4^q+3^q)/2^{q+1}
=2^{q-1}+3^q/2^{q+1}
\)；
从而 \(\rho_q>\max_{0\le i\le q}d_i=\varphi^q/2\)，并对每个 \(i\) 有绝对收敛展开
\[
\frac{1}{\rho_q-d_i}
=\frac{1}{\rho_q}\sum_{m=0}^{\infty}\left(\frac{d_i}{\rho_q}\right)^{m}.
\]
将其代入 \eqref{eq:xi-replica-softcore-perron-secular} 并由一致收敛交换求和次序，得到
\[
2
=\frac{1}{\rho_q}\sum_{m=0}^{\infty}\rho_q^{-m}\sum_{i=0}^{q} w_i^{(q)}\,d_i^{\,m}.
\]
由定理 \ref{thm:xi-replica-softcore-fibonacci-moment-collapse}，
\(
\sum_i w_i^{(q)}d_i^{\,m}=2^{-m}F_{m+3}^q
\)，故
\[
2
=\frac{1}{\rho_q}\sum_{m=0}^{\infty}\rho_q^{-m}\,2^{-m}F_{m+3}^{\,q}.
\]
令 \(\rho_q=2^{q-1}x_q\)，并将 \(m=0\) 项（\(F_3=2\)）分离，即得
\[
2
=\frac{2}{x_q}+ \frac{2}{x_q}\sum_{m=1}^{\infty}\left(\frac{F_{m+3}}{2^{m+1}}\right)^{q}x_q^{-m}.
\]
两边同乘 \(x_q/2\) 得到固定点方程 \eqref{eq:xi-replica-softcore-perron-fixed-point}。

右端作为 \(x\) 的函数在 \((1,\infty)\) 上严格递减，而左端 \(x\mapsto x\) 严格递增，故解唯一。证毕。
\end{proof}

\begin{corollary}[Perron 根的两项主导展开与可控余项]\label{cor:xi-replica-softcore-perron-two-term-expansion}
令
\[
b_m:=\frac{F_{m+3}}{2^{m+1}}\qquad(m\ge 1),
\qquad
\varepsilon_q:=\sum_{m=1}^{\infty}b_m^{\,q}.
\]
则 \(x_q\in(1,1+\varepsilon_q)\)，并且
\begin{equation}\label{eq:xi-replica-softcore-perron-two-term-expansion}
\rho_q
=2^{q-1}+\frac{3^q}{2^{q+1}}+R_q,
\end{equation}
其中余项满足显式上界
\begin{equation}\label{eq:xi-replica-softcore-perron-remainder-bound}
0<R_q
<2^{q-1}\sum_{m=2}^{\infty}b_m^{\,q}
\le 2^{q-1}\cdot \frac{b_2^{\,q}}{1-(\varphi/2)^q}
=\frac{5^q}{2^{2q+1}}\cdot\frac{1}{1-(\varphi/2)^q}.
\end{equation}
因此
\(
\rho_q
=2^{q-1}+\frac{3^q}{2^{q+1}}+O\!\bigl(5^q/2^{2q}\bigr)
\)
当 \(q\to\infty\)。
\end{corollary}
\begin{proof}
由定理 \ref{thm:xi-replica-softcore-perron-fibonacci-fixed-point}，
\[
x_q-1=\sum_{m=1}^{\infty} b_m^{\,q}x_q^{-m}.
\]
因 \(x_q>1\)，即得 \(0<x_q-1<\sum_{m\ge 1}b_m^q=\varepsilon_q\)，从而 \(x_q\in(1,1+\varepsilon_q)\)。
将首项 \(m=1\) 分离并乘以 \(2^{q-1}\) 得
\[
\rho_q
=2^{q-1}+2^{q-1}\frac{b_1^{\,q}}{x_q}
2^{q-1}\sum_{m=2}^{\infty}b_m^{\,q}x_q^{-m}.
\]
由于 \(b_1=F_4/2^2=3/4\) 且 \(x_q>1\)，有
\(
2^{q-1}b_1^q/x_q<2^{q-1}b_1^q=3^q/2^{q+1}
\)。
将差额并入余项并注意最后一项为正，得到展开式 \eqref{eq:xi-replica-softcore-perron-two-term-expansion}，且
\[
0<R_q<2^{q-1}\sum_{m=2}^{\infty}b_m^{\,q}x_q^{-m}\le 2^{q-1}\sum_{m=2}^{\infty}b_m^{\,q}.
\]
再由 Fibonacci 比值单调性 \(F_{n+1}/F_n<\varphi\)（\(n\ge 1\)）推出
\[
\frac{b_{m+1}}{b_m}=\frac{F_{m+4}}{2F_{m+3}}<\frac{\varphi}{2},
\]
故 \(\sum_{m\ge 2}b_m^q\) 由等比级数控制，得到 \eqref{eq:xi-replica-softcore-perron-remainder-bound}。证毕。
\end{proof}

\begin{theorem}[Fibonacci 幂 Hankel 行列式的完全乘积公式]\label{thm:xi-replica-softcore-fibonacci-power-hankel-determinant}
定义 \((q+1)\times(q+1)\) Hankel 矩阵
\[
H_q:=\bigl(F_{i+j+3}^{\,q}\bigr)_{0\le i,j\le q}.
\]
则对一切 \(q\ge 1\) 成立精确闭式
\begin{equation}\label{eq:xi-replica-softcore-fibonacci-power-hankel-determinant}
\det(H_q)
=\left(\prod_{i=0}^{q}\binom{q}{i}\right)
\left(\prod_{k=1}^{q}F_k^{\,2(q+1-k)}\right).
\end{equation}
右端为整数的纯乘积表达，且不含任何代数数域延拓。
\end{theorem}
\begin{proof}
由 Binet 形式与 \(\alpha=\varphi^3/\sqrt5\)、\(\beta=-\psi^3/\sqrt5\)，有
\(
F_{n+3}=\alpha\varphi^n+\beta\psi^n
\)；
从而对任意 \(n\ge 0\),
\[
F_{n+3}^{\,q}
=\sum_{i=0}^{q}\binom{q}{i}\alpha^{q-i}\beta^{i}(\varphi^{q-i}\psi^{i})^{n}
=\sum_{i=0}^{q}c_i r_i^{\,n},
\]
其中 \(r_i:=\varphi^{q-i}\psi^{i}\) 且 \(c_i:=\binom{q}{i}\alpha^{q-i}\beta^{i}\)。
令 \(V:=(r_i^{\,j})_{0\le i,j\le q}\) 为 Vandermonde 型矩阵，并令 \(C:=\mathrm{diag}(c_0,\dots,c_q)\)。
则 \(H_q=V^{\mathsf T}CV\)，从而
\begin{equation}\label{eq:xi-replica-softcore-hankel-det-vandermonde-step}
\det(H_q)=(\det V)^2\prod_{i=0}^{q}c_i
=\left(\prod_{0\le i<j\le q}(r_j-r_i)^2\right)\left(\prod_{i=0}^{q}c_i\right).
\end{equation}

先算系数乘积。
由 \(\sum_{i=0}^{q} i=\sum_{i=0}^{q}(q-i)=q(q+1)/2\) 与
\(
\alpha\beta=(\varphi^3/\sqrt5)(-\psi^3/\sqrt5)=-(\varphi\psi)^3/5=1/5
\),
得到
\[
\prod_{i=0}^{q}c_i
=\left(\prod_{i=0}^{q}\binom{q}{i}\right)(\alpha\beta)^{q(q+1)/2}
=\left(\prod_{i=0}^{q}\binom{q}{i}\right)5^{-q(q+1)/2}.
\]

再算 Vandermonde 因子。
注意 \(r_i=\varphi^q(\psi/\varphi)^i=\varphi^q r^i\)，其中 \(r:=\psi/\varphi=-\varphi^{-2}\)。
于是
\[
\prod_{0\le i<j\le q}(r_j-r_i)^2
=(\varphi^q)^{q(q+1)}\prod_{0\le i<j\le q}(r^j-r^i)^2.
\]
而几何节点的 Vandermonde 恒等式给出
\[
\prod_{0\le i<j\le q}(r^j-r^i)^2
=\prod_{k=1}^{q}(r^k-1)^{2(q+1-k)}\cdot r^{\,2\sum_{0\le i<j\le q}i}.
\]
对每个 \(k\ge 1\)，有
\[
1-r^k
=1-(-\varphi^{-2})^k
=\frac{\varphi^{2k}-(-1)^k}{\varphi^{2k}}
=\frac{\varphi^k(\varphi^k-\psi^k)}{\varphi^{2k}}
=\frac{\sqrt5\,F_k}{\varphi^k},
\]
故
\(
(r^k-1)^2=(1-r^k)^2=5F_k^2/\varphi^{2k}
\)。
将其代回并整理 \(\varphi\) 的幂次，可验证所有 \(\varphi\) 因子完全相消，从而得到
\[
\prod_{0\le i<j\le q}(r_j-r_i)^2
=5^{q(q+1)/2}\prod_{k=1}^{q}F_k^{\,2(q+1-k)}.
\]
与 \eqref{eq:xi-replica-softcore-hankel-det-vandermonde-step} 中的 \(5^{-q(q+1)/2}\) 系数因子精确抵消，得到 \eqref{eq:xi-replica-softcore-fibonacci-power-hankel-determinant}。证毕。
\end{proof}

\begin{corollary}[交错几何点的 Vandermonde 平方塌缩为 Fibonacci 乘积]\label{cor:xi-replica-softcore-alternating-geom-vandermonde}
令
\[
\xi_i:=(-1)^i\varphi^{\,q-2i}\qquad(0\le i\le q).
\]
则成立严格恒等式
\begin{equation}\label{eq:xi-replica-softcore-alternating-geom-vandermonde}
\prod_{0\le i<j\le q}(\xi_j-\xi_i)^2
=5^{q(q+1)/2}\prod_{k=1}^{q}F_k^{\,2(q+1-k)}\in\ZZ.
\end{equation}
因此左端虽表面上属于 \(\QQ(\sqrt5)\) 的代数整数乘积，其平方判别式却在 \(\ZZ\) 中塌缩为完全显式的 Fibonacci 乘积。
\end{corollary}
\begin{proof}
在定理 \ref{thm:xi-replica-softcore-fibonacci-power-hankel-determinant} 的证明中，节点 \(r_i=\varphi^{q-i}\psi^i\) 满足
\(
r_i=(-1)^i\varphi^{q-2i}=\xi_i
\)。
因此其中已得的 Vandermonde 平方闭式即为 \eqref{eq:xi-replica-softcore-alternating-geom-vandermonde}。证毕。
\end{proof}

\begin{theorem}[严格交错与全矩阵惯性平衡]\label{thm:xi-replica-softcore-full-inertia-balance}
设
\[
d_i:=\frac{\mu_i^{(q)}}2\qquad(0\le i\le q),
\]
按降序重排为
\(
d_{(0)}>d_{(1)}>\cdots>d_{(q)}
\)。
则例外谱可重排为
\(
\lambda_{(0)}>\lambda_{(1)}>\cdots>\lambda_{(q)}
\)
并满足严格交错
\[
\lambda_{(0)}>d_{(0)}>\lambda_{(1)}>d_{(1)}>\cdots>\lambda_{(q)}>d_{(q)}.
\]
进一步有：
\begin{enumerate}
  \item 例外谱负特征值个数为 \(\lfloor(q+1)/2\rfloor\)，正特征值个数为 \(\lceil(q+1)/2\rceil\)；
  \item 全矩阵惯性满足
  \[
  \operatorname{In}\!\bigl(T^{(q)}\bigr)=\bigl(2^{q-1},\,2^{q-1},\,0\bigr).
  \]
\end{enumerate}
\end{theorem}
\begin{proof}
交错性由定理 \ref{thm:xi-replica-softcore-exceptional-weighted-secular} 的割线函数直接推出：
在每个相邻极点区间与顶端区间，方程恰有一根。
例外谱符号计数与零点排除亦可由推论 \ref{cor:xi-replica-softcore-exceptional-inertia} 直接得到。

体谱中奇 \(i\) 给出负特征值，偶 \(i\) 给出正特征值，故
\[
N_-^{\mathrm{bulk}}
=
\sum_{\substack{0\le i\le q\\ i\ \mathrm{odd}}}
\left(\binom qi-1\right)
=
2^{q-1}-\left\lfloor\frac{q+1}{2}\right\rfloor,
\]
\[
N_+^{\mathrm{bulk}}
=
\sum_{\substack{0\le i\le q\\ i\ \mathrm{even}}}
\left(\binom qi-1\right)
=
2^{q-1}-\left\lceil\frac{q+1}{2}\right\rceil.
\]
再加上例外谱计数
\[
N_-^{\mathrm{exc}}=\left\lfloor\frac{q+1}{2}\right\rfloor,\qquad
N_+^{\mathrm{exc}}=\left\lceil\frac{q+1}{2}\right\rceil,
\]
得到
\[
N_-=2^{q-1},\qquad N_+=2^{q-1}.
\]
由前述零点排除，\(N_0=0\)。故惯性公式成立。证毕。
\end{proof}

\begin{proposition}[体谱特征空间的零和超平面方程]\label{prop:xi-replica-softcore-bulk-eigenspace-zero-sum-hyperplane}
取 \(K\) 的单位正交本征向量 \(e_+,e_-\)，并对任意 \(S\subseteq[q]\) 定义
\[
v_S:=\bigotimes_{r=1}^{q}\xi_r,\qquad
\xi_r=
\begin{cases}
e_-,&r\in S,\\
e_+,&r\notin S.
\end{cases}
\]
当 \(|S|=i\) 时，\(\{v_S\}\) 构成 \(E_i\) 的一组正交基。
并且
\[
\ker\!\left(T^{(q)}-\frac{\mu_i^{(q)}}2I\right)
=
\left\{
\sum_{|S|=i}a_Sv_S:\ 
\sum_{|S|=i}a_S\langle v_S,\mathbf1\rangle=0
\right\}
=E_i\cap\mathbf1^\perp,
\]
其维数为 \(\binom qi-1\)。
特别地，由 \(\langle v_S,\mathbf1\rangle\) 仅依赖于 \(|S|=i\) 且非零，上式等价为
\[
\sum_{|S|=i}a_S=0.
\]
\end{proposition}
\begin{proof}
张量积谱给出 \(\{v_S:|S|=i\}\) 为 \(E_i\) 的正交基。
对任意 \(x\in E_i\)，
\[
T^{(q)}x=\frac12\bigl(D^{(q)}x+J_qx\bigr)=\frac{\mu_i^{(q)}}2x+\frac12J_qx.
\]
因此
\[
\left(T^{(q)}-\frac{\mu_i^{(q)}}2I\right)x=0
\iff J_qx=0
\iff x\perp\mathbf1.
\]
写 \(x=\sum_{|S|=i}a_Sv_S\) 即得线性约束式。
由于 \(\mathbf1\) 在 \(E_i\) 上投影非零，该约束为一条独立线性方程，故维数为 \(\binom qi-1\)。证毕。
\end{proof}

\endinput


\paragraph{二级退化五次的最大对称性、子结式闭式与特征坍缩门槛}
沿用结论 \ref{con:xi-terminal-zm-leyang-delta-quartic-discriminant-reflexive} 的四次 \(\delta\)-生成多项式
\[
F_\delta(T;y):=T^{4}+(8y-3)T^{3}+(2y^{2}-34y-4)T^{2}-y(y-1)\,P_{\mathrm{LY}}(y)\in\ZZ[y][T],
\qquad
P_{\mathrm{LY}}(y)=256y^{3}+411y^{2}+165y+32,
\]
及其判别式分解
\[
\mathrm{Disc}_{T}\bigl(F_\delta(T;y)\bigr)=-y(y-1)\,P_{\mathrm{LY}}(y)\cdot Q_{5}(y)^{2},
\]
其中
\[
Q_{5}(y)=4096y^{5}+5376y^{4}-464y^{3}-2749y^{2}-723y+80\in\ZZ[y].
\]

\begin{theorem}[二级退化五次的全对称伽罗瓦群]\label{thm:xi-q5-galois-s5}\leanverified{paper\_xi\_q5\_galois\_s5}
多项式 \(Q_5(y)\) 在 \(\QQ\) 上不可约，且其分裂域 \(L_5/\QQ\) 的伽罗瓦群满足
\[
\operatorname{Gal}(L_5/\QQ)\cong S_5.
\]
\end{theorem}
\begin{proof}
不可约性可由一次良性模素约化读出：在 \(\FF_7[y]\) 上，\(Q_5\bmod 7\) 仍为次数 \(5\) 的不可约多项式，
故 \(Q_5\) 在 \(\QQ[y]\) 不可约。

由结论 \ref{con:xi-terminal-zm-leyang-q5-discriminant-727}，
\[
\mathrm{Disc}(Q_5)=-2^{28}\cdot 3^{9}\cdot 11^{2}\cdot 31^{3}\cdot 727,
\]
从而 \(7\nmid \mathrm{Disc}(Q_5)\)，素数 \(7\) 在分裂域中不分歧。
由于 \(Q_5\bmod 7\) 为不可约 \(5\) 次因子，Dedekind 分解型与 Frobenius 循环型对应给出
\(\operatorname{Gal}(L_5/\QQ)\) 含一个 \(5\)-循环。

另一方面，\(\nu_{727}(\mathrm{Disc}(Q_5))=1\)。
并且在 \(\FF_{727}[y]\) 上出现唯一二重根型分解：
存在 \(a\neq b\in\FF_{727}\) 与不可约二次因子 \(h_2\in\FF_{727}[y]\) 使
\[
Q_5(y)\equiv (y-a)^2(y-b)\,h_2(y)\quad(\bmod 727).
\]
因此 \(727\) 处的惯性群包含一个换位。
在 \(S_5\) 中，一个 \(5\)-循环与一个换位生成 \(S_5\)，故伽罗瓦群为 \(S_5\)。
可复现核验脚本：\texttt{scripts/exp\_xi\_leyang\_q5\_a7\_subresultant\_galois\_audit.py}。证毕。
\end{proof}

\begin{corollary}[二级退化参数的不可根式性]\label{cor:xi-q5-unsolvable-by-radicals}
\leanverified{paper\_xi\_q5\_unsolvable\_by\_radicals}
若 \(\alpha\) 为 \(Q_5\) 的任一根，则 \(\QQ(\alpha)/\QQ\) 不为可解扩张，因而 \(\alpha\) 不可由有理数经有限次四则运算与开方表示。
\end{corollary}
\begin{proof}
由定理 \ref{thm:xi-q5-galois-s5}，分裂域伽罗瓦群为不可解群 \(S_5\)，与根式可解性判别矛盾。证毕。
\end{proof}

\begin{theorem}[一阶子结式闭式与唯一二重根刚性]\label{thm:xi-fdelta-subresultant-unique-double-root}
令 \(f(T):=F_\delta(T;y)\)、\(g(T):=\partial_TF_\delta(T;y)\)。
则二者关于 \(T\) 的第一子结式满足严格闭式
\[
\mathrm{SRes}_1(f,g)(T)
=
-4A_7(y)\,T
-y(y-1)(8y-3)(16y+11)(32y+19)\,P_{\mathrm{LY}}(y),
\]
其中
\[
A_7(y):=45056 y^7 + 60160 y^6 - 11984 y^5 - 19765 y^4 + 18906 y^3 + 14723 y^2 + 2216 y + 200.
\]
并且在特征 \(0\) 上成立：
\begin{enumerate}
  \item 若 \(y_0\in\overline{\QQ}\) 使 \(\mathrm{Disc}_T(F_\delta(\cdot;y_0))=0\)，则 \(F_\delta(T;y_0)\) 恰有且仅有一个二重根，其余两根为单根；
  \item 该唯一二重根 \(\delta_\ast(y_0)\) 满足（当 \(A_7(y_0)\neq 0\) 时）
  \[
  \delta_\ast(y)
  =
  -\frac{y(y-1)(8y-3)(16y+11)(32y+19)\,P_{\mathrm{LY}}(y)}{4A_7(y)},
  \]
  且当 \(y\in\{0,1\}\cup\{P_{\mathrm{LY}}(y)=0\}\) 时 \(\delta_\ast(y)=0\)。
\end{enumerate}
\end{theorem}
\begin{proof}
子结式理论给出：当 \(\gcd(f,g)\) 在 \(\overline{\QQ}(y)[T]\) 中次数为 \(1\) 时，
\(\mathrm{SRes}_1(f,g)\) 与该一次 \(\gcd\) 在 \(\overline{\QQ}(y)^\times\) 意义下成比例；
因此其唯一零点即为二重根。
所示闭式为直接子结式计算的恒等式。

若存在 \(y_0\) 使 \(F_\delta(\cdot;y_0)\) 发生更高碰撞（出现三重根或两处不同二重根），则 \(\gcd(f,g)\) 的次数至少为 \(2\)，
从而 \(\mathrm{SRes}_1(f,g)\) 作为 \(T\)-一次多项式必须恒等为零，等价于
\[
A_7(y_0)=0,
\qquad
y_0(y_0-1)(8y_0-3)(16y_0+11)(32y_0+19)P_{\mathrm{LY}}(y_0)=0.
\]
但在 \(\ZZ[y]\) 中有
\[
\gcd\!\bigl(A_7,\ y(y-1)P_{\mathrm{LY}}\bigr)=1,
\qquad
\gcd(A_7,Q_5)=1,
\]
故上述同时为零不可能发生。
于是在判别式为零处 \(\gcd(f,g)\) 只能为一次，得到“唯一二重根”结论。
解一次方程给出 \(\delta_\ast\) 的闭式；当 \(y\in\{0,1\}\cup\{P_{\mathrm{LY}}=0\}\) 时分子为零，故 \(\delta_\ast=0\)。
可复现核验脚本：\texttt{scripts/exp\_xi\_leyang\_q5\_a7\_subresultant\_galois\_audit.py}。证毕。
\end{proof}

\begin{theorem}[\texorpdfstring{$P_{\mathrm{LY}}$}{PLY} 与 \texorpdfstring{$Q_5$}{Q5} 的干涉素数刻画]\label{thm:xi-ply-q5-resultant-interference-primes}
\leanverified{paper\_xi\_ply\_q5\_resultant\_interference\_primes}
记 \(\mathrm{PLY}(y):=P_{\mathrm{LY}}(y)\)。
则关于 \(y\) 的 Sylvester 结果式满足
\[
\operatorname{Res}_{y}\!\bigl(\mathrm{PLY}(y),Q_{5}(y)\bigr)
=2^{20}\cdot 3^{12}\cdot 31^{4}\cdot 73.
\]
因此对任意素数 \(p\nmid 2\cdot 3\cdot 31\cdot 73\)，有
\[
\gcd\!\bigl(\overline{\mathrm{PLY}},\overline{Q_{5}}\bigr)=1
\quad\text{于}\quad \FF_p[y],
\]
从而零点方案 \(V(\overline{\mathrm{PLY}})\) 与 \(V(\overline{Q_5})\) 在 \(\PP^1_{\FF_p}\) 上不相交。
\end{theorem}
\begin{proof}
由结果式判别准则，
\(
\gcd(\overline{\mathrm{PLY}},\overline{Q_5})\neq 1
\)
当且仅当
\(
p\mid \operatorname{Res}_{y}(\mathrm{PLY},Q_5)
\)。
直接计算 Sylvester 行列式并分解即得所示恒等式。
可复现核验脚本：\texttt{scripts/exp\_xi\_leyang\_ply\_q5\_a7\_interference\_geometry\_audit.py}。证毕。
\end{proof}

\begin{corollary}[奇素数上的公共根：仅 \texorpdfstring{$31$}{31} 与 \texorpdfstring{$73$}{73} 非平凡]\label{cor:xi-ply-q5-common-roots-prime31-73}
\leanverified{paper\_xi\_ply\_q5\_common\_roots\_prime31\_73}
令 \(p\ge 5\) 为素数。
则存在 \(\bar y\in\FF_p\) 使
\[
\mathrm{PLY}(\bar y)\equiv Q_5(\bar y)\equiv 0\pmod p
\]
当且仅当 \(p\in\{31,73\}\)。
并且：
\begin{enumerate}
  \item 若 \(p=73\)，唯一公共根为 \(\bar y\equiv -10\pmod{73}\)，且在两者中均为一次根；
  \item 若 \(p=31\)，唯一公共根为 \(\bar y\equiv 12\pmod{31}\)，且在两者中均为二次根。
\end{enumerate}
\end{corollary}
\begin{proof}
由定理 \ref{thm:xi-ply-q5-resultant-interference-primes}，对 \(p\ge 5\) 的公共根只可能出现在 \(p=31,73\)。
直接在有限域中分解可得：
\[
\overline{\mathrm{PLY}}=(y+10)(y+13)(y-4),\quad
\overline{Q_5}=(y+10)\cdot(\text{四次因子})\quad(\bmod 73),
\]
\[
\overline{\mathrm{PLY}}=(y-6)(y-12)^2,\quad
\overline{Q_5}=(y-4)(y-12)^2(y^2+8y-8)\quad(\bmod 31).
\]
故结论成立。证毕。
\end{proof}

\begin{theorem}[二级退化除子与主分支除子的交叉素数全集]\label{thm:xi-dprim-dsec-intersection-primes}
定义
\[
D_{\mathrm{prim}}(y):=y(y-1)\mathrm{PLY}(y),
\qquad
D_{\mathrm{sec}}(y):=Q_5(y).
\]
则存在 \(\bar y\in\FF_p\) 使
\[
D_{\mathrm{prim}}(\bar y)\equiv D_{\mathrm{sec}}(\bar y)\equiv 0\pmod p
\]
当且仅当
\[
p\in\{2,3,5,13,31,73\}.
\]
其中对奇素数 \(p\in\{5,13,31,73\}\) 的交点可取为
\[
(p,\bar y)=(5,0),\ (13,1),\ (31,12),\ (73,-10).
\]
\end{theorem}
\begin{proof}
首先
\[
Q_5(0)=80=2^4\cdot 5,\qquad
Q_5(1)=5616=2^4\cdot 3^3\cdot 13.
\]
故由 \(y=0,1\) 两条分支可得素数 \(2,3,5,13\)。
其次，由推论 \ref{cor:xi-ply-q5-common-roots-prime31-73}，\(\mathrm{PLY}\) 与 \(Q_5\) 的交叉只在 \(31,73\) 出现，对应根分别为 \(12,-10\)。
合并即得全集。证毕。
\end{proof}

\leanverified{paper\_xi\_fdelta\_char73\_transverse\_node}
\begin{theorem}[特征 \texorpdfstring{$73$}{73} 的横截干涉与普通节点]\label{thm:xi-fdelta-char73-transverse-node}
在 \(\FF_{73}\) 中取 \(\bar y\equiv -10\)，则
\[
\overline{F_\delta}(T;\bar y)=T^2(T-5)^2.
\]
令 \(\overline{C_\delta}\subset\mathbb A^2_{\FF_{73}}\) 由 \(\overline{F_\delta}(T;y)=0\) 定义，则：
\begin{enumerate}
  \item 点 \((T,y)=(5,-10)\) 为普通二重点（\(A_1\)-节点）；
  \item 点 \((T,y)=(0,-10)\) 为光滑点。
\end{enumerate}
\end{theorem}
\begin{proof}
代入 \(y=-10\) 直接得到分解式。
在点 \((5,-10)\) 处取局部坐标 \(t:=T-5,\ s:=y+10\)，则
\[
\overline{F_\delta}(t+5;s-10)
=25t^2+6ts-5s^2+\text{(总次数 }\ge 3\text{ 项)}.
\]
其二次主部判别式为
\[
6^2-4\cdot 25\cdot(-5)\equiv 25\not\equiv 0\pmod{73},
\]
故二次主部对应两条不同切线，奇点为普通节点。
另一方面在 \((0,-10)\) 处，
\[
\partial_T\overline{F_\delta}(0,-10)=0,\qquad
\partial_y\overline{F_\delta}(0,-10)\equiv 47\not\equiv 0\pmod{73},
\]
由 Jacobian 判据得该点光滑。证毕。
\end{proof}

\begin{theorem}[特征 \texorpdfstring{$31$}{31} 的唯一总坍缩纤维与几何 \texorpdfstring{$A_3$}{A3} 触节点]\label{thm:xi-fdelta-char31-total-collapse}
\leanverified{paper\_xi\_fdelta\_char31\_total\_collapse}
对奇素数 \(p\neq 2\)，考虑 \(F_\delta(T;y)\) 在 \(\FF_p\) 上的约化。
则：
\begin{enumerate}
  \item 存在 \(y_0\in\FF_p\) 使
  \[
  F_\delta(T;y_0)\equiv T^4\quad\text{于}\ \FF_p[T]
  \]
  当且仅当 \(p=31\)，且唯一解为 \(y_0\equiv 12\pmod{31}\)；
  \item 在 \(\overline{C_\delta}:\overline{F_\delta}(T;y)=0\subset\mathbb A^2_{\FF_{31}}\) 上，点 \((T,y)=(0,12)\) 的几何解析类型为 \(A_3\)（触节点）。
\end{enumerate}
\end{theorem}
\begin{proof}
第一部分与系数消零判别相同：\(F_\delta(T;y_0)\equiv T^4\) 等价于
\[
8y_0-3\equiv 0,\qquad
2y_0^2-34y_0-4\equiv 0,\qquad
y_0(y_0-1)P_{\mathrm{LY}}(y_0)\equiv 0\ (\bmod p).
\]
前两式给出 \(y_0\equiv 3/8\) 且 \(p\in\{17,31\}\)。
再用
\[
P_{\mathrm{LY}}\!\left(\frac{3}{8}\right)=\frac{10571}{64}=\frac{11\cdot 31^2}{2^6}
\]
可知仅 \(p=31\) 成立；由 \(8^{-1}\equiv 4\pmod{31}\) 得 \(y_0\equiv 12\)。

第二部分取局部坐标 \(t:=T,\ s:=y-12\)。在 \(\FF_{31}\) 上展开：
\[
\overline{F_\delta}(t;12+s)
=-12s^2+14t^2s+t^4+8t^3s+2t^2s^2+10s^3-15s^4-8s^5.
\]
令 \(s_1:=s+2t^2\)（可逆坐标变换），得到
\[
\overline{F_\delta}
=-12s_1^2-13t^4+\text{(加权次数 }\ge 5,\ \text{其中 }w(t)=1,w(s_1)=2\text{)}.
\]
故其加权初始主部为非退化型 \(u^2-\lambda v^4\)（\(\lambda\neq 0\)），
即与标准形 \(u^2=v^4\) 同解析类型，因而为 \(A_3\) 触节点。证毕。
\end{proof}

\begin{theorem}[\texorpdfstring{$\operatorname{Res}_y(Q_5,A_7)$}{Res\_y(Q5,A7)} 分解与多重碰撞障碍素数]\label{thm:xi-q5-a7-resultant-double-collision-primes}
关于 \(y\) 的结果式满足
\[
\operatorname{Res}_{y}\!\bigl(Q_5(y),A_7(y)\bigr)
=-2^{47}\cdot 3^{16}\cdot 5\cdot 13\cdot 31^{8}\cdot 73.
\]
因此对任意素数 \(p\nmid 2\cdot 3\cdot 5\cdot 13\cdot 31\cdot 73\)，有
\[
\gcd(\overline{Q_5},\overline{A_7})=1
\quad\text{于}\quad \FF_p[y].
\]
于是任意 \(\bar y\) 满足 \(\overline{Q_5}(\bar y)=0\) 时必有 \(\overline{A_7}(\bar y)\neq 0\)，
从而 \(\overline{F_\delta}(T;\bar y)\) 不能出现“两处不同二重根”。
另一方面：
\[
\overline{F_\delta}(T;0)=T^2(T+1)^2\ \text{于}\ \FF_5,\qquad
\overline{F_\delta}(T;1)=T^2(T-4)^2\ \text{于}\ \FF_{13},
\]
\[
\overline{F_\delta}(T;-10)=T^2(T-5)^2\ \text{于}\ \FF_{73},
\]
而在 \(p=31,\ \bar y=12\) 处有
\[
\overline{F_\delta}(T;12)=T^4,
\]
其几何类型由定理 \ref{thm:xi-fdelta-char31-total-collapse} 给出为 \(A_3\)。
\end{theorem}
\begin{proof}
结果式分解与互素判别由 Sylvester 理论直接给出。
对 \(p\nmid 2\cdot 3\cdot 5\cdot 13\cdot 31\cdot 73\)，\(\overline{Q_5}\) 与 \(\overline{A_7}\) 互素，
配合定理 \ref{thm:xi-fdelta-subresultant-unique-double-root} 的一阶子结式公式，
可知在 \(\overline{Q_5}(\bar y)=0\) 处不存在“二处不同二重根”情形。
对 \(p=5,13,73,31\) 的显式分解由逐点代入 \(y=0,1,-10,12\) 直接核验。证毕。
\end{proof}

\begin{theorem}[主分歧三次与二级退化五次的线性互不相交与联合密度]\label{thm:xi-leyang-ply-q5-linearly-disjoint-chebotarev}
令 \(L_3\) 为 \(P_{\mathrm{LY}}(y)\) 的分裂域，\(L_5\) 为 \(Q_5(y)\) 的分裂域。
则：
\begin{enumerate}
  \item \(L_3\cap L_5=\QQ\)，从而 \(\operatorname{Gal}(L_3L_5/\QQ)\cong S_3\times S_5\)；
  \item 对一切不分歧素数 \(p\nmid \mathrm{Disc}(P_{\mathrm{LY}})\mathrm{Disc}(Q_5)\)，两者模 \(p\) 的分解型联合分布为直积分布；
  \item 特别地，
  \[
  \delta\bigl(P_{\mathrm{LY}}\ \text{全分裂且}\ Q_5\ \text{全分裂}\bigr)=\frac{1}{6}\cdot\frac{1}{120}=\frac{1}{720},
  \]
  \[
  \delta\bigl(P_{\mathrm{LY}}\ \text{在}\ \FF_p\ \text{上有根且}\ Q_5\ \text{在}\ \FF_p\ \text{上有根}\bigr)
  =\frac{2}{3}\cdot\frac{19}{30}=\frac{19}{45}.
  \]
\end{enumerate}
\end{theorem}
\begin{proof}
由结论 \ref{con:xi-terminal-37-squarefree-coupling}，
\[
\mathrm{Disc}(P_{\mathrm{LY}})=-3^9\cdot 31^2\cdot 37,
\]
故 \(L_3\) 的唯一二次子域为 \(\QQ(\sqrt{-111})\)。
由结论 \ref{con:xi-terminal-zm-leyang-q5-discriminant-727}，
\[
\mathrm{Disc}(Q_5)=-2^{28}\cdot 3^{9}\cdot 11^{2}\cdot 31^{3}\cdot 727,
\]
故 \(L_5\) 的唯一二次子域为 \(\QQ(\sqrt{-67611})\)。
两二次子域分歧素集不同，因而不相等；而 \(L_3/\QQ\) 的真非平凡子域仅有该二次子域，故 \(L_3\cap L_5=\QQ\)。
直积伽罗瓦群同构随之得到。

联合分布与密度由 Chebotarev 定理在直积群上的等分布给出。
其中 \(P_{\mathrm{LY}}\) 的三类分解型密度由结论 \ref{con:xi-terminal-zm-elliptic-normalization} 给出为
\((1)(1)(1)\) 密度 \(1/6\)、\((1)(2)\) 密度 \(1/2\)、\((3)\) 密度 \(1/3\)，故“有根”密度为 \(2/3\)。
对 \(S_5\) 的共轭类计数给出 \(Q_5\) 的分解型密度：全分裂对应恒等类密度 \(1/120\)；
“无一次因子”等价于无不动点，仅可能为 \(5\)-循环（密度 \(1/5\)）或 \((3)(2)\)（密度 \(1/6\)），故“有一次因子”密度为 \(1-(1/5+1/6)=19/30\)。
乘积即得所示密度。证毕。
\end{proof}

\begin{conclusion}[\texorpdfstring{$\delta$}{delta}-分支域与二级退化五次分裂域的线性互不相交]\label{con:xi-delta-b5-q5-linearly-disjoint-s5xs5}
令 \(L_B\) 为 \(B_5(\delta)\) 的分裂域（命题 \ref{prop:xi-terminal-zm-delta-discy-factorization-a5b5}），令 \(L_5\) 为 \(Q_5(y)\) 的分裂域（定理 \ref{thm:xi-q5-galois-s5}）。
则
\[
L_B\cap L_5=\QQ,
\qquad
\Gal(L_BL_5/\QQ)\cong S_5\times S_5.
\]
\end{conclusion}
\begin{proof}
由定理 \ref{thm:xi-terminal-zm-delta-a5-b5-galois-s5} 与定理 \ref{thm:xi-q5-galois-s5}，两者均为 \(S_5\)-伽罗瓦扩张。
因此任一交域 \(L_B\cap L_5\) 必为两者的正规子扩张；而 \(S_5\) 的非平凡真正规子群仅有 \(A_5\)，故交域只能是 \(\QQ\) 或各自唯一的二次分辨子域。
\par
由结论 \ref{con:xi-terminal-zm-delta-b5-discriminant-ramification-quadratic}，\(L_B\) 的二次分辨子域为 \(\QQ(\sqrt{-140267})\)；
由定理 \ref{thm:xi-leyang-ply-q5-linearly-disjoint-chebotarev} 的证明，\(L_5\) 的唯一二次子域为 \(\QQ(\sqrt{-67611})\)。
两二次域分歧素集不同，因而不相等，遂 \(L_B\cap L_5=\QQ\)。
线性互不相交推出直积伽罗瓦群同构。证毕。
\end{proof}

\begin{conclusion}[四阶核判别式五次域与 \texorpdfstring{$Q_5$}{Q5} 分裂域的线性互不相交]\label{con:xi-p7-q5-linearly-disjoint-s5xs5}
令 \(L_{p_7}\) 为五次多项式
\[
p_7(x)=x^5-2x^4-7x^3-2x+2
\]
的分裂域（命题 \ref{prop:pom-a4-galois-s5}），令 \(L_5\) 为 \(Q_5(y)\) 的分裂域（定理 \ref{thm:xi-q5-galois-s5}）。
则
\[
L_{p_7}\cap L_5=\QQ,
\qquad
\Gal(L_{p_7}L_5/\QQ)\cong S_5\times S_5.
\]
\end{conclusion}
\begin{proof}
由命题 \ref{prop:pom-a4-galois-s5} 与定理 \ref{thm:xi-q5-galois-s5}，两者均为 \(S_5\)-伽罗瓦扩张。
因此任一交域 \(L_{p_7}\cap L_5\) 必为两者的正规子扩张；而 \(S_5\) 的非平凡真正规子群仅有 \(A_5\)，故交域只能是 \(\QQ\) 或各自唯一的二次分辨子域。
\par
由命题 \ref{prop:pom-a4-galois-s5}，\(L_{p_7}\) 的二次分辨子域为
\(
\QQ\!\left(\sqrt{\Disc_x(p_7)}\right)=\QQ(\sqrt{-985219})
\)；
由结论 \ref{con:xi-terminal-zm-delta-a5-discriminant-quadratic}，\(L_5\) 的唯一二次子域为 \(\QQ(\sqrt{-67611})\)。
两二次域不相等，遂 \(L_{p_7}\cap L_5=\QQ\)。
线性互不相交推出直积伽罗瓦群同构。证毕。
\end{proof}

\paragraph{\texorpdfstring{$p_7$}{p7} 的 \texorpdfstring{$S_5$}{S5} 分裂域：三类典型中间域的 \texorpdfstring{$\zeta$}{zeta}--Artin 分解与 Chebotarev 全剖面}
令 \(L:=L_{p_7}\) 为 \(p_7\) 的分裂域（命题 \ref{prop:pom-a4-galois-s5}），并固定同构 \(\Gal(L/\QQ)\cong S_5\)。
令
\[
K_{5}:=L^{S_4},\qquad
K_{10}:=L^{S_2\times S_3},\qquad
K_{20}:=L^{S_3},
\]
其中 \(S_4\) 为一点稳定子，\(S_2\times S_3\) 为二子集稳定子，\(S_3\) 为有序对稳定子。
记 \(S_5\) 的三个有理不可约表示（Specht 模）为
\[
\rho_{4}:=S^{(4,1)},\qquad
\rho_{5}:=S^{(3,2)},\qquad
\rho_{6}:=S^{(3,1,1)},
\]
并以 \(L(s,\rho_i)\) 表示相应 Artin \(L\)-函数。

\begin{theorem}[三类中间域的 Dedekind \texorpdfstring{$\zeta$}{zeta} 的 Artin--Dedekind 分解]\label{thm:xi-p7-s5-minimal-artin-dedekind-zeta-factorization}
\leanverified{paper\_xi\_p7\_s5\_minimal\_artin\_dedekind\_zeta\_factorization}
在上述记号下成立完全显式的分解恒等式
\[
\zeta_{K_{5}}(s)=\zeta(s)\,L(s,\rho_4),
\qquad
\zeta_{K_{10}}(s)=\zeta(s)\,L(s,\rho_4)\,L(s,\rho_5),
\]
\[
\zeta_{K_{20}}(s)=\zeta(s)\,L(s,\rho_4)^{2}\,L(s,\rho_5)\,L(s,\rho_6).
\]
从而得到不可约分量的最短比值表达
\[
L(s,\rho_5)=\frac{\zeta_{K_{10}}(s)}{\zeta_{K_{5}}(s)},
\qquad
L(s,\rho_6)=\frac{\zeta(s)\,\zeta_{K_{20}}(s)}{\zeta_{K_{10}}(s)\,\zeta_{K_{5}}(s)}.
\]
\end{theorem}
\begin{proof}
对任意中间域 \(K=L^H\)（\(H\le S_5\)），Artin 形式给出
\[
\zeta_K(s)=L\!\left(s,\Ind_{H}^{S_5}\mathbf 1\right),
\]
其中 \(\Ind_{H}^{S_5}\mathbf 1\) 为 \(S_5\) 在左陪集集 \(S_5/H\) 上的置换表示。
因此只需对三类 \(H\) 给出置换表示的 Specht 分解。
\par
当 \(H=S_4\) 时，\(S_5/H\) 可识别为 \(5\) 点集合，置换表示分解为
\(
\Ind_{S_4}^{S_5}\mathbf 1\cong \mathbf 1\oplus \rho_4
\)。
当 \(H=S_2\times S_3\) 时，\(S_5/H\) 可识别为 \(10\) 个二子集，置换表示分解为
\(
\Ind_{S_2\times S_3}^{S_5}\mathbf 1\cong \mathbf 1\oplus \rho_4\oplus \rho_5
\)。
\par
当 \(H=S_3\) 时，\(S_5/H\) 可识别为 \(20\) 个有序对。
将有序对集合按“无向边”与“定向翻转”分解，可得
\[
\Ind_{S_3}^{S_5}\mathbf 1
\cong
\Ind_{S_2\times S_3}^{S_5}\mathbf 1\ \oplus\ \Ind_{S_2\times S_3}^{S_5}\mathrm{sgn}_{S_2},
\]
其中 \(\mathrm{sgn}_{S_2}\) 为二点交换的符号表示。
后一项同构于 \(\bigwedge^2(\mathbf 1\oplus \rho_4)\cong \rho_4\oplus \rho_6\)，从而得到
\[
\Ind_{S_3}^{S_5}\mathbf 1\cong \mathbf 1\oplus 2\rho_4\oplus \rho_5\oplus \rho_6.
\]
将以上分解代入 Artin \(L\)-函数的乘法性即得所示三条 Dedekind \(\zeta\) 分解与两条比值恒等式。
可复现核验脚本：\texttt{scripts/exp\_xi\_p7\_s5\_artin\_dedekind\_chebotarev\_certificate.py}。证毕。
\end{proof}

\begin{theorem}[三类中间域的 Chebotarev 分裂型全剖面]\label{thm:xi-p7-s5-k5-k10-k20-splitting-type-table}
\leanverified{paper\_xi\_p7\_s5\_k5\_k10\_k20\_splitting\_type\_table}
令 \(p\nmid 2\Disc_x(p_7)\) 为不分歧素数，并令 \(\mathrm{Frob}_p\subset S_5\) 为其 Frobenius 共轭类。
则 \(p\) 在 \(K_5,K_{10},K_{20}\) 中的分解型完全由 \(\mathrm{Frob}_p\) 的循环型决定，且对应自然密度为该共轭类在 \(S_5\) 中所占比例：
\begin{enumerate}
  \item \((1^5)\)（密度 \(1/120\)）：
  \[
  K_5:1^5;\qquad K_{10}:1^{10};\qquad K_{20}:1^{20}.
  \]
  \item \((2\cdot 1^3)\)（密度 \(1/12\)）：
  \[
  K_5:2\cdot 1^3;\qquad K_{10}:2^3\cdot 1^4;\qquad K_{20}:2^7\cdot 1^6.
  \]
  \item \((2^2\cdot 1)\)（密度 \(1/8\)）：
  \[
  K_5:2^2\cdot 1;\qquad K_{10}:2^4\cdot 1^2;\qquad K_{20}:2^{10}.
  \]
  \item \((3\cdot 1^2)\)（密度 \(1/6\)）：
  \[
  K_5:3\cdot 1^2;\qquad K_{10}:3^3\cdot 1;\qquad K_{20}:3^6\cdot 1^2.
  \]
  \item \((3\cdot 2)\)（密度 \(1/6\)）：
  \[
  K_5:3\cdot 2;\qquad K_{10}:6\cdot 3\cdot 1;\qquad K_{20}:6^2\cdot 3^2\cdot 2.
  \]
  \item \((4\cdot 1)\)（密度 \(1/4\)）：
  \[
  K_5:4\cdot 1;\qquad K_{10}:4^2\cdot 2;\qquad K_{20}:4^5.
  \]
  \item \((5)\)（密度 \(1/5\)）：
  \[
  K_5:5;\qquad K_{10}:5^2;\qquad K_{20}:5^4.
  \]
\end{enumerate}
\end{theorem}
\begin{proof}
对 \(p\nmid \Disc(L)\)，\(p\) 在 \(K=L^H\) 中的素分解剩余度序列等于 \(\mathrm{Frob}_p\) 在 \(S_5/H\) 上的循环长度多重集。
分别取 \(H=S_4\)（\(5\) 点）、\(H=S_2\times S_3\)（\(10\) 个二子集）、\(H=S_3\)（\(20\) 个有序对），对 \(S_5\) 的七个循环型代表元计算其在三组集合上的循环分解，即得所列分解型。
密度由 Chebotarev 定理给出，等于共轭类大小除以 \(|S_5|=120\)。
可复现核验脚本同定理 \ref{thm:xi-p7-s5-minimal-artin-dedekind-zeta-factorization}。证毕。
\end{proof}

\begin{theorem}[二次塔 \texorpdfstring{$K_{20}/K_{10}$}{K20/K10} 的定向翻转判别与惰性统计]\label{thm:xi-p7-s5-k20-over-k10-oriented-flip-inertia-profile}
在上述设定下，扩张 \(K_{20}/K_{10}\) 为二次扩张。
令 \(p\nmid 2\Disc_x(p_7)\)，取 \(\sigma\in \mathrm{Frob}_p\subset S_5\)，并令 \(\mathfrak P\) 为 \(K_{10}\) 中位于 \(p\) 上方的一素理想。
把 \(\mathfrak P\) 视为 \(\sigma\) 在二子集集合上的某一轨道，并记该轨道长度为 \(m\)。
则在二次扩张 \(K_{20}/K_{10}\) 中：
\begin{itemize}
  \item \(\mathfrak P\) 分裂当且仅当 \(\sigma^{m}\) 在该二子集上点态固定；
  \item \(\mathfrak P\) 惰性当且仅当 \(\sigma^{m}\) 交换该二子集的两端点。
\end{itemize}
并且按 \(\mathrm{Frob}_p\) 的循环型，惰性素理想的“数目—剩余度”剖面精确为：
\begin{itemize}
  \item \((1^5)\)：无惰性；
  \item \((2\cdot 1^3)\)：恰有 \(1\) 个剩余度 \(1\) 的素理想惰性；
  \item \((2^2\cdot 1)\)：恰有 \(2\) 个剩余度 \(1\) 的素理想惰性；
  \item \((3\cdot 1^2)\)：无惰性；
  \item \((3\cdot 2)\)：恰有 \(1\) 个剩余度 \(1\) 的素理想惰性；
  \item \((4\cdot 1)\)：恰有 \(1\) 个剩余度 \(2\) 的素理想惰性；
  \item \((5)\)：无惰性。
\end{itemize}
\end{theorem}
\begin{proof}
由稳定子链 \(S_3\le S_2\times S_3\le S_5\) 且 \([S_2\times S_3:S_3]=2\)，可知 \(K_{20}/K_{10}\) 为二次扩张。
集合论上，\(K_{20}\) 对应有序对而 \(K_{10}\) 对应无序对；有序对到无序对为二重覆盖，其 deck 变换为“翻转方向”。
\par
对 \(\sigma\) 在无序对轨道（长度 \(m\)）上的作用，提升到有序对后：若绕轨道一圈返回时方向保持，则得到两条长度 \(m\) 的定向循环（对应分裂）；
若返回时方向翻转，则得到一条长度 \(2m\) 的定向循环（对应惰性）。
这正等价于判别式中 \(\sigma^m\) 是否交换该二子集两端点。
将该判别应用于定理 \ref{thm:xi-p7-s5-k5-k10-k20-splitting-type-table} 的循环剖面逐类核验，即得惰性统计。
可复现核验脚本同定理 \ref{thm:xi-p7-s5-minimal-artin-dedekind-zeta-factorization}。证毕。
\end{proof}
\input{para__xi-p7-dyadic-ramification-and-signatures}

\begin{theorem}[唯一奇分歧素数处的判别式指数与 Artin 导子]\label{thm:xi-p7-s5-odd-ramification-discriminant-and-conductor}
\leanverified{paper\_xi\_p7\_s5\_odd\_ramification\_discriminant\_and\_conductor}
记
\[
q:=985219.
\]
则 \(L/\QQ\) 的奇素分歧集合恰为 \(\{q\}\)，且 \(q\) 在 \(L/\QQ\) 为温和分歧，惯性群由一换位 \(\tau\in S_5\) 生成。
对三类中间域的判别式奇部分指数满足
\[
\nu_q\!\bigl(\Disc(K_{5})\bigr)=1,\qquad
\nu_q\!\bigl(\Disc(K_{10})\bigr)=3,\qquad
\nu_q\!\bigl(\Disc(K_{20})\bigr)=7.
\]
并且对 \(\rho_4,\rho_5,\rho_6\) 的 Artin 导子指数满足
\[
a_q(\rho_4)=1,\qquad a_q(\rho_5)=2,\qquad a_q(\rho_6)=3.
\]
\end{theorem}
\begin{proof}
由命题 \ref{prop:pom-a4-galois-s5}，\(\Disc_x(p_7)=-2^4\cdot 985219=-2^4\cdot q\)，其平方自由部分为 \(-q\)，故 \(L/\QQ\) 的分歧素数仅可能在 \(\{2,q\}\) 内，从而奇分歧仅为 \(q\)。
对 \(q\) 的模分解出现恰一重根且其余根单纯，故局部分歧指数为 \(2\)，从而为温和分歧，惯性群同构于 \(C_2\) 并可取生成元为换位 \(\tau\)。
\par
温和情形下，对任意 \(K=L^H\) 的判别式指数满足
\[
\nu_q(\Disc(K))=\dim\Ind_H^{S_5}\mathbf 1-\dim(\Ind_H^{S_5}\mathbf 1)^{\langle\tau\rangle}
=\#(S_5/H)-\#\bigl(\tau\text{ 在 }S_5/H\text{ 上的轨道数}\bigr).
\]
对 \(H=S_4\)（\(5\) 点作用）换位轨道数为 \(4\)，故得 \(5-4=1\)；
对 \(H=S_2\times S_3\)（二子集作用）轨道数为 \(7\)，故得 \(10-7=3\)；
对 \(H=S_3\)（有序对作用）轨道数为 \(13\)，故得 \(20-13=7\)。
\par
对 \(C_2=\langle\tau\rangle\) 的温和导子公式为
\[
a_q(\rho)=\dim\rho-\dim\rho^{\langle\tau\rangle}=\frac{\dim\rho-\chi_\rho(\tau)}{2}.
\]
由置换表示分解 \(\Ind_{S_4}^{S_5}\mathbf 1\cong \mathbf 1\oplus\rho_4\) 可得
\(
\chi_{\rho_4}(\tau)=\chi_{\Ind_{S_4}^{S_5}\mathbf 1}(\tau)-1
\)；
而换位在 \(5\) 点上固定 \(3\) 点，故 \(\chi_{\rho_4}(\tau)=2\)，从而 \(a_q(\rho_4)=(4-2)/2=1\)。
同理由 \(\Ind_{S_2\times S_3}^{S_5}\mathbf 1\cong \mathbf 1\oplus\rho_4\oplus\rho_5\) 及换位在二子集上固定 \(4\) 个二子集，得 \(\chi_{\rho_5}(\tau)=4-1-2=1\)，从而 \(a_q(\rho_5)=(5-1)/2=2\)。
再由 \(\Ind_{S_3}^{S_5}\mathbf 1\cong \mathbf 1\oplus 2\rho_4\oplus\rho_5\oplus\rho_6\) 及换位在有序对上固定 \(6\) 个有序对，得
\(\chi_{\rho_6}(\tau)=6-1-2\cdot 2-1=0\)，从而 \(a_q(\rho_6)=(6-0)/2=3\)。
可复现核验脚本同定理 \ref{thm:xi-p7-s5-minimal-artin-dedekind-zeta-factorization}。证毕。
\end{proof}
\input{para__xi-p7-conductor-discriminant-rayclass}

\begin{theorem}[分歧素数处局部 Euler 因子的三重同形]\label{thm:xi-p7-s5-local-euler-factor-triple-isomorphism-at-q}
在定理 \ref{thm:xi-p7-s5-odd-ramification-discriminant-and-conductor} 的设定下，
在 \(\FF_q\) 上 \(p_7\) 完全分裂且仅出现一个二重根。
因此 \(q\) 在分裂域 \(L\) 的任一素处的剩余度皆为 \(1\)，从而分解群与惯性群相等并满足
\[
D_q(L/\QQ)=I_q(L/\QQ)\cong C_2.
\]
于是对任意 Artin 表示 \(\rho\)，其在 \(q\) 处的局部 Euler 因子由惯性不变子空间维数完全决定：
\[
L_q(s,\rho)=\bigl(1-q^{-s}\bigr)^{-\dim \rho^{I_q}}.
\]
特别地，
\[
\dim \rho_4^{I_q}=\dim \rho_5^{I_q}=\dim \rho_6^{I_q}=3,
\]
从而三条互不同构的不可约表示在同一分歧素处给出同形局部因子
\[
L_q(s,\rho_4)=L_q(s,\rho_5)=L_q(s,\rho_6)=\bigl(1-q^{-s}\bigr)^{-3}.
\]
相应地，三条 Dedekind \(\zeta\) 的 \(q\)-局部因子为
\[
\zeta_{K_5,q}(s)=\bigl(1-q^{-s}\bigr)^{-4},\qquad
\zeta_{K_{10},q}(s)=\bigl(1-q^{-s}\bigr)^{-7},\qquad
\zeta_{K_{20},q}(s)=\bigl(1-q^{-s}\bigr)^{-13}.
\]
\end{theorem}
\begin{proof}
由于 \(p_7\bmod q\) 在 \(\FF_q\) 上完全分裂且仅一处二重根，分裂域的剩余域不扩张，故任一 \(q\)-上素处的剩余度为 \(1\)，从而分解群等于惯性群。
当 \(D_q=I_q\) 时局部 Frobenius 在 \(\rho^{I_q}\) 上恒等，故局部因子为 \((1-q^{-s})^{-\dim\rho^{I_q}}\)。
\par
由定理 \ref{thm:xi-p7-s5-odd-ramification-discriminant-and-conductor} 的特征标计算，
对 \(C_2=\langle\tau\rangle\) 有
\[
\dim\rho^{I_q}=\dim\rho^{\langle\tau\rangle}=\frac{\dim\rho+\chi_\rho(\tau)}{2},
\]
从而
\(
\dim\rho_4^{I_q}=(4+2)/2=3
\)，
\(
\dim\rho_5^{I_q}=(5+1)/2=3
\)，
\(
\dim\rho_6^{I_q}=(6+0)/2=3
\)。
最后对 \(K=L^H\) 的 Dedekind \(\zeta\) 由定理 \ref{thm:xi-p7-s5-minimal-artin-dedekind-zeta-factorization} 的 Artin 分解取局部因子并相乘，得到所列指数 \(4,7,13\)（亦等于定理 \ref{thm:xi-p7-s5-odd-ramification-discriminant-and-conductor} 中的 \(\tau\)-轨道数）。证毕。
\end{proof}

\begin{conclusion}[三条 \texorpdfstring{$S_3$}{S3} 三次分裂域与四阶核五次 \texorpdfstring{$S_5$}{S5} 分裂域的线性互不相交及完全分裂密度]\label{con:xi-p7-three-s3-linearly-disjoint-chebotarev}
令
\[
P_{37}(y):=y^3-37y^2+3y+25\in\ZZ[y],\qquad
Q_{\mathrm{LY}}(y):=256y^3+195y^2+93y-48\in\ZZ[y],
\]
以及
\[
f_T(t):=48114t^3+7263t^2-506t-55\in\ZZ[t].
\]
记其分裂域分别为 \(L_{37},L_{Q},L_{T}\)，并记 \(L_{p_7}\) 为 \(p_7\) 的分裂域（命题 \ref{prop:pom-a4-galois-s5}）。
则：
\begin{enumerate}
  \item 三个三次多项式均在 \(\QQ\) 上不可约，且其判别式分别为
  \[
  \Disc(P_{37})=2^8\cdot 23^2\cdot 37,\qquad
  \Disc(Q_{\mathrm{LY}})=-3^3\cdot 37\cdot 2677^2,\qquad
  \Disc(f_T)=2^{2}\cdot 3^{6}\cdot 5^{2}\cdot 11\cdot 23\cdot 5894101,
  \]
  其中 \(\Disc(Q_{\mathrm{LY}})\) 见结论 \ref{con:xi-terminal-zm-leyang-badprime-37-common-double-root}，
  \(\Disc(f_T)\) 见式 \eqref{eq:window6_edge_flux_skeleton_arithmetic}。
  因而
  \[
  \Gal(L_{37}/\QQ)\cong \Gal(L_{Q}/\QQ)\cong \Gal(L_{T}/\QQ)\cong S_3.
  \]
  \item \(L_{37},L_{Q},L_{T},L_{p_7}\) 两两线性互不相交；因此其复合域
  \[
  L:=L_{p_7}L_{37}L_{Q}L_{T}
  \]
  满足
  \[
  \Gal(L/\QQ)\cong S_5\times S_3\times S_3\times S_3.
  \]
  \item 由 Chebotarev 密度定理，满足 \(p\) 在四个扩张中均完全分解的素数集合具有狄利克雷密度
  \[
  \frac{1}{|S_5|\cdot |S_3|^3}=\frac{1}{120\cdot 6^3}=\frac{1}{25920}.
  \]
  更一般地，对任意指定的四重共轭类
  \(
  (C_{p_7},C_{37},C_Q,C_T)\subset S_5\times S_3^3
  \)
  均存在无穷多素数 \(p\) 使相应 Frobenius 元落入该共轭类，且其密度等于
  \(
  |C_{p_7}||C_{37}||C_Q||C_T|/(120\cdot 6^3)
  \)。
\end{enumerate}
\end{conclusion}
\begin{proof}
三次不可约性可由有理根判别与一次良性模素约化给出。
三次不可约多项式的伽罗瓦群为 \(A_3\) 当且仅当判别式为平方；由上式判别式分解可知三者判别式均非平方，故其伽罗瓦群皆为 \(S_3\)。
\par
为证线性互不相交，任取两者之交 \(E:=L_i\cap L_j\)。
因 \(L_i/\QQ\) 与 \(L_j/\QQ\) 皆为 Galois 扩张，\(E/\QQ\) 亦 Galois。
对 \(S_3\)-扩张，其非平凡真 Galois 子扩张唯有唯一二次子域；
对 \(S_5\)-扩张亦同理仅有唯一二次分辨子域（对应符号表示）。
因此任意交域若非 \(\QQ\)，必为二次域。
\par
由判别式平方自由部分可识别相应二次子域：
\[
L_{37}\text{之二次子域 }=\QQ(\sqrt{37}),\quad
L_{Q}\text{之二次子域 }=\QQ(\sqrt{-3\cdot 37}),\quad
L_{T}\text{之二次子域 }=\QQ(\sqrt{11\cdot 23\cdot 5894101}),
\]
而 \(L_{p_7}\) 之二次分辨子域为
\(
\QQ(\sqrt{\Disc_x(p_7)})=\QQ(\sqrt{-985219})
\)
（命题 \ref{prop:pom-a4-galois-s5}）。
上述四个二次域两两不同，故任意两域之交只能为 \(\QQ\)，从而两两线性互不相交，复合域伽罗瓦群为直积。
\par
最后，直积伽罗瓦群下的 Chebotarev 密度为共轭类所占比例；完全分解对应恒等共轭类，密度即为群阶倒数。
可复现核验脚本：\texttt{scripts/exp\_xi\_p7\_three\_s3\_disjointness\_chebotarev\_certificate.py}。证毕。
\end{proof}

\begin{conclusion}[两类 \texorpdfstring{$S_5$}{S5} 分解型的直积分布与显式密度]\label{con:xi-delta-b5-q5-chebotarev-product-densities}
沿用结论 \ref{con:xi-delta-b5-q5-linearly-disjoint-s5xs5} 的记号。
对任意不分歧素数
\[
p\nmid \Disc(B_5)\Disc(Q_5),
\]
设 \(\sigma_p\in S_5\times S_5\) 为 \(L_BL_5/\QQ\) 上的 Frobenius 共轭类。
则 \(\sigma_p\) 在直积群中按共轭类等分布，从而 \(B_5\bmod p\) 与 \(Q_5\bmod p\) 的分解型联合分布为乘积分布。
特别地，
\[
\delta\bigl(B_5\ \text{在}\ \FF_p\ \text{上不可约且}\ Q_5\ \text{在}\ \FF_p\ \text{上不可约}\bigr)
=\left(\frac{24}{120}\right)^2=\frac{1}{25},
\]
\[
\delta\bigl(B_5\ \text{全分裂且}\ Q_5\ \text{全分裂}\bigr)
=\left(\frac{1}{120}\right)^2=\frac{1}{14400},
\]
\[
\delta\bigl(B_5\ \text{分解型 }(2)(1)(1)(1)\ \text{且}\ Q_5\ \text{不可约}\bigr)
=\frac{10}{120}\cdot\frac{24}{120}=\frac{1}{60}.
\]
\end{conclusion}
\begin{proof}
由结论 \ref{con:xi-delta-b5-q5-linearly-disjoint-s5xs5}，\(\Gal(L_BL_5/\QQ)\cong S_5\times S_5\)。
对不分歧素数 \(p\)，Chebotarev 定理给出 Frobenius 共轭类在直积群上的等分布；
并且五次多项式在 \(\FF_p\) 的分解型与 \(S_5\) 中 Frobenius 的循环型一一对应。
将 \(S_5\) 各共轭类大小（恒等类 \(1\)、换位 \(10\)、\(5\)-循环 \(24\)）代入即可得到上述密度。证毕。
\end{proof}

\begin{conclusion}[双五次分解型的乘积分布与三组基本密度]\label{con:xi-p7-q5-chebotarev-product-densities}
沿用结论 \ref{con:xi-p7-q5-linearly-disjoint-s5xs5} 的记号。
对任意不分歧素数
\[
p\nmid \Disc(p_7)\Disc(Q_5),
\]
设 \(\sigma_p\in S_5\times S_5\) 为 \(L_{p_7}L_5/\QQ\) 上的 Frobenius 共轭类。
则 \(\sigma_p\) 在直积群中按共轭类等分布，从而 \(p_7\bmod p\) 与 \(Q_5\bmod p\) 的分解型联合分布为乘积分布。
特别地，令“不可约模 \(p\)”对应 \(S_5\) 中的 \(5\)-循环共轭类（大小 \(24\)），令“全分裂模 \(p\)”对应恒等共轭类（大小 \(1\)），则
\[
\delta\bigl(p_7\ \text{在}\ \FF_p\ \text{上不可约且}\ Q_5\ \text{在}\ \FF_p\ \text{上不可约}\bigr)
=\left(\frac{24}{120}\right)^2=\frac{1}{25},
\]
\[
\delta\bigl(p_7\ \text{不可约且}\ Q_5\ \text{全分裂}\bigr)
=\frac{24}{120}\cdot\frac{1}{120}=\frac{1}{600},
\]
\[
\delta\bigl(p_7\ \text{全分裂且}\ Q_5\ \text{全分裂}\bigr)
=\left(\frac{1}{120}\right)^2=\frac{1}{14400}.
\]
\end{conclusion}
\begin{proof}
由结论 \ref{con:xi-p7-q5-linearly-disjoint-s5xs5}，\(\Gal(L_{p_7}L_5/\QQ)\cong S_5\times S_5\)。
对不分歧素数 \(p\)，Chebotarev 定理给出 Frobenius 共轭类在直积群上的等分布；
并且五次多项式在 \(\FF_p\) 的分解型与 \(S_5\) 中 Frobenius 的循环型一一对应。
将 \(S_5\) 各共轭类大小代入即可得到所述密度。证毕。
\end{proof}

\begin{conclusion}[三条 \texorpdfstring{$S_5$}{S5} 分裂域的三重直积与极大阿贝尔子扩张]\label{con:xi-p7-q5-b5-triple-product-s5cube-abelianization}
令 \(L_{p_7}\) 为 \(p_7\) 的分裂域（命题 \ref{prop:pom-a4-galois-s5}），令 \(L_5\) 为 \(Q_5\) 的分裂域（定理 \ref{thm:xi-q5-galois-s5}），令 \(L_B\) 为 \(B_5\) 的分裂域（定理 \ref{thm:xi-terminal-zm-delta-a5-b5-galois-s5}）。
则
\[
L_{p_7}\cap L_5\cap L_B=\QQ,
\qquad
\Gal(L_{p_7}L_5L_B/\QQ)\cong S_5\times S_5\times S_5,
\]
从而
\[
[L_{p_7}L_5L_B:\QQ]=120^3.
\]
并且其极大阿贝尔子扩张为三二次合成域
\[
\QQ(\sqrt{-985219},\ \sqrt{-67611},\ \sqrt{-140267}),
\]
其伽罗瓦群同构于 \((\ZZ/2\ZZ)^3\)。
\end{conclusion}
\begin{proof}
结论 \ref{con:xi-p7-q5-linearly-disjoint-s5xs5} 与结论 \ref{con:xi-delta-b5-q5-linearly-disjoint-s5xs5} 分别给出
\(
L_{p_7}\cap L_5=\QQ
\)
与
\(
L_B\cap L_5=\QQ
\)。
同理，由命题 \ref{prop:pom-a4-galois-s5} 与结论 \ref{con:xi-terminal-zm-delta-b5-discriminant-ramification-quadratic} 可读出 \(L_{p_7}\)、\(L_B\) 的唯一二次分辨子域分别为 \(\QQ(\sqrt{-985219})\)、\(\QQ(\sqrt{-140267})\)，两者不相等，故 \(L_{p_7}\cap L_B=\QQ\)。
\par
由此 \(\Gal(L_{p_7}L_5/\QQ)\cong S_5\times S_5\)。
其二次子域对应于三个非平凡符号表征（两坐标符号与其乘积），故仅可能为
\[
\QQ(\sqrt{-985219}),\qquad
\QQ(\sqrt{-67611}),\qquad
\QQ\!\left(\sqrt{\mathrm{sf}\bigl((-985219)(-67611)\bigr)}\right),
\]
其中 \(\mathrm{sf}(\cdot)\) 表示平方自由核。
而 \(L_B\) 的唯一二次子域为 \(\QQ(\sqrt{-140267})\)（结论 \ref{con:xi-terminal-zm-delta-b5-discriminant-ramification-quadratic}），不在上述三者之列，
故 \((L_{p_7}L_5)\cap L_B=\QQ\)。
于是
\(
[L_{p_7}L_5L_B:\QQ]=[L_{p_7}:\QQ][L_5:\QQ][L_B:\QQ]
\)
并得到三重直积伽罗瓦群同构。
\par
最后，每个 \(S_5\) 的交换化同构于 \(\ZZ/2\ZZ\)（符号表示），三重直积的交换化为 \((\ZZ/2\ZZ)^3\)，对应的极大阿贝尔子扩张由三条二次分辨子域生成，即为所示三二次合成域。证毕。
\end{proof}

\begin{theorem}[两个五次域的合成次数与原始元]\label{thm:xi-p7-q5-primitive-element-degree25}
\leanverified{paper\_xi\_p7\_q5\_primitive\_element\_degree25}
取 \(p_7\) 的一根 \(\alpha\) 与 \(Q_5\) 的一根 \(\beta\)，并令
\[
K_1:=\QQ(\alpha),\qquad K_2:=\QQ(\beta).
\]
则
\[
[K_1K_2:\QQ]=25,
\]
并且 \(\theta:=\alpha+\beta\) 为 \(K_1K_2\) 的原始元，因而其极小多项式次数为 \(25\)。
\end{theorem}
\begin{proof}
由结论 \ref{con:xi-p7-q5-linearly-disjoint-s5xs5} 得 \(L_{p_7}\cap L_5=\QQ\)，从而其任意子域交亦为 \(\QQ\)，特别地 \(K_1\cap K_2=\QQ\)。
因此
\(
[K_1K_2:\QQ]=[K_1:\QQ][K_2:\QQ]=25
\)。
\par
为证 \(\theta\) 为原始元，取 \(\alpha\) 的共轭为 \(\{\alpha_i\}_{i=1}^5\)，取 \(\beta\) 的共轭为 \(\{\beta_j\}_{j=1}^5\)。
若存在
\(
\alpha_i+\beta_j=\alpha_{i'}+\beta_{j'}
\)，则
\[
\alpha_i-\alpha_{i'}=\beta_{j'}-\beta_j\in K_1\cap K_2=\QQ.
\]
另一方面，\(\Gal(L_{p_7}/\QQ)\cong S_5\)（命题 \ref{prop:pom-a4-galois-s5}）在五根上 \(2\)-重传递。
若对某 \(i\neq i'\) 有 \(\alpha_i-\alpha_{i'}\in\QQ\)，记 \(\alpha_i-\alpha_{i'}=q\in\QQ\)。
由于 \(q\) 在 \(\Gal(L_{p_7}/\QQ)\) 作用下不变，而 \(2\)-重传递意味着对任意有序对 \((k,\ell)\)（\(k\neq \ell\)）存在自同构把 \((i,i')\) 送至 \((k,\ell)\)，从而
\[
\alpha_k-\alpha_\ell=g(\alpha_i-\alpha_{i'})=q.
\]
取 \((k,\ell)=(i',i)\) 得 \(\alpha_{i'}-\alpha_i=q\)，与 \(\alpha_{i'}-\alpha_i=-(\alpha_i-\alpha_{i'})=-q\) 矛盾。
故必有 \(i=i'\)。
同理由 \(\Gal(L_5/\QQ)\cong S_5\)（定理 \ref{thm:xi-q5-galois-s5}）得到 \(j=j'\)。
于是 \(\{\alpha_i+\beta_j\}_{i,j}\) 两两不同，恰有 \(25\) 个共轭，从而 \(\deg m_\theta=25\)，并生成 \(K_1K_2\)。证毕。
\end{proof}

\begin{theorem}[三个五次域的合成次数与原始元]\label{thm:xi-p7-q5-b5-primitive-element-degree125}
\leanverified{paper\_xi\_p7\_q5\_b5\_primitive\_element\_degree125}
取 \(B_5\) 的一根 \(\gamma\) 并令 \(K_3:=\QQ(\gamma)\)，再置 \(\Theta:=\alpha+\beta+\gamma\)。
则
\[
[K_1K_2K_3:\QQ]=125,
\]
并且 \(\Theta\) 为 \(K_1K_2K_3\) 的原始元，因而其极小多项式次数为 \(125\)。
\end{theorem}
\begin{proof}
由结论 \ref{con:xi-p7-q5-b5-triple-product-s5cube-abelianization} 得 \((L_{p_7}L_5)\cap L_B=\QQ\)，从而 \((K_1K_2)\cap K_3=\QQ\)。
因此
\(
[K_1K_2K_3:\QQ]=[K_1K_2:\QQ][K_3:\QQ]=25\cdot 5=125
\)。
\par
设 \(\gamma\) 的共轭为 \(\{\gamma_k\}_{k=1}^5\)。
若
\(
\alpha_i+\beta_j+\gamma_k=\alpha_{i'}+\beta_{j'}+\gamma_{k'}
\)，则
\[
(\alpha_i+\beta_j)-(\alpha_{i'}+\beta_{j'})=-(\gamma_k-\gamma_{k'})\in K_3.
\]
左边属于 \(K_1K_2\)，由 \((K_1K_2)\cap K_3=\QQ\) 得两边皆为有理数。
记
\(
(\alpha_i+\beta_j)-(\alpha_{i'}+\beta_{j'})=q\in\QQ
\)，则 \(\gamma_k-\gamma_{k'}=-q\in\QQ\)。
由 \(\Gal(L_B/\QQ)\cong S_5\)（定理 \ref{thm:xi-terminal-zm-delta-a5-b5-galois-s5}）在五根上的 \(2\)-重传递性，若 \(k\neq k'\) 则同样推出所有有序对差 \(\gamma_u-\gamma_v\)（\(u\neq v\)）恒等于同一有理数，从而与 \(\gamma_{k'}-\gamma_k=-(\gamma_k-\gamma_{k'})\) 矛盾。
故 \(k=k'\)，从而 \(q=0\) 且 \(\alpha_i+\beta_j=\alpha_{i'}+\beta_{j'}\)。
由定理 \ref{thm:xi-p7-q5-primitive-element-degree25} 的论证可得 \((i,j)=(i',j')\)，从而 \(k=k'\)。
因此 \(\{\alpha_i+\beta_j+\gamma_k\}_{i,j,k}\) 两两不同，恰有 \(125\) 个共轭，从而 \(\deg m_\Theta=125\) 且 \(\Theta\) 生成 \(K_1K_2K_3\)。证毕。
\end{proof}

\begin{remark}[可复核的消元证书：次数 \(25\) 与 \(125\) 的整系数极小多项式]\label{rem:xi-p7-q5-b5-primitive-element-minpoly-certificates}
定理 \ref{thm:xi-p7-q5-primitive-element-degree25} 与定理 \ref{thm:xi-p7-q5-b5-primitive-element-degree125} 的原始元极小多项式可由整系数消元（resultant）得到。
相应的整系数多项式证书由脚本 \texttt{scripts/exp\_pom\_xi\_s5\_primitive\_element\_minpoly\_certificate.py} 生成并写入
\texttt{artifacts/export/pom\_xi\_s5\_primitive\_element\_minpoly\_certificate.txt}。
\end{remark}

\begin{conclusion}[\texorpdfstring{$S_3\times S_5\times S_5$}{S3xS5xS5} 的三重直积与指定分解型的正密度存在]\label{con:xi-delta-b5-triple-product-chebotarev-density}
令 \(L_3\) 为 \(P_{\mathrm{LY}}(y)\) 的分裂域，\(L_5\) 为 \(Q_5(y)\) 的分裂域，\(L_B\) 为 \(B_5(\delta)\) 的分裂域。
则三者两两线性互不相交，且
\[
\Gal(L_3L_5L_B/\QQ)\cong S_3\times S_5\times S_5.
\]
并且存在一素数集合 \(\mathcal P\) 具有自然密度
\[
\delta(\mathcal P)=\frac{1}{3}\cdot\frac{1}{5}\cdot\frac{1}{12}=\frac{1}{180},
\]
使得对任意 \(p\in\mathcal P\) 同时满足：
\begin{enumerate}
  \item \(P_{\mathrm{LY}}\bmod p\) 在 \(\FF_p[y]\) 上不可约；
  \item \(Q_5\bmod p\) 在 \(\FF_p[y]\) 上不可约；
  \item \(B_5\bmod p\) 的分解型为 \((2)(1)(1)(1)\)。
\end{enumerate}
更一般地，对任意给定的共轭类三元组 \((C_3,C_5,C_5')\subset S_3\times S_5\times S_5\)，对应“模 \(p\) 分解型同时落在该三元组”的素数集合具有自然密度
\[
\frac{|C_3|}{|S_3|}\cdot\frac{|C_5|}{|S_5|}\cdot\frac{|C_5'|}{|S_5|}.
\]
\end{conclusion}
\begin{proof}
由定理 \ref{thm:xi-leyang-ply-q5-linearly-disjoint-chebotarev} 得 \(L_3\cap L_5=\QQ\)，由结论 \ref{con:xi-delta-b5-q5-linearly-disjoint-s5xs5} 得 \(L_B\cap L_5=\QQ\)。
又由结论 \ref{con:xi-terminal-37-squarefree-coupling} 与结论 \ref{con:xi-terminal-zm-delta-b5-discriminant-ramification-quadratic} 可分别读出 \(L_3\)、\(L_B\) 的唯一二次子域为 \(\QQ(\sqrt{-111})\)、\(\QQ(\sqrt{-140267})\)，二者不等；
而 \(\Gal(L_3/\QQ)\cong S_3\) 的真正规子扩张仅可能为该二次子域，故 \(L_3\cap L_B=\QQ\)。
从而三者两两线性互不相交，直积伽罗瓦群同构成立。
\par
对不分歧素数，Chebotarev 定理在直积群上给出共轭类等分布。
取 \(S_3\) 中 \(3\)-循环类大小为 \(2\)（比例 \(1/3\)），取 \(S_5\) 中 \(5\)-循环类大小为 \(24\)（比例 \(1/5\)），取换位类大小为 \(10\)（比例 \(1/12\)），即得密度 \(1/180\)。
一般三元组的密度公式同理由共轭类大小乘积给出。证毕。
\end{proof}

\endinput


\begin{theorem}[\texorpdfstring{$A_7$}{A7} 的满对称伽罗瓦群与极点不可根式性]\label{thm:xi-a7-galois-s7}
令 \(A_7(y)\) 如定理 \ref{thm:xi-fdelta-subresultant-unique-double-root}。
则 \(A_7\) 在 \(\QQ\) 上不可约，且其分裂域 \(L_7/\QQ\) 满足
\[
\operatorname{Gal}(L_7/\QQ)\cong S_7.
\]
因此由 \(\delta_\ast\) 的分母诱导的极点参数集合在算术上达到最大对称性破缺，并不可由根式表示。
\end{theorem}
\begin{proof}
首先，\(A_7\bmod 23\) 在 \(\FF_{23}[y]\) 上不可约（分解型为 \([7]\)），且 \(23\nmid \mathrm{Disc}(A_7)\)，
故 Frobenius 元给出一个 \(7\)-循环。

其次，判别式素因子分解为
\[
\mathrm{Disc}(A_7)=2^{38}\cdot 3^{19}\cdot 31^{7}\cdot 5651\cdot 12045409\cdot 258006961,
\]
从而 \(\nu_{5651}(\mathrm{Disc}(A_7))=1\)。
并且 \(A_7\bmod 5651\) 亦出现唯一二重根型分解：
存在 \(a\neq b\in\FF_{5651}\) 与不可约三次因子 \(h_3\in\FF_{5651}[y]\) 使
\[
A_7(y)\equiv (y-a)^2(y-b)\,h_3(y)\quad(\bmod 5651).
\]
故 \(5651\) 处的惯性群包含换位。
在 \(S_7\) 中，一个 \(7\)-循环与一个换位生成 \(S_7\)，故伽罗瓦群为 \(S_7\)。
可复现核验脚本：\texttt{scripts/exp\_xi\_leyang\_q5\_a7\_subresultant\_galois\_audit.py}。证毕。
\end{proof}

\begin{theorem}[次级退化七次层对三/五次层的完全独立性与直积伽罗瓦结构]\label{thm:xi-leyang-a7-disjoint-from-ply-q5}
令 \(L_3:=\mathrm{Spl}(P_{\mathrm{LY}})\)、\(L_5:=\mathrm{Spl}(Q_5)\)、\(L_7:=\mathrm{Spl}(A_7)\)。
则
\[
L_7\cap (L_3L_5)=\QQ,
\qquad
\Gal(L_3L_5L_7/\QQ)\cong \Gal(L_3L_5/\QQ)\times \Gal(L_7/\QQ)\cong (S_3\times S_5)\times S_7.
\]
\end{theorem}
\begin{proof}
由定理 \ref{thm:xi-leyang-ply-q5-linearly-disjoint-chebotarev} 已知 \(L_3\cap L_5=\QQ\) 且 \(\Gal(L_3L_5/\QQ)\cong S_3\times S_5\)。
由定理 \ref{thm:xi-a7-galois-s7}，扩张 \(L_7/\QQ\) 为 \(S_7\)-伽罗瓦扩张。
因此交域 \(L_7\cap (L_3L_5)\) 亦为 \(\QQ\) 上伽罗瓦扩张。
\par
由于 \(S_7\) 的非平凡真正规子群仅有 \(A_7\)，\(L_7\) 的非平凡真正规子扩张仅有其二次分辨子域
\(
K_7=\QQ(\sqrt{\Disc(A_7)})
\)。
由定理 \ref{thm:xi-a7-galois-s7} 的判别式分解，
\[
\Disc(A_7)=2^{38}\cdot 3^{19}\cdot 31^{7}\cdot 5651\cdot 12045409\cdot 258006961,
\]
故 \(K_7\) 在素数 \(5651,12045409,258006961\) 处必分歧。
另一方面，合成域 \(L_3L_5\) 的分歧素数集合包含于
\(\mathrm{Supp}(\Disc(P_{\mathrm{LY}})\Disc(Q_5))\)，
其素支撑由定理 \ref{thm:xi-leyang-ply-q5-linearly-disjoint-chebotarev} 的判别式数据可见不含上述三素数。
因此 \(K_7\not\subseteq L_3L_5\)，从而交域只能为 \(\QQ\)。
直积伽罗瓦群同构随之得到。证毕。
\end{proof}

\begin{theorem}[节点双根与次级退化双根的同调消元恒等式：六次商多项式]\label{thm:xi-leyang-b6-from-node-vs-double-root}
\leanverified{paper\_xi\_leyang\_b6\_from\_node\_vs\_double\_root}
令
\[
P_4(y):=8192y^4+5888y^3-3680y^2-2848y+152,
\qquad
N(y):=y(y-1)(8y-3)(16y+11)(32y+19)\,P_{\mathrm{LY}}(y).
\]
则存在唯一六次多项式
\[
B_6(y)=360448y^6+267264y^5-221824y^4+18600y^3+149481y^2+25423y+1520\in\ZZ[y]
\]
使得在 \(\ZZ[y]\) 中有严格恒等式
\[
93\,N(y)+4\,A_7(y)\,P_4(y)=Q_5(y)\,B_6(y).
\]
\end{theorem}
\begin{proof}
由定理 \ref{thm:xi-fdelta-subresultant-unique-double-root}，在 \(\Disc_T(F_\delta(\cdot;y))=0\) 的情形下，\(F_\delta(T;y)\) 的唯一二重根满足
\[
\delta_\ast(y)=-\frac{N(y)}{4A_7(y)}.
\]
另一方面，节点消元给出（参见定理 \ref{thm:xi-terminal-zm-delta-node-preimage-elimination-r10} 的证明要点）：
在结点条件 \(Q_5(y)=0\) 上，节点对应的 \(\delta\)-坐标满足
\[
93\,\delta=P_4(y).
\]
在 \(Q_5(y)=0\) 的每个根处，判别式退化的双根与节点 \(\delta\)-坐标一致，
从而 \(\delta_\ast(y)=P_4(y)/93\) 在 \(Q_5(y)=0\) 上成立。
于是
\[
93\,\delta_\ast(y)-P_4(y)\equiv 0\quad(\bmod\ Q_5(y)).
\]
代入 \(\delta_\ast(y)=-N(y)/(4A_7(y))\) 并清分母得到
\[
93\,N(y)+4A_7(y)P_4(y)\equiv 0\quad(\bmod\ Q_5(y)).
\]
因此 \(Q_5\) 在 \(\QQ[y]\) 中整除左端多项式。
由于 \(Q_5\) 在 \(\ZZ[y]\) 中为本原多项式，Gauss 引理推出商多项式属于 \(\ZZ[y]\)；
再由次数比较（\(\deg(93N+4A_7P_4)=11\)、\(\deg Q_5=5\)）知商为唯一的六次式。
对该商作显式多项式除法即得所列 \(B_6\)。可复现核验脚本：\texttt{scripts/exp\_xi\_leyang\_b6\_s6\_discriminant\_audit.py}。证毕。
\end{proof}

\begin{theorem}[六次商多项式的满对称伽罗瓦群与敏感素数 \(59\)]\label{thm:xi-leyang-b6-galois-s6}
令 \(L_6:=\mathrm{Spl}(B_6)\)。
则：
\begin{enumerate}
  \item \(B_6\) 在 \(\QQ\) 上不可约；
  \item \(\Gal(L_6/\QQ)\cong S_6\)；
  \item 判别式满足
  \[
  \Disc(B_6)=2^{41}\cdot 3^{9}\cdot 31^{2}\cdot 59\cdot p_\ast,
  \qquad
  p_\ast=16381178872545177926239240820863,
  \]
  因而二次分辨子域 \(\QQ(\sqrt{\Disc(B_6)})\) 在素数 \(59\) 处必分歧。
\end{enumerate}
\end{theorem}
\begin{proof}
取素数 \(97\)。可检验 \(B_6\bmod 97\) 在 \(\FF_{97}[y]\) 上保持六次不可约，故 \(B_6\) 在 \(\QQ[y]\) 上不可约。
\par
记 \(G:=\Gal(L_6/\QQ)\)。
由不可约性，\(G\subseteq S_6\) 为传递子群，且存在某个不分歧素数处 Frobenius 元为 \(6\)-循环（由模素不可约分解型给出）。
由判别式分解可知 \(\Disc(B_6)\notin \QQ^{\times 2}\)，从而 \(G\nsubseteq A_6\)。
进一步地，模 \(59\) 分解出现唯一二重根型分解：
\[
B_6(y)\equiv (y-23)^2\cdot h_2(y)\cdot h_2'(y)\quad(\bmod 59),
\]
其中 \(h_2,h_2'\) 为不可约二次因子。
因此在 \(59\) 处的惯性群含有换位。
在 \(S_6\) 中，\(6\)-循环与换位生成 \(S_6\)，故 \(G=S_6\)。
判别式恒等式为直接计算并分解得到。
可复现核验脚本：\texttt{scripts/exp\_xi\_leyang\_b6\_s6\_discriminant\_audit.py}。证毕。
\end{proof}

\begin{theorem}[四层非可解对称群的直积化与联合分解型的素数密度]\label{thm:xi-leyang-s3s5s6s7-product-and-densities}
令 \(L_3,L_5,L_7\) 如上，\(L_6\) 如定理 \ref{thm:xi-leyang-b6-galois-s6}。
则
\[
L_6\cap (L_3L_5L_7)=\QQ,
\qquad
\Gal(L_3L_5L_6L_7/\QQ)\cong S_3\times S_5\times S_6\times S_7.
\]
并且除去有限坏素数后，对素数 \(p\) 的自然密度成立以下精确公式：
\begin{enumerate}
  \item \(P_{\mathrm{LY}},Q_5,B_6,A_7\) 在 \(\FF_p\) 上同时存在一次因子的素数密度为
  \[
  \frac{2}{3}\cdot\frac{19}{30}\cdot\frac{91}{144}\cdot\frac{531}{840}
  =\frac{14573}{86400};
  \]
  \item 四者在 \(\FF_p\) 上同时完全分解的素数密度为
  \[
  \frac{1}{|S_3||S_5||S_6||S_7|}=\frac{1}{2612736000}.
  \]
\end{enumerate}
\end{theorem}
\begin{proof}
由定理 \ref{thm:xi-leyang-a7-disjoint-from-ply-q5} 已知 \(L_7\cap(L_3L_5)=\QQ\)，且 \(\Gal(L_3L_5L_7/\QQ)\cong S_3\times S_5\times S_7\)。
由定理 \ref{thm:xi-leyang-b6-galois-s6}，扩张 \(L_6/\QQ\) 为 \(S_6\)-伽罗瓦扩张，其非平凡真正规子扩张仅有二次分辨子域
\(
K_6=\QQ(\sqrt{\Disc(B_6)})
\)。
由判别式分解，\(K_6\) 在素数 \(59\) 处必分歧。
\par
另一方面，\(\Gal(L_3L_5L_7/\QQ)\cong S_3\times S_5\times S_7\) 的一切二次子域由各因子的符号特征张成；
因此任何二次子域的分歧素数集合包含于
\(\mathrm{Supp}(\Disc(P_{\mathrm{LY}})\Disc(Q_5)\Disc(A_7))\)。
由定理 \ref{thm:xi-leyang-ply-q5-linearly-disjoint-chebotarev} 与定理 \ref{thm:xi-a7-galois-s7} 的判别式数据可见，上述支撑不含 \(59\)，故 \(K_6\not\subseteq L_3L_5L_7\)，从而 \(L_6\cap (L_3L_5L_7)=\QQ\)。
直积伽罗瓦结构随之得到。
\par
对不分歧素数，Chebotarev 定理给出 Frobenius 共轭类在直积群上的乘法性分布。
事件“在 \(\FF_p\) 上有一次因子”等价于对应对称群因子中的置换有不动点，
其在 \(S_3,S_5,S_6,S_7\) 上的比例分别为 \(2/3,19/30,91/144,531/840\)，相乘即得第一式；
“完全分解”对应 Frobenius 为恒等元，比例为 \(1/|S_n|\)，相乘得到第二式。证毕。
\end{proof}

\begin{theorem}[Perron 三阶 Fuchs 算子的视在奇点五次层与二级退化五次层的独立性]\label{thm:xi-leyang-qapp-s5-disjoint-from-q5}
令
\[
Q_{\mathrm{app}}(y)=168y^5+1260y^4-422y^3-183y^2-639y-76\in\ZZ[y],
\]
并记 \(L_{\mathrm{app}}:=\mathrm{Spl}(Q_{\mathrm{app}})\)、\(L_5:=\mathrm{Spl}(Q_5)\)。
则：
\begin{enumerate}
  \item \(\Gal(L_{\mathrm{app}}/\QQ)\cong S_5\)；
  \item \(L_{\mathrm{app}}\cap L_5=\QQ\)，从而 \(\Gal(L_{\mathrm{app}}L_5/\QQ)\cong S_5\times S_5\)；
  \item \(Q_{\mathrm{app}}\) 与 \(Q_5\) 在 \(\FF_p\) 上同时具有一次因子的素数密度为
  \[
  \left(\frac{19}{30}\right)^2=\frac{361}{900}.
  \]
\end{enumerate}
\end{theorem}
\begin{proof}
取不分歧素数 \(11\)，可检验 \(Q_{\mathrm{app}}\bmod 11\) 在 \(\FF_{11}[y]\) 上不可约，故 Frobenius 含 \(5\)-循环。
另一方面，判别式含素数 \(109\) 的奇次幂因子，且模 \(109\) 分解出现单个二重根型因子，从而惯性群含换位。
传递子群中由 \(5\)-循环与换位生成者必为 \(S_5\)，故 \(\Gal(L_{\mathrm{app}}/\QQ)\cong S_5\)。
\par
两者均为 \(S_5\)-伽罗瓦扩张，若交域非平凡，则二者的二次分辨子域必须相同。
但 \(\QQ(\sqrt{\Disc(Q_{\mathrm{app}})})\) 的分歧素支撑包含 \(109\)，而 \(L_5\) 的唯一二次子域为 \(\QQ(\sqrt{\Disc(Q_5)})\)，其分歧支撑由 \(\Disc(Q_5)\) 的素因子集合给出并不含 \(109\)。
因此二次分辨子域不相等，交域只能为 \(\QQ\)，从而直积伽罗瓦群同构成立。
密度结论由线性互不相交下 Chebotarev 乘法性给出，且 \(S_5\) 上“有不动点”的比例为 \(19/30\)。可复现核验脚本：\texttt{scripts/exp\_xi\_leyang\_qapp\_s5\_disjointness\_audit.py}。证毕。
\end{proof}

\begin{corollary}[由 Lee--Yang 算术分层诱导的非可解 Artin 表示族：受控分歧与 Frobenius 独立性]\label{cor:xi-leyang-artin-family-s3s5s6s7}
在定理 \ref{thm:xi-leyang-s3s5s6s7-product-and-densities} 的设定下，置 \(K:=L_3L_5L_6L_7\)。
取各因子 \(S_n\) 的标准表示 \(\mathrm{Std}_n\)（维数 \(n-1\)），并作直和
\[
\rho:=\mathrm{Std}_3\oplus \mathrm{Std}_5\oplus \mathrm{Std}_6\oplus \mathrm{Std}_7.
\]
则 \(\rho\) 给出一个维数 \(17\) 的 Artin 表示，其分歧素数集合包含于
\[
\mathrm{Supp}\bigl(\Disc(P_{\mathrm{LY}}),\Disc(Q_5),\Disc(B_6),\Disc(A_7)\bigr),
\]
并且对几乎所有素数 \(p\)，\(\rho(\mathrm{Frob}_p)\) 在四个对称群因子上的投影在 Chebotarev 意义下相互独立。
\end{corollary}
\begin{proof}
由定理 \ref{thm:xi-leyang-s3s5s6s7-product-and-densities}，\(\Gal(K/\QQ)\cong S_3\times S_5\times S_6\times S_7\)，取各因子投影并与标准表示复合即得 \(\rho\)。
分歧支撑包含于各层判别式素支撑的并集为判别式基本性质的直接推论；
独立性由直积伽罗瓦群下 Chebotarev 分布的乘法性给出。证毕。
\end{proof}

\endinput

\subsubsection{Toeplitz--scan--Hankel 接口的复杂度分离、采样阈值与账本维数闭环}\label{subsubsec:xi-time-protocol-conclusions-part45}

\begin{conclusion}[Toeplitz--PSD 层级与共动扫描的严格复杂度分离律]\label{con:xi-toeplitz-psd-vs-comoving-scan-strict-separation}
设高度窗长度为 \(T\)，固定 \(\delta_0>0\) 与图表尺度 \(t>0\)。
记
\[
N_{\mathrm{Toep}}(T,\delta_0)
\]
为固定图表的 Toeplitz--PSD 阶数提升证书的最小资源阶，
\[
N_{\mathrm{scan}}(T,t,\delta_0)
\]
为允许共动地址的扫描覆盖证书的最小资源阶。
若采用扫描覆盖最优标度
\[
N_{\mathrm{scan}}(T,t,\delta_0)\asymp \frac{T}{t+\delta_0},
\]
则当 \(T\to\infty\) 时存在绝对常数 \(c>0\) 使
\[
N_{\mathrm{Toep}}(T,\delta_0)\ \ge\
c\,\frac{(t+\delta_0)^2}{\delta_0}\cdot
\frac{N_{\mathrm{scan}}(T,t,\delta_0)^2}
{\log\!\bigl(1+(t+\delta_0)^2N_{\mathrm{scan}}(T,t,\delta_0)^2\bigr)}.
\]
特别地，
\[
N_{\mathrm{Toep}}\ \gtrsim\ \frac{N_{\mathrm{scan}}^2}{\log N_{\mathrm{scan}}},
\]
即两类可审计证书在高度窗意义下出现严格多项式指数分离：共动扫描把主阶由近二次压缩到一次。
\end{conclusion}
\begin{proof}
由结论 \ref{con:xi-terminal-horizon-certificate-quadratic-linear-lb}，
\[
N_{\mathrm{Toep}}(T,\delta_0)\gtrsim \frac{T^2}{\delta_0}\log(1+T^2)\ge \frac{T^2}{\delta_0}.
\]
由 \(N_{\mathrm{scan}}(T,t,\delta_0)\asymp T/(t+\delta_0)\)，存在常数 \(c_1>0\) 使
\[
T\ge c_1\,(t+\delta_0)\,N_{\mathrm{scan}}(T,t,\delta_0),
\]
故
\[
N_{\mathrm{Toep}}(T,\delta_0)\ \ge\
c_2\,\frac{(t+\delta_0)^2}{\delta_0}\,N_{\mathrm{scan}}(T,t,\delta_0)^2.
\]
再利用 \(\log(1+x)\ge 1\)（\(x\ge 0\)）即得
\[
N_{\mathrm{Toep}}(T,\delta_0)\ \ge\
c_2\,\frac{(t+\delta_0)^2}{\delta_0}\cdot
\frac{N_{\mathrm{scan}}(T,t,\delta_0)^2}
{\log\!\bigl(1+(t+\delta_0)^2N_{\mathrm{scan}}(T,t,\delta_0)^2\bigr)}.
\]
证毕。
\end{proof}

\begin{conclusion}[严格化 counterterm 的不可观测性与零维账本必需性]\label{con:xi-counterterm-unobservability-zero-dimensional-ledger-necessity}
设 \(C\) 为 Carath\'eodory 接口，定义核
\[
K(w,w'):=\frac{C(w)+C(w')^\ast}{1-ww'^\ast}.
\]
对任意 \(\eta\in\RR\)，令严格化变换
\[
C^\sharp(w):=C(w)+i\eta.
\]
则成立：
\begin{enumerate}
  \item \textnormal{核不变性：}
  \[
  K^\sharp(w,w')\equiv K(w,w').
  \]
  因而所有由 \(K\) 生成的 Toeplitz 证书数据（任意阶惯性、负指标及由其构造的铅笔）均不携带 \(\eta\) 信息。
  \item \textnormal{谱读出不变性：}记
  \[
  \Lambda_N(C):=\Spec\!\bigl(\mathcal P_N(C)\bigr)\cap\mathbb D,
  \]
  其中 \(\mathcal P_N(C)\) 表示由 \(N\) 级 Toeplitz 数据构造的纤维谱读出铅笔。
  则
  \[
  \Lambda_N(C^\sharp)=\Lambda_N(C)\qquad(\forall N).
  \]
\end{enumerate}
因此 \(\eta\) 在该审计口径内是不可识别参数；若要恢复 \(\eta\)，必须引入外置寄存信息，该信息在 \(\Sigma^{\circ,\mathrm{pro}}\) 语义下属于连续圆维不增的零维核。
\end{conclusion}
\begin{proof}
由定义直接计算：
\[
C^\sharp(w)+C^\sharp(w')^\ast
=C(w)+i\eta+C(w')^\ast-i\eta
=C(w)+C(w')^\ast,
\]
故 \(K^\sharp\equiv K\)。
于是任意由 \(K\) 生成的 Toeplitz 截断与铅笔对象保持不变，\(\mathcal P_N(C^\sharp)=\mathcal P_N(C)\)，从而
\(
\Lambda_N(C^\sharp)=\Lambda_N(C)
\)。
结论成立。证毕。
\end{proof}

\begin{conclusion}[scan \texorpdfstring{$\to$}{->} Toeplitz 负证书的最短编译定理]\label{con:xi-scan-toeplitz-negative-certificate-shortest-compiler}
在有限缺陷口径下，设稳定缺陷阶为 \(\kappa\)，并由共动扫描指纹 \(\{u_n\}\) 定义 Hankel 矩阵
\[
\mathsf H_{\kappa-1}(u):=(u_{p+q})_{0\le p,q\le \kappa-1}.
\]
则对每个充分大 \(N\)（特别地 \(N\ge 2\kappa-2\)）存在可逆矩阵 \(R_N\) 与 \(P_N\succeq 0\)，满足
\[
R_N^\ast\,\mathbf T_N(C)\,R_N
=
\begin{pmatrix}
P_N & 0\\
0 & -\,\mathsf H_{\kappa-1}(u)^\ast\mathsf H_{\kappa-1}(u)
\end{pmatrix}.
\]
并且对任意 \(v\in\CC^\kappa\)，令
\[
x:=R_N\binom{0}{v}\in\CC^{N+1},
\]
则
\[
x^\ast\mathbf T_N(C)x
=-\,\|\mathsf H_{\kappa-1}(u)v\|_2^2.
\]
因此存在最短编译器：仅输入截断指纹
\[
(u_0,\dots,u_{2\kappa-2}),
\]
即可构造 Toeplitz 负向量见证 \(x\)。
且该阈值不可降低：一般位置下，不存在只依赖 \((u_0,\dots,u_{2\kappa-3})\) 的通用算法，能同时完成“判定缺陷阶 \(\kappa\)”与“输出保证负性的 Toeplitz 见证”。
\end{conclusion}
\begin{proof}
合同分解与负向量 Rayleigh 恒等式由推论 \ref{cor:xi-toeplitz-scan-hankel-congruence-decomposition} 与命题 \ref{prop:xi-toeplitz-negative-witness-from-scan-hankel} 直接给出。
又 \(\mathsf H_{\kappa-1}(u)\) 的全部条目只依赖 \(u_0,\dots,u_{2\kappa-2}\)，故可由该前缀直接构造见证。
\par
最短性由定理 \ref{thm:xi-toeplitz-negative-inertia-minimal-sampling} 给出：\(\det(\mathsf H_{\kappa-1})\) 对 \(u_{2\kappa-2}\) 为仿射函数，一般位置下其系数非零，故同一前缀 \((u_0,\dots,u_{2\kappa-3})\) 存在可逆/奇异两类延拓，导致缺陷判定与负证书存在性分岔。证毕。
\end{proof}

\begin{conclusion}[unitary-slice 交叉审计的采样下界]\label{con:xi-unitary-slice-cross-audit-sampling-lower-bound}
沿用命题 \ref{prop:xi-unitary-slice-spectral-crossing} 的记号，固定 \(L>1\)，令
\[
\theta\mapsto U_L(\theta)
\]
为对应的幺正算子族，满足
\[
\Xi_{\Delta,L}\!\left(\tfrac12+it\right)=0
\iff
1\in\Spec\!\bigl(U_L(\theta)\bigr),\qquad \theta=2t\log L.
\]
设在区间 \([\theta_0,\theta_1]\) 上有 Lipschitz 约束
\[
\|U_L(\theta)-U_L(\theta')\|_{\mathrm{op}}
\le
\mathsf M\,|\theta-\theta'|.
\]
定义谱距函数
\[
d(\theta):=\min_{\lambda\in\Spec(U_L(\theta))}|\lambda-1|.
\]
取均匀网格 \(\theta_j=\theta_0+jh\)。若
\[
\min_j d(\theta_j)>\mathsf M h,
\]
则 \([\theta_0,\theta_1]\) 内不存在谱穿越点，因而不存在 Hardy 截面零点。
其逆否命题给出采样下界：若区间内存在穿越（等价于存在零点），则任何步长为 \(h\) 的均匀扫描都必须满足
\[
\min_j d(\theta_j)\le \mathsf M h.
\]
于是若审计协议以阈值 \(\varepsilon\) 触发（需观测到 \(d(\theta_j)\le\varepsilon\)），则必须取
\[
h\le \frac{\varepsilon}{\mathsf M},
\qquad
N_\theta:=\left\lceil\frac{\theta_1-\theta_0}{h}\right\rceil
\ge
\left\lceil\frac{\mathsf M(\theta_1-\theta_0)}{\varepsilon}\right\rceil.
\]
\end{conclusion}
\begin{proof}
幺正算子为正规算子，故谱的 Hausdorff 距离满足
\[
\mathrm{dist}_H\!\bigl(\Spec(U),\Spec(V)\bigr)\le \|U-V\|_{\mathrm{op}}.
\]
对任意 \(\theta\in[\theta_j,\theta_{j+1}]\)，有
\[
d(\theta)
\ge
d(\theta_j)-\|U_L(\theta)-U_L(\theta_j)\|_{\mathrm{op}}
\ge
d(\theta_j)-\mathsf M h.
\]
若所有网格点都满足 \(d(\theta_j)>\mathsf M h\)，则 \(d(\theta)>0\) 对全区间成立，故 \(1\notin\Spec(U_L(\theta))\)，再由命题 \ref{prop:xi-unitary-slice-spectral-crossing} 即得无零点。
取逆否命题并代入触发阈值 \(\varepsilon\) 即得采样下界。证毕。
\end{proof}

\begin{conclusion}[深度谱不适定性与素数账本信息量下界]\label{con:xi-depth-illposedness-prime-ledger-information-lower-bound}
在有限缺陷口径下，记深度 Hankel 判别式
\[
\Delta_{M-1}:=\det(\mu_{p+q})_{0\le p,q\le M-1}.
\]
由 Vandermonde 分解（定理 \ref{thm:xi-dethankel-vandermonde-square-finite-blaschke}）可知：深度碰撞/逃逸对应 \(\Delta_{M-1}\to 0\)，且该不适定性由深度谱本身决定。
\par
另一方面，令 \(H_0\in M_d(\ZZ)\) 满秩（可取为上述 Hankel 主块的整数实现），并记
\[
H_p(k):=\sum_{i=1}^d\min\bigl(v_p(s_i),k\bigr),
\]
其中 \(s_i\) 为 \(H_0\) 的 Smith 不变量。
若同余方程 \(H_0B\equiv H_1'\pmod{p^k}\) 可解，定义解集
\[
S_k:=\{B\in M_d(\ZZ/p^k\ZZ):\ H_0B\equiv H_1'\pmod{p^k}\},
\]
则
\[
\#S_k=p^{dH_p(k)},
\qquad
\lim_{k\to\infty}H_p(k)=v_p(\det H_0).
\]
因此，任何声称“无额外账本即可在模 \(p^k\) 下唯一重建”的协议均不成立；
唯一化所需外置账本信息量（bit）满足不可突破下界
\[
I_{\mathrm{ledger}}(k)\ \ge\ \log_2\#S_k
=d\,H_p(k)\log_2 p,
\]
并在 \(k\to\infty\) 时逼近
\[
d\,v_p(\det H_0)\log_2 p.
\]
当深度谱趋于碰撞/逃逸而使相应 \(p\)-进赋值累积时，该下界沿含 \(p\)-因子方向刚性增大，形成“深度谱不适定性 \(\Rightarrow\) 素数账本信息量刚性”的定量桥梁。
\end{conclusion}
\begin{proof}
解数公式与极限式分别由定理 \ref{thm:xi-hankel-normalform-crt-adelic-multiplicity} 及定理 \ref{thm:xi-smith-entropy-inverts-vp-invariants} 直接给出。
若无外置信息，确定性重建在 \(\#S_k>1\) 时无法区分全部候选；按信息论计数，下界为 \(\log_2\#S_k\)。
极限下界由 \(H_p(k)\to v_p(\det H_0)\) 立得。
关于不适定性与 \(\Delta_{M-1}\) 的对应由定理 \ref{thm:xi-dethankel-vandermonde-square-finite-blaschke} 的 Vandermonde 平方因子给出。证毕。
\end{proof}

\begin{conclusion}[\texorpdfstring{$\Sigma^{\circ,\mathrm{pro}}$}{Sigma^{circ,pro}} 的度量刻画与有限维化闭环]\label{con:xi-sigma-circ-pro-metric-characterization-finite-dimensional-closure}
设拓扑有限生成紧阿贝尔群
\[
G\simeq \TT^{\,r}\times\widehat{\ZZ}^{\,d}\times F,
\]
在 \(\widehat{\ZZ}^{\,d}\) 上取 primorial 超度量 \(d_{\#}\)，在 \(\TT^r\) 上取欧氏度量 \(d_{\mathrm{Euc}}\)，并在 \(G\) 上取 max-积度量 \(d=\max\{d_{\mathrm{Euc}},d_{\#}\}\)。
则（定理 \ref{thm:conclusion-adelic-hausdorff-dimension-cdim-pcdim}）
\[
\dim_{\mathrm H}(G,d)=r+d.
\]
且 \(r=\cdim(G)\)，\(d=\operatorname{pcdim}(\widehat{\ZZ}^{\,d})\)。
因此可由纯度量数据恢复二元签名
\[
\Sigma^{\circ,\mathrm{pro}}(G)
=
\Bigl(\dim_{\mathrm{top}}G,\ \dim_{\mathrm H}(G,d)-\dim_{\mathrm{top}}G\Bigr),
\]
其中 \(\dim_{\mathrm{top}}G=r\)。
\par
进一步地，对任意拓扑有限生成交换 profinite 群 \(K\)，令
\(
r:=\operatorname{pcdim}(K)
\)。
则存在紧连通交换群 \(\Sigma_K\) 使
\[
0\longrightarrow K\longrightarrow \Sigma_K\longrightarrow \TT^{\,r}\longrightarrow 0
\]
正合，且该 \(r\) 为最小可能值（定理 \ref{thm:conclusion-prime-fiber-min-circle-internalization}）。
于是 \(\operatorname{pcdim}\) 可解释为“把零维账本内化为连通相位门所需的最小圆因子数”，从而给出无穷素数支撑信息在有限维相位体系内闭合的严格原理化表述。
\end{conclusion}
\begin{proof}
第一部分由定理 \ref{thm:conclusion-adelic-hausdorff-dimension-cdim-pcdim} 立即得到：Hausdorff 维数加法与 \((\cdim,\operatorname{pcdim})\) 识别同时成立，故 \(\Sigma^{\circ,\mathrm{pro}}\) 可由 \((\dim_{\mathrm{top}},\dim_{\mathrm H})\) 纯度量恢复。
第二部分直接应用定理 \ref{thm:conclusion-prime-fiber-min-circle-internalization} 的存在性与最小性结论。证毕。
\end{proof}

\endinput

\subsubsection{折叠纤维的 Bayes--InfoNCE 二阶曲率、可逆提升预算与一变量全息守恒}\label{subsubsec:xi-time-protocol-conclusions-part46}

\begin{theorem}[Bayes--InfoNCE 的二阶渐近展开与倒纤维和不变量]\label{thm:xi-fold-bayes-infonce-second-order-reciprocal-fiber}
固定 \(m\ge 1\)，令 \(\Omega_m=\{0,1\}^m\)，\(\Fold_m:\Omega_m\twoheadrightarrow X_m\)。
定义纤维多重度与推前分布
\[
d_m(x):=\bigl|\Fold_m^{-1}(x)\bigr|,
\qquad
w_m(x):=\frac{d_m(x)}{2^m}\quad(x\in X_m).
\]
对 \(K\ge 2\) 记 \(n:=K-1\)，并令
\[
L_{K,m}:=L_m^{K,\star}
\]
表示 Bayes 最优 InfoNCE 值（定理 \ref{thm:pom-infonce-partition-optimal-closed-form} 的 Fold 口径）。
则有二阶展开
\[
L_{K,m}
=\log n-H(w_m)
+\frac{|X_m|+1}{2n}
+\frac{\displaystyle \sum_{x\in X_m}\frac{2^m}{d_m(x)}-6|X_m|-1}{12\,n^2}
+O_m(n^{-3}),
\]
其中
\[
H(w_m):=-\sum_{x\in X_m}w_m(x)\log w_m(x).
\]
等价地，若
\[
\kappa_m(w_m):=\EE_{X\sim w_m}[\log d_m(X)],
\]
则
\[
H(w_m)=m\log 2-\kappa_m(w_m),
\]
故主项也可写为 \(\log n+\kappa_m(w_m)-m\log 2\)。
\end{theorem}
\begin{proof}
由定理 \ref{thm:pom-infonce-bayes-optimal-third-order-largeK} 在 \(g=\Fold_m\) 下，
\[
L_{K,m}
=\log K-H(w_m)
+\frac{|X_m|-1}{2n}
+\frac{\mathsf S_{-1}(w_m)-6|X_m|+5}{12n^2}
+O\!\left(\frac{\mathsf S_{-2}(w_m)}{n^3}\right),
\]
其中
\[
\mathsf S_{-1}(w_m)=\sum_{x\in X_m}\frac{1}{w_m(x)}
=\sum_{x\in X_m}\frac{2^m}{d_m(x)}.
\]
又 \(\log K=\log(n+1)=\log n+\frac1n-\frac{1}{2n^2}+O(n^{-3})\)，代回并并合同阶项，得到
\[
\frac1n\text{ 系数 }=\frac{|X_m|-1}{2}+1=\frac{|X_m|+1}{2},
\]
\[
\frac1{n^2}\text{ 系数 }=\frac{\mathsf S_{-1}(w_m)-6|X_m|+5}{12}-\frac12
=\frac{\mathsf S_{-1}(w_m)-6|X_m|-1}{12}.
\]
由于 \(m\) 固定且 \(w_m(x)\ge 2^{-m}\)，\(\mathsf S_{-2}(w_m)\) 为仅依赖 \(m\) 的有限常数，故余项为 \(O_m(n^{-3})\)。
最后一条恒等式由命题 \ref{prop:fold-logmultiplicity-pressure-derivative} 给出。证毕。
\end{proof}

\begin{corollary}[二阶曲率识别离散度指数与反向 KL 证书]\label{cor:xi-fold-infonce-second-order-curvature-dispersion}
在定理 \ref{thm:xi-fold-bayes-infonce-second-order-reciprocal-fiber} 的设定下，定义
\[
\mathfrak D_m:=\mathfrak D(w_m)
:=\frac{1}{|X_m|^2}\sum_{x\in X_m}\frac{1}{w_m(x)}
=\frac{1}{|X_m|^2}\sum_{x\in X_m}\frac{2^m}{d_m(x)}.
\]
则有极限恒等式
\[
\lim_{K\to\infty}\left(
\frac{12(K-1)^2}{|X_m|^2}
\left[
L_{K,m}-\log(K-1)+H(w_m)-\frac{|X_m|+1}{2(K-1)}
\right]
+\frac{6|X_m|+1}{|X_m|^2}
\right)=\mathfrak D_m.
\]
并且
\[
D_{\mathrm{KL}}(U_m\|w_m)\le \log \mathfrak D_m,
\qquad
U_m:=\mathrm{Unif}(X_m).
\]
\end{corollary}
\begin{proof}
由定理 \ref{thm:xi-fold-bayes-infonce-second-order-reciprocal-fiber} 的 \(n^{-2}\) 系数比较，且
\(
\sum_x 2^m/d_m(x)=|X_m|^2\mathfrak D_m
\)，可得极限式。
反向 KL 上界由命题 \ref{prop:pom-reverse-kl-bound-by-dispersion} 直接给出。证毕。
\end{proof}

\begin{theorem}[熵驱动的可逆提升长度下界]\label{thm:xi-fold-reversible-lift-length-lowerbound-kappa}
\leanverified{paper\_xi\_fold\_reversible\_lift\_length\_lowerbound\_kappa}
设 \(|\Sigma|=M\ge 2\)。
称
\[
\Psi_m:\Omega_m\to X_m\times\Sigma^{\le L_m}
\]
为 \(\Fold_m\) 的可逆提升，若 \(\Psi_m\) 单射且
\[
\mathrm{pr}_{X_m}\circ\Psi_m=\Fold_m.
\]
则必有
\[
L_m\ \ge\ \frac{\kappa_m(w_m)}{\log M}-1
=\frac{m\log 2-H(w_m)}{\log M}-1.
\]
进一步由 \(H(w_m)\le \log|X_m|\) 得
\[
L_m\ \ge\ \frac{m\log 2-\log|X_m|}{\log M}-1.
\]
当 \(w_m\) 非均匀时，第一式严格强于仅由平均纤维规模得到的粗下界。
\end{theorem}
\begin{proof}
固定 \(x\in X_m\)。
由可逆提升定义，\(\Fold_m^{-1}(x)\) 中的不同微态在第二坐标上必须可区分，故
\[
d_m(x)\le |\Sigma^{\le L_m}|=\sum_{\ell=0}^{L_m}M^\ell< M^{L_m+1}.
\]
取对数并对 \(X\sim w_m\) 取期望：
\[
\kappa_m(w_m)=\EE_{w_m}[\log d_m(X)]\le (L_m+1)\log M.
\]
整理即得
\(
L_m\ge \kappa_m(w_m)/\log M-1
\)。
再用 \(H(w_m)=m\log 2-\kappa_m(w_m)\)（命题 \ref{prop:fold-logmultiplicity-pressure-derivative}）得到等价形式。
由 \(H(w_m)\le \log|X_m|\) 立得第二条。
当 \(w_m\) 非均匀时 \(H(w_m)<\log|X_m|\)，故第一条严格强于粗化下界。证毕。
\end{proof}

\begin{proposition}[边界膨胀与信息缺口的同尺度性]\label{prop:xi-fold-boundary-info-gap-same-scale}
将 \(X_m\subseteq\Omega_m=\{0,1\}^m\) 视作 \(m\)-维超立方体图 \(Q_m\) 的顶点子集（边为 Hamming 距离 \(1\)）。
令边界边集
\[
\partial X_m:=\bigl\{\{u,v\}\in E(Q_m):\ u\in X_m,\ v\notin X_m\bigr\}.
\]
则成立边缘等周下界
\[
\boxed{\;
|\partial X_m|
\ \ge\
|X_m|\log_2\!\Bigl(\frac{2^m}{|X_m|}\Bigr)\; }.
\]
特别地，对黄金均值（Zeckendorf）稳定集合 \(|X_m|=F_{m+2}=\Theta(\varphi^m)\)，有
\[
\frac{|\partial X_m|}{|X_m|}
\ \ge\
m\log_2\!\Bigl(\frac{2}{\varphi}\Bigr)+O(1).
\]
并且定理 \ref{thm:xi-fold-reversible-lift-length-lowerbound-kappa} 的粗化下界
\(
L_m\ge (m\log2-\log|X_m|)/\log M-1
\)
与上式共享同一主导信息缺口 \(m\log(2/\varphi)\)，从而任何可逆提升都必要求线性级别的侧信息长度。
\end{proposition}
\begin{proof}
超立方体的标准边缘等周不等式（Harper 型）给出：对任意 \(S\subseteq\Omega_m\)，其边界满足
\[
|\partial S|\ \ge\ |S|\log_2\!\Bigl(\frac{2^m}{|S|}\Bigr).
\]
取 \(S=X_m\) 得第一式。
\par
当 \(|X_m|=F_{m+2}\) 时，由 \(F_{m+2}=\Theta(\varphi^m)\) 得
\[
\log_2\!\Bigl(\frac{2^m}{|X_m|}\Bigr)
=m\log_2\!\Bigl(\frac{2}{\varphi}\Bigr)+O(1),
\]
代回即得第二式。
最后一段由定理 \ref{thm:xi-fold-reversible-lift-length-lowerbound-kappa} 的第二条不等式直接代入 \(|X_m|=F_{m+2}\) 并取主阶比较即可。证毕。
\end{proof}

\begin{theorem}[一变量全息守恒的纤维独立集分解]\label{thm:xi-fold-univariate-holographic-conservation}
沿用定理 \ref{thm:pom-fiber-multivariate-holographic-conservation} 的记号。对每个 \(x\in X_m\)，记
\[
I_{\Gamma_x}(z):=\sum_{S\in\mathrm{Ind}(\Gamma_x)} z^{|S|},
\qquad
|x|:=\sum_{i=1}^m x_i.
\]
则对任意标量 \(z\) 成立严格恒等式
\[
\sum_{x\in X_m} z^{|x|}I_{\Gamma_x}(z)=(1+z)^m.
\]
\end{theorem}
\begin{proof}
在定理 \ref{thm:pom-fiber-multivariate-holographic-conservation} 中取
\(
\boldsymbol z=(z,\dots,z)
\)。
此时
\[
w_{x,v}(\boldsymbol z)=\frac{z\cdot z}{z}=z,
\]
从而
\[
\mathcal H_x(\boldsymbol z)
=\left(\prod_{i=1}^m z^{x_i}\right)\sum_{S\in\mathrm{Ind}(\Gamma_x)}z^{|S|}
=z^{|x|}I_{\Gamma_x}(z).
\]
代入
\(
\sum_x \mathcal H_x(\boldsymbol z)=\prod_i(1+z_i)
\)
即得结论。证毕。
\end{proof}

\begin{corollary}[在 \texorpdfstring{$z=-1$}{z=-1} 处的消去与二项卷积约束]\label{cor:xi-fold-zminusone-cancellation-binomial-convolution}
在定理 \ref{thm:xi-fold-univariate-holographic-conservation} 的设定下：
\begin{enumerate}
  \item 有
  \[
  \sum_{x\in X_m}(-1)^{|x|}I_{\Gamma_x}(-1)=0.
  \]
  \item 令
  \[
  N_x(r):=\bigl|\{S\in\mathrm{Ind}(\Gamma_x): |S|=r\}\bigr|
  \]
  并约定 \(N_x(r)=0\)（\(r<0\)），则对每个 \(0\le k\le m\) 有
  \[
  \sum_{x\in X_m}N_x(k-|x|)=\binom{m}{k}.
  \]
\end{enumerate}
\end{corollary}
\begin{proof}
第一式由定理 \ref{thm:xi-fold-univariate-holographic-conservation} 取 \(z=-1\) 得
\(
(1-1)^m=0
\)。
第二式把两端按 \(z\) 展开并比较 \(z^k\) 系数即可。证毕。
\end{proof}

\begin{proposition}[右逆空间的指数下界与二阶曲率可审计证书]\label{prop:xi-fold-right-inverse-count-dispersion-lowerbound}
记
\[
\mathcal S(\Fold_m):=\{s:X_m\to\Omega_m:\ \Fold_m\circ s=\mathrm{id}_{X_m}\}.
\]
则
\[
|\mathcal S(\Fold_m)|=\prod_{x\in X_m} d_m(x),
\qquad
\log |\mathcal S(\Fold_m)|=\sum_{x\in X_m}\log d_m(x).
\]
并且若 \(\mathfrak D_m\) 如推论 \ref{cor:xi-fold-infonce-second-order-curvature-dispersion}，则
\[
\log |\mathcal S(\Fold_m)|
\ \ge\
|X_m|\left(\log\frac{2^m}{|X_m|}-\log\mathfrak D_m\right),
\]
\[
|\mathcal S(\Fold_m)|
\ \ge\
\left(\frac{2^m}{|X_m|\,\mathfrak D_m}\right)^{|X_m|}.
\]
此外，由推论 \ref{cor:xi-fold-infonce-second-order-curvature-dispersion}，\(\mathfrak D_m\) 可由 \(L_{K,m}\) 的二阶曲率极限读出，因此上述下界可由 Bayes--InfoNCE 曲线进行外部审计。
\end{proposition}
\begin{proof}
计数闭式由定理 \ref{thm:pom-microcanonical-section-count}（取 \(F=\Fold_m\)）直接给出。
令 \(U_m=\mathrm{Unif}(X_m)\)，则
\[
D_{\mathrm{KL}}(U_m\|w_m)
=\frac1{|X_m|}\sum_{x\in X_m}\log\frac{1/|X_m|}{w_m(x)}
=m\log2-\log|X_m|-\frac{1}{|X_m|}\sum_x\log d_m(x),
\]
故
\[
\sum_{x\in X_m}\log d_m(x)
=|X_m|\left(\log\frac{2^m}{|X_m|}-D_{\mathrm{KL}}(U_m\|w_m)\right).
\]
再由命题 \ref{prop:pom-reverse-kl-bound-by-dispersion}，
\(
D_{\mathrm{KL}}(U_m\|w_m)\le \log\mathfrak D_m
\)，即得对数下界；指数化得到第二式。证毕。
\end{proof}

\endinput

\subsubsection{Schur--Poisson 二体完备、Toeplitz 簇极限与 Smith--Ward 审计闭环}\label{subsubsec:xi-time-protocol-conclusions-part47}

\begin{theorem}[Schur--Poisson 互能同一性：局部冗余的链式守恒分解]\label{thm:xi-schur-poisson-mutual-energy-ledger-identity}
\leanverified{paper\_xi\_schur\_poisson\_mutual\_energy\_ledger\_identity}
设 \(A=\{a_1,\dots,a_\kappa\}\subset\mathbb D\) 为有限点状缺陷多重集，\(\mathsf P=\mathsf P(A)\) 为对应 Pick--Poisson Gram 矩阵。
记
\[
p_m:=\mathsf P_{mm}=P_{a_m}(-1),\qquad
\rho_{jk}:=\rho(a_j,a_k).
\]
对固定点序 \(1,\dots,\kappa\)，令 Schur 补通量
\[
\pi_m:=\mathsf P_{mm}-\mathsf P_{m,1:m-1}\,(\mathsf P^{(m-1)})^{-1}\,\mathsf P_{1:m-1,m}\qquad(1\le m\le \kappa),
\]
则成立严格恒等式
\[
2\sum_{1\le j<k\le \kappa}\bigl(-\log\rho_{jk}\bigr)
=\sum_{m=1}^{\kappa}\log\frac{p_m}{\pi_m}.
\]
特别地，每一项 \(\log(p_m/\pi_m)\ge 0\)，总和与点序无关；并且存在某个 \(m_\star\) 使
\[
\log\frac{p_{m_\star}}{\pi_{m_\star}}
\ge
\frac{2}{\kappa}\sum_{1\le j<k\le \kappa}\bigl(-\log\rho_{jk}\bigr).
\]
\end{theorem}
\begin{proof}
由推论 \ref{cor:xi-pick-poisson-energy-coulomb-decomposition}，
\[
-\log\det\mathsf P
=\sum_{m=1}^{\kappa}(-\log p_m)
+2\sum_{j<k}(-\log\rho_{jk}).
\]
由结论 \ref{con:xi-terminal-pick-poisson-schur-ledger}，
\[
-\log\det\mathsf P=\sum_{m=1}^{\kappa}(-\log\pi_m),\qquad 0<\pi_m\le p_m.
\]
两式相减即得
\[
2\sum_{j<k}(-\log\rho_{jk})
=\sum_{m=1}^{\kappa}\bigl(-\log\pi_m+\log p_m\bigr)
=\sum_{m=1}^{\kappa}\log\frac{p_m}{\pi_m}.
\]
由 \(\pi_m\le p_m\) 得每项非负。
右端虽依赖点序分配到各步，但其总和等于左端的对称二体互能，故与点序无关。
最后由 \(\max_m x_m\ge \frac1\kappa\sum_m x_m\)（取 \(x_m=\log(p_m/\pi_m)\)）即得所示下界。证毕。
\end{proof}

\begin{theorem}[二体完备性：\texorpdfstring{$-\log\det\mathsf P$}{-log det P} 的三集合二阶 Möbius 恒等式]\label{thm:xi-pick-poisson-three-set-mobius-identity}
\leanverified{paper\_xi\_pick\_poisson\_three\_set\_mobius\_identity}
设 \(A,B,C\subset\mathbb D\) 为按重数计两两不交的有限点集。
定义
\[
\Phi(S):=-\log\det\mathsf P(S).
\]
则有严格恒等式
\[
\Phi(A\sqcup B\sqcup C)-\Phi(A\sqcup B)-\Phi(B\sqcup C)+\Phi(B)
=2\sum_{a\in A}\sum_{c\in C}\bigl(-\log\rho(a,c)\bigr).
\]
等价地，
\[
\det\mathsf P(A\sqcup B\sqcup C)\cdot \det\mathsf P(B)
=\det\mathsf P(A\sqcup B)\cdot \det\mathsf P(B\sqcup C)\cdot
\prod_{a\in A,c\in C}\rho(a,c)^2.
\]
\end{theorem}
\begin{proof}
由推论 \ref{cor:xi-pick-poisson-energy-coulomb-decomposition}，对任意有限 \(S\) 有
\[
\Phi(S)=\sum_{x\in S}(-\log p_x)+2\sum_{\{x,y\}\subset S}(-\log\rho(x,y)).
\]
将该展开代入
\[
\Phi(A\sqcup B\sqcup C)-\Phi(A\sqcup B)-\Phi(B\sqcup C)+\Phi(B)
\]
后，所有单点项、块内点对项与经由 \(B\) 的交叉项逐项抵消，仅剩 \(A\)--\(C\) 交叉点对，遂得第一式。
对第一式取指数并整理即得第二式。证毕。
\end{proof}

\begin{corollary}[任意分块的严格簇分解：并集体积势的有限 Mayer 展开]\label{cor:xi-pick-poisson-finite-mayer-cluster-factorization}
\leanverified{paper\_xi\_pick\_poisson\_finite\_mayer\_cluster\_factorization}
设 \(A=\bigsqcup_{r=1}^{R}A_r\) 为按重数计的不交并。
则
\[
\det\mathsf P(A)
=\left(\prod_{r=1}^{R}\det\mathsf P(A_r)\right)
\prod_{1\le r<s\le R}\ \prod_{a\in A_r,\ b\in A_s}\rho(a,b)^2,
\]
等价地，
\[
\Phi(A)
=\sum_{r=1}^{R}\Phi(A_r)
+2\sum_{1\le r<s\le R}\ \sum_{a\in A_r,\ b\in A_s}(-\log\rho(a,b)).
\]
因此 \(\Phi\) 的并集复合严格止于二体项，不存在三体及更高阶簇修正。
\end{corollary}
\begin{proof}
由命题 \ref{prop:xi-pick-poisson-det-factorization}，
\[
\det\mathsf P(A)
=\left(\prod_{a\in A}p_a\right)
\left(\prod_{\{x,y\}\subset A}\rho(x,y)^2\right).
\]
将点对乘积按“块内对 + 跨块对”分解，即得所示乘积闭式；取负对数得到加法形式。
全部修正项仅由点对贡献，故高阶簇项不存在。证毕。
\end{proof}

\begin{theorem}[Toeplitz 负谱正则化行列式的几何簇极限律]\label{thm:xi-toeplitz-negative-regularized-det-geometric-cluster-limit}
在稳定显影区间与半经典极限 \(N\to\infty\)、\(N(1-|a_j|)\to\infty\) 下，定义
\[
\mathrm{Det}_{-}\!\bigl(\mathbf T_N(A)\bigr)
:=\prod_{j=1}^{\kappa}\bigl(-\lambda_j^{-}(\mathbf T_N(A))\bigr).
\]
设 \(A=\bigsqcup_{r=1}^{R}A_r\)，并对每个 \(A_r\) 取同一构造口径下的稳定 Toeplitz 证书序列 \(\mathbf T_N^{(r)}(A_r)\)。
则成立极限簇分解
\[
\lim_{N\to\infty}
\frac{\mathrm{Det}_{-}\!\bigl(\mathbf T_N(A)\bigr)}
{\prod_{r=1}^{R}\mathrm{Det}_{-}\!\bigl(\mathbf T_N^{(r)}(A_r)\bigr)}
=
\prod_{1\le r<s\le R}\ \prod_{a\in A_r,\ b\in A_s}\rho(a,b)^2.
\]
\end{theorem}
\begin{proof}
由推论 \ref{cor:xi-toeplitz-negative-regularized-determinant-limit}，
\[
\lim_{N\to\infty}\mathrm{Det}_{-}\!\bigl(\mathbf T_N(A)\bigr)=\det\mathsf P(A),
\qquad
\lim_{N\to\infty}\mathrm{Det}_{-}\!\bigl(\mathbf T_N^{(r)}(A_r)\bigr)=\det\mathsf P(A_r).
\]
故极限比值等于 \(\det\mathsf P(A)/\prod_r\det\mathsf P(A_r)\)。
再由推论 \ref{cor:xi-pick-poisson-finite-mayer-cluster-factorization} 即得右端跨块伪双曲 Vandermonde 因子。证毕。
\end{proof}

\begin{proposition}[阿贝尔多通道否证的单通道局域化与候选极点并集定理]\label{prop:xi-abelian-single-channel-localization-and-pencil-union}
\leanverified{paper\_xi\_abelian\_single\_channel\_localization\_and\_pencil\_union}
设 \(H\) 为有限阿贝尔群，\(\mathcal A:=\mathbb C[H]\)。
对解析 \(F:\mathbb D\to\mathcal A\)，记角色通道分量为 \(\{F_\chi\}_{\chi\in\widehat H}\)。
则对每个 \(N\ge 0\)，
\[
\mathbf T_N^{\mathcal A}(F)\sim_{\mathrm U}\bigoplus_{\chi\in\widehat H}\mathbf T_N(F_\chi),
\qquad
\nu_-\!\bigl(\mathbf T_N^{\mathcal A}(F)\bigr)
=\sum_{\chi\in\widehat H}\nu_-\!\bigl(\mathbf T_N(F_\chi)\bigr).
\]
特别地，若 \(\nu_-(\mathbf T_N^{\mathcal A}(F))>0\)，则必存在某个 \(\chi\) 使 \(\mathbf T_N(F_\chi)\nsucceq 0\)，并且可取否证向量在通道直和分解下完全支撑于该单一通道。
\par
进一步，定义通道候选集合
\[
\Lambda_{N,\chi}
:=\Spec\!\bigl(\mathbf T_N(F_\chi),\mathbf T_N^{(1)}(F_\chi)\bigr)\cap\mathbb D,
\]
并记 \(\mathbf T_N^{\mathcal A,(1)}(F)\) 为对应的一阶移位块 Toeplitz 矩阵。
以及块候选集合
\[
\Lambda_N^{\mathrm{block}}
:=\Spec\!\bigl(\mathbf T_N^{\mathcal A}(F),\mathbf T_N^{\mathcal A,(1)}(F)\bigr)\cap\mathbb D.
\]
则按重数计有
\[
\Lambda_N^{\mathrm{block}}=\bigcup_{\chi\in\widehat H}\Lambda_{N,\chi}.
\]
由结论 \ref{con:xi-finite-falsifier-offline-fiber-readout} 的 Cayley 读出规则，离线零点纤维多重集亦按同一并集律分解。
\end{proposition}
\begin{proof}
第一部分由推论 \ref{cor:xi-abelian-block-toeplitz-directsum} 与定理 \ref{thm:xi-abelian-channel-kappaCar-additivity} 直接得到；
单通道局域化即推论 \ref{cor:xi-abelian-channel-negativity-localization}。
\par
对候选极点，移位块 Toeplitz 同样在 Fourier--Pontryagin 分解下按通道直和对角化，故
\[
\bigl(\mathbf T_N^{\mathcal A}(F),\mathbf T_N^{\mathcal A,(1)}(F)\bigr)
\sim_{\mathrm U}
\bigoplus_{\chi\in\widehat H}
\bigl(\mathbf T_N(F_\chi),\mathbf T_N^{(1)}(F_\chi)\bigr).
\]
广义谱在有限直和下按重数并合，遂得
\(\Lambda_N^{\mathrm{block}}=\bigcup_\chi \Lambda_{N,\chi}\)。
最后把结论 \ref{con:xi-finite-falsifier-offline-fiber-readout} 的候选纤维构造逐通道应用并并合即可。证毕。
\end{proof}

\begin{theorem}[多通道有限型的稳定阈值：全局稳定阶等于通道稳定阶的最大值]\label{thm:xi-abelian-global-stability-threshold-max-channel}
\leanverified{paper\_xi\_abelian\_global\_stability\_threshold\_max\_channel}
在命题 \ref{prop:xi-abelian-single-channel-localization-and-pencil-union} 的设定下，进一步假设每个通道 \(F_\chi\) 处于有限型口径（盘内极点有限且边界无极点）。
记
\[
\kappa_\chi:=\mathrm{sq}_-(F_\chi),
\]
并定义通道稳定阶
\[
N_0(\chi):=\min\Bigl\{N:\ \nu_-\!\bigl(\mathbf T_n(F_\chi)\bigr)=\kappa_\chi,\ \forall n\ge N\Bigr\}.
\]
则块 Toeplitz 的全局稳定阶
\[
N_0^{\mathrm{block}}
:=\min\Bigl\{N:\ \nu_-\!\bigl(\mathbf T_n^{\mathcal A}(F)\bigr)=\sum_{\chi}\kappa_\chi,\ \forall n\ge N\Bigr\}
\]
满足
\[
N_0^{\mathrm{block}}=\max_{\chi\in\widehat H}N_0(\chi),
\]
并且对所有 \(N\ge N_0^{\mathrm{block}}\) 有
\[
\nu_-\!\bigl(\mathbf T_N^{\mathcal A}(F)\bigr)=\sum_{\chi\in\widehat H}\kappa_\chi.
\]
\end{theorem}
\begin{proof}
由定理 \ref{thm:xi-toeplitz-inertia-stabilizes-to-kappa}，每个通道存在稳定阶 \(N_0(\chi)\)，且
\(\nu_-(\mathbf T_N(F_\chi))=\kappa_\chi\)（\(N\ge N_0(\chi)\)）。
令 \(N_\star:=\max_\chi N_0(\chi)\)，则当 \(N\ge N_\star\) 时所有通道同时稳定。
由定理 \ref{thm:xi-abelian-channel-kappaCar-additivity}，
\[
\nu_-\!\bigl(\mathbf T_N^{\mathcal A}(F)\bigr)
=\sum_{\chi}\nu_-\!\bigl(\mathbf T_N(F_\chi)\bigr)
=\sum_{\chi}\kappa_\chi,
\]
故 \(N_0^{\mathrm{block}}\le N_\star\)。
\par
反向地，若 \(N<N_\star\)，取 \(\chi_0\) 使 \(N<N_0(\chi_0)\)。
由 \(N_0(\chi_0)\) 的最小性，存在 \(n\ge N\) 使
\(\nu_-(\mathbf T_n(F_{\chi_0}))\neq \kappa_{\chi_0}\)；
而同一定理给出 \(\nu_-(\mathbf T_n(F_{\chi_0}))\le \kappa_{\chi_0}\)，故该值严格小于 \(\kappa_{\chi_0}\)。
再用逐通道可加性得
\[
\nu_-\!\bigl(\mathbf T_n^{\mathcal A}(F)\bigr)
=\sum_{\chi}\nu_-\!\bigl(\mathbf T_n(F_\chi)\bigr)
<\sum_{\chi}\kappa_\chi.
\]
因此任何 \(N<N_\star\) 均不能成为全局稳定阶，故 \(N_0^{\mathrm{block}}\ge N_\star\)。
综上 \(N_0^{\mathrm{block}}=N_\star\)。证毕。
\end{proof}

\begin{theorem}[Hankel--Smith 模审计的 primorial 最优性：最小 Gödel 成本的素数实现]\label{thm:xi-hankel-smith-mod-audit-primorial-optimality}
设 \(H_0,H_1\in M_d(\mathbb Z)\) 且 \(\det(H_0)\neq 0\)，令
\[
A:=H_0^{-1}H_1\in M_d(\mathbb Q).
\]
由定理 \ref{thm:pom-hankel-smith-minimal-denominator}，存在最小整化因子 \(N_\ast\) 使 \(N_\ast A\in M_d(\mathbb Z)\)，并满足 \(N_\ast\mid \det(H_0)\)。
给定候选 \(\widehat A\in M_d(\mathbb Q)\) 满足 \(N_\ast\widehat A\in M_d(\mathbb Z)\)，定义
\[
F:=N_\ast(H_0\widehat A-H_1)\in M_d(\mathbb Z),\qquad M:=2\|F\|_\infty.
\]
若取有限素数集 \(\mathcal P\) 使
\[
\gcd\!\Bigl(\prod_{p\in\mathcal P}p,\ N_\ast\Bigr)=1,\qquad
F\equiv 0\pmod p\ \ (\forall p\in\mathcal P),\qquad
\prod_{p\in\mathcal P}p>M,
\]
则必有 \(\widehat A=A\)。
\par
进一步，定义对数成本
\[
\mathcal C(\mathcal P):=\sum_{p\in\mathcal P}\log p=\log\!\Bigl(\prod_{p\in\mathcal P}p\Bigr).
\]
在无互素约束的情形，令 \(n\) 为满足 \(p_n^\#>M\) 的最小整数（\(p_n^\#:=\prod_{i=1}^{n}p_i\)），则
\[
\min_{\mathcal P:\ \prod_{p\in\mathcal P}p>M}\ \mathcal C(\mathcal P)=\log(p_n^\#).
\]
在互素约束 \(\gcd(\prod_{p\in\mathcal P}p,N_\ast)=1\) 下，极小集合可由“删去与 \(N_\ast\) 冲突的有限素数后，按从小到大取素数”的贪心方案实现。
\end{theorem}
\begin{proof}
第一部分正是定理 \ref{thm:pom-hankel-deterministic-mod-audit-threshold} 在 \(n=d\) 情形的直接应用。
\par
对最优性，先看无约束情形。
取任意由 \(n\) 个互异素数组成的集合 \(\mathcal P\)，按升序写为 \(q_1<\cdots<q_n\)。
由于 \(q_i\ge p_i\)（第 \(i\) 小素数的最小性），有
\[
\prod_{p\in\mathcal P}p=\prod_{i=1}^{n}q_i\ge \prod_{i=1}^{n}p_i=p_n^\#,
\qquad
\mathcal C(\mathcal P)\ge \log(p_n^\#).
\]
故要使乘积超过 \(M\)，所需最小规模由 \(p_n^\#>M\) 的最小 \(n\) 决定，且该 \(n\) 下最小成本即 \(\log(p_n^\#)\)。
\par
在互素约束下，只需把可选素数集限制为
\(
\mathbb P_{N_\ast}:=\{p\in\mathbb P:\ p\nmid N_\ast\}
\)。
按升序记其元素 \(\ell_1<\ell_2<\cdots\)。
与上同理，任何可行集合 \(\mathcal P\subset\mathbb P_{N_\ast}\) 的升序元素逐项不小于 \(\ell_i\)，故在给定元素个数下，\(\prod_{p\in\mathcal P}p\) 与 \(\mathcal C(\mathcal P)\) 同时由前缀 \(\{\ell_1,\dots,\ell_n\}\) 最小化。
因此约束问题的极小解由该递增前缀贪心得到，这恰是“去除有限冲突素数后的素数前缀”。证毕。
\end{proof}

\begin{proposition}[\texorpdfstring{$p$}{p}-进分支熵的饱和律：\texorpdfstring{$\mathrm{Stiff}_p$}{Stiffp} 作为模解空间增长率]\label{prop:xi-padic-branch-entropy-saturation-law}
\leanverified{paper\_xi\_padic\_branch\_entropy\_saturation\_law}
设 \(H_0\in M_d(\mathbb Z)\) 满秩，Smith 正规形不变量为
\[
s_1\mid s_2\mid\cdots\mid s_d.
\]
对素数 \(p\) 与 \(k\ge 1\)，定义 Smith 精度势
\[
H_p(k):=\sum_{i=1}^{d}\min\bigl(v_p(s_i),k\bigr).
\]
则有
\[
\#\ker\!\Bigl(H_0:(\mathbb Z/p^k\mathbb Z)^d\to(\mathbb Z/p^k\mathbb Z)^d\Bigr)=p^{H_p(k)}.
\]
并且对矩阵同余方程
\[
H_0X\equiv H_1\pmod{p^k},
\]
若可解，则解集为仿射空间且基数严格为
\[
\#\{X\bmod p^k:\ H_0X\equiv H_1\}=p^{dH_p(k)}.
\]
因此在可解层级上有增长极限
\[
\lim_{k\to\infty}\frac1d\log_p\#\{X\bmod p^k:\ H_0X\equiv H_1\}
=\lim_{k\to\infty}H_p(k)
=v_p(\det H_0)
=:\mathrm{Stiff}_p.
\]
\end{proposition}
\begin{proof}
核规模公式由命题 \ref{prop:xi-hankel-normalform-hp-discrete-concave-invert}（取 \(H_1=0\)）给出。
同余方程解数公式由命题 \ref{prop:xi-smith-pk-congruence-solvability-count} 给出。
最后 \(\min(v,k)\to v\)（\(k\to\infty\)）逐项收敛，故
\[
\lim_{k\to\infty}H_p(k)=\sum_{i=1}^{d}v_p(s_i)=v_p\!\Bigl(\prod_{i=1}^{d}s_i\Bigr)=v_p(\det H_0),
\]
遂得所示极限。证毕。
\end{proof}

\begin{theorem}[Toeplitz 层级的一致收敛可证伪性：固定阶谱裕度给出显式扰动半径]\label{thm:xi-toeplitz-fixed-order-margin-perturbation-radius}
令 \(0<r<1\)、\(N\ge 0\)。
设 \(C,\widetilde C\in\mathrm{Hol}(\{|w|<1\})\)，且
\[
\|C-\widetilde C\|_{L^\infty(\{|w|\le r\})}\le \varepsilon.
\]
记 \(\mathbf T_N(C),\mathbf T_N(\widetilde C)\) 为对应系数生成的 Toeplitz 主子矩阵（在本文 Carath\'eodory 口径下为 Hermitian）。
则存在显式常数
\[
L_N(r):=1+2\sum_{m=1}^{N}r^{-m}
\]
使
\[
\|\mathbf T_N(C)-\mathbf T_N(\widetilde C)\|_{\mathrm{op}}
\le L_N(r)\,\varepsilon.
\]
因而：
\begin{enumerate}
  \item 若 \(\mathbf T_N(C)\succeq 0\) 且 \(\lambda_{\min}(\mathbf T_N(C))\ge \delta>0\)，当 \(\varepsilon<\delta/L_N(r)\) 时必有 \(\mathbf T_N(\widetilde C)\succeq 0\)。
  \item 若 \(\lambda_{\min}(\mathbf T_N(C))\le -\delta<0\)，当 \(\varepsilon<\delta/L_N(r)\) 时必有 \(\lambda_{\min}(\mathbf T_N(\widetilde C))<0\)。
\end{enumerate}
故固定阶稳定负谱裕度给出可审计的有限否证半径，与猜想 \ref{conj:discussion-kernel-zeta-toeplitz-convergence} 的“有限证书可证伪接口”一致。
\end{theorem}
\begin{proof}
写
\[
C(w)=\sum_{n\ge 0}c_n w^n,\qquad
\widetilde C(w)=\sum_{n\ge 0}\widetilde c_n w^n.
\]
由 Cauchy 系数估计，
\[
|c_n-\widetilde c_n|
\le \varepsilon\,r^{-n}\qquad(n\ge 0).
\]
故对 \(\Delta_N:=\mathbf T_N(C)-\mathbf T_N(\widetilde C)\)，其偏移 \(|j-k|=m\) 处条目满足
\[
|(\Delta_N)_{jk}|\le \varepsilon\,r^{-m}\le \varepsilon\,r^{-N}.
\]
逐行求和得
\[
\|\Delta_N\|_{\infty}
\le \varepsilon\!\left(1+2\sum_{m=1}^{N}r^{-m}\right)
=L_N(r)\varepsilon.
\]
于是
\[
\|\Delta_N\|_{\mathrm{op}}\le \|\Delta_N\|_{\infty}\le L_N(r)\varepsilon.
\]
再由 Weyl 特征值扰动不等式，
\[
\lambda_{\min}\!\bigl(\mathbf T_N(\widetilde C)\bigr)
\ge \lambda_{\min}\!\bigl(\mathbf T_N(C)\bigr)-\|\Delta_N\|_{\mathrm{op}},
\]
\[
\lambda_{\min}\!\bigl(\mathbf T_N(\widetilde C)\bigr)
\le \lambda_{\min}\!\bigl(\mathbf T_N(C)\bigr)+\|\Delta_N\|_{\mathrm{op}}.
\]
代入上界 \(\|\Delta_N\|_{\mathrm{op}}\le L_N(r)\varepsilon\) 即得两条结论。证毕。
\end{proof}

\begin{proposition}[尺度中心化账本的 Cholesky 释义：\texorpdfstring{$\eta_m$}{eta_m} 为归一化条件体积的对数密度]\label{prop:xi-centered-ledger-cholesky-conditional-volume-density}
在统一伸缩流 \(\phi_m\) 下，Schur 补通量满足
\[
\pi_r(m)=m^{-1}\pi_r(1).
\]
定义尺度中心化增量
\[
\eta_r:=-\log\pi_r+V_\partial\qquad(1\le r\le\kappa),
\]
则 \(\eta_r(m)\equiv \eta_r(1)\)，且
\[
\sum_{r=1}^{\kappa}\eta_r=I_\partial
\]
为严格不变量。
\par
此外，若 \(\mathsf P\succ 0\) 且 \(\mathsf P=LL^\ast\) 为与点序一致的 Cholesky 分解，则
\[
L_{rr}^2=\pi_r\qquad(1\le r\le\kappa),
\]
从而
\[
\eta_r=-2\log L_{rr}+V_\partial.
\]
即 \(\eta_r\) 精确刻画第 \(r\) 步条件体积元相对边界势的归一化对数密度，并在统一伸缩下保持不变。
\end{proposition}
\begin{proof}
伸缩律 \(\pi_r(m)=m^{-1}\pi_r(1)\) 见结论 \ref{con:xi-terminal-pick-poisson-schur-ledger-homothetic-scaling}；
\(\eta_r\) 的尺度不变性与总和分解见结论 \ref{con:xi-terminal-pick-poisson-centered-ledger-invariants}。
对 Cholesky 部分，由定理 \ref{thm:xi-schur-ledger-subprobability-offdiagonal-gap}(3) 有 \(\pi_r=|L_{rr}|^2\)。
在正定情形可取 \(L_{rr}>0\)，故 \(L_{rr}^2=\pi_r\)，代入定义即得 \(\eta_r=-2\log L_{rr}+V_\partial\)。证毕。
\end{proof}

\begin{theorem}[Ward 型双线性递推的符号因子不可规约：\texorpdfstring{$(-1)^{\lfloor n/2\rfloor}$}{(-1)^{floor(n/2)} 为非平凡二阶扭结}]\label{thm:xi-ward-sign-factor-gauge-irreducible}
\leanverified{paper\_xi\_ward\_sign\_factor\_gauge\_irreducible}
沿用结论 \ref{con:xi-terminal-zm-elliptic-denominator-ward-bilinear-recurrence} 的分母序列 \((v_n)_{n\ge 0}\)：
\[
v_{n+2}v_{n-2}=29\,v_n^2+(-1)^{\lfloor n/2\rfloor}\,25\,v_{n+1}v_{n-1}\qquad(n\ge 2).
\]
不存在任意符号规约 \(\varepsilon_n\in\{\pm1\}\) 使
\[
u_n:=\varepsilon_n v_n
\]
满足无符号递推
\[
u_{n+2}u_{n-2}=29\,u_n^2+25\,u_{n+1}u_{n-1}\qquad(n\ge 2).
\]
换言之，因子 \((-1)^{\lfloor n/2\rfloor}\) 不能被逐项重定相消去。
\end{theorem}
\begin{proof}
反设存在 \((\varepsilon_n)\) 使上述无符号递推成立。
把 \(u_n=\varepsilon_n v_n\) 代入并与原递推逐项比较，得到对一切 \(n\ge 2\)
\[
\varepsilon_{n+2}\varepsilon_{n-2}=1,
\qquad
(-1)^{\lfloor n/2\rfloor}\varepsilon_{n+2}\varepsilon_{n-2}
=\varepsilon_{n+1}\varepsilon_{n-1}.
\]
第一式给出 \(\varepsilon_{n+2}=\varepsilon_{n-2}\)，故 \((\varepsilon_n)\) 在模 \(4\) 意义下周期。
令 \(e_j:=\varepsilon_j\)（\(j=0,1,2,3\)）。
由第二式：
\[
n\equiv 0\pmod 4\ \Longrightarrow\ e_1e_3=+1,
\]
\[
n\equiv 2\pmod 4\ \Longrightarrow\ e_1e_3=-1,
\]
矛盾。
故不存在此类规约。证毕。
\end{proof}

\begin{corollary}[标准高度的二阶差分读出：分母增长率的可审计极限公式]\label{cor:xi-elliptic-denominator-second-difference-height-readout}
\leanverified{paper\_xi\_elliptic\_denominator\_second\_difference\_height\_readout}
在结论 \ref{con:xi-terminal-zm-elliptic-denominator-growth-nt-height} 的设定下，分母增长满足
\[
\log v_n=\hat h(P)\,n^2+o(n^2).
\]
若进一步记
\[
r_n:=\log v_n-\hat h(P)\,n^2
\]
并满足离散二阶可差分条件
\[
r_{n+1}+r_{n-1}-2r_n\to 0\qquad(n\to\infty),
\]
则有
\[
\lim_{n\to\infty}\Bigl(\log v_{n+1}+\log v_{n-1}-2\log v_n\Bigr)=2\hat h(P),
\]
等价地，
\[
\lim_{n\to\infty}\frac{v_{n+1}v_{n-1}}{v_n^2}=e^{2\hat h(P)}.
\]
因此 \(\hat h(P)\) 可由分母序列的离散二阶差分直接读出；结合定理 \ref{thm:xi-ward-sign-factor-gauge-irreducible}，该读出对符号规约保持刚性。
\end{corollary}
\begin{proof}
由定义，
\[
\log v_n=\hat h(P)\,n^2+r_n.
\]
于是
\[
\log v_{n+1}+\log v_{n-1}-2\log v_n
=\hat h(P)\bigl((n+1)^2+(n-1)^2-2n^2\bigr)
+(r_{n+1}+r_{n-1}-2r_n).
\]
由假设 \(r_{n+1}+r_{n-1}-2r_n\to 0\)，得
\[
\log v_{n+1}+\log v_{n-1}-2\log v_n\to 2\hat h(P).
\]
取极限即得第一式。
第二式由指数化立即推出。证毕。
\end{proof}

\paragraph{Toeplitz 深度、尺度流与全息读出的补充刚性命题}

\begin{theorem}[缺陷熵与边界深度给出的 Toeplitz 证书阶数硬下界]\label{thm:xi-toeplitz-depth-hard-lower-bound-defect-entropy}
设 \(A_{\kappa_0}=\{a_1,\dots,a_{\kappa_0}\}\subset\mathbb D\)（\(\kappa_0\ge 1\)），
\[
\mathsf P:=\mathsf P(A_{\kappa_0})\succ0,\qquad
E(A_{\kappa_0}):=-\log\det\mathsf P,\qquad
h_{\min}:=\min_{1\le j\le \kappa_0}(1-|a_j|),
\]
并记条件数
\[
\mathrm{cond}(\mathsf P):=\frac{\lambda_{\max}(\mathsf P)}{\lambda_{\min}(\mathsf P)}.
\]
若 \(N\) 达到命题 \ref{prop:xi-toeplitz-threshold-via-anchored-interpolation} 的稳定门槛
\[
N\ \ge\ C\,\frac{\log \mathrm{cond}(\mathsf P)}{h_{\min}},
\]
则必有更强显式下界
\[
N\ \ge\ \frac{C}{h_{\min}}
\left(
\frac{1}{\kappa_0}E(A_{\kappa_0})
+\log\frac{\operatorname{tr}\mathsf P}{\kappa_0}
\right).
\]
并且
\[
\operatorname{tr}\mathsf P=\sum_{j=1}^{\kappa_0}P_{a_j}(-1),
\]
故右端完全由缺陷几何量给定。
\end{theorem}
\begin{proof}
设 \(\lambda_1\ge\cdots\ge\lambda_{\kappa_0}>0\) 为 \(\mathsf P\) 的特征值。
由均值不等式，
\[
\lambda_{\max}=\lambda_1\ge \frac{1}{\kappa_0}\sum_{j=1}^{\kappa_0}\lambda_j
=\frac{\operatorname{tr}\mathsf P}{\kappa_0},
\]
\[
\lambda_{\min}=\lambda_{\kappa_0}\le
\left(\prod_{j=1}^{\kappa_0}\lambda_j\right)^{1/\kappa_0}
=(\det\mathsf P)^{1/\kappa_0}.
\]
故
\[
\mathrm{cond}(\mathsf P)
=\frac{\lambda_{\max}}{\lambda_{\min}}
\ge
\frac{\operatorname{tr}\mathsf P/\kappa_0}{(\det\mathsf P)^{1/\kappa_0}}.
\]
取对数并用 \(E(A_{\kappa_0})=-\log\det\mathsf P\) 得
\[
\log\mathrm{cond}(\mathsf P)\ge
\log\frac{\operatorname{tr}\mathsf P}{\kappa_0}
+\frac{1}{\kappa_0}E(A_{\kappa_0}).
\]
再代回门槛不等式即得结论。
\(\operatorname{tr}\mathsf P=\sum_j P_{a_j}(-1)\) 来自对角项定义。证毕。
\end{proof}

\begin{theorem}[尺度流下有限内点缺陷不变集的不可能性]\label{thm:xi-cayley-scaling-no-finite-interior-invariant-set}
\leanverified{paper\_xi\_cayley\_scaling\_no\_finite\_interior\_invariant\_set}
令 \(m>1\)，并取 Cayley 共轭尺度流
\[
\phi_m:=\mathcal C\circ(z\mapsto mz)\circ\mathcal C^{-1}:\mathbb D\to\mathbb D
\]
（见命题 \ref{prop:xi-pick-poisson-det-scaling-law}）。
若有限集合 \(A\subset\mathbb D\) 满足 \(\phi_m(A)=A\)，则必有 \(A=\varnothing\)。
\end{theorem}
\begin{proof}
令 \(Z:=\mathcal C^{-1}(A)\subset\mathbb H\)。
由定义，
\[
\phi_m(A)=A
\iff
\mathcal C(mZ)=\mathcal C(Z)
\iff
mZ=Z.
\]
若 \(Z\neq\varnothing\)，取 \(z\in Z\)。
则对任意 \(n\ge 0\)，\(m^n z\in Z\)。
因 \(m>1\) 且 \(z\in\mathbb H\)（特别 \(z\neq 0\)），点列 \(\{m^n z\}_{n\ge 0}\) 两两不同，故 \(Z\) 无穷，与 \(A\) 有限矛盾。
故 \(Z=\varnothing\)，从而 \(A=\varnothing\)。证毕。
\end{proof}

\begin{theorem}[锚点 Gram 子集全息分布的 \texorpdfstring{$PSL(2,\RR)$}{PSL(2,R)} 严格不变性]\label{thm:xi-anchor-subset-holographic-distribution-psl-invariant}
设 \(A=\{x_1,\dots,x_n\}\subset\RR\)（\(n\le \kappa\)），极点集 \(\{p_k\}_{k=1}^{\kappa}\subset\CC_-\)，权重 \(w_k>0\)。
对任意 \(I\subset[\kappa]\) 且 \(|I|=n\)，定义
\[
C_{A,I}:=\left(\frac{1}{x_i-p_k}\right)_{1\le i\le n,\ k\in I}.
\]
由命题 \ref{prop:xi-basepoint-scan-anchor-det-cauchy-vandermonde} 的 Cauchy--Binet 展开可定义
\[
\mathbb P_A(I):=
\frac{\left(\prod_{k\in I}w_k\right)\,|\det C_{A,I}|^2}{\det G_A},
\]
其中 \(G_A\) 为锚点 Gram 矩阵。
则 \(\mathbb P_A\) 为 \(\{|I|=n\}\) 上概率分布。
\par
对任意 Möbius 变换
\(
m(x)=\frac{ax+b}{cx+d}\in PSL(2,\RR)
\),
令
\[
x_i':=m(x_i),\qquad p_k':=m(p_k),\qquad
w_k':=w_k\,\frac{(ad-bc)^2}{|cp_k+d|^2},
\]
并用 \((x_i',p_k',w_k')\) 构造 \(\mathbb P_{A'}^{\,\prime}\)。
则对所有 \(|I|=n\) 成立
\[
\mathbb P_{A'}^{\,\prime}(I)=\mathbb P_A(I).
\]
\end{theorem}
\begin{proof}
由恒等式
\[
\frac{1}{m(x)-m(p)}
=
\frac{(cx+d)(cp+d)}{ad-bc}\cdot\frac{1}{x-p},
\]
可得
\[
\det C_{A',I}^{\,\prime}
=
\left(\prod_{i=1}^{n}(cx_i+d)\right)
\left(\prod_{k\in I}(cp_k+d)\right)
(ad-bc)^{-n}\det C_{A,I}.
\]
故
\[
\left(\prod_{k\in I}w_k'\right)\,|\det C_{A',I}^{\,\prime}|^2
=
\left(\prod_{i=1}^{n}|cx_i+d|^2\right)
\left(\prod_{k\in I}w_k\right)\,|\det C_{A,I}|^2,
\]
其中 \(\prod_{k\in I}|cp_k+d|^2\) 与 \((ad-bc)^{2n}\) 由 \(w_k'\) 与 \(|\det C_{A',I}^{\,\prime}|^2\) 的反幂严格抵消。
另一方面，由命题 \ref{prop:xi-basepoint-scan-psl2r-covariance}，
\[
\det G_{A'}^{\,\prime}
=
\left(\prod_{i=1}^{n}|cx_i+d|^2\right)\det G_A.
\]
分子分母同乘同一因子，故 \(\mathbb P_{A'}^{\,\prime}(I)=\mathbb P_A(I)\)。证毕。
\end{proof}

\begin{theorem}[统一伸缩仅耗散 Poisson 自能而守恒互能]\label{thm:xi-scaling-mutual-energy-conservation-selfenergy-dissipation}
\leanverified{paper\_xi\_scaling\_mutual\_energy\_conservation\_selfenergy\_dissipation}
设 \(A,B\subset\mathbb D\) 为有限且按重数计不交的点集。
定义
\[
E(A):=-\log\det\mathsf P(A),\qquad
I(A;B):=2\sum_{a\in A,\ b\in B}\bigl(-\log\rho(a,b)\bigr).
\]
令 \(m>1\) 且 \(\phi_m\) 同定理 \ref{thm:xi-cayley-scaling-no-finite-interior-invariant-set}。
则有：
\begin{enumerate}
  \item 互能守恒：
  \[
  I(\phi_m(A);\phi_m(B))=I(A;B).
  \]
  \item 若 \(|A|=\kappa_0\)，则
  \[
  E(\phi_m(A))=E(A)+\kappa_0\log m.
  \]
  且该平移完全来自 Poisson 自能项；伪双曲互能项逐项不变。
\end{enumerate}
\end{theorem}
\begin{proof}
由命题 \ref{prop:xi-pick-poisson-det-scaling-law}，
\[
\rho(\phi_m(a),\phi_m(b))=\rho(a,b),\qquad
P_{\phi_m(a)}(-1)=m^{-1}P_a(-1),\qquad
\det\mathsf P(\phi_m(A))=m^{-|A|}\det\mathsf P(A).
\]
第一条直接给出互能守恒。
\par
由推论 \ref{cor:xi-pick-poisson-energy-coulomb-decomposition}，
\[
E(A)=\sum_{a\in A}-\log P_a(-1)
+2\sum_{\{a,b\}\subset A}-\log\rho(a,b).
\]
在 \(\phi_m\) 下，
\(
-\log P_{\phi_m(a)}(-1)=-\log P_a(-1)+\log m
\)，
而 \(\rho\)-项逐项不变，故总量增量为 \(|A|\log m\)。
等价地，也可由
\(
\det\mathsf P(\phi_m(A))=m^{-|A|}\det\mathsf P(A)
\)
直接得到 \(E(\phi_m(A))=E(A)+|A|\log m\)。证毕。
\end{proof}

\begin{theorem}[同一全纯矩列的实线确定性与复平面非确定性]\label{thm:xi-joukowsky-holomorphic-moments-real-determinate-complex-indeterminate}
对每个 \(r>1\)，定义
\[
J_r(e^{i\theta})=r e^{i\theta}+r^{-1}e^{-i\theta},\qquad
E_r:=J_r(\mathbb S^1),\qquad
\mu_r:=(J_r)_{\#}\mu_{\mathrm{Haar}}.
\]
则对任意 \(n\ge 0\)，全纯矩
\[
m_n:=\int_{\CC} w^n\,d\mu_r(w)
\]
与 \(r\) 无关，且
\[
m_{2k}=\binom{2k}{k},\qquad m_{2k+1}=0\qquad(k\ge 0).
\]
此外成立：
\begin{enumerate}
  \item 在 \(\RR\) 上，矩列 \(\{m_n\}\) 唯一对应测度
  \[
  d\nu(x)=\frac{1}{\pi\sqrt{4-x^2}}\mathbf 1_{[-2,2]}(x)\,dx.
  \]
  \item 在 \(\CC\) 上，\(\{\mu_r\}_{r>1}\) 构成不可数族两两互异且两两奇异的概率测度，并共享同一矩列 \(\{m_n\}\)。
\end{enumerate}
\end{theorem}
\begin{proof}
先算矩：
\[
\int_{\CC}w^n\,d\mu_r(w)
=\frac{1}{2\pi}\int_{0}^{2\pi}\left(r e^{i\theta}+r^{-1}e^{-i\theta}\right)^n\,d\theta
=\sum_{j=0}^{n}\binom{n}{j}r^{2j-n}\frac{1}{2\pi}\int_{0}^{2\pi}e^{i(2j-n)\theta}\,d\theta.
\]
仅当 \(2j-n=0\) 积分非零，故奇次为 \(0\)，偶次 \(n=2k\) 时取 \(j=k\) 得 \(m_{2k}=\binom{2k}{k}\)。
\par
对 (1)，取 \(\Theta\sim\mathrm{Unif}[0,2\pi]\) 与 \(X:=2\cos\Theta\)。
则 \(\EE[X^{2k+1}]=0\)、\(\EE[X^{2k}]=\binom{2k}{k}\)，且 \(X\) 分布即 \([-2,2]\) 上反正弦测度。
由于支撑紧，矩问题在 \(\RR\) 上确定。
\par
对 (2)，每个 \(E_r\) 为方程
\[
\frac{x^2}{(r+r^{-1})^2}+\frac{y^2}{(r-r^{-1})^2}=1
\]
给出的椭圆边界。
若 \(r>s>1\)，则 \(r+r^{-1}>s+s^{-1}\) 且 \(r-r^{-1}>s-s^{-1}\)。
故任意 \((x,y)\in E_s\) 满足
\[
\frac{x^2}{(r+r^{-1})^2}+\frac{y^2}{(r-r^{-1})^2}
<
\frac{x^2}{(s+s^{-1})^2}+\frac{y^2}{(s-s^{-1})^2}
=1,
\]
即 \((x,y)\notin E_r\)，从而 \(E_r\cap E_s=\varnothing\)。
又 \(\mu_r(E_r)=1\)，因此 \(r\neq s\) 时 \(\mu_r\perp\mu_s\)。
参数 \(r>1\) 不可数，故得到不可数个两两奇异测度，且它们共享已算出的同一矩列。证毕。
\end{proof}

\begin{theorem}[有限样本相等模式分布在 \texorpdfstring{$t=2n$}{t=2n} 的完备可识别性]\label{thm:xi-equality-pattern-2n-complete-identifiability}
设 \(X\) 为有限集，\(|X|=n\)，\(w\in\Delta(X)\)。
取
\[
X_1,\dots,X_t\stackrel{\mathrm{i.i.d.}}{\sim}w,
\]
并令 \(\Pi_t\) 为 \([t]\) 的随机划分：\(i,j\) 同块当且仅当 \(X_i=X_j\)。
则：
\begin{enumerate}
  \item 由 \(\Pi_{2n}\) 的全分布可唯一恢复权重多重集 \(\{w(x):x\in X\}\)，因此 \(w\) 在重标号意义下唯一确定。
  \item 若 \(n\) 未知，但 \(t\ge 2n\) 且处于非退化情形（\(\{w(x):x\in X\}\) 两两不同），则 \(n\) 亦由 \(\Pi_t\) 的分布唯一确定，进而可恢复 \(w\)。
\end{enumerate}
\end{theorem}
\begin{proof}
记幂和
\[
p_k(w):=\sum_{x\in X}w(x)^k\qquad(k\ge 1).
\]
对任意 \(k\le t\)，由 \(\Pi_t\) 的分布可得到其限制 \(\Pi_t|_{[k]}\) 的分布。
令 \(\mathbf 1_k:=\{\{1,\dots,k\}\}\)（单块划分），则
\[
p_k(w)=\mathbb P\!\left(\Pi_t|_{[k]}=\mathbf 1_k\right),
\]
故 \(\Pi_t\) 的全分布唯一决定 \(\{p_k(w)\}_{k\le t}\)（与划分格 Möbius 反演接口等价）。
\par
当 \(t=2n\) 时，已知 \(p_1,\dots,p_n\)（其中 \(p_1=1\)）。
令
\[
Q(z):=\prod_{x\in X}(z-w(x))
=z^n-e_1z^{n-1}+e_2z^{n-2}-\cdots+(-1)^n e_n.
\]
Newton 恒等式把 \((e_1,\dots,e_n)\) 与 \((p_1,\dots,p_n)\) 互相唯一确定，
故 \(Q\) 唯一，从而根多重集 \(\{w(x):x\in X\}\) 唯一。
这给出第 (1) 条。
\par
对第 (2) 条，在非退化假设下，原子测度
\(
\nu:=\sum_{x\in X}\delta_{w(x)}
\)
恰有 \(n\) 个互异原子，矩为 \(\int u^k\,d\nu(u)=p_k(w)\)。
故 Hankel 矩阵 \((p_{i+j}(w))_{i,j\ge 0}\) 的秩为 \(n\)。
当 \(t\ge 2n\) 时，已知矩足以读出该秩，因而可先恢复 \(n\)，再由第 (1) 条恢复权重多重集。证毕。
\end{proof}

\input{subsubsec__xi-time-protocol-conclusions-part47c-fold-bias-infonce-hadamard}
\input{subsubsec__xi-time-protocol-conclusions-part47ca-infonce-hausdorff-prony-capacity-renyi-shift}

\endinput


\endinput


